\documentclass[11pt,openany]{book}
\setcounter{secnumdepth}{3}
\usepackage[margin=1.1in]{geometry}
\usepackage[T1]{fontenc}
\usepackage[utf8]{inputenc}
\usepackage{lmodern}
\usepackage[french]{babel}
\addto\captionsfrench{\renewcommand{\tablename}{Tableau}}
\usepackage{amsmath,amssymb,amsthm}
\usepackage{graphicx}
\usepackage{tikz}
\usetikzlibrary{babel} % French babel makes : ; ! ? active
\usetikzlibrary{arrows.meta,decorations.markings,decorations.pathmorphing,decorations.pathreplacing,calc}
\input{figures/icons}
\usepackage[unicode,colorlinks=true,linkcolor=blue!60!black,citecolor=blue!60!black]{hyperref}
\usepackage{microtype}
\setlength{\emergencystretch}{3em} % long names (licences, references) in narrow lines
\usepackage{empheq}
\usepackage{array}
\usepackage{booktabs}
\usepackage{multirow}
\usepackage{longtable}
\usepackage{circuitikz}
\definecolor{keyyellow}{rgb}{1.0,0.97,0.78}
% key equations: \begin{empheq}[box=\keybox]{equation} ... \end{empheq} -- light-yellow box, number stays outside
\newcommand{\keybox}[1]{{\setlength{\fboxsep}{7pt}\colorbox{keyyellow}{#1}}}
\usepackage[most]{tcolorbox}
% key statements (propositions etc.): the same light-yellow background, breakable across pages
\newtcolorbox{keyblock}{enhanced,breakable,colback=keyyellow,colframe=keyyellow,boxrule=0pt,arc=0pt,left=6pt,right=6pt,top=2pt,bottom=2pt,boxsep=0pt}
\renewcommand{\chaptermark}[1]{\markboth{\small\chaptername~\thechapter.\ #1}{}}
\renewcommand{\sectionmark}[1]{\markright{\small\thesection.\ #1}}
\renewcommand{\topfraction}{0.95}
\renewcommand{\textfraction}{0.05}
\renewcommand{\floatpagefraction}{0.75}

% part pages: title at the top third, and text may follow on the same page (used for the part introductions)
\makeatletter
\renewcommand\part{\if@openright\cleardoublepage\else\clearpage\fi\thispagestyle{plain}\if@twocolumn\onecolumn\@tempswatrue\else\@tempswafalse\fi\null\vspace{0.18\textheight}\secdef\@part\@spart}
\renewcommand\@endpart{\par\vspace{3em}\if@tempswa\twocolumn\fi}
\makeatother
\newtheorem{theorem}{Théorème}[chapter]
\newtheorem{proposition}{Proposition}[chapter]
\newtheorem{lemma}{Lemme}[chapter]
\theoremstyle{remark}
\newtheorem{remark}{Remarque}[chapter]
% exercises: \begin{exercise}[\lvlB] ... \end{exercise}, numbered "Exercise 4.3"; difficulty marks
\theoremstyle{definition}
\newtheorem{exercise}{Exercice}[chapter]
\newcommand{\lvlA}{$\star$}
\newcommand{\lvlB}{$\star\star$}
\newcommand{\lvlC}{$\star\star\star$}
% answer to an exercise in the Answers appendix: \answer{ex:NN:k}{text}
\newcommand{\answer}[2]{\par\noindent\textbf{\ref{#1}.}~#2\par\smallskip}
\newcommand{\dd}{\mathrm{d}}
% neuron type code: first four letters of the primary mediator
\newcommand{\nt}[1]{\textsf{\textbf{#1}}}
\newcommand{\mV}{\,\text{mV}}
\newcommand{\ms}{\,\text{ms}}
\newcommand{\mM}{\,\text{mM}}
\newcommand{\um}{\,\text{\textmu m}}
\newcommand{\ion}[2]{\mathrm{#1}^{#2}}
\newcommand{\Na}{\ion{Na}{+}}
\newcommand{\K}{\ion{K}{+}}
\newcommand{\Cl}{\ion{Cl}{-}}
\newcommand{\Ca}{\ion{Ca}{2+}}
\newcommand{\conc}[1]{[\mathrm{#1}]}

\title{\Huge L'ordinateur de 20 watts\\[16pt] \Large Comment calcule le cerveau: une approche quantitative pour ingénieurs}
\hypersetup{pdftitle={L'ordinateur de 20 watts. Comment calcule le cerveau : une approche quantitative pour ingénieurs},pdfauthor={Konstantin Kladko}}
\author{\Large Konstantin~Kladko}
% cover: a salad of the book's figures around the title card (figures/bookcover.tex)
\newcommand{\covertitle}{L'ordinateur de 20 watts}
\newcommand{\coversubtitle}{Comment calcule le cerveau: une approche quantitative pour ingénieurs}
\newcommand{\coverauthor}{Konstantin~Kladko}
\date{}

\begin{document}
\frontmatter
\input{figures/bookcover}
\thispagestyle{empty}
\vspace*{\fill}
\noindent{\small \copyright~2026 Konstantin~Kladko.\\[4pt]
Le texte, les figures, les exercices et les réponses sont diffusés sous licence Creative Commons Attribution -- Pas d'utilisation commerciale -- Partage dans les mêmes conditions 4.0 International (CC BY-NC-SA 4.0): \url{https://creativecommons.org/licenses/by-nc-sa/4.0/deed.fr}. Le livre peut être librement copié, utilisé dans l'enseignement, traduit et adapté à des fins non commerciales, à condition de citer l'auteur et de diffuser les œuvres dérivées sous la même licence.\\[4pt]
Les programmes des modèles et de compilation sont diffusés sous licence MIT.\\[4pt]
Sources, modèles et nouvelles versions: \url{https://github.com/kladkogex/20-watt-computer}.}
\clearpage

\tableofcontents
\chapter*{Comment lire ce livre}
\addcontentsline{toc}{chapter}{Comment lire ce livre}

Il n'est pas nécessaire de lire ce livre dans l'ordre. Les chapitres~1 à~4 forment le socle commun: les types de neurones, ce que calculent \nt{GLUT} et \nt{GABA}, et la langue qu'ils parlent. Ensuite, le livre se ramifie. La carte de la figure~\ref{fig:roadmap} montre quels chapitres s'appuient sur quels autres: une flèche du chapitre $A$ vers le chapitre $B$ signifie que $B$ utilise des notions et des formules de $A$. La carte ne montre que les dépendances principales; les petits renvois en arrière ne gênent pas la lecture.

\begin{figure}[h]
\centering
\begin{tikzpicture}[>=Stealth,scale=0.95,
  ch/.style={draw,thick,rounded corners=3pt,align=center,font=\scriptsize,text width=1.85cm,minimum height=0.62cm,inner sep=2pt},
  base/.style={ch,fill=blue!8}, bus/.style={ch,fill=orange!14}, sum/.style={ch,fill=gray!18},
  io/.style={ch,fill=green!10}, sort/.style={ch,fill=cyan!8}, lab/.style={ch,fill=violet!10},
  dep/.style={->,thick,gray!60!black}]
% foundation
\node[base] (c1) at (0,0) {1 Types de neurones};
\node[base] (c2) at (2.3,0) {2 \nt{GLUT}};
\node[base] (c3) at (4.6,0) {3 \nt{GLUT} et \nt{GABA}};
\node[base] (c4) at (6.9,0) {4 Code Morse};
% broadcast channels and the synthesis
\node[bus] (c5) at (4.6,-1.5) {5 \nt{SERO}};
\node[bus] (c6) at (7.3,-1.5) {6 \nt{DOPA}};
\node[bus] (c7) at (9.2,-0.9) {7 Théories de la dopamine};
\node[bus] (c8) at (9.2,-2.1) {8 \nt{NORA}};
\node[bus] (c9) at (11.5,-2.1) {9 \nt{ACET}};
\node[sum] (c10) at (13.8,-0.9) {10 Toute la machine};
% inputs and outputs
\node[io] (c12) at (2.3,-4.1) {12 Recon-\\naissance};
\node[io] (c11) at (4.6,-4.1) {11 Entrées};
\node[io] (c13) at (6.9,-4.1) {13 Muscles};
\node[io] (c14) at (9.2,-3.05) {14 Douleur};
\node[io] (c15) at (9.2,-4.1) {15 Apprentissage moteur};
\node[io] (c16) at (9.2,-5.15) {16 Chimie\\du corps};
% lab, sorting, sleep
\node[lab] (c17) at (11.5,-4.1) {17 Débogage};
\node[lab] (c21) at (13.8,-4.1) {21 Machines vivantes};
\node[lab] (c22) at (13.8,-5.15) {22 DishBrain};
\node[sort] (c18) at (11.5,-5.15) {18 Neurones-\\instruments};
\node[sort] (c19) at (13.8,-6.2) {19 Nombre de dendrites};
\node[bus] (c20) at (11.5,-6.2) {20 Sommeil};
% dependencies
\draw[dep] (c1) -- (c2); \draw[dep] (c2) -- (c3); \draw[dep] (c3) -- (c4);
\draw[dep] (c1) -- (c5); \draw[dep] (c4) -- (c6); \draw[dep] (c5) -- (c6);
\draw[dep] (c6) -- (c7); \draw[dep] (c6) -- (c8); \draw[dep] (c8) -- (c9);
\draw[dep] (c7) -- (c10); \draw[dep] (c9) -- (c10);
\draw[dep] (c4) -- (c11); \draw[dep] (c11) -- (c12); \draw[dep] (c11) -- (c13); \draw[dep] (c5) -- (c13);
\draw[dep] (c13) -- (c14); \draw[dep] (c13) -- (c15); \draw[dep] (c13) -- (c16);
\draw[dep] (c15) -- (c17); \draw[dep] (c9) -- (c17); \draw[dep] (c17) -- (c21); \draw[dep] (c21) -- (c22);
\draw[dep] (c16) -- (c18); \draw[dep] (c18) -- (c19); \draw[dep] (c16) -- (c20);
% legend
\begin{scope}[yshift=-7.25cm]
\foreach \s/\t [count=\i] in {base/socle, bus/canaux et régimes, sum/synthèse, io/entrées et sorties, sort/classements des neurones, lab/laboratoire}{
  \pgfmathtruncatemacro{\col}{mod(\i-1,3)}\pgfmathtruncatemacro{\row}{(\i-1)/3}
  \node[\s,text width=0.3cm,minimum height=0.32cm] at ({\col*4.8+0.2},{-\row*0.55}) {};
  \node[anchor=west,font=\scriptsize] at ({\col*4.8+0.55},{-\row*0.55}) {\t};}
\end{scope}
\end{tikzpicture}
\caption{Carte des chapitres. Une flèche va d'un chapitre vers celui qui s'appuie sur lui. Les autres renvois entre chapitres sont des pointeurs, pas des dépendances.}
\label{fig:roadmap}
\end{figure}

\section*{Parcours de lecture}

\begin{center}
\footnotesize
\setlength{\tabcolsep}{4pt}
\renewcommand{\arraystretch}{1.2}
\begin{tabular}{>{\raggedright}p{3.0cm}>{\raggedright}p{3.9cm}>{\raggedright\arraybackslash}p{6.4cm}}
\toprule
Parcours & Pour qui & Chapitres dans l'ordre \\
\midrule
cours d'un trimestre & enseignant, 10 semaines & semaine 1: ch.~1--2; 2: ch.~3; 3: ch.~4; 4: ch.~5--6; 5: ch.~7--8; 6: ch.~9; 7: ch.~10; 8: ch.~11--12; 9: ch.~13 et~15; 10: ch.~17 \\
IA et apprentissage automatique & qui connaît les réseaux de neurones et veut comprendre comment le cerveau apprend & 1--4, 5 (survol), 6, 7, 8, 9, 11 (survol), 12, 13 (survol), 15 \\
électronique et matériel & concepteur de circuits, développeur de puces neuromorphiques & 1--4, 5, 11, 13, 16, 18, 19, 17 (chapitres 9 et 15 en survol) \\
neuro-interfaces et calculateurs vivants & qui construit des interfaces cerveau-ordinateur et des puces à neurones vivants & 1, 2, 4, 11, 13, 17 (chapitres 9 et 15 en survol), 21, 22 (chapitre 12 en survol) \\
aperçu rapide & ingénieur qui dispose d'un week-end & 1, 2, 3, 6, 10; les renvois du chapitre~6 aux chapitres~4--5 se comprennent par le contexte \\
\bottomrule
\end{tabular}
\end{center}

\og Survol\fg{} signifie qu'il suffit de lire le début du chapitre, ses énoncés dans les encadrés jaunes et son tableau récapitulatif.

Chaque chapitre se termine par des exercices: estimations, dérivations, modélisation et conception; les réponses brèves sont réunies dans l'annexe~\ref{sec:answers}. Toutes les notations se trouvent dans l'annexe~\ref{sec:notation}, et les expériences sur lesquelles reposent les principaux énoncés de chaque chapitre, dans l'annexe~\ref{sec:supplement}. Les nombres du livre proviennent de petits programmes du dossier \texttt{models/}; on peut les exécuter et les modifier.

\mainmatter

\chapter{Types de neurones}
\label{sec:neurontypes}

\section{Le neurone comme boîte noire}

Supposons qu'on vous remette un sac de quatre-vingt-six milliards de neurones~\cite{Azevedo2009} en vous demandant de les trier en tas. Comment vous y prendriez-vous?

Un ingénieur à qui l'on donne un sac de circuits intégrés inconnus ne les trierait pas selon la forme du boîtier. Il demanderait: combien la pièce a-t-elle d'entrées, combien de sorties, et que fournit-elle en sortie? Dessinons le neurone de la même façon, comme un bloc sur un schéma (fig.~\ref{fig:blackbox}).

\begin{figure}[t]
\centering
\begin{tikzpicture}[>=Stealth,x=1cm,y=1cm]
% the box
\node[draw,very thick,minimum width=2.6cm,minimum height=2.6cm,fill=blue!6] (N) at (0,0) {};
\node at (0,0.55) {\small neurone};
\node at (0,-0.15) {$\displaystyle u=\sum_i w_i x_i$};
\node at (0,-0.8) {\small $y=\Theta(u-\theta)$};
% inputs
\foreach \k/\y/\lab in {1/1.05/{x_1},2/0.55/{x_2},3/0.05/{x_3}}{
  \draw[->,thick,blue!60!black] (-4.2,\y) -- (-1.3,\y);
  \node[left] at (-4.2,\y) {$\lab$};
  \node[above,blue!60!black] at (-2.1,\y-0.06) {\scriptsize $w_{\k}$};}
\node at (-4.45,-0.4) {$\vdots$};
\node at (-2.1,-0.4) {$\vdots$};
\draw[->,thick,blue!60!black] (-4.2,-1.05) -- (-1.3,-1.05);
\node[left] at (-4.2,-1.05) {$x_n$};
\node[above,blue!60!black] at (-2.1,-1.11) {\scriptsize $w_{n}$};
\draw[decorate,decoration={brace,amplitude=5pt},gray!60!black] (-4.9,-1.2) -- (-4.9,1.2);
\node[left,align=right,gray!40!black] at (-5.1,0) {\small $n\sim10^3$--$10^5$\\[-2pt]\small entrées};
% single output wire, then fan-out
\draw[very thick,red!70!black] (1.3,0) -- (3.2,0);
\foreach \x in {1.6,2.2,2.34,2.9}{\draw[thick,red!70!black] (\x,0) -- (\x,0.4);}
\node[above right,font=\tiny\itshape,gray!30!black,inner sep=1pt] at (1.42,0.45) {clic\dots\ clic-clic\dots};
\node[below,red!70!black,align=center] at (2.25,-0.05) {\scriptsize une sortie $y$};
\fill[red!70!black] (3.2,0) circle (0.06);
\foreach \y in {1.3,0.65,0,-0.65,-1.3}{
  \draw[->,thick,red!70!black] (3.2,0) -- (5.0,\y);}
\draw[decorate,decoration={brace,amplitude=5pt},gray!60!black] (5.3,1.45) -- (5.3,-1.45);
\node[right,align=left,gray!40!black] at (5.5,0) {\small copies de $y$\\[-2pt]\small vers $10^3$--$10^4$\\[-2pt]\small autres neurones};
\end{tikzpicture}
\caption{Le neurone vu par un ingénieur. De nombreuses entrées $x_i$, chacune avec son poids $w_i$; à l'intérieur, une somme pondérée $u$ et une comparaison au seuil $\theta$ ($\Theta$ est la fonction échelon: $1$ si l'argument est positif, $0$ sinon); une seule sortie $y$, une suite d'impulsions démultipliée par ramification vers des milliers de destinataires.}
\label{fig:blackbox}
\end{figure}

La sortie du bloc n'est pas un nombre, mais une suite de brèves impulsions de tension: \og clic\fg, une pause, \og clic-clic\fg, une longue pause. Toutes les impulsions ont la même hauteur, si bien que toute l'information réside dans le \emph{moment} où elles se produisent. À l'intérieur du bloc se trouvent, en première approximation grossière, un additionneur et un comparateur: les entrées sont sommées avec des poids, et si la somme franchit un seuil, le bloc émet une impulsion. Au chapitre suivant, nous le remplacerons par un modèle qui fonctionne honnêtement en temps continu.

La sortie est unique, mais le fil se ramifie, et chaque branche porte \emph{exactement la même} copie du signal. En électronique, on appelle cela la sortance, le fan-out. Pour un neurone du cortex, elle est de l'ordre de quelques milliers; une porte logique d'un circuit intégré ne pilote en général que quelques autres portes.

\emph{Que} transmet donc le fil de sortie au destinataire? L'impulsion court jusqu'au bout de chaque branche et s'y transforme en signal chimique: la branche libère dans l'étroite fente entre les cellules une pincée d'une substance déterminée, un \emph{neurotransmetteur}, qui modifie le courant d'entrée du destinataire. Voilà notre critère de tri: \emph{quel neurotransmetteur libère la sortie du neurone}.

\section{Une sortie, un type de signal}

Ce tri n'a de sens que si chaque neurone n'a qu'un seul type de sortie. Est-ce le cas? L'expérience dit que presque.

\begin{keyblock}
\begin{proposition}[Un neurone, un type de sortie]
\label{prop:dale}
Presque chaque neurone a un seul neurotransmetteur principal, et il le libère dans toutes les branches de sa sortie. Le type d'un neurone est donc défini par son neurotransmetteur principal~\cite{Strata1999}.
\end{proposition}
\end{keyblock}

Convenons de désigner le type d'un neurone par \emph{les quatre premières lettres du nom anglais de son neurotransmetteur principal}. Ainsi, un neurone dont le neurotransmetteur principal est le glutamate est un neurone \nt{GLUT}. Tous les types sont réunis dans le tableau~\ref{tab:types}.

\begin{table}[t]
\centering
\footnotesize
\setlength{\tabcolsep}{3pt}
\renewcommand{\arraystretch}{1.15}
\begin{tabular}{>{\raggedright}p{1.0cm}>{\raggedright}p{2.05cm}>{\raggedright}p{1.75cm}>{\raggedright}p{1.6cm}>{\centering}p{0.7cm}>{\raggedright}p{1.55cm}>{\raggedright}p{1.8cm}>{\raggedright\arraybackslash}p{3.2cm}}
\toprule
Type & Neurotrans\-metteur principal & Bus & Vitesse & Signe & Neurones dans le cerveau & Sorties par neurone & Rôle dans le calcul \\
\midrule
\nt{GLUT} & glutamate & adresses & rapide et lente & $+$ & $\sim8\cdot10^{10}$ & $\sim10^4$ & principal poids positif \\
\nt{GABA} & GABA & adresses & rapide et lente & $-$ & $\sim3\cdot10^{9}$ & $10^3$--$10^4$ & principal poids négatif \\
\nt{GLYC} & glycine & adresses & rapide & $-$ & $10^6$--$10^7$ & $10^2$--$10^3$ & poids négatif dans la moelle épinière et le tronc cérébral; seconde clé d'un des récepteurs du glutamate \\
\nt{ACET} & acétylcholine & les deux & rapide et lente & $\pm$ & $\sim10^6$ & muscle: $10$--$10^3$; cerveau: $\sim10^5$ & sortie vers le muscle; dans le cerveau, vitesse d'apprentissage (hypothèse) \\
\nt{DOPA} & dopamine & diffusion & lente & $\pm$ & $\sim5\cdot10^{5}$ & $10^5$--$10^6$ & erreur de prédiction de la récompense $\delta$ \\
\nt{SERO} & sérotonine & diffusion & lente (un type de récepteur rapide) & $\pm$ & $\sim3\cdot10^{5}$ & $\sim10^5$ & horizon de planification (hypothèse) \\
\nt{HIST} & histamine & diffusion & lente & $+$ & $\sim6\cdot10^{4}$ & pas d'estimation & signal global \og reste éveillé\fg{} \\
\nt{NORA} & noradrénaline & diffusion & lente & $\pm$ & $\sim5\cdot10^{4}$ & $\sim10^5$ & gain (hypothèse) \\
\nt{ADRE} & adrénaline & diffusion & lente & $\pm$ & $\sim10^4$ & pas d'estimation & peu de neurones; rôle dans le cerveau obscur \\
\nt{ASPA} & aspartate & adresses & rapide & $+$ & pas d'estimation & pas d'estimation & faible poids positif; rôle contesté \\
\bottomrule
\end{tabular}
\caption{Types de neurones, par nombre décroissant de neurones dans le cerveau. \emph{Bus}: adresses, le neurotransmetteur est libéré dans une fente étroite vers un destinataire déterminé; diffusion, d'un petit groupe de neurones sources vers de vastes régions (section~\ref{sec:buses}). \emph{Vitesse}: rapide, par des récepteurs ionotropes, en millisecondes; lente, par des récepteurs métabotropes, en centaines de millisecondes ou plus. \emph{Signe}: le signe habituel; c'est le destinataire qui le fixe en dernier ressort (section~\ref{sec:receiver}), et $\pm$ signifie que l'on rencontre les deux signes. \emph{Neurones dans le cerveau}: combien de neurones de ce type compte le cerveau humain entier, en ordre de grandeur; au total, environ $8.6\cdot10^{10}$ neurones~\cite{Azevedo2009}. Presque tous les neurones \nt{GLUT}, environ $7\cdot10^{10}$, sont les petits neurones d'entrée du cervelet; les nombres de \nt{GLUT} et de \nt{GABA} découlent de ce comptage et de la proportion de neurones \nt{GABA} dans le cortex~\cite{Hendry1987}. Parmi les types peu nombreux, seuls \nt{HIST}~\cite{Airaksinen1991} et le principal des deux groupes de neurones \nt{DOPA}~\cite{Pakkenberg1991} ont été comptés directement, si bien que pour \nt{DOPA} il s'agit d'une borne inférieure; pour \nt{SERO}, c'est la somme de deux comptages partiels~\cite{Baker1990,Baker1991}. Les nombres pour \nt{NORA}, \nt{GLYC}, \nt{ACET} et \nt{ADRE} sont des estimations grossières en ordre de grandeur, sans comptage direct. \emph{Sorties par neurone}: combien un neurone possède de terminaisons synaptiques, c'est-à-dire de sites de libération du neurotransmetteur sur les branches du fil de sortie (fig.~\ref{fig:blackbox}), en ordre de grandeur. Pour les types à diffusion, les terminaisons n'ont pas été comptées directement: pour \nt{DOPA}, le nombre est déduit de la part de sa cible que couvre l'arborisation d'un seul fil de sortie~\cite{Matsuda2009}; pour les autres, les estimations sont encore plus grossières. Pour \nt{ACET}, deux cas: le neurone qui commande un muscle et le neurone de diffusion dans le cerveau.}
\label{tab:types}
\end{table}

% the pie is a table-type float with a figure caption, so it can never float ahead of Table~\ref{tab:types}
\begin{table}[tb]
\makeatletter\def\@captype{figure}\makeatother
\centering
\definecolor{vizglut}{HTML}{2A78D6}
\definecolor{vizgaba}{HTML}{E34948}
\definecolor{vizrest}{HTML}{8A8986}
\begin{tikzpicture}[>=Stealth]
% pie: GLUT 96.4 %, GABA 3.6 % (13.0 degrees), the rest 0.006 % (a hairline at 0 degrees)
\fill[vizglut] (0,0) -- (13:1.9) arc (13:360:1.9) -- cycle;
\fill[vizgaba] (0,0) -- (0:1.9) arc (0:13:1.9) -- cycle;
\draw[white,line width=1pt] (0,0) -- (13:1.9);
\draw[vizrest,line width=1.2pt] (0,0) -- (0:1.9);
\node[white,align=center,font=\small] at (190:0.95) {\nt{GLUT}\\$96.4\,\%$};
\draw[gray!60!black] (6.5:1.75) -- (2.05,2.1);
\node[anchor=west,font=\small] at (2.05,2.1) {\nt{GABA} $3.6\,\%$};
% zoom box for the remaining types
\draw[densely dashed,vizrest] (1.9,0) -- (3.1,1.65);
\draw[densely dashed,vizrest] (1.9,0) -- (3.1,-2.2);
\draw[vizrest,rounded corners=3pt] (3.1,-2.2) rectangle (10.1,1.65);
\node[anchor=west,font=\scriptsize] at (3.2,1.4) {tous les autres: $\approx0.006\,\%$; échelle logarithmique};
\foreach \code/\lg/\y/\val/\lx in {GLYC/6.477/1.0/{$10^6$--$10^7$}/4.05,ACET/6.0/0.6/{$\sim10^6$}/3.0,DOPA/5.699/0.2/{$\sim5\cdot10^5$}/2.699,SERO/5.477/-0.2/{$\sim3\cdot10^5$}/2.477,HIST/4.778/-0.6/{$\sim6\cdot10^4$}/1.778,NORA/4.699/-1.0/{$\sim5\cdot10^4$}/1.699,ADRE/4.0/-1.4/{$\sim10^4$}/1.0}{
  \node[anchor=east,font=\scriptsize] at (4.35,\y) {\nt{\code}};
  \fill[vizrest] (4.45,\y-0.13) rectangle ({4.45+\lg-3},\y+0.13);
  \node[anchor=west,font=\scriptsize] at ({4.5+\lx},\y) {\val};}
\draw[vizrest,thick] (7.45,1.0) -- (8.45,1.0);
\draw[vizrest,thick] (7.45,0.92) -- (7.45,1.08);
\draw[vizrest,thick] (8.45,0.92) -- (8.45,1.08);
\draw[gray!60!black] (4.45,-1.72) -- (8.45,-1.72);
\foreach \e in {3,4,5,6,7}{
  \draw[gray!60!black] ({4.45+\e-3},-1.72) -- ({4.45+\e-3},-1.66);
  \node[below,font=\scriptsize] at ({4.45+\e-3},-1.72) {$10^{\e}$};}
\end{tikzpicture}
\caption{Parts des types de neurones dans le cerveau, d'après les estimations du tableau~\ref{tab:types}. Le disque est presque entièrement occupé par \nt{GLUT}; \nt{GABA} forme un secteur étroit, et tous les autres types ensemble ne sont qu'un cheveu sur la frontière. À droite, ce cheveu est agrandi, sur une échelle logarithmique du nombre de neurones; pour \nt{GLYC}, la barre s'arrête au milieu de l'intervalle $10^6$--$10^7$, et le segment montre l'intervalle entier. \nt{ASPA} n'a pas d'estimation et ne figure pas sur le dessin.}
\label{fig:typepie}
\end{table}

La figure~\ref{fig:typepie} montre à quel point ces tas sont inégaux: sur mille neurones, $964$ sont \nt{GLUT}, $36$ sont \nt{GABA}, et tous les autres types réunis représentent environ six neurones sur cent mille.

Le nom GABA compte déjà quatre lettres. Les substances qui ne sont presque jamais le neurotransmetteur principal (neuropeptides, gaz, lipides, purines) ne reçoivent pas de code propre; il en est question dans l'annexe~\ref{sec:othermediators}. Le mot \og presque\fg{} de la proposition~\ref{prop:dale} compte. Beaucoup de neurones libèrent, en plus du neurotransmetteur principal, des substances auxiliaires~\cite{Yao2023}. Certains libèrent aussi un second neurotransmetteur rapide, parfois de signe opposé: une partie des neurones \nt{DOPA} libère également du GABA~\cite{Tritsch2012}. Et chez certains neurones, des neurotransmetteurs différents sortent de branches différentes d'une même sortie~\cite{Zhang2015}. Tout cela est mesuré: \og un neurone, un neurotransmetteur\fg{} est donc une règle avec des exceptions, pas une loi. Mais comme première division, elle tient bon: dans l'atlas cellulaire du cerveau de la souris, où les neurones sont répartis selon les gènes qu'ils expriment en plus de cinq mille groupes, les neurones se divisent toujours en grandes classes d'abord selon leur neurotransmetteur principal~\cite{Yao2023}.

Cette règle a une conséquence qui distingue d'emblée le cerveau d'un réseau de neurones artificiel. Dans un réseau artificiel, un même neurone peut avoir un poids positif vers un destinataire et négatif vers un autre: les poids $w_{ij}$ ne sont que des nombres dans une matrice. Dans le cerveau, le signe de l'action est fixé surtout par le neurotransmetteur, et un neurone n'en a qu'un. Donc, si le neurone $j$ excite un destinataire, il excite aussi tous les autres. Toute la \emph{colonne} de la matrice des poids issue du neurone $j$ est de même signe:
\begin{empheq}[box=\keybox]{equation}
\operatorname{sign} w_{ij} \;=\; s_j \quad\text{pour tous les destinataires } i, \qquad s_j\in\{+1,-1\}.
\label{eq:dalesign}
\end{empheq}
Les neurones se divisent en excitateurs ($s_j=+1$) et inhibiteurs ($s_j=-1$), et c'est une propriété du neurone lui-même, non de la connexion. Si le réseau a besoin d'une connexion négative issue d'un neurone excitateur, il doit passer par un intermédiaire: le neurone excitateur excite un neurone inhibiteur, qui inhibe la cible. On obtient un poids négatif retardé d'un maillon.

On connaît plus d'une centaine de neurotransmetteurs, mais presque tout le travail du cerveau repose sur une demi-douzaine d'entre eux, qui se répartissent en deux classes entièrement différentes.

\section{Deux bus: d'adresses et de diffusion}
\label{sec:buses}

Les réseaux informatiques connaissent deux façons de livrer un message. La première est \emph{adressée}, l'unicast: le paquet va d'un expéditeur à un destinataire déterminé, rapidement, et personne d'autre ne le voit. La seconde est la \emph{diffusion}, le broadcast: l'expéditeur transmet un seul message à tous les nœuds du réseau à la fois. Les neurones utilisent les deux, et les neurotransmetteurs se répartissent exactement selon ce critère (fig.~\ref{fig:phone_radio}).

\begin{figure}[t]
\centering
\begin{tikzpicture}[>=Stealth,scale=0.95]
% ---- unicast: point-to-point wiring
\node[above] at (2.0,3.3) {\small \textbf{Bus d'adresses}: glutamate et GABA};
\foreach \i in {0,...,4}{
  \fill[blue!60!black] (0.2+0.9*\i,2.5) circle (0.1);
  \fill[blue!60!black] (0.2+0.9*\i,0.6) circle (0.1);}
\foreach \a/\b in {0/0,0/1,1/1,1/3,2/2,3/2,3/4,4/4,2/0}{
  \draw[->,blue!60!black] (0.2+0.9*\a,2.38) -- (0.2+0.9*\b,0.72);}
\fill[red!70!black] (1.1,1.55) circle (0.09);
\pic[scale=1.2] at (1.75,1.9) {letter};
\draw[->,red!70!black,thick] (1.1,1.55) -- (0.28,0.7);
\draw[->,red!70!black,thick] (1.1,1.55) -- (1.95,0.72);
\node[align=center,below] at (2.0,0.2) {\small des milliards de neurones,\\[-2pt]\small chacun ses destinataires,\\[-2pt]\small temps: millisecondes};
% ---- broadcast: one small source, global targets
\begin{scope}[xshift=7.2cm]
\node[above] at (1.8,3.3) {\small \textbf{Bus de diffusion}: dopamine, \dots};
\foreach \i in {0,...,8}{\fill[blue!60!black] (-0.2+0.5*\i,2.75) circle (0.08);}
\fill[orange!85!black] (1.8,0.35) ellipse (0.35 and 0.18);
\pic[scale=1.5] at (2.7,0.75) {mast};
\foreach \x in {-0.2,0.3,0.8,1.3,1.8,2.3,2.8,3.3,3.8}{
  \draw[->,orange!85!black] (1.8,0.5) .. controls (1.8,1.3) and (\x,1.6) .. (\x,2.62);}
\node[align=center,below] at (1.8,0.1) {\small des milliers à des centaines de milliers\\[-2pt]\small de neurones, un seul message pour tous,\\[-2pt]\small temps: centaines de millisecondes et plus};
\end{scope}
\end{tikzpicture}
\caption{Deux modes de transmission. À gauche, le bus d'adresses, semblable au courrier avec l'adresse sur l'enveloppe; en rouge, un neurone et deux de ses destinataires. À droite, la diffusion, comme une station de radio: un petit groupe de neurones sources, et tous reçoivent le même message.}
\label{fig:phone_radio}
\end{figure}

Le \emph{bus d'adresses}, ce sont le glutamate et le GABA. L'écrasante majorité des neurones les libèrent. Chaque sortie mène à des cellules déterminées, le neurotransmetteur n'agit que dans la fente étroite du contact, et il agit vite: en une milliseconde, les canaux ioniques du destinataire s'ouvrent; en quelques millisecondes, tout est fini. C'est le bus de \emph{données}: il transporte ce que le cerveau calcule.

Le \emph{bus de diffusion}, ce sont la dopamine, la sérotonine, la noradrénaline et en partie l'acétylcholine. Ils sont libérés par de minuscules groupes de neurones, mais la sortie de chacun d'eux se ramifie de façon incroyablement large. Le neurotransmetteur est souvent libéré non dans une fente étroite, mais dans l'espace extracellulaire, où il se répand et agit sur tous ceux qui se trouvent à proximité. Il agit lentement: il n'ouvre pas les canaux directement, mais déclenche à l'intérieur du destinataire une chaîne de réactions chimiques. Ce bus ne transporte pas des données, mais des \emph{paramètres de contrôle}, communs à de grandes portions du réseau, qui en changent le mode de fonctionnement: le gain, la vitesse d'apprentissage, le signal \og c'était bien\fg. C'est pourquoi on appelle ces neurotransmetteurs des \emph{neuromodulateurs}. On a longtemps pensé que chacun de ces canaux était un seul nombre pour tout le cerveau. Les mesures modernes ont précisé le tableau: un canal n'est pas un fil unique, mais un petit faisceau de lignes parallèles. Les neurones sources d'un même neurotransmetteur se divisent en sous-systèmes qui ont leurs propres destinataires et leur propre message, et la quantité de neurotransmetteur libérée sur place est en outre ajustée par des signaux locaux~\cite{Ren2018,Poe2020}. Ces lignes se comptent en unités ou en dizaines, pas en milliards: le canal reste donc un canal de diffusion.

S'il ne faut retenir qu'une image de ce chapitre, c'est celle-ci. Le cerveau est un immense réseau à transmission adressée des données, surmonté de quelques canaux de diffusion porteurs de signaux de contrôle. Un programmeur reconnaîtra une architecture familière: du calcul, plus un petit ensemble de variables de contrôle, chacune commune à une grande portion du réseau.

\section{\nt{GLUT}: le poids positif}

Le glutamate est le principal neurotransmetteur excitateur du cerveau. Quand la sortie d'un neurone libère du glutamate, des canaux s'ouvrent chez le destinataire, le sodium entre, la tension de sa membrane monte, et sa chance d'émettre sa propre impulsion augmente. Dans le langage de la fig.~\ref{fig:blackbox}, c'est une entrée de poids \emph{positif}, $w_i>0$.

Qui libère du glutamate? Presque tous ceux qui transmettent des données sur de longues distances: de trois quarts à quatre cinquièmes des neurones du cortex et la plupart des neurones qui relient différentes régions du cerveau. Le réseau du cerveau est avant tout un réseau de connexions glutamatergiques.

Un détail d'ingénierie. Après la libération, la concentration de glutamate dans la fente bondit d'un facteur de plusieurs milliers, jusqu'à environ une millimole par litre, puis retombe avec une constante de temps d'environ une milliseconde~\cite{Clements1992}: le glutamate diffuse hors de la fente, et des protéines-pompes spécialisées le réabsorbent dans les cellules. Pourquoi? Pour la même raison que, dans une ligne de communication, on ramène le signal à zéro après chaque symbole. Si le neurotransmetteur n'est pas éliminé, l'impulsion suivante arrive sur une ligne \og occupée\fg, et des impulsions espacées de quelques millisecondes se confondent.

\section{\nt{GABA}: le poids négatif}

Imaginez un réseau dont tous les poids sont positifs. Si chaque impulsion engendre en moyenne plus d'une impulsion suivante, l'activité croît exponentiellement et envahit tout le réseau en quelques pas. Un électronicien reconnaîtra une boucle de rétroaction positive de gain supérieur à un: le circuit part en saturation. Pour l'éviter, il faut des poids négatifs. Leur source principale est l'acide gamma-aminobutyrique, en abrégé GABA.

Le GABA ouvre chez le destinataire des canaux perméables aux ions chlorure. Le potentiel d'équilibre du chlorure se situe près du potentiel de repos ou un peu en dessous; les canaux chlorure ouverts maintiennent donc la membrane près du repos et empêchent la tension de monter jusqu'au seuil. C'est une entrée de poids \emph{négatif}, $w_i<0$. Les neurones GABA représentent d'un cinquième à un quart des neurones du cortex~\cite{Hendry1987}. Leurs sorties sont courtes, et ils inhibent leurs proches voisins.

Dans le cerveau, l'inhibition fait bien plus qu'une simple soustraction. Elle agit comme une rétroaction négative qui maintient tout le réseau à son point de fonctionnement. On considère que, dans un cortex qui fonctionne normalement, les courants excitateur et inhibiteur qui arrivent sur un neurone s'équilibrent presque exactement: ce régime est prédit par la théorie~\cite{vanVreeswijk1996} et visible dans les mesures de courants dans le cortex vivant~\cite{Haider2006}, et le neurone reste en permanence près du seuil. Un circuit stabilisé par une forte rétroaction négative autour d'un point instable réagit vite et finement aux petits signaux: c'est une vieille astuce d'ingénieur.

\section{\nt{DOPA}: le signal d'erreur}

Le premier canal de diffusion est la dopamine. L'être humain possède de l'ordre d'un demi-million de neurones dopaminergiques, environ un sur cent soixante-dix mille; c'est le nombre compté dans le principal de leurs deux groupes~\cite{Pakkenberg1991}, il s'agit donc d'une borne inférieure. Leurs sorties se dirigent vers les régions qui choisissent les actions.

Que transmet ce canal? La réponse vient d'expériences où l'on enregistre les impulsions des neurones dopaminergiques d'un animal qui reçoit une récompense~\cite{Schultz1997}. De temps en temps, on donne à l'animal une goutte de jus, sans prévenir, et les neurones dopaminergiques y répondent par une brève bouffée d'impulsions. Puis on se met à allumer une lampe avant le jus, toujours une seconde avant. Quand l'animal a appris cette association, les neurones se mettent à répondre par une bouffée à la \emph{lampe}, et ne répondent plus du tout au jus lui-même s'il arrive à l'heure. Et si, après la lampe, le jus ne vient pas, les neurones se \emph{taisent} au moment où il aurait dû arriver, et leur fréquence tombe sous son niveau habituel.

La règle qui explique les trois cas (fig.~\ref{fig:rpe}): les neurones n'annoncent pas à quel point c'est bien, mais à quel point c'est \emph{mieux ou pire que prévu}.

\begin{figure}[t]
\centering
\begin{tikzpicture}[>=Stealth,x=3.4cm,y=1cm]
\foreach \y/\lab in {2.8/{jus inattendu},1.4/{jus après la lampe},0/{lampe, mais pas de jus}}{
  \draw[gray!70!black] (0,\y) -- (3,\y);
  \draw[densely dotted,gray!60!black] (0,\y+0.3) -- (3,\y+0.3);
  \node[left,align=right,font=\small] at (-0.03,\y+0.35) {\lab};}
% row 1: juice without a cue -> burst at the juice
\draw[very thick,orange!85!black] plot[domain=0:3,samples=120] (\x,{2.8+0.3+0.75*exp(-((\x-2.08)/0.07)^2)});
\pic[scale=0.8] at (2,2.58) {juice};
% row 2: cue then juice -> burst at the cue, nothing at the juice
\draw[very thick,orange!85!black] plot[domain=0:3,samples=120] (\x,{1.4+0.3+0.75*exp(-((\x-1.08)/0.07)^2)});
\pic[scale=0.85] at (1,1.15) {lamp}; \pic[scale=0.85] at (2,1.15) {juice};
% row 3: cue, juice omitted -> burst at the cue, dip when the juice was due
\draw[very thick,orange!85!black] plot[domain=0:3,samples=120] (\x,{0.3+0.75*exp(-((\x-1.08)/0.07)^2)-0.28*exp(-((\x-2.1)/0.12)^2)});
\pic[scale=0.85] at (1,-0.25) {lamp}; \pic[scale=0.85] at (2,-0.25) {nojuice};
\node[right,font=\scriptsize] at (2.36,0.1) {creux};
\node[right,font=\scriptsize] at (2.13,3.75) {surprise!};
\node[left,font=\scriptsize] at (0.97,2.25) {pic à la lampe};
\node[above,font=\scriptsize] at (2,1.75) {rien};
\draw[->,gray!70!black] (0,-0.75) -- (3.05,-0.75) node[right,black,font=\small] {$t$};
\draw[gray!60!black] (1,-0.8) -- (1,-0.7) node[below=2pt,font=\scriptsize] {lampe, $1\,\text{s}$};
\draw[gray!60!black] (2,-0.8) -- (2,-0.7) node[below=2pt,font=\scriptsize] {jus, $2\,\text{s}$};
\end{tikzpicture}
\caption{Fréquence des impulsions des neurones \nt{DOPA} dans trois expériences (schéma d'après~\cite{Schultz1997}). Le pointillé est la fréquence de fond, constante. La lampe est le signal, la goutte est le jus, la goutte barrée est un jus promis mais non donné. La réponse est l'erreur de prédiction~\eqref{eq:rpe}.}
\label{fig:rpe}
\end{figure}

Un programmeur qui connaît l'apprentissage par renforcement~\cite{SuttonBarto2018} reconnaîtra tout de suite cette grandeur. Un agent qui apprend à choisir ses actions garde une estimation $V(s)$: combien de récompense il attend quand il se trouve dans l'état $s$. Après chaque pas, il compare cette attente à ce qu'il a reçu:
\begin{empheq}[box=\keybox]{equation}
\delta_t \;=\; r_t \;+\; \gamma\,V(s_{t+1}) \;-\; V(s_t) ,
\label{eq:rpe}
\end{empheq}
et corrige son estimation: $V(s_t)\leftarrow V(s_t)+\alpha\,\delta_t$. Ici $r_t$ est la récompense reçue, $\gamma<1$ le facteur d'actualisation (combien une récompense future vaut moins qu'une récompense immédiate), $\alpha$ la vitesse d'apprentissage. La grandeur $\delta_t$ s'appelle l'\emph{erreur de différence temporelle}.

Elle se comporte exactement comme les neurones dopaminergiques. Jus inattendu: $V$ vaut encore zéro, $r>0$, $\delta>0$. Jus appris: la lampe a prédit la récompense, $V(s_t)$ avant le jus vaut déjà $r$, et $\delta=0$. Le jus n'est pas venu: $V(s_t)=r$, mais $r_t=0$, $\delta<0$. Et le pic se déplace vers la lampe parce que c'est au moment de la lampe que l'attente a bondi: $\gamma V(s_{t+1})-V(s_t)>0$. Les pas $t,t+1,\dots$ ne sont ici qu'une convention de l'algorithme: l'organisme n'a pas d'horloge qui égrène des pas, et nous obtiendrons la même grandeur en temps continu au chapitre~\ref{sec:dofa}.

La dopamine est donc la diffusion d'un signal d'erreur scalaire $\delta$ vers tous ceux qui doivent apprendre d'après lui. En première approximation, c'est un seul fil pour tout le cerveau et un seul nombre à chaque instant; à quel point ce fil est en réalité un faisceau, le chapitre~\ref{sec:dofaalt} l'examine.

\section{Les autres canaux de diffusion}

Pour les autres neuromodulateurs, on ne dispose que d'hypothèses de travail, qui traduisent leurs rôles dans le même langage de paramètres globaux~\cite{Doya2002}. Tenez-les bien pour des hypothèses.

\emph{\nt{NORA}, la noradrénaline: le gain.} Les neurones noradrénergiques sont encore moins nombreux que les dopaminergiques, quelques dizaines de milliers selon une estimation grossière, et leurs sorties se répandent pratiquement dans tout le cerveau. Ils s'activent brusquement quand il se passe quelque chose d'inattendu ou d'important. La noradrénaline modifie la pente de la caractéristique de transfert des neurones destinataires: les entrées faibles deviennent plus faibles, les fortes plus fortes. On le mesure ainsi: les impulsions de fond du destinataire sont plus supprimées que sa réponse à l'entrée, et le rapport signal sur bruit augmente~\cite{Foote1975NE}. Dans un modèle simple, on décrit cela par un seul nombre: si le neurone répond par une fréquence $f=F\bigl(g\cdot u\bigr)$ à un courant d'entrée $u$, la noradrénaline augmente le coefficient $g$~\cite{AstonJones2005} (détails au chapitre~\ref{sec:nora}). Un réseau à grand $g$ fonctionne de façon plus \og contrastée\fg{} et décide plus vite; à petit $g$, de façon \og floue\fg, en essayant davantage d'options, comme un agent d'apprentissage par renforcement en mode exploration.

\emph{\nt{ACET}, l'acétylcholine: la vitesse d'apprentissage.} L'acétylcholine règle la part d'écoute que le réseau accorde à l'entrée courante et à ses propres données mémorisées. Quand son niveau est élevé, les entrées venues des organes des sens sont renforcées, les connexions internes \og de rappel\fg{} affaiblies, et en même temps la modification des poids est facilitée~\cite{Hasselmo2006}. En gros, c'est un commutateur \og écriture/lecture\fg{} et, avec lui, la vitesse d'apprentissage $\alpha$.

\emph{\nt{SERO}, la sérotonine: l'horizon de planification.} Ce que transmet la sérotonine est le moins bien compris. Une hypothèse la relie au facteur d'actualisation $\gamma$ de~\eqref{eq:rpe}: un niveau élevé de sérotonine correspondrait à la disposition à attendre une grosse récompense plutôt qu'à saisir tout de suite une petite. Les neurones sérotoninergiques travaillent de façon régulière et lente, quelques impulsions par seconde pendant l'éveil, et se taisent presque pendant l'une des phases du sommeil~\cite{JacobsAzmitia1992}.

Les six neurotransmetteurs sont réunis dans le tableau~\ref{tab:transmitters}.

\begin{table}[t]
\centering
\small
\begin{tabular}{>{\raggedright}p{2.4cm}>{\raggedright}p{3.0cm}>{\raggedright}p{2.0cm}>{\raggedright\arraybackslash}p{5.2cm}}
\toprule
Type de neurone & Part des neurones & Bus & Rôle dans le calcul \\
\midrule
\nt{GLUT} (glutamate) & majorité & adresses & poids positif, $w>0$ \\
\nt{GABA} (GABA) & 20--25\% dans le cortex & adresses & poids négatif, $w<0$; stabilisation \\
\nt{DOPA} (dopamine) & $\sim 6\cdot10^{-6}$ & diffusion & erreur de prédiction $\delta$ \\
\nt{NORA} (noradrénaline) & $\sim 6\cdot10^{-7}$ & diffusion & gain $g$ (hypothèse) \\
\nt{ACET} (acétylcholine) & faible & les deux & vitesse d'apprentissage $\alpha$; \og écriture/lecture\fg{} (hypothèse) \\
\nt{SERO} (sérotonine) & $\sim 3\cdot10^{-6}$ & diffusion & actualisation $\gamma$ (hypothèse) \\
\bottomrule
\end{tabular}
\caption{Les principaux neurotransmetteurs du cerveau dans le langage du calcul. La colonne de droite est une forte simplification: fiable pour le glutamate, le GABA et la dopamine, hypothétique pour les autres.}
\label{tab:transmitters}
\end{table}

\section{Le destinataire fixe le signe}
\label{sec:receiver}

Je vous ai un peu trompé en disant que le glutamate est un poids positif et le GABA un poids négatif. Aucune molécule ne porte de signe en elle-même. Une molécule n'est qu'une clé. Ce qui se passera quand elle sera captée, c'est la \emph{serrure} qui en décide: le récepteur de la cellule destinataire.

Le meilleur exemple est la dopamine. Elle a des récepteurs de deux familles~\cite{Beaulieu2011}. Sur les uns, elle facilite l'excitation de la cellule; sur les autres, elle la freine. Dans la région qui choisit les actions, certains neurones portent des récepteurs du premier type, d'autres du second, et un même signal dopaminergique renforce un groupe tout en affaiblissant l'autre. C'est comme un seul paquet diffusé que différents nœuds du réseau interprètent différemment.

\begin{keyblock}
\begin{proposition}[Le destinataire fixe le sens du message]
\label{prop:receptor}
Le signe et la vitesse de l'action d'un neurotransmetteur ne sont pas fixés par le neurotransmetteur lui-même, mais par le récepteur auquel il se lie. Les récepteurs sont de deux sortes. Les récepteurs \emph{ionotropes} sont eux-mêmes des canaux ioniques et s'ouvrent en une fraction de milliseconde: c'est une commande directe du courant, comme la grille d'un transistor. Les récepteurs \emph{métabotropes} déclenchent à l'intérieur de la cellule une chaîne de réactions chimiques et agissent en centaines de millisecondes ou plus: c'est plutôt une écriture dans un registre de configuration, qui modifie le travail ultérieur de la cellule. Le bus d'adresses passe surtout par les premiers, le bus de diffusion surtout par les seconds.
\end{proposition}
\end{keyblock}

Pourquoi peut-on alors dire que \og le glutamate excite\fg? Parce que presque tous les récepteurs du glutamate que l'on rencontre en pratique sont excitateurs, et presque tous les récepteurs du GABA inhibiteurs. Ce n'est pas une loi de la nature, mais une convention très fiable, et c'est sur elle que repose la règle~\eqref{eq:dalesign}.

\section{Ce qu'il faut retenir}

L'écrasante majorité des neurones forme le bus d'adresses qui porte les données, avec des poids d'un seul signe par neurone; une infime minorité forme des canaux de diffusion qui transmettent à de grandes portions du cerveau quelques signaux de contrôle communs. Environ deux neurones sur cent mille, et c'est d'eux que dépendent la façon dont apprend tout le reste du réseau et le régime dans lequel il fonctionne. Pour un ingénieur, rien d'étonnant: dans tout ordinateur, les registres de contrôle sont incomparablement moins nombreux que les cellules de mémoire, et pourtant la machine ne fonctionne pas sans eux.

Au chapitre suivant, nous verrons ce que peut calculer, en temps continu, un réseau fait uniquement de neurones \nt{GLUT}, le type le plus nombreux.

\section{Exercices}

\begin{exercise}[\lvlA]
\label{ex:01:1}
\emph{Tas et fils.} D'après le tableau~\ref{tab:types} (milieu de l'intervalle en échelle logarithmique; \nt{ACET}: $10^5$ sorties): (a)~si un neurone était une personne, combien de populations de la Terre feraient les neurones \nt{GLUT}, et quelle taille de ville tous ceux de diffusion? (b)~Combien de fois la part des terminaisons de \nt{DOPA}, \nt{SERO}, \nt{NORA}, \nt{ACET} dépasse-t-elle la part de leurs neurones?
\end{exercise}

\begin{exercise}[\lvlA]
\label{ex:01:2}
\emph{Retour à zéro.} Le glutamate dans la fente décroît comme $c_0e^{-t/\tau_c}$, $c_0=1\mM$, $\tau_c=1\ms$; le destinataire détecte une concentration supérieure à $10$~\textmu M. (a)~Combien de temps la ligne reste-t-elle \og occupée\fg{} après une libération? (b)~Les pompes sont en partie bloquées, $\tau_c=10\ms$. Jusqu'à quelle fréquence de libérations la ligne a-t-elle encore le temps de se libérer?
\end{exercise}

\begin{exercise}[\lvlB]
\label{ex:01:3}
\emph{Où se déplace le pic.} Dans l'expérience de la fig.~\ref{fig:rpe}, la lampe est suivie des états $s_0\to\dots\to s_{K-1}$ au pas $\Delta$, $T=K\Delta$; la récompense $r$ arrive à la dernière transition, l'instant de la lampe est imprévisible. Trouvez les $V(s_k)$ appris et la réponse $\delta$ à la lampe. En posant $\gamma=e^{-\Delta/\tau_\gamma}$, trouvez $\tau_\gamma$ pour $T=1$~s, $\Delta=0.1$~s, $\gamma=0.95$.
\end{exercise}

\begin{exercise}[\lvlB]
\label{ex:01:4}
\emph{Modélisation: apprendre par l'erreur.} Simulez la chaîne de l'exercice~\ref{ex:01:3} avec $K=10$, $\gamma=0.95$, $r=1$ et la règle $V(s_t)\leftarrow V(s_t)+\alpha\delta_t$. À quel essai $n^*$ la réponse à la lampe dépasse-t-elle pour la première fois la réponse au jus, pour $\alpha=0.05$, $0.1$, $0.2$, $0.5$? Comment $n^*$ dépend-il de $\alpha$?
\end{exercise}

\begin{exercise}[\lvlB]
\label{ex:01:5}
\emph{Conception: diffuser un seul nombre.} Il faut livrer le nombre $\delta$ aux $10^{10}$ neurones; chaque destinataire, relais compris, doit recevoir $10$ terminaisons de la couche précédente, et chaque maillon ajoute $1\ms$. Combien de neurones et de maillons compte l'arbre de relais le plus économe avec $10^4$ sorties par neurone (\nt{GLUT}) et avec $10^{5.5}$ (\nt{DOPA})?
\end{exercise}

\begin{exercise}[\lvlB]
\label{ex:01:6}
\emph{Pourquoi l'équilibre est sensible.} Réseau de $80\,\%$ de \nt{GLUT} et $20\,\%$ de \nt{GABA}, tous connectés à tous, poids $J_E/N$ et $-J_I/N$; $\tau\dot r=-r+Wr+h$. Pour $J_E=10$, la seule valeur propre non nulle de $W$ vaut $0.5$. Trouvez le rapport des courants inhibiteur et excitateur, et de combien de pour cent affaiblir l'inhibition pour déclencher l'avalanche.
\end{exercise}

\begin{exercise}[\lvlC]
\label{ex:01:7}
\emph{Recherche: \nt{SERO} est-il $\gamma$?} D'après l'exercice~\ref{ex:01:3}, la réponse des neurones \nt{DOPA} à la lampe vaut $re^{-T/\tau_\gamma}$. Délais $T=1,\dots,5$~s, $n$ présentations chacun, $\ln\delta$ mesuré avec une dispersion de $0.5$. Quel $n$ faut-il pour distinguer $\tau_\gamma=2$~s de $2.4$~s (niveau $0.05$, puissance $0.8$)? Quel mécanisme, autre que $\gamma$, donnerait la même décroissance avec $T$?
\end{exercise}

\chapter{Ce que calculent les neurones \nt{GLUT}}
\label{sec:glutcomputer}

\section{Les règles du jeu}

Au chapitre~\ref{sec:neurontypes}, nous avons vu que l'écrasante majorité des neurones du cerveau sont \nt{GLUT}: ils ne font qu'exciter, leurs poids ne sont que positifs. Prenons autant de ces neurones que nous voulons et demandons-nous: que peut calculer un tel réseau si l'on ne s'autorise que ce qui existe dans un organisme vivant?

Or un organisme vivant n'a pas de générateur d'horloge qui ferait basculer tous les éléments en même temps. Chaque neurone travaille pour son compte, en temps continu: les impulsions arrivent et repartent quand elles peuvent, avec des retards variés. Il y a des \emph{entrées}, les neurones sensoriels qui réagissent à la lumière, au son ou à la pression, et des \emph{sorties}, les neurones qui commandent les muscles. Entre les deux se trouve le réseau, et personne ne lui dit quand commencer ni quand terminer le calcul. Pour un électronicien, c'est un circuit \emph{asynchrone} dont les éléments ont des retards mal connus à l'avance.

Prenons le modèle de neurone le plus simple parmi ceux qui fonctionnent honnêtement en temps continu. Au repos, la tension $V$ sur la membrane du neurone est nulle. Chaque impulsion reçue du neurone $j$ la fait monter d'un saut $w_j\ge0$, après quoi la tension redescend d'elle-même vers zéro avec la constante de temps $\tau$:
\begin{empheq}[box=\keybox]{equation}
\tau\,\frac{\dd V}{\dd t} \;=\; -V \;+\; \tau\sum_j w_j \sum_k \delta\bigl(t-t_j^{(k)}-d_j\bigr), \qquad w_j\ge 0 .
\label{eq:glutneuron}
\end{empheq}
Ici $t_j^{(k)}$ sont les instants des impulsions du neurone $j$, $d_j$ le retard sur le chemin qui en vient, et $\delta$ la fonction delta, qui traduit le saut instantané. Quand $V$ atteint le seuil $\theta$, le neurone émet une impulsion, $V$ est remis à zéro, et pendant environ une milliseconde le neurone ne peut pas se déclencher de nouveau. Valeurs typiques: $\tau$ va d'une fraction de milliseconde à vingt millisecondes selon les neurones, les retards d'une fraction de milliseconde à quelques dizaines de millisecondes. C'est le \emph{neurone intégrateur à fuite}: un circuit RC avec un seuil. C'est une simplification délibérée: chez les neurones du cortex, les branches d'entrée émettent elles-mêmes des impulsions locales~\cite{Smith2013}, et un neurone se comporte plutôt comme un petit réseau multicouche (chapitre~\ref{sec:synthesis}). Pour les questions de ce chapitre, un seul circuit RC suffit.

La principale différence avec une porte logique est la fuite: le neurone ne se souvient pas d'une impulsion éternellement, mais pendant environ $\tau$, et seules s'additionnent les impulsions arrivées dans cette fenêtre.

\section{OU et ET}

Commençons par les portes. Si une seule impulsion porte la tension au-dessus du seuil, $w\ge\theta$, le neurone se déclenche à chaque impulsion de n'importe quelle entrée. C'est un \emph{OU}.

Le ET est plus intéressant. Supposons $w<\theta<2w$: une impulsion ne suffit pas, il en faut deux. Mais la question \og sont-elles arrivées toutes les deux?\fg{} n'a pas de sens tant qu'on n'a pas dit \emph{dans quel délai}. Supposons que la première impulsion arrive à l'instant $0$, et la seconde au bout d'un temps $\Delta$. À l'arrivée de la seconde, il reste de la première $we^{-\Delta/\tau}$, et le neurone se déclenche si $we^{-\Delta/\tau}+w\ge\theta$. D'où la fenêtre de coïncidence:
\begin{empheq}[box=\keybox]{equation}
\Delta \;\le\; \Delta_{\max} \;=\; \tau\,\ln\frac{w}{\theta-w} .
\label{eq:window}
\end{empheq}
Ce n'est pas un ET logique, mais un \emph{détecteur de coïncidences}: un ET dans le temps (fig.~\ref{fig:coincidence}). Pour $\theta=1.5w$, la fenêtre vaut $\tau\ln2\approx0.7\tau$. Pour un neurone du cortex avec $\tau=10\ms$, cela fait environ sept millisecondes. Chez les neurones du système auditif, où $\tau$ est inférieur à une milliseconde, la fenêtre se réduit à une fraction de milliseconde.

\begin{figure}[t]
\centering
\begin{tikzpicture}[>=Stealth,x=0.9cm,y=1.6cm]
\foreach \dx/\lab/\c [count=\i] in {0.5/{coïncidence}/red!70!black, 2.6/{pas de coïncidence}/blue!60!black}{
 \begin{scope}[xshift={(\i-1)*7.2cm}]
  \draw[->,gray!70!black] (0,0) -- (6.3,0) node[below left,black] {\small $t$};
  \draw[->,gray!70!black] (0,0) -- (0,1.75) node[left,black] {\small $V$};
  \draw[densely dashed,gray!60!black] (0,1.5) -- (6.0,1.5) node[right,black] {\small $\theta$};
  \draw[very thick,\c] (0,0) -- (1,0) -- (1,1)
      plot[domain=1:1+\dx,samples=30] (\x,{exp(-(\x-1)/1.5)});
  \pgfmathsetmacro{\v}{exp(-\dx/1.5)+1}
  \draw[very thick,\c] (1+\dx,{exp(-\dx/1.5)}) -- (1+\dx,{min(\v,1.5)});
  \ifdim\v pt<1.5pt
    \draw[very thick,\c] plot[domain=1+\dx:6,samples=40] (\x,{\v*exp(-(\x-1-\dx)/1.5)});
  \else
    \draw[very thick,\c] (1+\dx,1.5) -- (1+\dx,0) -- (6,0);
    \draw[->,thick,\c] (1+\dx,1.55) -- (1+\dx,1.72) node[above,font=\scriptsize] {impulsion};
  \fi
  \draw[->,thick] (1,-0.35) -- (1,-0.08); \draw[->,thick] (1+\dx,-0.35) -- (1+\dx,-0.08);
  \draw[decorate,decoration={brace,amplitude=3pt,mirror}] (1,-0.42) -- (1+\dx,-0.42) node[midway,below=3pt] {\scriptsize $\Delta$};
  \node[\c] at (3.5,-0.95) {\small \lab};
 \end{scope}}
\end{tikzpicture}
\caption{Détecteur de coïncidences, seuil $\theta=1.5w$. À gauche, la seconde impulsion est arrivée au bout de $\Delta<\Delta_{\max}$, la somme a atteint le seuil, le neurone s'est déclenché. À droite, $\Delta>\Delta_{\max}$: il ne reste presque rien de la première impulsion, et le neurone est resté muet.}
\label{fig:coincidence}
\end{figure}

Une autre façon d'obtenir un ET passe par la durée plutôt que par l'instant précis. Tant que la lumière est allumée, un neurone sensoriel émet des impulsions fréquentes et régulières; si une telle entrée ne suffit pas au seuil mais que deux suffisent, le neurone travaille tant que les deux sont actives. Le réseau calcule alors par \emph{niveaux}, et l'instant précis d'arrivée des impulsions n'a pas d'importance. La plus grande partie de ce qui suit fonctionnera ainsi.

\section{Ce qu'on ne peut pas faire}

Essayons maintenant de faire un NON: un neurone qui travaille quand l'entrée se tait. Impossible: sans impulsions d'entrée, il n'y a dans~\eqref{eq:glutneuron} aucun saut, et la tension reste à zéro, sous le seuil. Et un réseau de nombreux neurones? Non plus, et cela se démontre d'un coup pour tous les réseaux.

Le plus commode est de raisonner en fréquences. Soit $r_i$ la fréquence des impulsions du neurone $i$, $I_i$ l'apport des entrées, et $f_i$ sa caractéristique de transfert: la façon dont la fréquence dépend de l'apport total. Pour un neurone intégrateur, cette caractéristique est non décroissante: plus d'apport, plus d'impulsions. Les fréquences d'un réseau parvenu à l'équilibre vérifient les équations
\begin{equation}
r_i \;=\; f_i\Bigl(\sum_j w_{ij}\,r_j + I_i\Bigr), \qquad w_{ij}\ge 0 .
\label{eq:ratenet}
\end{equation}

\begin{keyblock}
\begin{proposition}[Monotonie]
\label{prop:monotone}
Supposons que le réseau~\eqref{eq:ratenet} ne soit fait que de neurones \nt{GLUT} et parte du silence complet. Si l'on renforce les entrées, la fréquence établie d'aucun neurone ne diminue. Tout ce que calcule un tel réseau, ce sont des fonctions \emph{monotones} des entrées.
\end{proposition}
\end{keyblock}

Démonstration. Partons du silence, $r^{(0)}=0$, et substituons les fréquences dans le membre de droite de~\eqref{eq:ratenet}: $r^{(k+1)}=F\bigl(r^{(k)}\bigr)$. Le membre de droite $F$ est non décroissant en chacun de ses arguments, puisque les poids sont positifs ou nuls et que les $f_i$ sont non décroissantes. Comme $r^{(1)}\ge0=r^{(0)}$, par récurrence $r^{(k+1)}\ge r^{(k)}$: les fréquences ne font que croître et atteignent la plus petite solution de~\eqref{eq:ratenet}. Si l'on remplace les entrées $I$ par des entrées plus grandes $I'$, alors à chaque pas $r^{(k)}(I)\le r^{(k)}(I')$, et à la limite aussi. Un vrai réseau atteint l'équilibre non par pas, mais continûment; pour lui, la conclusion est la même: avec des connexions positives, une inégalité entre deux exécutions, vraie au départ, reste vraie tout le temps.

Pour des impulsions isolées, l'énoncé n'est pas vrai à la lettre: une impulsion d'entrée supplémentaire peut faire se déclencher le neurone un peu plus tôt, et à cause de la milliseconde réfractaire son impulsion suivante disparaît. Mais l'effet est faible et peu fiable, on ne peut pas calculer avec. Pour les fréquences, la monotonie est stricte.

Les conséquences sont celles de tout circuit sans inverseurs. Pas de NON. Pas de OU exclusif: ajouter une seconde entrée ne peut pas éteindre la sortie. Et le plus gênant: \emph{on ne peut pas arrêter une activité par une activité}; on peut seulement retirer ce qui l'entretient.

\section{Deux fils: ON et OFF}

L'organisme lui-même indique la sortie de l'impasse. Dans beaucoup d'organes des sens, un stimulus est transmis par deux ensembles de neurones: dans la vision, certaines cellules répondent à l'augmentation de la lumière, d'autres à sa diminution~\cite{Schiller1992}. Appelons-les neurones \emph{d'allumage} (ON) et \emph{d'extinction} (OFF). Au lieu d'un fil $x$, le réseau reçoit une paire $(x,\bar x)$, et l'absence qu'il ne sait pas calculer lui est signalée par le capteur lui-même.

Avec une paire de fils, le NON est gratuit: il suffit d'échanger les fils. Le ET et le OU se font chacun avec deux neurones:
\begin{empheq}[box=\keybox]{equation}
\begin{aligned}
x\wedge y \;&\to\; \bigl(\;x\wedge y,\;\; \bar x\vee\bar y\;\bigr),\\
x\vee y \;&\to\; \bigl(\;x\vee y,\;\; \bar x\wedge\bar y\;\bigr).
\end{aligned}
\label{eq:dualrail}
\end{empheq}
À droite, toutes les opérations sont monotones, donc la proposition~\ref{prop:monotone} n'est pas violée.

Pour un circuit asynchrone, le code à deux fils (le codage double rail des électroniciens) a un second avantage, plus important encore. Une paire de fils a trois états: \og oui\fg, \og non\fg{} et, quand les deux se taisent, \emph{\og pas encore de réponse\fg}. Le destinataire apprend que le calcul est terminé non par une horloge, mais par le signal lui-même: dans chaque paire, exactement un fil s'est activé. Entre deux stimuli, les capteurs se taisent, et le réseau revient au silence, prêt pour la fois suivante. En électronique asynchrone, on construit sur ce principe des circuits qui fonctionnent correctement quels que soient les retards des éléments~\cite{Sparso2001}.

\begin{keyblock}
\begin{proposition}[Calcul asynchrone]
\label{prop:dualrail}
Si chaque entrée arrive sous forme d'une paire de fils ON/OFF, un réseau de neurones \nt{GLUT} peut calculer n'importe quelle fonction logique des entrées. La réponse arrive elle aussi par paires de fils, et on voit qu'elle est prête quand un fil de chaque paire s'est activé. Aucune horloge n'est nécessaire. Le prix: environ deux fois plus de neurones.
\end{proposition}
\end{keyblock}

Exemple: le OU exclusif, \og exactement un des deux stimuli a changé\fg. Le fil direct de la réponse est $(x\wedge\bar y)\vee(\bar x\wedge y)$, le fil inverse $(x\wedge y)\vee(\bar x\wedge\bar y)$. Cela fait six neurones, dont aucun n'a besoin d'inhibition (fig.~\ref{fig:dualxor}).

\begin{figure}[t]
\centering
\begin{tikzpicture}[>=Stealth,x=1cm,y=1cm,
  nrn/.style={circle,draw,very thick,fill=blue!6,minimum size=0.75cm,inner sep=0pt,font=\small},
  inp/.style={font=\small}]
\node[inp] (X)  at (0,3.3) {$x$};
\node[inp] (Xb) at (0,2.2) {$\bar x$};
\node[inp] (Y)  at (0,1.1) {$y$};
\node[inp] (Yb) at (0,0)   {$\bar y$};
\node[nrn] (A1) at (3.2,3.3) {$2$};
\node[nrn] (A2) at (3.2,2.2) {$2$};
\node[nrn] (A3) at (3.2,1.1) {$2$};
\node[nrn] (A4) at (3.2,0)   {$2$};
\node[nrn] (O)  at (7.2,2.75) {$1$};
\node[nrn] (Ob) at (7.2,0.55) {$1$};
\foreach \s/\t in {X/A1,Yb/A1,Xb/A2,Y/A2,X/A3,Y/A3,Xb/A4,Yb/A4}{\draw[->,thick,blue!60!black] (\s) -- (\t);}
\draw[->,thick,blue!60!black] (A1) -- node[pos=0.45,above,sloped,font=\scriptsize,black] {$x\wedge\bar y$} (O);
\draw[->,thick,blue!60!black] (A2) -- node[pos=0.45,below,sloped,font=\scriptsize,black] {$\bar x\wedge y$} (O);
\draw[->,thick,blue!60!black] (A3) -- node[pos=0.45,above,sloped,font=\scriptsize,black] {$x\wedge y$} (Ob);
\draw[->,thick,blue!60!black] (A4) -- node[pos=0.45,below,sloped,font=\scriptsize,black] {$\bar x\wedge\bar y$} (Ob);
\draw[->,very thick,red!70!black] (O) -- (8.6,2.75) node[right,black] {$x\oplus y$: \og oui\fg};
\draw[->,very thick,red!70!black] (Ob) -- (8.6,0.55) node[right,black] {$\overline{x\oplus y}$: \og non\fg};
\node[below,font=\small,gray!40!black] at (3.2,-0.55) {ET: seuil $2$};
\node[below,font=\small,gray!40!black] at (7.2,-0.55) {OU: seuil $1$};
\node[below,font=\small,gray!40!black] at (0,-0.55) {paires de fils};
\end{tikzpicture}
\caption{OU exclusif en code à deux fils, avec des neurones \nt{GLUT} seulement. Le nombre dans le cercle est le seuil, en unités du poids d'une entrée. Quatre neurones ET reconnaissent les quatre combinaisons d'entrées, deux neurones OU les rassemblent en une paire de fils de réponse.}
\label{fig:dualxor}
\end{figure}

\section{Des retards au lieu d'une horloge}

Pour calculer avec le temps, il suffit de retards dans les fils et de détecteurs de coïncidences. Le meilleur exemple est la façon dont le cerveau détermine d'où vient un son.

Le son d'une source placée sur le côté arrive à l'oreille proche plus tôt qu'à l'oreille éloignée. La différence dépend de la direction: de zéro quand la source est droit devant, jusqu'à $0.22\,\text{m}/343\,\text{m/s}\approx0.6\ms$ chez l'homme quand elle est exactement sur le côté. Le cerveau distingue des différences d'environ dix \emph{microsecondes}~\cite{KlumppEady1956,Grothe2010}, alors que les neurones eux-mêmes ne se déclenchent pas avec une précision meilleure qu'une fraction de milliseconde. Comment?

Menons depuis l'oreille gauche un fil qui longe une rangée de détecteurs de coïncidences de gauche à droite, et depuis l'oreille droite un fil de droite à gauche (fig.~\ref{fig:itd}). Le signal parcourt le fil à la vitesse $v$. Jusqu'au détecteur placé à la distance $x$ du bord gauche d'une rangée de longueur $L$, le signal gauche met le temps $x/v$, le signal droit $(L-x)/v$. Si le son arrive à l'oreille droite avec un retard $\delta t$, les deux signaux se rencontrent simultanément au point où
\begin{empheq}[box=\keybox]{equation}
\frac{x}{v} \;=\; \frac{L-x}{v} + \delta t
\quad\Longrightarrow\quad
x \;=\; \frac{L}{2} + \frac{v\,\delta t}{2} .
\label{eq:itd}
\end{empheq}
Seul se déclenche le détecteur auquel les deux signaux sont arrivés dans la fenêtre~\eqref{eq:window}. Le numéro du détecteur déclenché donne la direction de la source. Une différence de temps s'est changée en \emph{position}.

\begin{figure}[t]
\centering
\begin{tikzpicture}[>=Stealth,
  nrn/.style={circle,draw,very thick,fill=blue!6,minimum size=0.7cm,inner sep=0pt}]
\foreach \k in {0,...,6}{\node[nrn] (D\k) at (1.4*\k,0) {\scriptsize $2$};}
\node[nrn,fill=red!15] at (D4) {\scriptsize $2$};
\draw[->,thick,blue!60!black] (-1.4,0.9) node[left,black] {\small oreille gauche} -- (9.2,0.9);
\draw[->,thick,orange!85!black] (9.8,-0.9) node[right,black] {\small oreille droite} -- (-0.8,-0.9);
\foreach \k in {0,...,6}{
  \draw[->,blue!60!black] (1.4*\k,0.9) -- (D\k.north);
  \draw[->,orange!85!black] (1.4*\k,-0.9) -- (D\k.south);}
\draw[->,thick,red!70!black] (D4.north east) ++(0.1,0.05) -- ++(0.5,0.45) node[above right,font=\small,black] {déclenché};
\node[font=\small,gray!40!black] at (4.2,-1.5) {le retard croît $\longrightarrow$ \hspace{2.2cm} $\longleftarrow$ le retard croît};
\pic[scale=1.4] at (-2.3,1.75) {speaker};
\node[font=\scriptsize,align=center] at (-2.3,2.35) {source à gauche};
\end{tikzpicture}
\caption{Déterminer la direction d'un son. Les signaux des deux oreilles courent l'un vers l'autre; chaque cercle est un détecteur de coïncidences de seuil $2$. Si le son arrive plus tard à l'oreille droite, les signaux se rencontrent à droite du milieu, selon la formule~\eqref{eq:itd}. La sortie du réseau est le numéro du neurone déclenché. Ici la source est à gauche: le son est arrivé plus tôt à l'oreille gauche.}
\label{fig:itd}
\end{figure}

Il n'y a ici ni horloge, ni inhibition, ni même de code à deux fils: le réseau ne répond pas \og oui\fg{} ou \og non\fg, il désigne un neurone parmi beaucoup. Un tel circuit de lignes à retard et de détecteurs de coïncidences fonctionne apparemment bel et bien dans le tronc cérébral auditif des oiseaux~\cite{Carr1990}. La précision de l'ordre de la microseconde ne vient pas de la précision de chaque neurone, mais de ce que les détecteurs sont nombreux et que chacun surveille son propre retard. Chez les mammifères, on n'a pas trouvé de telle rangée de retards: la même tâche y est apparemment résolue par des détecteurs de coïncidences réglés par une inhibition précisément calée dans le temps, venue de neurones \nt{GLYC}~\cite{Brand2002,Grothe2010}. Comment l'inhibition rétrécit la fenêtre de coïncidence, nous le verrons au chapitre~\ref{sec:gaba} sur l'exemple de \nt{GABA}; la glycine inhibe de la même façon, en ouvrant des canaux chlorure chez le destinataire.

\section{La mémoire}

Un ordinateur a besoin de mémoire. Dans un tel réseau, la mémoire est une activité qui s'entretient elle-même. Prenons un groupe de neurones \nt{GLUT} fortement connectés entre eux et décrivons-le par une seule fréquence $r$ (fig.~\ref{fig:recurrentgroup}a). À l'équilibre, $r=f(wr+I)$, où $w$ est la force de la rétroaction du groupe sur lui-même, et $I$ l'apport extérieur. Résolvons cette équation graphiquement, en coupant la courbe $f(wr+I)$ par la droite $r$ (fig.~\ref{fig:bistable}).

\begin{figure}[t]
\centering
% (b): tau dr/dt = -r + f(w r + I), f as in fig:bistable, w=1, I=0.6; Euler dt=0.001 tau.
\begin{tikzpicture}[>=Stealth,
  nrn/.style={circle,draw,thick,fill=blue!6,minimum size=0.5cm,inner sep=0pt}]
\begin{scope}
\node[font=\small] at (-0.5,1.55) {(a)};
\foreach \k in {1,...,6}{\node[nrn] (n\k) at ({1.4+1.0*cos(60*\k)},{1.0*sin(60*\k)}) {};}
\foreach \a/\b in {1/2,2/3,3/4,4/5,5/6,6/1,1/4,2/5,3/6}{\draw[<->,blue!60!black] (n\a) -- (n\b);}
\foreach \k in {2,3,4}{\draw[->,thick,green!45!black] ($(n\k)+(-0.95,0)$) -- (n\k);}
\node[font=\small,green!45!black] at (-0.35,-0.45) {$I$};
\node[font=\Large] at (3.1,0) {$\approx$};
\node[nrn,minimum size=0.85cm,font=\small] (R) at (4.35,-0.1) {$r$};
\draw[->,thick,blue!60!black] (R) edge[loop above,looseness=6] node[above,font=\small] {$w$} (R);
\draw[->,thick,green!45!black] (3.55,-0.95) node[left,font=\small] {$I$} -- (R);
\end{scope}
\begin{scope}[xshift=6.3cm,y=2.6cm,x=1.05cm]
\node[font=\small] at (-0.6,1.25) {(b)};
\draw[->,gray!70!black] (0,0) -- (7.3,0) node[below left,font=\small,black] {$t/\tau$};
\draw[->,gray!70!black] (0,0) -- (0,1.12) node[left,font=\small,black] {$r$};
\foreach \t in {0,2,4,6}{\draw[gray!60!black] (\t,-0.02) -- (\t,0.02); \node[below,font=\scriptsize] at (\t,0) {\t};}
\draw[densely dotted,gray!60!black] (0,0.5) -- (7,0.5) node[right,font=\scriptsize,black,align=left] {seuil\\d'écriture};
\fill[green!45!black,opacity=0.25] (1,-0.2) rectangle (2,-0.13);
\fill[gray!60!black,opacity=0.35] (1,-0.29) rectangle (1.5,-0.22);
\draw[very thick,blue!60!black] plot coordinates {(0.0,0.002) (0.1,0.002) (0.2,0.002) (0.3,0.002) (0.4,0.002) (0.5,0.002) (0.6,0.002) (0.7,0.002) (0.8,0.002) (0.9,0.002) (1.0,0.002) (1.1,0.083) (1.2,0.165) (1.3,0.242) (1.4,0.313) (1.5,0.378) (1.6,0.437) (1.7,0.491) (1.8,0.539) (1.9,0.583) (2.0,0.623) (2.1,0.643) (2.2,0.665) (2.3,0.688) (2.4,0.71) (2.5,0.732) (2.6,0.753) (2.7,0.773) (2.8,0.792) (2.9,0.81) (3.0,0.826) (3.2,0.855) (3.4,0.88) (3.6,0.9) (3.8,0.917) (4.0,0.931) (4.4,0.953) (4.8,0.967) (5.2,0.977) (5.6,0.984) (6.0,0.988) (6.5,0.992) (7.0,0.994)};
\draw[very thick,gray!60!black,densely dashed] plot coordinates {(0.0,0.002) (0.1,0.002) (0.2,0.002) (0.3,0.002) (0.4,0.002) (0.5,0.002) (0.6,0.002) (0.7,0.002) (0.8,0.002) (0.9,0.002) (1.0,0.002) (1.1,0.083) (1.2,0.165) (1.3,0.242) (1.4,0.313) (1.5,0.378) (1.6,0.358) (1.7,0.336) (1.8,0.313) (1.9,0.291) (2.0,0.269) (2.2,0.228) (2.4,0.191) (2.6,0.159) (2.8,0.133) (3.0,0.11) (3.2,0.091) (3.4,0.076) (3.6,0.063) (3.8,0.052) (4.0,0.043) (4.4,0.03) (4.8,0.021) (5.2,0.015) (5.6,0.011) (6.0,0.008) (6.5,0.006) (7.0,0.004)};
\node[font=\scriptsize,blue!60!black,anchor=west] at (3.3,0.8) {apport $1\tau$: écrit};
\node[font=\scriptsize,gray!50!black,anchor=west] at (2.3,0.3) {apport $0.5\tau$: éteint};
\end{scope}
\end{tikzpicture}
\caption{Un groupe fortement connecté à lui-même. (a)~De nombreux neurones \nt{GLUT} qui s'excitent mutuellement se comportent comme un seul neurone de fréquence $r$ qui s'excite lui-même avec la force $w$; de l'extérieur arrive l'apport $I$. (b)~Fréquence du groupe au cours du temps pour $w=1$ et la même courbe $f$ que sur la fig.~\ref{fig:bistable}. L'apport $I=0.6$ est appliqué pendant la durée marquée par la bande sous l'axe. S'il dure $1\tau$, le groupe a le temps de franchir le point instable $r=0.5$, s'embrase et reste allumé une fois l'apport terminé. S'il ne dure que $0.5\tau$, le groupe n'atteint pas ce point et s'éteint: pour écrire un bit, l'apport doit durer plus d'environ $0.7\tau$.}
\label{fig:recurrentgroup}
\end{figure}

\begin{figure}[t]
\centering
\begin{tikzpicture}[>=Stealth,x=5cm,y=3.2cm]
\draw[->,gray!70!black] (0,0) -- (1.1,0) node[below left,black] {\small $r$};
\draw[->,gray!70!black] (0,0) -- (0,1.1);
\draw[thick,gray!60!black] (0,0) -- (1.05,1.05) node[right,black] {\small $r$};
\draw[very thick,blue!60!black] plot[domain=0:1.05,samples=80] (\x,{1/(1+exp(-(\x-0.5)/0.08))});
\draw[very thick,red!70!black,densely dashed] plot[domain=0:1.05,samples=80] (\x,{1/(1+exp(-(\x+0.3-0.5)/0.08))});
\fill[blue!60!black] (0.002,0.002) circle (0.018);
\draw[blue!60!black,fill=white,thick] (0.5,0.5) circle (0.018);
\fill[blue!60!black] (0.998,0.998) circle (0.018);
\node[below right,font=\small] at (0.02,0.0) {éteint};
\node[right,font=\small] at (0.52,0.46) {instable};
\node[below,font=\small] at (0.99,0.96) {allumé};
\node[anchor=east,font=\small,red!70!black] at (0.19,0.62) {$f(wr+I)$};
\node[anchor=west,font=\small,blue!60!black] at (0.45,0.1) {$f(wr)$};
\end{tikzpicture}
\caption{Une mémoire faite d'un groupe de neurones \nt{GLUT} à forte rétroaction. Courbe pleine: sans apport extérieur, deux intersections stables avec la droite $r$, \og éteint\fg{} et \og allumé\fg, et une instable entre les deux. Un bref apport $I$ (courbe en tirets) ne laisse que l'intersection supérieure.}
\label{fig:bistable}
\end{figure}

Si la rétroaction est forte, il y a trois intersections: celle du bas est le silence, celle du haut le groupe qui brûle à plein régime, celle du milieu est instable. On obtient un élément à deux états stables, une bascule. Y écrire un \og 1\fg{} est facile: un bref apport relève la courbe, l'intersection du bas disparaît, et le groupe s'embrase, puis reste allumé une fois l'apport terminé (fig.~\ref{fig:recurrentgroup}b). Un apport trop bref n'a pas le temps d'amener le groupe au point instable, et le groupe s'éteint. Une telle activité auto-entretenue, qui persiste des secondes après la disparition du stimulus, s'observe bel et bien dans le cerveau et passe pour la base de la mémoire à court terme~\cite{Wang2001}. Mais ce n'est pas la seule façon de se souvenir pendant quelques secondes. L'information peut aussi être gardée dans une variable lente des synapses elles-mêmes: après une bouffée d'impulsions, la facilitation, comme celle de la synapse-\og trait\fg{} du chapitre~\ref{sec:code}, dure environ une seconde, et le réseau peut presque se taire entre de brèves bouffées: non plus une bascule, mais un condensateur chargé~\cite{Mongillo2008}. De tels intervalles \og silencieux\fg{} pendant le maintien en mémoire s'observent~\cite{Stokes2015}, et un stockage sans impulsions permanentes coûte moins d'énergie. Laquelle des deux façons le cortex utilise dans quel cas reste débattu.

Et comment effacer? D'après la proposition~\ref{prop:monotone}, impossible par une activité; il reste à \emph{retirer} quelque chose. Première façon: retirer le soutien; si le groupe ne brûle qu'avec l'apport d'un autre groupe, la mémoire s'éteint avec lui. Seconde façon: la fatigue. Dans les vraies synapses, la réserve de neurotransmetteur s'épuise lors d'un travail prolongé~\cite{Tsodyks1997}, et chez les neurones montent des courants lents qui réduisent l'excitabilité. La rétroaction s'affaiblit, l'intersection supérieure disparaît, et le groupe s'éteint de lui-même au bout de quelques secondes. On obtient une mémoire qui s'allume par un signal et ne s'éteint qu'avec le temps.

\section{Les séquences}

Au lieu d'une horloge, on peut avoir une \emph{minuterie déclenchée par un événement}. Relions des groupes de neurones \nt{GLUT} en chaîne: le premier excite le deuxième, le deuxième le troisième, et ainsi de suite. Un signal extérieur allume le premier groupe, et une vague d'activité parcourt la chaîne, chaque maillon un ou deux retards après le précédent, et une seule fois: juste après une impulsion, les neurones sont réfractaires, et la vague est déjà passée plus loin.

Une telle chaîne mesure le temps écoulé depuis un événement: si le maillon numéro $k$ s'est déclenché, $k$ retards se sont écoulés depuis le lancement. On peut envoyer ses sorties vers les muscles et obtenir une séquence de mouvements qui se déroule d'elle-même. Chez les oiseaux chanteurs, pendant le chant, des neurones \nt{GLUT} isolés d'une des régions du cerveau se déclenchent strictement à tour de rôle, chacun à son moment du chant, ce qui ressemble beaucoup à une telle chaîne~\cite{Hahnloser2002}.

Contrairement à une horloge, la chaîne se tait tant qu'on ne l'a pas lancée, ne s'exécute qu'une fois et ne mesure le temps que pour ceux qui y sont branchés. Le réseau n'a toujours pas de cadence commune.

\section{Ce qui manque}

Avec des neurones \nt{GLUT} seuls, sans inhibition, on obtient beaucoup. Mais trois choses manquent, et toutes trois à cause de la proposition~\ref{prop:monotone}.

\emph{Le choix.} Le réseau doit choisir une direction, une action. Si deux stimuli sont presque égaux, dans le réseau de la fig.~\ref{fig:itd} les deux sorties se déclenchent, et rien ne permet de supprimer la plus faible au profit de la plus forte.

\emph{La remise à zéro.} On ne peut pas effacer la mémoire par un signal.

\emph{La stabilité.} Dans un réseau dont les connexions ne suivent pas le plan d'un ingénieur, les boucles de rétroaction positive sont inévitables, et l'activité peut y croître en avalanche (chapitre~\ref{sec:neurontypes}). Le seul frein est cette même fatigue, lente et grossière.

Les trois maux ont un seul remède: un neurone qui éteint une impulsion par une impulsion. C'est \nt{GABA}, et on en a besoin non pour calculer, mais pour choisir, remettre à zéro et garder le calcul sous contrôle.

\section{Exercices}

\begin{exercise}[\lvlA]
\label{ex:02:1}
\emph{Une rangée de détecteurs pour le son.} Les signaux des oreilles parcourent les fils de la fig.~\ref{fig:itd} à $5$~m/s, la plus grande différence d'arrivée est de $0.64\ms$. Les détecteurs, des neurones~\eqref{eq:glutneuron} avec $\tau=0.5\ms$, $\theta=1.5w$, sont placés de sorte que deux voisins diffèrent de $10$~\textmu s. Combien de détecteurs compte la rangée, et combien se déclenchent pour un son?
\end{exercise}

\begin{exercise}[\lvlB]
\label{ex:02:2}
\emph{Des entrées inégales décalent la fenêtre.} Un détecteur~\eqref{eq:glutneuron} avec $\tau=0.5\ms$, $\theta=1.5$ reçoit une impulsion de chacune de deux entrées de poids $1$ et $0.8$. Trouvez la fenêtre de coïncidence et son centre. De combien de microsecondes se trompe la rangée de la fig.~\ref{fig:itd} si l'entrée d'une oreille est $20\,\%$ plus faible?
\end{exercise}

\begin{exercise}[\lvlB]
\label{ex:02:3}
\emph{Quels réseaux n'oscillent pas.} Le réseau $\tau\dot r_i=-r_i+f_i\bigl(\textstyle\sum_jw_{ij}r_j+I_i\bigr)$ avec des $f_i$ non décroissantes et $w_{ij}\ge0$ pour $i\ne j$ est monotone et n'oscille pas durablement. Montrez que le changement $r_i\to\pm r_i$ y ramène tout réseau à nombre pair de connexions inhibitrices dans chaque cycle. Deux groupes à inhibition mutuelle oscillent-ils? Et la boucle \nt{GLUT}--\nt{GABA}?
\end{exercise}

\begin{exercise}[\lvlB]
\label{ex:02:4}
\emph{Conception: additionneur à deux fils.} Chaque bit est une paire de fils $(x,\bar x)$. Construisez avec des neurones \nt{GLUT} un additionneur complet de trois bits donnant les paires $(S,\bar S)$ de la somme et $(C,\bar C)$ de la retenue, avec poids et seuils, en huit neurones au plus. Combien en faudrait-il avec des ET et des OU, comme sur la fig.~\ref{fig:dualxor}?
\end{exercise}

\begin{exercise}[\lvlB]
\label{ex:02:5}
\emph{Modélisation: combien dure une écriture.} Groupe de la fig.~\ref{fig:recurrentgroup}: $\tau\dot r=-r+f(r+I)$, $f(u)=1/\bigl(1+e^{-(u-0.5)/0.08}\bigr)$, dans l'état bas avant l'écriture; l'apport $I$ est appliqué pendant le temps $T$. Trouvez numériquement le plus petit $T$ qui écrit le bit pour $I=0.6$, et le seuil $I_c$ en dessous duquel le bit ne s'écrit pour aucun $T$.
\end{exercise}

\begin{exercise}[\lvlA]
\label{ex:02:6}
\emph{La chaîne comme minuterie.} Un maillon: $100$ neurones \nt{GLUT}, retard $2\ms$, dispersion indépendante $0.2\ms$. (a)~Combien de neurones pour mesurer $1$~s, et quelle erreur relative? Comment dépend-elle de l'intervalle? (b)~Une vraie minuterie a une erreur de $5\,\%$ quel que soit l'intervalle. Quelle source de dispersion donne cela?
\end{exercise}

\begin{exercise}[\lvlC]
\label{ex:02:7}
\emph{Recherche: bascule ou condensateur.} $100$ neurones gardent un bit $10$~s. Bascule: $20$ impulsions par seconde chacun. Condensateur: la facilitation $u$ décroît avec $\tau_F=1$~s, le bit se lit si $u\ge u_{\max}/2$, on rafraîchit par une bouffée de trois impulsions de chacun. Combien de fois le condensateur est-il plus économe? Et le bit, si le groupe se tait $200\ms$?
\end{exercise}

\chapter{\nt{GLUT} et \nt{GABA} ensemble}
\label{sec:gaba}
\begin{chaptergoals}
\item comment un veto $a\wedge\bar b$ rend \nt{GLUT} et \nt{GABA} complets avec un seul fil par bit;
\item pourquoi un shunt divise sans soustraire, et comment l'inhibition réduit la fenêtre de coïncidence;
\item comment l'inhibition mutuelle à $\beta>1$ élit un gagnant, et comment un signal efface la mémoire;
\item quand une inhibition rapide garde le réseau stable, et quand une inhibition lente crée un rythme.
\end{chaptergoals}

\section{Les règles du jeu}

Un réseau fait seulement de neurones \nt{GLUT} ne sait ni choisir, ni effacer sa mémoire, ni empêcher l'activité de partir en avalanche, et tout cela à cause de la monotonie (proposition~\ref{prop:monotone}). Ajoutons-y le deuxième type de neurones par le nombre, \nt{GABA}. Les règles restent les mêmes: seulement ce qui existe dans un organisme vivant, et pas de signal d'horloge.

Le modèle de neurone est le même que~\eqref{eq:glutneuron}, mais une impulsion venue d'un neurone \nt{GABA} fait baisser la tension du destinataire au lieu de la faire monter. Les poids ont maintenant deux signes, et c'est l'expéditeur qui fixe le signe, règle~\eqref{eq:dalesign}.

Quelques paramètres du circuit. Les neurones \nt{GABA} représentent d'un cinquième à un quart des neurones du cortex~\cite{Hendry1987}. Leurs sorties sont courtes: ils inhibent leurs voisins, pas des régions lointaines. L'inhibition arrive au bout d'une milliseconde et dure quelques millisecondes. Sans entrée, un neurone \nt{GABA} se tait: pour inhiber quelqu'un, il doit lui-même recevoir une excitation, en général de neurones \nt{GLUT} du même réseau.

Une connexion négative d'un neurone \nt{GLUT} $A$ vers un neurone $B$ doit donc passer par un intermédiaire: $A$ excite un neurone \nt{GABA}, qui inhibe $B$. Le changement de signe coûte un neurone et un retard, environ une milliseconde.

Pour l'inhibition, il importe non seulement \emph{de qui} vient la connexion, mais aussi \emph{où} elle se pose: le site d'arrivée détermine en grande partie ce que fait la connexion~\cite{Markram2004}. Les sorties \nt{GLUT} se posent surtout sur les branches d'entrée du destinataire, ses \emph{dendrites}, et portent des données: chacune de ces entrées est un terme de la somme. Une sortie posée sur le corps du neurone agit sur la somme tout entière: c'est ainsi que beaucoup de neurones \nt{GABA} rapides divisent la réponse du destinataire et fixent le moment où il se déclenche. Une sortie posée à l'endroit où naît l'impulsion du destinataire a le dernier mot: c'est un veto sur toute la sortie à la fois. Et une sortie posée sur la terminaison du fil d'un autre neurone change la probabilité de libération $p$ d'une seule ligne d'entrée, sans toucher aux autres~\cite{Rudomin1999}; c'est ainsi, en particulier, que fonctionne le portillon de la douleur du chapitre~\ref{sec:pain}. Hors du cerveau, les sorties se posent sur les fibres musculaires (chapitre~\ref{sec:muscles}) et sur les organes (chapitre~\ref{sec:chemoutput}), et certaines ne se posent nulle part et libèrent leur neurotransmetteur dans le volume du tissu ou dans le sang (chapitres~\ref{sec:serocomputer} et~\ref{sec:chemoutput}).

\section{NON et veto}

Peut-on maintenant faire un NON? Presque. Un neurone qui travaille quand l'entrée se tait doit tirer son excitation de quelque part. Dans le cerveau vivant, elle existe: chaque neurone reçoit en permanence une activité de fond venue de milliers d'autres. Supposons que le fond maintienne le neurone $Y$ en activité, et que l'entrée $x$ l'inhibe par un intermédiaire \nt{GABA}. C'est un NON.

Plus utile, cependant, est le \emph{veto}: \og $a$, mais pas si $b$\fg:
\begin{empheq}[box=\keybox]{equation}
y \;=\; a \wedge \bar b .
\label{eq:veto}
\end{empheq}
Le neurone $Y$ reçoit une excitation de $a$ et une inhibition du neurone \nt{GABA} qu'excite $b$. Le fond n'est pas nécessaire ici: l'excitation vient de $a$ lui-même. Avec des vetos et des OU, on construit tout. Par exemple, le OU exclusif: $x\oplus y=(x\wedge\bar y)\vee(\bar x\wedge y)$, soit deux neurones-vetos, deux intermédiaires \nt{GABA} et un neurone OU.

\begin{keyblock}
\begin{proposition}[Complétude de \nt{GLUT}+\nt{GABA}]
\label{prop:gabacomplete}
Un réseau de neurones \nt{GLUT} et \nt{GABA} peut calculer n'importe quelle fonction logique des entrées, et chaque bit est transmis par un seul fil. Les capteurs ON/OFF du chapitre~\ref{sec:glutcomputer} ne sont plus nécessaires. Si la fonction doit valoir un quand toutes les entrées se taisent, il faut une source d'activité de fond.
\end{proposition}
\end{keyblock}

Démonstration: le veto~\eqref{eq:veto} associé au OU est fonctionnellement complet si l'on dispose d'une constante égale à un, car NON $x$ est le veto \og un, mais pas si $x$\fg. Et les fonctions qui valent zéro quand toutes les entrées se taisent se construisent avec des vetos et des OU, même sans constante.

\section{La direction du mouvement}

Un veto dans le temps, comme le ET dans le temps du chapitre~\ref{sec:glutcomputer}, donne un détecteur de direction du mouvement.

Soient deux capteurs voisins, $A$ et $B$, séparés d'une distance $\Delta x$. Le neurone de sortie $Y$ reçoit une excitation de $B$ et une inhibition de $A$ par un intermédiaire \nt{GABA}, avec un retard $d$ (fig.~\ref{fig:direction}). Une tache de lumière qui court à la vitesse $v$ de $A$ vers $B$ passe sur $A$ à l'instant $0$, et l'inhibition atteint $Y$ à l'instant $d$; elle passe sur $B$ à l'instant $\Delta x/v$, et c'est alors qu'arrive l'excitation. Si ces instants coïncident à la fenêtre~\eqref{eq:window} près, l'inhibition éteint l'excitation, et $Y$ se tait. Cela se produit pour une vitesse d'environ
\begin{empheq}[box=\keybox]{equation}
v^* \;=\; \frac{\Delta x}{d} .
\label{eq:vstar}
\end{empheq}
Pour un mouvement de $B$ vers $A$, l'excitation venue de $B$ arrive à l'instant $0$, et l'inhibition venue de $A$ seulement à l'instant $\Delta x/v+d$, quand $Y$ s'est déjà déclenché.

\begin{figure}[t]
\centering
\begin{tikzpicture}[>=Stealth,
  nrn/.style={circle,draw,very thick,fill=blue!6,minimum size=0.9cm,inner sep=0pt},
  gab/.style={nrn,fill=red!10},
  inh/.style={-{Bar[width=7pt]},thick,red!70!black}]
\node[nrn] (A) at (0,2) {$A$};
\node[nrn] (B) at (3,2) {$B$};
\node[gab] (G) at (0.9,0.6) {\small \nt{GABA}};
\node[nrn] (Y) at (3,-0.6) {$Y$};
\draw[->,thick] (A) -- (G);
\draw[inh] (G) -- node[below left=-2pt] {\scriptsize retard $d$} (Y);
\draw[->,thick] (B) -- (Y);
\draw[->,very thick,red!70!black] (Y) -- (4.6,-0.6) node[right,black] {\small sortie};
\draw[->,thick,orange!85!black] (-0.6,2.9) -- (3.6,2.9) node[right,black] {\small se tait};
\draw[<-,thick,orange!85!black] (-0.6,3.4) -- (3.6,3.4) node[right,black] {\small répond};
\foreach \yy/\dd in {2.9/-1,3.4/1}{
  \foreach \k in {1,2,3}{\draw[orange!70!yellow,line width=0.8pt] (1.5+\dd*0.12+\dd*0.13*\k,\yy+0.1-0.1*\k+0.1) -- ++(\dd*0.25,0);}
  \fill[yellow!70!orange,draw=orange!70!black] (1.5,\yy) circle (0.14);}
\node[font=\scriptsize,align=center] at (1.5,3.95) {tache de lumière};
\end{tikzpicture}
\caption{Détecteur de direction du mouvement. $B$ excite $Y$, $A$ l'inhibe par un intermédiaire \nt{GABA} avec un retard $d$. Pour un mouvement de $A$ vers $B$ à une vitesse proche de $\Delta x/d$, l'inhibition éteint l'excitation; pour le mouvement inverse, elle arrive trop tard. La tache jaune avec ses traits est la lumière qui court au-dessus des capteurs.}
\label{fig:direction}
\end{figure}

Dans la rétine, la sélectivité à la direction du mouvement semble bien naître ainsi: par une inhibition asymétrique venue du côté d'où le mouvement ne doit pas provoquer de réponse~\cite{Barlow1965}.

\section{Soustraction ou division}

Jusqu'ici, notre inhibition soustrayait. En réalité, elle est plus subtile.

Une synapse ouvre un canal, c'est-à-dire qu'elle branche une conductance. Pour une synapse excitatrice, le potentiel d'équilibre se situe bien au-dessus du seuil, environ $70\mV$ au-dessus du repos: le canal ouvert tire la tension vers le haut. Pour une synapse \nt{GABA} inhibitrice, le canal laisse passer le chlorure, dont le potentiel d'équilibre est proche du repos, et le canal ouvert tire la tension non pas vers le bas, mais \emph{vers le repos}. Si l'on compte la tension à partir du repos, la tension établie d'une membrane de conductance de fuite $g_L$, de conductance excitatrice $g_E$ et de conductance inhibitrice $g_I$ vaut
\begin{empheq}[box=\keybox]{equation}
V_\infty \;=\; \frac{g_E\,E_E}{g_L+g_E+g_I} ,
\label{eq:shunt}
\end{empheq}
où $E_E\approx70\mV$ est le potentiel d'équilibre de la synapse excitatrice. C'est la loi d'Ohm pour un diviseur de tension: $g_E$ tire vers $E_E$, $g_L$ et $g_I$ vers zéro.

$g_I$ ne figure pas au numérateur avec un signe moins, mais au dénominateur: l'inhibition ne \emph{soustrait} pas de l'excitation, elle la \emph{divise} (fig.~\ref{fig:shunt}). On appelle cette inhibition shuntante: les canaux chlorure ouverts court-circuitent la membrane. Pour le réseau, c'est un réglage du gain. Si $g_I$ est proportionnelle à l'activité totale des voisins, chaque neurone ne répond pas à la force absolue de son entrée, mais à sa part dans l'activité totale. Cette division par l'activité totale se rencontre dans beaucoup de régions du cerveau et passe pour l'une de ses opérations standard~\cite{Carandini2012}: un même réseau fonctionne avec une entrée faible comme avec une entrée forte.

Une réserve: la formule~\eqref{eq:shunt} concerne un neurone qui n'émet pas d'impulsions. Chez un neurone qui se déclenche, la tension reste en permanence près du seuil, le courant à travers la conductance inhibitrice est presque constant, et le shunt agit sur la \emph{fréquence} des impulsions surtout par soustraction~\cite{HoltKoch1997}. La fréquence est divisée si l'on ajoute en même temps au neurone des conductances de fond excitatrice et inhibitrice équilibrées, comme dans le régime fluctuant et équilibré du cortex vivant~\cite{Chance2002}. Le cerveau n'obtient donc pas la division d'un shunt seul, mais de l'inhibition associée au bruit de fond~\cite{Carandini2012}.

\begin{figure}[t]
\centering
\begin{tikzpicture}[>=Stealth,x=1.25cm,y=0.05cm]
\draw[->,gray!70!black] (0,0) -- (5.4,0) node[right,black] {\small $g_E/g_L$};
\draw[->,gray!70!black] (0,0) -- (0,68) node[above,black] {\small $V_\infty$, mV};
\foreach \v in {20,40,60}{\draw[gray!60!black] (-0.05,\v) -- (0.05,\v) node[left=3pt,font=\scriptsize] {\v};}
\foreach \x in {1,2,3,4,5}{\draw[gray!60!black] (\x,-1.5) -- (\x,1.5) node[below=5pt,font=\scriptsize] {\x};}
\draw[very thick,blue!60!black] plot[domain=0:5,samples=60] (\x,{70*\x/(1+\x)});
\draw[very thick,red!70!black] plot[domain=0:5,samples=60] (\x,{70*\x/(3+\x)});
\draw[thick,densely dashed,gray!50!black] plot[domain=0.56:5,samples=60] (\x,{70*\x/(1+\x)-25});
\node[right,font=\small,blue!60!black] at (5.02,58.3) {sans inhibition};
\node[right,font=\small,red!70!black] at (5.02,45.5) {division: $g_I=2g_L$};
\node[right,font=\small,gray!40!black] at (5.02,32.5) {soustraction de $25\mV$};
\end{tikzpicture}
\caption{L'inhibition shuntante divise la tension au lieu d'en soustraire. Courbe bleue: tension établie~\eqref{eq:shunt} sans inhibition; courbe rouge: avec une conductance inhibitrice $g_I=2g_L$. La rouge est la même que la bleue, étirée trois fois à l'horizontale: tout se passe comme si l'entrée était divisée par $1+g_I/g_L=3$. En tirets: une inhibition qui soustrait $25\mV$ constants; la courbe descend en bloc, sa pente ne change pas.}
\label{fig:shunt}
\end{figure}

Il y a une seconde conséquence. La constante de temps de la membrane vaut $\tau_{\text{eff}}=C/(g_L+g_E+g_I)$, et l'inhibition la raccourcit. Or la fenêtre de coïncidence~\eqref{eq:window} est proportionnelle à $\tau$; une inhibition arrivée en même temps que l'excitation \emph{rétrécit donc la fenêtre}. Le circuit dans lequel une entrée excite un neurone directement et, en même temps, avec une milliseconde de retard, l'inhibe par un intermédiaire \nt{GABA} est l'un des plus répandus du cerveau: le neurone n'a le temps de répondre qu'à ce qui est arrivé pendant la première ou la deuxième milliseconde~\cite{Pouille2001}.

\section{Le choix}

Passons à l'essentiel de ce qui manquait au réseau \nt{GLUT}: choisir un élément parmi plusieurs. Prenons deux groupes de neurones \nt{GLUT} d'entrées $I_1$ et $I_2$. Chacun inhibe l'autre par ses propres intermédiaires \nt{GABA}, avec la force $\beta$ (fig.~\ref{fig:wta}). Pour les fréquences $r_1,r_2$, on obtient
\begin{equation}
\tau\,\frac{\dd r_1}{\dd t} = -r_1 + \bigl[I_1-\beta r_2\bigr]_+ ,\qquad
\tau\,\frac{\dd r_2}{\dd t} = -r_2 + \bigl[I_2-\beta r_1\bigr]_+ ,
\label{eq:wta}
\end{equation}
où $[u]_+=\max(u,0)$: une fréquence n'est jamais négative.

\begin{figure}[t]
\centering
\begin{tikzpicture}[>=Stealth,
  nrn/.style={circle,draw,very thick,fill=blue!6,minimum size=0.9cm,inner sep=0pt},
  gab/.style={nrn,fill=red!10,minimum size=0.75cm},
  inh/.style={-{Bar[width=7pt]},thick,red!70!black}]
\node[nrn] (E1) at (0,0) {$r_1$};
\node[nrn] (E2) at (4,0) {$r_2$};
\node[gab] (G1) at (1.4,1.1) {};
\node[gab] (G2) at (2.6,-1.1) {};
\draw[->,thick] (E1) -- (G1); \draw[inh] (G1) -- (E2);
\draw[->,thick] (E2) -- (G2); \draw[inh] (G2) -- (E1);
\draw[->,thick,blue!60!black] (-1.6,0) node[left,black] {$I_1$} -- (E1);
\draw[->,thick,blue!60!black] (5.6,0) node[right,black] {$I_2$} -- (E2);
\end{tikzpicture}
\caption{Choisir entre deux. Deux groupes de neurones \nt{GLUT} (en bleu) s'inhibent mutuellement par des intermédiaires \nt{GABA} (en rouge).}
\label{fig:wta}
\end{figure}

Tant que les deux neurones travaillent, le système~\eqref{eq:wta} est linéaire, et son comportement est fixé par deux valeurs propres: $-(1+\beta)/\tau$ pour des écarts identiques des deux fréquences et $-(1-\beta)/\tau$ pour des écarts opposés. Avec une inhibition faible, $\beta<1$, les deux valeurs sont négatives: les deux neurones travaillent, l'inhibition ne fait qu'atténuer le plus faible. Avec une inhibition forte, $\beta>1$, la seconde valeur devient positive: toute différence entre $r_1$ et $r_2$ grandit jusqu'à ce qu'un neurone se taise complètement.

\begin{keyblock}
\begin{proposition}[Le gagnant rafle tout]
\label{prop:wta}
Si l'inhibition mutuelle est supérieure à un, $\beta>1$, l'état où les deux groupes travaillent est instable. Le réseau aboutit à un état où un groupe travaille et l'autre se tait. Le gagnant, de fréquence $r_1=I_1$, conserve sa victoire tant que $I_2<\beta I_1$.
\end{proposition}
\end{keyblock}

Nous l'avons vérifié sur le modèle. Pour $\beta=0.5$ et des entrées $1.0$ et $0.9$, les deux groupes travaillent, aux fréquences $0.73$ et $0.53$. Pour $\beta=2$ et les mêmes entrées, le second groupe se tait complètement. Ce choix a une mémoire: si l'on échange ensuite les entrées, l'ancien gagnant reste gagnant, car $1.0<2\cdot0.9$.

Avec cent groupes, c'est pareil: chacun inhibe les autres, un seul survit.

\section{Effacer la mémoire}

La mémoire du chapitre~\ref{sec:glutcomputer} ne s'éteignait que par la fatigue. Ajoutons-lui un neurone \nt{GABA} qu'excite un signal de \og remise à zéro\fg{} et qui inhibe le groupe. L'inhibition abaisse la courbe $f(wr+I)$ de la fig.~\ref{fig:bistable}, l'intersection supérieure disparaît, et le groupe s'éteint, puis reste dans l'état bas une fois le signal terminé. La mémoire s'écrit maintenant par un signal et s'efface par un autre, comme une bascule à entrées de mise à un et de remise à zéro.

Plus élégant encore: relier deux de ces groupes par une inhibition mutuelle, comme sur la fig.~\ref{fig:wta}, et donner à chacun une rétroaction sur lui-même. Un signal à l'entrée d'un des groupes fait passer la cellule dans l'état de ce groupe, et la cellule se souvient de la dernière décision. C'est l'analogue direct du verrou formé de deux portes NON-OU croisées.

\section{Stabilité et rythme}

Reste le troisième mal du réseau \nt{GLUT}: l'avalanche. Prenons un groupe de neurones \nt{GLUT} avec une rétroaction sur lui-même $w_{EE}$ et un groupe de neurones \nt{GABA} que le premier excite avec la force $w_{IE}$ et qui inhibent le premier avec la force $w_{EI}$. Pour les fréquences $E$ et $I$, on obtient
\begin{equation}
\tau_E\frac{\dd E}{\dd t} = -E + w_{EE}E - w_{EI}I + h, \qquad
\tau_I\frac{\dd I}{\dd t} = -I + w_{IE}E ,
\label{eq:ei}
\end{equation}
où $h$ est l'entrée extérieure~\cite{WilsonCowan1972}. Sans inhibition, pour $w_{EE}>1$, l'activité croît sans limite: chaque impulsion en engendre plus d'une suivante. Avec l'inhibition, la fréquence établie
\begin{equation}
E^* \;=\; \frac{h}{1-w_{EE}+w_{EI}w_{IE}}
\label{eq:eistar}
\end{equation}
est finie si le dénominateur est positif. Mais cela ne suffit pas: il faut encore que l'équilibre soit attractif. L'analyse linéaire de~\eqref{eq:ei} donne deux conditions.

\begin{keyblock}
\begin{proposition}[Conditions de stabilité]
\label{prop:ei}
L'équilibre~\eqref{eq:eistar} est stable si l'inhibition est
\begin{enumerate}
\item[(i)] \emph{assez forte}: $w_{EI}\,w_{IE} > w_{EE}-1$;
\item[(ii)] \emph{assez rapide}: $\tau_I < \tau_E/(w_{EE}-1)$.
\end{enumerate}
Si (i) n'est pas satisfaite, l'activité croît en avalanche. Si (i) est satisfaite mais pas (ii), l'inhibition arrive en retard, et l'activité se met à osciller.
\end{proposition}
\end{keyblock}

La condition~(i) porte sur le déterminant de la matrice du système~\eqref{eq:ei}, la condition~(ii) sur sa trace. Le sens de la seconde est clair pour quiconque a réglé un régulateur: une rétroaction en retard n'amortit pas l'écart, elle l'amplifie.

\begin{figure}[t]
\centering
\begin{tikzpicture}[>=Stealth]
\draw[->,gray!70!black] (0,0) -- (10.9,0) node[right,black] {\small $t$, ms};
\draw[->,gray!70!black] (0,0) -- (0,3.3) node[left,black] {\small $E$};
\foreach \t in {0,50,100,150}{\draw[gray!70!black] (\t*0.07,0.06) -- (\t*0.07,-0.06) node[below,font=\scriptsize] {\t};}
\input{figures/ei_traces}
\node[font=\small,gray!40!black,anchor=west] at (1.3,3.45) {sans inhibition: avalanche};
\node[font=\small,blue!60!black,anchor=west] at (7.4,1.25) {inhibition rapide};
\node[font=\small,red!70!black,anchor=west] at (7.4,2.55) {inhibition lente};
\end{tikzpicture}
\caption{Fréquence d'un groupe \nt{GLUT} avec une rétroaction $w_{EE}=1.5$, $\tau_E=10\ms$, après l'application de l'entrée. Sans inhibition (tirets), l'activité croît sans limite. Avec une inhibition de force $w_{EI}w_{IE}=2$ et $\tau_I=2\ms$ (courbe bleue), elle se stabilise au niveau $E^*=h/(1-1.5+2)=0.67h$. Avec la même inhibition, mais lente, $\tau_I=30\ms$ (courbe rouge), l'activité oscille.}
\label{fig:ei}
\end{figure}

On considère que le cortex fonctionne justement dans ce régime: sa propre rétroaction est assez forte pour partir en avalanche sans inhibition, et seule une inhibition rapide le maintient à son point de fonctionnement~\cite{Tsodyks1997isn}.

Ce régime a une signature paradoxale qui permet de le reconnaître de l'extérieur~\cite{Tsodyks1997isn}. Ajoutons dans~\eqref{eq:ei} une entrée extérieure $h_I$ directement au groupe \nt{GABA}. Les fréquences établies deviennent
\begin{equation}
E^* \;=\; \frac{h-w_{EI}\,h_I}{D},\qquad
I^* \;=\; \frac{w_{IE}\,h+(1-w_{EE})\,h_I}{D},\qquad
D \;=\; 1-w_{EE}+w_{EI}w_{IE} .
\label{eq:isn}
\end{equation}
Si $w_{EE}>1$, le coefficient de $h_I$ dans la seconde formule est négatif: poussez les neurones \nt{GABA}, et leur fréquence \emph{baisse}. La cause est la boucle. Le surcroît d'inhibition atténue le groupe \nt{GLUT}, qui excite moins les neurones \nt{GABA}, et ceux-ci perdent plus qu'ils n'ont reçu. Avec les valeurs de la fig.~\ref{fig:ei} ($w_{EE}=1.5$, $w_{EI}w_{IE}=2$, $D=1.5$), chaque unité d'entrée ajoutée abaisse $I^*$ d'un tiers d'unité. C'est précisément cette signature qu'on a trouvée dans le cortex: quand la réponse des neurones du cortex visuel est supprimée par un stimulus environnant, le courant inhibiteur qu'ils reçoivent n'augmente pas, mais baisse en même temps que le courant excitateur~\cite{Ozeki2009}. Plus tard, la même signature a été trouvée dans différentes aires et couches du cortex~\cite{Sanzeni2020}.

Les oscillations dues à la violation de la condition~(ii) se rencontrent elles aussi dans le cerveau. La boucle \og \nt{GLUT} excite \nt{GABA}, \nt{GABA} inhibe \nt{GLUT}\fg{} avec un retard de quelques millisecondes engendre un rythme de période de l'ordre de $12$--$50\ms$ (fréquence de $20$--$80$\,Hz), et ces rythmes rapides du cortex sont bien connus~\cite{Cardin2009,Buzsaki2012}. Ce n'est pas l'horloge d'un ordinateur: le rythme est local et n'apparaît que quand le réseau est actif. Mais pendant quelques cycles, il découpe le temps en fenêtres dans lesquelles les neurones peuvent se déclencher.

\section{Bilan}

\begin{table}[t]
\centering
\small
\begin{tabular}{>{\raggedright}p{3.3cm}>{\raggedright}p{4.4cm}>{\raggedright\arraybackslash}p{4.4cm}}
\toprule
& \nt{GLUT} seul & \nt{GLUT} + \nt{GABA} \\
\midrule
NON & seulement par des capteurs ON/OFF & veto $a\wedge\bar b$; NON avec activité de fond \\
fils par bit & deux & un \\
direction du mouvement & non & veto retardé \\
réglage du gain & non & division: shunt associé au bruit de fond \\
fenêtre de coïncidence & $\sim\tau$ & rétrécie par l'inhibition \\
choix d'un parmi plusieurs & non & inhibition mutuelle, $\beta>1$ \\
effacement de la mémoire & seulement par fatigue & par un signal via \nt{GABA} \\
stabilité & seulement par fatigue & rétroaction négative rapide \\
rythme & inutile & apparaît avec une inhibition lente \\
\bottomrule
\end{tabular}
\caption{Ce que \nt{GABA} ajoute à un réseau de neurones \nt{GLUT}.}
\label{tab:gaba}
\end{table}

\nt{GABA} n'ajoute presque pas de nouveaux \emph{calculs} au réseau (tout cela pouvait se calculer aussi avec des capteurs à deux fils), mais il ajoute la \emph{commande} (tableau~\ref{tab:gaba}): le choix, la remise à zéro, le réglage du gain et la stabilité. Dans le cerveau, l'excitation porte les données, et l'inhibition décide quelles données passent, quand et à quel volume. Pour un neurone du cortex sur quatre ou cinq, c'est une division du travail très raisonnable.

\section{Exercices}

\begin{exercise}[\lvlA]
\label{ex:03:1}
\emph{Le shunt en chiffres.} $E_E=70\mV$. D'après~\eqref{eq:shunt}, trouvez $V_\infty$ pour $g_E=g_L$ et $4g_L$, sans inhibition et avec un shunt $g_I=2g_L$. Quel $V_\infty$ pour $g_E=4g_L$ donnerait une soustraction qui retire autant que le shunt pour $g_E=g_L$? De quel facteur le shunt rétrécit-il la fenêtre de coïncidence?
\end{exercise}

\begin{exercise}[\lvlB]
\label{ex:03:2}
\emph{Dériver les conditions de stabilité.} Trouvez la trace et le déterminant de la matrice du système linéaire~\eqref{eq:ei} et déduisez-en les deux conditions de la proposition~\ref{prop:ei}. À la limite de la condition~(ii), l'équilibre perd sa stabilité de façon oscillante. Trouvez la période de ces oscillations avec les valeurs de la fig.~\ref{fig:ei}.
\end{exercise}

\begin{exercise}[\lvlB]
\label{ex:03:3}
\emph{Un rythme né d'un retard.} Remplaçons l'inhibition lente du modèle~\eqref{eq:ei} par une inhibition rapide avec un retard $d$: $\tau_E\dot E=-E(t)-GE(t-d)+h$. Pour $\tau_E=10\ms$, trouvez le plus petit retard $d_c$ qui donne des oscillations, et leur période pour $G=2$ et $G=5$. Tombent-elles dans les rythmes rapides du cortex, $12$--$50\ms$, contrairement à l'exercice~\ref{ex:03:2}?
\end{exercise}

\begin{exercise}[\lvlB]
\label{ex:03:4}
\emph{Modélisation: un choix avec mémoire.} Intégrez~\eqref{eq:wta} pour $\beta=2$, $\tau=1$. (a)~Avec $I_1=1$, faites monter lentement $I_2$ de $0$ à $3$ puis redescendre: pour quels $I_2$ le réseau bascule-t-il? (b)~De combien augmente le temps de décision à partir du silence quand l'écart des entrées $\varepsilon$ est divisé par dix?
\end{exercise}

\begin{exercise}[\lvlB]
\label{ex:03:5}
\emph{Conception: un détecteur pour une bande de vitesses.} Le détecteur de la fig.~\ref{fig:direction} doit se taire pour un mouvement de $A$ vers $B$ parcouru en $5$--$20\ms$. Un intermédiaire \nt{GABA} de retard $d_k$ interdit $Y$ sur $[d_k,\,d_k+5\ms]$ après une impulsion de $A$; l'excitation venue de $B$ est immédiate. Combien d'intermédiaires faut-il, et avec quels retards?
\end{exercise}

\begin{exercise}[\lvlB]
\label{ex:03:6}
\emph{Vetos et fond.} Combien de neurones demande la parité $x\oplus y\oplus z$ dans le schéma général \og un intermédiaire \nt{GABA} par entrée, un veto par combinaison où $f=1$, un OU commun\fg? Construisez-la avec deux neurones, un \nt{GABA} et un \nt{GLUT}, recevant les entrées directement, en donnant poids et seuils.
\end{exercise}

\begin{exercise}[\lvlC]
\label{ex:03:7}
\emph{Conception: choisir parmi cent.} Chacun des $100$ groupes \nt{GLUT} doit inhiber également tous les autres, mais pas lui-même. Des intermédiaires \nt{GABA} linéaires sont excités par les groupes et les inhibent; les groupes peuvent s'exciter entre eux, mais pas eux-mêmes. Combien faut-il d'intermédiaires au minimum? Et sans connexion excitatrice directe?
\end{exercise}

\begin{exercise}[\lvlC]
\label{ex:03:8}
\emph{Recherche: quand l'inhibition est paradoxale.} Dans le modèle de la fig.~\ref{fig:ei} ($w_{EI}=2$, $w_{IE}=1$), le groupe \nt{GLUT} reçoit la fonction de transfert $f(u)=[u]_+^2$, le groupe \nt{GABA} reste linéaire. Au-dessus de quelle entrée $h_c$ un apport supplémentaire au groupe \nt{GABA} abaisse-t-il sa propre fréquence? Vérifiez par simulation.
\end{exercise}

\chapter{Le code Morse: ce que l'on sait de la langue des neurones \nt{GLUT} et \nt{GABA}}
\chaptermark{Le code Morse}
\label{sec:code}

\section{Un alphabet d'un seul signe}

Le code Morse a deux signes, le point et le trait, et des pauses de trois longueurs entre eux. L'alphabet du neurone, nous l'avons vu au chapitre~\ref{sec:neurontypes}, est encore plus pauvre: il n'a qu'un signe. Une impulsion dure environ une milliseconde, et toutes les impulsions d'un même neurone ont la même hauteur et la même forme. Tout le message est donc dans les pauses, c'est-à-dire dans la suite des instants $t_1,t_2,t_3,\dots$ C'est un Morse sans traits, où seul le rythme des points porte le sens (fig.~\ref{fig:morse}).

\begin{figure}[t]
\centering
\begin{tikzpicture}[>=Stealth,x=0.08cm,y=1cm]
% Morse letter: dot, dash, dot
\node[left,font=\small] at (-6,2.55) {Morse:};
\fill[black] (0,2.45) rectangle (6,2.65);
\fill[black] (14,2.45) rectangle (32,2.65);
\fill[black] (40,2.45) rectangle (46,2.65);
\node[right,font=\small,gray!40!black] at (52,2.55) {point, trait, point: le sens est dans les durées et les pauses};
% stimulus onset
\node[left,font=\small] at (-6,0.4) {neurone:};
\draw[->,thick,green!45!black] (0,-0.55) -- (0,-0.05);
\node[below,font=\scriptsize,green!45!black] at (0,-0.55) {stimulus};
\draw[gray!70!black] (0,0) -- (152,0) node[right,black] {\small $t$};
% spikes: single ones blue, the burst red
\foreach \s in {14,26,41,88,112,131}{\draw[very thick,blue!60!black] (\s,0) -- (\s,0.8);}
\foreach \s in {60,63,66,69}{\draw[very thick,red!70!black] (\s,0) -- (\s,0.8);}
% latency
\draw[<->,thick] (0,1.1) -- node[above,font=\scriptsize] {$t_1$: instant de la 1\textsuperscript{re} impulsion} (14,1.1);
% burst
\draw[decorate,decoration={brace,amplitude=4pt},red!70!black] (58,0.95) -- (71,0.95) node[midway,above=4pt,font=\scriptsize] {bouffée: un \og trait\fg};
% counting windows
\foreach \a/\b/\n in {0/50/3,50/100/5,100/150/2}{
  \draw[decorate,decoration={brace,amplitude=4pt,mirror},gray!50!black] (\a+1,-0.2) -- (\b-1,-0.2) node[midway,below=4pt,font=\small] {$n=\n$};}
\node[font=\scriptsize,gray!40!black] at (75,-1.05) {code de fréquence: nombre d'impulsions $n$ dans chaque fenêtre};
\foreach \x in {0,50,100,150}{\draw[gray!60!black] (\x,-0.06) -- (\x,0.06);}
\end{tikzpicture}
\caption{Le code Morse et l'alphabet du neurone. Une même suite d'impulsions peut se lire de plusieurs façons: compter les impulsions dans des fenêtres de $50\ms$ (section~\ref{sec:rate}), mesurer le délai $t_1$~\eqref{eq:latency} ou repérer une bouffée, un \og trait\fg~\eqref{eq:burst}.}
\label{fig:morse}
\end{figure}

Quelles pauses sont possibles? Par en dessous, elles sont limitées par la période réfractaire: après une impulsion, le neurone ne peut pas se déclencher de nouveau pendant environ une milliseconde, si bien qu'on ne dépasse pas quelques centaines d'impulsions par seconde. Par au-dessus, rien ne les limite: un neurone peut se taire pendant des minutes. Chez les neurones \nt{GLUT} du cortex, les fréquences vont d'une fraction d'impulsion à quelques dizaines d'impulsions par seconde, et la plupart des neurones travaillent la plupart du temps près de la borne inférieure.

Aux chapitres~\ref{sec:glutcomputer} et~\ref{sec:gaba}, nous avons adopté au passage tantôt une façon, tantôt une autre d'écrire des nombres avec des impulsions. Posons maintenant la question directement: dans quelle langue les neurones \nt{GLUT} et \nt{GABA} se parlent-ils? Il n'y a pas de réponse complète, cette langue n'est pas déchiffrée. Mais on en sait quelques choses avec certitude, et presque tout ce qu'on en sait peut s'obtenir en raisonnant en ingénieur des télécommunications: capacité du canal, bruit, récepteur, coût de la transmission.

\section{Combien peut-on transmettre}

Commençons par la borne supérieure. Supposons que le destinataire distingue les instants des impulsions avec une précision $\Delta t$: les décalages inférieurs à $\Delta t$ sont pour lui indiscernables. Découpons le temps en cases de longueur $\Delta t$, plus courtes que la période réfractaire, de sorte que chaque case contienne une impulsion ou aucune (fig.~\ref{fig:cells}). À une fréquence moyenne $r$, une impulsion tombe dans une case avec la probabilité $p=r\Delta t$. Le flux porte le plus d'information quand les cases sont indépendantes, et chaque case donne alors une entropie de $-p\log_2p-(1-p)\log_2(1-p)\approx p\log_2(e/p)$ bits pour $p\ll1$. En divisant par $\Delta t$, on obtient le débit maximal et sa part par impulsion~\cite{MacKay1952}:
\begin{empheq}[box=\keybox]{equation}
R_{\max} \;\approx\; r\,\log_2\frac{e}{r\,\Delta t}\ \ \text{bits/s},
\qquad
\frac{R_{\max}}{r} \;\approx\; \log_2\frac{e}{r\,\Delta t}\ \ \text{bits par impulsion}.
\label{eq:spikeentropy}
\end{empheq}
Pour $r=10$ impulsions par seconde et $\Delta t=1\ms$, cela fait environ huit bits par impulsion et quatre-vingts bits par seconde.

\begin{figure}[t]
\centering
\begin{tikzpicture}[>=Stealth,x=0.5cm,y=1cm]
% time cut into cells of width Delta t; a cell holds one impulse or none
\foreach \k in {0,...,23}{\draw[gray!55] (\k,0) rectangle (\k+1,0.7);}
\foreach \k in {2,7,8,15,21}{
  \fill[red!10] (\k,0) rectangle (\k+1,0.7);
  \draw[very thick,red!70!black] (\k+0.5,0.05) -- (\k+0.5,0.65);}
\foreach \k in {0,...,23}{
  \pgfmathtruncatemacro{\bit}{(\k==2)||(\k==7)||(\k==8)||(\k==15)||(\k==21)}
  \ifnum\bit=1 \def\bitcol{red!70!black}\else \def\bitcol{gray!60!black}\fi
  \node[font=\scriptsize,text=\bitcol] at (\k+0.5,-0.25) {\bit};}
\draw[<->,gray!60!black] (0,0.95) -- (1,0.95) node[midway,above,font=\scriptsize,black] {$\Delta t$};
\node[font=\small,anchor=east] at (-0.3,0.35) {impulsions};
\node[font=\small,anchor=east] at (-0.3,-0.25) {bits};
\draw[decorate,decoration={brace,amplitude=5pt,mirror},gray!60!black] (0,-0.55) -- (24,-0.55)
  node[midway,below=6pt,font=\scriptsize,black,align=center] {un destinataire qui ne fait que compter: \og $n=5$\fg;\\ un destinataire qui voit les cases: lesquelles des $24$ contiennent les $5$, soit $\log_2\binom{24}{5}\approx15$ bits};
\end{tikzpicture}
\caption{D'où vient la capacité~\eqref{eq:spikeentropy}. Le temps est découpé en cases de longueur $\Delta t$; chaque case contient une impulsion ($1$) ou n'en contient pas ($0$): le flux d'impulsions est une chaîne binaire. Une impulsion tombe dans une case avec la probabilité $p=r\Delta t$. Plus les cases sont fines, plus la chaîne est longue et plus il y a de bits; un compteur d'impulsions perd presque tout ce qui est inscrit dans la \emph{position} des uns.}
\label{fig:cells}
\end{figure}

Les bits par impulsion ne dépendent de la précision temporelle que de façon logarithmique: avec une précision de $10\ms$, il reste environ cinq bits par impulsion. Mais si le destinataire ne fait que compter les impulsions sur une seconde, pour $r=10$ il reçoit moins de deux bits pour toute la seconde (section~\ref{sec:rate}).

\begin{keyblock}
\begin{proposition}[Capacité d'un canal à impulsions]
\label{prop:capacity}
Un flux d'impulsions de fréquence moyenne $r$, lu avec une précision temporelle $\Delta t$, ne porte pas plus que~\eqref{eq:spikeentropy}. Presque toute cette capacité réside dans l'\emph{instant} des impulsions: un destinataire qui ne fait que compter les impulsions tire du même flux des dizaines de fois moins qu'un destinataire qui distingue leurs instants à la milliseconde près.
\end{proposition}
\end{keyblock}

C'est une borne supérieure, pas une mesure. Les mesures sur des neurones sensoriels dont on connaît le stimulus donnent d'un à quelques bits par impulsion; dans les meilleurs cas, cela atteint la moitié de la borne calculée pour leur précision~\cite{Rieke1997}. Donc, au moins à l'entrée du système nerveux, l'instant des impulsions est utilisé et non perdu.

\section{Un canal bruité}

La capacité~\eqref{eq:spikeentropy} a été calculée sans bruit. Sur l'endroit où il naît, trois faits sont établis.

\emph{Le générateur d'impulsions lui-même est précis.} Si l'on injecte de nombreuses fois le même courant, variant vite dans le temps, dans un neurone du cortex prélevé avec un morceau de tissu, les impulsions se répètent à environ une milliseconde près~\cite{Mainen1995}. Si le courant est constant, les instants des impulsions se dispersent d'un essai à l'autre: c'est la variation rapide de l'entrée, et non le neurone lui-même, qui produit la précision temporelle.

\emph{La synapse n'est pas fiable.} Une impulsion qui atteint une synapse \nt{GLUT} ne libère une dose de neurotransmetteur qu'avec une certaine probabilité $p$. Aux synapses de l'hippocampe, plus de la moitié des impulsions n'atteignent tout simplement pas le destinataire, et la cause en est précisément le caractère aléatoire de la libération, non des défauts de conduction~\cite{AllenStevens1994}. Pour un ingénieur des télécommunications, c'est un canal à effacements; plusieurs synapses entre deux neurones sauvent en partie la situation.

\emph{Dans le cortex vivant, le flux d'impulsions paraît aléatoire.} Les intervalles entre impulsions sont dispersés à peu près comme dans un flux aléatoire de Poisson, et quand on répète le même stimulus, le nombre d'impulsions dans une fenêtre varie de sorte que sa variance est proche de sa moyenne~\cite{Softky1993}.

Comment un élément précis produit-il un flux aléatoire? Le neurone reçoit des milliers d'entrées, chacune peu fiable, et l'excitation et l'inhibition s'équilibrent, comme au chapitre~\ref{sec:gaba}. La tension de membrane fluctue près du seuil, et ce sont ces fluctuations qui fixent les instants des impulsions. Est-ce du bruit, ou un message que nous ne savons pas lire? La question reste largement ouverte. Une partie du \og bruit\fg{} s'est déjà révélée être un message: dans le cortex visuel, une part importante de la variabilité suit les mouvements de l'animal lui-même~\cite{Stringer2019}. La difficulté est ici de principe: pour qui ne connaît pas le code, un message bien compressé est indiscernable d'un flux aléatoire.

\section{La fréquence: lente, mais sûre}
\label{sec:rate}

Le code le plus simple est le \emph{code de fréquence}: le nombre est inscrit dans le nombre d'impulsions que le neurone émet pendant une fenêtre $T$. Ni les effacements ni la gigue des instants ne menacent ce code. Mais il a un prix. Si le flux est poissonien, le nombre d'impulsions $n$ dans la fenêtre a pour moyenne $rT$ et pour écart type $\sqrt{rT}$:
\begin{empheq}[box=\keybox]{equation}
\frac{\sigma_n}{\bar n} \;=\; \frac{1}{\sqrt{r\,T}} .
\label{eq:rateerror}
\end{empheq}
Pour connaître une fréquence à dix pour cent près, il faut cent impulsions, soit cinq secondes à vingt impulsions par seconde. Des niveaux voisins sont discernables s'ils sont séparés de deux écarts types environ; jusqu'à la fréquence $r$, on distingue donc en un temps $T$ environ $\sqrt{rT}$ niveaux: pour $r=10$ et $T=1\,\text{s}$, trois ou quatre niveaux, moins de deux bits.
Le cerveau ne travaille pas si lentement. Un homme reconnaît un animal sur une image montrée un instant, et la réponse apparaît dans le cerveau au bout d'environ $150\ms$~\cite{Thorpe1996}. Selon une estimation grossière, le signal traverse pendant ce temps une dizaine d'étapes de traitement successives, à raison de $10$--$20\ms$ par étape. Aux fréquences habituelles du cortex, chaque neurone n'a le temps d'émettre que zéro ou une impulsion par étape, et sa fréquence n'est pas définie sur une telle fenêtre. Deux issues: prendre beaucoup de neurones, ou regarder le temps.

\emph{Beaucoup de neurones.} Supposons que $N$ neurones portent le même nombre et que le destinataire additionne leurs impulsions. Si leurs bruits sont indépendants, $rT$ est remplacé dans~\eqref{eq:rateerror} par $NrT$: cent neurones à $r=20$ donnent en $15\ms$ trente impulsions et une précision d'environ vingt pour cent. L'espace remplace le temps. Mais les bruits de neurones voisins du cortex ne sont pas indépendants: le coefficient de corrélation des nombres d'impulsions d'une paire de voisins est en général d'environ $0.1$~\cite{Zohary1994}. S'il vaut $c$, la variance de la moyenne sur $N$ neurones vaut
\begin{empheq}[box=\keybox]{equation}
\sigma^2_N \;=\; \sigma^2_1\,\frac{1+(N-1)\,c}{N}
\;\xrightarrow[N\to\infty]{}\; c\,\sigma^2_1 .
\label{eq:correlated}
\end{empheq}

\begin{keyblock}
\begin{proposition}[Limite du moyennage]
\label{prop:pooling}
Avec un bruit corrélé, le moyennage sur un groupe ne réduit pas l'erreur de plus d'un facteur $1/\sqrt{c}$. Pour $c=0.1$, c'est un facteur trois, et ce gain est atteint dès quelques dizaines de neurones; en ajouter davantage est inutile.
\end{proposition}
\end{keyblock}

Toute corrélation n'est pas également nuisible. Seule limite la précision la part du bruit commun qui \emph{ressemble au signal lui-même}, c'est-à-dire qui déplace tout le groupe comme le déplacerait une variation de la grandeur mesurée~\cite{MorenoBote2014}. Combien le cerveau contient réellement d'un tel bruit, on cherche encore à l'établir.

Il existe une autre façon d'utiliser beaucoup de neurones: inscrire le nombre dans l'identité du neurone qui s'est déclenché, comme dans le détecteur de direction d'un son (fig.~\ref{fig:itd}). C'est le code \emph{de position}, le plus fiable que l'on connaisse: chez les neurones sensoriels, le numéro du fil indique quel point de la peau est touché ou quelle est la hauteur d'un son, et cela a été établi maintes fois. Il est coûteux en neurones, mais rapide: un message ne prend qu'un seul retard.

\section{Le temps: rapide, mais à partir de quand?}

La seconde issue consiste à écrire le nombre dans l'instant de l'impulsion. Prenons le neurone~\eqref{eq:glutneuron} et appliquons-lui, à partir de l'instant $t=0$, une entrée constante qui porterait la tension au niveau $U$. La tension croît comme $V=U\bigl(1-e^{-t/\tau}\bigr)$ et atteint le seuil $\theta$ à l'instant
\begin{empheq}[box=\keybox]{equation}
t_1 \;=\; \tau\,\ln\frac{U}{U-\theta}, \qquad U>\theta .
\label{eq:latency}
\end{empheq}
Entrée forte, impulsion précoce; entrée faible, impulsion tardive; entrée sous le seuil, pas d'impulsion. Pour $\tau=10\ms$, l'entrée $U=4\theta$ donne une première impulsion au bout de $2.9\ms$, $U=2\theta$ au bout de $6.9\ms$, $U=1.2\theta$ au bout de $17.9\ms$. Une seule impulsion porte un nombre.

Mais à partir de quel instant compter $t_1$? Il n'y a pas d'horloge, et le destinataire ne sait pas quand le stimulus a commencé. Il y a trois réponses.

\emph{Le délai relatif.} Deux neurones voient le même stimulus, et la différence de leurs délais $t_1^A-t_1^B$ ne dépend que de leurs entrées, non de l'instant du début. Des détecteurs de coïncidences avec des retards comparent les instants (fig.~\ref{fig:itd}). Si $N$ neurones ont émis chacun une impulsion, l'\emph{ordre} même de leurs arrivées porte jusqu'à $\log_2N!$ bits: pour dix neurones, environ vingt-deux bits, alors que la réponse \og déclenché ou non\fg{} n'en donnerait pas plus de dix. Dans la rétine, on a montré qu'on peut reconstruire une image vite et de façon fiable à partir des délais relatifs des premières impulsions de ses neurones de sortie~\cite{Gollisch2008}.

\emph{La salve comme début de mot.} Un changement brusque du stimulus fait se déclencher presque simultanément un grand nombre de neurones. Cette salve sert elle-même de marque de début, comme la longue pause entre les mots en Morse, et le destinataire compte les délais à partir d'elle.

\emph{Le rythme local.} Au chapitre~\ref{sec:gaba}, la boucle \nt{GLUT}--\nt{GABA} engendrait un rythme local de période $12$--$50\ms$. Ce n'est pas une horloge, mais à l'intérieur de son cycle un neurone à entrée forte a le temps de se déclencher avant un neurone à entrée faible, et~\eqref{eq:latency} devient un code \emph{de phase}: le nombre est inscrit dans la partie du cycle où l'impulsion est arrivée. On a observé ce code: des neurones associés à un lieu précis de l'espace se déclenchent de plus en plus tôt dans le cycle d'un rythme local lent à mesure que l'animal traverse \og leur\fg{} lieu~\cite{OKeefe1993}. Dans quelle mesure le cerveau se sert de la phase, on l'ignore encore.

\section{Le destinataire choisit le code}

Une suite d'impulsions contient à la fois une fréquence, des instants et un ordre. Ce qui en sera lu, c'est le destinataire qui le décide, et il le décide par un seul paramètre: sa constante de temps $\tau$.

Moyennons l'équation~\eqref{eq:glutneuron} dans le temps. Si le neurone reçoit des flux indépendants de fréquences $r_j$, la tension moyenne et sa fluctuation relative valent
\begin{empheq}[box=\keybox]{equation}
\bar V \;=\; \tau\sum_j w_j\,r_j ,
\qquad
\frac{\sigma_V}{\bar V} \;\approx\; \frac{1}{\sqrt{2\,r_{\text{ent}}\,\tau}} ,
\label{eq:reader}
\end{empheq}
où la seconde égalité est écrite pour des poids égaux, et $r_{\text{ent}}$ est la fréquence totale de toutes les entrées. Quand $\tau$ est grande, la tension fluctue à peine et suit la somme pondérée des fréquences: le destinataire est un \emph{intégrateur}, il lit les fréquences et oublie les instants. Quand $\tau$ est petite, la tension a le temps de retomber entre les impulsions, et le seuil n'est atteint que lorsque plusieurs impulsions arrivent dans la fenêtre~\eqref{eq:window}: le destinataire est un \emph{détecteur de coïncidences}, il lit les instants (fig.~\ref{fig:reader}). Mille entrées à cinq impulsions par seconde avec $\tau=10\ms$ donnent une fluctuation d'environ dix pour cent: c'est presque une intégration pure. Si l'inhibition, comme au chapitre~\ref{sec:gaba}, a raccourci $\tau$ à $2\ms$, la fluctuation dépasse vingt pour cent, et le neurone commence à remarquer les coïncidences.

\begin{figure}[t]
\centering
\begin{tikzpicture}[>=Stealth,x=0.115cm,y=1cm]
% common input: the same spike train for both receivers
\foreach \s in {6,21,22,38,52,55.5,59,62.5,66,69.5,73,76.5,80,83.5,87,90.5,94,97.5}{\draw[thick,gray!40!black] (\s,5.25) -- (\s,5.65);}
\draw[gray!70!black] (0,5.25) -- (105,5.25);
\node[left,font=\small] at (-1,5.45) {entrée};
\node[above,font=\scriptsize,gray!40!black] at (13.5,5.65) {rares};
\node[above,font=\scriptsize,gray!40!black] at (21.5,6.0) {paire};
\node[above,font=\scriptsize,gray!40!black] at (75,5.65) {fréquentes};
% axes and thresholds
\foreach \y/\th in {2.7/2.5,0/1.5}{
  \draw[gray!70!black] (0,\y) -- (106,\y) node[right,black] {\small $t$};
  \draw[densely dashed,gray!60!black] (0,\y+0.6*\th) -- (104,\y+0.6*\th) node[right,black,font=\small] {$\theta$};
}
\node[anchor=south west,font=\small] at (0,4.3) {$\tau=10\ms$: intégrateur};
\node[anchor=south east,font=\small] at (104,1.0) {$\tau=2\ms$: détecteur de coïncidences};
% integrator: tau=10 ms, theta=2.5w
\begin{scope}[yshift=2.7cm]
\foreach \a/\b/\v in {0/6/0.000,6/21/1.000,21/22/1.223,22/38/2.107,38/52/1.425,52/55.5/1.351,55.5/59/1.952,59/62.5/2.376,62.5/66/0.000,66/69.5/1.000,69.5/73/1.705,73/76.5/2.201,76.5/80/0.000,80/83.5/1.000,83.5/87/1.705,87/90.5/2.201,90.5/94/0.000,94/97.5/1.000,97.5/104/1.705}{\draw[very thick,blue!60!black] plot[domain=\a:\b,samples=14] (\x,{0.6*\v*exp(-(\x-\a)/10)});}
\foreach \s/\p/\q in {6/0.000/1.000,21/0.223/1.223,22/1.107/2.107,38/0.425/1.425,52/0.351/1.351,55.5/0.952/1.952,59/1.376/2.376,62.5/1.674/2.500,62.5/2.500/0.000,66/0.000/1.000,69.5/0.705/1.705,73/1.201/2.201,76.5/1.551/2.500,76.5/2.500/0.000,80/0.000/1.000,83.5/0.705/1.705,87/1.201/2.201,90.5/1.551/2.500,90.5/2.500/0.000,94/0.000/1.000,97.5/0.705/1.705}{\draw[very thick,blue!60!black] (\s,0.6*\p) -- (\s,0.6*\q);}
\foreach \s in {62.5,76.5,90.5}{\draw[->,thick,red!70!black] (\s,1.58) -- (\s,1.95);}
\end{scope}
% coincidence: tau=2 ms, theta=1.5w
\begin{scope}[yshift=0cm]
\foreach \a/\b/\v in {0/6/0.000,6/21/1.000,21/22/1.001,22/38/0.000,38/52/1.000,52/55.5/1.001,55.5/59/1.174,59/62.5/1.204,62.5/66/1.209,66/69.5/1.210,69.5/73/1.210,73/76.5/1.210,76.5/80/1.210,80/83.5/1.210,83.5/87/1.210,87/90.5/1.210,90.5/94/1.210,94/97.5/1.210,97.5/104/1.210}{\draw[very thick,blue!60!black] plot[domain=\a:\b,samples=14] (\x,{0.6*\v*exp(-(\x-\a)/2)});}
\foreach \s/\p/\q in {6/0.000/1.000,21/0.001/1.001,22/0.607/1.500,22/1.500/0.000,38/0.000/1.000,52/0.001/1.001,55.5/0.174/1.174,59/0.204/1.204,62.5/0.209/1.209,66/0.210/1.210,69.5/0.210/1.210,73/0.210/1.210,76.5/0.210/1.210,80/0.210/1.210,83.5/0.210/1.210,87/0.210/1.210,90.5/0.210/1.210,94/0.210/1.210,97.5/0.210/1.210}{\draw[very thick,blue!60!black] (\s,0.6*\p) -- (\s,0.6*\q);}
\foreach \s in {22}{\draw[->,thick,red!70!black] (\s,0.98) -- (\s,1.35);}
\end{scope}
\foreach \x in {0,50,100}{\draw[gray!60!black] (\x,-0.06) -- (\x,0.06) node[below,font=\scriptsize] {\x\,ms};}
\end{tikzpicture}
\caption{Une même entrée, deux destinataires différents. Chaque impulsion fait monter la tension de $w$, après quoi elle retombe, comme dans le modèle de neurone~\eqref{eq:glutneuron}; les flèches rouges sont les impulsions du destinataire. Le destinataire lent ($\tau=10\ms$, $\theta=2.5w$) se déclenche là où les impulsions sont \emph{nombreuses} et ne remarque pas la paire. Le rapide ($\tau=2\ms$, $\theta=1.5w$) ne se déclenche que pour une paire dans la fenêtre~\eqref{eq:window}, et laisse passer un flux fréquent mais sans coïncidences.}
\label{fig:reader}
\end{figure}

Les mesures montrent que, dans le cortex actif, la conductance de la membrane est plusieurs fois plus grande qu'au repos~\cite{Destexhe2003}. La constante $\tau$ y est donc plus courte, et les neurones du cortex en régime de travail sont plus proches de détecteurs de coïncidences que d'intégrateurs.

\begin{keyblock}
\begin{proposition}[Le destinataire fixe le code]
\label{prop:reader}
Une même suite d'impulsions porte une fréquence pour un destinataire lent, des instants pour un destinataire rapide, et, pour un volume dans lequel le neurotransmetteur est libéré, un niveau de concentration lissé sur des secondes (chapitre~\ref{sec:serocomputer}). Le code utilisé n'est pas fixé par l'expéditeur, mais par la constante de temps du destinataire, et l'inhibition l'ajuste.
\end{proposition}
\end{keyblock}
\section{Points et traits: la synapse comme filtre}

Jusqu'ici, notre synapse était un fil avec un poids. En réalité, sa réponse dépend des impulsions précédentes, et cela donne à la langue des neurones un second signe.

\emph{La dépression: la synapse transmet les changements.} Chaque libération consomme une fraction $U$ de la réserve de neurotransmetteur prête à être libérée~\cite{Tsodyks1997}, et la réserve se reconstitue avec une constante de temps $\tau_{\text{rec}}$ de l'ordre de centaines de millisecondes. À la fréquence $r$, l'amplitude de la réponse à une impulsion s'établit approximativement au niveau $A(r)=A_0/(1+U r\tau_{\text{rec}})$, où $A_0$ est la réponse d'une synapse reposée. Le courant que reçoit le destinataire par unité de temps vaut
\begin{empheq}[box=\keybox]{equation}
r\,A(r) \;=\; \frac{A_0\,r}{1+U r\,\tau_{\text{rec}}}
\;\xrightarrow[r\to\infty]{}\; \frac{A_0}{U\tau_{\text{rec}}} .
\label{eq:depression}
\end{empheq}
Pour $U=0.5$ et $\tau_{\text{rec}}=0.8\,\text{s}$, la saturation commence dès $1/(U\tau_{\text{rec}})=2.5$ impulsions par seconde. Si la fréquence passe de cinq à vingt, le courant établi n'augmente que d'un tiers. En revanche, les premières impulsions à la nouvelle fréquence arrivent tant que la réserve est encore l'ancienne, et le courant bondit un instant d'un facteur quatre, puis retombe en $\tau_{\text{rec}}/(1+Ur\tau_{\text{rec}})\approx90\ms$ (fig.~\ref{fig:depression}).

Une telle synapse ne transmet pas la fréquence, mais sa \emph{variation}, et au fond la variation relative $\Delta r/r$~\cite{Abbott1997depression}. Pour un électronicien, c'est un filtre passe-haut placé directement dans le fil.

\begin{figure}[t]
\centering
\begin{tikzpicture}[>=Stealth,x=12.5cm,y=1cm]
% input rate
\draw[->,gray!70!black] (0,2.6) -- (0.9,2.6) node[right,black] {\small $t$};
\draw[->,gray!70!black] (0,2.6) -- (0,4.3) node[above,black] {\small $r$};
\draw[very thick,blue!60!black] (0,2.9) -- (0.2,2.9) -- (0.2,3.8) -- (0.8,3.8);
\node[left,font=\scriptsize] at (0,2.9) {$5$};
\node[left,font=\scriptsize] at (0,3.8) {$20$};
\node[right,font=\small,blue!60!black] at (0.4,4.1) {fréquence à l'entrée de la synapse};
% output current
\draw[->,gray!70!black] (0,0) -- (0.9,0) node[right,black] {\small $t$};
\draw[->,gray!70!black] (0,0) -- (0,2.45) node[left,black] {\small $rA$};
\draw[densely dashed,gray!60!black] (0,0.667) -- (0.8,0.667);
\draw[very thick,red!70!black] (0,0.5) -- (0.2,0.5) -- (0.2,2.0)
   plot[domain=0.2:0.8,samples=60] (\x,{0.667+(2.0-0.667)*exp(-(\x-0.2)/0.089)});
\node[left,font=\scriptsize] at (0,0.5) {$1.7$};
\node[left,font=\scriptsize] at (0,2.0) {$6.7$};
\node[right,font=\scriptsize] at (0.8,0.667) {$2.2$};
\node[right,font=\small,red!70!black] at (0.28,1.6) {courant reçu: bond, puis retombée en $\approx90\ms$};
\draw[gray!60!black] (0.2,-0.08) -- (0.2,0.08); \node[below,font=\scriptsize] at (0.2,0) {$0.2\,\text{s}$};
\draw[gray!60!black] (0.8,-0.08) -- (0.8,0.08); \node[below,font=\scriptsize] at (0.8,0) {$0.8\,\text{s}$};
\end{tikzpicture}
\caption{Synapse à dépression d'après la formule~\eqref{eq:depression}, pour $U=0.5$, $\tau_{\text{rec}}=0.8\,\text{s}$; courant en unités de $A_0$ par seconde; en tirets, le courant établi: il n'est passé que de $1.7$ à $2.2$, mais a bondi à $6.7$ au moment du saut. La synapse informe le destinataire que la fréquence a changé, non de sa valeur.}
\label{fig:depression}
\end{figure}

\emph{La facilitation: la bouffée comme trait.} D'autres synapses ont une probabilité de libération faible, mais qui augmente après chaque impulsion pendant des dizaines de millisecondes. Une impulsion isolée se perd le plus souvent à travers une telle synapse. Mais une \emph{bouffée}, quelques impulsions à la suite à quelques millisecondes d'intervalle, passe presque à coup sûr. Même sans facilitation, la probabilité que les $k$ impulsions d'une bouffée soient toutes perdues vaut
\begin{empheq}[box=\keybox]{equation}
P_{\text{perte}} \;=\; (1-p)^k ,
\label{eq:burst}
\end{empheq}
pour $p=0.2$, c'est $0.8$ pour une impulsion isolée et $0.41$ pour une bouffée de quatre; la facilitation la réduit encore davantage. Le neurone acquiert ainsi un second signe: l'impulsion isolée est un point, la bouffée un trait. Une synapse à dépression transmet le mieux la première impulsion après une pause; une synapse à facilitation laisse passer les bouffées et étouffe les points isolés. Que les bouffées soient transmises de façon plus fiable est mesuré; que le cerveau s'en serve réellement comme d'un signe à part est une hypothèse plausible~\cite{Lisman1997}.

\section{Comment parlent les neurones \nt{GABA}}

Les plus nombreux d'entre eux dans le cortex sont les neurones rapides~\cite{Tremblay2016}. Leurs impulsions sont plus courtes que celles des neurones \nt{GLUT}, et ils peuvent les émettre régulièrement et longtemps, à des centaines par seconde, presque sans se fatiguer. Leurs synapses sont plus fiables: la connexion avec chaque destinataire comporte de nombreux sites de libération, et une impulsion ne se perd presque jamais. Leurs sorties se posent près de l'endroit où naît l'impulsion du destinataire, si bien qu'ils ne commandent pas la façon dont le destinataire additionne ses entrées, mais le fait qu'il se déclenche ou non, et quand.

Eux-mêmes reçoivent l'excitation de nombreux neurones \nt{GLUT} voisins et signalent l'activité totale alentour: exactement ce qu'il faut pour la division du chapitre~\ref{sec:gaba}. Enfin, les neurones \nt{GABA} rapides sont connectés entre eux et se déclenchent facilement ensemble; ce sont eux qui fixent le rythme local dont il a été question plus haut~\cite{Cardin2009}.

Dans le langage du Morse, les neurones \nt{GABA} ne sont pas responsables des lettres, mais des \emph{pauses}. Ils découpent le flux continu des impulsions \nt{GLUT} en fenêtres à l'intérieur desquelles le destinataire peut se déclencher, rétrécissent les fenêtres de coïncidence et baissent le volume de tout le flux quand il devient trop fort.

\begin{keyblock}
\begin{proposition}[Partage des rôles]
\label{prop:punctuation}
Les neurones \nt{GLUT} portent le contenu du message: \emph{quoi} et \emph{combien}. Les neurones \nt{GABA} rapides en portent le balisage: \emph{quand} on peut parler et \emph{à quel volume}; une autre classe encore, en inhibant les inhibiteurs, décide \emph{pour qui} le portillon est ouvert. Certaines parties de ce partage sont mesurées; comme règle générale, c'est une hypothèse de travail.
\end{proposition}
\end{keyblock}

Il existe aussi d'autres neurones \nt{GABA}, plus lents, dont les sorties se posent sur des entrées particulières du destinataire. Ils fonctionnent comme le veto~\eqref{eq:veto} du chapitre~\ref{sec:gaba}: ils ferment un flux d'impulsions \nt{GLUT} sans toucher aux autres.

La division en neurones \nt{GABA} rapides et lents est un tableau ancien, et incomplet. On distingue aujourd'hui dans le cortex au moins trois grandes classes~\cite{Tremblay2016}, et la troisième est étonnante: elle n'inhibe pas les neurones \nt{GLUT}, mais d'autres neurones \nt{GABA}, avant tout ceux qui ferment les entrées~\cite{Pfeffer2013}. L'inhibition de l'inhibition est une double négation. Un tel neurone \emph{ouvre} le portillon sans envoyer au destinataire une seule impulsion excitatrice. Il est activé par des signaux venus d'autres aires du cortex et par les canaux de diffusion, si bien qu'il fonctionne comme une entrée d'autorisation pour toute une portion du réseau: tant qu'il est actif, les vetos sont levés et le flux passe~\cite{Pi2013}. L'annexe~\ref{sec:othermediators} appelle ce procédé la désinhibition.
\section{Une langue économe}

La dernière chose que l'on sait avec certitude de la langue des neurones, c'est qu'elle coûte cher. Une impulsion, avec le travail des synapses qu'elle provoque, coûte environ $10^9$ molécules d'ATP (estimation pour le cortex des rongeurs~\cite{Attwell2001}), et une molécule fournit environ $10^{-19}$ joule. Une impulsion coûte donc de l'ordre de $E_{\text{imp}}\sim10^{-10}$ joule. Le cerveau entier consomme environ $20$ watts; si la moitié va à la transmission des signaux, $P\sim10$ watts, la fréquence moyenne d'un neurone vaut
\begin{empheq}[box=\keybox]{equation}
\bar r \;\approx\; \frac{P}{N\,E_{\text{imp}}}
\;\sim\; \frac{10\ \text{W}}{8.6\cdot10^{10}\cdot10^{-10}\ \text{J}}
\;\approx\; 1\ \text{impulsion par seconde}.
\label{eq:energy}
\end{empheq}
Des estimations plus détaillées pour le cortex donnent encore moins~\cite{Lennie2003}. Le code du cerveau est forcément \emph{parcimonieux}: seule une petite fraction des neurones peut travailler vite.

On voit d'après~\eqref{eq:spikeentropy} que ce n'est pas si grave. Avec une précision de $1\ms$, une impulsion d'un neurone qui travaille à $100$ par seconde ne porte pas plus de cinq bits, et celle d'un neurone qui travaille une fois par seconde, jusqu'à onze. Si l'on paie à l'impulsion, il est avantageux de parler rarement. Le cortex échange d'ailleurs bien de la précision contre de l'énergie: quand la nourriture manque, les neurones conservent leur fréquence d'impulsions, mais dépensent moins en courants synaptiques, et leurs réponses deviennent moins précises~\cite{Padamsey2022}.

En Morse, les lettres fréquentes sont courtes: la lettre la plus fréquente de l'anglais, E, est un seul point. Le code des neurones, pour autant qu'on le sache, suit la même règle sous une autre forme: ce qui est attendu coûte peu d'impulsions~\cite{Barlow1961}. Face à un stimulus constant, un neurone réduit peu à peu sa fréquence, les capteurs du chapitre~\ref{sec:glutcomputer} répondent aux changements et non aux niveaux, et les neurones \nt{DOPA} du chapitre~\ref{sec:neurontypes} ne signalent que l'erreur de prédiction de la récompense.

\begin{keyblock}
\begin{proposition}[Économie]
\label{prop:economy}
L'énergie limite la fréquence moyenne d'un neurone à environ une impulsion par seconde. La langue des neurones \nt{GLUT} est donc parcimonieuse et dépense ses impulsions pour les nouvelles: pour ce qui a changé ou n'avait pas été prédit. La limite énergétique est mesurée; que le code du cerveau soit ajusté pour maximiser le nombre de bits par joule est une hypothèse bien fondée.
\end{proposition}
\end{keyblock}

\section{Bilan}

\begin{table}[t]
\centering
\footnotesize
\setlength{\tabcolsep}{4pt}
\renewcommand{\arraystretch}{1.15}
\begin{tabular}{>{\raggedright}p{2.6cm}>{\raggedright}p{2.3cm}>{\raggedright}p{3.0cm}>{\raggedright}p{2.0cm}>{\raggedright\arraybackslash}p{3.6cm}}
\toprule
Mode d'écriture & Ce qu'il porte & Qui le lit & Durée d'un message & Ce que l'on sait \\
\midrule
numéro du neurone déclenché (code de position) & quelle valeur & tout destinataire; détecteurs de coïncidences & un retard & établi chez les neurones sensoriels et dans la localisation des sons \\
fréquence d'un neurone & une grandeur & intégrateur à grande $\tau$; volume (ch.~\ref{sec:serocomputer}) & centaines de ms à secondes & établi; trop lent pour une étape de traitement de $10$--$20\ms$ \\
fréquence d'un groupe & une grandeur & neurone à milliers d'entrées & $10$--$50\ms$ & établi; gain limité par les corrélations~\eqref{eq:correlated} \\
instant de la première impulsion, ordre & une grandeur, un ordre & détecteurs de coïncidences à retards & un retard & établi à l'entrée du système nerveux; débattu dans le cortex \\
phase dans un rythme local & une grandeur, une position & neurones à rythme commun & cycle de $12$--$50\ms$ ou plus lent & observé dans certaines régions; rôle général: hypothèse \\
bouffée (\og trait\fg) & \og important\fg: livraison fiable & synapse à facilitation & $\sim10\ms$ & fiabilité mesurée; signe à part: hypothèse \\
variation de fréquence & une nouvelle & synapse à dépression & $\sim0.1\,\text{s}$ & établi \\
impulsions des neurones \nt{GABA} rapides & quand et à quel volume & tous les neurones \nt{GLUT} voisins & $1$--$30\ms$ & rythme et division mesurés; \og ponctuation\fg{} comme règle: hypothèse \\
\bottomrule
\end{tabular}
\caption{Les façons dont les neurones \nt{GLUT} et \nt{GABA} écrivent des messages avec des impulsions. La colonne de droite sépare ce qui est mesuré de ce qui est supposé.}
\label{tab:code}
\end{table}

Le tableau~\ref{tab:code} rassemble ce que nous savons; on y voit aussi ce que nous ne savons pas. Nous ne savons pas dans quelle mesure le cortex utilise l'instant précis des impulsions, et pas seulement les fréquences de groupes. Nous ne savons pas si les neurones ont des \og mots\fg, des combinaisons stables d'impulsions de nombreux neurones qui se lisent comme un tout. Et nous ne savons pas si la langue est la même dans les différentes parties du cerveau. On peut apprendre le Morse en une soirée, parce que sa table est connue. La table de la langue des neurones reste à dresser, et ce seront très probablement des ingénieurs des télécommunications qui la dresseront.

\section{Exercices}

\begin{exercise}[\lvlA]
\label{ex:04:1}
\emph{La capacité en chiffres.} $\Delta t=1\ms$. (a)~Combien de bits par joule fournit un neurone de fréquence $1$ impulsion par seconde, au coût d'impulsion de~\eqref{eq:energy}? (b)~Pour $r=10$, combien de fois moins que la borne~\eqref{eq:spikeentropy} tire un destinataire qui ne fait que compter les impulsions sur une seconde?
\end{exercise}

\begin{exercise}[\lvlB]
\label{ex:04:2}
\emph{Fréquence optimale.} Un neurone dépense la puissance $P_0+rE_{\text{imp}}$, où $P_0$ est l'autre moitié des $20$~W, répartie sur tous les neurones, $E_{\text{imp}}=10^{-10}$~J. Pour quelle fréquence $r$ le nombre de bits par joule $R_{\max}(r)/(P_0+rE_{\text{imp}})$ selon~\eqref{eq:spikeentropy} avec $\Delta t=1\ms$ est-il maximal?
\end{exercise}

\begin{exercise}[\lvlB]
\label{ex:04:3}
\emph{Quelle corrélation est nuisible.} $N=1000$ neurones, variance $\sigma^2=1$, corrélation par paire $c=0.1$. Calculez l'information de Fisher $\mathcal I=f'^{\mathsf T}\Sigma^{-1}f'$ ($\Sigma$ est la matrice de covariance) pour des pentes égales $f'=\mathbf 1$ et pour des pentes $\pm1$ avec $\mathbf 1^{\mathsf T}f'=0$. Dans quel rapport diffèrent-elles?
\end{exercise}

\begin{exercise}[\lvlB]
\label{ex:04:4}
\emph{Modélisation: détecteur de coïncidences.} Le neurone~\eqref{eq:glutneuron} reçoit $1000$ entrées poissoniennes à $5$ impulsions par seconde, $w=1$; le seuil $\theta$ est tel que la tension libre le franchit une fois par seconde. Trouvez par simulation combien d'entrées simultanées $m$ déclenchent le neurone avec la probabilité $1/2$ pour $\tau=10$ et $2\ms$.
\end{exercise}

\begin{exercise}[\lvlB]
\label{ex:04:5}
\emph{Conception: détecteur de changements relatifs.} Choisissez une synapse~\eqref{eq:depression}: le courant établi à $40$ impulsions par seconde ne dépasse pas de plus de $10\,\%$ celui à $10$, et après un saut à $40$ le courant atteint son nouveau niveau avec une constante de temps d'au plus $50\ms$. Quel est le plus petit $\tau_{\text{rec}}$, et pour quel $U$?
\end{exercise}

\begin{exercise}[\lvlA]
\label{ex:04:6}
\emph{Instant de la première impulsion et bouffées.} (a)~D'après~\eqref{eq:latency}, trouvez l'entrée pour laquelle la première impulsion arrive exactement au bout d'une constante de temps. De quoi dépend-elle? (b)~Avec une probabilité de libération $p=0.2$, combien d'impulsions faut-il dans une bouffée pour que la probabilité de toutes les perdre soit inférieure à $5\,\%$?
\end{exercise}

\begin{exercise}[\lvlC]
\label{ex:04:7}
\emph{Recherche: combien de bits dans la disposition.} Le neurone est fait de cases indépendantes~\eqref{eq:spikeentropy}, $r=10$, $\Delta t=1\ms$. (a)~Décomposez l'entropie d'un \og mot\fg{} de $20$ cases en entropie du nombre d'impulsions et reste. (b)~Simulez l'estimation de l'entropie d'après les fréquences des mots sur $N=30$ et $10^4$ enregistrements. De combien est-elle sous-estimée?
\end{exercise}

\chapter{Ce que calculent les neurones \nt{SERO}}
\label{sec:serocomputer}
\begin{chaptergoals}
\item comment un pacemaker à récepteur inhibiteur fait NON et NON-OU avec un seul neurone \nt{SERO};
\item pourquoi la libération en volume ne laisse que quelques fils indépendants, fixés par $\ell\sim\sqrt{D\tau}$;
\item comment une concentration lisse les impulsions en un niveau, que l'auto-inhibition pourrait tenir;
\item pourquoi \nt{SERO} est candidat à l'horizon de planification, et comment l'horizon fixe la patience.
\end{chaptergoals}

\section{Le premier canal de diffusion}

Jusqu'ici, tous les neurones du livre parlaient par le bus d'adresses: \nt{GLUT} et \nt{GABA} envoyaient leurs impulsions à des destinataires précis, et nous avons établi ce que calcule un tel réseau et dans quelle langue il parle. Passons maintenant au second bus de la section~\ref{sec:buses}, le bus de diffusion. Son premier canal est \nt{SERO}. Comme avant, nous ne nous autorisons que ce qui existe dans un organisme vivant.

Ce chapitre introduit trois dispositifs dont se serviront ensuite tous les canaux de diffusion: un neurone qui, alimenté par un apport de fond régulier, fonctionne tout seul tant qu'on ne l'arrête pas; un fil qui est en réalité un volume de tissu; une concentration lente au lieu d'impulsions isolées. \nt{DOPA}, \nt{NORA} et \nt{ACET}, aux chapitres~\ref{sec:dofa}--\ref{sec:acet}, seront construits avec les mêmes pièces.

Du point de vue de l'ingénieur, le neurone \nt{SERO} diffère sur quatre points importants.

\emph{Premièrement: le signe est choisi par le destinataire.} La sérotonine a des récepteurs des deux signes, et la règle~\eqref{eq:dalesign} ne s'applique pas ici: le signe n'est pas fixé par l'émetteur mais par le récepteur du destinataire (proposition~\ref{prop:receptor})~\cite{BarnesSharp1999}.

\emph{Deuxièmement: le neurone \nt{SERO} fonctionne tout seul si le fond est régulier.} À l'éveil, il émet régulièrement plusieurs impulsions par seconde, sans avoir besoin d'une poussée pour chacune: c'est un stimulateur (pacemaker). Mais il lui faut un apport de fond constant. Dans une tranche de cerveau, où cet apport manque, la plupart de ces cellules se taisent; le rythme régulier revient si l'on applique de la noradrénaline par les récepteurs $\alpha_1$, et leur blocage l'éteint~\cite{VandermaelenAghajanian1983}. Dans l'organisme, ce fond vient de \nt{NORA}. Pour le schéma, c'est un oscillateur qui a besoin d'une alimentation: tant qu'elle est là, il fonctionne tout seul, jusqu'à ce qu'une inhibition l'arrête.

\emph{Troisièmement: il est lent.} Presque tous les récepteurs de la sérotonine sont métabotropes: ils n'ouvrent pas eux-mêmes un canal, mais déclenchent une chaîne de réactions dans la cellule. La réponse est lente: le courant inhibiteur qui passe par le récepteur principal monte et retombe en une seconde environ~\cite{CourtneyFord2016}, des dizaines de fois plus longtemps que la réponse d'une synapse rapide \nt{GLUT}.

\emph{Quatrièmement: il libère son neurotransmetteur non dans un fil, mais dans un volume.} Dans les régions où projettent les neurones \nt{SERO}, la sérotonine est souvent libérée non dans une fente étroite, mais dans l'espace extracellulaire, où elle se répand~\cite{Bunin1998}. Le destinataire perçoit une concentration et ne peut pas savoir de quel émetteur vient une molécule.

À ces échelles de temps, les impulsions isolées ne comptent plus; nous décrirons donc le réseau par des fréquences $r_j$ et des concentrations. La concentration $c$ de sérotonine dans le volume où libèrent les neurones $j$ croît avec la libération et décroît parce que le neurotransmetteur est recapté:
\begin{empheq}[box=\keybox]{equation}
\tau_c\,\frac{\dd c}{\dd t} \;=\; -c \;+\; \kappa\sum_j r_j ,
\qquad
r_i \;=\; f_i\bigl(b_i + s_i\,g\,c\bigr), \quad s_i=\pm1 .
\label{eq:seroneuron}
\end{empheq}
C'est une manière de plus de lire un train d'impulsions, qui complète la proposition~\ref{prop:reader}: le volume ne voit ni les impulsions isolées ni leur ordre, seulement la fréquence moyennée sur le temps $\tau_c$. La première équation est le même circuit RC que la membrane: $\tau_c$ est la durée de vie du neurotransmetteur dans le volume (elle n'a pas été mesurée directement, et nous la prenons de l'ordre de la seconde); $\kappa$ est la quantité de neurotransmetteur apportée par une impulsion. La seconde décrit le destinataire: $b_i$ est son apport propre, positif pour un pacemaker (il inclut l'entrée de fond régulière venue de \nt{NORA}); $g$ est la sensibilité du récepteur; le signe $s_i$ est choisi par le destinataire; $f_i$ est une caractéristique de transfert non décroissante.

\section{NON avec un seul neurone}

Prenons un pacemaker à récepteur inhibiteur, $s=-1$. Tant qu'il n'y a pas de sérotonine autour, il fonctionne sur son apport propre $b$. Quand la sérotonine apparaît, l'apport diminue de $gc$, et si $gc>b$, le neurone se tait. C'est un NON: la sortie est active quand l'entrée ne l'est pas. Il a coûté un seul neurone (fig.~\ref{fig:serogates}). La proposition~\ref{prop:monotone} ne s'applique pas ici: elle supposait des poids positifs.

Si le même pacemaker reçoit la sérotonine de deux volumes différents, il se tait dès qu'il y en a dans au moins l'un d'eux. C'est un NON-OU, et avec des NON-OU on peut construire n'importe quelle fonction logique. Le code à deux fils du réseau \nt{GLUT} est ici inutile: l'absence de signal est calculée par des neurones qui fonctionnent tant qu'on ne les arrête pas.

\begin{figure}[t]
\centering
\begin{tikzpicture}[>=Stealth,
  nrn/.style={circle,draw,very thick,fill=blue!6,minimum size=0.95cm,inner sep=0pt},
  pace/.style={nrn,double,double distance=1.2pt},
  inh/.style={-{Bar[width=7pt]},thick,red!70!black}]
\node[pace] (N1) at (0,0) {};
\draw[inh] (-1.6,0) node[left,black] {$c$} -- (N1);
\draw[->,very thick,red!70!black] (N1) -- (1.3,0) node[right,black] {$\bar c$};
\node[below=0.55cm] at (0,0) {\small NON};
\node[pace] (N2) at (5,0) {};
\draw[inh] (3.4,0.45) node[left,black] {$c_1$} -- (N2);
\draw[inh] (3.4,-0.45) node[left,black] {$c_2$} -- (N2);
\draw[->,very thick,red!70!black] (N2) -- (6.3,0) node[right,black] {$\overline{c_1\vee c_2}$};
\node[below=0.55cm] at (5,0) {\small NON-OU};
\end{tikzpicture}
\caption{Logique à neurones \nt{SERO}. Double cercle: pacemaker; avec un apport de fond régulier, il fonctionne tout seul tant qu'on ne l'inhibe pas. Trait barré: récepteur inhibiteur, $s=-1$; $c$, $c_1$, $c_2$ sont les concentrations de sérotonine dans les volumes autour du destinataire. À gauche, NON; à droite, NON-OU.}
\label{fig:serogates}
\end{figure}

\begin{keyblock}
\begin{proposition}[Complétude de la logique \nt{SERO}]
\label{prop:serocomplete}
Si le réseau contient des pacemakers à récepteurs inhibiteurs et que les libérations dans des volumes différents ne se mélangent pas, le réseau~\eqref{eq:seroneuron} peut calculer n'importe quelle fonction logique, et chaque bit passe par un seul fil.
\end{proposition}
\end{keyblock}

Sur le plan logique, le réseau \nt{SERO} est plus puissant que \nt{GLUT}, mais la condition \og les libérations dans des volumes différents ne se mélangent pas\fg{} n'est pas remplie dans le cerveau (section~\ref{sec:volume}).

\section{Rétroaction négative}

Les neurones à sérotonine portent des récepteurs inhibiteurs sur eux-mêmes~\cite{Sprouse1987}. La sérotonine qu'ils ont libérée leur revient et les freine. Pour un ingénieur, le schéma est familier: un amplificateur à rétroaction négative (fig.~\ref{fig:serofeedback}).

Ce qui est mesuré et ce qui relève du modèle. Dans le noyau du raphé lui-même, où se trouvent les neurones \nt{SERO}, la sérotonine inhibe les voisins non par un volume commun, mais point à point, par des synapses: le transporteur l'élimine vite et l'empêche de s'accumuler hors de la fente. Une libération isolée ne produit qu'une courte pause dans la décharge, et la fréquence moyenne n'en est pas changée~\cite{CourtneyFord2016}. La boucle ci-dessous, fermée par un volume commun, est donc un modèle: nous calculons ce que ferait une telle rétroaction si elle agissait au niveau des fréquences moyennes.

\begin{figure}[t]
\centering
\begin{tikzpicture}[>=Stealth,
  nrn/.style={circle,draw,very thick,fill=blue!6,minimum size=0.95cm,inner sep=0pt},
  pace/.style={nrn,double,double distance=1.2pt},
  blk/.style={draw,very thick,rounded corners=2pt,fill=orange!10,minimum height=0.9cm,minimum width=2.2cm,align=center,font=\small},
  inh/.style={-{Bar[width=7pt]},thick,red!70!black}]
\node[pace] (N) at (0,0) {};
\node[blk] (V) at (3.6,0) {volume\\$\tau_c$};
\draw[->,very thick] (N) -- node[above] {\small $r$} (V);
\draw[->,very thick] (V) -- (6.0,0) node[right] {\small $c$ vers tous les destinataires};
\draw[inh] (V.south) -- ++(0,-0.9) -| node[pos=0.25,below] {\small $-g\,c$} (N.south);
\draw[->,thick,blue!60!black] (-1.7,0) node[left,black] {\small $b$} -- (N);
\end{tikzpicture}
\caption{Neurone \nt{SERO} rétrocontrôlé par son propre neurotransmetteur: la concentration $c$ va à tous les destinataires et inhibe le neurone lui-même. Dans le modèle, c'est un régulateur de niveau; dans le cerveau, ce rôle de la boucle est une hypothèse.}
\label{fig:serofeedback}
\end{figure}

Calculons ce que fait cette boucle. Supposons la caractéristique de transfert linéaire, $f(u)=au$ pour $u>0$. À l'équilibre, $c=\kappa r$ et $r=a(b-gc)$. En substituant l'une dans l'autre, on obtient
\begin{empheq}[box=\keybox]{equation}
r \;=\; \frac{a\,b}{1+G}, \qquad G \;=\; a\,g\,\kappa .
\label{eq:feedback}
\end{empheq}
La grandeur $G$ est le gain de boucle. C'est la formule de tout amplificateur à rétroaction négative, et elle donne les deux mêmes avantages.

\emph{Robustesse à la dispersion.} Supposons que l'excitabilité $a$ du neurone change d'une fraction $\varepsilon$. Sans rétroaction, la fréquence changerait de la même fraction. Avec rétroaction, pour $G\gg1$, la fréquence $r\approx b/(g\kappa)$ ne dépend presque plus de $a$: sa variation est $1+G$ fois plus petite. Pour $G=9$, la dispersion est divisée par dix.

\emph{Rapidité.} Portons $r=a(b-gc)$ dans la première équation~\eqref{eq:seroneuron}: on obtient $\tau_c\,\dd c/\dd t=-(1+G)\,c+\kappa ab$. La concentration rejoint son niveau avec la constante de temps $\tau_c/(1+G)$, soit $1+G$ fois plus vite que sans rétroaction.

Voici le premier calcul que pourrait effectuer le système \nt{SERO}: maintenir le niveau global de sérotonine dans le cerveau à une valeur de consigne, comme un thermostat maintient la température. C'est une hypothèse: expérimentalement, l'autoinhibition n'apparaît pour l'instant que sous forme de courtes pauses, qui ne changent pas la fréquence moyenne.

\section{Des impulsions au niveau}

Le deuxième calcul est caché dans la première équation~\eqref{eq:seroneuron}. Les neurones émettent des impulsions isolées, mais le destinataire perçoit une concentration qui les lisse sur le temps $\tau_c$. Si la fréquence de la source varie, la concentration la suit avec retard:
\begin{equation}
c(t) \;=\; \frac{\kappa}{\tau_c}\int_{-\infty}^{t} e^{-(t-t')/\tau_c}\,\sum_j r_j(t')\,\dd t' .
\label{eq:lowpass}
\end{equation}
C'est une moyenne glissante de l'activité des sources sur les derniers $\tau_c$, un filtre passe-bas (fig.~\ref{fig:lowpass}). Le convertisseur numérique-analogique le plus simple fait passer des impulsions dans un circuit RC et transforme ainsi leur fréquence en niveau. Le système \nt{SERO} fait de même avec des centaines de milliers de neurones à la fois.

Cette grandeur est aussi une mémoire. Une fois les sources silencieuses, la concentration persiste encore un temps de l'ordre de $\tau_c$. Ce n'est pas un bit, mais une trace qui s'efface doucement.

\begin{figure}[t]
\centering
\begin{tikzpicture}[>=Stealth,x=1.45cm,y=1cm]
\draw[gray!70!black] (0,3.2) -- (8.1,3.2);
\node[left,font=\small] at (-0.05,3.4) {impulsions};
\node[above,font=\scriptsize,gray!40!black] at (1,3.65) {$2$ par seconde};
\node[above,font=\scriptsize,gray!40!black] at (3.5,3.65) {$8$ par seconde};
\node[above,font=\scriptsize,gray!40!black] at (6.5,3.65) {$2$ par seconde};
\draw[->,gray!70!black] (0,0) -- (8.3,0) node[right,black] {\small $t$};
\draw[->,gray!70!black] (0,0) -- (0,2.75) node[left,black] {\small $c$};
\foreach \s in {0.136,0.529,0.760,1.223,1.714,1.748,1.755,2.663,2.701,2.734,3.414,3.493,3.719,3.800,3.928,3.948,4.074,4.327,4.420,4.589,4.728,4.736,4.914,5.025,5.205,5.220,6.224,6.544,7.178}{\draw[thick,blue!60!black] (\s,3.2) -- (\s,3.6);}
\foreach \a/\b/\v in {0.000/0.136/0.000,0.136/0.529/1.000,0.529/0.760/1.675,0.760/1.223/2.330,1.223/1.714/2.466,1.714/1.748/2.509,1.748/1.755/3.425,1.755/2.663/4.402,2.663/2.701/2.775,2.701/2.734/3.672,2.734/3.414/4.553,3.414/3.493/3.306,3.493/3.719/4.055,3.719/3.800/4.235,3.800/3.928/4.905,3.928/3.948/5.316,3.948/4.074/6.211,4.074/4.327/6.476,4.327/4.420/6.028,4.420/4.589/6.493,4.589/4.728/6.483,4.728/4.736/6.642,4.736/4.914/7.589,4.914/5.025/7.351,5.025/5.205/7.579,5.205/5.220/7.331,5.220/6.224/8.221,6.224/6.544/4.012,6.544/7.178/3.914,7.178/8.000/3.076}{\draw[very thick,orange!85!black] plot[domain=\a:\b,samples=10] (\x,{0.25*\v*exp(-(\x-\a))});}
\foreach \s/\p/\q in {0.136/0.000/1.000,0.529/0.675/1.675,0.760/1.330/2.330,1.223/1.466/2.466,1.714/1.509/2.509,1.748/2.425/3.425,1.755/3.402/4.402,2.663/1.775/2.775,2.701/2.672/3.672,2.734/3.553/4.553,3.414/2.306/3.306,3.493/3.055/4.055,3.719/3.235/4.235,3.800/3.905/4.905,3.928/4.316/5.316,3.948/5.211/6.211,4.074/5.476/6.476,4.327/5.028/6.028,4.420/5.493/6.493,4.589/5.483/6.483,4.728/5.642/6.642,4.736/6.589/7.589,4.914/6.351/7.351,5.025/6.579/7.579,5.205/6.331/7.331,5.220/7.221/8.221,6.224/3.012/4.012,6.544/2.914/3.914,7.178/2.076/3.076}{\draw[very thick,orange!85!black] (\s,0.25*\p) -- (\s,0.25*\q);}
\draw[thick,densely dashed,black] plot[domain=0:2,samples=30] (\x,{0.25*2*(1-exp(-\x))})
  plot[domain=2:5,samples=40] (\x,{0.25*(8-6.271*exp(-(\x-2)))})
  plot[domain=5:8,samples=40] (\x,{0.25*(2+5.688*exp(-(\x-5)))});
\foreach \x in {0,2,5,8}{\draw[gray!60!black] (\x,-0.05) -- (\x,0.05) node[below=3pt,font=\scriptsize] {\x\,s};}
\end{tikzpicture}
\caption{Des impulsions au niveau. En haut, les impulsions des neurones \nt{SERO}: deux secondes de décharge rare, trois secondes de décharge dense, puis de nouveau rare. En bas, la concentration de sérotonine~\eqref{eq:lowpass} pour $\tau_c=1\,\text{s}$. En tirets, la même chose si les impulsions arrivaient parfaitement régulièrement.}
\label{fig:lowpass}
\end{figure}

\section{Le fil, c'est le volume}
\label{sec:volume}

Revenons à la condition de la proposition~\ref{prop:serocomplete}. La sérotonine libérée à proximité par différents émetteurs s'additionne en une seule concentration.

\emph{Le fil n'est pas ici une fibre, mais un volume de tissu.} Deux signaux ne restent distincts que s'ils sont libérés dans des régions plus éloignées que la distance sur laquelle le neurotransmetteur a le temps de se répandre. En un temps $\tau$, une molécule de coefficient de diffusion $D$ parcourt
\begin{empheq}[box=\keybox]{equation}
\ell \;\sim\; \sqrt{D\,\tau}\, .
\label{eq:diffusionlength}
\end{empheq}
La tortuosité des espaces extracellulaires ralentit la diffusion d'une petite molécule d'un facteur $2.6$ environ~\cite{Sykova2008}, et le coefficient de diffusion de la sérotonine elle-même dans le tissu n'a pas été mesuré; d'après la taille de la molécule, nous l'estimons à quelques unités de $10^{-6}$~cm$^2$/s. Si le signal vit $\tau\sim1$~s (encore une estimation), alors $\ell\sim\sqrt{3\cdot10^{-6}}$~cm $\approx20$~\textmu m. Dans un tel réseau, l'adresse est un \emph{lieu}: le destinataire ne sait de qui vient le signal que par l'endroit où il se trouve lui-même.

Les vrais neurones \nt{SERO} sont construits de telle sorte qu'il ne reste presque rien de ces adresses distinctes. La sortie de chacun est répartie sur d'immenses portions du cerveau: les branches d'une seule cellule atteignent de nombreuses régions éloignées les unes des autres~\cite{GagnonParent2014}. On estime à environ $10^5$ les sites de libération par cellule (tableau~\ref{tab:types}); ils n'ont pas été comptés directement. Les \og fils\fg{} des différents neurones \nt{SERO} se chevauchent et se mélangent. Ils ne se fondent pourtant pas en un fil unique: les neurones \nt{SERO} se répartissent en quelques sous-systèmes, qui projettent vers des régions différentes et répondent différemment à un même événement~\cite{Ren2018}. Mais ces sous-systèmes se comptent par unités, et non par centaines de milliers.

\begin{keyblock}
\begin{proposition}[Nombre de fils indépendants]
\label{prop:volumewires}
Dans un réseau à transmission volumique, le nombre de signaux indépendants n'est pas limité par le nombre de neurones, mais par le nombre de régions disjointes de taille $\ell\sim\sqrt{D\tau}$ dans lesquelles ils libèrent le neurotransmetteur. Si la sortie de chaque neurone est répartie sur un grand volume, tout le réseau ne transmet que quelques variables communes.
\end{proposition}
\end{keyblock}

Dans le cerveau, il n'y a donc pas de quoi construire les circuits logiques de la proposition~\ref{prop:serocomplete}. En revanche, le régulateur~\eqref{eq:feedback} et le filtre~\eqref{eq:lowpass} n'ont pas besoin de fils distincts: une seule grandeur commune leur suffit. Le filtre découle directement de la libération volumique mesurée dans les régions cibles~\cite{Bunin1998}; le régulateur de niveau est une hypothèse.

\section{L'horizon de planification}
\label{sec:horizon}

À quoi peut servir une seule grandeur commune et lente? Un bon candidat est un paramètre qui doit être le même pour tout le réseau et ne varier qu'en quelques secondes ou minutes. Au chapitre~\ref{sec:neurontypes}, un tel paramètre est déjà apparu: le coefficient $\gamma$ de la formule de l'erreur~\eqref{eq:rpe}, qui dit combien une récompense future vaut moins qu'une récompense immédiate. En temps continu, sa place est prise par l'\emph{horizon} $\tau_\gamma$: une récompense $R$ qu'il faut attendre pendant un temps $D$ vaut aujourd'hui $R\,e^{-D/\tau_\gamma}$.

L'horizon décide s'il vaut la peine d'attendre. Supposons qu'on puisse prendre tout de suite une petite récompense $r_0$ ou attendre une grande récompense $R$. Attendre est avantageux tant que
\begin{empheq}[box=\keybox]{equation}
D \;<\; \tau_\gamma\,\ln\frac{R}{r_0} .
\label{eq:patience}
\end{empheq}
Pour $\tau_\gamma=2\,\text{s}$ et une récompense différée trois fois plus grande, il vaut la peine d'attendre jusqu'à $2.2\,\text{s}$; si l'horizon est deux fois plus long, jusqu'à $4.4\,\text{s}$. La patience est proportionnelle à l'horizon.

La formule~\eqref{eq:patience} repose sur une dépréciation exponentielle. La dépréciation mesurée pour des délais de quelques secondes est plus proche de l'hyperbole $R/(1+D/\tau_\gamma)$: c'est ainsi que décroissent chez le singe aussi bien la valeur d'une récompense différée dans le choix que la réponse des neurones \nt{DOPA} au signal qui l'annonce~\cite{KobayashiSchultz2008}. Pour l'hyperbole, la condition~\eqref{eq:patience} devient $D<\tau_\gamma\,(R/r_0-1)$: la dépendance au rapport des récompenses n'est plus logarithmique mais linéaire, et avec les mêmes nombres la limite d'attente se déplace presque du double. L'hyperbole s'obtient à partir de la même exponentielle si le délai est aléatoire (exercice~\ref{ex:05:7}), et la proportionnalité de la patience à l'échelle de temps est conservée.

Selon une hypothèse, c'est \nt{SERO} qui fixe l'horizon~\cite{Doya2002}. Il a tout ce qu'il faut pour ce rôle: un seul niveau commun pour tout le réseau, qui varie en quelques secondes. Il existe aussi des expériences directes: si l'on active artificiellement les neurones à sérotonine, l'animal attend plus longtemps une récompense différée avant d'abandonner~\cite{Miyazaki2014}, et s'obstine plus longtemps à chercher une récompense là où elle est devenue rare~\cite{Lottem2018}. Comment la concentration de sérotonine se transforme en $\tau_\gamma$ à l'intérieur du réseau, on ne le sait pas. Au chapitre suivant, l'horizon entrera dans l'erreur que diffuse \nt{DOPA}.

\section{Bilan}

\begin{table}[t]
\centering
\small
\begin{tabular}{>{\raggedright}p{3.6cm}>{\raggedright}p{4.3cm}>{\raggedright\arraybackslash}p{4.3cm}}
\toprule
& \nt{GLUT} & \nt{SERO} \\
\midrule
temps de réaction & millisecondes & fraction de seconde à une seconde \\
signe de l'action & seulement $+$ & $+$ et $-$, choisi par le destinataire \\
sans entrée & se tait & fonctionne seul avec un apport de fond régulier \\
adressage & point à point & volume; l'adresse est un lieu \\
NON & seulement par des détecteurs d'\og extinction\fg{} & un seul neurone \\
mémoire & activité auto-entretenue, s'éteint par fatigue & concentration, s'efface en un temps $\tau_c$ \\
calculs principaux & coïncidences, délais, logique, séquences & lissage; régulation de niveau et horizon $\tau_\gamma$ (hypothèses) \\
\bottomrule
\end{tabular}
\caption{Ce que calculent les réseaux faits uniquement de neurones \nt{GLUT} et uniquement de neurones \nt{SERO}.}
\label{tab:glutvssero}
\end{table}

Comme élément logique (tableau~\ref{tab:glutvssero}), le neurone \nt{SERO} est plus riche que \nt{GLUT}, mais comme système de communication c'est une diffusion générale: lente et presque sans adresses. Il ne cherche donc pas à être un processeur; il lisse et diffuse quelques grandeurs communes, celles dont il a été question au chapitre~\ref{sec:neurontypes}, parmi lesquelles, selon l'hypothèse, l'horizon de planification. Le pacemaker, le volume et la concentration lente nous attendent encore trois fois: chez \nt{DOPA}, \nt{NORA} et \nt{ACET}.

\section{Exercices}

\begin{exercise}[\lvlA]
\label{ex:05:1}
\emph{Le fil, c'est le volume.} Prenez pour la sérotonine $D=3\cdot10^{-6}$~cm$^2$/s (section~\ref{sec:volume}). Trouvez par~\eqref{eq:diffusionlength} $\ell$ pour un signal qui vit $1\,\text{s}$, et le temps qu'il faut au neurotransmetteur pour se répandre sur $5$~cm. Qu'est-ce qui assure alors la diffusion générale de \nt{SERO}: la diffusion moléculaire ou la ramification d'un fil à $\sim10^5$ sites de libération?
\end{exercise}

\begin{exercise}[\lvlA]
\label{ex:05:2}
\emph{Régulateur de niveau.} Montrez que dans~\eqref{eq:seroneuron}, pour $f(u)=au$, la sensibilité relative $S=(a/r)\,\partial r/\partial a$ vaut exactement $1/(1+G)$, et pas seulement pour $G\gg1$. Pour $G=9$, l'excitabilité $a$ augmente de $20\,\%$. De combien de pour cent $r$ change-t-il, sans linéarisation?
\end{exercise}

\begin{exercise}[\lvlB]
\label{ex:05:3}
\emph{Récepteur lent dans la boucle.} Le récepteur \nt{SERO} répond avec retard: $r(t)=a\bigl(b-gc(t-T)\bigr)$, le volume obéissant à~\eqref{eq:seroneuron}. Montrez que pour $G>1$ la boucle n'est stable que si $T<T^*=\tau_c\arccos(-1/G)/\sqrt{G^2-1}$. Trouvez $T^*$ pour $\tau_c=1\,\text{s}$, avec $G=2$ et $G=9$. Quelle réduction de la dispersion reste possible avec un retard de quelques centaines de millisecondes?
\end{exercise}

\begin{exercise}[\lvlB]
\label{ex:05:4}
\emph{Bruit de grenaille du niveau.} Chaque impulsion ajoute à $c$, selon~\eqref{eq:lowpass}, une exponentielle de $\tau_c=1\,\text{s}$. Trouvez le coefficient de variation de $c$ si les $3\cdot10^5$ neurones \nt{SERO} libèrent tous dans un même volume des impulsions poissoniennes à $2$ par seconde. Quel $\tau_c$ donnerait $\mathrm{CV}=10\,\%$ pour un seul neurone à $4$ impulsions par seconde?
\end{exercise}

\begin{exercise}[\lvlB]
\label{ex:05:5}
\emph{Simulation: rétroaction contre bruit.} Simulez $N=50$ pacemakers poissoniens de fréquences $r_i=a_i[b-gc]_+$, les $a_i$ uniformes dans $[2.8,\,5.2]$, et le volume~\eqref{eq:seroneuron} avec $\kappa=1/N$, $\tau_c=1\,\text{s}$. De quel facteur la rétroaction ($b=1$, $g=2.25$) réduit-elle le CV du niveau $c$ par rapport à $b=0.1$, $g=0$? Égalise-t-elle les fréquences des neurones?
\end{exercise}

\begin{exercise}[\lvlB]
\label{ex:05:6}
\emph{Conception: XOR en logique \nt{SERO}.} La porte est un pacemaker qui émet $r_0$ si $b-g\sum_kc_k>0$, et $0$ sinon, dans son propre volume $\tau_c\,\dd c/\dd t=-c+\kappa r$; elle calcule un NON-OU. Construisez un XOR avec cinq de ces portes et trouvez son temps d'établissement pour $\tau_c=1\,\text{s}$ et $G'=g\kappa r_0/b=2$.
\end{exercise}

\begin{exercise}[\lvlB]
\label{ex:05:7}
\emph{Patience face à un délai inconnu.} (a)~Un animal accepte d'attendre une récompense trois fois plus grande si l'attente ne dépasse pas $6\,\text{s}$. Quel horizon $\tau_\gamma$ cela implique-t-il? (b)~Supposons le délai $D$ distribué exponentiellement, de moyenne $\bar D$. Montrez qu'attendre est avantageux si $\bar D<\tau_\gamma(R/r_0-1)$, et comparez au délai fixe pour $\tau_\gamma=2\,\text{s}$, $R=3r_0$.
\end{exercise}

\begin{exercise}[\lvlC]
\label{ex:05:8}
\emph{Comment la concentration fixe l'horizon.} Un neurone \nt{DOPA} reçoit $V$ directement avec le poids $k_1$ et, par un intermédiaire \nt{GABA} qui lisse avec $\tau_d=0.1\,\text{s}$, avec le poids $-k_2$; pour $V$ lent, la somme vaut $(k_1-k_2)V+k_2\tau_d\dot V$. Choisissez $k_1$, $k_2$ pour obtenir~\eqref{eq:ctd} avec $\tau_\gamma=2\,\text{s}$. \nt{SERO} multiplie $k_1$ par $m$: de combien faut-il changer $m$ pour doubler l'horizon?
\end{exercise}

\chapter{Des réseaux de neurones qui savent apprendre: ajoutons \nt{DOPA}}
\chaptermark{Des réseaux qui apprennent}
\label{sec:dofa}

\section{Les règles du jeu}

Dans tous les schémas des chapitres précédents, les poids étaient fixés par l'ingénieur. Dans l'organisme, il n'y a pas d'ingénieur. Les poids doivent s'établir d'eux-mêmes, en cours de fonctionnement, et de façon que le réseau fasse quelque chose d'utile. C'est ce qu'on appelle l'apprentissage.

Aux règles précédentes s'en ajoute une. La synapse modifie son poids elle-même, et elle ne peut connaître que ce qui se passe \emph{près d'elle}: l'activité de l'émetteur $x$, l'activité du destinataire $y$ et les substances qui lui parviennent. Personne ne peut envoyer à une synapse une correction personnelle.

Cela exclut la méthode principale des réseaux de neurones artificiels, la rétropropagation de l'erreur. Là, l'erreur de sortie traverse le réseau à rebours par les mêmes poids, dans une phase séparée après la passe avant, et chaque poids reçoit sa propre correction. Il faut pour cela des fils de retour qui reproduisent exactement les fils aller, et une horloge commune. Ni l'un ni l'autre n'a été trouvé dans le cerveau, du moins sous la forme qu'ils ont dans les réseaux artificiels~\cite{Lillicrap2020,WhittingtonBogacz2019}.

Nous écrirons la variation d'un poids comme un processus continu et lent:
\begin{equation}
\frac{\dd w}{\dd t} \;=\; \eta\,\Phi(\text{ce que sait la synapse}),
\label{eq:plasticity}
\end{equation}
où $\eta$ est la vitesse d'apprentissage, si faible que le poids ne change presque pas pendant la durée d'une impulsion. Tout le chapitre porte sur ce que peut être la fonction $\Phi$.

\section{Où est stocké le poids}
\label{sec:weightstore}

Qu'est-ce qu'une synapse qui \og change de poids\fg{}? Une connexion a deux côtés: la terminaison du fil de l'émetteur (côté \emph{présynaptique}) et une portion de la membrane du destinataire portant les récepteurs (côté \emph{postsynaptique}), séparées par une fente étroite. Dans les connexions \nt{GLUT}, la membrane du destinataire se trouve en général sur une \emph{épine}, petite saillie d'une branche d'entrée du destinataire. Il est commode de décomposer le poids de la connexion en trois facteurs~\cite{delCastillo1954}:
\begin{empheq}[box=\keybox]{equation}
w \;=\; N_s\,p\,A_1 .
\label{eq:quantal}
\end{empheq}
Ici, $N_s$ est le nombre de sites de libération entre les deux neurones, $p$ la probabilité qu'une impulsion libère un quantum de neurotransmetteur en un site (le même $p$ qu'au chapitre~\ref{sec:code}), $A_1$ la réponse du destinataire à un quantum, c'est-à-dire le nombre de récepteurs qui le captent. Le facteur $p$ réside chez l'émetteur, $A_1$ chez le destinataire, et $N_s$ exige les deux côtés: de nouveaux sites de libération apparaissent, d'anciens disparaissent, et cela dure toute la vie.

\emph{C'est le destinataire qui décide.} Dans une synapse \nt{GLUT} typique, l'un des récepteurs du glutamate ne conduit le courant que si deux conditions sont réunies: le glutamate est arrivé de l'émetteur, et la membrane du destinataire est déjà excitée~\cite{Nowak1984}. C'est la règle de Hebb inscrite dans une seule molécule: le détecteur de coïncidences du chapitre~\ref{sec:glutcomputer}, placé à l'entrée même de la synapse. En s'ouvrant, le récepteur laisse entrer du calcium dans le destinataire, et le calcium déclenche la modification du poids. Le destinataire est le seul endroit où l'on voit à la fois l'entrée et la sortie; c'est donc chez lui que la décision est prise.

\emph{N'importe quel côté peut exécuter la décision.} La potentialisation à long terme la mieux étudiée a lieu chez le destinataire: de nouveaux récepteurs s'insèrent dans la membrane~\cite{Malinow2002}, et l'épine grossit~\cite{Matsuzaki2004}: c'est $A_1$ qui augmente. Mais le destinataire peut aussi modifier $p$: il libère une substance qui retraverse la fente vers l'émetteur (le canal de retour de l'annexe~\ref{sec:othermediators}), et celui-ci se met à libérer moins de neurotransmetteur~\cite{WilsonNicoll2001}. Quant aux variations de poids à court terme, la dépression et la facilitation du chapitre~\ref{sec:code}, ce sont des variations de $p$ que l'émetteur gère lui-même d'après son activité récente.

Pour l'ingénieur, le registre du poids est réparti entre deux puces. La règle d'apprentissage est exécutée par le destinataire, et il peut en écrire le résultat chez lui comme chez l'émetteur, par le canal de retour. Dans les règles de ce chapitre, le mot \og local\fg{} ne désigne donc pas un seul côté de la connexion, mais la connexion tout entière.

\section{Apprendre sans professeur}

La chose la plus simple que puisse faire une synapse, c'est se renforcer quand l'émetteur et le destinataire sont actifs ensemble:
\begin{equation}
\frac{\dd w}{\dd t} \;=\; \eta\,x\,y .
\label{eq:hebb}
\end{equation}
C'est la règle de Hebb~\cite{Hebb1949}. Elle est locale et n'a besoin d'aucune horloge. Mais elle souffre du même mal que le réseau \nt{GLUT} du chapitre~\ref{sec:glutcomputer}: la rétroaction positive. Un poids fort augmente $y$, un grand $y$ renforce encore le poids, et les poids croissent sans limite. Le remède est une rétroaction négative sur le poids: que le poids décroisse aussi proportionnellement à $y^2$. Pour un neurone d'entrées $\mathbf x$ et de sortie linéaire $y=\mathbf w\cdot\mathbf x$, on obtient
\begin{empheq}[box=\keybox]{equation}
\frac{\dd\mathbf w}{\dd t} \;=\; \eta\,y\,\bigl(\mathbf x - y\,\mathbf w\bigr) .
\label{eq:oja}
\end{empheq}
La règle reste locale: la synapse connaît son entrée $x_i$, la sortie $y$ et son poids $w_i$.

\begin{keyblock}
\begin{proposition}[Ce que trouve la règle de Hebb]
\label{prop:oja}
La règle~\eqref{eq:oja} amène le vecteur des poids à une longueur unité selon la direction où les entrées varient le plus, c'est-à-dire vers le vecteur propre principal de la matrice de corrélation des entrées $C=\langle\mathbf x\mathbf x^{\mathsf T}\rangle$~\cite{Oja1982}. Le neurone apprend à émettre la grandeur unique la plus expressive en laquelle on puisse comprimer ses entrées.
\end{proposition}
\end{keyblock}

Démonstration. Moyennons~\eqref{eq:oja} sur les entrées: $\dd\mathbf w/\dd t=\eta\bigl(C\mathbf w-(\mathbf w^{\mathsf T}C\mathbf w)\,\mathbf w\bigr)$. Les points fixes sont les vecteurs propres unitaires de $C$. Si $\mathbf w$ est proche du vecteur propre $\mathbf e_k$ de valeur propre $\lambda_k$, une petite composante selon un autre vecteur $\mathbf e_j$ croît comme $e^{\eta(\lambda_j-\lambda_k)t}$. Seul est stable le point pour lequel tous les $\lambda_j<\lambda_k$, c'est-à-dire le vecteur principal.

Il existe aussi une variante de la règle de Hebb qui tient compte de l'instant des impulsions. Si l'impulsion de l'émetteur arrive $s$ millisecondes \emph{avant} celle du destinataire, le poids augmente de $A_+e^{-s/\tau_+}$; si elle arrive $s$ millisecondes \emph{après}, il diminue de $A_-e^{-s/\tau_-}$, avec $\tau_\pm$ de l'ordre de vingt millisecondes. Cette règle a été mesurée sur les synapses des neurones \nt{GLUT}~\cite{BiPoo1998}. Elle renforce les connexions qui \emph{ont pu causer} la réponse et affaiblit celles qui arrivent en retard (fig.~\ref{fig:stdp}). Faites passer de nombreuses fois une même séquence d'excitations dans un réseau, et des connexions de chaque maillon vers le suivant y pousseront d'elles-mêmes: la chaîne du chapitre~\ref{sec:glutcomputer}, assemblée sans ingénieur.

\begin{figure}[t]
\centering
\begin{tikzpicture}[>=Stealth,x=0.07cm,y=1.3cm]
\draw[->,gray!70!black] (-66,0) -- (68,0) node[right,black] {\small $s$};
\draw[->,gray!70!black] (0,-1.2) -- (0,1.35) node[above,black] {\small variation du poids};
\draw[very thick,blue!60!black] plot[domain=0.5:60,samples=50] (\x,{exp(-\x/20)});
\draw[very thick,red!70!black] plot[domain=-60:-0.5,samples=50] (\x,{-0.6*exp(\x/20)});
\draw[blue!60!black,thick] (0.5,0) -- (0.5,1);
\draw[red!70!black,thick] (-0.5,0) -- (-0.5,-0.6);
\foreach \x in {-40,-20,20,40}{\draw[gray!60!black] (\x,-0.04) -- (\x,0.04);}
\node[below,font=\scriptsize] at (20,-0.04) {$20\ms$};
\node[above,font=\scriptsize] at (-20,0.04) {$-20\ms$};
\node[align=left,font=\small,blue!60!black,anchor=west] at (22,0.95) {émetteur en avance:\\ la connexion se renforce};
\node[align=right,font=\small,red!70!black,anchor=east] at (-12,-0.95) {émetteur en retard:\\ la connexion s'affaiblit};
\end{tikzpicture}
\caption{La règle de Hebb temporelle. En abscisse, $s=t_{\text{dest}}-t_{\text{ém}}$, le retard de l'impulsion du destinataire sur celle de l'émetteur; $\tau_\pm=20\ms$, $A_-=0.6A_+$. La fenêtre, large de quelques dizaines de millisecondes, est du même ordre que la fenêtre de coïncidence~\eqref{eq:window}.}
\label{fig:stdp}
\end{figure}

Mais toutes les règles de cette section apprennent ce qui est \emph{fréquent}, et non ce qui est \emph{utile}. Une synapse qui ne connaît que $x$ et $y$ ne sait pas si l'action a apporté du jus ou de la douleur. Pour apprendre d'après le résultat, il faut à la synapse un troisième signal.

\section{Le troisième facteur}

Le troisième signal est apporté par \nt{DOPA}: les neurones à dopamine diffusent à tous un seul nombre, l'erreur de prédiction de la récompense $\delta$~\eqref{eq:rpe} (chapitre~\ref{sec:neurontypes}). Supposons la variation du poids proportionnelle au produit de trois facteurs: l'activité de l'émetteur, celle du destinataire et $\delta$.

La difficulté est que la récompense n'arrive pas au moment où le réseau choisissait l'action, mais des centaines de millisecondes ou des secondes plus tard. Quand $\delta$ arrive, le produit $xy$ est nul depuis longtemps, et la synapse a besoin de se souvenir qu'elle a participé. La mémoire la plus simple est notre vieux circuit RC:
\begin{empheq}[box=\keybox]{equation}
\tau_e\,\frac{\dd e}{\dd t} \;=\; -e \;+\; x\,y ,
\qquad
\frac{\dd w}{\dd t} \;=\; \eta\,\delta(t)\,e(t) .
\label{eq:threefactor}
\end{empheq}
La grandeur $e$ est appelée \emph{trace d'éligibilité}~\cite{Fremaux2016}. L'activité conjointe laisse dans la synapse une marque qui s'efface en un temps $\tau_e$. Si de la dopamine arrive dans ce délai, le poids change avec le signe de $\delta$; sinon, la marque disparaît sans laisser de trace (fig.~\ref{fig:trace}). Sur les synapses \nt{GLUT}, on a mesuré que la dopamine arrivée dans la seconde ou les deux secondes qui suivent l'activité conjointe transforme celle-ci en renforcement de la connexion, alors que la dopamine arrivée plus tard n'a pas cet effet~\cite{Yagishita2014}. Des fenêtres de plasticité longues de plusieurs secondes ont aussi été trouvées là où le troisième signal n'est pas la dopamine, mais une forte excitation du destinataire lui-même~\cite{Bittner2017}.

\begin{figure}[t]
\centering
\begin{tikzpicture}[>=Stealth,x=7cm,y=1cm]
\fill[orange!12] (0.7,-0.2) rectangle (0.8,5.2);
% rows: x*y, delta, traces, weights
\foreach \y/\lab in {4.4/{$x\,y$},3.1/{$\delta$},1.4/{trace $e$},0/{poids $w$}}{
  \draw[gray!70!black] (0,\y) -- (1.42,\y);
  \node[left,font=\small] at (-0.02,\y+0.3) {\lab};}
\draw[very thick,blue!60!black] (0,4.4) -- (0,4.9) -- (0.2,4.9) -- (0.2,4.4);
\node[right,font=\scriptsize,blue!60!black] at (0.21,4.8) {émetteur et destinataire actifs ensemble};
\draw[very thick,orange!85!black] (0,3.1) -- (0.7,3.1) -- (0.7,3.7) -- (0.8,3.7) -- (0.8,3.1) -- (1.4,3.1);
\node[right,font=\scriptsize,orange!85!black] at (0.81,3.6) {arrivée de dopamine};
% traces, each in units of its own maximum
\draw[very thick,blue!60!black] plot[domain=0:0.2,samples=20] (\x,{1.4+1.1*(1-exp(-\x))/(1-exp(-0.2))})
  plot[domain=0.2:1.4,samples=60] (\x,{1.4+1.1*exp(-(\x-0.2))});
\draw[thick,densely dashed,gray!50!black] plot[domain=0:0.2,samples=30] (\x,{1.4+1.1*(1-exp(-\x/0.05))/(1-exp(-4))})
  plot[domain=0.2:1.4,samples=80] (\x,{1.4+1.1*exp(-(\x-0.2)/0.05)});
\node[right,font=\scriptsize,blue!60!black] at (1.0,2.15) {$\tau_e=1\,\text{s}$};
\node[right,font=\scriptsize,gray!40!black] at (0.3,1.65) {$\tau_e=50\ms$};
% weights
\draw[very thick,blue!60!black] (0,0.25) -- (0.7,0.25) -- (0.8,0.8) -- (1.4,0.8) node[right,font=\scriptsize] {$\tau_e=1\,\text{s}$: le poids a crû};
\draw[thick,densely dashed,gray!50!black] (0,0.2) -- (1.4,0.2) node[right,font=\scriptsize,gray!40!black] {$\tau_e=50\ms$: inchangé};
% time axis marks
\foreach \x/\l in {0/0,0.2/0.2,0.7/0.7,1.4/1.4}{\draw[gray!60!black] (\x,-0.05) -- (\x,0.05) node[below=3pt,font=\scriptsize] {\l\,s};}
\draw[<->,thick,gray!50!black] (0.2,5.35) -- node[above,font=\scriptsize,black] {délai de la récompense $0.5\,\text{s}$} (0.7,5.35);
\end{tikzpicture}
\caption{La trace selon la règle~\eqref{eq:threefactor}. L'activité conjointe dure $0.2\,\text{s}$; la dopamine arrive une demi-seconde après sa fin. Les traces sont représentées en fraction de leur maximum. La trace lente ($\tau_e=1\,\text{s}$) est encore vivante à ce moment, la trace rapide ($\tau_e=50\ms$) s'est effacée depuis longtemps.}
\label{fig:trace}
\end{figure}

\begin{keyblock}
\begin{proposition}[Apprentissage par un signal diffusé]
\label{prop:threefactor}
La règle~\eqref{eq:threefactor} ne modifie que les synapses qui ont participé au fonctionnement du réseau pendant les derniers $\tau_e$, dans le sens fixé par le signal commun $\delta$. Le signal diffusé, combiné à la trace locale, produit une correction adressée. Un délai entre l'action et le résultat est admissible tant qu'il ne dépasse pas $\tau_e$.
\end{proposition}
\end{keyblock}

\section{Pourquoi cela marche: le bruit comme recherche}
\label{sec:dofagradient}

Pourquoi la règle~\eqref{eq:threefactor} mène-t-elle au mieux? La réponse vient du bruit du chapitre~\ref{sec:code}. Soit la sortie du neurone $y=f(u)+\xi$, où $u=\sum_jw_jx_j$ et $\xi$ est une fluctuation aléatoire de moyenne nulle et de variance $\sigma^2$. La récompense $R$ dépend de $y$. Prenons comme trace non pas simplement $x_jy$, mais sa partie aléatoire $x_j\xi$, et comme troisième facteur l'erreur $\delta=R-\bar R$, où $\bar R$ est la récompense attendue. Pour un bruit faible, $R\approx R\bigl(f(u)\bigr)+R'\,\xi$, donc
\begin{empheq}[box=\keybox]{equation}
\bigl\langle \delta\,\xi\,x_j \bigr\rangle \;=\; \sigma^2\,R'\,x_j \;=\; \frac{\sigma^2}{f'(u)}\,\frac{\partial\langle R\rangle}{\partial w_j} .
\label{eq:perturbation}
\end{empheq}
Le pas moyen suit le gradient de la récompense attendue~\cite{Williams1992}. Personne n'a calculé de dérivées: le réseau a dévié au hasard, la dopamine a signalé si c'était mieux, et les synapses qui ont participé se sont déplacées dans le sens favorable (fig.~\ref{fig:perturbation}). Le bruit qui, au chapitre~\ref{sec:code}, gênait la transmission des messages devient ici une recherche.

\begin{figure}[t]
\centering
\begin{tikzpicture}[>=Stealth,x=2cm,y=3cm]
\draw[->,gray!70!black] (0,0) -- (4.2,0) node[below left,font=\small,black] {sortie $y$};
\draw[->,gray!70!black] (0,0) -- (0,1.2) node[left,font=\small,black] {récompense $R$};
% reward hill
\draw[very thick,blue!60!black] plot[domain=0:4,samples=60] (\x,{max(0,1-((\x-3)/2.2)^2)});
\node[font=\scriptsize,blue!60!black,anchor=south] at (3,1.02) {meilleure sortie};
% noise around f(u)
\draw[thick,gray!60!black] plot[domain=0.6:2.4,samples=50] (\x,{0.28*exp(-((\x-1.5)/0.3)^2/2)});
\draw[densely dashed,gray!60!black] (1.5,0) -- (1.5,0.535);
\node[below,font=\small] at (1.5,0) {$f(u)$};
\node[font=\scriptsize,gray!50!black] at (1.5,-0.2) {fluctuation $\xi$};
% expected reward
\draw[densely dotted,gray!60!black] (0.5,0.535) node[left,font=\scriptsize,black] {$\bar R$} -- (2.1,0.535);
% two random deviations
\fill[green!45!black] (2.1,0.833) circle (0.035);
\draw[very thick,green!45!black] (2.1,0.535) -- (2.1,0.833);
\node[font=\scriptsize,green!45!black,anchor=west,align=left] at (2.18,0.6) {$\xi>0$: mieux,\\ $\delta>0$};
\fill[red!70!black] (0.9,0.089) circle (0.035);
\draw[very thick,red!70!black] (0.9,0.535) -- (0.9,0.089);
\node[font=\scriptsize,red!70!black,anchor=east,align=right] at (0.83,0.27) {$\xi<0$: moins bien,\\ $\delta<0$};
% resulting step
\draw[->,very thick,orange!85!black] (1.55,0.4) -- (1.95,0.4) node[right,font=\scriptsize,black] {pas};
\end{tikzpicture}
\caption{Le bruit comme recherche~\eqref{eq:perturbation}. La sortie du neurone fluctue autour de $f(u)$. Un écart aléatoire vers la droite (en vert) a rapporté plus que prévu, $\delta>0$; vers la gauche (en rouge), moins, $\delta<0$. Dans les deux cas, le produit $\delta\,\xi$ est positif, et le poids se déplace vers la droite, là où la récompense croît. En moyenne, le pas est proportionnel à la pente $R'$: au sommet de la colline, les écarts des deux côtés sont également mauvais, et les pas s'annulent.}
\label{fig:perturbation}
\end{figure}

\begin{keyblock}
\begin{proposition}[Descente de gradient sans fils de retour]
\label{prop:gradient}
Si le réseau présente des écarts d'activité aléatoires, et que le signal diffusé est égal à l'erreur de prédiction de la récompense et nul en moyenne dans chaque situation, la règle~\eqref{eq:threefactor} modifie en moyenne les poids selon le gradient de la récompense attendue. Il n'y faut ni fils de retour ni phase d'apprentissage séparée.
\end{proposition}
\end{keyblock}

On comprend maintenant pourquoi la dopamine ne signale pas la récompense, mais l'\emph{erreur} de sa prédiction. Si le troisième facteur était la récompense elle-même, $R\ge0$, la moyenne dans~\eqref{eq:perturbation} ne changerait pas, mais deux maux apparaîtraient. D'abord, un terme $\bar R\,\xi x_j$ s'ajouterait à la partie utile: nul en moyenne, mais très bruité. Ensuite, et c'est plus grave, une vraie synapse ne sait pas séparer, dans la trace $x_jy$, la partie aléatoire de la partie non aléatoire. Pour des entrées données, $x_jf(u)$ est le même nombre d'un essai à l'autre; si $\delta$ est nul en moyenne dans cette situation, cette partie de la trace ne change rien, mais avec $\delta\ge0$ elle renforce, essai après essai, tout ce que faisait le réseau, le bon comme le mauvais. C'est pourquoi $\bar R$ doit être l'attente \emph{dans la situation donnée}, et non la moyenne sur toute la vie. La calculer est l'affaire du critique, dont il sera question plus bas.

Nous l'avons vérifié sur un modèle. Deux groupes de neurones \nt{GLUT} choisissent l'une de deux actions par inhibition mutuelle, comme sur la fig.~\ref{fig:wta}; les poids du signal \og départ\fg{} vers les groupes sont d'abord égaux, et c'est le bruit qui décide du vainqueur. L'action $A$ apporte du jus avec la probabilité $0.8$, l'action $B$ avec la probabilité $0.2$, et le jus arrive une demi-seconde après le choix. Avec $\tau_e=1\,\text{s}$ et $\delta=R-\bar R$, la part des choix de $A$ sur cinq cents essais est passée de $0.65$--$0.8$ dans la première centaine à $0.96$--$1.0$ dans la dernière, et les poids sont restés entre $0.85$ et $1.35$. Avec $\tau_e=50\ms$, moins que le délai de la récompense, le réseau n'a rien appris: la part des choix de $A$ est restée autour de la moitié. Avec $\delta=R$, sans soustraire l'attente, le réseau s'est aussi mis à choisir $A$, mais le poids du vainqueur a été multiplié par mille en ces mêmes cinq cents essais et continuait de croître; et avec le même $\delta=R$ et une vitesse d'apprentissage dix fois plus grande, le réseau a figé pour toujours, dans l'une de trois simulations, son premier choix aléatoire, le plus mauvais.

\section{Le prix d'un fil unique}

Le signal $\delta$ est unique pour tous, alors que le bruit qu'il évalue est la somme des fluctuations des $N$ neurones qui influent sur le résultat. Chaque synapse voit dans $\delta$ sa partie utile et le bruit étranger de ses $N-1$ voisins. Pour isoler sa part, il faut moyenner sur de nombreux essais, et le temps d'apprentissage croît proportionnellement à $N$~\cite{Werfel2005}. La rétropropagation n'exige pas ce prix.

\begin{keyblock}
\begin{proposition}[Le prix de la diffusion]
\label{prop:broadcastcost}
Le temps d'apprentissage par un seul signal commun croît proportionnellement au nombre de neurones dont le bruit influe sur le résultat. Un tel apprentissage convient à de petits circuits qui choisissent parmi peu d'options, mais pas au réglage de milliards de poids.
\end{proposition}
\end{keyblock}

Le cerveau semble bien répartir le travail ainsi: les sorties des neurones \nt{DOPA} vont avant tout aux régions qui choisissent les actions (chapitre~\ref{sec:neurontypes}), tandis que l'immense réseau \nt{GLUT} du cortex, qui contient l'écrasante majorité des poids, apprend surtout par des règles locales, sans professeur extérieur. On les ramenait autrefois aux règles de Hebb~\eqref{eq:oja}. Aujourd'hui, de plus en plus de données suggèrent que le cortex a son propre objectif, \emph{prédire} sa prochaine entrée; l'erreur est alors calculée sur place, comme la différence entre ce qui était prédit et ce qui est arrivé~\cite{Keller2018}: le professeur est intégré au réseau lui-même. Les fenêtres de plasticité longues de plusieurs secondes, où le troisième signal est une forte excitation du destinataire lui-même~\cite{Bittner2017}, montrent qu'un signal d'erreur local existe bien au niveau des synapses. Qu'il s'agisse d'une règle générale du cortex reste une hypothèse.

\section{D'où vient l'erreur: le critique en temps continu}

Reste à savoir comment les neurones \nt{DOPA} calculent eux-mêmes $\delta$. Dans les programmes d'apprentissage par renforcement, l'erreur de prédiction se calcule par pas $t,t+1,\dots$ Dans l'organisme, il n'y a pas de pas; écrivons-la donc en temps continu~\cite{Doya2000}. Soit $r(t)$ le flux de récompense et $V(t)$ l'attente de toute la récompense future, la plus lointaine valant moins:
\begin{equation}
V(t) \;=\; \int_t^{\infty} e^{-(s-t)/\tau_\gamma}\,r(s)\,\dd s .
\end{equation}
En dérivant, on obtient $\dd V/\dd t=V/\tau_\gamma-r$. Le résidu de cette égalité est précisément l'erreur de prédiction:
\begin{empheq}[box=\keybox]{equation}
\delta(t) \;=\; r(t) \;+\; \frac{\dd V}{\dd t} \;-\; \frac{V}{\tau_\gamma} .
\label{eq:ctd}
\end{empheq}
La constante $\tau_\gamma$ est l'horizon de planification, l'analogue continu du coefficient $\gamma$; selon l'hypothèse, il est fixé par \nt{SERO} (section~\ref{sec:horizon}), et~\eqref{eq:patience} est alors la règle qui dit combien de temps il vaut la peine d'attendre. Vérifions trois cas (fig.~\ref{fig:ctdcases}). Une lampe qui promet du jus s'allume: $V$ monte brusquement, $\dd V/\dd t$ est grand, c'est un pic. Le jus promis arrive: $r>0$, mais l'attente s'épuise au même instant, $\dd V/\dd t\approx-r$, et il n'y a pas de pic. Le jus n'arrive pas: l'attente disparaît sans récompense, $\dd V/\dd t<0$, c'est un creux.

\begin{figure}[t]
\centering
\begin{tikzpicture}[>=Stealth,x=1.15cm,y=1cm]
\foreach \col/\ttl/\lampon/\juice in {0/{jus inattendu}/0/1, 1/{jus promis}/1/1, 2/{pas de jus}/1/0}{
 \begin{scope}[xshift=\col*4.6cm]
  \node[font=\small] at (1.5,3.55) {\ttl};
  \foreach \yy in {2.4,1.2,0}{\draw[gray!60!black] (0,\yy) -- (3.1,\yy);}
  % reward r
  \ifnum\juice=1
    \fill[green!45!black] (1.95,2.4) rectangle (2.05,2.85);
    \pic[scale=0.55] at (2.0,3.1) {juice};
  \else
    \pic[scale=0.55] at (2.0,3.1) {nojuice};
  \fi
  % expectation V and error delta
  \ifnum\lampon=1
    \pic[scale=0.55] at (1.0,3.1) {lamp};
    \draw[very thick,blue!60!black] (0,1.2) -- (1,1.2) -- (1,1.45)
      plot[domain=1:2,samples=20] (\x,{1.2+0.25*exp((\x-1)/1.5)}) -- (2,{1.2+0.25*exp(1/1.5)}) -- (2,1.2) -- (3.1,1.2);
    \draw[very thick,red!70!black] (0,0) -- (0.95,0) -- (0.95,0.5) -- (1.05,0.5) -- (1.05,0) -- (3.1,0);
    \ifnum\juice=0
      \draw[very thick,red!70!black] (1.95,0) -- (1.95,-0.45) -- (2.05,-0.45) -- (2.05,0);
      \node[font=\scriptsize,red!70!black,anchor=west] at (2.1,-0.35) {creux};
    \else
      \node[font=\scriptsize,gray!50!black,anchor=north,align=center] at (2.0,-0.08) {$r$ et la chute de $V$\\ s'annulent};
    \fi
  \else
    \draw[very thick,blue!60!black] (0,1.2) -- (3.1,1.2);
    \draw[very thick,red!70!black] (0,0) -- (1.95,0) -- (1.95,0.5) -- (2.05,0.5) -- (2.05,0) -- (3.1,0);
  \fi
  \ifnum\col=0
    \node[font=\small,anchor=east] at (-0.1,2.55) {$r$};
    \node[font=\small,anchor=east] at (-0.1,1.35) {$V$};
    \node[font=\small,anchor=east] at (-0.1,0.1) {$\delta$};
  \fi
 \end{scope}}
\end{tikzpicture}
\caption{Trois cas pour l'erreur~\eqref{eq:ctd}; de haut en bas, la récompense $r$, l'attente $V$ et l'erreur $\delta$. À gauche, il n'y a pas d'attente, et le jus est une pure surprise. Au milieu, la lampe fait monter $V$ d'un coup: c'est un saut de $\dd V/\dd t$, et le pic de $\delta$ tombe sur la lampe. Ensuite $V$ croît exactement comme l'exige l'horizon, et $\delta=0$; au moment du jus, la récompense $r$ et la chute de $V$ s'annulent. À droite, $V$ chute sans récompense: c'est un creux. Ce sont le pic, le transfert et le creux de la fig.~\ref{fig:rpe}, mais on voit maintenant d'où ils viennent.}
\label{fig:ctdcases}
\end{figure}

Que faut-il pour~\eqref{eq:ctd} parmi les schémas dont nous disposons déjà? Trois choses.

\emph{L'estimation $V$ elle-même.} Elle est fournie par un groupe de neurones \nt{GLUT}, le critique, dont les poids apprennent par la même règle~\eqref{eq:threefactor} avec le même $\delta$: si $\delta>0$, l'attente était sous-estimée, et les connexions actives juste avant se renforcent.

\emph{La dérivée $\dd V/\dd t$.} Elle est calculée par tout schéma qui répond aux variations et non au niveau: une excitation directe combinée à une inhibition retardée par un intermédiaire \nt{GABA}, comme dans le détecteur de direction du chapitre~\ref{sec:gaba}, ou une synapse à dépression~\eqref{eq:depression} du chapitre~\ref{sec:code}.

\emph{Le signe dans un seul fil.} L'erreur peut être positive ou négative, alors que la fréquence des impulsions n'est que positive. La solution est celle du neurone \nt{SERO} du chapitre~\ref{sec:serocomputer}: le neurone \nt{DOPA} est un pacemaker. Ses quelques impulsions régulières par seconde signifient $\delta=0$, une bouffée $\delta>0$, une pause $\delta<0$. L'erreur négative dispose de moins de place: la fréquence ne peut pas descendre sous zéro. Chez les vrais neurones \nt{DOPA}, les creux sont effectivement moins profonds que les pics~\cite{Bayer2005}.

Que la dopamine se comporte comme $\delta$ est mesuré. Comment exactement le cerveau calcule $\dd V/\dd t$ est une hypothèse.

\section{L'acteur et le critique}

Assemblons le tout (fig.~\ref{fig:actorcritic}). Les neurones de contexte excitent les groupes d'actions, qui choisissent un vainqueur par \nt{GABA} (proposition~\ref{prop:wta}) --- c'est l'\emph{acteur} ---, ainsi que le critique, qui fournit $V$~\cite{SuttonBarto2018}. Le neurone \nt{DOPA} reçoit la récompense $r$ et $V$, calcule~\eqref{eq:ctd} et diffuse $\delta$ à toutes les synapses de l'acteur et du critique.

\begin{figure}[t]
\centering
\begin{tikzpicture}[>=Stealth,
  nrn/.style={circle,draw,very thick,fill=blue!6,minimum size=0.9cm,inner sep=0pt,font=\small},
  gab/.style={circle,draw,very thick,fill=red!10,minimum size=0.6cm,inner sep=0pt},
  dofa/.style={circle,draw,very thick,double,double distance=1.2pt,fill=orange!15,minimum size=1.35cm,inner sep=0pt,font=\scriptsize},
  inh/.style={-{Bar[width=7pt]},thick,red!70!black},
  bc/.style={->,thick,densely dashed,orange!85!black}]
\node[nrn] (S1) at (0,2.4) {$x_1$};
\node[nrn] (S2) at (0,1.2) {$x_2$};
\node[nrn] (S3) at (0,0) {$x_3$};
\node[left=0.15cm,align=right,font=\small] at (-0.5,1.2) {contexte};
\node[nrn] (A) at (4,2.6) {$A$};
\node[nrn] (B) at (4,1.0) {$B$};
\node[gab] (G) at (5.2,1.8) {};
\draw[->,thick] (A) -- (G); \draw[->,thick] (B) -- (G);
\draw[inh] (G) to[bend left=35] (A);
\draw[inh] (G) to[bend right=35] (B);
\node[above,font=\small] at (4.6,3.05) {acteur: choix de l'action};
\node[nrn] (V) at (4,-1.1) {$V$};
\node[below,font=\small] at (4,-1.6) {critique};
\foreach \s in {S1,S2,S3}{\foreach \t in {A,B,V}{\draw[->,gray!60!black] (\s) -- (\t);}}
\node[dofa] (D) at (8.0,-1.1) {\nt{DOPA}};
\draw[->,thick] (V) -- node[above,font=\scriptsize] {$V,\ \dd V/\dd t$} (D);
\draw[->,thick,green!45!black] (10.0,-1.1) node[right,black,font=\small] {récompense $r$} -- (D);
\pic[scale=1.1] at (10.6,-0.5) {juice};
\draw[bc] (D.north) -- (8.0,4.0) -- node[above,font=\small,black] {$\delta$ vers toutes les synapses de l'acteur et du critique} (2.2,4.0) -- (2.2,3.0);
\draw[bc] (D.south) -- (8.0,-2.4) -- (2.2,-2.4) -- (2.2,-0.6);
\end{tikzpicture}
\caption{Un réseau qui apprend à choisir ses actions. Flèches grises: synapses \nt{GLUT} à poids modifiables; cercle rouge: neurone \nt{GABA}; double cercle: pacemaker \nt{DOPA}, qui calcule $\delta$ selon~\eqref{eq:ctd}; flèches en tirets: diffusion de $\delta$.}
\label{fig:actorcritic}
\end{figure}

Le signe de l'action du neurotransmetteur est choisi par le destinataire (proposition~\ref{prop:receptor}), et dans la région qui choisit les actions, une partie des neurones porte des récepteurs de la dopamine d'une famille, une autre partie ceux de l'autre famille. En tranche, la dopamine associée à l'activité conjointe renforce les synapses des premiers et affaiblit celles des seconds~\cite{Shen2008}; dans le modèle, cela signifie que chez les seconds le signe de $\delta$ dans~\eqref{eq:threefactor} est inversé. Les premiers apprennent ce qu'il \emph{vaut la peine de faire}, les seconds ce qu'il \emph{ne faut pas faire}.

\section{Bilan}

\begin{table}[t]
\centering
\footnotesize
\setlength{\tabcolsep}{4pt}
\renewcommand{\arraystretch}{1.15}
\begin{tabular}{>{\raggedright}p{2.6cm}>{\raggedright}p{2.7cm}>{\raggedright}p{3.1cm}>{\raggedright\arraybackslash}p{5.0cm}}
\toprule
Règle & Ce que sait la synapse & Ce qu'apprend le réseau & Ce qu'on sait \\
\midrule
Hebb en fréquence~\eqref{eq:oja} & $x$, $y$, son poids & comprimer les entrées: directions principales & le renforcement par activité conjointe est mesuré; le mécanisme qui borne les poids est objet de recherche \\
Hebb temporel & ordre des impulsions de $x$ et $y$ à $\sim20\ms$ près & liens de cause, séquences & mesuré sur les synapses \nt{GLUT} \\
trois facteurs~\eqref{eq:threefactor} & $x$, $y$, trace $e$, signal commun $\delta$ & les actions qui apportent une récompense & l'effet de la dopamine dans les secondes qui suivent l'activité est mesuré; comme mécanisme principal de l'apprentissage du choix, hypothèse bien fondée \\
critique~\eqref{eq:ctd} & idem & l'attente de récompense $V$ & la ressemblance de la dopamine avec $\delta$ est mesurée; le schéma qui calcule $\dd V/\dd t$ est une hypothèse \\
rétropropagation de l'erreur & une correction personnelle pour chaque poids & tout, mais il faut des fils de retour et une horloge & non trouvée dans le cerveau sous cette forme; on en cherche des approximations \\
\bottomrule
\end{tabular}
\caption{Règles d'apprentissage dans un réseau de neurones \nt{GLUT}, \nt{GABA} et \nt{DOPA}.}
\label{tab:learning}
\end{table}

Les règles d'apprentissage sont rassemblées dans le tableau~\ref{tab:learning}. \nt{GLUT} porte les données, \nt{GABA} commande le flux, et \nt{DOPA} décide quelles modifications du réseau conserver.

\section{Exercices}

\begin{exercise}[\lvlA]
\label{ex:06:1}
\emph{Combien de temps vit la trace.} L'activité conjointe $xy=1$ dure $T_a=0.2\,\text{s}$ à partir de $e=0$, et la dopamine arrive $d=0.5\,\text{s}$ après sa fin. (a)~Combien de fois la trace~\eqref{eq:threefactor} est-elle plus grande à cet instant pour $\tau_e=1\,\text{s}$ que pour $\tau_e=50\ms$? (b)~Pour quel $\tau_e$ est-elle maximale?
\end{exercise}

\begin{exercise}[\lvlA]
\label{ex:06:2}
\emph{Dérive de la règle de Hebb.} L'émetteur et le destinataire émettent des trains poissoniens indépendants de $10$ impulsions par seconde, et chaque paire d'impulsions modifie le poids selon la fenêtre de la fig.~\ref{fig:stdp}, avec $\tau_\pm=20\ms$, $A_-=0.6A_+$. De combien de $A_+$ le poids croît-il en une heure d'activité parfaitement aléatoire? Pour quel $\tau_-$ la dérive disparaît-elle?
\end{exercise}

\begin{exercise}[\lvlB]
\label{ex:06:3}
\emph{Corrélations, et non covariances.} La règle de la proposition~\ref{prop:oja} trouve le vecteur propre principal de $C=\langle\mathbf x\mathbf x^{\mathsf T}\rangle$. Les fréquences sont positives: $\mathbf x=\mathbf m+\mathbf n$, $\mathbf m=(\mu,0)$, covariance du bruit $\operatorname{diag}(s_1^2,s_2^2)$. À quelle condition le neurone apprend-il $\mathbf e_1$? Qu'apprend-il pour $\mu=1$, $s_1=0.1$, $s_2=0.5$?
\end{exercise}

\begin{exercise}[\lvlA]
\label{ex:06:4}
\emph{L'horizon dans le pic.} La lampe s'allume à un instant imprévisible, et un jus de taille $R$ arrive $L=1\,\text{s}$ plus tard. Montrez que pour l'attente apprise selon~\eqref{eq:ctd}, $\delta\equiv0$ entre la lampe et le jus, et trouvez l'aire du pic sur la lampe pour $\tau_\gamma=2\,\text{s}$ et $\tau_\gamma=4\,\text{s}$.
\end{exercise}

\begin{exercise}[\lvlB]
\label{ex:06:5}
\emph{D'où vient le prix de la diffusion.} $N$ neurones fluctuent indépendamment, $\xi_i\sim\mathcal N(0,\sigma^2)$, l'erreur est $\delta=a\sum_i\xi_i$, et la synapse $j$ estime le gradient par $\delta\,\xi_jx_j$. Trouvez le rapport signal sur bruit d'un essai. Un neurone apprend en $100$ essais de $5\,\text{s}$: combien de temps prendra l'apprentissage pour $N=10^4$ et $N=10^{10}$?
\end{exercise}

\begin{exercise}[\lvlB]
\label{ex:06:6}
\emph{Conception: un schéma qui calcule $\delta$.} Le neurone \nt{DOPA} reçoit $r$, $V$ avec le poids $k_1$ et une copie de $V$ lissée avec $\tau_d$, avec le poids $-k_2$. (a)~Choisissez $k_1$, $k_2$ pour que, pour $V$ lent, la sortie soit~\eqref{eq:ctd}. (b)~Pour $V=V_0e^{t/\tau_\gamma}$, le vrai $\delta=0$. Pour quel $\tau_d$ la fausse erreur du schéma ne dépasse-t-elle pas $5\,\%$ de $V/\tau_\gamma$, si $\tau_\gamma=2\,\text{s}$?
\end{exercise}

\begin{exercise}[\lvlB]
\label{ex:06:7}
\emph{Simulation: délai limite de la récompense.} Ajoutez un délai de récompense $d$ dans \texttt{run()} d'une copie de \texttt{models/\allowbreak dofa\_learning.py}. Trouvez la part des choix $A$ dans les $100$ derniers des $500$ essais pour $d=0.5\,\text{s}$ et $\tau_e=0.05$, $0.2$, $1$, $4\,\text{s}$, et pour $\tau_e=1\,\text{s}$, $d=3\,\text{s}$. Pour quelle part de trace restant à la récompense (exercice~\ref{ex:06:1}) l'apprentissage disparaît-il?
\end{exercise}

\begin{exercise}[\lvlC]
\label{ex:06:8}
\emph{Peut-on contourner le prix de la diffusion?} $N$ neurones: $y_i=w_i+\xi_i$, $\xi_i\sim\mathcal N(0,\sigma^2)$, récompense $R=-\sum_i(y_i-w_i^*)^2$, règle $\Delta w_i=\eta\,(R-\langle R\rangle)\,\xi_i$. En combien d'essais, pour le meilleur $\eta$, l'erreur $\sum_i(w_i-w_i^*)^2$ est-elle divisée par $e$? Et si seuls $m$ neurones aléatoires sur $N$ fluctuent? La recherche creuse aide-t-elle?
\end{exercise}

\chapter{Théories modernes alternatives de l'apprentissage par la dopamine}
\chaptermark{Autres théories de \nt{DOPA}}
\label{sec:dofaalt}

\section{Ce qui passe vraiment dans le fil}

Au chapitre~\ref{sec:dofa}, le système \nt{DOPA} était pour nous un seul fil portant un seul nombre, l'erreur de prédiction de la récompense $\delta$~\eqref{eq:ctd}. Cette image explique beaucoup, mais plus on mesure la dopamine avec précision, plus on trouve de choses qui n'y entrent pas. Certains neurones répondent à un événement désagréable comme à une récompense; la concentration de dopamine monte pendant que l'animal court vers le but, alors que rien d'inattendu ne se produit; différents neurones \nt{DOPA} se comportent différemment.

Pour l'ingénieur, toutes ces découvertes posent une seule question: \emph{quel est le format du signal sur le bus de diffusion?} Un nombre ou un vecteur? Un seul signal dans le fil, ou plusieurs, séparés en fréquence? Parle-t-il de l'avenir, comme $\delta$, ou du passé? Voici six réponses modernes; pour chacune, nous séparerons ce qui est mesuré de ce qui est supposé.

\section{Non pas un nombre, mais une distribution}

L'erreur~\eqref{eq:ctd} apprend au critique l'attente \emph{moyenne} de la récompense. Mais une demi-goutte de jus certaine et une goutte avec une probabilité d'un demi ont la même moyenne. Pour les distinguer, il faut toute la distribution de la récompense.

Prenons le problème le plus simple: après un indice annonciateur arrive une récompense de taille aléatoire $r$. Supposons qu'il y ait beaucoup de neurones \nt{DOPA} et que chacun, le $i$-ième, entretienne sa propre estimation $V_i$, de sorte qu'au moment de la récompense son erreur vaut $\delta_i=r-V_i$. Et supposons que chacun compte l'erreur positive avec le poids $\alpha_i^+$ et l'erreur négative avec le poids $\alpha_i^-$. Après chaque récompense, l'estimation se déplace de
\begin{equation}
\Delta V_i \;=\; \eta\,\bigl(\alpha_i^+\,[\delta_i]_+ \;-\; \alpha_i^-\,[-\delta_i]_+\bigr).
\label{eq:asymmetric}
\end{equation}
L'estimation cesse de changer quand, en moyenne, les corrections positives équilibrent les négatives:
\begin{empheq}[box=\keybox]{equation}
\alpha_i^+\,\bigl\langle[r-V_i]_+\bigr\rangle \;=\; \alpha_i^-\,\bigl\langle[V_i-r]_+\bigr\rangle ,
\qquad
\kappa_i \;=\; \frac{\alpha_i^+}{\alpha_i^++\alpha_i^-} .
\label{eq:expectile}
\end{empheq}
Pour $\kappa_i=1/2$, c'est la moyenne ordinaire. Pour $\kappa_i>1/2$, le neurone est \og optimiste\fg{}: son estimation est décalée vers le haut de la distribution; pour $\kappa_i<1/2$, il est \og pessimiste\fg{}.

Exemple: le jus arrive avec la probabilité $1/2$, la goutte a la taille $1$. Alors~\eqref{eq:expectile} donne $\kappa_i\cdot\tfrac12(1-V_i)=(1-\kappa_i)\cdot\tfrac12V_i$, c'est-à-dire $V_i=\kappa_i$ (fig.~\ref{fig:expectiles}). Un ensemble de neurones de $\kappa_i$ différents enregistre la distribution de la récompense comme un ensemble de quantiles enregistre une distribution en statistique.

\begin{figure}[t]
\centering
\begin{tikzpicture}[>=Stealth,x=7cm,y=3cm]
\draw[->,gray!70!black] (-0.08,0) -- (1.15,0) node[right,black] {\small récompense};
\draw[->,gray!70!black] (-0.08,0) -- (-0.08,0.75) node[left,black] {\small probabilité};
\fill[blue!25] (-0.03,0) rectangle (0.03,0.5);
\fill[blue!25] (0.97,0) rectangle (1.03,0.5);
\node[below,font=\small] at (0,0) {$0$};
\node[below,font=\small] at (1,0) {$1$};
\node[left,font=\scriptsize] at (-0.08,0.5) {$1/2$};
\foreach \k/\c/\lab/\h in {0.2/blue!60!black/{pessimiste, $\kappa=0.2$}/0.62, 0.5/gray!50!black/{moyenne, $\kappa=0.5$}/0.42, 0.8/red!70!black/{optimiste, $\kappa=0.8$}/0.22}{
  \draw[very thick,\c] (\k,0) -- (\k,\h);
  \node[above,font=\scriptsize,\c] at (\k,\h) {\lab};}
\end{tikzpicture}
\caption{La distribution de la récompense enregistrée par un ensemble de neurones \nt{DOPA}. Récompense $0$ ou $1$ avec la probabilité $1/2$ (barres bleues); un neurone de fraction $\kappa$ aboutit à l'estimation $V=\kappa$~\eqref{eq:expectile}.}
\label{fig:expectiles}
\end{figure}

Ce qui est mesuré: les différents neurones \nt{DOPA} ont des \og points d'inversion\fg{} différents (la taille de récompense pour laquelle la réponse passe du pic au creux), et les neurones dont les réponses aux erreurs positives et négatives sont plus asymétriques s'inversent pour de plus grandes récompenses, comme l'exige~\eqref{eq:expectile}~\cite{Dabney2020}. À partir des réponses d'un groupe de neurones, on parvient à reconstituer la forme de la distribution. Que ce soit la fonction principale de cette dispersion, et non un effet secondaire, est une hypothèse.

\begin{keyblock}
\begin{proposition}[La distribution par l'asymétrie]
\label{prop:expectile}
Si les neurones \nt{DOPA} diffèrent par la fraction $\kappa_i$ avec laquelle ils comptent les erreurs positives, leurs estimations convergent vers différents points~\eqref{eq:expectile} de la distribution de la récompense. Le groupe de neurones transmet non pas la moyenne, mais la distribution, et par là même le risque.
\end{proposition}
\end{keyblock}

\section{Non seulement l'erreur, mais l'importance}

D'après la formule de l'erreur~\eqref{eq:ctd}, un événement désagréable (un choc, un goût amer) devrait provoquer un creux: c'est pire que prévu. Une partie des neurones \nt{DOPA} fait bien cela, mais une autre répond au désagréable par un \emph{pic}, comme à une récompense~\cite{Matsumoto2009}. Ces neurones ne signalent pas le signe de l'erreur, mais le fait que quelque chose d'\emph{important} s'est produit: d'inattendu, de fort, qui demande de l'attention~\cite{BrombergMartin2010}.

Pour l'ingénieur, il est commode d'y voir deux signaux. Le premier est l'erreur signée $\delta$: dans quel sens modifier les poids. Le second est son module $|\delta|$ ou la nouveauté de l'événement: \emph{de combien} les modifier, c'est-à-dire la vitesse d'apprentissage; après un événement inattendu, il faut apprendre plus vite. Par son sens, il est proche du gain noradrénergique du chapitre~\ref{sec:neurontypes}.

Il existe une troisième variante: un fil séparé pour un axe séparé. Les neurones \nt{DOPA} qui projettent vers l'une des régions de choix des actions ne répondent pas à la récompense, mais à la nouveauté et à l'intensité du stimulus, et ils sont nécessaires pour que l'animal apprenne à éviter le danger~\cite{Menegas2018}: ce qui passe dans le fil n'est pas \og mieux ou pire\fg{}, mais \og dangereux ou non\fg{}.

Ce qui est mesuré: différents groupes de neurones \nt{DOPA} répondent différemment au désagréable et au nouveau, et différentes sorties enseignent des choses différentes. La manière dont le cerveau utilise le signal d'importance (comme vitesse d'apprentissage, comme signal d'attention ou comme erreur séparée sur l'axe du danger) est en discussion.

\section{Non seulement enseigner, mais aussi presser}

La dopamine a aussi une action immédiate: plus il y en a, plus l'animal agit volontiers et vite. Comment concilier cela avec l'apprentissage?

La réponse de l'ingénieur est la \emph{séparation en fréquence}. Les pics et les creux rapides, d'une durée de quelques centaines de millisecondes, portent $\delta$ et enseignent. Le niveau lent, qui varie en quelques minutes, porte une autre grandeur, la récompense moyenne par unité de temps $\bar R$~\cite{Niv2007}. Supposons qu'une action exécutée en un temps $\tau_a$ exige un effort $C/\tau_a$: plus c'est rapide, plus c'est cher. En revanche, chaque seconde de retard coûte $\bar R$, la récompense qu'on aurait pu obtenir pendant ce temps. Le coût total $C/\tau_a+\bar R\,\tau_a$ est minimal pour
\begin{empheq}[box=\keybox]{equation}
\tau_a^{*} \;=\; \sqrt{\frac{C}{\bar R}} .
\label{eq:vigor}
\end{empheq}
Plus l'environnement est riche, plus le temps coûte cher et plus il est avantageux d'agir vite: si $\bar R$ double, le temps d'action diminue d'un facteur $\sqrt2\approx1.4$. Remarquez que $\bar R$ est la récompense attendue que l'on soustrayait au chapitre~\ref{sec:dofa}.

La libération de dopamine sur place peut varier sans que change la fréquence des impulsions: des signaux locaux l'ajustent au niveau même des sorties~\cite{Mohebi2019}. Mais les impulsions elles-mêmes contiennent aussi des composantes lentes: dans les expériences de la section suivante, la montée progressive à l'approche de la récompense se voit aussi dans la fréquence des neurones \nt{DOPA}~\cite{Kim2020}. Le niveau lent a donc deux sources, la fréquence des impulsions et la régulation locale de la libération, et laquelle domine dépend apparemment de la sortie et de la tâche.

\begin{keyblock}
\begin{proposition}[Deux signaux sur un même fil]
\label{prop:multiplex}
Le système \nt{DOPA} transmet au moins deux grandeurs séparées dans le temps: l'erreur rapide $\delta$, qui enseigne, et le niveau lent, qui fixe le prix du temps~\eqref{eq:vigor} et par là la vitesse des actions. Il est mesuré que les composantes rapide et lente se comportent différemment; que la composante lente soit précisément $\bar R$ est une hypothèse.
\end{proposition}
\end{keyblock}

\section{La montée à l'approche}

Quand un animal court dans un couloir vers une récompense lointaine, la concentration de dopamine dans la région de choix des actions monte progressivement pendant plusieurs secondes, à mesure que le but se rapproche~\cite{Hamid2016}. D'après le chapitre~\ref{sec:dofa}, cela ne devrait pas arriver: tout se déroule comme prévu, et $\delta$ devrait être nul.

Première explication: cette montée n'est pas $\delta$, mais l'attente $V$ elle-même, le signal \og le but est proche, cela vaut la peine de faire un effort\fg{} de la section précédente~\cite{Hamid2016}. La seconde explication conserve $\delta$~\cite{Kim2020}. Si l'attente est juste, alors, avec l'horizon $\tau_\gamma$, elle croît sur le chemin vers la récompense $R$ comme $V=Re^{-(T-t)/\tau_\gamma}$, et d'après la formule de l'erreur~\eqref{eq:ctd}, $\dd V/\dd t-V/\tau_\gamma=0$. Mais supposons que, de loin, l'attente soit sous-estimée et que l'estimation se précise à l'approche. L'attente croît alors plus vite, disons comme $V=Re^{-(T-t)/\tau_v}$ avec $\tau_v<\tau_\gamma$, et
\begin{empheq}[box=\keybox]{equation}
\delta(t) \;=\; \frac{\dd V}{\dd t}-\frac{V}{\tau_\gamma} \;=\; V\Bigl(\frac{1}{\tau_v}-\frac{1}{\tau_\gamma}\Bigr) \;>\; 0 .
\label{eq:ramp}
\end{empheq}
L'erreur est positive et croît avec $V$: c'est bien la montée progressive (fig.~\ref{fig:ramp}). Pour $\tau_\gamma=4\,\text{s}$ et $\tau_v=1\,\text{s}$, on obtient $\delta=0.75\,V$ par seconde. Cette version a été testée dans un couloir virtuel, en téléportant l'animal plus près du but: la dopamine a répondu par un pic, comme il convient à une dérivée, et non par un niveau progressif~\cite{Kim2020}.

\begin{figure}[t]
\centering
\begin{tikzpicture}[>=Stealth,x=1.6cm,y=2.6cm]
\draw[->,gray!70!black] (-4.2,0) -- (0.5,0) node[below left,black] {\small $t$};
\draw[->,gray!70!black] (-4.2,0) -- (-4.2,1.2);
\foreach \x/\l in {-4/{$T-4$},-2/{$T-2$},0/{$T$}}{\draw[gray!60!black] (\x,-0.03) -- (\x,0.03); \node[below,font=\scriptsize] at (\x,0) {\l\,s};}
\draw[thick,densely dashed,gray!60!black] plot[domain=-4:0,samples=40] (\x,{exp(\x/4)});
\draw[very thick,blue!60!black] plot[domain=-4:0,samples=60] (\x,{exp(\x)});
\draw[very thick,red!70!black] plot[domain=-4:0,samples=60] (\x,{0.75*exp(\x)});
\node[anchor=south east,font=\small,gray!50!black] at (-1.3,0.77) {$V$ pour $\tau_v=\tau_\gamma$: $\delta=0$};
\node[right,font=\small,blue!60!black] at (0.05,1.0) {$V$ pour $\tau_v=1\,\text{s}$};
\node[right,font=\small,red!70!black] at (0.05,0.72) {$\delta$ selon~\eqref{eq:ramp}};
\end{tikzpicture}
\caption{La montée de la dopamine à l'approche de la récompense vue comme une erreur de prédiction, $\tau_\gamma=4\,\text{s}$. En tirets, l'attente qui croît exactement comme l'exige l'horizon ($\delta=0$); en bleu, l'attente qui croît plus vite; en rouge, l'erreur~\eqref{eq:ramp}.}
\label{fig:ramp}
\end{figure}

Ce qui est mesuré: la montée progressive existe, à la fois dans la concentration de dopamine et dans la fréquence des impulsions des neurones \nt{DOPA}~\cite{Kim2020}, et les sauts instantanés de situation provoquent des sauts de dopamine. Laquelle des deux explications est la bonne (ou si les deux le sont, sur des sorties différentes) est en discussion.

\section{Le regard en arrière}

Toutes les théories ci-dessus regardent \emph{en avant}. L'une des théories modernes inverse le sens~\cite{Jeong2022}. Supposons que le cerveau n'apprenne pas à prédire la récompense à partir de l'indice, mais, une fois la récompense reçue, \emph{regarde en arrière}: quel événement l'a précédée plus souvent que les autres? La mesure de ce lien est une différence de deux probabilités:
\begin{empheq}[box=\keybox]{equation}
M \;=\; P(\text{indice avant la récompense}) \;-\; P(\text{indice dans une fenêtre au hasard}) .
\label{eq:retro}
\end{empheq}
Si l'indice survient souvent sans récompense, le second terme est grand, et le lien est faible. Dans cette théorie, la dopamine ne signale pas une erreur de prédiction, mais la mesure dans laquelle un événement donné est une \emph{cause} probable de la récompense.

Le mécanisme du regard en arrière exige la même chose que la règle à trois facteurs du chapitre~\ref{sec:dofa}: chaque événement doit laisser une trace qui s'efface, pour qu'au moment de la récompense on puisse lire ce qui l'a précédée. La différence n'est pas dans la synapse, mais dans ce que transmet \nt{DOPA}.

Dans les expériences où la théorie~\eqref{eq:retro} et l'erreur~\eqref{eq:ctd} prédisent des choses différentes, la libération de dopamine a suivi la première dans plusieurs cas~\cite{Jeong2022}. Mais dans d'autres expériences, la réponse dopaminergique s'est déplacée progressivement, au cours de l'apprentissage, du moment de la récompense vers celui de l'indice, comme l'exige l'apprentissage par l'erreur~\eqref{eq:ctd}~\cite{Amo2022}. On ne sait pas encore quelle théorie décrit le cerveau.

\section{Une erreur qui ne porte pas que sur la récompense}

La dernière extension porte sur l'\emph{objet} de l'erreur. Si l'on remplace le jus attendu par un autre, de même valeur mais de goût différent, la valeur n'a pas changé, et l'erreur~\eqref{eq:ctd} est nulle. Pourtant, les neurones \nt{DOPA} répondent à cette substitution~\cite{Takahashi2017}. Et des pics de dopamine provoqués artificiellement peuvent apprendre à l'animal un lien entre deux stimuli neutres, sans aucune récompense~\cite{Sharpe2017}. Il semble que \nt{DOPA} signale l'erreur de prédiction non seulement de la récompense, mais aussi de \emph{ce qui} va se produire.

Formellement, c'est une généralisation simple de~\eqref{eq:ctd}: au lieu d'un seul flux de récompense $r$, on prend un ensemble de caractéristiques de ce qui se passe, $\varphi_k$ (goût, lieu, son), et pour chacune sa propre attente $M_k$ et sa propre erreur~\cite{Gardner2018}:
\begin{empheq}[box=\keybox]{equation}
\delta_k(t) \;=\; \varphi_k(t) \;+\; \frac{\dd M_k}{\dd t} \;-\; \frac{M_k}{\tau_\gamma} ,
\qquad k=1,\dots,K .
\label{eq:vectortd}
\end{empheq}
La récompense n'est qu'une des caractéristiques, et le vecteur des erreurs n'enseigne pas seulement ce qui est bon, mais aussi ce qui suit quoi: un modèle du monde.

L'hétérogénéité des neurones \nt{DOPA} va dans le même sens: dans une même tâche, différents neurones portent différentes grandeurs, les uns liés à la récompense, d'autres au mouvement, d'autres encore à la direction de l'attention~\cite{Engelhard2019}.

\begin{keyblock}
\begin{proposition}[Un faisceau au lieu d'un fil]
\label{prop:bundle}
Le signal du système \nt{DOPA} n'est pas un nombre unique. Il est réparti entre les neurones (distribution~\eqref{eq:expectile}, caractéristiques différentes~\eqref{eq:vectortd}, axe séparé du danger) et dans le temps (erreur rapide et niveau lent~\eqref{eq:vigor}). L'erreur scalaire~\eqref{eq:ctd} est une première approximation de ce signal, comme le sommateur à seuil est une première approximation du neurone.
\end{proposition}
\end{keyblock}

\section{Bilan}

\begin{table}[t]
\centering
\footnotesize
\setlength{\tabcolsep}{4pt}
\renewcommand{\arraystretch}{1.15}
\begin{tabular}{>{\raggedright}p{2.5cm}>{\raggedright}p{3.2cm}>{\raggedright}p{3.6cm}>{\raggedright\arraybackslash}p{4.2cm}}
\toprule
Théorie & Ce qui passe dans le fil & Mesure principale & Ce qu'on sait \\
\midrule
erreur de prédiction (chap.~\ref{sec:dofa}) & scalaire $\delta$~\eqref{eq:ctd} & pic, transfert sur l'indice, creux & mesuré; première approximation \\
distribution & ensemble de $\delta_i$ d'asymétries différentes & points d'inversion différents selon les neurones & accord avec l'expérience; rôle de la dispersion: hypothèse \\
importance & $|\delta|$, nouveauté; axe du danger & pics au désagréable chez une partie des neurones & mesuré; son usage est en discussion \\
prix du temps & niveau lent $\bar R$ & la dopamine accélère les actions; composante lente dans la libération aux sorties et dans la fréquence des impulsions & séparation rapide/lent mesurée; $\bar R$: hypothèse \\
montée à l'approche & $V$, ou $\delta$ quand l'estimation se précise~\eqref{eq:ramp} & montée progressive, sauts lors de la téléportation dans le couloir & montée mesurée; explication en discussion \\
regard en arrière & mesure du lien causal~\eqref{eq:retro} & expériences où les deux théories divergent & résultats contradictoires \\
vecteur d'erreurs & $\delta_k$ par caractéristique~\eqref{eq:vectortd} & réponse au changement de goût; apprentissage sans récompense & mesuré; forme vectorielle: hypothèse \\
\bottomrule
\end{tabular}
\caption{Ce que transmet le système \nt{DOPA}: l'image standard du chapitre~\ref{sec:dofa} et ses extensions modernes.}
\label{tab:dofaalt}
\end{table}

Les théories du tableau~\ref{tab:dofaalt} ne s'excluent pas: presque toutes gardent l'erreur de prédiction comme base et en élargissent le format. Seul le regard en arrière lui fait sérieusement concurrence, et ce débat sera tranché par l'expérience. Pour l'ingénieur, la conclusion est simple: le prix de la diffusion de la proposition~\ref{prop:broadcastcost} baisse si le fil est en réalité un ensemble de fils aux destinataires différents.

\section{Exercices}

\begin{exercise}[\lvlA]
\label{ex:07:1}
\emph{Les points de la distribution pour une seule goutte.} Un jus de taille $1$ arrive avec la probabilité $p=0.2$, sinon la récompense vaut $0$. Trouvez par~\eqref{eq:expectile} $V_i$ pour $\kappa_i=0.2$, $0.5$, $0.8$. Que vaut la médiane de cette récompense, et en quoi le point~\eqref{eq:expectile} diffère-t-il d'un quantile?
\end{exercise}

\begin{exercise}[\lvlB]
\label{ex:07:2}
\emph{La distribution d'après deux neurones.} La récompense vaut $0$ ou $x$, avec la probabilité $p$ pour $x$. Deux neurones à règle~\eqref{eq:asymmetric} ont signalé $V(1/2)=0.25$ et $V(0.8)=0.5$. Retrouvez $p$, $x$ et l'écart type de la récompense. Que signalera un neurone avec $\kappa=0.2$?
\end{exercise}

\begin{exercise}[\lvlB]
\label{ex:07:3}
\emph{Simulation: la distribution par l'asymétrie.} Simulez neuf neurones \nt{DOPA} avec $\kappa_i=0.1,\dots,0.9$ et la règle~\eqref{eq:asymmetric}, $\alpha_i^+=2\kappa_i$, $\alpha_i^-=2(1-\kappa_i)$, pour une récompense uniforme sur $[0,1]$. Comment la dispersion des $V_i$ dépend-elle de $\eta$, et quel $\eta$ faut-il pour distinguer cette distribution de deux gouttes équiprobables $0.25$ et $0.75$?
\end{exercise}

\begin{exercise}[\lvlA]
\label{ex:07:4}
\emph{Le prix du temps.} (a)~Établissez~\eqref{eq:vigor} et trouvez le coût total à l'optimum. (b)~Le cerveau a estimé $\bar R$ avec une erreur d'un facteur $k$. Trouvez le surcoût relatif pour $k=2$ et $k=4$. Avec quelle précision le niveau lent de dopamine doit-il signaler $\bar R$?
\end{exercise}

\begin{exercise}[\lvlB]
\label{ex:07:5}
\emph{Pic lors de la téléportation.} L'attente croît selon~\eqref{eq:ramp} avec $\tau_v=1\,\text{s}$, $\tau_\gamma=4\,\text{s}$; l'animal est téléporté instantanément du point $T-2\,\text{s}$ au point $T-1\,\text{s}$. Trouvez l'aire du pic de $\delta$ en fraction de $R$, et combien de secondes de rampe avant le saut il remplace. Que montrera la téléportation si la dopamine signale $V$ lui-même?
\end{exercise}

\begin{exercise}[\lvlB]
\label{ex:07:6}
\emph{Conception: deux signaux sur un fil.} Le fil \nt{DOPA} porte un niveau qui croît de $s=1/60$~imp/s chaque seconde, et des événements de $A=3$ impulsions. Une synapse à trace $\tau_e=1\,\text{s}$ soustrait de la fréquence sa copie lissée avec $\tau_f$. Pour quels $\tau_f$ perd-on au plus $10\,\%$ de l'événement, en gardant l'apport du niveau pendant la trace sous $10\,\%$ de celui de l'événement?
\end{exercise}

\begin{exercise}[\lvlB]
\label{ex:07:7}
\emph{Conception d'expérience: regard en arrière contre erreur.} Dans une fenêtre, l'indice survient avec la probabilité $0.1$, et après lui la récompense avec la probabilité $0.8$. Comparez $\Delta P=P(\text{réc.}\mid\text{ind.})-P(\text{réc.}\mid\text{sans})$ et $M$ de~\eqref{eq:retro} si (a)~l'on ajoute des récompenses sans indice, $0.08$ par fenêtre; (b)~l'on double la fréquence de l'indice sans récompenses.
\end{exercise}

\begin{exercise}[\lvlC]
\label{ex:07:8}
\emph{Le faisceau et le prix de la diffusion.} Dans le modèle de l'exercice~\ref{ex:06:8}, les $N$ neurones sont répartis en $K$ groupes, chacun avec son fil, et dans chaque fil s'infiltre une fraction $\rho$ des erreurs des autres. Montrez que la constante d'apprentissage vaut $N/K+2+\rho^2N(1-1/K)$. Combien d'essais donne-t-elle pour $K\to\infty$, $N=10^{10}$ et $\rho=0.01$?
\end{exercise}

\chapter{Quand chercher, quand décider: ajoutons \nt{NORA}}
\chaptermark{Ajoutons \nt{NORA}}
\label{sec:nora}
\begin{chaptergoals}
\item comment un seul bouton, le gain, fait passer un neurone d'amplificateur analogique à comparateur;
\item pourquoi augmenter le gain n'aide que contre le bruit qui naît après l'amplificateur;
\item comment le gain fait passer un réseau de choix de la réflexion à la décision sans changer un poids;
\item pourquoi le gain joue l'inverse d'une température, et pourquoi la récompense suit un U inversé.
\end{chaptergoals}

\section{Ce qui manque}

Au chapitre~\ref{sec:dofa}, le réseau apprenait grâce au bruit: les écarts aléatoires de l'activité constituaient la recherche, et \nt{DOPA} signalait s'ils amélioraient les choses. Mais un paramètre y est resté non fixé: \emph{combien} chercher. Si le bruit est trop fort, le réseau s'agite et n'utilise pas ce qu'il sait déjà. S'il est trop faible, le réseau choisit sans cesse la même chose et ne remarque pas que, tout près, les choses se sont améliorées. En apprentissage par renforcement, on appelle cela le compromis entre exploitation et exploration, et chaque algorithme possède un bouton pour le régler. Au chapitre~\ref{sec:neurontypes}, nous avons supposé que, dans le cerveau, c'est \nt{NORA} qui tourne ce bouton.

L'homme possède quelques dizaines de milliers de neurones à noradrénaline (tableau~\ref{tab:types}), presque tous réunis en un seul petit groupe, et leurs sorties se répartissent dans pratiquement tout le cerveau, bien que différentes parties de ce groupe desservent en partie, on l'a découvert, des régions différentes~\cite{Poe2020}. C'est un canal de diffusion encore plus étroit que \nt{DOPA}. Il fonctionne selon deux modes~\cite{AstonJones2005}. En mode \emph{tonique}, de fond, les neurones émettent des impulsions régulières, à une fréquence allant d'une fraction d'impulsion à quelques impulsions par seconde, et cette fréquence varie lentement avec l'éveil. En mode \emph{phasique}, ils répondent par une courte bouffée à un événement important pour la tâche en cours, à peu près au moment de la décision. Un point de mesure commode, mais imprécis, est le diamètre de la pupille: il varie avec l'activité de ces neurones, mais aussi avec celle d'autres régions~\cite{Joshi2016} et des sorties cholinergiques dans le cortex~\cite{Reimer2016}. La pupille indique le niveau général d'éveil, et non \nt{NORA} seul.

\section{Le gain}

Voici ce qui est mesuré: la noradrénaline supprime l'activité de fond des neurones plus fortement que leur réponse au stimulus, et le rapport du signal au fond augmente~\cite{Foote1975NE}. Que la pente de la caractéristique d'un neurone individuel soit alors littéralement multipliée par un coefficient ne découle pas de l'expérience. Comme au chapitre~\ref{sec:neurontypes}, prenons un modèle simple qui rend compte de cet effet: la noradrénaline modifie la pente de la caractéristique de transfert du destinataire~\cite{ServanSchreiber1990}. Les entrées faibles, sous le seuil (le fond), donnent alors encore moins, et les fortes encore plus. Supposons que le neurone réponde par la fréquence
\begin{empheq}[box=\keybox]{equation}
r \;=\; F\bigl(g\,(u-\theta)\bigr), \qquad F(z)=\frac{r_{\max}}{1+e^{-z}} ,
\label{eq:gain}
\end{empheq}
où $u$ est le courant d'entrée, $\theta$ le seuil, et $g$ le gain, piloté dans le modèle par \nt{NORA}. Pour $g$ petit, la fréquence suit l'entrée continûment sur une large plage: le neurone est un \emph{amplificateur analogique}. Pour $g$ grand, presque toute entrée au-dessus du seuil donne la fréquence maximale, et au-dessous, zéro: le neurone est un \emph{comparateur}, un élément presque numérique (fig.~\ref{fig:gaincurves}). Un seul bouton fait basculer le même neurone entre deux modes qui, en électronique, sont réalisés par des circuits différents.

\begin{figure}[t]
\centering
\begin{tikzpicture}[>=Stealth,x=1.1cm,y=2.4cm]
\draw[->,gray!70!black] (-3.3,0) -- (3.4,0) node[below left,black] {\small $u$};
\draw[->,gray!70!black] (-3.3,0) -- (-3.3,1.2) node[left,black] {\small $r$};
\draw[densely dashed,gray!60!black] (0,0) -- (0,1.1);
\node[below,font=\small] at (0,0) {$\theta$};
\draw[densely dotted,gray!60!black] (-3.3,1) -- (3.2,1) node[right,black,font=\small] {$r_{\max}$};
\draw[very thick,blue!60!black] plot[domain=-3.2:3.2,samples=60] (\x,{1/(1+exp(-0.8*\x))});
\draw[very thick,orange!85!black] plot[domain=-3.2:3.2,samples=60] (\x,{1/(1+exp(-2*\x))});
\draw[very thick,red!70!black] plot[domain=-3.2:3.2,samples=120] (\x,{1/(1+exp(min(-8*\x,9)))});
\node[blue!60!black,font=\small,anchor=west] at (0.5,0.12) {petit $g$: amplificateur};
\node[red!70!black,font=\small,anchor=east] at (-0.3,0.85) {grand $g$: comparateur};
\end{tikzpicture}
\caption{La caractéristique de transfert~\eqref{eq:gain} pour trois valeurs du gain $g$.}
\label{fig:gaincurves}
\end{figure}

Un grand gain aide-t-il contre le bruit? Pas toujours. Supposons que deux entrées diffèrent de $\Delta u$ et qu'il faille dire laquelle est la plus grande. Si le bruit $\sigma_{\text{avant}}$ s'ajoute à l'entrée \emph{avant} l'amplificateur, le signal et le bruit sont amplifiés ensemble, et leur rapport ne change pas. Si en revanche le bruit $\sigma_{\text{après}}$ s'ajoute \emph{après} l'amplificateur (celui des synapses peu fiables et des impulsions aléatoires du réseau lui-même, comme au chapitre~\ref{sec:code}), le rapport signal sur bruit
\begin{empheq}[box=\keybox]{equation}
d' \;=\; \frac{g\,\Delta u}{\sqrt{g^2\sigma_{\text{avant}}^2+\sigma_{\text{après}}^2}}
\label{eq:gainsnr}
\end{empheq}
croît avec $g$ jusqu'à ce que $g\sigma_{\text{avant}}$ égale $\sigma_{\text{après}}$~\cite{ServanSchreiber1990}.

\begin{keyblock}
\begin{proposition}[Quand le gain aide]
\label{prop:gainnoise}
En augmentant le gain, \nt{NORA} n'améliore la discrimination des entrées que contre le bruit qui naît \emph{après} l'amplificateur, dans le réseau lui-même. Le bruit arrivé avec l'entrée, le gain ne l'élimine pas.
\end{proposition}
\end{keyblock}

\section{Le gain dans le choix entre deux options}

Revenons au choix du chapitre~\ref{sec:gaba}: deux groupes de neurones \nt{GLUT}, d'entrées $I_1$, $I_2$, s'inhibent mutuellement avec la force $\beta$. Chaque groupe a désormais un gain $g$: dans les équations du choix~\eqref{eq:wta}, l'expression entre crochets est multipliée par $g$. Tant que les deux groupes sont actifs, les fréquences stationnaires se déduisent de $r_1=g(I_1-\beta r_2)$, $r_2=g(I_2-\beta r_1)$. En additionnant et en soustrayant ces égalités, on obtient la somme $r_1+r_2=g(I_1+I_2)/(1+g\beta)$ et la différence $r_1-r_2=g(I_1-I_2)/(1-g\beta)$. Leur rapport dit avec quelle netteté le réseau préfère une entrée à l'autre:
\begin{empheq}[box=\keybox]{equation}
\frac{r_1-r_2}{r_1+r_2} \;=\; \varepsilon\,\frac{1+g\beta}{1-g\beta}, \qquad \varepsilon=\frac{I_1-I_2}{I_1+I_2}, \quad g\beta<1 .
\label{eq:wtagain}
\end{empheq}
Une petite différence d'entrées $\varepsilon$ est amplifiée d'un facteur $(1+g\beta)/(1-g\beta)$, et ce facteur croît sans limite quand $g\beta$ approche de un. Pour $g\beta>1$, l'état où les deux groupes sont actifs est instable, comme dans la proposition~\ref{prop:wta}: l'un gagne, l'autre se tait (fig.~\ref{fig:pitchfork}).

\begin{figure}[t]
\centering
\begin{tikzpicture}[>=Stealth,x=3.2cm,y=1.5cm]
\draw[->,gray!70!black] (0,0) -- (2.15,0) node[below left,black] {\small $g\beta$};
\draw[->,gray!70!black] (0,-1.25) -- (0,1.35) node[above,black] {\small $(r_1-r_2)/(r_1+r_2)$};
\draw[gray!60!black] (1,-0.04) -- (1,0.04); \node[below right,font=\small] at (1,0) {$1$};
\node[left,font=\scriptsize] at (0,1) {$1$}; \node[left,font=\scriptsize] at (0,-1) {$-1$};
\draw[very thick,blue!60!black] plot[domain=0:0.905,samples=60] (\x,{0.05*(1+\x)/(1-\x)}) -- (1,1) -- (2.05,1);
\draw[very thick,blue!60!black] (1.105,-1) -- (2.05,-1);
\draw[thick,densely dashed,gray!60!black] (1,0) -- (2.05,0);
\node[font=\small,align=center] at (0.45,-0.62) {réfléchit:\\les deux groupes\\sont actifs};
\node[font=\small,align=center] at (1.55,0.45) {a décidé:\\un seul est actif};
\node[font=\small,blue!60!black,anchor=south west] at (1.05,-0.97) {l'entrée faible a gagné avant: elle tient};
\end{tikzpicture}
\caption{Choix entre deux options pour différents gains; les entrées diffèrent de $\varepsilon=0.05$. Pour $g\beta>1$, les deux décisions sont stables (traits pleins; la victoire de l'entrée faible pour $g\beta>(1+\varepsilon)/(1-\varepsilon)\approx1.1$), et le réseau maintient celle qu'il a déjà prise; l'état symétrique (en tirets) est instable.}
\label{fig:pitchfork}
\end{figure}

\begin{keyblock}
\begin{proposition}[Le gain change le régime du choix]
\label{prop:gainwta}
Dans un réseau de choix à inhibition mutuelle $\beta$, le gain $g$ fait franchir au réseau la frontière $g\beta=1$. En dessous, le réseau \og réfléchit\fg{}: les deux groupes sont actifs, et la différence des entrées est amplifiée selon~\eqref{eq:wtagain}. Au-dessus, le réseau \og a décidé\fg{}: un seul groupe est actif, et la décision tient. En modifiant $g$, \nt{NORA} peut faire passer le réseau d'un régime à l'autre sans toucher à un seul poids.
\end{proposition}
\end{keyblock}

\section{Le gain comme température}

Ajoutons maintenant au choix le bruit qui naît dans le réseau lui-même, après l'amplificateur. Le groupe gagnant est déterminé par le signe de la grandeur $g(u_A-u_B)+\eta$, où $u_A-u_B$ est la différence des entrées, c'est-à-dire la différence des poids appris au chapitre~\ref{sec:dofa}, et $\eta$ un bruit d'échelle $\sigma$. Si le bruit suit une loi logistique (pour un bruit gaussien, la courbe est presque la même), la probabilité de choisir $A$ vaut
\begin{empheq}[box=\keybox]{equation}
P(A) \;=\; \frac{1}{1+e^{-g\,(u_A-u_B)/\sigma}} .
\label{eq:softmax}
\end{empheq}
Le programmeur reconnaîtra ici le choix par softmax à la température $T=\sigma/g$. Le gain est l'inverse de la température. Pour $g$ petit, même une option nettement meilleure n'est guère plus souvent choisie que la moins bonne: le réseau explore beaucoup. Pour $g$ grand, la meilleure option connue est choisie presque toujours: le réseau exploite ce qu'il sait.

Précisons ce qu'est $g$ dans~\eqref{eq:wtagain} et~\eqref{eq:softmax}: c'est le gain \emph{effectif} appliqué à la différence des entrées de la tâche en cours, au moment du choix. Les deux modes de \nt{NORA} agissent sur lui en sens opposés. La bouffée phasique au moment de la décision l'augmente, et le réseau consolide son choix. Une forte activité tonique va au contraire de pair avec l'exploration: chez l'homme, la pupille de base est plus large avant un choix au hasard~\cite{JepmaNieuwenhuis2011}, et chez le rat, l'entrée \nt{NORA} dans le cortex cingulaire antérieur fait passer le choix en mode aléatoire~\cite{Tervo2014stochastic}. Une forte \nt{NORA} tonique correspond donc à un $g$ effectif \emph{petit}, et la bouffée à un $g$ grand.

\section{Combien chercher}

Quel gain vaut le mieux? Nous l'avons vérifié sur un modèle. Le réseau choisit l'une de deux actions selon~\eqref{eq:softmax}; chacune rapporte une récompense avec une certaine probabilité, que le réseau estime d'après les récompenses de l'action choisie, comme au chapitre~\ref{sec:dofa}. De temps en temps, le monde change: les deux probabilités sont tirées à nouveau. Ce peut être l'action choisie qui change, ou celle que le réseau n'a pas essayée depuis longtemps; de la seconde, le réseau ne peut apprendre le changement qu'en explorant. La fig.~\ref{fig:invertedU} montre quelle part de la meilleure récompense possible le réseau obtient pour différents $g$ constants.

\begin{figure}[t]
\centering
\begin{tikzpicture}[>=Stealth,x=1.2cm,y=1cm]
\draw[->,gray!70!black] (0,0) -- (7.6,0) node[below left,black] {\small $g$};
\draw[->,gray!70!black] (0,0) -- (0,4.4) node[above,black,font=\small] {part de la meilleure récompense};
\foreach \x/\l in {0/0.5,1/1,2/2,3/4,4/8,5/16,6/32,7/64}{\draw[gray!60!black] (\x,-0.06) -- (\x,0.06); \node[below,font=\scriptsize] at (\x,0) {\l};}
\foreach \v/\y in {0.8/0.8,0.9/2.4,1.0/4.0}{\draw[gray!60!black] (-0.06,\y) -- (0.06,\y); \node[left,font=\scriptsize] at (0,\y) {\v};}
\draw[very thick,blue!60!black] plot[mark=*,mark size=1.3pt] coordinates {(0,0.368) (1,0.880) (2,1.728) (3,2.784) (4,3.456) (5,3.616) (6,3.488) (7,3.264)};
\draw[very thick,red!70!black] plot[mark=*,mark size=1.3pt] coordinates {(0,0.368) (1,0.672) (2,1.184) (3,1.840) (4,2.176) (5,1.920) (6,1.456) (7,1.104)};
\node[blue!60!black,font=\small,anchor=south] at (5,3.7) {monde calme};
\node[red!70!black,font=\small,anchor=north] at (5.1,1.25) {monde qui change vite};
\draw[->,thick,blue!60!black] (5,3.25) -- (5,3.55);
\draw[->,thick,red!70!black] (4,1.8) -- (4,2.1);
\node[font=\small\itshape,gray!35!black,anchor=west] at (0.15,2.3) {à l'aveugle};
\node[font=\small\itshape,gray!35!black] at (6.45,2.55) {rate les changements};
\end{tikzpicture}
\caption{Part de la meilleure récompense possible pour un gain $g$ constant (échelle de $g$ logarithmique). Courbe bleue: le monde change en moyenne une fois tous les mille essais; courbe rouge: une fois tous les vingt. Les flèches marquent le meilleur $g$: dans un monde qui change vite, il est deux fois plus petit. Moyenne sur 60 simulations de 4000 essais.}
\label{fig:invertedU}
\end{figure}

Pour $g=0.5$, le réseau n'utilise presque pas ce qu'il sait et obtient environ les trois quarts du possible; pour $g$ très grand, il rate les changements. La meilleure valeur: $g=16$ quand le monde change une fois tous les mille essais (part $0.976$), et $g=8$ quand il change une fois tous les vingt ($0.886$).

\begin{keyblock}
\begin{proposition}[Le U inversé]
\label{prop:explore}
La récompense que recueille le réseau dépend du gain comme un U inversé: un $g$ trop petit gaspille des essais en exploration, un $g$ trop grand ne remarque pas les changements. Le meilleur $g$ est d'autant plus petit que le monde change vite.
\end{proposition}
\end{keyblock}

Le U inversé s'observe aussi dans les expériences: les animaux réussissent le mieux une tâche quand l'activité tonique des neurones \nt{NORA} est modérée et que leurs réponses phasiques sont nettes, alors qu'avec une forte activité tonique ils se laissent distraire et passent d'une tâche à l'autre~\cite{AstonJones2005}. Cela s'accorde avec la distinction des modes de la section précédente: la bouffée phasique au moment de la décision donne un grand $g$ effectif précisément quand il faut choisir, alors qu'une forte activité tonique sans bouffées amplifie tout indistinctement, y compris les entrées étrangères à la tâche, si bien que le $g$ effectif pour la tâche en cours diminue: c'est l'exploration.

\section{Qui tourne le bouton}

Puisque le meilleur gain dépend de la vitesse à laquelle le monde change, il serait bon de l'ajuster. Mais comment les neurones \nt{NORA} sauraient-ils que le monde a changé? L'indice naturel est la \emph{surprise}: les erreurs de prédiction sont devenues plus grandes que d'habitude. Il importe de distinguer deux sortes d'incertitude. Si la récompense est simplement aléatoire, les erreurs sont toujours grandes, et c'est attendu. Si en revanche les erreurs ont soudain augmenté par rapport à leur taille habituelle, c'est que quelque chose a changé. Selon une hypothèse, c'est précisément cette incertitude inattendue que signale \nt{NORA}, l'incertitude attendue étant signalée par \nt{ACET}~\cite{YuDayan2005}.

Que faire de ce signal? On peut baisser le gain et explorer davantage. On peut augmenter la vitesse d'apprentissage pour oublier plus vite les estimations périmées: les expériences montrent qu'après un changement brusque, les humains apprennent plus vite, et que leur pupille se dilate alors~\cite{Nassar2012}; rappelons que la pupille n'indique pas seulement \nt{NORA}, mais aussi \nt{ACET}~\cite{Reimer2016}. On peut enfin faire les deux à la fois: selon une autre hypothèse, la bouffée phasique de \nt{NORA} fonctionne comme une \emph{interruption}, un signal par lequel le réseau remet à zéro son état courant et se reconfigure pour la nouvelle situation~\cite{BouretSara2005}, comme la ligne d'interruption commune d'un processeur.

Disons honnêtement ce que nous n'avons pas trouvé. Dans le modèle, nous avons essayé les deux manières simples d'ajuster: baisser $g$, ou augmenter la vitesse d'apprentissage, quand l'erreur lissée dépasse sa moyenne à long terme. Dans un monde qui tantôt reste longtemps stable, tantôt change souvent, les meilleurs réglages des deux méthodes ont donné $0.926$--$0.927$ de la meilleure récompense, presque autant que le meilleur $g$ constant ($0.925$ pour $g=8$) ou qu'une vitesse d'apprentissage constante mais assez grande ($0.928$), et les réglages malheureux, moins. Les courbes de la fig.~\ref{fig:invertedU} sont plates près du sommet, et dans un tel monde il n'y a presque rien à ajuster. L'ajustement devrait être rentable là où des régimes différents exigent des réglages très différents. Quelle loi de commande \nt{NORA} met réellement en œuvre n'est pas encore établi.

\section{Bilan}

\begin{table}[t]
\centering
\footnotesize
\setlength{\tabcolsep}{4pt}
\renewcommand{\arraystretch}{1.15}
\begin{tabular}{>{\raggedright}p{3.0cm}>{\raggedright}p{4.4cm}>{\raggedright\arraybackslash}p{6.0cm}}
\toprule
Ce que fait \nt{NORA} & Comment & Ce qu'on sait \\
\midrule
change la pente de la réponse des neurones & le gain $g$ dans~\eqref{eq:gain} & l'augmentation du rapport signal sur bruit est mesurée; le modèle multiplicatif est une simplification \\
fait basculer le choix & fait franchir au réseau $g\beta=1$ & découle du modèle~\eqref{eq:wtagain}; pas de mesure directe du seuil \\
fixe combien chercher & $g$ est l'inverse de la température~\eqref{eq:softmax}; tonique: petit $g$ effectif (exploration), bouffée: grand (décision) & le U inversé et les deux modes de fonctionnement sont mesurés; le lien avec l'exploration est une hypothèse bien fondée \\
signale les changements & incertitude inattendue & la pupille et la vitesse d'apprentissage augmentent après un changement: mesuré; qu'il s'agisse d'un signal \nt{NORA} est une hypothèse \\
interrompt et remet à zéro & bouffée phasique sur un événement inattendu & les bouffées sur les événements importants sont mesurées; l'\og interruption\fg{} est une hypothèse \\
\bottomrule
\end{tabular}
\caption{Ce que \nt{NORA} ajoute à un réseau de \nt{GLUT}, \nt{GABA} et \nt{DOPA}.}
\label{tab:nora}
\end{table}

\nt{DOPA} dit au réseau \emph{dans quel sens} modifier les poids. \nt{NORA} ne modifie aucun poids, mais décide \emph{comment} le réseau utilise ce qu'il sait déjà (tableau~\ref{tab:nora}): un seul bouton, le gain, transforme le neurone d'amplificateur en comparateur, le réseau de choix de \og réfléchissant\fg{} en \og décidé\fg{}, et le choix d'exploration en exploitation. Comment \nt{NORA} apprend à quelle vitesse change le monde est une question ouverte.

\section{Exercices}

\begin{exercise}[\lvlA]
\label{ex:08:1}
\emph{Un seul bouton pour tout le cerveau.} (a)~D'après le tableau~\ref{tab:types}, estimez combien de neurones du cerveau reviennent à un neurone \nt{NORA}. (b)~Le bruit avant l'amplificateur est $\sigma_{\text{avant}}=0.5$, après, $\sigma_{\text{après}}=4$. Pour quel $g$ la grandeur $d'$ de~\eqref{eq:gainsnr} atteint-elle $90\,\%$ de sa limite?
\end{exercise}

\begin{exercise}[\lvlB]
\label{ex:08:2}
\emph{Choix avec gain.} $\tau\,\dd r_1/\dd t=-r_1+g[I_1-\beta r_2]_+$, $\tau\,\dd r_2/\dd t=-r_2+g[I_2-\beta r_1]_+$, $\tau=20\ms$. (a)~En combien de temps $r_1-r_2$ s'amortit-il pour $g\beta=0.5$ et $0.9$? (b)~Pour $\varepsilon=0.05$, trouvez le $g\beta$ en dessous duquel les deux groupes sont actifs et celui au-dessus duquel la victoire de l'entrée faible est stable.
\end{exercise}

\begin{exercise}[\lvlA]
\label{ex:08:3}
\emph{Le softmax à partir du bruit.} Chacun des $K$ groupes reçoit son propre bruit de Gumbel, $P(\eta_k<z)=\exp(-e^{-z/\sigma})$, et c'est le plus grand $gu_k+\eta_k$ qui gagne. Trouvez $P(k)$. Pour quel $g$ la meilleure de deux options, avec $u_A-u_B=0.1$, $\sigma=1$, est-elle choisie dans $90\,\%$ des cas?
\end{exercise}

\begin{exercise}[\lvlB]
\label{ex:08:4}
\emph{Estimation du meilleur gain.} Dans le modèle de la fig.~\ref{fig:invertedU}, les probabilités de récompense après un changement sont uniformes sur $[0,1]$. Le réseau essaie la moins bonne action une fois tous les $e^{g\Delta}$ essais; en égalant cela à $1/h$, on obtient $g^*\approx\ln(1/h)/\bar\Delta$. Trouvez $\bar\Delta=\langle|p_A-p_B|\rangle$ et $g^*$ pour $h=0.001$, $0.01$, $0.05$.
\end{exercise}

\begin{exercise}[\lvlB]
\label{ex:08:5}
\emph{Simulation: gain et vitesse des changements.} Trouvez $g^*(h)$ pour $h$ de $0.001$ à $0.2$ avec une copie de \texttt{models/\allowbreak nora\_explore.py} (elle écrit des fichiers dans le dossier courant). Où l'estimation de l'exercice~\ref{ex:08:4} est-elle juste, et pourquoi $g^*$ cesse-t-il de décroître quand les changements sont rapides?
\end{exercise}

\begin{exercise}[\lvlB]
\label{ex:08:6}
\emph{Conception: une interruption pour le réseau de choix.} Le réseau de l'exercice~\ref{ex:08:2}, avec $\beta=2$, $\tau=20\ms$, $g=1$, maintient $r_1=0.9$, $r_2=0$, bien que les entrées soient désormais $I_1=0.9$, $I_2=1.0$. \nt{NORA} abaisse $g$ à $g_{\text{lo}}$ pendant un temps $T_p$. Trouvez le plus petit $T_p$ analytiquement pour $g_{\text{lo}}=0$, et par simulation pour $g_{\text{lo}}=0.25$.
\end{exercise}

\begin{exercise}[\lvlC]
\label{ex:08:7}
\emph{Quand l'ajustement est rentable.} Un oracle sait si l'époque est calme ou changeante et règle $g$ en conséquence. (a)~Simulez son gain sur le meilleur $g$ constant dans le monde du chapitre (époques de $1000$ essais, $h=0.001$ et $0.05$). (b)~Construisez un monde où ce gain est plus grand, et trouvez quelle part en récupère l'ajustement par la surprise.
\end{exercise}

\chapter{Quand écrire, quand lire: ajoutons \nt{ACET}}
\chaptermark{Ajoutons \nt{ACET}}
\label{sec:acet}

\section{Ce qui manque}

Une puce de mémoire possède, outre les fils d'adresse et de données, une ligne de plus: la validation d'écriture. Tant qu'elle est inactive, la mémoire est seulement lue; quand elle est active, ce qui se trouve sur les lignes de données est écrit. Sans cette ligne, la mémoire ne fonctionne pas: si chaque lecture écrit en même temps quelque chose, ce qui a été lu, y compris ce qui a été lu avec une erreur, est aussitôt réécrit.

Dans le cerveau, ce problème est plus aigu que dans une puce. Au chapitre~\ref{sec:glutcomputer}, la mémoire était un groupe de neurones \nt{GLUT} maintenu actif par ses propres connexions (fig.~\ref{fig:bistable}). Au chapitre~\ref{sec:dofa}, nous avons vu que ces connexions apprennent par la règle de Hebb en plein fonctionnement. Les mêmes synapses conservent donc l'ancien et reçoivent le nouveau, et il n'y a pas de phase d'écriture séparée, pas plus que d'horloge. Comment le réseau distingue-t-il le moment où il faut mémoriser ce que l'on voit du moment où il faut se rappeler ce que l'on sait? Au chapitre~\ref{sec:neurontypes}, nous avons supposé que cette ligne est tenue par \nt{ACET}. Dans ce chapitre, nous verrons comment une telle ligne doit fonctionner et pourquoi on ne peut pas s'en passer.

\section{Les trois actions d'un seul fil}

Les neurones à acétylcholine qui travaillent à l'intérieur du cerveau sont peu nombreux et réunis en quelques petits groupes, et leurs sorties se répartissent dans tout le cortex. C'est un canal de diffusion, comme \nt{DOPA} et \nt{NORA}. Le niveau d'acétylcholine dans le cortex est élevé pendant l'éveil attentif et bas pendant le sommeil lent~\cite{Hasselmo1999}. Comme pour la dopamine, le sens du signal est fixé par le récepteur du destinataire (proposition~\ref{prop:receptor}): l'acétylcholine a des récepteurs rapides, qui ouvrent un canal, et des récepteurs lents, qui déclenchent des réactions dans la cellule. Trois actions nous importent.

\emph{Premièrement: affaiblir les connexions internes.} Des expériences sur tranches de cortex olfactif ont montré que l'acétylcholine supprime fortement la transmission par les connexions entre neurones d'une même région et touche à peine les entrées qui y arrivent de l'extérieur~\cite{HasselmoBower1992}.

\emph{Deuxièmement: renforcer les entrées.} L'acétylcholine augmente l'excitabilité des neurones et facilite la transmission sur certaines voies d'entrée; le signal venu des organes des sens devient plus fort que l'activité propre du réseau~\cite{Hasselmo2006}.

\emph{Troisièmement: faciliter la modification des poids.} Quand le niveau d'acétylcholine est élevé, les connexions changent plus facilement~\cite{Hasselmo2006}.

Pour l'ingénieur, ces trois actions forment une seule opération: le passage du mode \og lire\fg{} au mode \og écrire\fg{}. Le réseau cesse de s'écouter lui-même, se met à écouter l'entrée et mémorise ce qu'il entend.

\section{Un réseau qui complète}

Prenons $N$ neurones \nt{GLUT} reliés entre eux. Stockons dans leurs connexions plusieurs motifs (des ensembles de $pN$ neurones actifs simultanément) par la règle de Hebb~\cite{Hebb1949}. Si l'on présente à ce réseau la moitié d'un motif, les connexions excitent la moitié manquante: le réseau \emph{complète} le motif à partir d'une partie, comme le groupe de la fig.~\ref{fig:bistable} s'entretenait lui-même. C'est une mémoire associative: l'adresse est une partie du contenu. \nt{GABA}, comme au chapitre~\ref{sec:gaba}, maintient le nombre total de neurones actifs autour de $pN$, pour que deux motifs ne puissent pas s'allumer ensemble.

Notons $a$ le niveau d'acétylcholine, de $0$ à $1$, et inscrivons directement ses trois actions dans le modèle:
\begin{empheq}[box=\keybox]{equation}
\tau\frac{\dd r_i}{\dd t} = -r_i + F\Bigl(\tfrac{1+a}{2}\,x_i + (1-a)\sum_j W_{ij}\,r_j - g\Bigr),
\qquad
\frac{\dd W_{ij}}{\dd t} = \eta\,a\,(r_i-p)(r_j-p) .
\label{eq:acetnet}
\end{empheq}
Ici, $r_i$ est la fréquence du neurone $i$, $x_i$ son entrée venue des organes des sens, $F$ une caractéristique de transfert à seuil, $g$ l'inhibition globale due à \nt{GABA}, qui croît quand plus de $pN$ neurones sont actifs. Le facteur $\frac{1+a}{2}$ est l'amplification de l'entrée, $1-a$ l'affaiblissement des connexions internes, $\eta a$ la vitesse d'apprentissage. La seconde équation est la règle de Hebb~\eqref{eq:hebb} sous une forme adaptée aux motifs clairsemés~\cite{Tsodyks1988}: le poids croît quand les deux neurones sont plus actifs que d'habitude, et diminue quand un seul l'est. La partie négative des poids est, comme toujours, réalisée par des intermédiaires \nt{GABA}. Il n'y a pas d'horloge: les neurones comme les poids changent continûment.

\begin{figure}[t]
\centering
\begin{tikzpicture}[>=Stealth,
  nrn/.style={circle,draw,very thick,fill=blue!6,minimum size=0.8cm,inner sep=0pt,font=\small},
  acet/.style={circle,draw,very thick,double,double distance=1.2pt,fill=orange!15,minimum size=1.35cm,inner sep=0pt,font=\scriptsize},
  bc/.style={->,thick,densely dashed,orange!85!black}]
\foreach \i/\y in {1/2.4,2/1.2,3/0}{
  \node[nrn] (N\i) at (3,\y) {$r_\i$};
  \draw[->,thick,blue!60!black] (0,\y) node[left,black] {\small $x_\i$} -- (N\i);}
\node[left,align=right,font=\small] at (-0.5,1.2) {des organes\\des sens};
\draw[->,thick,gray!60!black] (N1) to[bend left=30] (N2);
\draw[->,thick,gray!60!black] (N2) to[bend left=30] (N1);
\draw[->,thick,gray!60!black] (N2) to[bend left=30] (N3);
\draw[->,thick,gray!60!black] (N3) to[bend left=30] (N2);
\draw[->,thick,gray!60!black] (N1.east) .. controls (4.6,2.0) and (4.6,0.4) .. (N3.east);
\node[right,font=\small,gray!40!black,align=left] at (4.7,1.2) {connexions\\internes $W$};
\node[acet] (A) at (1.2,-1.9) {\nt{ACET}};
\draw[bc] (A) -- (1.6,-0.15);
\node[font=\small,orange!60!black,anchor=east] at (0.95,-0.9) {$+$ entrée};
\draw[bc] (A) -- (4.3,0.6);
\node[font=\small,orange!60!black,anchor=west] at (4.3,0.1) {$-$ connexions internes, $+$ vitesse d'apprentissage};
\node[right,font=\small,align=left] at (2.2,-1.9) {niveau haut: \og écrire\fg{}\\niveau bas: \og lire\fg{}};
\end{tikzpicture}
\caption{\nt{ACET} comme ligne de validation d'écriture. Les neurones $r_i$ reçoivent l'entrée $x_i$ des organes des sens (flèches bleues) et s'excitent mutuellement par les connexions internes $W$ (en gris), qui stockent les motifs. L'acétylcholine (flèches en tirets) renforce l'entrée, affaiblit les connexions internes et facilite la modification des poids, comme dans~\eqref{eq:acetnet}.}
\label{fig:acetnet}
\end{figure}

\section{L'écriture et la lecture se gênent}

Testons le modèle~\eqref{eq:acetnet} sur une tâche où l'écriture et la lecture sont réellement en conflit. Un réseau de $200$ neurones stocke déjà cinq motifs de $20$ neurones actifs. Il faut maintenant en mémoriser un sixième, qui recouvre à moitié l'un des anciens: ils ont dix neurones en commun. Cela arrive constamment: une nouvelle pièce ressemble à une pièce connue, un nouveau visage à un ancien.

\emph{Écriture avec $a$ bas.} Séparons d'abord les deux premières actions de l'acétylcholine de la troisième: gardons la vitesse d'apprentissage maximale, et seuls l'amplification de l'entrée et l'affaiblissement des connexions internes varieront. Pour $a=0$, l'entrée allume les dix neurones communs, et ceux-ci, par les connexions internes, allument la moitié restante de l'\emph{ancien} motif. \nt{GABA} empêche les deux de brûler en même temps, et l'ancien motif, soutenu par ses propres connexions, évince le nouveau, qui ne tient que par l'entrée. Le réseau ne voit pas ce qu'il a devant lui, mais ce dont il se souvient, et la règle de Hebb enregistre consciencieusement ce qu'il a vu.

Au bout du compte, le nouveau motif n'a pas du tout été mémorisé: à partir de sa moitié distinctive, le réseau ne se rappelle rien. Et l'ancien est abîmé: rappelé par sa moitié distinctive, il entraîne en moyenne six neurones étrangers sur dix. Pour $a=0.1$, le nouveau motif est déjà mémorisé, mais il fusionne avec l'ancien, et les dégâts sont même plus grands: chacun des deux, rappelé par sa moitié, entraîne environ huit neurones étrangers; à partir de $a=0.2$, l'écriture est propre: moins d'un neurone superflu au rappel (moyenne sur dix simulations).

\emph{Écriture avec une vraie vitesse d'apprentissage.} Supposons maintenant que l'acétylcholine fixe aussi la vitesse d'apprentissage, comme dans~\eqref{eq:acetnet}. En une seule présentation, le nouveau motif n'est entièrement mémorisé que pour $a\ge0.9$; pour $a=0.75$, aux huit dixièmes; pour $a\le0.5$, pas du tout.

\emph{Lecture.} Présentons à un réseau qui stocke proprement cinq motifs la moitié d'un motif, puis retirons-la aussi. Pour $a\le0.2$, le réseau complète tout le motif et le maintient seul; pour $a=0.3$, presque tout; pour $a\ge0.4$, les connexions internes sont trop affaiblies, et une fois l'indice disparu, l'activité s'éteint (fig.~\ref{fig:acetsim}).

\begin{figure}[t]
\centering
\begin{tikzpicture}[>=Stealth,x=9cm,y=3.2cm]
\draw[->,gray!70!black] (0,0) -- (1.1,0) node[right,black] {\small $a$};
\draw[->,gray!70!black] (0,0) -- (0,1.2);
\foreach \x in {0,0.2,0.4,0.6,0.8,1.0}{\draw[gray!60!black] (\x,-0.02) -- (\x,0.02); \node[below,font=\scriptsize] at (\x,0) {$\x$};}
\foreach \y in {0,0.5,1}{\draw[gray!60!black] (-0.01,\y) -- (0.01,\y); \node[left,font=\scriptsize] at (0,\y) {$\y$};}
\node[rotate=90,font=\small] at (-0.1,0.6) {part du motif restituée};
\draw[very thick,blue!60!black] plot[mark=*,mark size=1.6pt] coordinates {(0,1.00) (0.1,1.00) (0.2,1.00) (0.3,0.96) (0.4,0.04) (0.5,0.00) (0.75,0.00) (1.0,0.00)};
\draw[very thick,red!70!black] plot[mark=*,mark size=1.6pt] coordinates {(0.1,0.00) (0.2,0.00) (0.3,0.00) (0.5,0.00) (0.75,0.80) (0.9,1.00) (1.0,1.00)};
\node[font=\small,blue!60!black,anchor=west,align=left] at (0.03,0.5) {lecture:\\motif d'après\\sa moitié};
\node[font=\small,red!70!black,anchor=east,align=right] at (1.0,0.35) {écriture:\\nouveau motif\\en une fois};
\end{tikzpicture}
\caption{Le modèle~\eqref{eq:acetnet}: $200$ neurones, motifs de $20$ neurones actifs, moyenne sur dix simulations. Points bleus: la lecture, c'est-à-dire la part du motif restituée d'après sa moitié une fois l'indice retiré. Points rouges: l'écriture, c'est-à-dire la part d'un nouveau motif, à moitié semblable à un ancien, dont le réseau se souvient après une seule présentation, si l'acétylcholine fixe aussi la vitesse d'apprentissage. La lecture fonctionne pour $a\le0.3$, l'écriture pour $a\ge0.9$.}
\label{fig:acetsim}
\end{figure}

\begin{keyblock}
\begin{proposition}[L'écriture et la lecture exigent des modes différents]
\label{prop:writeread}
Dans le modèle~\eqref{eq:acetnet}, où les mêmes connexions conservent l'ancien et apprennent le nouveau, et où l'acétylcholine fixe aussi la vitesse d'apprentissage, il n'existe pas de niveau $a$ auquel l'écriture et la lecture marchent aussi bien l'une que l'autre. Pour écrire, il faut affaiblir les connexions internes, sinon le réseau écrit non pas l'entrée, mais ce dont il s'est souvenu; pour lire, il faut les renforcer, sinon il ne complète pas le motif à partir d'une partie. Il faut un commutateur, et un seul fil de diffusion peut en tenir lieu.
\end{proposition}
\end{keyblock}

Si l'acétylcholine ne modifiait pas la vitesse d'apprentissage, une bande étroite $0.2\le a\le0.3$ conviendrait aux deux tâches. Mais le réseau apprendrait alors à chaque lecture, c'est-à-dire qu'il réécrirait ce qu'il a lu, avec ses erreurs. L'idée même de supprimer les connexions internes pendant l'apprentissage vient de modèles construits à partir de mesures sur tranches~\cite{HasselmoAndersonBower1992}. Les nombres ci-dessus concernent notre modèle simplifié; où se trouvent exactement les frontières des modes dans le cerveau, on ne le sait pas.

\section{Jusqu'où croire ses yeux}

Le commutateur \og écrire/lire\fg{} est un cas extrême d'une question plus générale: jusqu'à quel point croire l'entrée, et jusqu'à quel point la mémoire? Les ingénieurs ont une réponse exacte à cette question, et elle explique pourquoi un seul fil commande à la fois le mode et la vitesse d'apprentissage.

Supposons que le réseau suive une grandeur $s(t)$ qui dérive lentement: en une seconde, sa variance croît de $q$. Les organes des sens rapportent $y(t)=s(t)+$bruit d'intensité $R$. La meilleure estimation $\hat s$ en temps continu est donnée par le filtre de Kalman--Bucy~\cite{KalmanBucy1961}:
\begin{empheq}[box=\keybox]{equation}
\frac{\dd\hat s}{\dd t} \;=\; \frac{P}{R}\,\bigl(y-\hat s\bigr),
\qquad
\frac{\dd P}{\dd t} \;=\; q-\frac{P^2}{R} ,
\label{eq:kalman}
\end{empheq}
où $P$ est la variance de l'erreur d'estimation, c'est-à-dire l'incertitude du réseau sur ce qu'il sait. La première équation est le même circuit RC que celui de la membrane dans le modèle du neurone~\eqref{eq:glutneuron}, mais avec une constante de temps $R/P$ variable. Quand le réseau n'est pas sûr, $P$ est grand, la constante de temps est petite, et l'estimation suit vite l'entrée: c'est \og écrire\fg{}. Quand il est sûr, $P$ est petit, et l'estimation n'écoute presque plus l'entrée, elle s'en tient à la mémoire: c'est \og lire\fg{}.

En régime établi, $P=\sqrt{qR}$ et la constante de temps de la mémoire vaut $\sqrt{R/q}$: si le monde dérive d'une unité en cent secondes, $q=0.01$, et si le bruit des organes des sens est $R=1$, la mémoire doit retenir environ dix secondes.

L'essentiel est ici qu'un seul nombre, $P/R$, fixe deux choses à la fois: le poids de l'entrée par rapport à la mémoire, et la vitesse à laquelle change l'estimation stockée. C'est exactement ce que fait l'acétylcholine dans~\eqref{eq:acetnet}. Selon une hypothèse~\cite{YuDayan2005}, l'acétylcholine communique justement au réseau une grandeur semblable à $P$: l'incertitude \emph{attendue}, celle que le réseau connaît d'avance. Ce canal n'est pas toujours lent: une partie des neurones \nt{ACET} répond à la récompense et à la punition en une vingtaine de millisecondes, d'autant plus fort que le renforcement est inattendu~\cite{Hangya2015}.

\begin{keyblock}
\begin{proposition}[Un seul bouton pour le poids de l'entrée et la vitesse d'apprentissage]
\label{prop:kalman}
Dans le filtre optimal~\eqref{eq:kalman}, le poids avec lequel l'entrée se mêle à la mémoire et la vitesse d'apprentissage sont une seule et même grandeur, $P/R$. Il est donc raisonnable de les commander par un seul signal. Si les propriétés du monde ne changent pas, ce signal atteint un niveau constant $\sqrt{q/R}$, et il ne faut l'ajuster que lorsque les propriétés du monde changent.
\end{proposition}
\end{keyblock}

La seconde phrase de la proposition explique le résultat du chapitre~\ref{sec:nora}, où l'ajustement de la vitesse d'apprentissage par la surprise n'a pas battu la meilleure vitesse constante: dans un monde dont les propriétés de changement sont fixes, la meilleure vitesse d'apprentissage est justement constante. Le gain apparaît quand le monde devient soudain différent. L'estimation se retrouve alors tout à coup loin de l'entrée, et $P$ l'ignore. La bonne action est d'augmenter brusquement $P$, et donc de passer en mode écriture. Selon l'hypothèse discutée au chapitre~\ref{sec:nora}, c'est \nt{NORA} qui signale un tel changement \emph{inattendu}~\cite{YuDayan2005}. La division du travail est naturelle: \nt{NORA} remarque que le monde a changé, \nt{ACET} maintient le réseau en mode écriture tant que la nouvelle situation n'est pas apprise.

\section{Le sommeil comme mode lecture}

Si un niveau élevé d'acétylcholine est le mode écriture, que fait le réseau quand ce niveau est bas? Pendant le sommeil lent, le niveau d'acétylcholine chute, les connexions internes sont libérées de la suppression, et l'entrée venue des organes des sens est presque nulle. Le réseau, en mode lecture sans indice, rejoue ce qui y est inscrit. Les enregistrements montrent que les neurones qui ont travaillé ensemble pendant la journée se déclenchent de nouveau ensemble pendant le sommeil lent~\cite{WilsonMcNaughton1994}, et même dans le même ordre~\cite{Skaggs1996}. Selon une hypothèse~\cite{Hasselmo1999}, ces répétitions recopient dans la mémoire à long terme ce qui a été appris vite pendant la journée (allonger artificiellement les courtes bouffées de répétitions améliore la mémoire~\cite{FernandezRuiz2019}): le jour, écriture depuis l'entrée; la nuit, réécriture depuis l'intérieur. Pour l'ingénieur, cela ressemble à l'écriture dans un tampon rapide pendant la journée et à son vidage dans une mémoire permanente lente pendant la nuit.

\section{Bilan}

\begin{table}[t]
\centering
\footnotesize
\setlength{\tabcolsep}{4pt}
\renewcommand{\arraystretch}{1.15}
\begin{tabular}{>{\raggedright}p{3.2cm}>{\raggedright}p{4.2cm}>{\raggedright\arraybackslash}p{6.0cm}}
\toprule
Ce que fait \nt{ACET} & Comment & Ce qu'on sait \\
\midrule
affaiblit les connexions internes & facteur $1-a$ dans~\eqref{eq:acetnet} & mesuré sur tranches \\
renforce l'entrée & facteur $\frac{1+a}{2}$ & l'augmentation de l'excitabilité et de la transmission sur les voies d'entrée est mesurée \\
facilite l'apprentissage & vitesse $\eta a$ & mesuré \\
commute \og écrire/lire\fg{} & proposition~\ref{prop:writeread} & découle des modèles; pas de mesure directe des modes \\
signale l'incertitude attendue & $P$ dans~\eqref{eq:kalman} & hypothèse \\
libère le réseau pour les répétitions du sommeil & $a$ bas en sommeil lent & le niveau bas pendant le sommeil et les répétitions sont mesurés; la recopie en mémoire à long terme est une hypothèse \\
\bottomrule
\end{tabular}
\caption{Ce que \nt{ACET} ajoute à un réseau de \nt{GLUT}, \nt{GABA}, \nt{DOPA} et \nt{NORA}.}
\label{tab:acet}
\end{table}

\nt{DOPA} dit au réseau dans quel sens modifier les poids, \nt{NORA} avec quelle détermination utiliser ce qu'il sait, et \nt{ACET} s'il doit écouter le monde ou lui-même (tableau~\ref{tab:acet}). C'est la dernière ligne de commande diffusée du tableau~\ref{tab:types}: la validation d'écriture, sans laquelle une mémoire qui stocke ses motifs dans les connexions mêmes par lesquelles elle les lit s'abîmerait à chaque lecture.

\section{Exercices}

\begin{exercise}[\lvlA]
\label{ex:09:1}
\emph{Combien de temps retenir.} Le filtre~\eqref{eq:kalman} avec $q=0.01$, $R=1$ retient environ $10$~s. Trouvez $P_\infty$ et la mémoire $\sqrt{R/q}$ si (a)~l'amplitude du bruit du capteur double; (b)~le monde dérive cent fois plus vite. (c)~De quel facteur doit croître le bruit pour que le réseau retienne une minute?
\end{exercise}

\begin{exercise}[\lvlB]
\label{ex:09:2}
\emph{Mode écriture après un changement.} Le monde a changé, et l'incertitude du filtre~\eqref{eq:kalman} avec $q=0.01$, $R=1$ a bondi à $P_0\to\infty$. (a)~Avec quelle constante de temps $P$ revient-il à $P_\infty$? (b)~Au bout de combien de secondes la vitesse d'apprentissage ne dépasse-t-elle plus la valeur établie de plus de $10\,\%$?
\end{exercise}

\begin{exercise}[\lvlB]
\label{ex:09:3}
\emph{Vitesse d'apprentissage constante.} Le réseau apprend selon $\dd\hat s/\dd t=(y-\hat s)/T$; le monde dérive, $\dd s/\dd t=w$, $y=s+n$, où $w$ et $n$ sont des bruits blancs d'intensités $q$ et $R$. Trouvez le meilleur $T$ et la variance de l'erreur correspondante. De quel facteur augmente-t-elle si $T$ est faux d'un facteur $2$, puis $10$?
\end{exercise}

\begin{exercise}[\lvlB]
\label{ex:09:4}
\emph{Où est la frontière de l'écriture.} Dans \texttt{models/\allowbreak acet\_memory.py}, le seuil vaut $\theta=0.35$, le poids à l'intérieur d'un motif $w=0.041$ (idéalement $0.045$). Un nouveau motif partage $m=10$ neurones avec un ancien. Pour quel $a$ dans~\eqref{eq:acetnet} les autres neurones de l'ancien motif ne s'allument-ils pas à partir de ces $m$ (seuil franc)?
\end{exercise}

\begin{exercise}[\lvlB]
\label{ex:09:5}
\emph{Simulation: capacité de la mémoire.} Dans \texttt{models/\allowbreak acet\_memory.py}, modifiez la fonction \texttt{read\_sweep}: pour $a=1$, écrivez $M$ motifs ($M$ de $5$ à $100$), puis pour $a=0$ lisez les dix premiers d'après leur moitié. Pour quel $M$ apparaissent les erreurs, et comment la limite $M_{\max}$ dépend-elle de $N$ et de $p$?
\end{exercise}

\begin{exercise}[\lvlB]
\label{ex:09:6}
\emph{Conception: écrire sans fuite.} Avec la vitesse d'apprentissage $\eta a$, le réseau apprend aussi pendant la lecture. Proposez une fonction monotone $\eta(a)\le\eta$, nulle pour $a\le0.3$ et pas moins bonne que $\eta a$ pour $a\ge0.9$. Vérifiez: après $20$ lectures à $a=0.2$ avec un indice contenant trois neurones étrangers, combien de ceux-ci sont entrés dans le motif?
\end{exercise}

\begin{exercise}[\lvlC]
\label{ex:09:7}
\emph{Recopier le tampon la nuit.} (a)~La nuit, $a=0.05$ et une répétition dure $0.1\,\text{s}$, contre $0.2\,\text{s}$ le jour pour $a=1$. Combien de répétitions valent en poids une présentation diurne, pour la vitesse $\eta a$ et pour $\eta(a)$ de l'exercice~\ref{ex:09:6}? (b)~La mémoire permanente, $100$ fois plus lente que le tampon, entend le motif $20$ fois par nuit, $0.1\,\text{s}$ chaque fois. Combien de nuits?
\end{exercise}

\chapter{Toute la machine}
\chaptermark{Toute la machine}
\label{sec:synthesis}

\section{Ce que nous avons assemblé}

Nous avons pris un par un les types de neurones du tableau~\ref{tab:types} et, pour chacun, nous nous sommes demandé ce qu'il apporte à une machine de calcul. Les règles étaient toujours les mêmes: seulement ce qui existe dans un organisme vivant, et aucune horloge commune. Assemblons maintenant les pièces en une seule machine et regardons comment elles travaillent ensemble.

L'image du chapitre~\ref{sec:neurontypes} s'est confirmée et précisée. Un immense réseau adressé de neurones \nt{GLUT} porte les données. Les neurones \nt{GABA} commandent le flux de ces données. Au-dessus du réseau se tiennent quatre canaux de diffusion, chacun avec son propre bouton: \nt{DOPA} dit quels changements de poids conserver, \nt{SERO} maintient le niveau général et, par hypothèse, l'horizon de planification, \nt{NORA} choisit entre chercher et décider, \nt{ACET} entre écrire et lire (fig.~\ref{fig:machine}).

\begin{figure}[t]
\centering
\begin{tikzpicture}[>=Stealth,
  blk/.style={draw,very thick,rounded corners=3pt,align=center,font=\small},
  reg/.style={draw,very thick,double,double distance=1.2pt,circle,minimum size=1.25cm,inner sep=0pt,font=\scriptsize,fill=orange!15},
  bc/.style={->,thick,densely dashed,orange!85!black}]
% sensors, bus, actions
\node[blk,fill=green!8,text width=1.7cm] (S) at (0.9,0) {capteurs\\\scriptsize ON/OFF};
\draw[very thick,rounded corners=4pt,fill=blue!5] (2.6,-1.35) rectangle (9.0,1.35);
\node[font=\small] at (5.8,0.95) {bus \nt{GLUT}: données};
\node[font=\scriptsize,align=center] at (4.3,0.25) {coïncidences, délais,\\logique (chap.~\ref{sec:glutcomputer})};
\node[font=\scriptsize,align=center] at (7.4,0.25) {mémoire, code\\(chap.~\ref{sec:code}, \ref{sec:acet})};
\draw[thick,red!70!black,fill=red!8,rounded corners=2pt] (2.8,-1.15) rectangle (8.8,-0.55);
\node[font=\scriptsize,red!60!black] at (5.8,-0.85) {\nt{GABA} (ch.~\ref{sec:gaba}): veto, choix, division, rythme};
\node[blk,fill=blue!5,text width=2.0cm] (A) at (10.9,0) {choix de\\l'action};
\draw[->,very thick] (S) -- (2.6,0);
\draw[->,very thick] (9.0,0) -- (A);
\draw[->,very thick] (A) -- (12.6,0) node[right,font=\small] {muscles};
\pic[scale=1.1] at (13.2,-0.5) {spring};
\pic[scale=1.15] at (0.4,0.85) {eyepic};
\pic[scale=1.05] at (1.35,0.85) {speaker};
% registers
\node[reg] (D) at (3.4,3.2) {\nt{DOPA}};
\node[reg] (C) at (5.6,3.2) {\nt{SERO}};
\node[reg] (N) at (7.8,3.2) {\nt{NORA}};
\node[reg] (H) at (10.0,3.2) {\nt{ACET}};
\node[font=\scriptsize,align=center,above=1pt] at (D.north) {erreur $\delta$};
\node[font=\scriptsize,align=center,above=1pt] at (C.north) {niveau, $\tau_\gamma$};
\node[font=\scriptsize,align=center,above=1pt] at (N.north) {gain $g$};
\node[font=\scriptsize,align=center,above=1pt] at (H.north) {écriture $a$};
\draw[bc] (D.south) -- (3.4,1.35);
\draw[bc] (C.south) -- (5.6,1.35);
\draw[bc] (N.south) -- (7.8,1.35);
\draw[bc] (H.south) -- (10.0,2.1) -- (8.8,1.35);
\draw[bc] (H.south east) to[bend left=20] (A.north east);
\draw[->,thick] (2.85,1.35) -- (2.85,2.75) -- (D.west) node[pos=0.25,left,font=\scriptsize] {$V$};
\draw[->,thick,green!45!black] (0.4,3.2) node[left,black,font=\small] {récompense $r$} -- (2.2,3.2) -- (D.west);
\pic[scale=0.9] at (-0.45,3.7) {juice};
\draw[->,thick,densely dotted,gray!50!black] ([yshift=0.45cm]D.north) to[bend left=18] node[above,font=\scriptsize,black] {surprise} ([yshift=0.45cm]N.north);
\end{tikzpicture}
\caption{Toute la machine. Le bus \nt{GLUT} porte les données des capteurs jusqu'au choix de l'action; \nt{GABA} commande le flux à l'intérieur du bus. Les doubles cercles sont les canaux de diffusion; les flèches en tirets diffusent leurs signaux à tout le réseau, chacun avec son propre bouton. \nt{DOPA} reçoit la récompense $r$ et l'attente $V$ d'un critique situé dans le bus (chap.~\ref{sec:dofa}); la flèche en pointillés est l'hypothèse selon laquelle \nt{NORA} apprend la surprise par des erreurs anormalement grandes (chap.~\ref{sec:nora}).}
\label{fig:machine}
\end{figure}

\section{Une seule formule pour tous les boutons}

On peut réunir les boutons dans une seule formule schématique. Ce n'est pas un nouveau modèle, mais un résumé des modèles des chapitres~\ref{sec:glutcomputer}--\ref{sec:acet}: chaque symbole vient de son chapitre.
\begin{empheq}[box=\keybox]{equation}
\begin{aligned}
\tau\,\frac{\dd r_i}{\dd t} \;&=\; -r_i + F\Bigl(g\,\Bigl[\tfrac{1+a}{2}\,x_i + (1-a)\sum_j W_{ij}\,r_j - \sum_k M_{ik}\,q_k - \theta\Bigr]\Bigr),\\
\frac{\dd W_{ij}}{\dd t} \;&=\; \eta\,a\,\delta(t)\,e_{ij}(t),
\qquad \tau_e\,\frac{\dd e_{ij}}{\dd t} = -e_{ij} + r_i\,r_j,
\qquad \delta = r + \frac{\dd V}{\dd t} - \frac{V}{\tau_\gamma} .
\end{aligned}
\label{eq:machine}
\end{empheq}
Ici, $r_i$ sont les fréquences des neurones \nt{GLUT}, $x_i$ l'entrée venant des capteurs, $W_{ij}\ge0$ leurs poids mutuels (chapitre~\ref{sec:glutcomputer}); $q_k$ sont les fréquences des neurones \nt{GABA}, $M_{ik}\ge0$ la force de leur inhibition (chapitre~\ref{sec:gaba}); $e_{ij}$ est la trace d'éligibilité, $\delta$ l'erreur \nt{DOPA} (chapitre~\ref{sec:dofa}); $\tau_\gamma$ l'horizon, fixé par hypothèse par \nt{SERO}; $g$ le gain \nt{NORA} (chapitre~\ref{sec:nora}); $a$ le niveau \nt{ACET} (chapitre~\ref{sec:acet}).

Regardez ce qui manque dans~\eqref{eq:machine}. Il n'y a pas de cadence: tout est écrit en équations différentielles, et chaque neurone évolue de lui-même. Il n'y a pas de mémoire séparée: ce sont les mêmes $W_{ij}$ qui servent au calcul qui retiennent, ainsi que la même activité $r_i$. Il n'y a pas de corrections individuelles pour l'immense majorité des synapses: leur apprentissage est le produit d'une trace locale et d'un seul nombre commun; un fil d'erreur propre à chaque sortie n'existe que dans le circuit qui apprend les mouvements (chapitre~\ref{sec:motorlearning}). Et les signaux de commande ne sont que de quatre sortes, $\delta$, $\tau_\gamma$, $g$, $a$, pour quatre-vingts milliards de neurones. Chacun d'eux n'est pas en réalité un seul nombre, mais un petit faisceau de lignes parallèles (proposition~\ref{prop:bundle}); ces lignes se comptent toutefois par unités.

\begin{keyblock}
\begin{proposition}[Architecture]
\label{prop:architecture}
La machine du cerveau, telle que nous la comprenons aujourd'hui, se décrit en trois couches. Les données circulent sur le bus d'adresses \nt{GLUT}. Des neurones \nt{GABA} adressés orientent, limitent et normalisent ce flux. Le mode de fonctionnement est fixé par quelques canaux de diffusion, chacun formé d'un petit faisceau de lignes parallèles communes à de grandes zones du réseau. Les deux premières couches sont solidement établies; la répartition des rôles entre les canaux de diffusion est pour l'essentiel une hypothèse.
\end{proposition}
\end{keyblock}

\section{Tous les types dans un seul tableau}

Le tableau~\ref{tab:synthesis} rassemble tous les types de neurones du livre. Les quatre types qui n'ont pas eu leur propre chapitre y occupent chacun une ligne; deux d'entre eux ont trouvé un rôle dans les chapitres sur les sorties de la machine: \nt{GLYC} dans les boucles de la moelle épinière (chapitre~\ref{sec:muscles}), \nt{ADRE} dans la sortie chimique (chapitre~\ref{sec:chemoutput}). Le rôle des deux autres dans le calcul est soit simple, soit inconnu. Les substances qui ne définissent pas de type de neurone sont réunies dans l'annexe~\ref{sec:othermediators}.

\begin{table}[t]
\centering
\footnotesize
\setlength{\tabcolsep}{3pt}
\renewcommand{\arraystretch}{1.15}
\begin{tabular}{>{\raggedright}p{1.3cm}>{\raggedright}p{1.0cm}>{\raggedright}p{3.9cm}>{\raggedright}p{1.5cm}>{\raggedright}p{1.8cm}>{\raggedright\arraybackslash}p{3.8cm}}
\toprule
Type & Chap. & Ce qu'il calcule & Temps & Bus & Ce qu'on sait \\
\midrule
\nt{GLUT} & \ref{sec:glutcomputer}, \ref{sec:code}, \ref{sec:sensors}, \ref{sec:motorlearning}, \ref{sec:dendrites} & données: coïncidences, délais, logique, mémoire, séquences & ms & adressé & établi \\
\nt{GABA} & \ref{sec:gaba}, \ref{sec:code}, \ref{sec:pain}, \ref{sec:motorlearning} & commande du flux: veto, choix, division, stabilité, rythme; portillon de la douleur & ms & adressé & établi; division de la fréquence --- avec le bruit de fond \\
\nt{SERO} & \ref{sec:serocomputer}, \ref{sec:pain}, \ref{sec:sleep} & régulateur de niveau, lissage; horizon $\tau_\gamma$; volume de la douleur; modes du sommeil & secondes & volume & auto-inhibition mesurée; horizon --- hypothèse \\
\nt{DOPA} & \ref{sec:dofa}, \ref{sec:dofaalt} & erreur de prédiction $\delta$; niveau lent --- prix du temps & $0.1$\,s et minutes & diffusion & $\delta$ mesurée; extensions --- chap.~\ref{sec:dofaalt} \\
\nt{NORA} & \ref{sec:nora}, \ref{sec:pain}, \ref{sec:chemoutput}, \ref{sec:sleep} & gain $g$: chercher ou décider; surprise; \og accélérateur\fg{} des organes & secondes & diffusion & hausse du rapport signal/bruit mesurée; rôle dans la recherche --- hypothèse \\
\nt{ACET} & \ref{sec:acet}, \ref{sec:muscles}, \ref{sec:chemoutput}, \ref{sec:sleep} & écriture ou lecture; vitesse d'apprentissage; sortie vers le muscle; \og frein\fg{} des organes & secondes & les deux & trois effets mesurés; bascule des modes --- hypothèse \\
\nt{GLYC} & \ref{sec:muscles}, \ref{sec:pain} & poids négatif dans la moelle et le tronc cérébral: veto sur le muscle antagoniste & ms & adressé & établi \\
\nt{HIST} & --- & signal global \og reste éveillé\fg{} & lent & diffusion & rôle dans le calcul non étudié \\
\nt{ADRE} & \ref{sec:chemoutput} & \og accélérateur\fg{} de tout le corps par le sang; dans le cerveau, inconnu & lent & diffusion & établi hors du cerveau; rôle dans le cerveau obscur \\
\nt{ASPA} & --- & faible poids positif & ms & adressé & rôle discuté \\
\bottomrule
\end{tabular}
\caption{Tous les types de neurones du livre: ce que chacun calcule et à quel point on le sait.}
\label{tab:synthesis}
\end{table}

\section{Comment les boutons travaillent ensemble}

Jusqu'ici, chaque chapitre tournait un seul bouton. Suivons maintenant un seul événement à travers toute la machine (fig.~\ref{fig:timeline}). Un animal a appris depuis longtemps que l'action $A$ rapporte du jus. À un moment donné, le monde change: $A$ ne rapporte plus de jus, c'est l'action $B$ qui en rapporte, et l'animal ne l'essaie presque jamais.

\emph{Erreur.} Le jus n'arrive pas après $A$, et les neurones \nt{DOPA} répondent par un creux, $\delta<0$ selon la formule de l'erreur~\eqref{eq:ctd}. Les synapses qui viennent de choisir $A$ portent une trace et s'affaiblissent.

\emph{Surprise.} Un seul creux n'est qu'un hasard ordinaire. Mais les creux se répètent, et les erreurs deviennent plus grandes que d'habitude. Selon l'hypothèse du chapitre~\ref{sec:nora}, c'est précisément le signal pour \nt{NORA}: le monde a changé. Le gain $g$ baisse, le choix devient souple, et le bruit mène plus souvent à un essai de $B$.

\emph{Écriture.} Selon l'hypothèse du chapitre~\ref{sec:acet}, après le changement \nt{ACET} monte et fait passer le réseau en mode écriture: les connexions internes qui gardent l'ancienne habitude sont affaiblies, l'entrée des capteurs est renforcée, les poids changent facilement.

\emph{Nouvelle habitude.} L'essai de $B$ rapporte un jus que personne n'attendait: une bouffée, $\delta>0$, et les synapses qui ont choisi $B$ se renforcent. Quelques essais de ce genre, et l'attente $V$ rattrape le nouveau monde; les erreurs redeviennent petites.

\emph{Retour.} La surprise est passée, \nt{NORA} rétablit un gain élevé, le choix redevient ferme; \nt{ACET} redescend, et le réseau, en mode lecture, se sert de ce qu'il a appris. Pendant tout ce temps, \nt{SERO}, par hypothèse, fixe l'horizon $\tau_\gamma$: jusqu'à quel point un jus lointain vaut encore la peine d'être attendu.

\begin{figure}[t]
\centering
\begin{tikzpicture}[>=Stealth,x=1.1cm,y=0.95cm]
\foreach \y/\lab in {4/{$\delta$ (\nt{DOPA})},3/{$g$ (\nt{NORA})},2/{$a$ (\nt{ACET})},1/{choix de $B$}}{
  \draw[gray!60!black] (0,\y-0.35) -- (10,\y-0.35);
  \node[left,font=\scriptsize] at (0,\y) {\lab};}
\draw[densely dashed,gray!60!black] (2,0.4) -- (2,4.55) node[above,font=\scriptsize,black] {le monde a changé};
\draw[densely dashed,gray!60!black] (5,0.4) -- (5,4.55) node[above,font=\scriptsize,black] {premier jus pour $B$};
\pic[scale=0.85] at (2,5.25) {nojuice};
\pic[scale=0.85] at (5,5.25) {juice};
% delta: dips after the change, burst at the first juice for B, then back to zero
\draw[thick,red!70!black] (0,3.65) -- (2.3,3.65) -- (2.3,3.3) -- (2.5,3.3) -- (2.5,3.65) -- (3.2,3.65) -- (3.2,3.3) -- (3.4,3.3) -- (3.4,3.65) -- (4.1,3.65) -- (4.1,3.35) -- (4.3,3.35) -- (4.3,3.65) -- (5.0,3.65) -- (5.0,4.35) -- (5.2,4.35) -- (5.2,3.65) -- (5.8,3.65) -- (5.8,4.1) -- (6.0,4.1) -- (6.0,3.65) -- (6.6,3.65) -- (6.6,3.85) -- (6.8,3.85) -- (6.8,3.65) -- (10,3.65);
% gain: high, drops after surprise, recovers
\draw[thick,blue!60!black] (0,3.2) -- (3.0,3.2) -- (3.4,2.75) -- (5.8,2.75) -- (7.2,3.2) -- (10,3.2);
% ACh: low, rises after the change, falls after learning
\draw[thick,green!45!black] (0,1.75) -- (2.8,1.75) -- (3.3,2.25) -- (6.8,2.25) -- (7.8,1.75) -- (10,1.75);
% choice of B
\draw[thick,black] (0,0.7) -- (3.0,0.7) plot[domain=3:10,samples=40] (\x,{0.7+0.6/(1+exp(-(\x-5.3)*1.8))});
\draw[->,gray!70!black] (0,0.2) -- (10.2,0.2) node[below left,font=\scriptsize,black] {temps};
\end{tikzpicture}
\caption{Un seul événement suivi à travers toute la machine. C'est un schéma qualitatif fondé sur les modèles et les hypothèses des chapitres~\ref{sec:dofa}--\ref{sec:acet}, non le résultat d'un calcul. Après le changement, \nt{DOPA} répond par des creux; les creux répétés font monter, par hypothèse, le signal de surprise, \nt{NORA} baisse le gain, \nt{ACET} active l'écriture; le premier jus pour $B$ donne une bouffée, le choix passe à $B$, et les boutons reviennent à leur place.}
\label{fig:timeline}
\end{figure}

Remarquez qu'aucun bouton ne sait ce que font les autres. Chacun réagit à ses propres indices --- l'erreur, sa taille inhabituelle, l'incertitude --- et l'ordre des opérations s'établit de lui-même, comme dans un circuit asynchrone où chaque bloc se déclenche quand ses entrées sont prêtes. À notre connaissance, ce fonctionnement concerté n'a pas encore été vérifié dans son ensemble: les boutons ont été mesurés séparément, leur action conjointe reste une hypothèse.

\section{En quoi cette machine diffère d'un ordinateur}

\begin{table}[t]
\centering
\footnotesize
\setlength{\tabcolsep}{4pt}
\renewcommand{\arraystretch}{1.15}
\begin{tabular}{>{\raggedright}p{2.2cm}>{\raggedright}p{4.2cm}>{\raggedright\arraybackslash}p{6.6cm}}
\toprule
& Ordinateur ordinaire & Cerveau \\
\midrule
temps & horloge commune, de l'ordre du gigahertz & pas d'horloge; chaque neurone évolue de lui-même, les fenêtres sont fixées par les événements et des rythmes locaux (chap.~\ref{sec:glutcomputer}, \ref{sec:gaba}, \ref{sec:code}) \\
éléments & portes binaires fiables & éléments à seuil analogiques; une synapse perd la plupart des impulsions (chap.~\ref{sec:code}) \\
signe d'une connexion & quelconque & fixé par le neurone émetteur (chap.~\ref{sec:neurontypes}) \\
mémoire & séparée du processeur & dans les connexions mêmes qui calculent, et dans l'activité auto-entretenue (chap.~\ref{sec:glutcomputer}, \ref{sec:acet}) \\
apprentissage & rétropropagation de l'erreur & règles locales, un seul nombre diffusé $\delta$ (chap.~\ref{sec:dofa}) et des fils d'erreur vers les sorties du circuit moteur (chap.~\ref{sec:motorlearning}) \\
commande & registres et bus de commandes & quatre canaux de diffusion, chacun un faisceau de quelques lignes (chap.~\ref{sec:serocomputer}, \ref{sec:dofa}--\ref{sec:acet}) \\
énergie & des dizaines à des centaines de watts par processeur & environ $20$ watts pour tout le cerveau, en moyenne environ une impulsion par seconde et par neurone (chap.~\ref{sec:code}) \\
\bottomrule
\end{tabular}
\caption{Le cerveau et un ordinateur ordinaire.}
\label{tab:computer}
\end{table}

Le tableau~\ref{tab:computer} résume les principales différences. Aucune n'est un défaut que la nature n'aurait pas eu le temps de corriger: chacune fait partie d'une même solution. Sans horloge commune, il faut des capteurs ON et OFF et des détecteurs de coïncidences. Des éléments peu fiables exigent la redondance de nombreux neurones. Une mémoire confondue avec le calcul a besoin d'\nt{ACET} pour que l'écriture ne gâche pas la lecture. Un apprentissage sans fils de retour a besoin de \nt{DOPA} et de bruit. Et la limite énergétique d'une impulsion par seconde oblige tout le langage à être économe et à ne dépenser des impulsions que pour les nouvelles.

\section{Ce que nous ne savons pas}

Le plus honnête est de terminer ce résumé par la liste de ce que nous ne savons pas. Chaque point est un problème pour un ingénieur.

\emph{Le code.} Nous connaissons la capacité du canal impulsionnel et nous savons que le destinataire choisit quelle partie du message il lit (chapitre~\ref{sec:code}). Mais on ignore dans quelle mesure le cortex se sert du temps précis des impulsions et si les neurones ont des \og mots\fg{}.

\emph{L'apprentissage à travers de nombreuses couches.} Un signal diffusé enseigne lentement, et d'autant plus lentement qu'il y a de neurones influant sur le résultat (proposition~\ref{prop:broadcastcost}). On ignore donc comment le cortex règle des milliards de poids dans des circuits à plusieurs couches; il existe des modèles où ce sont des bouffées qui transportent l'erreur d'une couche à l'autre~\cite{Payeur2021}. Une partie de la réponse se trouve peut-être dans les calculs à l'intérieur des branches du neurone lui-même, où un seul neurone se comporte comme un petit réseau multicouche: pour reproduire la réponse d'un modèle de neurone cortical, un réseau artificiel a besoin de cinq à huit couches~\cite{Beniaguev2021}, et les branches des neurones humains savent même calculer le OU exclusif~\cite{Gidon2020}. Un autre candidat est l'apprentissage auto-supervisé: si le cortex apprend à prédire ses propres entrées, l'erreur naît sur place, dans chaque couche, et aucun signal commun n'est nécessaire (chapitre~\ref{sec:dofa})~\cite{Keller2018}.

\emph{Ce que transmettent réellement les canaux de diffusion.} Pour \nt{DOPA}, il existe une image standard et six extensions (chapitre~\ref{sec:dofaalt}); pour \nt{NORA} et \nt{ACET}, des hypothèses plausibles (chapitres~\ref{sec:nora}, \ref{sec:acet}); pour \nt{SERO}, surtout des conjectures.

\emph{La coordination dans le temps.} Nous avons vu comment calculer une fonction, mesurer le temps écoulé depuis un événement et découper le flux en fenêtres par un rythme local. Mais comment un immense réseau sans horloge commune coordonne le travail de régions éloignées, on ne le comprend qu'en partie.

\emph{L'énergie.} Vingt watts pour tout. Nous comprenons pourquoi le code doit être parcimonieux (proposition~\ref{prop:economy}), mais personne ne sait encore construire une machine qui ferait la même chose avec la même puissance~\cite{MehonicKenyon2022}.

Le cerveau est une machine de calcul qui n'a pas été conçue par un ingénieur, mais selon des lois que tout ingénieur reconnaîtra: circuits RC, seuils, rétroactions, filtres, bus et registres de diffusion. Nous n'avons lu son schéma qu'en partie. Ceux qui savent lire les schémas en achèveront la lecture.

\section{La suite du livre}

Ce chapitre a assemblé le noyau de calcul de la machine. Les chapitres suivants en sortent. Les chapitres~\ref{sec:sensors} et~\ref{sec:recognition} portent sur les entrées: comment les capteurs transforment la lumière, le son et les molécules en impulsions, et comment la machine y reconnaît des objets. Les chapitres~\ref{sec:muscles}--\ref{sec:chemoutput} portent sur les sorties et sur ce qui les sert: les muscles, la douleur comme signal d'alarme, l'apprentissage des mouvements par des fils d'erreur séparés et la lente sortie chimique vers le corps. Le chapitre~\ref{sec:debugging} explique comment on débogue la machine de l'extérieur, notamment par des interfaces cerveau-ordinateur. Les chapitres~\ref{sec:eetypes} et~\ref{sec:dendrites} trient à nouveau les neurones, cette fois selon leur fonction électrique et selon le nombre de leurs entrées. Le chapitre~\ref{sec:sleep} décrit ce que fait la machine lorsqu'elle est coupée du monde, et les chapitres~\ref{sec:livingmachine} et~\ref{sec:dishbrain} des machines faites de neurones vivants sur une puce.

\section{Exercices}

Les exercices de ce chapitre relient les résultats de plusieurs chapitres: chacun s'appuie sur au moins deux d'entre eux.

\begin{exercise}[\lvlA]
\label{ex:10:1}
\emph{Le bus et les registres.} Les quatre canaux de diffusion sont des faisceaux d'environ cinq lignes (proposition~\ref{prop:bundle}); chaque ligne change une fois par seconde et a quatre niveaux distincts. En combien d'années la commande transmettrait-elle ce que le bus \nt{GLUT} ($8.6\cdot10^{10}$ neurones à la fréquence~\eqref{eq:energy}, précision de $1\ms$ selon~\eqref{eq:spikeentropy}) transmet en une seconde?
\end{exercise}

\begin{exercise}[\lvlB]
\label{ex:10:2}
\emph{Deux boutons sur une même boucle.} Dans la boucle~\eqref{eq:ei}, les boutons de~\eqref{eq:machine} agissent sur l'excitation: $\tau_E\dot E=-E+g[(1-a)w_{EE}E-w_{EI}I+h]$; ils ne commandent pas \nt{GABA}. Avec les valeurs de la fig.~\ref{fig:ei}, une bouffée de \nt{NORA} porte $g$ à $6$. Pour quel $a$ minimal la boucle est-elle stable? Peut-on tourner les boutons \nt{NORA} et \nt{ACET} indépendamment?
\end{exercise}

\begin{exercise}[\lvlB]
\label{ex:10:3}
\emph{D'où vient le prix de la diffusion.} Les sorties de $N$ neurones sont $y_k=f(u_k)+\xi_k$, avec un bruit gaussien indépendant de variance $\sigma^2$; $R\approx R_0+\sum_kR'_k\xi_k$ avec des $R'_k$ égaux, et \nt{DOPA} diffuse $\delta=R-R_0$. Trouvez le rapport signal/bruit de la correction $\delta\,\xi_jx_j$ en un essai. Comment le nombre d'essais croît-il avec $N$?
\end{exercise}

\begin{exercise}[\lvlB]
\label{ex:10:4}
\emph{Lier un son et un éclair.} Le son atteint le bus en $5\ms$, l'éclair en $30\ms$ plus le délai de la première impulsion~\eqref{eq:latency} ($\tau=10\ms$, $U$ de $1.2\theta$ à $4\theta$); fond de $5$~imp./s par entrée. Choisissez le délai de la ligne sonore et la plus petite fenêtre de coïncidence valable pour tous les contrastes. Combien de fausses liaisons par seconde le fond donne-t-il?
\end{exercise}

\begin{exercise}[\lvlB]
\label{ex:10:5}
\emph{Simulation: qui remarque le changement.} Le critique voit $y=s+\text{bruit}$ ($R=1$); avec une probabilité de $1/200$, $s$ saute à une nouvelle valeur de variance $J^2=4$. Simulez un filtre de Kalman avec $q=0$ qui porte $P$ à $J^2$ lorsque $E\leftarrow0.9E+\delta^2/(P+R)$ dépasse $20$. De combien son erreur est-elle plus petite qu'avec le meilleur $\alpha$ constant?
\end{exercise}

\begin{exercise}[\lvlA]
\label{ex:10:6}
\emph{Combien d'essais pour changer d'habitude.} Le réseau a appris $Q_A=0.8$, $Q_B=0.2$ et choisit selon~\eqref{eq:softmax}, $\sigma=1$, $g=8$. Désormais $A$ donne du jus avec une probabilité de $0.2$; $Q_A\leftarrow Q_A+\alpha(r-Q_A)$, $Q_B$ ne change pas. Après combien de choix de $A$ la probabilité de choisir $A$ tombe-t-elle à $0.6$ pour $\alpha=0.1$ et $0.5$? Et si \nt{NORA} abaisse $g$ à $2$?
\end{exercise}

\begin{exercise}[\lvlC]
\label{ex:10:7}
\emph{Un faisceau au lieu d'un fil.} Les sorties sont $y_k=w_k+\xi_k$, la récompense $R=-\sum_ky_k^2$; une ligne du faisceau diffuse à un groupe de $G$ neurones l'erreur de ses sorties, $w_k\leftarrow w_k+\eta\,\delta_l\xi_k$. Montrez qu'avec le meilleur $\eta$ il faut $\approx(G+2)\ln(1/\varepsilon)$ essais pour atteindre $\sum w_k^2=\varepsilon\sum w_k^2(0)$. Combien de temps apprend un circuit de $10^6$ neurones sur cinq lignes, à raison d'un essai toutes les trois secondes, pour $\varepsilon=0.01$?
\end{exercise}

\chapter{Les entrées de la machine: comment s'y branchent les yeux, les oreilles et les autres capteurs}
\chaptermark{Les entrées de la machine}
\label{sec:sensors}

\section{Le capteur comme transducteur}

Sur la fig.~\ref{fig:machine}, les données entraient dans la machine par la gauche, depuis le bloc \og capteurs\fg{}. Ouvrons maintenant ce bloc. Pour un ingénieur, chaque organe des sens est un \emph{transducteur}, avec un étage d'entrée analogique et un convertisseur analogique-numérique à la fin; seulement, les chiffres en sortie ne sont pas des nombres, mais des impulsions.

Le schéma est le même pour tous les organes des sens (fig.~\ref{fig:transducer}). Une grandeur physique --- lumière, pression, vibration, molécule --- agit sur une protéine réceptrice de la membrane d'une cellule sensorielle. Cette protéine fonctionne comme une porte: elle ouvre un canal ionique, directement ou par une courte chaîne de réactions. Il naît un \emph{courant de récepteur}, et avec lui une tension analogique sur la membrane. Cette tension passe ensuite par un réglage du gain et arrive dans un codeur d'impulsions, le neurone intégrateur déjà connu~\eqref{eq:glutneuron}, qui transforme le niveau en fréquence et en instants d'impulsions. La sortie est presque toujours la même: les neurones sensoriels de l'œil, de l'oreille, de l'odorat et de la peau libèrent du glutamate, si bien que les entrées se branchent directement sur le bus \nt{GLUT}.

\begin{figure}[t]
\centering
\begin{tikzpicture}[>=Stealth,
  blk/.style={draw,very thick,rounded corners=2pt,align=center,font=\footnotesize,minimum height=1.0cm,text width=1.9cm,fill=blue!5}]
\node[align=center,font=\small] (P) at (0.1,0) {lumière, son,\\pression,\\molécule};
\node[blk] (R) at (3.0,0) {récepteur-\\porte};
\node[blk] (G) at (6.5,0) {réglage\\du gain};
\node[blk] (E) at (10.0,0) {codeur\\d'impulsions};
\node[align=center,font=\small] (B) at (13.4,0) {bus\\\nt{GLUT}};
\draw[->,very thick] (P) -- (R);
\foreach \ic/\xx in {sun/-1.1,speaker/-0.3,press/0.5,molecule/1.3}{\pic[scale=0.85] at (\xx,1.05) {\ic};}
\draw[->,very thick] (R) -- node[above,font=\scriptsize] {courant} (G);
\draw[->,very thick] (G) -- node[above,font=\scriptsize] {niveau} (E);
\draw[->,very thick,red!70!black] (E) -- node[above,font=\scriptsize,black] {impulsions} (B);
\draw[decorate,decoration={brace,amplitude=5pt,mirror},gray!60!black] (1.8,-0.8) -- (7.7,-0.8) node[midway,below=5pt,font=\scriptsize,black] {partie analogique};
\draw[decorate,decoration={brace,amplitude=5pt,mirror},gray!60!black] (8.8,-0.8) -- (11.2,-0.8) node[midway,below=5pt,font=\scriptsize,black] {impulsions};
\end{tikzpicture}
\caption{Tout organe des sens comme chaîne de transformations. La protéine réceptrice ouvre un canal et produit un courant; le réglage du gain l'adapte aux conditions; le codeur~\eqref{eq:glutneuron} transforme le niveau en impulsions, qui partent sur le bus \nt{GLUT}. Dans l'œil et l'oreille, les premiers étages ne produisent aucune impulsion: le signal reste analogique tant qu'il n'a pas à voyager loin.}
\label{fig:transducer}
\end{figure}

Un détail est bien connu de l'électronicien. Dans l'œil et dans l'oreille, les toutes premières cellules ne produisent pas d'impulsions du tout: elles transmettent à leurs voisines une tension continue, comme un étage analogique sur une carte. Les impulsions n'apparaissent que là où le signal doit aller loin. Un signal analogique s'atténue et se charge de bruit sur quelques millimètres, alors qu'une impulsion, nous l'avons vu au chapitre~\ref{sec:neurontypes}, arrive au bout du fil avec la même amplitude. On numérise là où commence la ligne longue.

\section{À la limite de la physique}

Les capteurs du cerveau travaillent près des limites physiques de la sensibilité. Dans l'obscurité, le capteur de lumière répond à \emph{un seul photon}: l'absorption d'une seule particule de lumière donne un courant d'environ un picoampère, visible sur le fond de son propre bruit~\cite{Baylor1979}. Le capteur de son détecte un déplacement de son faisceau sensible inférieur au nanomètre~\cite{Hudspeth1989}.

Quand la lumière est faible, la précision n'est pas limitée par le capteur, mais par la lumière elle-même: les photons arrivent au hasard, en flux de Poisson. Si, pendant le temps d'accumulation $T$, le capteur reçoit en moyenne $N=\Phi T$ photons, leur nombre fluctue de $\sqrt N$. Pour qu'un supplément $\Delta N$ soit perceptible, il doit être plusieurs fois plus grand que cette fluctuation, et la fraction discernable du changement de luminosité vaut
\begin{empheq}[box=\keybox]{equation}
\frac{\Delta I}{I} \;\approx\; \frac{k}{\sqrt{\Phi\,T}} .
\label{eq:photonnoise}
\end{empheq}
C'est la même formule que pour le nombre d'impulsions~\eqref{eq:rateerror}: le bruit de grenaille des photons et celui des impulsions obéissent à la même loi. Dans l'obscurité, l'œil n'a qu'un seul moyen d'améliorer sa précision: accumuler les photons plus longtemps, et le temps d'accumulation augmente effectivement jusqu'à quelques dixièmes de seconde. C'est pourquoi, au crépuscule, nous voyons plus lentement.

\section{Plage dynamique: le réglage du gain}

Le problème d'ingénierie le plus difficile pour les capteurs est la plage. L'éclairement varie d'environ dix ordres de grandeur entre une nuit étoilée et la neige au soleil, l'intensité sonore de douze entre le seuil d'audition et le seuil de douleur. Or la fréquence des impulsions va de zéro à quelques centaines par seconde, moins de trois ordres de grandeur. Impossible de faire tenir linéairement une telle entrée dans une telle sortie.

La solution est celle de tout récepteur radio: le \emph{contrôle automatique de gain}. Les réponses des cellules sensorielles de l'œil sont bien décrites par la formule
\begin{empheq}[box=\keybox]{equation}
r \;=\; r_{\max}\,\frac{I}{I+\sigma} ,
\label{eq:nakarushton}
\end{empheq}
où $\sigma$ est la luminosité à laquelle la réponse atteint la moitié de son maximum~\cite{NakaRushton1966}. À elle seule, \eqref{eq:nakarushton} ne couvre que deux ou trois ordres de grandeur. Mais $\sigma$ n'est pas constante: lors d'un travail prolongé sur un fond lumineux, elle suit la luminosité moyenne. La courbe de réponse glisse le long de l'axe logarithmique des luminosités et place chaque fois sa partie raide là où se trouve l'entrée (fig.~\ref{fig:adaptation}). C'est la même division par l'activité totale que l'inhibition shuntante du chapitre~\ref{sec:gaba}~\cite{Carandini2012}; seulement, le dénominateur est ici la luminosité moyenne des dernières secondes.

\begin{figure}[t]
\centering
\begin{tikzpicture}[>=Stealth,x=1.25cm,y=2.8cm]
\draw[->,gray!70!black] (-2,0) -- (6.4,0) node[below left,font=\small,black] {$\log_{10} I$};
\draw[->,gray!70!black] (-2,0) -- (-2,1.15) node[left,font=\small,black] {$r/r_{\max}$};
\foreach \u in {-2,0,2,4,6}{\draw[gray!60!black] (\u,-0.02) -- (\u,0.02); \node[below,font=\scriptsize] at (\u,0) {$\u$};}
\draw[densely dashed,gray!60!black] (-2,0.5) -- (6.2,0.5);
\foreach \s/\c/\lab in {0/blue!60!black/{fond sombre},2/green!45!black/{moyen},4/red!70!black/{clair}}{
  \pgfmathsetmacro{\lo}{max(-2,\s-3)}
  \draw[very thick,\c] (-2,0) -- (\lo,0) plot[domain=\lo:6.2,samples=80] (\x,{1/(1+10^(\s-\x))});
  \node[font=\scriptsize,\c,anchor=west] at (\s+0.15,0.42) {\lab};}
\foreach \s/\ic in {0/moon,2/cloud,4/sun}{\pic at (\s+0.6,0.64) {\ic};}
\end{tikzpicture}
\caption{Réglage du gain d'un capteur de lumière selon~\eqref{eq:nakarushton}. Chaque courbe couvre deux ou trois ordres de grandeur de luminosité; quand on passe à un fond plus clair, $\sigma$ augmente et la courbe glisse le long de l'axe logarithmique. Ensemble, les courbes glissantes couvrent toute la plage.}
\label{fig:adaptation}
\end{figure}

Ce réglage a un prix, déjà rencontré au chapitre~\ref{sec:code}: le capteur ne signale pas la luminosité, mais la luminosité \emph{par rapport au fond}. Une entrée constante cesse peu à peu de provoquer une réponse, alors que chaque changement en provoque une. Dès le départ, les entrées de la machine parlent de changements, et non de niveaux; c'est la même économie d'impulsions que dans la proposition~\ref{prop:economy}~\cite{Barlow1961}. D'où aussi les capteurs ON et OFF du chapitre~\ref{sec:glutcomputer}~\cite{Schiller1992}: les uns signalent que la scène est devenue plus claire, les autres qu'elle est devenue plus sombre.

Les ingénieurs ont repris ce procédé dans les \emph{caméras événementielles}. Une telle caméra ne prend pas d'images au rythme d'une horloge: chacun de ses pixels, de lui-même et sans cadence commune, émet un événement \og plus clair\fg{} ou \og plus sombre\fg{} lorsque le logarithme de la luminosité a changé d'une fraction donnée~\cite{Lichtsteiner2008}. C'est exactement le code ON/OFF à deux fils du chapitre~\ref{sec:glutcomputer}, gravé dans le silicium, et il donne à la caméra une plage et une rapidité hors de portée d'un capteur d'images ordinaire~\cite{Gallego2022}.

\section{L'œil: comprimer pour tenir dans un câble}

L'œil compte environ $10^8$ capteurs de lumière~\cite{Curcio1990}, mais seulement un million environ de neurones de sortie qui portent l'image au cerveau. Avant que le signal ne quitte l'œil, on le comprime donc d'un facteur cent environ. Les principaux procédés de compression nous sont familiers. Chaque neurone de sortie répond à la différence entre la luminosité d'une petite tache et celle de son pourtour: c'est la soustraction des voisins par inhibition, comme au chapitre~\ref{sec:gaba}. Les zones uniformes de l'image ne coûtent presque rien, alors que les bords et les changements sont transmis. Les différents neurones de sortie sont spécialisés: les uns signalent un éclaircissement, d'autres un assombrissement, d'autres encore un mouvement dans une direction donnée (fig.~\ref{fig:direction}); chez la souris, on a dénombré plus de trente canaux de sortie de ce genre~\cite{Baden2016}.

Quel est le débit du câble ainsi obtenu? Chez le cobaye, des enregistrements de neurones de sortie ont donné environ $875$ mille bits par seconde pour $\sim10^5$ de ces neurones; rapporté au million de neurones de sortie de l'homme, cela fait de l'ordre de $10^7$ bits par seconde~\cite{Koch2006}, autant qu'un vieux câble réseau à dix mégabits. Avec une fréquence moyenne de quelques impulsions par seconde par neurone, cela fait quelques bits par impulsion, comme nous l'avions trouvé au chapitre~\ref{sec:code}.

\section{L'oreille: un banc de filtres}

L'oreille résout un autre problème: il faut décomposer le son en fréquences. Un ingénieur installerait un banc de filtres passe-bande, et c'est exactement ce que fait l'oreille, mécaniquement. La vibration sonore court le long d'une cloison élastique dont les propriétés varient progressivement sur sa longueur, et chaque portion résonne à sa propre fréquence. Les capteurs placés sur une portion ne répondent qu'à leur bande (fig.~\ref{fig:filterbank}). Les fréquences sont réparties le long de la cloison de façon à peu près logarithmique: chaque octave reçoit une part semblable des capteurs. Le numéro du capteur qui a réagi, c'est la fréquence: le code de position du chapitre~\ref{sec:code}.

\begin{figure}[t]
\centering
\shorthandoff{:?}% pgfmath ternary below
\begin{tikzpicture}[>=Stealth,x=1.35cm,y=2.2cm]
\draw[->,gray!70!black] (-0.3,0) -- (7.6,0) node[above left,font=\small,black] {fréquence, Hz};
\draw[->,gray!70!black] (-0.3,0) -- (-0.3,1.15) node[left,font=\small,black] {réponse};
\foreach \k/\f in {0/125,1/250,2/500,3/1000,4/2000,5/4000,6/8000}{
  \draw[gray!60!black] (\k,-0.02) -- (\k,0.02); \node[below,font=\scriptsize] at (\k,0) {\f};
  \draw[thick,blue!60!black] plot[domain={max(\k-1.2,-0.3)}:{\k+0.7},samples=60] (\x,{1/(1+((\x-\k)/(\x<\k?0.45:0.2))^4)});}
% a piano keyboard: seven white keys per octave, like the octaves on the axis
\foreach \i in {0,...,48}{\draw[gray!50!black,fill=white,line width=0.3pt] ({\i/7},-0.44) rectangle ({(\i+1)/7},-0.26);}
\foreach \o in {0,...,6}{\foreach \j in {0,1,3,4,5}{\fill[black] ({\o+(\j+1)/7-0.035},-0.35) rectangle ({\o+(\j+1)/7+0.035},-0.26);}}
\end{tikzpicture}
\shorthandon{:?}
\caption{L'oreille comme banc de filtres passe-bande, schéma. Chaque courbe est la réponse des capteurs d'une portion à des sons de fréquences différentes; les fréquences centrales sont espacées d'une octave, comme des touches sur une échelle logarithmique. Les vrais filtres ont un flanc haute fréquence raide et un flanc basse fréquence doux. En bas, un clavier: comme dans l'oreille, chaque octave y occupe la même place.}
\label{fig:filterbank}
\end{figure}

Les résonateurs de l'oreille ne sont pas passifs. Une partie des capteurs fait elle-même vibrer la cloison au rythme du son et amplifie les vibrations faibles d'un facteur de plusieurs milliers en amplitude (de $66$ à $76$\,dB); quand le volume augmente, ce gain diminue~\cite{Ruggero1997,Robles2001}. Pour un ingénieur, c'est un amplificateur à réaction positive placé juste au bord de l'auto-oscillation: la même vie au bord de l'instabilité que celle du réseau du chapitre~\ref{sec:gaba}; près de la frontière, le système est particulièrement sensible aux petits signaux. La compression du volume est faite par ce même amplificateur: il amplifie beaucoup ce qui est faible, peu ce qui est fort.

L'oreille a aussi un second code. Aux basses fréquences, les impulsions des neurones sont en phase avec le son: le neurone se déclenche sur une portion déterminée de chaque onde, même s'il ne se déclenche pas à chaque onde. Chez le cobaye, ce verrouillage de phase commence à faiblir vers $600$\,Hz et reste perceptible jusqu'à environ $3.5$\,kHz~\cite{PalmerRussell1986}. Les instants des impulsions conservent ainsi le temps précis du son et, à partir de lui, la différence d'arrivée du son aux deux oreilles, dont le réseau du chapitre~\ref{sec:glutcomputer} tire la direction~\cite{Grothe2010}. Aux hautes fréquences, les impulsions ne suivent plus l'onde, et il ne reste que le code de position. Une seule oreille utilise les deux codes du chapitre~\ref{sec:code}: le temps là où les neurones suivent, la position là où ils ne suivent pas.

\section{Les autres capteurs}

\begin{table}[t]
\centering
\footnotesize
\setlength{\tabcolsep}{3pt}
\renewcommand{\arraystretch}{1.15}
\begin{tabular}{>{\raggedright}p{1.9cm}>{\raggedright}p{2.6cm}>{\raggedright}p{4.0cm}>{\raggedright\arraybackslash}p{4.6cm}}
\toprule
Sens & Ce qu'il mesure & Comment l'entrée est construite & Procédé du livre \\
\midrule
vue & lumière & $\sim10^8$ capteurs, compression d'un facteur $\sim100$, $\sim10^7$ bit/s & réglage du gain, soustraction des voisins, deux fils, direction du mouvement \\
ouïe & pression de l'air & banc de filtres mécanique, amplificateur actif & code de position et code temporel, délais \\
toucher & pression, vibration & capteurs à adaptation rapide et lente & paire \og filtre passe-haut --- filtre passe-bas\fg{} \\
température & chaleur & capteurs séparés du chaud et du froid & code à deux fils \\
odorat & molécules & $\sim400$ types de récepteurs, un seul type par capteur & code combinatoire: une odeur est l'ensemble des types activés \\
goût & substances des aliments & cinq qualités: sucré, amer, salé, acide, umami & lignes étiquetées; certaines cellules libèrent de l'ATP ou de la sérotonine, et non du glutamate \\
position du corps & étirement des muscles & capteurs de longueur et de vitesse d'étirement & rétroaction pour le régulateur du mouvement \\
\bottomrule
\end{tabular}
\caption{Comment les capteurs sont branchés. Tous les dispositifs de la troisième colonne sont mesurés; la colonne de droite les relie aux procédés des chapitres précédents.}
\label{tab:sensors}
\end{table}

Les autres sens sont réunis dans le tableau~\ref{tab:sensors}, et presque chaque ligne y fait reconnaître un procédé des chapitres précédents. Le toucher utilise une paire de filtres: certains capteurs de pression répondent au contact puis se taisent aussitôt, d'autres maintiennent leur réponse tant que la pression dure. La température est transmise par deux fils, l'un pour le chaud, l'autre pour le froid. L'entrée la plus insolite est l'odorat. L'homme possède environ quatre cents types de récepteurs olfactifs, et chaque neurone olfactif n'en porte qu'un seul type~\cite{Mombaerts2004}. Une odeur n'active pas un seul type, mais un ensemble, et le message, c'est \emph{quels} types ont réagi. Pour un ingénieur, c'est un vecteur binaire d'environ quatre cents composantes, dont le nombre de valeurs possibles est astronomique.

\begin{keyblock}
\begin{proposition}[Entrées de la machine]
\label{prop:sensors}
Chaque sens est un transducteur qui travaille près de la limite physique de sensibilité, comprime une immense plage dynamique par un contrôle automatique de gain et transmet sur le bus \nt{GLUT} avant tout des \emph{changements}. Pour cela, les capteurs emploient les mêmes procédés que la machine elle-même: code à deux fils, code de position, code temporel et division par le fond. La structure des capteurs est mesurée; que leurs codes soient ajustés pour porter le plus d'information par impulsion est une hypothèse bien fondée.
\end{proposition}
\end{keyblock}

Les entrées, comme tout le reste, fonctionnent de manière asynchrone. Chaque capteur émet une impulsion quand il a quelque chose à signaler, et les différents sens arrivent avec des retards différents: le son au bout de quelques millisecondes, la lumière au bout de dizaines de millisecondes. Comment la machine les rassemble en une seule image du monde sans cadence commune, c'est l'un des problèmes de coordination dans le temps du chapitre~\ref{sec:synthesis}.

\section{Exercices}

\begin{exercise}[\lvlA]
\label{ex:11:1}
\emph{Durée d'accumulation des photons.} Au crépuscule, un capteur reçoit $100$ photons par seconde; un supplément se perçoit s'il vaut $3$ fois la fluctuation~\eqref{eq:photonnoise}. (a)~En combien de temps un seul capteur perçoit-il un changement de luminosité de $10\,\%$? (b)~Combien de capteurs additionner pour y parvenir en $0.2\,\text{s}$, et de combien baisse la résolution spatiale?
\end{exercise}

\begin{exercise}[\lvlA]
\label{ex:11:2}
\emph{Combien de bits par impulsion.} (a)~L'œil transmet $\sim10^7$ bit/s par $10^6$ neurones de sortie, à $3$--$5$ impulsions par seconde. Quelle fraction de la capacité~\eqref{eq:spikeentropy} utilise-t-il pour $\Delta t=1\ms$? (b)~Une odeur active $20$ des $\approx400$ types de récepteurs. Combien de bits porte un tel ensemble, et combien de types deux odeurs aléatoires ont-elles en commun en moyenne?
\end{exercise}

\begin{exercise}[\lvlB]
\label{ex:11:3}
\emph{Ce que sait faire une seule courbe~\eqref{eq:nakarushton}.} (a)~Sur quelle plage de luminosités la réponse passe-t-elle de $10\,\%$ à $90\,\%$ de $r_{\max}$? Combien d'ordres de grandeur cela fait-il? (b)~Combien de positions de la courbe glissante faut-il pour couvrir dix ordres de grandeur de luminosité? Et douze ordres de volume sonore?
\end{exercise}

\begin{exercise}[\lvlB]
\label{ex:11:4}
\emph{Limite du code temporel dans l'oreille.} L'écart d'arrivée entre les deux oreilles atteint $0.22\,\text{m}/343\,\text{m/s}$. (a)~Les impulsions suivent la phase d'un son pur: jusqu'à quelle fréquence la direction s'en déduit-elle sans ambiguïté? (b)~Une impulsion fluctue de $0.5\ms$, et il faut distinguer $20$~\textmu s. Combien de neurones à $100$ imp./s faut-il moyenner sur $50\ms$?
\end{exercise}

\begin{exercise}[\lvlB]
\label{ex:11:5}
\emph{Une caméra événementielle sur le câble de l'œil.} Une caméra de $10^6$ pixels, de sortie $10^7$ bit/s, émet un événement (adresse et signe) quand $\ln I$ d'un pixel varie de $\ln1.2$. À quelle vitesse maximale peut défiler un bord long de $500$ pixels, de contraste $4$? Combien d'images par seconde une caméra ordinaire à $8$ bits par pixel passerait-elle par la même sortie?
\end{exercise}

\begin{exercise}[\lvlB]
\label{ex:11:6}
\emph{Simulation: réglage du gain.} Le capteur~\eqref{eq:nakarushton} ajuste $\sigma$ avec $\tau=1\,\text{s}$, en logarithmes, $\tau\,\dd\ln\sigma/\dd t=\ln I-\ln\sigma$, ou linéairement, $\tau\,\dd\sigma/\dd t=I-\sigma$; le fond saute $1\to100\to1$. Au bout de combien de secondes la réponse à un test de $+10\,\%$ revient-elle à la moitié de sa valeur adaptée, après le saut vers le haut et vers le bas, pour chaque loi?
\end{exercise}

\begin{exercise}[\lvlC]
\label{ex:11:7}
\emph{À quel monde la courbe est-elle ajustée.} Pour un petit bruit de sortie constant, l'information sur la luminosité est maximale quand $r(I)=r_{\max}\Phi_I(I)$, où $\Phi_I$ est la fonction de répartition des luminosités. (a)~Pour quelle distribution la courbe~\eqref{eq:nakarushton} est-elle optimale? Quel est l'écart de $\log_{10}I$? (b)~Et pour un bruit de sortie poissonnien?
\end{exercise}

\chapter{Reconnaissance d'images et de vidéos}
\chaptermark{Reconnaissance}
\label{sec:recognition}

\section{Le problème: la même chose avec des pixels différents}

Le chapitre~\ref{sec:sensors} a conduit l'image jusqu'à l'entrée de la machine: environ un million de neurones de sortie de l'œil, un flux comprimé qui transmet surtout des bords et des changements. Mais entre le capteur et l'action, il y a une étape dont nous n'avons encore rien dit. Un chat sur une image n'est pas un ensemble de pixels. Décalez-le d'une demi-image, tournez-le, éloignez-le, éteignez la moitié de la lumière: presque tous les pixels changent, et pourtant la réponse \og chat\fg{} doit rester la même. L'ingénieur appelle cela l'\emph{invariance}: la réponse du classifieur ne doit changer ni avec la translation, ni avec l'échelle, ni avec la rotation, ni avec l'éclairage.

Une expérience simple montre bien la difficulté. Prenons une \og rétine\fg{} unidimensionnelle de $24$ points et deux figures de cinq points qui peuvent se trouver n'importe où. Un classifieur qui compare l'entrée à un modèle mémorisé à une seule position devine juste dans $49$--$52\%$ des cas, autant qu'une pièce jetée à pile ou face. La translation détruit complètement la comparaison pixel par pixel. Il faut un autre schéma.

\section{Un pipeline d'étages}

Nous savons déjà, par le chapitre~\ref{sec:code}, à quelle vitesse le cerveau y parvient: un animal sur une image présentée en un éclair est reconnu en $150\ms$ environ, et le signal traverse une dizaine d'étages, à raison de $10$--$20\ms$ chacun~\cite{Thorpe1996}. À chaque étage, un neurone a le temps d'émettre zéro ou une impulsion. La reconnaissance rapide est donc un seul passage dans le pipeline, sans longue réflexion ni répétition. La question est de savoir ce que fait chaque étage.

La réponse, suggérée par les mesures sur les premiers étages de la vision~\cite{HubelWiesel1962}, consiste en deux opérations alternées (fig.~\ref{fig:conveyor}).

\emph{Sélectivité.} Un neurone de l'étage $S$ regarde une petite zone de l'entrée et ne se déclenche que si elle contient une combinaison précise de parties: tel bord, et tel autre à côté. C'est un ET sur les parties: le neurone à seuil du chapitre~\ref{sec:glutcomputer}, dont le seuil dépasse ce qu'apporte n'importe quelle partie seule.

\emph{Invariance.} Un neurone de l'étage $C$ rassemble les sorties de nombreux neurones $S$ qui cherchent la même combinaison à des endroits différents, et se déclenche si elle a été trouvée n'importe où. C'est un OU sur les positions, ou plus exactement le choix du maximum.

Pour le modèle unidimensionnel, les deux opérations s'écrivent ainsi:
\begin{empheq}[box=\keybox]{equation}
s_{f,p} \;=\; \Bigl[\sum_i t_{f,i}\,x_{p+i} - \theta\Bigr]_+ ,
\qquad
c_f \;=\; \max_p\, s_{f,p} ,
\label{eq:scstage}
\end{empheq}
où $x$ est l'entrée, $t_{f,i}$ le modèle du trait $f$, $p$ la position et $\theta$ le seuil. La première expression rend la réponse sélective, la seconde la rend indépendante de la position. Le réseau obtient le maximum de nombreuses entrées avec les moyens du chapitre~\ref{sec:gaba}: l'inhibition mutuelle ne laisse que l'entrée la plus forte (proposition~\ref{prop:wta}), et la division par l'activité totale égalise les réponses sous des luminosités différentes~\cite{Carandini2012}.

\begin{figure}[t]
\centering
\begin{tikzpicture}[>=Stealth,
  px/.style={draw,minimum size=0.36cm,inner sep=0pt},
  on/.style={px,fill=blue!55!black},
  snode/.style={circle,draw,very thick,minimum size=0.62cm,inner sep=0pt,font=\scriptsize,fill=blue!6},
  cnode/.style={circle,draw,very thick,minimum size=0.8cm,inner sep=0pt,font=\scriptsize,fill=red!10}]
% two inputs: the same shape at two positions
\foreach \row/\sh/\lab in {0/1/{image 1},1/7/{image 2}}{
  \begin{scope}[yshift=-\row*1.0cm]
  \node[left,font=\scriptsize] at (-0.25,0) {\lab};
  \foreach \i in {0,...,13}{\node[px] at (\i*0.4,0) {};}
  \foreach \d in {0,1,3,4}{\node[on] at ({(\sh+\d)*0.4},0) {};}
  \end{scope}}
% S stage: detectors of the shape at several positions
\foreach \k in {0,...,5}{
  \node[snode] (S\k) at ({0.8+0.8*\k},1.7) {$S$};}
\node[font=\scriptsize,left] at (-0.25,1.7) {étage $S$: ET sur les parties};
\draw[->,thick,blue!60!black] (1.2,0.25) -- (S1.south);
\draw[->,thick,orange!85!black] (3.6,0.25) -- (S4.south);
\node[font=\scriptsize,blue!60!black,anchor=east] at (1.25,0.85) {image 1};
\node[font=\scriptsize,orange!85!black,anchor=west] at (3.9,0.85) {image 2};
% C stage
\node[cnode] (C) at (2.8,3.25) {$C$};
\node[font=\scriptsize,left] at (-0.25,3.25) {étage $C$: OU sur les positions};
\foreach \k in {0,...,5}{\draw[->,gray!60!black] (S\k.north) -- (C);}
\draw[->,very thick,red!70!black] (C.north) -- ++(0,0.6) node[above,font=\scriptsize,black] {\og la figure est là\fg{} dans les deux images};
% next stages
\draw[decorate,decoration={brace,amplitude=5pt},gray!60!black] (6.3,1.4) -- (6.3,-1.3);
\node[right,font=\scriptsize,align=left] at (6.5,0.05) {le même\\objet, d'autres\\pixels};
\draw[->,thick,densely dashed,gray!60!black] (3.6,3.25) -- (5.9,3.25) node[right,font=\scriptsize,black,align=left] {paires $S$, $C$ suivantes:\\plus grandes, plus\\complexes, plus invariantes};
\end{tikzpicture}
\caption{Une paire d'étages du pipeline. Les neurones $S$ cherchent la même combinaison de parties à des endroits différents; le neurone $C$ répond si elle a été trouvée n'importe où~\eqref{eq:scstage}. La figure de l'image 1 et celle de l'image 2 excitent des neurones $S$ différents, mais le même neurone $C$. Les paires d'étages suivantes répètent la même chose sur des combinaisons de plus en plus grandes et complexes.}
\label{fig:conveyor}
\end{figure}

Dans le même modèle unidimensionnel, la paire d'étages~\eqref{eq:scstage} reconnaît la figure à n'importe quelle position sans erreur; si chaque point de l'entrée est inversé au hasard avec une probabilité de $0.1$, elle la reconnaît dans $86\%$ des cas, avec une probabilité de $0.2$, dans $70\%$. Pile ou face donne toujours la moitié. Empilez plusieurs de ces paires: la première cherche des bords, la suivante des combinaisons de bords, plus haut des parties d'objets, au sommet des objets entiers. À chaque paire, les combinaisons grandissent, et la réponse dépend de moins en moins de l'endroit et de la façon dont l'objet est placé~\cite{Riesenhuber1999}.

\begin{keyblock}
\begin{proposition}[Pipeline de sélectivité et d'invariance]
\label{prop:conveyor}
La reconnaissance rapide est un seul passage à travers une dizaine d'étages. Dans chaque paire d'étages, le ET sur les parties~\eqref{eq:scstage} rend la réponse sélective, et le OU sur les positions la rend indépendante de la translation. L'alternance de ces deux opérations donne une réponse qui distingue les objets et ne dépend presque pas de l'endroit ni de la façon dont on les voit. La structure des premiers étages et la vitesse du passage sont mesurées; que toute la chaîne se réduise à ces deux opérations est une simplification.
\end{proposition}
\end{keyblock}

Si ce dispositif vous semble familier, c'est normal. C'est précisément cette alternance --- chercher un trait à toutes les positions, puis regrouper les positions voisines --- que les ingénieurs ont transposée dans les réseaux artificiels \emph{convolutifs}~\cite{Fukushima1980}. Et récemment, le chemin a été parcouru en sens inverse: les réseaux convolutifs entraînés à reconnaître des objets se sont révélés le meilleur modèle connu des réponses des neurones des étages supérieurs de la vision~\cite{Yamins2014}. Le résultat, au sommet, se lit simplement: à partir des fréquences d'une centaine de neurones du dernier étage, un bloc de décision linéaire reconnaît l'objet un dixième de seconde après sa présentation~\cite{Hung2005}. C'est le code de fréquence d'un groupe du chapitre~\ref{sec:code}.

\section{Le professeur absent: apprendre sur la vidéo}

On entraîne un réseau convolutif sur des millions d'images étiquetées. Le cerveau, lui, ne reçoit presque aucune étiquette: un enfant voit des objets pendant des heures, mais n'entend que rarement le mot \og tasse\fg{}. Comment les étages $C$ savent-ils alors quels neurones $S$ regrouper? Pour regrouper \og cette figure ici\fg{} avec \og cette figure là\fg{}, il faut déjà savoir que c'est la même figure.

C'est la vidéo elle-même qui donne l'indice. Le monde change de façon continue: un objet dans le champ de vision reste le même objet pendant qu'il se déplace, tourne et change d'éclairage, et ce n'est que de temps en temps que le regard passe à un autre. Tout ce qui change \emph{lentement} se rapporte probablement à l'objet, et tout ce qui change vite, à sa position et à son aspect~\cite{WiskottSejnowski2002}. Il suffit donc au neurone $C$ de la règle de Hebb du chapitre~\ref{sec:dofa} munie d'une trace: lier ce qui s'est succédé, et pas seulement ce qui a été simultané~\cite{Foldiak1991}:
\begin{empheq}[box=\keybox]{equation}
\tau_{\text{tr}}\,\frac{\dd\bar y_k}{\dd t} \;=\; -\bar y_k + y_k ,
\qquad
\frac{\dd\mathbf w_k}{\dd t} \;=\; \eta\,\bar y_k\,\bigl(\mathbf x - \mathbf w_k\bigr) .
\label{eq:traceinvariance}
\end{empheq}
Ici, $y_k$ est l'activité du neurone $C$ numéro $k$, qui a gagné la compétition par l'inhibition, $\bar y_k$ sa trace (le même circuit RC que dans la règle~\eqref{eq:threefactor}) et $\mathbf x$ les sorties des neurones $S$. Le neurone qui a gagné sur le premier aspect de l'objet s'en souvient encore quand arrive le deuxième aspect, et l'attire lui aussi à lui.

\begin{figure}[t]
\centering
\begin{tikzpicture}[>=Stealth,x=1.0cm,y=1.0cm]
\foreach \i/\lab in {0/1,1/2,2/3,3/4}{
  \node[draw,thick,fill=blue!8,minimum width=0.9cm,minimum height=0.55cm,font=\scriptsize] at (\i+0.5,2.2) {$A$ à \lab};}
\foreach \i/\lab in {0/4,1/3,2/2,3/1}{
  \node[draw,thick,fill=orange!12,minimum width=0.9cm,minimum height=0.55cm,font=\scriptsize] at (\i+6.5,2.2) {$B$ à \lab};}
\node[font=\scriptsize,gray!40!black] at (5.0,2.2) {pause};
\draw[->,gray!70!black] (0,0) -- (10.4,0) node[below left,font=\scriptsize,black] {temps};
% traces
\draw[thick,blue!60!black] (0,0.05) .. controls (0.4,1.2) and (0.8,1.3) .. (1.2,1.35) -- (4,1.4) .. controls (4.6,0.5) and (5.2,0.15) .. (6,0.08) -- (10,0.05);
\draw[thick,orange!85!black] (0,0.05) -- (6,0.05) .. controls (6.4,1.2) and (6.8,1.3) .. (7.2,1.35) -- (10,1.4);
\node[font=\scriptsize,blue!60!black,anchor=west] at (1.4,1.65) {trace $\bar y_1$: dure tout le passage de $A$};
\node[font=\scriptsize,orange!85!black,anchor=west] at (7.2,1.65) {trace $\bar y_2$};
\draw[decorate,decoration={brace,amplitude=4pt},gray!60!black] (4.0,2.6) -- (0,2.6);
\draw[decorate,decoration={brace,amplitude=4pt,mirror},gray!60!black] (0,-0.35) -- (4,-0.35) node[midway,below=4pt,font=\scriptsize,black] {tous les aspects de $A$ se lient au neurone 1};
\end{tikzpicture}
\caption{Comment la vidéo enseigne l'invariance, schéma. L'objet $A$ passe par quatre positions, puis vient une pause, puis l'objet $B$. La trace~\eqref{eq:traceinvariance} du neurone qui a gagné sur le premier aspect de $A$ dure tout le passage, et tous les aspects de $A$ se retrouvent liés à un seul neurone. La pause efface la trace, et $B$ revient à un autre neurone.}
\label{fig:tracevideo}
\end{figure}

Nous l'avons vérifié sur un modèle (fig.~\ref{fig:tracevideo}). Deux neurones $C$ sont en compétition pour les entrées de $16$ neurones $S$: deux figures à huit positions. Une figure parcourt toutes les positions, $50\ms$ pour chacune, puis vient une pause de $0.3\,\text{s}$ et la figure aléatoire suivante. Aucune étiquette. Si la trace est courte, $\tau_{\text{tr}}=5\ms$, les neurones se partagent l'entrée selon la \emph{position}, et la réponse ne coïncide pas mieux avec la figure qu'à pile ou face: $52\%$ en moyenne sur dix essais. Si la trace est comparable au temps de présentation d'une position, à partir de $\tau_{\text{tr}}\approx20\ms$, chaque neurone rassemble \emph{toutes} les positions de sa figure (pour $20$, $50$ et $200\ms$, dans les dix essais): l'invariance est apprise sur la seule vidéo. La pause entre les objets est importante: sans elle, la trace déborde sur l'objet suivant et les confond.

On observe la même chose expérimentalement. Si l'on substitue discrètement un objet à un autre pendant que l'œil saute d'un endroit à un autre, les neurones des étages supérieurs se mettent, en une heure ou deux, à répondre à l'objet substitué comme au même objet: ils lient ce qui s'est succédé~\cite{LiDiCarlo2008}. Dans le cerveau, l'invariance s'apprend donc bien sur la continuité du temps. Que cette règle explique tout l'apprentissage visuel reste une hypothèse.

\section{La vidéo, c'est du mouvement}

La vidéo n'est pas seulement un flux d'images pour apprendre, c'est aussi la tâche elle-même: qu'est-ce qui bouge, dans quelle direction et à quelle vitesse. Le premier pas nous est familier: le détecteur de direction du chapitre~\ref{sec:gaba} (fig.~\ref{fig:direction}), où une inhibition retardée éteint le mouvement dans un sens. Ces détecteurs couvrent tout le champ visuel, et on les traite ensuite comme les traits de forme: les étages qui rassemblent de nombreux détecteurs d'une même direction à des endroits différents répondent au mouvement d'un grand objet, où qu'il se trouve. C'est la même paire ET et OU~\eqref{eq:scstage}; seulement, le trait n'est plus \og un bord\fg{}, mais \og un bord qui s'est déplacé en un temps $d$\fg{}.

La vidéo est aussi prévisible: l'image suivante est presque la même que la précédente. Selon l'hypothèse du \emph{codage prédictif}, les étages supérieurs envoient aux étages inférieurs une prédiction, et ce qui monte est surtout l'erreur, ce que la prédiction n'a pas pris en compte~\cite{RaoBallard1999}. C'est la même économie d'impulsions que dans la proposition~\ref{prop:economy}: payer pour les nouvelles, pas pour ce qu'on sait déjà. Certaines observations concordent avec cette hypothèse, mais elle n'est pas démontrée comme principe général du pipeline.

\section{Le cerveau et le réseau convolutif}

Jusqu'où va la ressemblance? Pour la reconnaissance rapide d'un seul objet, elle est vraiment profonde~\cite{Yamins2014}. Mais le cerveau possède aussi ce qu'un simple réseau convolutif n'a pas.

\emph{Rétroactions.} Le pipeline du cerveau n'est pas seulement direct: chaque étage a des connexions vers l'arrière et latérales. Pour les images difficiles --- avec occlusion, sous une mauvaise lumière ---, la réponse des étages supérieurs se forme plus tard, et ces réponses tardives sont mieux décrites par des réseaux récurrents que par des réseaux directs~\cite{Kar2019}. Un seul passage règle les cas faciles; les cas difficiles demandent plusieurs tours.

\emph{Robustesse.} On peut tromper un réseau artificiel par une falsification de pixels invisible à l'œil~\cite{Szegedy2014}: un changement qu'un humain ne voit pas transforme un \og panda\fg{} en \og gibbon\fg{}~\cite{Goodfellow2015}. La vision humaine est bien plus robuste à ces falsifications, mais pas invulnérable: des images qui trompent à la fois de nombreux réseaux infléchissent aussi, en présentation brève, le choix de l'humain~\cite{Elsayed2018}. Pourquoi la robustesse diffère tant reste une question ouverte; les candidats sont le bruit, la division par l'activité totale et les rétroactions.

\emph{Professeur.} Le réseau a besoin de millions d'exemples étiquetés; le cerveau, surtout de vidéo sans étiquettes~\eqref{eq:traceinvariance} et de quelques indications de \nt{DOPA} sur ce qui importe.

\emph{Énergie.} Le cerveau entier consomme environ $20$ watts (proposition~\ref{prop:economy}), et la reconnaissance n'en occupe qu'une partie; un réseau convolutif entraîné sur un accélérateur graphique dépense des centaines de watts.

\section{Illusions: où la machine se trompe, et pourquoi}
\label{sec:illusions}

Un ingénieur qui veut comprendre le circuit d'un autre lui envoie des signaux inhabituels et regarde où il se trompe. Pour la vision, ces signaux ont été inventés depuis longtemps: ce sont les illusions visuelles. Une illusion n'est pas une panne. Chacune désigne un mécanisme précis de la machine, qui fonctionne correctement dans des conditions ordinaires et se trahit sur une image choisie tout exprès.

\emph{Attentes raisonnables.} L'entrée est bruitée, et la machine la complète par ce qui se produit habituellement dans le monde. Les mouvements lents sont plus fréquents que les rapides; quand l'entrée est peu fiable --- par exemple quand l'image est peu contrastée ---, l'estimation de la vitesse se déplace vers les vitesses lentes, et un objet pâle semble se mouvoir plus lentement qu'un objet vif~\cite{Weiss2002}. On explique de même de nombreuses illusions de luminosité: nous ne voyons pas la luminosité du pixel, mais la surface à laquelle cette luminosité a le plus souvent correspondu dans le passé~\cite{Purves2011}. La photo de la robe que les uns voient blanc et or et les autres bleu et noir relève du même mécanisme: l'image est ambiguë, et les gens divergent sur l'éclairage qu'ils supposent inconsciemment~\cite{LaferSousa2015,Wallisch2017}. Les différences entre individus sont mesurées; que le cerveau combine ici l'entrée et l'attente selon les règles de la théorie des probabilités est une hypothèse bien fondée.

\emph{Réglage du gain.} Regardez une cascade pendant une minute, puis des rochers immobiles: les rochers se mettent à monter. Les canaux \og vers le bas\fg{} et \og vers le haut\fg{} forment une paire de fils, comme dans~\eqref{eq:dualrail}; le réglage du gain~\eqref{eq:nakarushton} a atténué le canal fatigué, et l'équilibre de la paire s'est déplacé. Les illusions d'inclinaison et de contraste, où l'entourage modifie l'aspect du centre, s'expliquent par la division par l'activité des voisins~\cite{Schwartz2009}, comme au chapitre~\ref{sec:gaba}. Il existe aussi un exemple instructif de prudence: les taches sombres aux croisements des bandes blanches d'une grille ont longtemps été expliquées par une simple inhibition des voisins, mais des expériences avec une grille modifiée ont montré que cette explication ne tient pas~\cite{SchillerCarvey2005}.

\emph{Compensation des délais.} Le trajet de l'œil aux étages supérieurs prend des dizaines de millisecondes, et un objet en mouvement a le temps de se déplacer entre-temps. Si un point immobile s'allume à côté de lui, l'objet semble en avance sur l'éclair, alors qu'ils coïncidaient~\cite{Nijhawan1994}. Selon une des explications, la machine prolonge le mouvement vers l'avant pour compenser le délai, comme au chapitre~\ref{sec:muscles}: avec une rétroaction retardée, il faut prédire. Cette illusion a plusieurs explications, et le débat n'est pas tranché.

\emph{Erreur dans le code temporel.} Une image immobile faite d'anneaux aux transitions de couleur tranchées semble tourner. Les zones contrastées sont traitées plus vite que les zones pâles, et les détecteurs de direction lisent la différence de délais~\eqref{eq:latency} comme un mouvement~\cite{Conway2005}. L'illusion est déclenchée par de petits sauts involontaires de l'œil et par les clignements: chaque saut renouvelle l'entrée, et les détecteurs voient de nouveau un \og mouvement\fg{}~\cite{OteroMillan2012}. C'est le code temporel du chapitre~\ref{sec:code}, mal lu.

\emph{Deux images en une.} Un cube en fil de fer qui nous présente tantôt une face, tantôt l'autre, ou des images différentes dans les deux yeux: la perception bascule toutes les quelques secondes. Les meilleurs modèles combinent une inhibition mutuelle entre les deux interprétations, une fatigue lente et du bruit~\cite{MorenoBote2007}; c'est exactement le schéma~\eqref{eq:halfcentre} qui engendrait, au chapitre~\ref{sec:muscles}, le rythme de la marche. Le moment de la bascule est fixé par le bruit et la fatigue; dans le modèle~\cite{MorenoBote2007}, le rôle principal revient au bruit, et la fatigue ne fait qu'aider.

Et les réseaux artificiels? Un réseau entraîné à prédire l'image suivante d'une vidéo voit la rotation sur les mêmes anneaux immobiles et ne la voit pas sur les images témoins~\cite{Watanabe2018}. Des réseaux entraînés à débruiter les images et à en augmenter la netteté cèdent aux mêmes illusions de couleur et de luminosité que les humains~\cite{GomezVilla2020}. Une partie des illusions découle donc de la tâche même qu'apprend la machine, et non des particularités du tissu nerveux. Savoir quelles illusions sont propres au seul cerveau est un bon test pour tout modèle de la vision.

\begin{keyblock}
\begin{proposition}[Les illusions, empreintes des mécanismes]
\label{prop:illusions}
Une illusion naît quand un mécanisme utile dans le monde ordinaire reçoit une entrée inhabituelle: l'attente l'emporte sur une entrée peu fiable, le réglage du gain déplace une paire de canaux, la compensation du délai porte l'objet en avant, une différence de délais est lue comme un mouvement, l'inhibition mutuelle avec bruit et fatigue bascule. Chaque illusion est un signal de test qui désigne son mécanisme. Les illusions elles-mêmes sont mesurées; leurs explications vont du bien fondé au discuté.
\end{proposition}
\end{keyblock}

\section{Bilan}

\begin{table}[t]
\centering
\footnotesize
\setlength{\tabcolsep}{4pt}
\renewcommand{\arraystretch}{1.15}
\begin{tabular}{>{\raggedright}p{3.0cm}>{\raggedright}p{4.6cm}>{\raggedright\arraybackslash}p{5.2cm}}
\toprule
Ce que fait la machine & Comment & Ce qu'on sait \\
\midrule
reconnaît en $150\ms$ & un seul passage à travers une dizaine d'étages & vitesse mesurée \\
rend la réponse sélective & ET sur les parties, étage $S$~\eqref{eq:scstage} & mesuré sur les premiers étages \\
rend la réponse invariante & OU sur les positions, étage $C$; inhibition et division & mesuré sur les premiers étages; pour les étages supérieurs, simplification \\
produit la réponse & fréquences d'une centaine de neurones du dernier étage, lecture linéaire & mesuré \\
apprend sans étiquettes & règle de Hebb avec trace~\eqref{eq:traceinvariance} sur une vidéo continue & la substitution d'objet modifie les réponses: mesuré; comme règle principale: hypothèse \\
voit le mouvement & détecteurs de direction, puis la même paire ET/OU & mesuré \\
économise sur le prévisible & l'erreur de prédiction monte & hypothèse \\
résout les cas difficiles & rétroactions, plusieurs tours & réponses tardives et rôle des rétroactions mesurés \\
se trompe de façon prévisible & illusions: attentes, réglage du gain, délais, inhibition mutuelle avec bruit et fatigue & illusions mesurées; explications allant du fondé au discuté \\
\bottomrule
\end{tabular}
\caption{Comment la machine reconnaît les images et les vidéos.}
\label{tab:recognition}
\end{table}

La reconnaissance est assemblée à partir de pièces que nous avions déjà (tableau~\ref{tab:recognition}): le seuil donne le ET sur les parties, l'inhibition mutuelle le OU sur les positions, la division l'indépendance vis-à-vis de la luminosité, la règle de Hebb avec trace le professeur absent. Les ingénieurs ont emprunté à la vision l'alternance de la sélectivité et de l'invariance, et en ont tiré les réseaux convolutifs; en retour, le cerveau n'a pas encore livré l'essentiel: savoir apprendre sur de la vidéo sans étiquettes, presque sans énergie.

\section{Exercices}

\begin{exercise}[\lvlA]
\label{ex:12:1}
\emph{La fréquence sur un seul étage.} Un étage du pipeline dure $15\ms$, les neurones tirent à $20$ impulsions par seconde, leurs bruits sont corrélés avec $c=0.1$~\eqref{eq:correlated}. Quelle est l'erreur relative sur la fréquence d'un groupe de cent neurones, et sa limite pour un groupe arbitrairement grand? Combien de niveaux de fréquence sont discernables sur un étage?
\end{exercise}

\begin{exercise}[\lvlA]
\label{ex:12:2}
\emph{L'énergie d'une reconnaissance.} En un passage de $150\ms$, dix étages de $10^7$ neurones émettent au plus une impulsion chacun, au prix de $10^{-10}$~J. (a)~Quelle part de l'énergie de tout le cerveau pendant ces $150\ms$ coûte la reconnaissance? (b)~Combien de fois plus dépense un accélérateur de $300$~W qui reconnaît une image en $10\ms$?
\end{exercise}

\begin{exercise}[\lvlB]
\label{ex:12:3}
\emph{Les seuils de l'étage $S$.} Une \og rétine\fg{} de $24$ points, les figures $A=11011$ et $B=10101$; le modèle~\eqref{eq:scstage} vaut $+1$ sur les uns de la figure et $-1$ sur les zéros. (a)~Pour quels seuils le neurone $S$ pardonne-t-il un point inversé tout en restant muet sur l'autre figure? (b)~Combien de neurones $S$ faut-il pour dix traits $5\times5$ sur une image $100\times100$?
\end{exercise}

\begin{exercise}[\lvlB]
\label{ex:12:4}
\emph{Combien de temps dure l'une des deux images.} Dans le schéma de bascule~\eqref{eq:halfcentre} avec $\tau\ll\tau_a$, tant que le groupe~1 gagne, $r_2=0$ et $r_1=I-a_1$. (a)~Trouvez la durée d'une victoire pour $I=1$, $\beta=2$, $b=2$, $\tau_a=0.4\,\text{s}$. (b)~Quel $\tau_a$ donne un changement d'image toutes les trois secondes?
\end{exercise}

\begin{exercise}[\lvlB]
\label{ex:12:5}
\emph{Simulation: le prix de l'invariance.} Avec la fonction \texttt{two\_stage} de \texttt{models/\allowbreak recognition\_models.py} ($4000$ essais), trouvez pour quelle probabilité d'inversion d'un point $p$ la paire d'étages $S$, $C$ se trompe une fois sur quatre pour $N=24$. Comment la proportion de bonnes réponses pour $p=0.1$ évolue-t-elle de $N=8$ à $N=192$?
\end{exercise}

\begin{exercise}[\lvlB]
\label{ex:12:6}
\emph{Simulation: la fenêtre de la trace.} Avec dix essais de \texttt{trace\_learning} de \texttt{models/\allowbreak recognition\_models.py} et une pause de $0.3\,\text{s}$, trouvez pour quels $\tau_{\text{tr}}$, de $5\ms$ à $2\,\text{s}$, tous les essais apprennent l'invariance sans erreur. D'où vient la borne supérieure de cette fenêtre?
\end{exercise}

\begin{exercise}[\lvlB]
\label{ex:12:7}
\emph{Un mouvement n'importe où.} Sur des points voisins d'une rangée espacés de $0.1^\circ$ se trouvent des détecteurs de direction du chapitre~\ref{sec:gaba}: l'inhibition venant de $A$ avec un délai $d$ éteint l'excitation venant de $B$ si elles sont séparées de $5\ms$ au plus; au-dessus, un étage $C$. Quels délais faut-il pour que $C$ reste muet pour un mouvement inverse de $5$ à $20^\circ$/s?
\end{exercise}

\begin{exercise}[\lvlC]
\label{ex:12:8}
\emph{Falsification de pixels.} Combien de points inversés faut-il pour changer la réponse du modèle \texttt{two\_stage} ($N=24$) sur une figure propre? Et celle du classifieur \og plus de $3.5$ uns en entrée, c'est $A$\fg{}, lui aussi invariant par translation? Quelle est sa proportion de bonnes réponses pour $p=0.1$, contre $0.86$ pour la paire d'étages?
\end{exercise}

\chapter{Les sorties de la machine: comment le cerveau commande les muscles}
\chaptermark{Les sorties de la machine}
\label{sec:muscles}

\section{La sortie rapide}

Tout ce que la machine fait vite dans le monde, elle le fait avec les muscles. Une parole, un regard, un pas, un sourire sont des contractions musculaires. La seconde sortie, lente, la chimie du corps, est examinée au chapitre~\ref{sec:chemoutput}. Sur la fig.~\ref{fig:machine}, la sortie était une flèche \og muscles\fg{}; voyons maintenant ce qu'il y a derrière cette flèche.

Le dernier maillon de la chaîne, ce sont les neurones \nt{ACET}, ceux-là mêmes dont le tableau~\ref{tab:types} indique \og sortie vers le muscle\fg{}. Chaque fibre musculaire reçoit son entrée d'un seul neurone de ce type, et un neurone commande un groupe de fibres, de quelques centaines à deux mille environ~\cite{Feinstein1955mu}. Dans les muscles qui demandent un réglage fin, les unités sont plus petites; dans les grands muscles des jambes, plus grosses. Le neurone et ses fibres forment une \emph{unité motrice}: le plus petit morceau de muscle que le cerveau puisse activer isolément.

La synapse du neurone sur le muscle est construite tout autrement que les synapses corticales du chapitre~\ref{sec:code}. Là, la plupart des impulsions se perdaient. Ici, chaque impulsion libère plusieurs fois plus d'acétylcholine qu'il n'en faut pour que la fibre se déclenche~\cite{WoodSlater2001}. C'est une synapse avec une \emph{marge de sécurité}: une impulsion du neurone donne exactement une contraction des fibres. À l'intérieur de la machine, on peut se permettre des connexions peu fiables et corriger les erreurs par la redondance; en sortie, une erreur coûte un mouvement, et la fiabilité est achetée directement dans le matériel.

\section{La force: fréquence et nombre d'unités}

Une impulsion donne une contraction isolée, une secousse: la force monte en quelques dizaines de millisecondes, puis retombe. En voici une bonne approximation~\cite{Fuglevand1993}:
\begin{equation}
h(t) \;=\; F_1\,\frac{t}{T_c}\,e^{\,1-t/T_c},
\label{eq:twitch}
\end{equation}
où $T_c$ est le temps de montée, de l'ordre de $50\ms$, et $F_1$ la force d'une secousse. Si les impulsions arrivent plus souvent, les secousses s'additionnent:
\begin{empheq}[box=\keybox]{equation}
F(t) \;=\; \sum_k h\bigl(t-t_k\bigr) .
\label{eq:forcesum}
\end{empheq}
C'est le même filtre passe-bas qui transformait les impulsions en concentration au chapitre~\ref{sec:serocomputer}; seulement, en sortie, ce n'est plus une concentration, mais une force. Le muscle est un convertisseur numérique-analogique des impulsions en force (fig.~\ref{fig:tetanus}). Dans notre modèle, à cinq impulsions par seconde, les secousses ne se recouvrent presque pas et la force saute de $0.2F_1$ à $1.1F_1$; à quinze, elle se maintient déjà entre $1.8F_1$ et $2.2F_1$; à quarante, elle devient presque lisse, autour de $5.4F_1$. Un vrai muscle sature quand les impulsions sont fréquentes, si bien que la somme linéaire~\eqref{eq:forcesum} ne vaut que jusqu'à quelques dizaines d'impulsions par seconde.

\begin{figure}[t]
\centering
\begin{tikzpicture}[>=Stealth,x=20cm,y=0.55cm]
\draw[->,gray!70!black] (0,0) -- (0.66,0) node[right,font=\small,black] {$t$, s};
\draw[->,gray!70!black] (0,0) -- (0,6.4) node[left,font=\small,black] {$F/F_1$};
\foreach \t in {0,0.2,0.4,0.6}{\draw[gray!60!black] (\t,-0.1) -- (\t,0.1); \node[below,font=\scriptsize] at (\t,0) {\t};}
\foreach \f in {1,3,5}{\draw[gray!60!black] (-0.004,\f) -- (0.004,\f); \node[left,font=\scriptsize] at (0,\f) {\f};}
\draw[thick,blue!60!black] plot file {figures/muscle_twitch_5.dat};
\draw[thick,green!45!black] plot file {figures/muscle_twitch_15.dat};
\draw[thick,red!70!black] plot file {figures/muscle_twitch_40.dat};
\node[font=\scriptsize,blue!60!black,anchor=west] at (0.605,0.9) {5 imp./s};
\node[font=\scriptsize,green!45!black,anchor=west] at (0.605,2.1) {15 imp./s};
\node[font=\scriptsize,red!70!black,anchor=west] at (0.605,5.4) {40 imp./s};
\end{tikzpicture}
\caption{Le muscle comme convertisseur numérique-analogique: force~\eqref{eq:forcesum} pour des impulsions régulières d'une seule unité motrice, $T_c=50\ms$. Des impulsions rares donnent des secousses séparées; des impulsions fréquentes fusionnent en une force lisse, de même que des impulsions fréquentes fusionnent en une concentration lisse sur la fig.~\ref{fig:lowpass}.}
\label{fig:tetanus}
\end{figure}

Le cerveau dispose de deux boutons pour la force: la fréquence des impulsions de chaque unité et le nombre d'unités activées. Le second bouton est très ingénieux. Les unités d'un muscle diffèrent: les plus faibles sont cent fois plus faibles que les plus fortes. Quand il faut une force de plus en plus grande, les unités s'activent \emph{par ordre de taille}, des petites aux grandes~\cite{Henneman1957}. Personne ne calcule cet ordre: les petits neurones s'excitent simplement plus facilement sous la même entrée commune, si bien que l'ordre d'activation est fixé par le matériel lui-même.

Qu'est-ce que cela apporte? Supposons que chaque unité soit $q$ fois plus forte que la précédente, $q$ étant un peu supérieur à un. La force de toutes les unités activées est la somme d'une progression géométrique, et presque toute cette somme revient aux dernières unités. Chaque nouvelle unité activée ajoute alors
\begin{empheq}[box=\keybox]{equation}
\frac{\Delta F}{F} \;\approx\; q-1 \;=\; \text{const} ,
\label{eq:sizeprinciple}
\end{empheq}
autrement dit, le pas de force est proportionnel à la force elle-même. Pour un effort faible, le cerveau le dose par petites portions; pour un effort fort, par grosses portions. C'est un \emph{nombre à virgule flottante}: l'exposant est donné par le nombre d'unités activées, la mantisse par la fréquence de leurs impulsions. C'est la même échelle logarithmique que celle des capteurs du chapitre~\ref{sec:sensors}: l'entrée et la sortie de la machine mesurent le monde en proportions relatives.

La virgule flottante a son revers. La dispersion de la force croît elle aussi proportionnellement à la force: un mouvement fort est inévitablement moins précis qu'un mouvement faible. Selon une hypothèse influente, c'est justement ce bruit dépendant du signal qui explique pourquoi nos mouvements sont fluides: une commande fluide donne la plus petite dispersion au point d'arrivée~\cite{HarrisWolpert1998}.

\section{Tirer sans pousser: deux fils par articulation}

Un muscle ne sait que tirer. Il ne peut pas pousser, pas plus qu'un neurone \nt{GLUT} ne peut émettre un signal négatif. La solution est connue depuis le chapitre~\ref{sec:glutcomputer}: \emph{deux fils}. Chaque articulation est servie par une paire de muscles qui tirent en sens opposés (fig.~\ref{fig:antagonists}). Si leurs forces sont $F_1$ et $F_2$ et le bras de levier $\ell$, le couple et la raideur de l'articulation valent
\begin{empheq}[box=\keybox]{equation}
M \;=\; \ell\,(F_1-F_2), \qquad K \;\approx\; \kappa_m\,(F_1+F_2),
\label{eq:antagonists}
\end{empheq}
où $\kappa_m$ est le coefficient qui relie la force d'un muscle à sa raideur: un muscle tendu fait un ressort plus raide qu'un muscle relâché. La différence des forces fixe le couple, leur somme la raideur~\cite{Hogan1984}. Avec deux fils, le cerveau commande deux grandeurs indépendantes: dans quel sens tourner l'articulation et avec quelle fermeté la tenir. Pour tenir une tasse sans renverser, nous contractons les deux muscles à la fois: le couple est le même, la raideur plus grande, et une poussée accidentelle dévie moins la main.

\begin{figure}[t]
\centering
\begin{tikzpicture}[>=Stealth,
  nrn/.style={circle,draw,very thick,minimum size=0.95cm,inner sep=0pt,font=\scriptsize},
  inh/.style={-{Bar[width=7pt]},thick,red!70!black}]
% the joint: a bar on a pivot, two springs pulling down on either side
\fill[gray!30] (4.6,-0.25) -- (5.4,-0.25) -- (5,0.15) -- cycle;
\draw[line width=3pt,gray!50!black] (2.4,0.2) -- (7.6,0.2);
\fill[black] (5,0.2) circle (0.08);
\draw[decorate,decoration={coil,segment length=5pt,amplitude=4pt},very thick,blue!60!black] (3.0,0.2) -- (3.0,-1.6);
\draw[decorate,decoration={coil,segment length=5pt,amplitude=4pt},very thick,orange!85!black] (7.0,0.2) -- (7.0,-1.6);
\draw[very thick] (2.5,-1.6) -- (3.5,-1.6); \draw[very thick] (6.5,-1.6) -- (7.5,-1.6);
\node[font=\small,blue!60!black] at (2.35,-0.75) {$F_1$};
\node[font=\small,orange!85!black] at (7.65,-0.75) {$F_2$};
\draw[->,thick] (5.9,0.75) arc (60:120:1.8) node[above,font=\scriptsize,pos=0.5] {$M=\ell(F_1-F_2)$};
% motor neurons and reflex circuit
\node[nrn,fill=green!12] (M1) at (1.6,-3.0) {\nt{ACET}};
\node[nrn,fill=green!12] (M2) at (8.4,-3.0) {\nt{ACET}};
\node[nrn,fill=red!10] (G) at (5.0,-3.0) {\nt{GLYC}};
\node[nrn,fill=blue!6] (S) at (1.6,-5.0) {\nt{GLUT}};
\draw[->,very thick] (M1) -- (3.0,-1.7);
\draw[->,very thick] (M2) -- (7.0,-1.7);
\draw[->,thick] (S) -- (M1);
\draw[->,thick] (S) -- (G);
\draw[inh] (G) -- (M2);
\draw[->,thick,densely dashed,gray!60!black] (3.15,-1.2) to[bend left=25] node[pos=0.35,right,font=\scriptsize,black] {étirement} (S.north east);
\draw[->,thick,blue!60!black] (-0.3,-3.0) node[left,font=\scriptsize,black,align=right] {commande\\du cerveau} -- (M1);
\draw[->,thick,blue!60!black] (10.3,-3.0) node[right,font=\scriptsize,black,align=left] {commande\\du cerveau} -- (M2);
\end{tikzpicture}
\caption{Deux muscles pour une articulation et la boucle de rétroaction la plus rapide. Les neurones \nt{ACET} reçoivent les commandes du cerveau et tirent l'articulation dans des sens opposés; la différence des forces la fait tourner, leur somme la rend plus raide. Le capteur d'étirement (\nt{GLUT}) du muscle de gauche excite son propre neurone \nt{ACET} et, par un intermédiaire \nt{GLYC}, inhibe le neurone du muscle antagoniste: le veto du chapitre~\ref{sec:gaba}.}
\label{fig:antagonists}
\end{figure}

\section{Rétroaction avec retard}

Les muscles ne travaillent pas à l'aveugle. Chacun contient des capteurs d'étirement, dont il a été question dans le tableau~\ref{tab:sensors}, et la boucle la plus courte se ferme directement dans la moelle épinière (fig.~\ref{fig:antagonists}). Si un muscle est étiré à l'improviste, le capteur d'étirement excite son propre neurone \nt{ACET}; le muscle se contracte et ramène l'articulation. En même temps, par un neurone intermédiaire inhibiteur, le muscle antagoniste est coupé pour ne pas gêner. Cet intermédiaire, c'est \nt{GLYC}, le type du tableau~\ref{tab:types} qui n'a pas eu son propre chapitre: sa place est précisément ici, dans les boucles rapides de la moelle épinière.

Chaque boucle a un retard $d$: environ $25$--$30\ms$ pour la boucle spinale du bras, une fois et demie à trois fois plus pour la boucle qui passe par le cortex, $100$--$200\ms$ pour la boucle qui passe par l'œil. Or une rétroaction retardée, nous l'avons vu dans la proposition~\ref{prop:ei}, fait osciller le système. Pour le régulateur le plus simple, qui corrige un écart $x$ à la vitesse $G$ d'après des informations vieilles de $d$ secondes, $\dd x/\dd t=-G\,x(t-d)$, la condition de stabilité est~\cite{Hayes1950}:
\begin{empheq}[box=\keybox]{equation}
G\,d \;<\; \frac{\pi}{2} ,
\label{eq:delaystable}
\end{empheq}
et à la limite, le système oscille à la fréquence $1/(4d)$. Nous l'avons vérifié sur un modèle: pour $d=30\ms$, l'écart s'amortit pour $G=40$ par seconde, et se met à osciller pour $G=55$ par seconde, juste au-dessus de la limite de $52$.

Il en découle deux conséquences. Premièrement, plus la boucle est longue, plus elle est obligée de corriger lentement. Par l'œil, avec un retard de cent cinquante millisecondes, on ne peut pas corriger la main plus vite qu'en un dixième de seconde environ, sinon elle se met à trembler. Deuxièmement, la limite de stabilité de la boucle spinale, $1/(4\cdot30\ms)\approx8$ oscillations par seconde, est proche de la fréquence du petit tremblement que présente toute personne en bonne santé, huit à douze oscillations par seconde. La boucle d'étirement est l'une des sources probables de ce tremblement~\cite{McAuleyMarsden2000}.

Comment le cerveau se meut-il alors vite et précisément si la rétroaction est en retard? Selon une hypothèse répandue, il \emph{prédit}: il entretient un modèle interne du corps et calcule le résultat d'une commande sans attendre les capteurs~\cite{WolpertGhahramani2000}. Les capteurs servent à corriger le modèle, et non chaque mouvement. Un ingénieur parlerait de commande anticipative avec observateur d'état.

\section{Les générateurs de rythme des mouvements}

La marche, la respiration, la mastication sont des mouvements rythmiques. Leur faut-il une horloge? Non. La moelle épinière et le tronc cérébral contiennent de petits réseaux qui produisent eux-mêmes un rythme, même si rien de rythmique ne leur parvient~\cite{MarderBucher2001}. Le plus simple de ces réseaux est fait de pièces que nous possédons déjà: deux groupes de neurones \nt{GLUT} qui s'inhibent mutuellement par des intermédiaires \nt{GABA} ou \nt{GLYC}, comme sur la fig.~\ref{fig:wta}, et qui se fatiguent lentement, comme la mémoire du chapitre~\ref{sec:glutcomputer}:
\begin{empheq}[box=\keybox]{equation}
\tau\,\frac{\dd r_i}{\dd t} = -r_i + \bigl[I - \beta\,r_j - a_i\bigr]_+ ,
\qquad
\tau_a\,\frac{\dd a_i}{\dd t} = -a_i + b\,r_i ,
\label{eq:halfcentre}
\end{empheq}
où $a_i$ est la fatigue du groupe $i$ et $j$ l'autre groupe. L'inhibition mutuelle avec $\beta>1$ choisit un gagnant (proposition~\ref{prop:wta}). Le gagnant travaille et se fatigue; quand la fatigue a mangé son entrée, le perdant, déjà reposé, prend le dessus et commence à son tour à se fatiguer. Les groupes travaillent à tour de rôle, et chacun commande son muscle de la paire (fig.~\ref{fig:halfcentre}).

\begin{figure}[t]
\centering
\begin{tikzpicture}[>=Stealth,x=3.2cm,y=2.4cm]
\draw[->,gray!70!black] (0,0) -- (4.2,0) node[right,font=\small,black] {$t$, s};
\draw[->,gray!70!black] (0,0) -- (0,1.2) node[left,font=\small,black] {$r$};
\foreach \t in {0,1,2,3,4}{\draw[gray!60!black] (\t,-0.03) -- (\t,0.03); \node[below,font=\scriptsize] at (\t,0) {\t};}
\draw[thick,blue!60!black] plot file {figures/halfcentre1.dat};
\draw[thick,orange!85!black] plot file {figures/halfcentre2.dat};
\node[font=\scriptsize,blue!60!black,anchor=west] at (1.6,1.28) {groupe 1: \og avant\fg{}};
\node[font=\scriptsize,orange!85!black,anchor=west] at (2.7,1.28) {groupe 2: \og arrière\fg{}};
\end{tikzpicture}
\caption{Générateur de rythme selon le modèle~\eqref{eq:halfcentre}: $I=1$, $\beta=2$, $b=2$, $\tau=20\ms$, $\tau_a=0.4\,\text{s}$. Deux groupes à inhibition mutuelle et fatigue travaillent à tour de rôle avec une période de $0.84\,\text{s}$, bien que l'entrée reçoive un signal constant.}
\label{fig:halfcentre}
\end{figure}

Dans notre modèle, sous une entrée constante, les groupes se relaient avec une période de $0.84\,\text{s}$, environ le double du temps de fatigue $\tau_a$. Une commande constante venue d'en haut, \og marche\fg{}, se transforme sur place en rythme. C'est un rythme local de plus, comme le rythme de la boucle \nt{GLUT}--\nt{GABA} du chapitre~\ref{sec:gaba}: il n'apparaît que lorsque le réseau travaille, et ne donne la cadence qu'aux muscles qu'il commande.

\begin{keyblock}
\begin{proposition}[Sortie de la machine]
\label{prop:muscles}
Le cerveau commande les muscles comme une hiérarchie de régulateurs. Au sommet, on choisit l'action (chapitre~\ref{sec:dofa}); plus bas, des générateurs locaux transforment des commandes constantes en rythme, et les boucles rapides de la moelle épinière tiennent les articulations. La force est codée en virgule flottante: l'exposant par le nombre d'unités, activées par ordre de taille, la mantisse par la fréquence des impulsions. Chaque articulation est commandée par une paire de muscles qui tirent; la différence de leurs forces fixe le couple, leur somme la raideur. Ces dispositifs sont mesurés; que le cerveau s'appuie sur un modèle interne du corps est une hypothèse bien fondée.
\end{proposition}
\end{keyblock}

\section{Bilan}

\begin{table}[t]
\centering
\footnotesize
\setlength{\tabcolsep}{4pt}
\renewcommand{\arraystretch}{1.15}
\begin{tabular}{>{\raggedright}p{3.0cm}>{\raggedright}p{4.6cm}>{\raggedright\arraybackslash}p{5.2cm}}
\toprule
Ce que fait la sortie & Comment & Ce qu'on sait \\
\midrule
transforme les impulsions en force & somme des secousses~\eqref{eq:forcesum}, filtre passe-bas & mesuré; la somme linéaire est une simplification \\
dose la force & activation des unités par taille, virgule flottante~\eqref{eq:sizeprinciple} & ordre d'activation mesuré \\
pousse en ne sachant que tirer & paire de muscles: couple = différence, raideur = somme~\eqref{eq:antagonists} & mesuré \\
tient l'articulation & boucle d'étirement, veto sur l'antagoniste par \nt{GLYC} & mesuré \\
ne tremble pas & $Gd<\pi/2$~\eqref{eq:delaystable} & découle de la théorie de la régulation; contribution au tremblement: cause probable \\
se meut vite & prédiction par un modèle interne & hypothèse bien fondée \\
marche et respire en rythme & générateur de rythme~\eqref{eq:halfcentre} & générateurs mesurés; le schéma le plus simple est un modèle \\
\bottomrule
\end{tabular}
\caption{Comment la machine commande les muscles.}
\label{tab:muscles}
\end{table}

La sortie de la machine est construite avec les mêmes procédés que tout le reste (tableau~\ref{tab:muscles}): deux fils au lieu d'un signe, un filtre passe-bas au lieu d'un CNA, une échelle logarithmique au lieu d'une échelle linéaire, l'inhibition mutuelle et la fatigue au lieu d'un générateur d'horloge. L'entrée (chapitre~\ref{sec:sensors}) et la sortie mesurent le monde en proportions relatives, et entre elles travaille la machine à laquelle ce livre est consacré.

\section{Exercices}

\begin{exercise}[\lvlA]
\label{ex:13:1}
\emph{La virgule flottante en chiffres.} Un muscle a $N=100$ unités motrices de forces $f_n=f_1q^{\,n-1}$; la plus forte est $100$ fois plus forte que la plus faible. (a)~Trouvez $q$, le pas relatif~\eqref{eq:sizeprinciple} et la plage dynamique en décibels. (b)~Quel est le pas à la force $10f_1$ pour un \og CNA linéaire\fg{} de $100$ unités identiques de même force totale?
\end{exercise}

\begin{exercise}[\lvlA]
\label{ex:13:2}
\emph{Le prix de la marge de sécurité.} À chaque impulsion, la synapse sur une fibre musculaire libère un nombre poissonnien de quanta de moyenne $m$, et la fibre se déclenche à partir de $n_{\text{seuil}}=20$ quanta; la marge de sécurité est $S=m/n_{\text{seuil}}$. Combien d'échecs par jour donne une unité qui tire à $20$ impulsions par seconde, pour $S=2$ et pour $S=4$?
\end{exercise}

\begin{exercise}[\lvlB]
\label{ex:13:3}
\emph{Le muscle comme filtre.} La secousse~\eqref{eq:twitch} avec $T_c=50\ms$ est la réponse impulsionnelle d'un système linéaire. (a)~Trouvez sa fonction de transfert $H(s)$ et sa fréquence de coupure. (b)~À quelle fréquence d'impulsions régulières l'amplitude crête à crête des ondulations de la force vaut-elle $10\%$ de la force moyenne?
\end{exercise}

\begin{exercise}[\lvlB]
\label{ex:13:4}
\emph{On ne corrige pas plus vite que le retard.} Régulateur $\dd x/\dd t=-G\,x(t-d)$. (a)~Montrez que l'écart s'amortit le plus vite sans oscillations pour $Gd=1/e$, à la vitesse $1/d$. (b)~Trouvez ce $G$ et la constante de temps d'amortissement pour la boucle spinale, $d=30\ms$, et pour la boucle par l'œil, $d=150\ms$.
\end{exercise}

\begin{exercise}[\lvlB]
\label{ex:13:5}
\emph{La période du générateur de rythme.} Pour $\tau\ll\tau_a$, la période $T$ du modèle~\eqref{eq:halfcentre} est donnée par l'équation $(\beta-1)(1-E^{2+b})/(\beta-E)=b\,(1-E^{1+b})/(1+b)$, où $E=e^{-T/(2\tau_a)}$. (a)~Trouvez $T$ pour $\beta=b=2$, $\tau_a=0.4$~s. (b)~Montrez que $T$ ne dépend pas de l'entrée $I$ et que le rythme n'est possible que pour $b>\beta-1$.
\end{exercise}

\begin{exercise}[\lvlB]
\label{ex:13:6}
\emph{Simulation du générateur.} Importez la fonction \texttt{halfcentre} de \texttt{models/\allowbreak muscle\_models.py} (sans exécuter le fichier lui-même) et mesurez la période, en écartant les deux premiers cycles, pour $b=0.8$, $1.05$, $1.2$, $1.5$, $2$, $4$ et pour $I=0.5$, $1$, $2$. La commande \og marche plus vite\fg{} peut-elle changer le tempo par $I$?
\end{exercise}

\begin{exercise}[\lvlB]
\label{ex:13:7}
\emph{Raideur contre réflexe.} Une articulation est prise dans une boucle d'étirement: $\dd\theta/\dd t=-a\theta(t)-G\theta(t-d)$, $a=K/B$, $d=30\ms$. Ramenez-la à l'exercice~\ref{ex:13:4} et trouvez le plus petit $a$ pour lequel l'écart s'amortit sans oscillations avec une constante de temps de $15\ms$, ainsi que le $G$ nécessaire. Comment modifier $G$ quand la co-contraction des muscles augmente?
\end{exercise}

\begin{exercise}[\lvlC]
\label{ex:13:8}
\emph{Le bouton du tempo.} Selon les exercices~\ref{ex:13:5} et~\ref{ex:13:6}, l'entrée $I$ du générateur~\eqref{eq:halfcentre} ne change que l'amplitude, pas le tempo. Modifiez au minimum le modèle, avec des pièces du livre: pas de rythme sous un $I_{\min}$, et au-dessus une période qui décroît avec $I$, au moins de moitié quand $I$ triple. Trouvez $I_{\min}$ pour $\tau\ll\tau_a$ et vérifiez par simulation.
\end{exercise}

\chapter{La douleur: un signal d'alarme}
\chaptermark{La douleur}
\label{sec:pain}

\section{Une alarme, pas une mesure}

Tous les capteurs du chapitre~\ref{sec:sensors} mesuraient le monde: combien de lumière, quelle hauteur de son, quelle pression. La douleur fonctionne autrement. Elle ne dit pas \og qu'y a-t-il là\fg{}, mais \og danger, lâche tout\fg{}. Pour un ingénieur, ce n'est pas un canal de mesure, mais un \emph{signal d'alarme}: une ligne de priorité maximale, qui doit percer à travers n'importe quel travail en cours de la machine.

Une alarme se distingue d'un canal de mesure par trois propriétés, et la douleur les a toutes les trois. Premièrement, l'accoutumance l'atténue difficilement: les capteurs de la douleur s'adaptent bien moins que les capteurs du toucher, et après une lésion légère ils deviennent même plus sensibles~\cite{LaMotte1982hyper}; un affaiblissement de la réponse après un fort échauffement n'a été mesuré que pour un seul de leurs types~\cite{Treede1995heat}. Le capteur de lumière se tait sur un fond constant (fig.~\ref{fig:adaptation}), mais une alarme qui se tairait au bout d'une minute serait inutile. Deuxièmement, elle a un seuil élevé: les capteurs de la douleur ne se déclenchent que lorsque l'agression approche du danger; pour l'échauffement, par exemple, les seuils vont de $41^\circ$ à $46$--$48^\circ$ selon le type de capteur~\cite{Treede1995heat, Weidner1999cfib}. Troisièmement --- et c'est l'essentiel pour ce chapitre ---, le volume de l'alarme n'est pas décidé par le seul capteur, mais aussi par la machine elle-même. Une même blessure fait plus ou moins mal selon la situation. Voyons de quelles pièces déjà connues tout cela est fait.

Les capteurs de la douleur sont des neurones \nt{GLUT}, comme les autres neurones sensoriels du chapitre~\ref{sec:sensors}. Certains d'entre eux libèrent, avec le glutamate, la substance~P de l'annexe~\ref{sec:othermediators}, qui renforce lentement la transmission.

\section{Deux fils: douleur rapide et douleur lente}

Cognez-vous le doigt, et vous sentirez la douleur deux fois: d'abord brève et aiguë, puis, une seconde plus tard, sourde et durable. Ce sont deux fils différents. Le fil rapide est isolé et conduit les impulsions à une vitesse $v$ de $5$ à $30$ mètres par seconde; le fil lent est fin et nu, $0.5$--$2$ mètres par seconde. Depuis le pied, le signal parcourt environ un mètre, et le temps d'arrivée
\begin{empheq}[box=\keybox]{equation}
t \;=\; \frac{L}{v}
\label{eq:paintime}
\end{empheq}
est de quelques dizaines de millisecondes pour le fil rapide et d'environ une seconde pour le fil lent. Les deux fils portent deux messages. Le rapide dit \emph{où} et \emph{quand}: ses impulsions arrivent plus tôt et plus précisément, comme dans le code temporel du chapitre~\ref{sec:code}. Le lent dit \emph{à quel point c'est grave}, et sa douleur sourde et durable maintient la machine en mode alarme tant que la lésion n'est pas guérie.

\section{Le portillon de la moelle épinière}

Le premier endroit où arrive le signal de douleur est la moelle épinière, et il y passe par un \emph{portillon}~\cite{MelzackWall1965}; les neurones concrets de ce portillon n'ont été trouvés que relativement récemment~\cite{Duan2014}. Sur le neurone \nt{GLUT} de sortie, qui transmet la douleur plus loin vers le cerveau, convergent le fil de la douleur et le fil du toucher. À côté se trouve un neurone intermédiaire inhibiteur, \nt{GABA} ou \nt{GLYC}, que le toucher excite (fig.~\ref{fig:gate}). Le toucher ferme le portillon. Dans la théorie originale du portillon~\cite{MelzackWall1965}, la douleur, à l'inverse, éteint cet intermédiaire, et une douleur forte ouvre le portillon; c'est une hypothèse de la théorie, qui n'a pas été montrée directement sur les neurones du portillon identifiés.

\begin{figure}[t]
\centering
\begin{tikzpicture}[>=Stealth,
  nrn/.style={circle,draw,very thick,minimum size=1.0cm,inner sep=0pt,font=\scriptsize},
  inh/.style={-{Bar[width=7pt]},thick,red!70!black},
  bc/.style={->,thick,densely dashed,orange!85!black}]
\node[nrn,fill=blue!6] (Z) at (6,0) {\nt{GLUT}};
\node[nrn,fill=red!10] (Q) at (3.6,-1.6) {\nt{GABA}};
\draw[->,very thick,red!70!black] (0,0.6) node[left,font=\small,black,align=right] {douleur $\nu$} -- (5.45,0.15);
\draw[->,very thick,blue!60!black] (0,-2.6) node[left,font=\small,black,align=right] {toucher $\mu$} -- (2.9,-2.0);
\draw[->,thick,blue!60!black] (1.8,-2.35) -- (5.5,-0.35) node[pos=0.8,below right=-2pt,font=\scriptsize] {$w_{z\mu}$};
\draw[inh] (1.6,0.5) -- (3.35,-1.1) node[pos=0.5,left=1pt,font=\scriptsize] {$w_{q\nu}$};
\draw[inh] (Q) -- (Z) node[pos=0.5,above left=-1pt,font=\scriptsize] {$\beta$};
\draw[->,very thick] (Z) -- (8.2,0) node[right,font=\small,align=left] {vers le cerveau: $z$};
\draw[bc] (6,2.1) node[above,font=\scriptsize,black,align=center] {d'en haut: \nt{SERO}, \nt{NORA}, opioïdes} -- (Z.north) node[pos=0.45,right,font=\scriptsize,black] {gain $g$};
\end{tikzpicture}
\caption{Le portillon de la douleur dans la moelle épinière selon le modèle~\eqref{eq:gate}. Le neurone \nt{GLUT} de sortie reçoit la douleur $\nu$ et une faible excitation du toucher $\mu$. L'intermédiaire inhibiteur (\nt{GABA} ou \nt{GLYC}) est excité par le toucher et, selon l'hypothèse de la théorie du portillon, éteint par la douleur. D'en haut, par les canaux de diffusion, arrive le gain $g$.}
\label{fig:gate}
\end{figure}

Écrivons cela pour les fréquences, comme au chapitre~\ref{sec:gaba}. La fréquence de l'intermédiaire $q$ et la fréquence de sortie $z$ valent
\begin{empheq}[box=\keybox]{equation}
q \;=\; \bigl[w_{q\mu}\,\mu + q_0 - w_{q\nu}\,\nu\bigr]_+ ,
\qquad
z \;=\; \bigl[\,g\,(\nu + w_{z\mu}\,\mu) - \beta\,q\,\bigr]_+ ,
\label{eq:gate}
\end{empheq}
où $\nu$ est l'entrée de douleur, $\mu$ l'entrée du toucher, $w$ les poids des connexions (premier indice: vers qui; second: de qui), $\beta$ la force de l'inhibition, $q_0$ l'activité de fond propre de l'intermédiaire et $g$ le gain fixé d'en haut (nous y reviendrons). C'est le veto~\eqref{eq:veto} du chapitre~\ref{sec:gaba}, \og douleur, sauf en cas de toucher\fg{}, avec une correction empruntée à la théorie du portillon à titre d'hypothèse: une douleur forte éteint elle-même le neurone qui pose le veto.

Nous avons calculé le modèle pour $w_{q\mu}=1$, $q_0=0.2$, $w_{q\nu}=0.5$, $w_{z\mu}=0.3$, $\beta=1.5$ (fig.~\ref{fig:gatecurves}). Sans toucher, la douleur commence à passer à partir de $\nu\approx0.18$. Avec le toucher, le seuil monte à $0.86$, et pour $\nu=1$ la sortie est divisée par quatre, de $1$ à $0.25$. Mais pour une douleur forte, $\nu=2$, le toucher ne change plus rien: le portillon est grand ouvert. C'est ce qu'on observe dans la vie: frotter une contusion soulage, mais pas en cas de blessure grave. Le toucher seul, sans douleur, ne passe pas le portillon: l'alarme ne se déclenche pas sur un simple contact.

\begin{figure}[t]
\centering
\begin{tikzpicture}[>=Stealth,x=5.2cm,y=1.35cm]
\draw[->,gray!70!black] (0,0) -- (2.12,0) node[right,font=\small,black] {douleur $\nu$};
\draw[->,gray!70!black] (0,0) -- (0,4.3) node[left,font=\small,black] {$z$};
\foreach \t in {0,0.5,1,1.5,2}{\draw[gray!60!black] (\t,-0.05) -- (\t,0.05); \node[below,font=\scriptsize] at (\t,0) {\t};}
\foreach \f in {1,2,3,4}{\draw[gray!60!black] (-0.01,\f) -- (0.01,\f); \node[left,font=\scriptsize] at (0,\f) {\f};}
\draw[very thick,red!70!black] plot file {figures/pain_gate_notouch.dat};
\draw[very thick,blue!60!black] plot file {figures/pain_gate_touch.dat};
\draw[very thick,orange!85!black,densely dashed] plot file {figures/pain_gate_sens.dat};
\draw[very thick,red!70!black] (0.08,3.9) -- (0.2,3.9) node[right,font=\scriptsize,black] {sans toucher};
\draw[very thick,blue!60!black] (0.08,3.45) -- (0.2,3.45) node[right,font=\scriptsize,black] {avec toucher};
\draw[very thick,orange!85!black,densely dashed] (0.08,3.0) -- (0.2,3.0) node[right,font=\scriptsize,black] {après sensibilisation, $g=2$};
\end{tikzpicture}
\caption{Sortie du portillon~\eqref{eq:gate} en fonction de l'intensité de la douleur. Le toucher relève le seuil et atténue la douleur faible et moyenne, mais pas la douleur forte. La sensibilisation (trait en tirets) double le gain: la même douleur est transmise deux fois plus fort, et le seuil s'abaisse.}
\label{fig:gatecurves}
\end{figure}

\section{Le bouton de volume, d'en haut}

Le portillon n'est pas commandé seulement d'en bas. Du tronc cérébral descendent vers lui des voies qui modifient le gain $g$ de~\eqref{eq:gate}: \nt{SERO}, \nt{NORA} et les peptides opioïdes de l'annexe~\ref{sec:othermediators}~\cite{Fields2004}. Ce sont les mêmes canaux de diffusion que dans les chapitres~\ref{sec:serocomputer}--\ref{sec:acet}; seulement, leur bouton est ici le volume de l'alarme. Ils savent le tourner dans les deux sens: les mêmes circuits descendants peuvent aussi bien atténuer la douleur que l'amplifier.

Pourquoi la machine baisserait-elle l'alarme? Parce que l'alarme est une interruption, et qu'une interruption n'est pas toujours opportune. Un animal qui fuit un prédateur ne doit pas s'arrêter à cause d'une patte blessée: d'abord courir, ensuite lécher la plaie. L'attente d'un soulagement atténue elle aussi la douleur, et une partie de cet effet disparaît si l'on bloque les récepteurs opioïdes~\cite{Levine1978}: l'attente calculée dans le cerveau descend vers le portillon par le canal opioïde. Ce tableau en soi est mesuré; la manière exacte dont le cerveau décide du volume à donner reste une hypothèse.

\section{La sensibilisation: une alarme qui apprend}

Après une lésion, les neurones de sortie du portillon deviennent plus sensibles: la même entrée provoque une sortie plus grande, et le seuil s'abaisse~\cite{Woolf1983}. Dans le modèle~\eqref{eq:gate}, c'est une hausse du gain $g$, et sa description la plus simple est un circuit RC lent, alimenté par la douleur elle-même:
\begin{empheq}[box=\keybox]{equation}
\tau_s\,\frac{\dd g}{\dd t} \;=\; -(g-1) + k_s\,z ,
\label{eq:sensitization}
\end{empheq}
où $\tau_s$ est le temps que met la sensibilité à revenir à la normale; il est plus long que la douleur elle-même, parce que la sensibilisation persiste après l'arrêt du stimulus nocif~\cite{Woolf1983}; $k_s$ est la force de l'alimentation. Tant que le tissu est lésé, la sortie du portillon maintient $g$ au-dessus de un, et tout ce qui se passe près de la blessure est transmis plus fort (trait en tirets sur la fig.~\ref{fig:gatecurves}: pour $g=2$, le seuil descend à $0.11$ et la sortie double). C'est un réglage raisonnable pour une alarme: tant que la lésion n'est pas guérie, mieux vaut être trop prudent.

Pour un ingénieur, c'est une réaction positive à fuite lente: la douleur augmente le gain, le gain augmente la douleur. Tant que la boucle est faible, $g$ revient à la normale une fois la blessure guérie. Mais si la boucle est forte, l'alarme peut rester bloquée en position active même quand la lésion a disparu, comme le groupe à forte rétroaction de la fig.~\ref{fig:bistable}. Le modèle~\eqref{eq:sensitization} est une simplification; l'existence de la sensibilisation centrale est mesurée.

\section{La douleur comme professeur et comme interruption}

Dans la machine, la douleur a deux fonctions en plus de l'alarme.

\emph{Professeur.} La douleur est une récompense négative: dans la formule de l'erreur~\eqref{eq:ctd}, elle entre comme $r<0$. Tout ce qui a été dit au chapitre~\ref{sec:dofa} s'applique ici: un signe annonciateur de la douleur se met à provoquer l'erreur à l'avance, et le réseau apprend à éviter ce qui mène à la douleur. Des expériences chez l'homme montrent que l'activité cérébrale pendant l'apprentissage par la douleur suit précisément l'erreur de différence temporelle~\cite{Seymour2004}, et une partie des neurones \nt{DOPA} répond aussi aux événements désagréables (chapitre~\ref{sec:dofaalt}).

\emph{Interruption.} La douleur arrache l'attention à la tâche en cours: c'est sa fonction, et non un effet secondaire~\cite{EcclestonCrombez1999}. Au chapitre~\ref{sec:nora}, ce même rôle d'\og interruption\fg{} a été discuté pour la bouffée phasique de \nt{NORA}. Et la réponse la plus rapide n'atteint même pas le cerveau: le retrait de la main devant un objet brûlant se referme directement dans la moelle épinière, par la même paire de muscles et le même veto sur l'antagoniste que sur la fig.~\ref{fig:antagonists}.

\begin{keyblock}
\begin{proposition}[Signal d'alarme]
\label{prop:pain}
La douleur n'est pas une mesure, mais un signal d'alarme de haute priorité. Elle s'adapte bien moins que les canaux de mesure, passe par deux fils (le rapide dit \og où\fg{}, le lent \og à quel point c'est grave\fg{}), traverse un portillon que le toucher ferme (et qu'une douleur forte, selon l'hypothèse de la théorie du portillon, ouvre), et son volume est fixé d'en haut par les canaux de diffusion. Ces dispositifs sont mesurés, sauf l'ouverture du portillon par la douleur; les modèles simples~\eqref{eq:gate} et~\eqref{eq:sensitization} sont des simplifications.
\end{proposition}
\end{keyblock}

\section{Bilan}

\begin{table}[t]
\centering
\footnotesize
\setlength{\tabcolsep}{4pt}
\renewcommand{\arraystretch}{1.15}
\begin{tabular}{>{\raggedright}p{3.0cm}>{\raggedright}p{4.9cm}>{\raggedright\arraybackslash}p{4.9cm}}
\toprule
Ce que fait la douleur & Comment & Ce qu'on sait \\
\midrule
donne l'alarme & seuil élevé, accoutumance plus faible que pour le toucher & seuil mesuré; comparaison de l'accoutumance: estimation \\
dit \og où\fg{} et \og à quel point c'est grave\fg{} & deux fils de vitesses différentes~\eqref{eq:paintime} & mesuré \\
passe le portillon & veto par \nt{GABA}/\nt{GLYC}~\eqref{eq:gate} & portillon mesuré; ouverture par la douleur: hypothèse de la théorie; modèle: simplification \\
change de volume & \nt{SERO}, \nt{NORA}, opioïdes descendants & mesuré; règle du choix du volume: hypothèse \\
apprend après une lésion & hausse du gain~\eqref{eq:sensitization} & sensibilisation mesurée; modèle: simplification \\
enseigne l'évitement & récompense négative dans~\eqref{eq:ctd} & ressemblance avec l'erreur de différence temporelle mesurée \\
interrompt le travail & capture de l'attention; réflexe de retrait & mesuré \\
\bottomrule
\end{tabular}
\caption{La douleur comme signal d'alarme.}
\label{tab:pain}
\end{table}

La douleur est assemblée à partir de pièces que nous avons déjà vues (tableau~\ref{tab:pain}): capteurs \nt{GLUT}, deux fils, veto par un intermédiaire inhibiteur, bouton de volume diffusé, circuit RC lent, erreur de prédiction et interruption. Une seule chose y est nouvelle: c'est la seule entrée de la machine dont la machine fixe elle-même le volume, une alarme que son propriétaire peut aussi bien atténuer que reconfigurer.

\section{Exercices}

\begin{exercise}[\lvlA]
\label{ex:14:1}
\emph{Deux fils.} Le signal part du pied, $L=1$~m. (a)~Une salve parcourt un faisceau de fils d'un même type, dont les vitesses remplissent la plage~\eqref{eq:paintime}. De combien s'étale-t-elle à l'arrivée dans la moelle, pour chaque type? (b)~Des vagues de douleur sur des fils à $15$ et $1$~m/s se distinguent si leur décalage atteint $0.1$~s. À partir de quelle distance fusionnent-elles?
\end{exercise}

\begin{exercise}[\lvlB]
\label{ex:14:2}
\emph{Quand le toucher fait mal.} Portillon~\eqref{eq:gate} avec les paramètres du texte; le gain $g$ croît avec la sensibilisation. Jusqu'à quel $g$ un toucher sans douleur ne passe-t-il le portillon pour aucune intensité $\mu$? À partir de quel $g$ un toucher $\mu=1$ commence-t-il à passer?
\end{exercise}

\begin{exercise}[\lvlB]
\label{ex:14:3}
\emph{Stabilité de la sensibilisation.} Les entrées sont constantes, la sortie du portillon est linéaire en $g$: $z=gA-C$, $A=\nu+w_{z\mu}\mu$, $C=\beta q\ge0$. (a)~Trouvez l'équilibre $g^*$ du modèle~\eqref{eq:sensitization} et sa condition de stabilité. (b)~Pour $k_s=0.5$, trouvez l'intensité de douleur au-delà de laquelle le modèle s'emballe, sans toucher et avec un toucher $\mu=1$.
\end{exercise}

\begin{exercise}[\lvlA]
\label{ex:14:4}
\emph{Le signe annonciateur de la douleur.} Un signal à l'instant $0$ prédit une douleur à l'instant $T=2$~s; dans~\eqref{eq:ctd}, $r(t)=-P\,\delta(t-T)$, $\tau_\gamma=2$~s. (a)~Trouvez l'aire de la bouffée de $\delta$ à l'instant du signal. (b)~Une action à l'instant $t_e=1$~s annule la douleur. Quelle est la bouffée de $\delta$ à cet instant, et qu'enseignera-t-elle au réseau?
\end{exercise}

\begin{exercise}[\lvlB]
\label{ex:14:5}
\emph{Une alarme qui se verrouille.} Complétez le portillon de \texttt{models/\allowbreak pain\_gate.py} par la dynamique~\eqref{eq:sensitization} et une saturation $z\le4$. Le toucher $\mu=1$ est toujours présent, la lésion $\nu=2$ dure un temps $T$, puis $\nu=0$. Trouvez par simulation le plus petit $T$ (en unités de $\tau_s$) après lequel l'alarme reste active, pour $k_s=4$, $4.6$, $5$, $6$, $8$.
\end{exercise}

\begin{exercise}[\lvlB]
\label{ex:14:6}
\emph{Concevoir le portillon.} Pour $\beta=1.5$, $w_{z\mu}=0.3$, $g=1$, choisissez dans~\eqref{eq:gate} $q_0$, $w_{q\nu}$, $w_{q\mu}$ de sorte que le seuil de douleur soit $0.2$ sans toucher et $1.0$ avec un toucher $\mu=1$, et que le toucher cesse d'atténuer la douleur à $\nu_\times=2.5$. Jusqu'à quel gain $g$ un toucher sans douleur ne passe-t-il pas le portillon?
\end{exercise}

\begin{exercise}[\lvlC]
\label{ex:14:7}
\emph{Le bouton de volume.} Le travail rapporte une valeur $V$ par unité de temps, travailler avec une lésion $\nu$ coûte $c\nu$, d'où l'intérêt d'interrompre pour $\nu>V/c$; la douleur interrompt quand $z>\theta=0.5$. (a)~Quel gain $g^*$ du portillon~\eqref{eq:gate} sans toucher donne ce seuil pour $V/c=0.25$, $0.5$, $2$? (b)~Par quels signaux du livre la machine estimerait-elle $V$ et $c$?
\end{exercise}

\chapter{Comment la machine apprend à bouger}
\chaptermark{Comment la machine apprend à bouger}
\label{sec:motorlearning}

\section{Trois professeurs}

Au chapitre~\ref{sec:muscles}, nous avons vu comment la machine actionne les muscles, et nous nous sommes arrêtés sur une hypothèse: les mouvements rapides et précis exigent un modèle interne du corps qui prédit le résultat d'une commande~\cite{WolpertGhahramani2000}. D'où vient ce modèle? Il faut l'apprendre, et l'apprendre autrement que tout ce que nous avons appris jusqu'ici.

Le livre a déjà présenté deux professeurs. Le premier, c'est \emph{aucun professeur}: la règle de Hebb (chapitre~\ref{sec:dofa}) renforce ce qui se produit souvent ensemble et comprime les entrées. Le deuxième, c'est la \emph{récompense}: \nt{DOPA} diffuse à tous un seul nombre, \og mieux ou moins bien que prévu\fg{}. Le mouvement a besoin d'un troisième professeur, l'\emph{erreur}. La main a manqué la tasse de deux centimètres vers la gauche: ce n'est pas \og moins bien\fg{}, c'est un vecteur qui dit à chaque sortie dans quel sens et de combien elle s'est trompée. Un ingénieur parlerait d'apprentissage supervisé.

La différence entre le deuxième et le troisième professeur est celle qui sépare un seul fil pour tous d'un fil propre à chaque sortie. Dans la proposition~\ref{prop:broadcastcost}, l'apprentissage par un seul nombre commun ralentissait avec chaque neurone dont le bruit influait sur le résultat. Si, au contraire, chaque sortie $i$ reçoit sa propre erreur $e_i$, on peut changer les poids selon la règle la plus simple:
\begin{empheq}[box=\keybox]{equation}
\frac{\dd w_{ij}}{\dd t} \;=\; -\eta\,e_i(t)\,x_j(t) .
\label{eq:lms}
\end{empheq}
C'est la règle des moindres carrés, connue depuis longtemps dans la technique des filtres adaptatifs~\cite{WidrowHoff1960}: un poids change proportionnellement au produit de son entrée par l'erreur de sa sortie et suit, en moyenne, le gradient du carré de l'erreur. Aucun bruit d'exploration, comme au chapitre~\ref{sec:dofa}, n'est nécessaire: la direction de la correction arrive directement par le fil d'erreur. Et la vitesse ne baisse pas avec le nombre de sorties, parce que les erreurs des autres n'atteignent pas la synapse.

\begin{keyblock}
\begin{proposition}[Trois professeurs]
\label{prop:teachers}
La règle de Hebb enseigne sans professeur ce qui est fréquent; le signal \nt{DOPA} enseigne par un seul nombre pour tout le réseau ce qui est avantageux, et d'autant plus lentement qu'il y a de neurones; le fil d'erreur vers chaque sortie enseigne par la règle~\eqref{eq:lms} ce qui est précis, à une vitesse indépendante du nombre de sorties. Le prix de la troisième méthode: un fil séparé pour chaque sortie, et quelqu'un qui connaît la bonne réponse.
\end{proposition}
\end{keyblock}

\section{Le fil d'erreur}

Qui peut connaître la bonne réponse pour un mouvement? Les capteurs eux-mêmes. Si la main est partie à gauche de la cible, l'œil et les capteurs d'étirement du chapitre~\ref{sec:sensors} le signalent. Il faut un circuit qui achemine cette erreur jusqu'aux bonnes sorties.

Le cerveau possède une région qui semble construite exactement pour cela~\cite{Marr1969,Albus1971}. Son schéma est d'une régularité inhabituelle, presque cristalline, et on y reconnaît facilement la règle~\eqref{eq:lms} (fig.~\ref{fig:errorwire}).

\emph{Couche d'expansion de l'entrée.} Les signaux sur la commande et sur l'état du corps arrivent sur une immense couche de petits neurones \nt{GLUT}. Ils sont environ soixante-dix milliards, plus que dans tout le reste du cerveau réuni~\cite{Azevedo2009}. Chacun mélange quelques entrées, et le signal se décompose en un vecteur très long. L'ingénieur connaît ce procédé: si l'on augmente la dimension de l'entrée, presque n'importe quelle dépendance voulue s'obtient, dans le nouvel espace, par une simple somme pondérée~\cite{Cover1965}.

\emph{Neurones de sortie.} La somme pondérée est calculée par une trentaine de millions de grands neurones de sortie~\cite{Andersen1992cereb}, qui reçoivent chacun de l'ordre de cent mille entrées de la couche d'expansion. Et ce sont des neurones \nt{GABA}: le résultat de l'apprentissage est transmis plus loin non par excitation, mais par inhibition, comme le veto du chapitre~\ref{sec:gaba}.

\emph{Fil d'erreur.} Chaque neurone de sortie reçoit encore une entrée, particulière: exactement un fil \nt{GLUT}, qui se déclenche rarement, environ une fois par seconde, mais provoque chaque fois dans le neurone de sortie une décharge puissante~\cite{Ito2001}. C'est précisément $e_i$: la probabilité de la décharge dépend de la direction de l'erreur de mouvement~\cite{Herzfeld2018}. La coïncidence de la décharge d'erreur avec l'activité d'une entrée affaiblit cette entrée, exactement le terme $-\eta\,e_i x_j$ de~\eqref{eq:lms}.

\begin{figure}[t]
\centering
\shorthandoff{:?}% pgfmath ternary below
\begin{tikzpicture}[>=Stealth,
  gran/.style={circle,draw,thick,fill=blue!6,minimum size=0.32cm,inner sep=0pt},
  outn/.style={circle,draw,very thick,fill=red!10,minimum size=1.1cm,inner sep=0pt,font=\scriptsize},
  inh/.style={-{Bar[width=7pt]},very thick,red!70!black}]
% inputs
\foreach \y/\lab in {1.6/{commande},0.4/{état du corps}}{
  \draw[->,thick,blue!60!black] (-0.6,\y) node[left,font=\scriptsize,black] {\lab} -- (0.35,\y);}
% expansion layer
\foreach \i in {0,...,11}{\node[gran] (g\i) at ({0.8+0.28*mod(\i,2)},{-0.3+0.23*\i}) {};}
\foreach \i in {0,...,11}{\pgfmathsetmacro{\yy}{mod(\i,3)>0 ? 1.6 : 0.4}\draw[gray!60!black] (0.35,\yy) -- (g\i);}
\node[font=\scriptsize,align=center] at (0.95,-1.05) {couche d'expansion\\$\sim7\cdot10^{10}$ \nt{GLUT}};
% parallel wires to the output neuron
\foreach \i in {0,...,11}{\draw[->,gray!70!black] (g\i.east) -- (4.1,{1.0+0.07*(\i-5.5)});}
\node[outn] (P) at (4.6,1.0) {\nt{GABA}};
\node[font=\scriptsize,align=center,above=2pt] at (P.north) {neurone de sortie\\$\sim10^5$ entrées $x_j$, poids $w_{ij}$};
\draw[inh] (P) -- (6.8,1.0) node[right,font=\scriptsize,align=left] {correction\\du mouvement};
% error wire
\draw[->,very thick,red!70!black,densely dashed] (4.6,-1.4) node[below,font=\scriptsize,black,align=center] {un seul fil d'erreur $e_i$\\\nt{GLUT}, $\sim1$ fois/s} -- (P.south);
\end{tikzpicture}
\shorthandon{:?}
\caption{Le circuit d'apprentissage supervisé, tel que le cerveau semble le réaliser. La commande et l'état du corps sont décomposés par une immense couche de neurones \nt{GLUT} en un long vecteur $x_j$; le neurone \nt{GABA} de sortie en fait la somme avec les poids $w_{ij}$. Un unique fil d'erreur apporte $e_i$; la coïncidence de l'erreur avec une entrée active affaiblit le poids de cette entrée, comme dans la règle~\eqref{eq:lms}.}
\label{fig:errorwire}
\end{figure}

Une difficulté nous est familière depuis le chapitre~\ref{sec:dofa}: l'erreur arrive plus tard que l'entrée qui l'a provoquée. Le raté se voit des dizaines ou des centaines de millisecondes après la commande. La synapse doit donc se souvenir qu'elle a été active: il faut une trace, comme dans la règle~\eqref{eq:threefactor}. Et la durée de cette trace, d'après les expériences, est ajustée au retard de la rétroaction propre à chaque tâche: là où l'erreur arrive au bout de cent millisecondes, c'est l'entrée active cent millisecondes avant elle qui s'affaiblit le plus~\cite{Suvrathan2016}.

\begin{keyblock}
\begin{proposition}[Le fil d'erreur]
\label{prop:errorwire}
Dans le circuit représenté sur la fig.~\ref{fig:errorwire}, chaque neurone de sortie reçoit exactement un fil d'erreur, et la coïncidence de l'erreur avec une entrée affaiblit cette entrée. C'est la règle des moindres carrés~\eqref{eq:lms}, appliquée à une entrée étendue à une dimension immense. La structure du circuit et l'affaiblissement par coïncidence sont mesurés; que le circuit entier fonctionne comme un apprentissage supervisé est l'hypothèse dominante.
\end{proposition}
\end{keyblock}

\section{Le modèle interne comme filtre adaptatif}

Qu'apprend un tel circuit? La réponse la plus naturelle: la prédiction. Supposons que l'entrée reçoive une commande $u(t)$ passée à travers un ensemble de filtres de constantes de temps différentes, comme dans la couche d'expansion: $x_j(t)=(g_j*u)(t)$. La sortie est leur somme pondérée, une prédiction de ce que signaleront les capteurs:
\begin{empheq}[box=\keybox]{equation}
\hat y(t) \;=\; \sum_j w_j\,(g_j*u)(t), \qquad e(t) \;=\; \hat y(t) - y(t) ,
\label{eq:adaptivefilter}
\end{empheq}
et les poids apprennent selon~\eqref{eq:lms}. C'est un \emph{filtre adaptatif}: un circuit qui ajuste lui-même sa fonction de transfert pour reproduire un système inconnu~\cite{Fujita1982}. Ici, le système inconnu est le corps lui-même, avec sa masse, son élasticité et ses retards. Le filtre appris est justement le modèle interne du chapitre~\ref{sec:muscles}: il dit ce que signaleront les capteurs, sans les attendre.

Un tel modèle est doublement utile. D'abord, sa prédiction peut remplacer la rétroaction retardée, et la condition de stabilité~\eqref{eq:delaystable} cesse alors de lier les mains: on peut corriger vite, d'après le modèle. Ensuite, la différence entre ce qui était prédit et ce qui est reçu est un signal de l'\emph{inattendu}. Tout ce que la machine a fait elle-même, le modèle l'a prédit et soustrait; il reste ce qu'a fait le monde. C'est pourquoi, selon une hypothèse, nous ne pouvons pas nous chatouiller nous-mêmes~\cite{Blakemore1998}.

\section{Quand le monde change: le rapide et le lent}

Mettons le circuit à l'épreuve d'une expérience facile à monter. Une personne tend la main vers une cible, et le monde change soudain: par exemple, une force latérale s'exerce sur le bras. Les premiers mouvements ratent la cible, puis le raté diminue. Si l'on supprime la force, la main rate d'abord de l'\emph{autre} côté: le modèle a appris la force et continue à la compenser. C'est la preuve directe qu'un modèle a été appris, et pas simplement que \og ça a fini par marcher\fg{}.

Les expériences montrent aussi une chose plus subtile: deux processus apprennent en même temps, un rapide, qui apprend vite et oublie vite, et un lent, qui apprend lentement et se souvient longtemps~\cite{Smith2006}. Les corrections sont mises à jour après chaque mouvement, par événement:
\begin{empheq}[box=\keybox]{equation}
e = p - (x_f + x_s), \qquad x_f \leftarrow A_f\,x_f + B_f\,e, \qquad x_s \leftarrow A_s\,x_s + B_s\,e ,
\label{eq:tworate}
\end{empheq}
où $p$ est la perturbation, $x_f$ et $x_s$ les corrections apprises par les deux processus, $A$ la part de la correction conservée jusqu'au mouvement suivant, $B$ la force avec laquelle elle apprend de l'erreur. Pour le processus rapide, $B$ est grand et $A$ nettement inférieur à un; pour le lent, c'est l'inverse.

\begin{figure}[t]
\centering
\begin{tikzpicture}[>=Stealth,x=0.034cm,y=1.9cm]
\draw[->,gray!70!black] (0,0) -- (355,0) node[right,font=\small,black] {mouvement};
\draw[->,gray!70!black] (0,-1.15) -- (0,1.25);
\foreach \v in {-1,1}{\draw[gray!60!black] (-3,\v) -- (3,\v); \node[left,font=\scriptsize] at (0,\v) {$\v$};}
\foreach \n in {100,200,300}{\draw[gray!60!black] (\n,-0.04) -- (\n,0.04); \node[below,font=\scriptsize] at (\n,0) {\n};}
\draw[thin,gray!70!black] plot file {figures/tworate_p.dat};
\draw[thick,orange!85!black] plot file {figures/tworate_fast.dat};
\draw[thick,blue!60!black] plot file {figures/tworate_slow.dat};
\draw[very thick,black] plot file {figures/tworate_net.dat};
\node[font=\scriptsize,gray!50!black] at (120,1.12) {perturbation $p$};
\node[font=\scriptsize,black,anchor=west] at (238,0.52) {somme: retour};
\node[font=\scriptsize,blue!60!black,anchor=west] at (150,0.39) {lent $x_s$};
\node[font=\scriptsize,orange!85!black,anchor=west] at (60,0.08) {rapide $x_f$};
\draw[densely dashed,gray!60!black] (226,-1.15) -- (226,1.25);
\node[font=\scriptsize,align=center,anchor=south west] at (228,-1.1) {fin des\\erreurs};
\end{tikzpicture}
\caption{Deux processus d'apprentissage selon le modèle~\eqref{eq:tworate}, $A_f=0.92$, $B_f=0.1$, $A_s=0.996$, $B_s=0.02$. Deux cents mouvements avec la perturbation $p=1$, puis six avec la perturbation inverse, $p=-1$, jusqu'à ce que la somme revienne à zéro. Après la ligne en tirets, les erreurs cessent d'arriver, et la somme revient d'elle-même vers l'ancienne correction: le processus rapide a oublié la perturbation inverse, le lent se souvient de la perturbation directe.}
\label{fig:tworate}
\end{figure}

Le modèle~\eqref{eq:tworate} explique une expérience difficile à comprendre sans lui (fig.~\ref{fig:tworate}). Dans notre calcul, après deux cents mouvements avec perturbation, la somme des corrections atteint $0.84$, et c'est le processus lent qui en porte la plus grande part, $0.63$. Ensuite, six mouvements avec la perturbation inverse ramènent la somme à zéro: le processus rapide est passé dans le négatif, $-0.53$, le lent est encore positif, $0.46$. Si l'on cesse alors de signaler les erreurs, le processus rapide oublie sa correction en quelques dizaines de mouvements, le lent bien plus lentement, et la somme remonte \emph{d'elle-même} à $0.37$ en quarante mouvements: l'ancienne habileté revient sans aucun entraînement. Cette récupération spontanée est mesurée chez l'homme; un modèle à un seul processus ne peut pas la produire: sans erreurs, il décroît simplement vers zéro.

Pourquoi deux processus? Le filtre optimal du chapitre~\ref{sec:acet} donne une réponse d'ingénieur plausible. Le corps change pour diverses raisons, à diverses vitesses: les muscles se fatiguent en quelques minutes, grossissent et s'affaiblissent en quelques mois. Si les perturbations sont rapides et lentes, la meilleure estimation se compose de deux filtres de constantes de temps différentes, chacun attrapant la sienne~\cite{Kording2007}. Le processus rapide est une correction provisoire, \og le bras est fatigué\fg{}; le lent, \og le bras a changé\fg{}. Ce n'est encore qu'une hypothèse, mais elle explique pourquoi une habileté apprise en un jour est en partie perdue le lendemain, et pourquoi elle s'apprend plus vite la deuxième fois.

\section{Bilan}

\begin{table}[t]
\centering
\footnotesize
\setlength{\tabcolsep}{4pt}
\renewcommand{\arraystretch}{1.15}
\begin{tabular}{>{\raggedright}p{3.0cm}>{\raggedright}p{4.6cm}>{\raggedright\arraybackslash}p{5.2cm}}
\toprule
Ce que fait le circuit & Comment & Ce qu'on sait \\
\midrule
apprend avec un professeur & fil d'erreur vers chaque sortie, règle~\eqref{eq:lms} & structure du circuit et affaiblissement par coïncidence mesurés; apprentissage supervisé: hypothèse dominante \\
étend l'entrée & soixante-dix milliards de neurones \nt{GLUT} & nombre de neurones mesuré; rôle de l'expansion: théorie \\
tient compte du retard de l'erreur & trace ajustée au retard & mesuré \\
construit un modèle interne & filtre adaptatif~\eqref{eq:adaptivefilter} & hypothèse bien fondée \\
apprend à deux vitesses & processus rapide et lent~\eqref{eq:tworate} & effet consécutif et retour spontané mesurés; deux processus: modèle \\
\bottomrule
\end{tabular}
\caption{Comment la machine apprend à bouger.}
\label{tab:motorlearning}
\end{table}

La machine a maintenant trois professeurs (tableau~\ref{tab:motorlearning}). Hebb comprime ce qui est fréquent; \nt{DOPA}, par un seul nombre, apprend à choisir ce qui est avantageux; le fil d'erreur enseigne la précision, et l'enseigne vite, parce que chaque sortie reçoit sa propre erreur. Le cerveau semble se servir des trois, chacun là où il coûte le moins: le signal diffusé pour un petit nombre de décisions, des fils séparés pour l'ajustement précis du corps, où la bonne réponse est fournie par les capteurs eux-mêmes.

\section{Exercices}

\begin{exercise}[\lvlA]
\label{ex:15:1}
\emph{Capacité du circuit à fil d'erreur.} Le circuit compte $3\cdot10^7$ neurones de sortie de $10^5$ entrées, chacun avec son fil d'erreur ($1$~imp./s). Un neurone à $n$ entrées apprend n'importe quel étiquetage de $\approx2n$ vecteurs~\cite{Cover1965}. (a)~Combien d'associations le circuit stocke-t-il? (b)~Estimez selon~\eqref{eq:spikeentropy}, pour $\Delta t=10\ms$, le temps minimal pour remplir la capacité d'un neurone.
\end{exercise}

\begin{exercise}[\lvlB]
\label{ex:15:2}
\emph{Vitesse de la règle des moindres carrés.} Les modes de la règle~\eqref{eq:lms} s'amortissent aux vitesses $\eta\lambda_k(C)$. Pour cinq filtres~\eqref{eq:adaptivefilter} sur un bruit blanc, $C_{jk}=2\sqrt{\tau_j\tau_k}/(\tau_j+\tau_k)$. Trouvez le rapport entre la vitesse la plus grande et la plus petite si les $\tau_j$ croissent d'un facteur $2$ ($10$--$160\ms$) et d'un facteur $4$.
\end{exercise}

\begin{exercise}[\lvlB]
\label{ex:15:3}
\emph{Analyse des deux processus.} Écrivez le modèle~\eqref{eq:tworate}, avec les paramètres de la fig.~\ref{fig:tworate}, sous la forme $\mathbf x_{n+1}=M\mathbf x_n+\mathbf b\,p$. (a)~Trouvez, à partir des valeurs propres de $M$, les constantes de temps en mouvements. (b)~Trouvez les $x_f$, $x_s$ stationnaires et leur somme pour un $p=1$ constant.
\end{exercise}

\begin{exercise}[\lvlB]
\label{ex:15:4}
\emph{Simulation: plus vite la deuxième fois.} Avec~\eqref{eq:tworate} et les paramètres de \texttt{models/\allowbreak motor\_learning.py}, trouvez le nombre de mouvements jusqu'à $x_f+x_s\ge0.8$ pour $p=1$: (a)~machine non entraînée; (b)~après $200$ mouvements avec $p=1$ et une perturbation inverse jusqu'à somme nulle, comme sur la fig.~\ref{fig:tworate}. Un modèle à un processus le peut-il?
\end{exercise}

\begin{exercise}[\lvlB]
\label{ex:15:5}
\emph{Régulateur à modèle interne.} L'objet $\dd x/\dd t=ku$ est observé avec un retard $d$; le régulateur $u=-G\hat x$ prédit $\hat x(t)=x(t-d)+\hat k\int_{t-d}^{t}u\,\dd s$. (a)~Pour $\hat k=k=1$, $d=150\ms$, trouvez $G$ donnant une constante de temps de $20\ms$; comparez à~\eqref{eq:delaystable}. (b)~Est-ce stable pour $\hat k=0.4k$?
\end{exercise}

\begin{exercise}[\lvlB]
\label{ex:15:6}
\emph{Un filtre adaptatif apprend le muscle.} Objet: la secousse~\eqref{eq:twitch} d'aire unité, $T_c=50\ms$; filtres~\eqref{eq:adaptivefilter}: circuits RC, $\tau_j=12.5$, $25$, $50$, $100$, $200\ms$; entrée: un nouveau nombre normal toutes les $10\ms$. En combien de secondes la règle~\eqref{eq:lms}, pour $\eta=50$ et $200$, réduit-elle l'erreur à $0.5\%$? Pourquoi les poids restent-ils loin de l'optimum même après $300$~s?
\end{exercise}

\begin{exercise}[\lvlC]
\label{ex:15:7}
\emph{Deux processus comme filtre optimal.} La perturbation est la somme de deux processus $p^{f,s}_{n+1}=A_{f,s}\,p^{f,s}_n+w^{f,s}_n$ de variances de bruit $q_f$, $q_s$; l'erreur est observée avec un bruit de variance $R$; le filtre de Kalman donne~\eqref{eq:tworate}. Pour quels $q_f/R$, $q_s/R$ obtient-on, avec $A_f=0.92$, $A_s=0.996$, les valeurs $B_f=0.1$, $B_s=0.02$? Et si l'on quadruple $q_s$?
\end{exercise}

\chapter{La seconde sortie: comment le cerveau pilote la chimie du corps}
\chaptermark{La sortie chimique}
\label{sec:chemoutput}

\section{Une sortie rapide et une sortie lente}

Au chapitre~\ref{sec:muscles}, nous avons vu la sortie rapide de la machine: les muscles. Mais elle en possède une seconde, lente. Le cerveau maintient dans les bonnes limites la température du corps, la fréquence cardiaque, la quantité d'eau et de sucre; il prépare le corps à la course ou au sommeil. Pour cela, il ne bouge pas les articulations: il émet des signaux chimiques. Il dispose de deux voies. La première, ce sont les nerfs qui vont directement aux organes internes: au cœur, aux vaisseaux, à l'intestin, aux glandes. La seconde, c'est le sang: certains neurones et certaines glandes libèrent leur substance non pas dans la fente qui les sépare d'un voisin, mais directement dans la circulation.

Le sang est le bus de diffusion le plus large de tout le livre. Les canaux de diffusion des chapitres~\ref{sec:serocomputer}--\ref{sec:acet} envoyaient un signal à de vastes régions du cerveau; le sang l'envoie à chaque cellule du corps. Le sang fait le tour du corps en une minute environ: le message atteint donc tout le monde en quelques dizaines de secondes et persiste des minutes, voire des heures, jusqu'à ce que la substance soit dégradée. Cette diffusion n'a aucune adresse. Comme dans la proposition~\ref{prop:receptor}, c'est le destinataire qui choisit le sens du message: seules répondent les cellules qui possèdent un récepteur à cette substance.

\section{L'accélérateur et le frein: deux fils vers les organes}

Le meilleur exemple est le cœur. Le cœur possède son propre stimulateur, comme le neurone \nt{SERO} du chapitre~\ref{sec:serocomputer}: il impose lui-même le rythme, et si l'on coupe tous les nerfs, le cœur bat environ cent fois par minute. Le cerveau ne crée pas ce rythme, il ne fait que l'ajuster, et cela par deux fils de signes opposés (fig.~\ref{fig:gasbrake}). Par l'un passe \nt{NORA}: c'est l'\emph{accélérateur}, qui accélère le rythme. Par l'autre passe \nt{ACET}: c'est le \emph{frein}, qui le ralentit. En gros,
\begin{empheq}[box=\keybox]{equation}
f \;\approx\; f_0 \;+\; a_S\,S \;-\; a_P\,P ,
\label{eq:gasbrake}
\end{empheq}
où $f_0\approx100$ battements par minute est le rythme propre du stimulateur, et $S$ et $P$ l'activité de l'accélérateur et du frein. Au repos, le frein est enfoncé, et le cœur bat en général de soixante à soixante-dix fois par minute; à l'effort, on relâche le frein et on appuie sur l'accélérateur.

\begin{figure}[t]
\centering
\begin{tikzpicture}[>=Stealth,
  blk/.style={draw,very thick,rounded corners=2pt,align=center,font=\small},
  pace/.style={circle,draw,very thick,double,double distance=1.2pt,minimum size=1.8cm,inner sep=0pt,font=\scriptsize,fill=red!8,align=center},
  inh/.style={-{Bar[width=7pt]},thick,red!70!black}]
\node[blk,fill=blue!5,text width=1.6cm] (B) at (0,0) {commande\\cérébrale};
\node[pace] (H) at (6.2,0) {stimula-\\teur};
\draw[->,very thick,orange!85!black] (B.north east) to[bend left=20] node[above,font=\scriptsize,black,align=center] {\nt{NORA}: accélérateur, secondes} (H.north west);
\draw[inh,very thick,blue!60!black] (B.south east) to[bend right=20] node[below,font=\scriptsize,black,align=center] {\nt{ACET}: frein, fractions de seconde} (H.south west);
\draw[->,very thick] (H) -- (8.4,0) node[right,font=\small] {fréquence $f$};
\node[blk,fill=orange!10,text width=1.6cm] (A) at (0,-2.6) {\nt{ADRE}\\dans le sang};
\draw[->,thick,orange!85!black] (B) -- (A);
\foreach \y/\lab in {-2.0/{cœur},-2.6/{vaisseaux},-3.2/{foie, muscles}}{
  \draw[->,thick,densely dashed,orange!85!black] (A.east) -- (4.3,\y) node[right,font=\scriptsize,black] {\lab};}
\end{tikzpicture}
\caption{Deux fils de signes opposés vers le stimulateur du cœur: l'accélérateur \nt{NORA} est lent, le frein \nt{ACET} est rapide. En bas, le même accélérateur sur le bus de diffusion: les cellules \nt{ADRE} libèrent l'adrénaline dans le sang, et tous les organes munis du récepteur adéquat reçoivent le signal (flèches en tirets).}
\label{fig:gasbrake}
\end{figure}

C'est le code à deux fils du chapitre~\ref{sec:glutcomputer} et la paire de muscles du chapitre~\ref{sec:muscles}, sauf qu'ici les fils ne commandent pas une articulation, mais la cadence d'un stimulateur. Et les deux fils n'ont pas la même vitesse. Tous deux se comportent comme des filtres passe-bas, mais le frein transmet les changements plusieurs fois plus vite que l'accélérateur~\cite{Berger1989}: le chemin d'\nt{ACET} est court, car la protéine liée au récepteur ouvre elle-même un canal potassique, sans second messager dans la cellule~\cite{Logothetis1987gbg}, tandis que \nt{NORA} agit par une chaîne de réactions intracellulaires. L'ingénieur reconnaît là l'astuce des deux actionneurs de vitesses différentes: le rapide assure l'ajustement fin et instantané, le lent les changements amples et durables.

C'est aussi là qu'\nt{ADRE} trouve son rôle, alors que le tableau~\ref{tab:types} indiquait que son rôle dans le cerveau restait obscur. Hors du cerveau, il est clair. Une partie des cellules de la voie de l'accélérateur n'envoie pas de fil vers un organe, mais libère l'adrénaline directement dans le sang. C'est le même accélérateur, mais sur le bus de diffusion: en quelques dizaines de secondes, il atteint à la fois le cœur, les vaisseaux, le foie et les muscles, et il persiste des minutes. L'accélérateur adressé \nt{NORA} ajuste un organe; l'accélérateur diffusé \nt{ADRE} fait passer tout le corps en mode course.

\section{Une cascade à rétroaction}

Beaucoup d'hormones ne sont pas émises par une seule glande, mais par une cascade. Un petit groupe de neurones à la base du cerveau libère dans le sang un premier signal; celui-ci fait libérer un deuxième signal par une glande intermédiaire; ce dernier fait libérer par la glande finale l'hormone qui agit sur le corps. Et l'hormone finale revient vers les premiers étages pour les freiner. On obtient un amplificateur à trois étages bouclé par une contre-réaction: un régulateur de niveau, comme celui du système \nt{SERO} dans la formule~\eqref{eq:feedback}.

Décrivons chaque étage par un circuit RC de constante de temps $\tau$ et de gain $g$, et la rétroaction par un coefficient $k$:
\begin{equation}
\tau\,\frac{\dd x_1}{\dd t} = -x_1 + g\bigl[u - k\,x_3\bigr]_+ ,\qquad
\tau\,\frac{\dd x_2}{\dd t} = -x_2 + g\,x_1 ,\qquad
\tau\,\frac{\dd x_3}{\dd t} = -x_3 + g\,x_2 ,
\label{eq:cascade}
\end{equation}
où $u$ est la commande du cerveau et $x_3$ le niveau de l'hormone finale. Le gain de boucle vaut $K=g^3k$. À l'équilibre, $x_3=g^3u/(1+K)$ et, comme pour tout régulateur à contre-réaction, une grande boucle rend le niveau insensible à la dispersion des gains des étages. Mais chaque étage décale la phase: à la pulsation $\omega=\sqrt3/\tau$, trois étages identiques donnent ensemble un déphasage d'un demi-tour et une atténuation d'un facteur huit. D'où le résultat classique de l'automatique:
\begin{empheq}[box=\keybox]{equation}
K \;<\; 8 ,
\label{eq:cascadestable}
\end{empheq}
faute de quoi le régulateur se met à osciller, avec une période d'environ $2\pi\tau/\sqrt3\approx3.6\,\tau$.

\begin{keyblock}
\begin{proposition}[Précision contre stabilité]
\label{prop:cascade}
Pour une cascade à trois étages avec rétroaction proportionnelle, l'erreur sur le niveau vaut $1/(1+K)$, et la cascade n'est stable que si $K<8$. Un tel régulateur ne tient donc pas le niveau à mieux qu'environ dix pour cent près; vouloir augmenter le gain le transforme en oscillateur.
\end{proposition}
\end{keyblock}

Nous l'avons vérifié sur le modèle~\eqref{eq:cascade} (fig.~\ref{fig:cascade}). Avec $K=4$, le niveau oscille un peu après la mise en marche, puis se stabilise. Avec $K=7$, il se stabilise aussi, mais lentement. Avec $K=12$, les oscillations ne s'amortissent pas et ne croissent pas sans limite non plus: un étage ne peut pas produire de signal négatif, et la cascade passe en pulsations régulières entre $0.83$ et $1.24$ fois le niveau moyen, avec une période de $3.7$ à $3.9\,\tau$.

\begin{figure}[t]
\centering
\begin{tikzpicture}[>=Stealth,x=0.3cm,y=1.2cm]
\draw[->,gray!70!black] (0,0) -- (41.5,0) node[right,font=\small,black] {$t/\tau$};
\draw[->,gray!70!black] (0,0) -- (0,2.8) node[left,font=\small,black] {$x_3/x_3^*$};
\foreach \t in {0,10,20,30,40}{\draw[gray!60!black] (\t,-0.05) -- (\t,0.05); \node[below,font=\scriptsize] at (\t,0) {\t};}
\foreach \f in {1,2}{\draw[gray!60!black] (-0.3,\f) -- (0.3,\f); \node[left,font=\scriptsize] at (0,\f) {\f};}
\draw[densely dashed,gray!60!black] (0,1) -- (40,1);
\draw[thick,blue!60!black] plot file {figures/cascade_K4.dat};
\draw[thick,red!70!black] plot file {figures/cascade_K12.dat};
\node[font=\scriptsize,blue!60!black,anchor=west] at (16,0.55) {$K=4$: se stabilise};
\node[font=\scriptsize,red!70!black,anchor=west] at (16,1.75) {$K=12$: pulse};
\end{tikzpicture}
\caption{Cascade à trois étages avec rétroaction selon~\eqref{eq:cascade}, niveau de l'hormone finale rapporté à sa valeur d'équilibre $x_3^*$. Quand le gain de boucle est inférieur à huit, le niveau se stabilise; au-delà, la cascade engendre d'elle-même des pulsations régulières.}
\label{fig:cascade}
\end{figure}

Les vraies hormones de cascade sortent effectivement par pulsations: le niveau de l'hormone du stress, par exemple, oscille avec une période d'environ une heure. Selon une hypothèse, ces pulsations naissent de la rétroaction retardée entre les étages elle-même, sans qu'il faille un générateur de rythme séparé~\cite{Walker2010}. C'est la même cause d'oscillation que pour la boucle \nt{GLUT}--\nt{GABA} de la proposition~\ref{prop:ei} et pour la boucle d'étirement du muscle de la condition~\eqref{eq:delaystable}: une contre-réaction qui arrive en retard.

\section{Le code impulsionnel des hormones}

Les pulsations ne sont pas seulement un effet secondaire; elles peuvent aussi être un code. Les neurones qui commandent la reproduction libèrent leur signal dans le sang par courtes bouffées, environ une fois par heure. Si l'on présente le même signal à la glande destinataire non plus par bouffées mais en continu, elle ne travaille pas à plein régime comme avec les pulsations: elle se tait; si l'on revient à une bouffée par heure, son activité reprend~\cite{Belchetz1978}. Le destinataire ne lit donc pas le niveau, mais la \emph{fréquence} des bouffées.

Pour un ingénieur, il n'y a là aucun mystère. Un récepteur qui se fatigue sous l'action prolongée d'une substance et cesse de répondre est un filtre passe-haut, comme la synapse qui s'épuise~\eqref{eq:depression} du chapitre~\ref{sec:code}. Il étouffe le signal continu et laisse passer les changements. À un tel destinataire, on ne peut rien communiquer autrement que par impulsions, et l'information passe du niveau à la fréquence et à la forme des bouffées. De plus, des fréquences de bouffées différentes agissent différemment sur les diverses cellules d'une même glande, si bien qu'une seule ligne chimique peut transmettre plus d'une grandeur, comme le fil unique \nt{DOPA} du chapitre~\ref{sec:dofaalt} portait une composante rapide et une composante lente.

\section{Le rythme circadien: un générateur lent de modes}

Le rythme le plus lent du corps est le rythme circadien. Il naît d'une contre-réaction retardée à l'intérieur des cellules: des gènes produisent des protéines qui, avec un retard de plusieurs heures, répriment leur propre production~\cite{TakahashiJS2017}. Encore une contre-réaction retardée (la quatrième de ce livre), et encore un rythme, cette fois d'une période d'environ un jour. Le générateur principal est un petit groupe de quelques dizaines de milliers de neurones à la base du cerveau, dont le rythme synchronise les autres cellules du corps.

Chez l'humain, la période propre de ce générateur n'est pas exactement d'un jour, mais de $24.18$ heures en moyenne~\cite{Czeisler1999}. Chaque jour, il est recalé par la lumière que lui transmettent les capteurs de l'œil (chapitre~\ref{sec:sensors}). Pour un électronicien, c'est une \emph{boucle à verrouillage de phase}: un oscillateur de pulsation propre $\omega_0$ se cale sur un signal externe de pulsation $\omega$, et le déphasage $\varphi$ obéit à l'équation
\begin{empheq}[box=\keybox]{equation}
\frac{\dd\varphi}{\dd t} \;=\; (\omega-\omega_0) \;-\; K_\varphi\,\sin\varphi .
\label{eq:pll}
\end{empheq}
Le verrouillage tient si $|\omega-\omega_0|<K_\varphi$. Un générateur de période $24.18$ heures doit se décaler d'environ onze minutes chaque jour; la lumière du matin suffit en général pour cela. Pendant la nuit polaire ou dans une grotte sans lumière, il n'y a plus de verrouillage, et le rythme dérive vers sa période propre.

Le générateur circadien n'est pas l'horloge d'un ordinateur. Il ne compte pas les pas de calcul et ne synchronise pas les impulsions des neurones. Il commute des \emph{modes}: sommeil et veille, niveaux hormonaux, température du corps, disposition à manger. C'est un registre de diffusion lent, de la même famille que les registres des chapitres~\ref{sec:serocomputer}--\ref{sec:acet}, sauf que sa valeur change selon un horaire calé sur le soleil.

\begin{keyblock}
\begin{proposition}[La sortie chimique]
\label{prop:chemoutput}
La sortie lente de la machine est construite avec les mêmes pièces que la sortie rapide. Les organes reçoivent deux fils de signes opposés et de vitesses différentes, l'accélérateur \nt{NORA} et le frein \nt{ACET}; l'organisme entier reçoit des signaux diffusés par le sang, dont l'adrénaline \nt{ADRE}. Les niveaux hormonaux sont tenus par des cascades à contre-réaction dont la précision est limitée par la stabilité; une partie des hormones est codée par la fréquence des bouffées; le rythme circadien est un générateur lent à verrouillage de phase sur la lumière. L'organisation des voies ainsi que les expériences sur les bouffées et sur la période sont mesurées; l'explication des pulsations de la cascade par la rétroaction elle-même est une hypothèse.
\end{proposition}
\end{keyblock}

\section{Bilan}

\begin{table}[t]
\centering
\footnotesize
\setlength{\tabcolsep}{4pt}
\renewcommand{\arraystretch}{1.15}
\begin{tabular}{>{\raggedright}p{3.1cm}>{\raggedright}p{4.8cm}>{\raggedright\arraybackslash}p{4.9cm}}
\toprule
Ce que fait la sortie chimique & Comment & Ce que l'on sait \\
\midrule
ajuste les organes & accélérateur \nt{NORA} et frein \nt{ACET}~\eqref{eq:gasbrake} & mesuré; le frein est plus rapide que l'accélérateur \\
fait passer tout le corps en mode course & adrénaline \nt{ADRE} dans le sang & mesuré \\
tient les niveaux hormonaux & cascade à rétroaction, $K<8$~\eqref{eq:cascadestable} & cascades et rétroaction mesurées; pulsations nées de la boucle elle-même: hypothèse \\
transmet par la fréquence des bouffées & récepteur fatigué comme filtre passe-haut & dépendance aux bouffées mesurée \\
commute les modes sur la journée & générateur à verrouillage de phase sur la lumière~\eqref{eq:pll} & période et recalage par la lumière mesurés \\
\bottomrule
\end{tabular}
\caption{Comment la machine pilote la chimie du corps.}
\label{tab:chemoutput}
\end{table}

La machine a deux sorties (tableau~\ref{tab:chemoutput}). La rapide, ce sont les muscles: des millisecondes, de façon adressée, vers chaque articulation. La lente, c'est la chimie du corps: des secondes, des minutes et des jours, par deux fils vers les organes et par le sang vers tous à la fois. Les deux sorties sont faites des mêmes pièces (deux fils de signes opposés, des stimulateurs, la rétroaction et son retard), celles-là mêmes dont la machine est construite à l'intérieur.

\section{Exercices}

\begin{exercise}[\lvlA]
\label{ex:16:1}
\emph{Le sang comme filtre.} Une hormone quitte le sang avec une demi-vie $t_{1/2}=10$~min et arrive par courtes bouffées toutes les $T=60$~min. De quel facteur le maximum de concentration dépasse-t-il le minimum, et de quel facteur le fondamental des bouffées est-il atténué? Pour quel $t_{1/2}$ le maximum ne dépasse-t-il le minimum que d'un facteur deux?
\end{exercise}

\begin{exercise}[\lvlA]
\label{ex:16:2}
\emph{Verrouillage circadien.} La période propre du générateur~\eqref{eq:pll} vaut $24.18$~h; la lumière décale la phase d'au plus une heure par jour: $K_\varphi=2\pi/24$~rad/j. Trouvez le déphasage stationnaire $\varphi^*$ en minutes et le temps de relaxation. Au bout de combien de jours un désaccord créé par un saut du signal externe de $\pm6$~h tombe-t-il sous une heure?
\end{exercise}

\begin{exercise}[\lvlB]
\label{ex:16:3}
\emph{Précision contre stabilité.} On allonge la cascade linéarisée~\eqref{eq:cascade} à $n$ étages identiques. Trouvez le gain critique $K_c(n)$ et la plus petite erreur accessible $1/(1+K_c)$ pour $n=3$ et $n=6$. Vers quoi tend cette erreur quand $n\to\infty$?
\end{exercise}

\begin{exercise}[\lvlB]
\label{ex:16:4}
\emph{Accélérateur et frein.} Dans~\eqref{eq:gasbrake}, les contributions de l'accélérateur et du frein sont les sorties de circuits RC avec $\tau_S=2$~s et $\tau_P=0.4$~s, les commandes étant limitées à $80$ et $60$ battements par minute. Rythme $f_0=100$; au repos, frein $35$, accélérateur $0$, $f=65$. En combien de temps peut-on atteindre $f=120$? Et sans relâcher le frein, avec l'accélérateur seul?
\end{exercise}

\begin{exercise}[\lvlB]
\label{ex:16:5}
\emph{Simulation: cascade avec retard.} Dans la rétroaction de la fonction \texttt{cascade} de \texttt{models/\allowbreak chem\_models.py} (copiez-la, n'exécutez pas le fichier), ajoutez un retard $d$: le premier étage voit $x_3(t-d)$. Trouvez $K_c$ et la période à la limite pour $d/\tau=0$, $0.5$, $1$, $2$, et vérifiez par simulation à $0.9K_c$ et $1.1K_c$.
\end{exercise}

\begin{exercise}[\lvlB]
\label{ex:16:6}
\emph{Un destinataire qui lit les bouffées.} Le récepteur se fatigue: $y=[c-a]_+$, $\tau_a\,\dd a/\dd t=c-a$, $\tau_a=30$~min. L'hormone arrive par bouffées de $6$~min avec une période $T$, à dose moyenne $\bar c$ constante. Trouvez la réponse moyenne pour $T=15$, $60$, $120$~min et $T\to\infty$. Que vaut-elle à concentration constante $\bar c$?
\end{exercise}

\begin{exercise}[\lvlC]
\label{ex:16:7}
\emph{Générateur ou rétroaction?} Proposez une expérience qui distingue les pulsations engendrées par la rétroaction retardée de la cascade~\eqref{eq:cascade} de celles d'un stimulateur séparé. Peut-on départager les hypothèses par le décalage de la période si l'on ne peut modifier $K$ que d'un facteur une fois et demie et que la période est mesurée à $5\%$ près?
\end{exercise}

\chapter{Déboguer la machine}
\chaptermark{Déboguer la machine}
\label{sec:debugging}

\section{Comment déboguer une machine sans schéma}

Un ingénieur à qui l'on confie une carte qui fonctionne, mais sans son schéma, la débogue avec quatre techniques. Il pose une \emph{sonde} et regarde ce qui passe dans un fil. Il \emph{injecte un signal de test} et regarde ce qui a changé. Il \emph{déconnecte} une partie du circuit et regarde ce qui a cessé de fonctionner. Et il \emph{lit les registres de contrôle} pour savoir dans quel mode tourne la carte. Tout ce livre est une tentative de lire le schéma du cerveau; dans ce chapitre, nous verrons lesquelles de ces quatre techniques les ingénieurs savent déjà employer, et ce que les chapitres précédents leur disent de ce qu'il faut attendre.

L'exemple le plus visible aujourd'hui d'un tel débogage, ce sont les interfaces cerveau-ordinateur: des appareils qui lisent les impulsions des neurones et les transforment en commandes pour des machines externes, ou, à l'inverse, envoient au cerveau des signaux venus de l'extérieur. C'est à elles qu'est consacrée la plus grande partie du chapitre, et d'abord à ce que fait la société Neuralink.

\section{La sonde: ce qu'entend une électrode}

La sonde du cerveau est une fine électrode insérée dans le tissu. Elle entend les impulsions des neurones situés dans un rayon d'une centaine de micromètres, en général d'un seul à quelques-uns. Sur l'électrode, une impulsion est une pointe de quelques dizaines à quelques centaines de microvolts qui dure environ une milliseconde, et sa forme permet de distinguer des neurones voisins les uns des autres. Ce sont les grands neurones \nt{GLUT} que l'on entend le mieux: ils sont majoritaires dans le cortex (tableau~\ref{tab:types}), et leurs courants sont plus forts.

La difficulté principale saute aux yeux. Une bonne sonde compte aujourd'hui des centaines ou des milliers d'électrodes, alors que le cerveau contient $8.6\cdot10^{10}$ neurones. Nous écoutons environ un cent-millionième de la machine. Pourquoi un échantillon aussi dérisoire permet-il de comprendre quoi que ce soit? Le chapitre~\ref{sec:code} a donné la réponse. Les neurones voisins sont bruités de façon corrélée, et la moyenne sur un groupe bute sur une limite (proposition~\ref{prop:pooling}): une centaine de neurones porte presque autant qu'un millier. De plus, la tâche que résout le cortex moteur est de faible dimension: un bras a peu d'articulations (chapitre~\ref{sec:muscles}), et l'activité des centaines de neurones qui le commandent tient dans un espace d'une dizaine ou deux de variables communes seulement~\cite{Gallego2017}. Une sonde qui écoute une centaine de neurones entend presque toutes ces variables. Si le cerveau stockait un message indépendant dans chaque neurone, une telle écoute serait inutile.

\section{Le flux de données: compresser sur la sonde}

La sonde produit un flux énorme. Pour voir la forme d'une impulsion, il faut numériser chaque électrode environ $20$ mille fois par seconde avec une précision d'une dizaine de bits. Mille électrodes donnent
\begin{empheq}[box=\keybox]{equation}
\begin{aligned}
R_{\text{brut}} \;&=\; N\,f_s\,b \;\approx\; 10^3\cdot2\cdot10^4\cdot10 \;\approx\; 2\cdot10^8\ \text{bit/s},\\
R_{\text{imp}} \;&\approx\; N\,r\,\log_2\frac{e}{r\,\Delta t} \;\approx\; 8\cdot10^4\ \text{bit/s}.
\end{aligned}
\label{eq:probedata}
\end{empheq}
La seconde estimation est la capacité du flux d'impulsions~\eqref{eq:spikeentropy} pour $r=10$ impulsions par seconde et une précision $\Delta t=1\ms$: tout ce qu'il contient réellement tient dans un volume deux mille cinq cents fois plus petit. Pour un appareil implanté, la différence est décisive: on ne peut pas transmettre le flux brut par radio à travers la peau avec l'alimentation que l'on peut se permettre à l'intérieur de la tête sans chauffer le tissu. Un appareil implanté doit donc détecter les impulsions directement sur la puce, à côté des électrodes, et ne transmettre vers l'extérieur que des événements: le numéro du canal et l'instant de l'impulsion. La première description du système Neuralink n'en était pas encore là: les puces amplifient et numérisent les signaux, tout le flux brut part par câble, et c'est un programme externe qui détecte les impulsions; la compression sur puce et la liaison sans fil y sont présentées comme un projet~\cite{Musk2019}.

C'est une technique familière. L'œil compresse le signal d'un facteur cent avant de l'envoyer par un long câble (chapitre~\ref{sec:sensors}); une caméra événementielle ne transmet pas des images, mais des changements. Pour tenir dans son fil, la sonde est obligée de faire ce que fait le cerveau lui-même: parler par impulsions, et seulement des nouveautés.

\section{Ce que fait Neuralink}

L'appareil de Neuralink est une sonde conçue sous ces contraintes~\cite{Musk2019}. Au lieu d'une grille rigide d'aiguilles, de nombreux fils souples plus fins qu'un cheveu, portant chacun des électrodes sur leur longueur: dans la première description du système, jusqu'à $3072$ électrodes sur $96$ fils, $32$ par fil; dans l'appareil implanté chez l'humain, selon la société, environ mille canaux sur $64$ fils. Les fils souples abîment moins le tissu et supportent mieux le fait que le cerveau bouge sans cesse légèrement dans le crâne. Un robot insère les fils en évitant les vaisseaux. Dans la première description, comme dit plus haut, le flux brut~\eqref{eq:probedata} sortait par câble. L'appareil destiné à l'humain, selon la société, est conçu autrement: le boîtier est entièrement recouvert par la peau, les impulsions sont détectées à l'intérieur, les données partent par radio et l'alimentation arrive sans fil.

La tâche du premier appareil est la lecture. On place les fils dans la partie du cortex qui planifie les mouvements du bras, et un décodeur transforme les impulsions en mouvement d'un curseur à l'écran. Une personne aux bras paralysés commande un ordinateur en imaginant le mouvement. L'idée elle-même n'est pas neuve: des interfaces à grilles rigides d'une centaine d'électrodes permettaient depuis longtemps de piloter un curseur~\cite{Hochberg2006}, d'écrire un texte en imaginant les gestes de l'écriture à la main~\cite{Willett2021}, et même de convertir des tentatives de parole en texte à l'écran à une soixantaine de mots par minute~\cite{Willett2023}. Ce que Neuralink apporte de neuf, c'est l'ingénierie: le nombre de canaux, l'absence de fil, le boîtier fermé et l'insertion robotisée, autrement dit la tentative de faire une sonde que l'on peut porter des années hors du laboratoire. À notre connaissance, les résultats des essais chez l'humain ont surtout été publiés par la société elle-même, et non dans des articles à comité de lecture; la société a notamment indiqué qu'une partie des fils s'était déplacée après l'insertion et que le nombre de canaux actifs avait diminué. Pour le débogage, c'est un désagrément ordinaire: une sonde qui bouge par rapport à la carte écoute déjà d'autres fils.

Le second axe de la société est l'écriture: un appareil pour la vision, qui envoie des signaux dans le cortex visuel d'une personne ayant perdu la vue. Nous en parlons plus bas, dans la section sur l'écriture.

\section{La lecture: le décodeur comme destinataire}

Le décodeur d'une interface est un destinataire au sens de la proposition~\ref{prop:reader}: c'est lui qui décide quelle partie du message lire. Pratiquement toutes les interfaces lisent les \emph{fréquences de groupes}: elles comptent les impulsions de chaque canal dans des fenêtres de quelques dizaines de millisecondes et les additionnent avec des poids choisis pour obtenir le mouvement voulu. C'est le code de fréquence de groupe de la section~\ref{sec:rate}, et ce choix n'est pas un hasard: avec quelques impulsions par fenêtre, la fréquence d'un seul neurone n'est pas définie, alors que la somme sur une centaine fonctionne déjà.

Combien d'information parvient-on à extraire? L'écriture \og par la main mentale\fg{} atteignait environ quatre-vingt-dix caractères par minute~\cite{Willett2021}, soit un caractère et demi par seconde: même à cinq bits par caractère, cela fait moins de huit bits par seconde. Le pilotage d'un curseur donne, selon Neuralink, de l'ordre de dix bits par seconde. Or la borne supérieure~\eqref{eq:spikeentropy} est de quelques dizaines de bits par seconde pour \emph{un seul} neurone, et de quelques milliers pour une centaine. L'interface extrait des neurones écoutés une petite fraction de ce qu'ils pourraient porter. Le goulot d'étranglement n'est pas le fil, mais le nombre de variables indépendantes que porte le groupe (proposition~\ref{prop:pooling}) et la façon dont le décodeur comprend un code qui, nous l'avons vu au chapitre~\ref{sec:code}, n'est pas entièrement déchiffré.

Il y a une seconde particularité, familière depuis le chapitre~\ref{sec:motorlearning}. Le décodeur et le cerveau apprennent l'un de l'autre. Le cerveau voit où est allé le curseur, reçoit une erreur et ajuste ses commandes: c'est le même apprentissage moteur par le fil d'erreur. Et les ingénieurs recalculent de temps en temps les poids du décodeur, parce que les neurones sous les fils changent et que les connexions du réseau réapprennent. On obtient deux élèves qui s'adaptent l'un à l'autre; pour qu'une telle paire ne s'emballe pas, il est commode de séparer leurs vitesses d'apprentissage: que l'un apprenne vite et l'autre lentement.

\begin{keyblock}
\begin{proposition}[Pourquoi une sonde sur mille neurones fonctionne]
\label{prop:probe}
Une interface qui écoute une part infime des neurones peut piloter un mouvement parce que l'activité du groupe est de faible dimension et que ses bruits sont corrélés: quelques centaines de neurones portent presque toutes les variables communes. Pour la même raison, l'interface n'extrait que quelques bits par seconde, autant qu'il y a de variables indépendantes dans la tâche, et non autant que peut en contenir le flux d'impulsions. Le fonctionnement des interfaces est mesuré; de combien le décodeur s'améliorera quand le code sera compris reste une question ouverte.
\end{proposition}
\end{keyblock}

\section{L'écriture: plus difficile que la lecture}

Injecter un signal dans la machine est plus difficile que le lire. Une électrode parcourue par un courant n'excite pas un seul neurone, mais tous les neurones et tous les fils qui l'entourent, sans distinction de type ni de signe. Elle ne peut pas écrire un code dans lequel compte \emph{quel} neurone a déchargé et \emph{quand} (chapitre~\ref{sec:code}): elle ne peut fixer que grossièrement \emph{où} et \emph{avec quelle force}.

Pour la vision, cela suffit en partie. Un courant envoyé par une électrode dans le cortex visuel est perçu comme un point lumineux à un endroit donné du champ visuel; quand on allumait des points simultanément, jusqu'à seize à la fois, une personne ayant perdu la vue reconnaissait certaines lettres et distinguait les contours des objets~\cite{Fernandez2021}. C'est un code de position, et on peut l'écrire. Mais rappelez-vous le chapitre~\ref{sec:sensors}: l'œil n'envoie pas au cerveau des pixels, mais une description compressée (bords, changements, éclaircissements et assombrissements, sur des fils séparés). Écrire des \og points de lumière\fg{} dans le cortex, c'est écrire des pixels dans une machine qui attend à l'entrée un tout autre code. C'est pourquoi la vision artificielle reste plus grossière qu'on ne l'attendrait d'après le nombre d'électrodes. Il existe aussi des signaux de test qui n'ont pas besoin d'électrode: les illusions visuelles présentent par l'œil une entrée inhabituelle et font sortir la machine de son régime normal, et chaque erreur désigne son propre mécanisme (section~\ref{sec:illusions}).

\section{Ce que la sonde ne voit pas}

L'électrode entend le bus d'adresses: les impulsions des neurones \nt{GLUT} et, moins bien, des petits \nt{GABA}. Mais aux chapitres~\ref{sec:serocomputer}--\ref{sec:acet}, nous avons vu que le mode de fonctionnement de toute la machine est fixé par des registres de diffusion: les concentrations de \nt{SERO}, \nt{DOPA}, \nt{NORA}, \nt{ACET}. L'électrode ne les entend pas: les neurones sources sont en nombre infime, et leur signal n'est pas une impulsion près de la sonde, mais une concentration dans un volume. Un débogueur qui voit le bus de données mais pas les registres de contrôle sera toujours surpris: une même entrée donnera des réponses différentes, parce que quelque part $g$, $a$ ou $\tau_\gamma$ aura changé. On peut mesurer les concentrations par des capteurs chimiques directement dans le tissu~\cite{Bunin1998}, mais les interfaces n'en utilisent pas encore. La seconde sortie de la machine, la lente, échappe elle aussi entièrement à la sonde: les hormones dans le sang (chapitre~\ref{sec:chemoutput}) se mesurent dans des prélèvements sanguins, pas avec une électrode.

La sonde ne voit pas non plus les poids, or c'est en eux que réside presque toute la mémoire et tout ce qui a été appris. On ne peut que les deviner d'après la façon dont les réponses changent. Et elle ne voit pas la douleur telle que la ressent une personne: au chapitre~\ref{sec:pain}, nous avons vu que la force du signal d'alarme est fixée par la machine elle-même (par le portillon de la moelle épinière et par les régulateurs d'en haut), si bien qu'on ne peut pas lire sur un capteur à quel point on a mal.

\section{Bilan: le débogage et les questions ouvertes}

\begin{table}[t]
\centering
\footnotesize
\setlength{\tabcolsep}{4pt}
\renewcommand{\arraystretch}{1.15}
\begin{tabular}{>{\raggedright}p{2.2cm}>{\raggedright}p{3.6cm}>{\raggedright}p{3.4cm}>{\raggedright\arraybackslash}p{4.0cm}}
\toprule
Technique de débogage & Ce que fait l'interface & Ce qu'explique le livre & Ce que l'on sait \\
\midrule
sonde & écoute des centaines à des milliers de neurones & faible dimension et bruit corrélé (chap.~\ref{sec:code}) & mesuré \\
compression sur la sonde & transmet des événements au lieu du signal brut & capacité du flux d'impulsions~\eqref{eq:spikeentropy} & résolu par l'ingénierie \\
lecture & fréquences de groupes $\to$ mouvement, texte, parole & le décodeur est un destinataire; code de fréquence de groupe & fonctionne; quelques bits par seconde \\
apprentissage conjoint & le cerveau et le décodeur s'adaptent l'un à l'autre & fil d'erreur (chap.~\ref{sec:motorlearning}) & observé; les règles d'adaptation sont un sujet de recherche \\
écriture & le courant allume des points dans le champ visuel & un code de position s'écrit, un code temporel non (chap.~\ref{sec:sensors}) & points mesurés; vision utile: pas encore \\
lecture des registres & presque jamais faite & canaux de diffusion (chap.~\ref{sec:serocomputer}--\ref{sec:acet}) & capteurs chimiques disponibles en expérience \\
\bottomrule
\end{tabular}
\caption{Déboguer la machine: ce que savent faire les interfaces cerveau-ordinateur et ce qu'en disent les chapitres du livre.}
\label{tab:debugging}
\end{table}

Le tableau~\ref{tab:debugging} résume le chapitre. Les interfaces cerveau-ordinateur savent déjà poser une sonde sur le bus d'adresses et y lire quelques bits par seconde, et le fait même que cela fonctionne s'explique par la faible dimension et le bruit corrélé du chapitre~\ref{sec:code}. Elles ne savent écrire dans la machine que grossièrement, parce que son code d'entrée est compressé et adressé à des fils séparés. Quant aux registres de contrôle et aux poids, ce qui rend la machine capable d'apprendre et de se régler, la sonde ne les voit encore presque pas.

Chaque question ouverte du chapitre~\ref{sec:synthesis} est aussi une tâche pour le débogueur. Quel code lire, on le saura quand les sondes seront plus denses et que les décodeurs sauront exploiter l'instant des impulsions, et pas seulement les fréquences. Comment apprend un réseau à plusieurs couches, on peut le vérifier en posant des sondes dans plusieurs couches à la fois et en regardant comment leurs réponses changent pendant l'apprentissage. Ce que transmettent les canaux de diffusion deviendra visible quand des sondes chimiques s'ajouteront aux sondes électriques. Ce livre est une tentative de lire le schéma de la machine d'après son comportement et d'après les lois que connaît tout ingénieur. Les débogueurs en construction aujourd'hui permettront de vérifier ce schéma à la sonde.

\section{Exercices}

\begin{exercise}[\lvlA]
\label{ex:17:1}
\emph{Budget de la liaison radio.} Transmettre à travers la peau coûte $1$~nJ par bit, et l'émetteur ne peut pas dépenser plus de $10$~mW. Quelle puissance faut-il pour le flux brut~\eqref{eq:probedata} avec $N=1000$ électrodes, et pour le flux d'impulsions? Quelle énergie par bit exigerait le flux brut?
\end{exercise}

\begin{exercise}[\lvlA]
\label{ex:17:2}
\emph{Combien de bits dans cent neurones.} Le décodeur additionne les impulsions de $N$ neurones sur des fenêtres de $50\ms$, avec $r=20$~imp/s, une corrélation de bruit par paires $c=0.1$, et un signal modulant la fréquence du groupe de $30\%$ (valeur efficace). Estimez le débit $\tfrac12\log_2(1+\text{SNR})$ par fenêtre pour $N=10$, $100$, $1000$ et $N\to\infty$. Et pour $c=0$, $N=1000$?
\end{exercise}

\begin{exercise}[\lvlB]
\label{ex:17:3}
\emph{Débit d'une interface-clavier.} Une fois et demie par seconde, l'interface choisit l'un de $N=32$ caractères: le caractère voulu avec une probabilité $P$, les erreurs étant équiprobables. Combien de bits par seconde transmet-elle pour $P=0.99$, $0.94$, $0.8$? Combien de caractères par seconde faut-il à $P=0.94$ pour $10$~bit/s?
\end{exercise}

\begin{exercise}[\lvlB]
\label{ex:17:4}
\emph{Simulation: décodeur de groupe.} $N$ neurones avec $\lambda_i=[b+m\,\mathbf v\cdot\mathbf p_i]_+$, $b=20$, $m=15$~imp/s, fenêtres de $50\ms$, $\mathbf v$ d'écart-type $0.5$ par composante; tous voient un bruit commun, $\mathbf v+\boldsymbol\xi$, $\sigma_\xi=0.2$. Simulez un décodeur non biaisé par vecteur de population. À quel $N$ les deux bruits s'égalent-ils, et combien de bits par composante et fenêtre pour $N\to\infty$?
\end{exercise}

\begin{exercise}[\lvlB]
\label{ex:17:5}
\emph{Deux élèves.} Le cerveau émet la commande $m=b\,v^*$, le décodeur produit $v=d\,m$, $\langle v^{*2}\rangle=1$. Après chaque essai, tous deux font un pas de gradient sur $\tfrac12(v-v^*)^2$ avec le taux $\eta$. Simulez à partir de $b=1.2$, $d=1$ pour $\eta=0.8$, $1.1$, $1.5$, $2$. Chacun seul est stable pour $\eta<2$; pour quel $\eta$ la paire est-elle stable?
\end{exercise}

\begin{exercise}[\lvlB]
\label{ex:17:6}
\emph{Conception: la chaîne de la sonde.} $1024$ canaux à $20$~k échantillons par seconde, neurones à $10$~imp/s, radio à $1$~nJ/bit. Quel seuil sur le bruit blanc donne au plus $0.1$~Hz de fausses impulsions par canal? On transmet le compte d'impulsions par tranche de $20\ms$ sur $2$~bits: quelle puissance, et combien de fois moins que pour des échantillons bruts sur $10$~bits?
\end{exercise}

\begin{exercise}[\lvlC]
\label{ex:17:7}
\emph{La limite de débit d'une interface.} Estimez $R\le DW\log_2(1+\text{SNR})$ pour $D=10$ variables de bande $W=5$~Hz avec $\text{SNR}\approx0.9$ (bruit semblable au signal) et $90$ (bruit indépendant de $1000$ neurones). On mesure environ $10$~bit/s. Quelle expérience distinguerait les deux explications de cet écart: un bruit semblable au signal, ou la méconnaissance du code?
\end{exercise}

\chapter{Les types de neurones selon leur fonction électrique}
\chaptermark{Les neurones comme composants}
\label{sec:eetypes}

Au chapitre~\ref{sec:neurontypes}, nous avons trié les neurones selon le neurotransmetteur que libère leur sortie. Un électronicien les trierait autrement: selon ce que le neurone fait de son signal, c'est-à-dire selon le composant qu'il réalise. Nous connaissons le composant principal du livre depuis le chapitre~\ref{sec:glutcomputer}: le neurone intégrateur à fuite~\eqref{eq:glutneuron}, un circuit RC suivi d'un comparateur. Mais le cerveau en contient d'autres. Le tableau~\ref{tab:eetypes} les réunit en quatre groupes, et presque tous ont déjà été rencontrés dans le livre.

\begin{center}
\footnotesize
\setlength{\tabcolsep}{3pt}
\renewcommand{\arraystretch}{1.15}
\begin{longtable}{>{\raggedright}p{3.1cm}>{\raggedright}p{3.3cm}>{\raggedright}p{4.7cm}>{\raggedright\arraybackslash}p{2.4cm}}
\caption{Les neurones comme composants: nom, équivalent en électronique, ce que fait le composant et où il apparaît dans le livre.}\label{tab:eetypes}\\
\toprule
Neurone & Composant & Ce qu'il fait & Où dans le livre \\
\midrule
\endfirsthead
\toprule
Neurone & Composant & Ce qu'il fait & Où dans le livre \\
\midrule
\endhead
\bottomrule
\endfoot
\multicolumn{4}{l}{\emph{Filtres et intégrateurs}} \\
neurone intégrateur à fuite & intégrateur RC avec comparateur & additionne les entrées sur une fenêtre $\tau$ et décharge au seuil & chap.~\ref{sec:glutcomputer}, \eqref{eq:glutneuron} \\
détecteur de coïncidences & ET avec fenêtre temporelle & ne décharge que si les entrées arrivent presque ensemble; $\tau$ très petit & chap.~\ref{sec:glutcomputer}, \ref{sec:code}, \eqref{eq:window} \\
neurone phasique & filtre passe-haut, détecteur de front & répond au changement de l'entrée, puis se tait & chap.~\ref{sec:sensors} \\
neurone à adaptation & filtre passe-haut avec contrôle automatique de gain & la fréquence baisse sous une entrée constante & chap.~\ref{sec:sensors}, \eqref{eq:nakarushton} \\
résonateur & circuit oscillant, filtre passe-bande & répond le mieux à une entrée de sa propre fréquence & section~\ref{sec:resonator} \\
intégrateur sans fuite & intégrateur à amplificateur opérationnel & transforme une vitesse en position et la maintient & section~\ref{sec:perfectint} \\
\midrule
\multicolumn{4}{l}{\emph{Convertisseurs}} \\
neurone tonique & convertisseur tension-fréquence & la fréquence suit l'entrée linéairement et se fatigue à peine & chap.~\ref{sec:code} \\
neurone gradué (sans impulsions) & amplificateur analogique, tampon & transmet une tension continue sur une courte distance & chap.~\ref{sec:sensors} \\
neurone redresseur & redresseur simple alternance & la fréquence n'est jamais négative, $[u]_+$ & chap.~\ref{sec:glutcomputer}, \ref{sec:gaba} \\
paire \og ON-OFF\fg{} & paire de redresseurs, code à deux fils & l'un répond à \og plus\fg{}, l'autre à \og moins\fg{} & chap.~\ref{sec:glutcomputer}, \eqref{eq:dualrail} \\
neurone à inhibition shuntante & diviseur analogique, réglage du gain & divise l'entrée au lieu de la soustraire & chap.~\ref{sec:gaba}, \eqref{eq:shunt} \\
\midrule
\multicolumn{4}{l}{\emph{Générateurs et temps}} \\
stimulateur & oscillateur à relaxation, multivibrateur & émet de lui-même des impulsions à fréquence constante & chap.~\ref{sec:serocomputer}, \ref{sec:dofa} \\
stimulateur à entrée & oscillateur commandé en tension & rythme propre que l'entrée accélère ou ralentit & chap.~\ref{sec:chemoutput} \\
neurone à bouffées & générateur de bouffées déclenché & émet des bouffées d'impulsions rapprochées, les \og traits\fg{} & chap.~\ref{sec:code}, \eqref{eq:burst} \\
neurone calé sur la phase & détecteur de phase, boucle à verrouillage de phase & décharge à une phase donnée de l'onde d'entrée & chap.~\ref{sec:sensors}, \ref{sec:chemoutput} \\
axone à retard & ligne à retard & retarde le signal d'un temps donné & chap.~\ref{sec:glutcomputer}, fig.~\ref{fig:itd} \\
\midrule
\multicolumn{4}{l}{\emph{Mémoire et commutation}} \\
neurone bistable & trigger de Schmitt, verrou & une brève entrée le fait passer dans un état durable & section~\ref{sec:bistable} \\
neurone à branches dendritiques & multiplexeur, petit réseau multicouche & une branche ne laisse passer l'entrée que si une entrée d'autorisation est présente & chap.~\ref{sec:synthesis} \\
\end{longtable}
\end{center}

Une réserve compte plus que tout le tableau: il ne s'agit pas de cellules différentes, mais de \emph{modes} différents. Un même neurone peut être un intégrateur quand l'entrée est lente et un détecteur de coïncidences quand l'inhibition a raccourci son $\tau$ (chapitre~\ref{sec:code}); un stimulateur pendant la veille et un récepteur silencieux pendant le sommeil. Le composant obtenu dépend des canaux ioniques de la membrane et des canaux de diffusion qui les ajustent.

\section{Quatre composants que le livre n'a pas encore rencontrés}

\subsection{Le résonateur}
\label{sec:resonator}

Certains neurones possèdent un courant lent qui s'oppose à toute variation de tension: si la tension monte, le courant la tire vers le bas, si elle baisse, il la tire vers le haut, mais avec retard. Pour un électronicien, un tel courant se comporte comme une inductance, et avec la capacité de la membrane il forme un circuit oscillant. Le neurone répond le plus fortement à une entrée d'une fréquence donnée, en général de quelques hertz à une vingtaine, et plus faiblement aux entrées plus lentes ou plus rapides~\cite{HutcheonYarom2000}. C'est un filtre passe-bande dans une seule cellule: dans une entrée bruitée, il sélectionne le rythme de sa fréquence, par exemple le rythme local du chapitre~\ref{sec:gaba}.

\subsection{L'intégrateur sans fuite}
\label{sec:perfectint}

Un intégrateur à fuite ne se souvient de son entrée que pendant un temps $\tau$, quelques dizaines de millisecondes. Mais pour rester immobile, l'œil doit se souvenir de sa position pendant des secondes: l'ordre de \og tourner\fg{} arrive comme une vitesse brève, alors qu'il faut ensuite tenir l'œil dans la nouvelle position. Le réseau résout ce problème par la même rétroaction que la mémoire du chapitre~\ref{sec:glutcomputer}. Si un groupe de neurones de constante de temps $\tau$ s'excite lui-même avec un poids $w$, alors $\tau\,\dd r/\dd t=-r+w\,r+I$, et la constante de temps du groupe vaut
\begin{empheq}[box=\keybox]{equation}
\tau_{\text{eff}} \;=\; \frac{\tau}{1-w} .
\label{eq:perfectint}
\end{empheq}
Avec $\tau=0.1\,\text{s}$ et $w=0.995$, on obtient $20\,\text{s}$: un intégrateur presque idéal fait de neurones dont chacun, seul, oublie en un dixième de seconde~\cite{Seung1996}. Le prix à payer est un réglage précis: pour $w$ un peu supérieur à un, l'intégrateur part en saturation; un peu inférieur, il oublie vite. Un tel intégrateur a donc besoin d'un réglage permanent par apprentissage, comme celui que décrit le chapitre~\ref{sec:motorlearning}.

\subsection{Le neurone bistable}
\label{sec:bistable}

Aux chapitres~\ref{sec:glutcomputer} et~\ref{sec:gaba}, la bascule était assemblée à partir d'un groupe de neurones. Mais il existe aussi des bascules dans une seule cellule. Certains neurones possèdent un courant qui, une fois enclenché, entretient lui-même l'excitation: une brève entrée fait passer le neurone dans une activité durable, un \emph{plateau}, et l'inhibition le ramène au repos. C'est un trigger de Schmitt: le neurone a deux états stables et un hystérésis entre eux. C'est le cas, par exemple, des neurones \nt{ACET} qui commandent les muscles, et cette propriété est activée par un canal de diffusion: le plateau apparaît quand la sérotonine atteint le neurone~\cite{Hounsgaard1988}. La même cellule est un verrou avec \nt{SERO} et un intégrateur ordinaire sans lui.

\subsection{Intégrateurs et résonateurs}

La théorie répartit les cellules excitables en deux classes~\cite{Izhikevich2007}. Les \emph{intégrateurs} ressemblent à un oscillateur commandé en tension qui part de zéro: juste au-dessus du seuil, le neurone décharge aussi lentement qu'on veut, et sa fréquence croît continûment avec l'entrée. Un tel neurone additionne tout ce qui arrive et transmet bien l'amplitude de l'entrée par sa fréquence. Les \emph{résonateurs} ressemblent à un oscillateur à fréquence minimale: ils ne savent pas décharger lentement, émettent d'emblée des impulsions à fréquence finie dès que l'entrée atteint le seuil, et préfèrent une entrée de leur propre fréquence. Les premiers sont les composants naturels du code de fréquence du chapitre~\ref{sec:code}, les seconds du code temporel.

\begin{keyblock}
\begin{proposition}[Un neurone, de nombreux composants]
\label{prop:eetypes}
Selon leur fonction électrique, les neurones fonctionnent comme des intégrateurs à fuite ou sans fuite, des détecteurs de coïncidences, des filtres passe-haut et passe-bande, des convertisseurs tension-fréquence, des redresseurs, des diviseurs, des oscillateurs, des lignes à retard et des bascules. Le composant que donne une cellule dépend de ses canaux ioniques et des canaux de diffusion qui les ajustent. Les modes eux-mêmes sont mesurés; la mesure dans laquelle le cerveau exploite cette reconfiguration pour calculer est surtout une hypothèse.
\end{proposition}
\end{keyblock}

Pour un ingénieur, cela ressemble à un réseau logique programmable: le matériel est le même, et c'est la configuration qui fixe le composant. Sauf qu'ici on ne change pas la configuration au moment de la programmation, mais en marche, par les canaux de diffusion lents dont il a été question aux chapitres~\ref{sec:serocomputer}--\ref{sec:acet}.

\section{Exercices}

\begin{exercise}[\lvlA]
\label{ex:18:1}
\emph{Tolérance de l'intégrateur sans fuite.} Des neurones avec $\tau=0.1$~s tiennent la position de l'œil selon~\eqref{eq:perfectint} pendant $T=1$~s avec une erreur d'au plus $\pm5\%$. (a)~Trouvez la plage admissible de $w$. (b)~$w$ est la somme des poids de $K$ synapses, chacun avec une erreur indépendante de $10\%$. Pour quel $K$ la dispersion de la somme tient-elle dans la tolérance?
\end{exercise}

\begin{exercise}[\lvlB]
\label{ex:18:2}
\emph{Le résonateur comme circuit.} Décrivons le courant lent de la section~\ref{sec:resonator} par une variable $W$: $C\,\dd V/\dd t=-g_LV-g_wW+I$, $\tau_w\,\dd W/\dd t=V-W$. Montrez que c'est un circuit $RC$ avec une branche parallèle $R_w=1/g_w$, $L_w=\tau_w/g_w$; pour $C/g_L=10\ms$, $\tau_w=100\ms$, $\alpha=g_w/g_L=4$, trouvez la fréquence de résonance et le rapport $|Z|/|Z(0)|$ en ce point.
\end{exercise}

\begin{exercise}[\lvlB]
\label{ex:18:3}
\emph{Comment un intégrateur se met à décharger.} (a)~Neurone $\tau\,\dd V/\dd t=-V+\mu$, $\tau=10\ms$, de seuil $\theta$ avec remise à zéro. Pour quel $\mu/\theta-1$ la fréquence vaut-elle $1$~Hz? (b)~Quelle fréquence donne le modèle $\tau\,\dd v/\dd t=v^2+\mu$ (impulsion: départ de $v$ vers $+\infty$, remise à $-\infty$) pour $\mu=0.01$?
\end{exercise}

\begin{exercise}[\lvlB]
\label{ex:18:4}
\emph{Simulation: intégrateur et résonateur.} Le modèle de l'exercice~\ref{ex:18:2} avec $g_L=1$, $\tau_m=10\ms$, $\tau_w=100\ms$, un seuil $1$ et une remise $V\to0$ ($W$ inchangé) reçoit $\mu=1.5\sin2\pi ft$, $f=1$--$40$~Hz. Dans quelle bande de fréquences le neurone émet-il des impulsions pour $\alpha=0$ et pour $\alpha=4$? Comparez avec la condition $1.5\,|Z(f)|>1$.
\end{exercise}

\begin{exercise}[\lvlB]
\label{ex:18:5}
\emph{Un verrou fait d'une seule cellule.} Neurone à courant de plateau: $\tau\,\dd r/\dd t=-r+F(wr+I)$, $F(x)=\min(1,\max(0,x))$, $\tau=10\ms$, $w=1.5$, fond $I=-0.2$. (a)~Quelle durée minimale d'une impulsion $I=0.3$ fait passer le neurone de $r=0$ au plateau? (b)~Quelle amplitude minimale $B$ d'une longue impulsion $I=-0.2-B$ le ramène au repos?
\end{exercise}

\begin{exercise}[\lvlB]
\label{ex:18:6}
\emph{Autoréglage de l'intégrateur.} Pendant la fixation du regard, l'entrée est $I=0$, un capteur de glissement donne la vitesse de dérive $e=\dd r/\dd t$, et $\dd w/\dd t=-\eta\,e\,r$. Pour $w=0.95$, $\tau=0.1$~s, $\langle r^2\rangle=1$, trouvez $\eta$ tel qu'après $60$~s de fixations $|1-w|$ entre dans la tolérance de l'exercice~\ref{ex:18:1}. Que se passe-t-il si l'on apprend aussi pendant les rotations de l'œil?
\end{exercise}

\begin{exercise}[\lvlC]
\label{ex:18:7}
\emph{Pourquoi reconfigurer le composant?} Un canal de diffusion commute l'$\alpha$ de l'exercice~\ref{ex:18:2} entre $0$ et $4$. De quel facteur le résonateur l'emporte-t-il sur l'intégrateur en rapport signal/bruit pour une faible sinusoïde dans un bruit blanc de courant à $11$~Hz et à $1$~Hz? Imaginez une tâche où la commutation de mode elle-même est avantageuse.
\end{exercise}

\chapter{Les neurones selon le nombre de dendrites}
\chaptermark{Le nombre de dendrites}
\label{sec:dendrites}

\section{Le nombre d'entrées, paramètre du circuit}

Au chapitre~\ref{sec:neurontypes}, nous avons trié les neurones selon ce que libère leur sortie, et au chapitre~\ref{sec:eetypes} selon le composant qu'ils réalisent. Il existe un troisième critère, que l'ingénieur remarquera aussitôt: \emph{combien d'entrées possède le neurone}. Les entrées sont collectées par les dendrites, des branches sur lesquelles se posent les axones d'autres neurones (chapitre~\ref{sec:dofa}, section~\ref{sec:weightstore}). Certains neurones n'ont aucune dendrite, d'autres en ont une ou deux, d'autres encore un arbre qui rassemble cent mille entrées. Pour un électronicien, c'est la sortance inverse, le fan-in, et elle correspond presque toujours à la tâche.

Une estimation simple montre pourquoi le nombre et la taille des dendrites comptent autant. La membrane est un condensateur de capacité spécifique $c_m\approx1$~µF/cm$^2$, et plus la surface $A$ du neurone et de ses dendrites est grande, plus il faut apporter de charge pour élever la tension de $\Delta V$. Le temps que met un courant d'entrée $I$ pour y parvenir, si l'on néglige la fuite, vaut
\begin{empheq}[box=\keybox]{equation}
t \;\approx\; \frac{c_m\,A\,\Delta V}{I} .
\label{eq:dendriteload}
\end{empheq}
Un petit neurone muni de quelques dendrites courtes a une surface de l'ordre de $2\cdot10^{-5}$~cm$^2$ et une capacité d'environ $20$~pF. Une seule grosse synapse délivrant quelques nanoampères l'élève de $20\mV$ en moins d'un dixième de milliseconde. Un neurone à grand arbre a une surface une quinzaine de fois plus grande et une capacité d'environ $300$~pF; une synapse ordinaire délivrant environ $0.1$~nA le chargerait en $60\ms$, bien plus que la constante de temps de fuite, c'est-à-dire qu'elle ne le chargerait jamais. Un tel neurone a besoin de dizaines d'entrées simultanées. Ces nombres sont des ordres de grandeur, mais la conclusion n'en dépend pas: \emph{peu d'entrées, c'est rapide et précis; beaucoup d'entrées, c'est lent et par la somme}.

\begin{figure}[t]
\centering
\begin{tikzpicture}[>=Stealth,
  soma/.style={circle,draw,very thick,fill=blue!6,minimum size=0.55cm,inner sep=0pt},
  ax/.style={->,very thick,red!70!black},
  den/.style={very thick,blue!60!black}]
% 0 dendrites: sensor on a line
\begin{scope}[xshift=0cm]
\draw[den,decorate,decoration={snake,amplitude=1pt,segment length=4pt}] (-0.9,0) -- (0,0);
\node[font=\scriptsize] at (-0.9,0.3) {capteur};
\draw[ax] (0,0) -- (1.0,0);
\draw[thick] (0.1,0) -- (0.1,0.45);
\node[soma,minimum size=0.4cm] at (0.1,0.65) {};
\node[font=\scriptsize,align=center] at (0.05,-0.75) {$0$\\ligne à capteur};
\end{scope}
% 1 dendrite
\begin{scope}[xshift=3.3cm]
\draw[den] (-0.9,0) -- (-0.28,0);
\node[soma] at (0,0) {};
\draw[ax] (0.28,0) -- (0.9,0);
\node[font=\scriptsize,align=center] at (0,-0.75) {$1$\\tampon, suiveur};
\end{scope}
% 2 dendrites: one per ear
\begin{scope}[xshift=6.2cm]
\draw[den] (-0.28,0) -- (-0.9,0); \draw[den] (0.28,0) -- (0.9,0);
\node[soma] at (0,0) {};
\draw[ax] (0,-0.28) -- (0,-0.6);
\node[font=\scriptsize] at (-0.9,0.25) {g.}; \node[font=\scriptsize] at (0.9,0.25) {d.};
\node[font=\scriptsize,align=center] at (0,-1.0) {$2$\\ET temporel};
\end{scope}
% ~4 short dendrites: mixer
\begin{scope}[xshift=9.0cm]
\foreach \a in {45,135,225,315}{\draw[den] (\a:0.28) -- (\a:0.7); \fill[blue!60!black] (\a:0.72) circle (0.06);}
\node[soma] at (0,0) {};
\draw[ax] (0.28,0) -- (1.0,0);
\node[font=\scriptsize,align=center] at (0,-1.0) {$\sim4$\\mélangeur};
\end{scope}
% many: summing tree
\begin{scope}[xshift=12.0cm]
\foreach \a in {60,90,120}{\draw[den] (0,0.28) -- ++(\a:0.55) coordinate (b); \foreach \d in {-20,20}{\draw[den] (b) -- ++(\a+\d:0.35);}}
\foreach \a in {-120,-60}{\draw[den] (0,-0.28) -- ++(\a:0.45);}
\node[soma] at (0,0) {};
\draw[ax] (0.28,0) -- (1.0,0);
\node[font=\scriptsize,align=center] at (0,-1.0) {$10^4$--$10^5$\\sommateur};
\end{scope}
\end{tikzpicture}
\caption{Cinq classes selon le nombre d'entrées. En bleu, les dendrites (entrées); en rouge, l'axone (sortie). Moins un neurone a d'entrées, plus il transmet vite et précisément un signal isolé; plus il en a, mieux il additionne et classe. Schéma de principe, et non dessin de la forme des cellules.}
\label{fig:dendriteclasses}
\end{figure}

\section{Zéro dendrite: un capteur au bout d'une ligne}

Les neurones sensoriels du toucher, de la douleur et de l'étirement musculaire (chapitres~\ref{sec:sensors}, \ref{sec:pain}, \ref{sec:muscles}) n'ont aucune dendrite sur leur corps cellulaire. L'axone sort du corps par une seule branche et se divise aussitôt en deux: une branche part vers la peau ou le muscle, et son extrémité sert elle-même de capteur; l'autre va dans la moelle épinière. L'impulsion court du capteur directement à la moelle, sans passer par le corps cellulaire. Pour un ingénieur, c'est une ligne de transmission dont le capteur est à l'extrémité distante, tandis que le corps cellulaire est branché sur le côté, comme une alimentation: il entretient la ligne, mais ne participe pas à la transmission. Le gain, c'est la vitesse et la fiabilité: il n'y a sur le trajet aucun endroit où il faudrait additionner le signal et décider à nouveau s'il faut le transmettre.

Les capteurs de la lumière et du son sont un cas encore plus extrême: ce ne sont pas des neurones collecteurs, mais les transducteurs eux-mêmes, dont l'entrée n'est pas une synapse, mais une protéine réceptrice (fig.~\ref{fig:transducer}).

\section{Une dendrite: le tampon}

Un neurone à une dendrite et un axone reçoit une seule ligne d'entrée et la transmet plus loin. C'est le cas des neurones olfactifs: leur unique dendrite atteint la surface du nez et porte un seul type de récepteur~\cite{Mombaerts2004}; des cellules relais de la rétine, entre les capteurs de lumière et les neurones de sortie de l'œil; des neurones qui relient les capteurs de l'oreille au cerveau. Un ingénieur les appellerait tampons ou suiveurs: ils ne calculent pas, ils acheminent un seul signal, en changeant parfois sa forme, par exemple en transformant la tension continue d'un capteur en impulsions, comme au chapitre~\ref{sec:sensors}.

\section{Deux dendrites: un ET temporel}

Chez les mammifères, les détecteurs de coïncidences qui déterminent la direction d'un son (fig.~\ref{fig:itd}) ont souvent exactement deux dendrites, orientées dans des directions opposées: l'une reçoit les entrées de l'oreille gauche, l'autre celles de l'oreille droite~\cite{Grothe2010}. Pourquoi séparer les entrées? Rappelons la formule~\eqref{eq:shunt}: quand de nombreux canaux excitateurs sont ouverts sur une même portion de membrane, la tension s'y rapproche du potentiel d'équilibre, et chaque entrée supplémentaire ajoute de moins en moins. Beaucoup d'entrées venant d'une même oreille, rassemblées sur une même dendrite, la saturent donc et ne peuvent pas, à elles seules, amener le neurone au seuil. En revanche, une entrée sur chaque dendrite s'additionnent de façon presque linéaire. Deux dendrites font du neurone un \og ET\fg{} plus strict: ne décharge que si les signaux sont arrivés des deux côtés, et ensemble. Des modèles ont montré qu'une telle séparation améliore effectivement la détection des coïncidences~\cite{AgmonSnir1998}.

\section{Quelques dendrites courtes: le mélangeur et le relais}

\emph{Les mélangeurs.} Dans le circuit d'apprentissage moteur du chapitre~\ref{sec:motorlearning}, environ soixante-dix milliards de petits neurones \nt{GLUT} déploient l'entrée en un vecteur très long. Chacun n'a qu'environ quatre dendrites courtes, et chacune se termine par une \og griffe\fg{} qui enserre la terminaison d'une seule entrée. Chaque mélangeur ne voit donc qu'environ quatre entrées parmi beaucoup, et fonctionne comme un petit élément \og ET\fg{} sur son quadruplet. Parmi $M$ lignes d'entrée, on peut choisir $\binom{M}{4}$ quadruplets différents: pour $M=100$, cela fait près de quatre millions. Pourquoi tant? Un élément à seuil à $n$ entrées peut séparer correctement en deux classes au plus
\begin{empheq}[box=\keybox]{equation}
P_{\max} \;\approx\; 2n
\label{eq:cover}
\end{empheq}
motifs aléatoires~\cite{Cover1965}. Le neurone de sortie du même circuit reçoit environ $10^5$ entrées venant des mélangeurs, et selon~\eqref{eq:cover}, sa borne supérieure est de l'ordre de $2\cdot10^5$ associations différentes. C'est la limite d'un élément idéal: si tous les poids sont positifs, comme pour les synapses excitatrices, et qu'il faut une marge contre le bruit, l'estimation pour ce même neurone est d'environ $4\cdot10^4$ associations~\cite{Brunel2004}. Les mélangeurs à peu d'entrées servent justement à fournir au grand sommateur de nombreuses entrées différentes et presque indépendantes~\cite{Marr1969,Albus1971}.

\emph{Les relais.} Sur le chemin qui va de l'oreille aux détecteurs de direction se trouvent des neurones à quelques dendrites courtes qui ne reçoivent que quelques synapses énormes, et certains n'en reçoivent qu'une en pratique. Selon~\eqref{eq:dendriteload}, c'est exactement ce qu'il faut: une petite capacité et un fort courant, et le neurone décharge une fraction de milliseconde après chaque impulsion d'entrée, en ne perdant presque rien de la précision temporelle. L'un de ces relais reçoit \nt{GLUT} et émet \nt{GLYC}: c'est un inverseur d'une précision de quelques dizaines de microsecondes, qui fournit une inhibition rapide aux détecteurs de direction~\cite{BorstSoria2012}.

\section{Des milliers d'entrées: le sommateur et le classifieur}

À l'autre bout de l'échelle se trouvent les neurones \nt{GLUT} du cortex, avec une dizaine de milliers d'entrées, et les neurones \nt{GABA} de sortie du circuit d'apprentissage moteur, avec cent mille. Selon~\eqref{eq:dendriteload}, un tel neurone est lent et ne peut pas répondre à une entrée isolée: il additionne. C'est le neurone intégrateur du chapitre~\ref{sec:glutcomputer}, le sommateur à partir duquel on construit les classifieurs. Un grand arbre apporte aussi du neuf: ses branches fonctionnent comme des sous-circuits distincts dotés de leurs propres seuils, si bien qu'un seul neurone se comporte comme un petit réseau multicouche~\cite{Beniaguev2021,Gidon2020}.

\section{Des dendrites qui servent aussi de sortie}

Certains neurones n'ont pas d'axone du tout, et leurs dendrites servent à la fois d'entrée et de sortie: elles reçoivent des signaux et libèrent elles-mêmes le neurotransmetteur. L'exemple le mieux étudié est celui des neurones \nt{GABA} de la rétine qui participent à la détection de la direction du mouvement (chapitre~\ref{sec:gaba}, fig.~\ref{fig:direction}). Dans un tel neurone, chaque branche répond d'elle-même au mouvement allant du corps cellulaire vers sa pointe et émet une inhibition depuis cette pointe~\cite{Euler2002}. Un neurone, ce sont plusieurs détecteurs de direction indépendants, un par branche. Pour un ingénieur, ce n'est pas un élément unique, mais un circuit intégré à plusieurs canaux dans un même boîtier.

\section{Bilan}

\begin{table}[t]
\centering
\footnotesize
\setlength{\tabcolsep}{3pt}
\renewcommand{\arraystretch}{1.15}
\begin{tabular}{>{\raggedright}p{1.9cm}>{\raggedright}p{3.4cm}>{\raggedright}p{2.6cm}>{\raggedright}p{1.8cm}>{\raggedright\arraybackslash}p{3.8cm}}
\toprule
Dendrites & Exemple & Composant & Entrées & Ce que l'on sait \\
\midrule
$0$ & capteurs du toucher, de la douleur, de l'étirement & ligne avec capteur à l'extrémité & --- & établi \\
$1$ & neurones olfactifs, cellules relais de la rétine, neurones de l'oreille & tampon, suiveur & $1$ à quelques-unes & établi \\
$2$ & détecteurs de direction du son & ET temporel & deux faisceaux & structure établie; gain dû à la séparation des entrées montré sur des modèles \\
$\sim4$ & mélangeurs du circuit d'apprentissage moteur & petit élément \og ET\fg{} & $\sim4$ & établi; rôle de déploiement de l'entrée: hypothèse bien fondée \\
quelques-unes, courtes & relais sur le chemin de l'oreille & relais, inverseur & $1$ à quelques énormes & établi \\
des milliers & neurones \nt{GLUT} du cortex, neurones de sortie de l'apprentissage moteur & sommateur, classifieur & $10^4$--$10^5$ & établi; branches comme sous-circuits mesurées sur quelques exemples \\
dendrites-sorties & neurones \nt{GABA} de direction de la rétine & circuit intégré multicanal & beaucoup & mesuré \\
\bottomrule
\end{tabular}
\caption{Les neurones selon le nombre de dendrites.}
\label{tab:dendrites}
\end{table}

\begin{keyblock}
\begin{proposition}[Le nombre d'entrées suit la tâche]
\label{prop:dendrites}
Là où il faut la vitesse et la précision d'un signal isolé (capteurs, relais, détecteurs de coïncidences), les neurones ont peu de dendrites, peu d'entrées et de grosses synapses: une petite capacité~\eqref{eq:dendriteload} se charge vite. Là où il faut additionner et classer, les neurones ont de grands arbres à des milliers d'entrées. La structure des classes est mesurée; l'explication computationnelle est bien fondée, mais reste une hypothèse pour de nombreuses classes.
\end{proposition}
\end{keyblock}

Le tableau~\ref{tab:dendrites} complète les deux tris précédents: par neurotransmetteur (tableau~\ref{tab:types}) et par composant (tableau~\ref{tab:eetypes}). Tous trois regardent le même élément sous des angles différents: ce qu'il émet, ce qu'il calcule et combien il écoute.

\section{Exercices}

\begin{exercise}[\lvlA]
\label{ex:19:1}
\emph{Combien d'entrées pour un grand neurone.} Un neurone avec $C=300$~pF et $\tau_m=20\ms$ reçoit des synapses, sources de courant de $0.1$~nA; seuil à $20\mV$ au-dessus du repos. (a)~Combien de synapses simultanées atteignent le seuil tout court; en $10\ms$; en $5\ms$? (b)~Et un neurone de $C=20$~pF sous une synapse de $4$~nA: en combien de temps?
\end{exercise}

\begin{exercise}[\lvlA]
\label{ex:19:2}
\emph{Pourquoi quatre.} Un mélangeur est un \og ET\fg{} sur $k$ lignes parmi $M=100$; dans un motif, la moitié des lignes est active; il y a $7\cdot10^{10}$ mélangeurs. (a)~Combien répondent au motif pour $k=2$; $4$; $40$? (b)~De quel facteur les mélangeurs augmentent-ils le nombre d'associations~\eqref{eq:cover} d'un neurone de sortie à $10^5$ entrées par rapport à $M=100$ lignes directes?
\end{exercise}

\begin{exercise}[\lvlB]
\label{ex:19:3}
\emph{D'où vient $2n$.} Un hyperplan passant par l'origine sépare $P$ points en position générale dans un espace à $n$ dimensions en deux classes de $C(P,n)=2\sum_{k=0}^{n-1}\binom{P-1}{k}$ façons. (a)~Démontrez que pour $P=2n$, exactement la moitié des $2^P$ partitions est séparable. (b)~Quelle fraction est séparable pour $n=100$ et $P=180$; $220$?
\end{exercise}

\begin{exercise}[\lvlB]
\label{ex:19:4}
\emph{Deux dendrites comme \og ET\fg{} strict.} Une dendrite est un compartiment à fuite $g_L$; chaque entrée y ouvre une conductance $g_0=0.5\,g_L$ de potentiel $E_E$; le corps voit $\kappa(V_1+V_2)$; seuil $\theta=1.1\,\kappa E_E$. Pourquoi aucun nombre d'entrées sur une seule dendrite ne déclenche-t-il le neurone, et combien d'entrées faut-il sur chacune?
\end{exercise}

\begin{exercise}[\lvlB]
\label{ex:19:5}
\emph{Concevoir un relais.} La gigue d'une impulsion vaut $\sigma_t=\sigma_V/\dot V$, où $\dot V$ est la pente de la tension au seuil. On veut $\sigma_t\le20$~µs avec un bruit $\sigma_V=1\mV$; négligez la fuite. Combien de synapses synchrones de $0.1$~nA faut-il pour cela à un neurone avec $C=300$~pF, et quelle est la gigue d'un neurone avec $C=20$~pF et une seule synapse de $4$~nA?
\end{exercise}

\begin{exercise}[\lvlB]
\label{ex:19:6}
\emph{Simulation: charge d'une longue dendrite.} Simulez un câble passif $\tau_m\partial_tV=\lambda^2\partial_x^2V-V+i/g$ de longueur $L=\ell/\lambda$ à extrémités fermées ($40$ compartiments, Euler). Une brève impulsion de courant entre au bout distant: quand, et jusqu'à quelle fraction de la tension qu'aurait la même charge étalée sur la dendrite, le bout proche monte-t-il pour $L=0.2$, $1$, $2$?
\end{exercise}

\begin{exercise}[\lvlC]
\label{ex:19:7}
\emph{Recherche: le nombre optimal d'entrées.} Capacité $C=C_0+nc_s$; $m$ entrées de courant $I_s$ agissent simultanément; le neurone mémorise $P\approx2n$ associations~\eqref{eq:cover}. Comment le temps de réponse dépend-il de $P$ à $m$ constant, et à fraction constante $m=\varphi n$? Proposez un critère de coût et trouvez le $n$ optimal pour un relais et pour un classifieur.
\end{exercise}

\chapter{Ce que font les neurones pendant le sommeil}
\chaptermark{Le sommeil}
\label{sec:sleep}
\begin{chaptergoals}
\item pourquoi le sommeil est un ensemble de modes commutés par \nt{ACET}, \nt{NORA} et \nt{SERO}, et non un arrêt;
\item comment pression de sommeil et deux seuils font un relais à hystérésis calé sur le rythme circadien;
\item comment la fatigue fait du cortex un oscillateur à relaxation, et \nt{ACET} arrête les oscillations;
\item comment la journée est rejouée en accéléré, et pourquoi le sommeil pourrait réduire les poids.
\end{chaptergoals}

\section{Le sommeil n'est pas un arrêt, mais un autre mode}

Devant un ordinateur éteint, un ingénieur dira: il ne fait rien. Avec un cerveau endormi, ce n'est pas le cas. Pendant le sommeil, les neurones ne se taisent pas: en sommeil profond, le cortex émet des impulsions; en sommeil paradoxal, presque autant que le jour. Ce qui change, ce n'est pas le fait que les neurones \emph{fonctionnent}, mais la \emph{manière} dont ils fonctionnent. Le sommeil est un ensemble de modes de la machine, et ce sont les mêmes canaux de diffusion que le jour qui les commutent.

Pour chaque mode, on peut écrire \og l'état des registres\fg{} (tableau~\ref{tab:sleepmodes}). Le jour, \nt{ACET}, \nt{NORA} et \nt{SERO} sont actifs, les capteurs sont connectés, les muscles obéissent. En sommeil lent, le sommeil profond, les trois canaux s'apaisent, et l'entrée venant des organes des sens est affaiblie~\cite{Hasselmo1999}: la libération d'acétylcholine dans le cortex chute, et les neurones \nt{NORA} et \nt{SERO} émettent moins d'impulsions que le jour~\cite{Marrosu1995,AstonJonesBloom1981,McGintyHarper1976}. En sommeil paradoxal, le tableau est étrange: \nt{ACET} remonte à son niveau diurne~\cite{Marrosu1995}, tandis que \nt{NORA} et \nt{SERO} se taisent presque complètement~\cite{AstonJonesBloom1981,McGintyHarper1976,JacobsAzmitia1992}. En sommeil paradoxal, les muscles du corps sont coupés: les neurones \nt{ACET} qui les commandent sont fortement inhibés par \nt{GABA} et \nt{GLYC}~\cite{BrooksPeever2012}, par le même veto qu'au chapitre~\ref{sec:gaba}, mais appliqué à toute la sortie à la fois.

\begin{table}[t]
\centering
\footnotesize
\setlength{\tabcolsep}{4pt}
\renewcommand{\arraystretch}{1.15}
\begin{tabular}{>{\raggedright}p{3.1cm}>{\raggedright}p{2.9cm}>{\raggedright}p{3.2cm}>{\raggedright\arraybackslash}p{3.4cm}}
\toprule
& veille & sommeil lent & sommeil paradoxal \\
\midrule
\nt{ACET} (écriture, chap.~\ref{sec:acet}) & élevé & bas & élevé \\
\nt{NORA} (gain, chap.~\ref{sec:nora}) & moyen, bouffées & bas & presque silencieux \\
\nt{SERO} (chap.~\ref{sec:serocomputer}) & moyen & bas & presque silencieux \\
entrée des capteurs & connectée & affaiblie & affaiblie \\
sortie vers les muscles & connectée & affaiblie & coupée par inhibition \\
activité du cortex & régulière & bascule lente: activité / silence & régulière, comme le jour \\
\bottomrule
\end{tabular}
\caption{L'état des registres de la machine dans les trois modes. Toutes les lignes sont mesurées; les niveaux d'\nt{ACET}, de \nt{NORA} et de \nt{SERO} par des enregistrements directs selon les stades du sommeil~\cite{Marrosu1995,AstonJonesBloom1981,McGintyHarper1976}.}
\label{tab:sleepmodes}
\end{table}

\section{Quand dormir: un relais à hystérésis}

Qu'est-ce qui décide du moment où la machine bascule dans le sommeil? Un schéma simple en deux parties fonctionne bien~\cite{Borbely1982}. La première partie est la \emph{pression de sommeil} $S$, qui s'accumule pendant la veille et se dissipe pendant le sommeil. C'est un condensateur qui se charge le jour et se décharge la nuit:
\begin{empheq}[box=\keybox]{equation}
\frac{\dd S}{\dd t} \;=\; \frac{1-S}{\tau_r}\ \ \text{(veille)},\qquad
\frac{\dd S}{\dd t} \;=\; -\frac{S}{\tau_d}\ \ \text{(sommeil)} .
\label{eq:sleeppressure}
\end{empheq}
La seconde partie, ce sont deux seuils: la machine s'endort quand $S$ atteint le seuil haut $H(t)$, et se réveille quand $S$ redescend au seuil bas $L(t)$. Les deux seuils oscillent doucement au rythme circadien du chapitre~\ref{sec:chemoutput}. On obtient un relais à hystérésis, le trigger de Schmitt de la section~\ref{sec:bistable}, et avec le condensateur, il forme un oscillateur à relaxation dont la phase est calée par le rythme circadien~\eqref{eq:pll}.

\begin{figure}[t]
\centering
\begin{tikzpicture}[>=Stealth,x=0.16cm,y=4.2cm]
\foreach \a/\b in {0/4.3,21.2/29.3,45.4/53.4,69.5/72}{\fill[blue!8] (\a,0) rectangle (\b,1);}
\draw[->,gray!70!black] (0,0) -- (75,0) node[right,font=\small,black] {$t$, h};
\draw[->,gray!70!black] (0,0) -- (0,1.08) node[left,font=\small,black] {$S$};
\foreach \t in {0,24,48,72}{\draw[gray!60!black] (\t,-0.02) -- (\t,0.02); \node[below,font=\scriptsize] at (\t,0) {\t};}
\draw[thick,densely dashed,red!70!black] plot file {figures/sleep_H.dat};
\draw[thick,densely dashed,green!45!black] plot file {figures/sleep_L.dat};
\draw[very thick,blue!60!black] plot file {figures/sleep_S.dat};
\node[font=\scriptsize,red!70!black,anchor=west] at (73,0.72) {$H$: s'endormir};
\node[font=\scriptsize,green!45!black,anchor=west] at (73,0.24) {$L$: se réveiller};
\node[font=\scriptsize,blue!60!black] at (25.25,0.93) {sommeil};
\node[font=\scriptsize,blue!60!black] at (49.4,0.93) {sommeil};
\end{tikzpicture}
\caption{Pression de sommeil $S$~\eqref{eq:sleeppressure} pour $\tau_r=18.2$~h, $\tau_d=4.2$~h. Le jour, $S$ se charge jusqu'au seuil haut $H$; la nuit (zones colorées), elle se décharge jusqu'au seuil bas $L$; les deux seuils oscillent avec le rythme circadien. La machine trouve d'elle-même un régime d'environ $16$ heures de veille et $8$ heures de sommeil.}
\label{fig:twoprocess}
\end{figure}

Dans notre modèle, avec les valeurs usuelles $\tau_r=18.2$~h et $\tau_d=4.2$~h, la machine trouve d'elle-même $16.1$ heures de veille et $8.0$ heures de sommeil (fig.~\ref{fig:twoprocess}). Le modèle explique aussi ce qui semble étrange: après quarante heures sans sommeil, la pression atteint $0.90$, mais le sommeil de récupération ne dure dans le modèle que $8.8$ heures, et non une journée de plus que d'habitude. La décharge est exponentielle, et une charge élevée s'écoule vite; la \og dette\fg{} se rembourse non par la durée, mais par la profondeur du sommeil.

Qu'est-ce qui joue physiquement le rôle de $S$? Un candidat sérieux est l'adénosine, le \og compteur de charge\fg{} de l'annexe~\ref{sec:othermediators}: sa concentration augmente pendant la veille et baisse pendant le sommeil~\cite{PorkkaHeiskanen1997,Dunwiddie2001}. Que la pression de sommeil soit l'adénosine et rien d'autre reste une hypothèse.

\section{Le sommeil lent: tout le réseau en oscillateur à relaxation}

En sommeil profond, les neurones du cortex fonctionnent en alternance, tout le groupe ensemble: moins d'une fois par seconde, ils s'allument ensemble pendant quelques centaines de millisecondes (l'\emph{état actif}) et se taisent ensemble (l'\emph{état de silence}); la fréquence de cette bascule est inférieure à $1$~Hz, le plus souvent $0.3$--$0.4$~Hz sous anesthésie~\cite{Steriade1993}. Cette bascule lente, nous pouvons l'assembler avec des pièces que nous avons déjà. Prenons un groupe de neurones \nt{GLUT} à forte rétroaction sur lui-même, la mémoire du chapitre~\ref{sec:glutcomputer}, qui possède deux états stables, et ajoutons une fatigue lente, comme dans le générateur de rythme du chapitre~\ref{sec:muscles}:
\begin{empheq}[box=\keybox]{equation}
\tau\,\frac{\dd r}{\dd t} = -r + F\bigl(w\,r + I - a\bigr),
\qquad
\tau_a\,\frac{\dd a}{\dd t} = -a + b\,r .
\label{eq:updown}
\end{empheq}
Le groupe s'allume, travaille et se fatigue; la fatigue $a$ ronge l'état haut, et le groupe tombe dans le silence; dans le silence, la fatigue se dissipe, l'état bas disparaît, et le groupe se rallume. On obtient le même oscillateur à relaxation que pour la pression de sommeil, mais environ quatre-vingt mille fois plus rapide.

\begin{figure}[t]
\centering
\begin{tikzpicture}[>=Stealth,x=2.9cm,y=1.0cm]
\begin{scope}[yshift=1.9cm]
\draw[->,gray!70!black] (0,0) -- (4.1,0);
\draw[->,gray!70!black] (0,0) -- (0,1.25) node[left,font=\small,black] {$r$};
\draw[thick,blue!60!black] plot file {figures/updown_sleep.dat};
\node[font=\scriptsize,anchor=west] at (4.15,0.5) {sommeil: $b=5$};
\end{scope}
\draw[->,gray!70!black] (0,0) -- (4.1,0) node[right,font=\small,black] {$t$, s};
\draw[->,gray!70!black] (0,0) -- (0,1.25) node[left,font=\small,black] {$r$};
\draw[thick,red!70!black] plot file {figures/updown_wake.dat};
\node[font=\scriptsize,anchor=west] at (4.15,0.5) {jour: $b=1$};
\foreach \t in {0,1,2,3,4}{\draw[gray!60!black] (\t,-0.05) -- (\t,0.05); \node[below,font=\scriptsize] at (\t,0) {\t};}
\end{tikzpicture}
\caption{Un même groupe~\eqref{eq:updown} dans deux modes: $w=5$, $I=2$, $\theta=2.5$, $\tau=10\ms$, $\tau_a=0.3$~s. Avec une forte fatigue ($b=5$, comme en sommeil lent), le groupe bascule entre activité et silence avec une période de $1.1$~s (valeur du modèle; dans le cortex, la période dépasse une seconde, et sous anesthésie elle est d'environ $3$~s). Affaiblissons la fatigue ($b=1$), comme le fait l'acétylcholine, et le même groupe fonctionne de façon régulière.}
\label{fig:updown}
\end{figure}

Et qu'est-ce qui fait passer le réseau de la bascule au fonctionnement régulier du jour? L'acétylcholine supprime dans les neurones les courants lents qui créent la fatigue~\cite{McCormick1992}. Dans notre modèle, avec une forte fatigue, $b=5$, le groupe bascule avec une période de $1.1$~s (plus vite que la bascule mesurée, mais du même ordre) et reste actif environ $35\%$ du temps, en moyenne sur un nombre entier de périodes; avec une fatigue faible, $b=1$, le même groupe fonctionne régulièrement (fig.~\ref{fig:updown}). Un seul registre, le niveau d'\nt{ACET}, transforme l'oscillateur en amplificateur. Par-dessus la bascule lente du sommeil passent aussi des bouffées plus rapides d'un rythme de $7$ à $15$ oscillations par seconde; elles naissent d'une boucle \nt{GLUT}--\nt{GABA} entre le cortex et un nœud relais intermédiaire~\cite{SteriadeMcCormickSejnowski1993}, parente du rythme du chapitre~\ref{sec:gaba}.

\section{Les répétitions: la journée en accéléré}

Au chapitre~\ref{sec:acet}, nous avons vu qu'en sommeil lent, les neurones qui ont travaillé ensemble le jour déchargent de nouveau ensemble, et dans le même ordre. Ajoutons un détail d'ingénieur: les répétitions se font \emph{en accéléré}. Une séquence qui occupait des secondes le jour est rejouée pendant le sommeil en quelques dixièmes de seconde~\cite{Nadasdy1999}: environ vingt fois plus vite dans l'hippocampe~\cite{LeeWilson2002}, six à sept fois plus vite dans le cortex~\cite{Euston2007}. Et elle n'est pas rejouée n'importe quand: les brèves bouffées de répétitions dans une région se calent sur la bascule activité--silence et sur les bouffées rapides de rythme du cortex. Si l'on renforce artificiellement cette coordination, la mémoire s'améliore~\cite{Maingret2016}: c'est là un lien causal, et non une coïncidence.

Pourquoi la machine repasse-t-elle la journée en accéléré? L'ingénieur connaît une difficulté qu'on appelle l'\emph{oubli catastrophique}: un réseau qui apprend de nouvelles choses à la suite, l'une après l'autre, abîme les anciennes. Le remède est l'\emph{entrelacement}: apprendre le nouveau mêlé à l'ancien, par petites portions, de nombreuses fois. L'hypothèse des deux systèmes d'apprentissage~\cite{McClelland1995} veut qu'une mémoire rapide enregistre la journée entière, puis, pendant le sommeil, la rejoue peu à peu, en l'entremêlant et de nombreuses fois, au cortex lent. C'est la même paire \og tampon rapide, stockage lent\fg{} qu'au chapitre~\ref{sec:acet}, et la même paire d'élèves rapide et lent qu'au chapitre~\ref{sec:motorlearning}. Les répétitions et leur coordination sont mesurées; que leur fonction soit l'apprentissage entrelacé de la mémoire lente est une hypothèse bien fondée.

\section{La renormalisation des poids}

Le jour, les règles de Hebb du chapitre~\ref{sec:dofa} renforcent surtout les connexions. Si l'on ne fait que renforcer, les poids croissent, et avec eux la dépense d'énergie, tandis que la sensibilité baisse: un neurone dont toutes les entrées sont fortes cesse de les distinguer. Selon l'hypothèse de l'\emph{homéostasie synaptique}~\cite{TononiCirelli2014}, le sommeil sert entre autres à ramener les poids à leur niveau de travail: toutes les connexions sont affaiblies d'un même facteur, si bien que leurs \emph{rapports}, c'est-à-dire ce qui a été appris, sont conservés. Pour un ingénieur, c'est une normalisation, comme dans la règle~\eqref{eq:oja}, mais faite non pas en marche, mais dans un créneau séparé.

Les mesures directes ne le contredisent pas: chez la souris, la surface de contact des synapses du cortex est après le sommeil inférieure d'environ $18\%$ à ce qu'elle est après la veille; la diminution est proportionnelle à la taille, et les contacts les plus gros et les plus stables ne changent presque pas~\cite{deVivo2017}. La mise à l'échelle des poids est mesurée; qu'elle soit la tâche principale du sommeil reste une hypothèse.

\section{Le sommeil paradoxal: une machine déconnectée de ses entrées et de ses sorties}

En sommeil paradoxal, le cortex fonctionne presque comme le jour, mais l'entrée venant des organes des sens est affaiblie et la sortie vers les muscles est coupée par inhibition. Pour un ingénieur, c'est un mode familier: l'\emph{essai sur banc}. Le circuit fonctionne à pleine puissance, mais ses entrées sont reliées à un générateur interne et ses sorties déconnectées des actionneurs, pour ne rien casser. Selon une hypothèse, les rêves sont justement un tel déroulement interne du modèle du monde construit pendant la journée; ce que fait alors le réseau et s'il apprend, on l'ignore. Seul est mesuré ici ce qui figure dans le tableau~\ref{tab:sleepmodes}: quel canal est actif, lequel se tait, et le fait que la sortie est bloquée.

\section{La maintenance et le sommeil local}

On a longtemps cru que, pendant le sommeil, les espaces intercellulaires s'élargissent et que le cerveau est plus vite lavé des déchets du métabolisme~\cite{Xie2013}. Des expériences récentes, qui mesuraient ce lavage d'une autre manière, ont obtenu l'inverse: pendant le sommeil, il est plus lent~\cite{Miao2024}. La question est aujourd'hui ouverte, et nous la laisserons ouverte.

Un autre fait est plus solide: le sommeil est une propriété des réseaux locaux, et non de toute la machine à la fois. Si l'on empêche longtemps un animal de dormir, certaines zones de son cortex, alors qu'il est éveillé et qu'il bouge, sombrent brièvement dans le silence comme en sommeil lent, et à ces moments-là l'animal se trompe plus souvent~\cite{Vyazovskiy2011}. Chaque réseau se fatigue et bascule de lui-même; le mode global apparaît quand les canaux de diffusion font basculer tous les réseaux à la fois.

\begin{keyblock}
\begin{proposition}[Le sommeil comme ensemble de modes]
\label{prop:sleep}
Pendant le sommeil, les neurones ne s'éteignent pas: les canaux de diffusion font passer la machine dans d'autres modes. Le moment de dormir est décidé par un relais à hystérésis~\eqref{eq:sleeppressure} calé par le rythme circadien. En sommeil lent, le réseau devient un oscillateur à relaxation~\eqref{eq:updown} et rejoue la journée en accéléré; en sommeil paradoxal, il fonctionne déconnecté de ses entrées et de ses sorties. Les modes, les répétitions et la mise à l'échelle des poids sont mesurés; leur fonction (apprentissage entrelacé et renormalisation) relève d'hypothèses bien fondées.
\end{proposition}
\end{keyblock}

\section{Bilan}

\begin{table}[t]
\centering
\footnotesize
\setlength{\tabcolsep}{4pt}
\renewcommand{\arraystretch}{1.15}
\begin{tabular}{>{\raggedright}p{3.2cm}>{\raggedright}p{4.4cm}>{\raggedright\arraybackslash}p{5.2cm}}
\toprule
Ce qui se passe & Comment & Ce que l'on sait \\
\midrule
changement de mode & registres \nt{ACET}, \nt{NORA}, \nt{SERO} (tabl.~\ref{tab:sleepmodes}) & mesuré \\
quand dormir & condensateur $S$ et relais à hystérésis~\eqref{eq:sleeppressure} & le modèle décrit bien les expériences; l'adénosine comme $S$: hypothèse \\
bascule lente & groupe à rétroaction et fatigue~\eqref{eq:updown} & bascule mesurée; le modèle est le schéma le plus simple \\
journée en accéléré & répétitions $6$ à $20$ fois plus rapides, coordonnées avec la bascule & mesuré; renforcer la coordination améliore la mémoire \\
pourquoi les répétitions & apprentissage entrelacé de la mémoire lente & hypothèse bien fondée \\
renormalisation des poids & mise à l'échelle des connexions pendant le sommeil & diminution des contacts mesurée; rôle principal du sommeil: hypothèse \\
sommeil paradoxal & fonctionnement avec entrées et sortie déconnectées & mode mesuré; fonction inconnue \\
lavage & élargissement des espaces intercellulaires & résultats contradictoires \\
sommeil local & certaines zones sombrent dans le silence & mesuré \\
\bottomrule
\end{tabular}
\caption{Ce que font les neurones pendant le sommeil.}
\label{tab:sleep}
\end{table}

Le tableau~\ref{tab:sleep} résume le chapitre. Le sommeil s'est révélé non pas une pause dans le travail de la machine, mais sa maintenance en marche: changement de mode par les mêmes canaux de diffusion, oscillateur fait des mêmes pièces que le rythme de la marche, journée rejouée en accéléré pour la mémoire lente, et renormalisation des poids. Le jour, la machine écrit; la nuit, elle réécrit et met de l'ordre dans ce qu'elle a écrit.

\section{Exercices}

\begin{exercise}[\lvlA]
\label{ex:20:1}
\emph{Un relais sans rythme circadien.} Seuils constants: $H=0.67$, $L=0.17$ (moyennes des seuils de la fig.~\ref{fig:twoprocess}). Déduisez de~\eqref{eq:sleeppressure} les durées de veille et de sommeil et la période de l'oscillateur pour $\tau_r=18.2$~h, $\tau_d=4.2$~h. Pourquoi le modèle à seuils oscillants donne-t-il pourtant $16.1+8.0=24$~h? À quel composant du chapitre~\ref{sec:chemoutput} cet effet se rattache-t-il?
\end{exercise}

\begin{exercise}[\lvlA]
\label{ex:20:2}
\emph{La dette de sommeil.} Un sommeil entamé à la pression $S_0$ dure $t_s=\tau_d\ln(S_0/L)$, $\tau_d=4.2$~h, $L$: seuil bas au réveil. (a)~Après $40$~h sans sommeil, $S_0=0.90$, et on a dormi $8.8$~h. Que valait $L$? Et avec l'habituel $L\approx0.092$? (b)~En une nuit, la surface des contacts synaptiques baisse de $18\%$: de combien les poids doivent-ils croître le jour?
\end{exercise}

\begin{exercise}[\lvlB]
\label{ex:20:3}
\emph{Concevoir les seuils.} Choisissez des seuils constants $H$ et $L$ pour lesquels le relais~\eqref{eq:sleeppressure} avec $\tau_r=18.2$~h et $\tau_d=4.2$~h donne exactement $16$~h de veille et $8$~h de sommeil. En quoi une telle machine est-elle moins bonne qu'une machine à seuils oscillants si, une nuit, on la réveille au bout de $4$~heures?
\end{exercise}

\begin{exercise}[\lvlB]
\label{ex:20:4}
\emph{La période de la bascule lente.} Dans le modèle~\eqref{eq:updown}, $F(x)=1/\bigl(1+e^{-(x-\theta)/k}\bigr)$, $w=5$, $I=2$, $\theta=2.5$, $k=0.25$, $b=5$, $\tau\ll\tau_a=0.3$~s. (a)~Trouvez les points de rebroussement $a_{\text{h}}$ et $a_{\text{b}}$ de la courbe $r=F(wr+I-a)$, où disparaissent l'état actif et le silence. (b)~En prenant $r\approx1$ en activité et $r\approx0$ en silence, trouvez la période et la fraction d'activité.
\end{exercise}

\begin{exercise}[\lvlB]
\label{ex:20:5}
\emph{Quand la bascule s'arrête.} Dans le modèle de l'exercice~\ref{ex:20:4}, le point fixe est à l'intersection de la courbe $a=wr+I-F^{-1}(r)$ et de la droite $a=br$. Il y a oscillation tant que l'intersection est sur la branche médiane, entre les points de rebroussement. Pour quels $b$? Comment \nt{ACET}, en modifiant $b$, transforme-t-il l'oscillateur en amplificateur?
\end{exercise}

\begin{exercise}[\lvlB]
\label{ex:20:6}
\emph{Simulation: dette ou phase.} Dans \texttt{sleep\_models.py}, prenez le premier réveil après $48$~h et maintenez la machine éveillée $D=24$, $32$, $40$, $48$, $64$~h (\texttt{stay\_awake}, \texttt{days=8}). Combien dure le sommeil de récupération pour chaque $D$? Qu'est-ce qui détermine sa durée?
\end{exercise}

\begin{exercise}[\lvlB]
\label{ex:20:7}
\emph{Simulation: oubli et entrelacement.} Un neurone linéaire $y=\mathbf w\cdot\mathbf x$, $n=40$, apprend selon la règle $\mathbf w\leftarrow\mathbf w-0.5\,(y-y^*)\,\mathbf x$. Les tâches A et B comptent chacune $10$ motifs, $x_j\sim\mathcal N(0,1/n)$, $y^*=\pm1$. Quelle est l'erreur quadratique moyenne sur A après $200$ époques de A puis $200$ époques de B? Et après $200$ époques de A et B entremêlées?
\end{exercise}

\begin{exercise}[\lvlC]
\label{ex:20:8}
\emph{Recherche: entrelacement ou renormalisation?} Le neurone de l'exercice~\ref{ex:20:7}, muni d'un seuil; deux procédures nocturnes: (i)~tous les poids sont multipliés par $0.82$; (ii)~les anciens motifs sont rejoués parmi les nouveaux. Que fait chacune aux rapports des poids et aux réponses? Comment distinguer les hypothèses par la taille des synapses au réveil?
\end{exercise}

\chapter{Une machine faite de neurones vivants}
\chaptermark{Une machine faite de neurones vivants}
\label{sec:livingmachine}
\begin{chaptergoals}
\item pourquoi une culture sans canaux de diffusion émet des salves, et ce qui fixe leur intervalle;
\item comment les réseaux vivants modifient leurs réponses sous rétroaction, et ce qui reste à prouver;
\item comment un organoïde sert de réservoir, sa lecture linéaire étant entraînée dans le silicium;
\item comment on a enseigné aux cultures jusqu'ici, et pourquoi le professeur reste hors de la boîte.
\end{chaptergoals}

\section{Un banc d'essai au schéma ouvert}

Au chapitre~\ref{sec:debugging}, nous avons débogué une machine dont on n'a pas le schéma: la sonde entend mille neurones sur quatre-vingts milliards, et il faut deviner tout le reste. Dans cette situation, un ingénieur ferait une chose simple: il monterait un petit morceau du circuit sur une plaque d'essai, où chaque composant est visible. Avec des neurones, c'est aussi possible.

On dissocie des neurones du cortex et on les cultive sur une puce munie d'une matrice d'électrodes, de quelques dizaines à plusieurs dizaines de milliers. En quelques semaines, les cellules émettent des prolongements et se connectent en un réseau de milliers ou de centaines de milliers de neurones, surtout des \nt{GLUT}, avec une part plus faible de \nt{GABA} qui dépend de la préparation: dans les cultures d'hippocampe, elle est d'environ $6\%$ des neurones~\cite{Benson1994}. Chaque électrode sait à la fois écouter, comme la sonde du chapitre~\ref{sec:debugging}, et injecter du courant, comme un signal de test. La seconde variante du banc, ce sont les \emph{organoïdes}: des boulettes tridimensionnelles de tissu nerveux d'environ un millimètre, cultivées à partir de cellules souches humaines~\cite{Lancaster2013}. Sur de tels bancs, les réseaux vivants fonctionnent pendant des mois; certaines plateformes permettent de mener des expériences sur des organoïdes à distance, par le réseau, plus de cent jours d'affilée~\cite{Jordan2024}. Ce domaine s'appelle le \emph{calcul biologique} ou l'\emph{intelligence organoïde}~\cite{Smirnova2023}.

Le banc diffère du cerveau sur deux points importants. Il est minuscule: il compte cent mille fois moins de neurones. Et il n'a pas de canaux de diffusion: une culture de cortex ne contient aucun neurone \nt{DOPA}, \nt{SERO}, \nt{NORA} ou \nt{ACET}. C'est une machine qui possède un bus de données et un contrôle de flux, mais pas de registres de contrôle. Voyons de quoi elle est capable.

\section{Ce que le réseau fait tout seul}

Laissée au repos, une culture ne se tait pas et ne produit pas non plus un bruit uniforme. De temps en temps, presque tout le réseau s'embrase d'un coup, en une \emph{salve}, puis retombe dans le silence; l'intervalle entre les salves va d'une seconde à plusieurs minutes selon la culture~\cite{Wagenaar2006}. Pour un ingénieur, le tableau est familier: un amplificateur à forte contre-réaction positive et sans charge devient un oscillateur.

Nous connaissons déjà le mécanisme. Les fortes connexions mutuelles des neurones \nt{GLUT} déclenchent une avalanche, comme au chapitre~\ref{sec:gaba} lorsque les conditions de la proposition~\ref{prop:ei} ne sont pas remplies. L'avalanche épuise vite la réserve de neurotransmetteur dans les synapses~\eqref{eq:depression}, et le réseau se tait. La réserve se reconstitue avec une constante de temps $\tau_{\text{rec}}$, et lorsque la fraction $x^*$ suffisante pour une nouvelle avalanche est rétablie, une nouvelle salve suit. L'intervalle entre salves vaut
\begin{empheq}[box=\keybox]{equation}
T_{\text{salve}} \;\approx\; \tau_{\text{rec}}\,\ln\frac{1}{1-x^*} .
\label{eq:burstinterval}
\end{empheq}
Avec $\tau_{\text{rec}}=0.8\,\text{s}$ du chapitre~\ref{sec:code} et $x^*=0.9$, on obtient environ deux secondes; avec $x^*=0.99$, environ quatre. C'est une estimation du modèle avec le $\tau_{\text{rec}}$ du livre, et elle tombe dans le bas de la plage des intervalles mesurés. Une mesure directe sur des microcultures montre que la réserve de neurotransmetteur se reconstitue après une salve en quelques dizaines de secondes~\cite{CohenSegal2011}: $\tau_{\text{rec}}$ y est donc plus grand, et avec un tel $\tau_{\text{rec}}$ la formule~\eqref{eq:burstinterval} donne des dizaines de secondes et des minutes, comme dans les cultures lentes. C'est un oscillateur de relaxation, cousin du générateur de rythme du chapitre~\ref{sec:muscles}: là, c'était la fatigue des groupes qui le rechargeait; ici, c'est la réserve de neurotransmetteur.

Les salves sont ce que fait le réseau quand il n'a rien à faire. Dans un cerveau vivant, on observe un tableau semblable pendant le sommeil profond (chapitre~\ref{sec:sleep}) et dans le cerveau en développement, avant que de vraies entrées n'y arrivent. La ressemblance est suggestive, mais l'identité des mécanismes reste une hypothèse.

\section{Enseigner sans professeur dopaminergique}

Puisque la culture ne contient pas de \nt{DOPA}, le signal \og mieux ou moins bien\fg{} doit être fourni par l'expérimentateur, au moyen des électrodes. Les expériences montrent que le réseau du banc modifie bel et bien ses réponses, et le plus intéressant est de savoir \emph{avec quoi} on l'enseigne.

\emph{Le professeur qui se tait.} On stimule le réseau par une électrode une fois toutes les une à trois secondes, jusqu'à ce que la réponse voulue apparaisse sur une autre électrode $50\pm10\ms$ après le stimulus. Dès qu'elle apparaît, on coupe la stimulation pendant cinq minutes. Après quelques cycles de ce genre, la réponse voulue suit le stimulus d'emblée~\cite{Shahaf2001}. La récompense, ici, c'est le silence: le stimulus secoue le réseau, et son arrêt laisse le réseau dans l'état qui a donné la bonne réponse. C'est la descente de gradient par perturbations aléatoires du chapitre~\ref{sec:dofa} (proposition~\ref{prop:gradient}), où le rôle de l'erreur est joué par le simple fait d'être secoué: on secoue jusqu'à ce que ce soit juste.

\emph{Le réseau qui joue au tennis.} Dans une expérience sur un banc à matrice d'électrodes dense, environ un million de neurones ont été branchés sur un jeu vidéo de tennis à une seule raquette~\cite{Kagan2022}. La position de la balle était injectée sous forme de courant dans des électrodes \og sensorielles\fg{}: le lieu de la balle par un code de place, sa distance par la fréquence. L'activité de deux zones \og motrices\fg{} déplaçait la raquette vers le haut ou vers le bas. La règle d'apprentissage était inhabituelle. Si la raquette renvoyait la balle, le réseau recevait un bref stimulus \emph{prévisible}; en cas de raté, plusieurs secondes de courant aléatoire \emph{imprévisible}. En vingt minutes de jeu, les séries de renvois s'allongeaient. Une hausse significative au cours d'une session n'apparaissait qu'avec la rétroaction complète; avec du silence à la place du stimulus, la culture jouait aussi mieux en fin de session que sans rétroaction, mais avec un effet plus faible, et il n'y avait pas d'apprentissage stable d'une session à l'autre sur plusieurs jours. L'analyse détaillée de cette expérience se trouve au chapitre~\ref{sec:dishbrain}.

Les auteurs ont expliqué le résultat par l'hypothèse que le réseau agit de façon à rendre son entrée prévisible~\cite{Friston2010}. C'est la même idée que l'économie d'impulsions du chapitre~\ref{sec:code} et la prédiction d'image du chapitre~\ref{sec:recognition}: la machine dépense moins quand l'entrée est attendue. Mais l'expérience elle-même, et surtout les propos des auteurs sur la \og sensibilité\fg{} de la culture, ont suscité de vives critiques: l'effet est faible et s'explique aussi plus simplement~\cite{Balci2023}. Le bilan honnête est le suivant: la culture du banc modifie son comportement sous l'effet de la rétroaction, c'est mesuré; qu'elle apprenne précisément à prédire son entrée reste une hypothèse.

\emph{Le réseau qui pilote un corps.} De façon analogue, on a relié une culture à un \og corps\fg{} virtuel simple et on lui a appris à se diriger vers une cible, en choisissant le motif spatio-temporel des stimulations d'entraînement~\cite{Bakkum2008}. Aucune de ces expériences ne fait intervenir la dopamine: le professeur, c'est le motif de courant que choisit l'ingénieur. Ce qui change quand le professeur devient un programme ou la dopamine elle-même est traité aux sections~\ref{sec:dishdopamine}--\ref{sec:cartpole}.

\begin{figure}[t]
\centering
\begin{tikzpicture}[>=Stealth,
  blk/.style={draw,very thick,rounded corners=2pt,align=center,font=\footnotesize,minimum height=0.8cm,fill=blue!5}]
% (a) closed loop
\node[font=\small] at (1.5,4.3) {(a) boucle fermée};
\node[blk,text width=2.6cm] (G) at (1.5,3.3) {jeu: balle et raquette};
\draw[very thick,fill=orange!10,rounded corners=3pt] (0.5,0.2) rectangle (2.5,1.6);
\foreach \x in {0.7,1.0,...,2.35}{\foreach \y in {0.4,0.7,1.0,1.3}{\fill[gray!60] (\x,\y) circle (0.04);}}
\draw[->,thick,blue!60!black] (G.west) -| (-0.5,0.9) -- (0.5,0.9);
\node[font=\scriptsize,blue!60!black,anchor=east,align=right] at (-0.6,2.1) {balle $\to$\\courant};
\draw[->,thick,red!70!black] (2.5,0.9) -- (3.5,0.9) |- (G.east);
\node[font=\scriptsize,red!70!black,anchor=west,align=left] at (3.6,2.1) {impulsions\\$\to$ raquette};
\node[font=\scriptsize] at (1.5,-0.2) {neurones sur matrice d'électrodes};
\node[font=\scriptsize,align=center] at (1.5,-0.85) {renvoi: stimulus prévisible\\raté: courant aléatoire};
% (b) reservoir
\begin{scope}[xshift=7.9cm]
\node[font=\small] at (1.6,4.3) {(b) réservoir};
\node[blk,text width=1.1cm] (X) at (-0.4,1.0) {entrée\\$u(t)$};
\draw[very thick,fill=orange!10] (1.6,1.0) circle (0.8);
\foreach \a/\r in {20/0.35,80/0.5,150/0.3,210/0.55,280/0.4,330/0.2,110/0.15}{\fill[red!60!black] ({1.6+\r*cos(\a)},{1.0+\r*sin(\a)}) circle (0.05);}
\node[font=\scriptsize] at (1.6,-0.2) {organoïde: sans professeur};
\node[blk,text width=1.8cm,fill=green!8] (W) at (4.0,1.0) {lecture\\$W_{\text{sort}}$};
\draw[->,thick] (X) -- (0.8,1.0);
\draw[->,thick] (2.4,1.0) -- node[above,font=\scriptsize] {$\mathbf r(t)$} (W);
\draw[->,thick] (W) -- (4.0,2.3) node[above,font=\scriptsize] {$y(t)$};
\node[font=\scriptsize,green!40!black] at (4.0,0.3) {apprentissage supervisé};
\end{scope}
\end{tikzpicture}
\caption{Deux façons de faire calculer un réseau vivant. (a) Boucle fermée: la position de la balle est injectée sous forme de courant, les impulsions des zones motrices déplacent la raquette, et un stimulus prévisible ou aléatoire après le coup sert de professeur. (b) Réservoir~\eqref{eq:reservoir}: l'organoïde ne reçoit aucun signal d'apprentissage; il décompose l'entrée en une multitude de réponses non linéaires dotées de mémoire, et seule une lecture linéaire externe apprend à partir de l'erreur. Les connexions de l'organoïde lui-même changent aussi, sans professeur.}
\label{fig:livingmachine}
\end{figure}

\section{Le réservoir vivant}

Il existe une autre façon d'utiliser un réseau vivant: ne pas l'entraîner du tout. En technique, on l'appelle le \emph{calcul par réservoir}~\cite{Maass2002,JaegerHaas2004}. On prend un système non linéaire à mémoire tout fait, n'importe lequel, pourvu que sa réponse à l'entrée soit riche et dépende du passé récent. L'entrée $u(t)$ le met en mouvement, on obtient un vecteur de réponses $\mathbf r(t)$, et seule la lecture linéaire est entraînée:
\begin{empheq}[box=\keybox]{equation}
y(t) \;=\; W_{\text{sort}}\,\mathbf r(t) ,
\label{eq:reservoir}
\end{empheq}
dont les poids sont ajustés par la règle des moindres carrés~\eqref{eq:lms} du chapitre~\ref{sec:motorlearning}. Le réservoir fait ce que faisaient les petits neurones \nt{GLUT} du chapitre~\ref{sec:motorlearning}: il déploie l'entrée en un long vecteur, dans lequel les classes voulues deviennent séparables par un plan (chapitre~\ref{sec:dendrites}).

Un organoïde sur matrice d'électrodes s'est révélé apte à jouer ce rôle de réservoir. Des sons de parole, injectés sous forme de motifs de courant, ont été distingués selon le locuteur avec une précision d'environ $78\%$ après l'entraînement d'une seule lecture linéaire~\cite{Cai2023}. Ce chiffre de $78\%$ est celui qu'annoncent les auteurs et que reprend la presse. Le résultat est utile, mais il faut comprendre où se trouve l'\og intelligence\fg{}: l'organoïde apporte la non-linéarité et une mémoire évanescente, l'apprentissage par l'erreur se fait à l'extérieur, dans le silicium (fig.~\ref{fig:livingmachine}). L'organoïde ne reste pas pour autant inchangé: selon les auteurs, sa propre connectivité fonctionnelle a évolué pendant l'entraînement, ce qu'ils appellent un apprentissage non supervisé dans l'organoïde lui-même~\cite{Cai2023}. La part de la précision due à ce changement et celle due à la lecture n'ont pas été séparées.

\section{Ce que le banc nous apprend sur la machine}

\emph{L'énergie.} D'après l'estimation~\eqref{eq:energy}, une impulsion coûte environ $10^{-10}$~J. Une culture d'un million de neurones, à une impulsion par seconde, dépense pour la transmission des signaux
\begin{empheq}[box=\keybox]{equation}
P \;\approx\; 10^{6}\times1\,\text{s}^{-1}\times10^{-10}\,\text{J} \;=\; 10^{-4}\,\text{W},
\label{eq:dishpower}
\end{empheq}
soit un dixième de milliwatt. L'électronique qui stimule, enregistre et traite ces impulsions dépense des ordres de grandeur de plus. Comme au chapitre~\ref{sec:debugging}, le goulot d'étranglement n'est pas le calcul, mais la liaison avec lui.

\emph{Les pièces manquantes.} Le banc montre ce que sait faire un bus \nt{GLUT} doté d'un contrôle de flux \nt{GABA}, sans registres de contrôle: s'auto-organiser, produire des salves, modifier ses réponses sous l'effet de la rétroaction et servir de bon réservoir. Ce qu'il ne sait pas faire sans eux, c'est décider lui-même de ce qui compte comme un succès. Sur le banc, ce signal est fourni par l'ingénieur; dans le cerveau, comme nous l'avons vu aux chapitres~\ref{sec:dofa}--\ref{sec:acet}, par les canaux de diffusion.

\emph{L'éthique.} Les organoïdes sont cultivés à partir de cellules humaines, et plus ils deviennent complexes, plus la question se pose sérieusement: un tel tissu peut-il ressentir quelque chose? Le regard de l'ingénieur suggère une réponse partielle: la douleur du chapitre~\ref{sec:pain} n'est pas le signal d'un capteur isolé, mais tout un système d'alarme, avec un portillon, un bouton de volume et des canaux de diffusion, qui n'existent pas sur le banc. Mais c'est un raisonnement, non une preuve, et le domaine lui-même propose de mener l'évaluation éthique en même temps que les expériences, dès le début~\cite{Smirnova2023}.

\begin{keyblock}
\begin{proposition}[Un réseau vivant sur le banc]
\label{prop:livingmachine}
Un réseau de neurones \nt{GLUT} et \nt{GABA} sur une matrice d'électrodes produit spontanément des salves~\eqref{eq:burstinterval}, modifie ses réponses sous l'effet de la rétroaction et peut servir de réservoir~\eqref{eq:reservoir} pour une lecture entraînée à l'extérieur. Tout cela est mesuré. Qu'un tel réseau apprenne selon une règle générale, par exemple prédire son entrée, reste une hypothèse, et la plupart des spécialistes rejettent les grands mots sur sa \og sensibilité\fg{}.
\end{proposition}
\end{keyblock}

\section{De la dopamine libérée par un éclair de lumière}
\label{sec:dishdopamine}

La seule expérience directe que nous connaissions où l'on a donné à une culture la dopamine pour professeur a été menée sur une plateforme d'organoïdes à laquelle les chercheurs se connectent à distance~\cite{Jordan2024}. Dans le liquide qui baigne les organoïdes, on a ajouté de la dopamine enfermée dans une \og cage\fg{} photosensible: la molécule liée n'agit pas sur les récepteurs. La concentration de dopamine en cage est de $1$~mM. Quand il faut récompenser le réseau, on l'éclaire aux ultraviolets: la cage se rompt, et la dopamine se répand dans le volume.

La boucle fonctionne ainsi. Le même stimulus est appliqué encore et encore; s'il a provoqué plus d'impulsions qu'auparavant, un éclair se déclenche et libère la dopamine. Sur $240$ répétitions, le nombre d'impulsions évoquées a globalement augmenté. Les auteurs écrivent eux-mêmes qu'une seule expérience de ce type ne permet pas de conclure que cette hausse est due à la dopamine~\cite{Jordan2024}.

Pour un ingénieur, c'est une première tentative de monter dans une boîte de culture la règle à trois facteurs~\eqref{eq:threefactor}: l'activité de l'émetteur et du destinataire laisse une trace, et un signal commun décide s'il faut consolider le changement. Mais le signal ne peut ici être que positif: il y a récompense ou non, et le montage ne produit aucune erreur négative $\delta<0$. De plus, la dopamine diffuse et n'est évacuée que lentement, comme la concentration de~\eqref{eq:lowpass}, si bien que des éclairs fréquents peuvent simplement élever son niveau global. Pour distinguer l'apprentissage de cet effet, il faut deux contrôles: autant d'éclairs à des instants aléatoires, sans lien avec la réponse du réseau, et les mêmes éclairs sans dopamine en cage. Il faut aussi un test après repos: le changement subsiste-t-il une fois la dopamine partie?

\section{La dopamine enseigne au silicium}
\label{sec:biohybrid}

Dans l'expérience inverse, des cellules vivantes enseignent à l'électronique~\cite{Keene2020}. Des cellules qui sécrètent de la dopamine sont cultivées dans un microcanal au-dessus d'un composant neuromorphique organique, un transistor dont la conductance du canal joue le rôle du poids synaptique. La dopamine libérée par les cellules s'oxyde sur l'électrode du composant, ce qui modifie la conductance. Le dispositif reproduit aussi le mécanisme d'élimination du neurotransmetteur de la fente, si bien que le poids peut non seulement être modifié durablement, mais aussi être ramené à sa valeur initiale.

En termes de circuit, c'est un intégrateur à fuite à entrée chimique: le poids accumule le flux de dopamine, et l'\og évacuation\fg{} fixe la vitesse de l'oubli; c'est le même circuit RC que~\eqref{eq:lowpass}. L'essentiel est ici le sens de la liaison. La sonde du chapitre~\ref{sec:debugging} ne voit pas les registres de diffusion, parce qu'ils sont chimiques. Ce composant lit justement un registre chimique et le convertit en poids électrique. Pour les interfaces cerveau-ordinateur, c'est la voie vers un capteur qui entend non seulement le bus de données, mais aussi les registres de contrôle.

\section{Des organoïdes dotés de leurs propres neurones dopaminergiques}
\label{sec:midbrainorg}

On sait cultiver, à partir de cellules souches humaines, des organoïdes qui ressemblent à la région du cerveau où se trouvent les neurones \nt{DOPA}. Ils contiennent des neurones dopaminergiques fonctionnels, qui sécrètent de la dopamine~\cite{Jo2016}. On les cultive surtout pour étudier des maladies. Nous n'avons trouvé aucune expérience où leur dopamine aurait servi de professeur à un autre réseau.

Pour la machine faite de neurones vivants, c'est la pièce manquante. Le banc du chapitre~\ref{sec:livingmachine}, c'est un bus de données et un contrôle de flux sans registres de contrôle; ici, la source du signal de diffusion existe déjà. Il manque le critique: un réseau qui calculerait l'erreur de prédiction~\eqref{eq:ctd} et piloterait ces neurones. Relier un organoïde de cortex à un tel organoïde de façon que la dopamine soit libérée en fonction de l'erreur est un problème ouvert; pour l'instant, c'est une hypothèse et non une expérience.

\section{Un organoïde tient un pendule en équilibre sans dopamine}
\label{sec:cartpole}

La plus complète des expériences récentes se passe entièrement de dopamine~\cite{Robbins2026}. Des organoïdes de cortex de souris ont été placés en boucle fermée sur la tâche du \og pendule inversé sur chariot\fg{}: l'activité de l'organoïde déplace le chariot, et la position du pendule lui revient sous forme de stimulation. Entre les essais, l'organoïde reçoit de brefs stimuli d'apprentissage à haute fréquence. Lesquels exactement, c'est un programme d'apprentissage par renforcement qui le choisit.

Les résultats sont mesurés:
\begin{itemize}
\item pour la plupart des organoïdes, les signaux d'apprentissage choisis par le programme ont donné de meilleurs résultats que des signaux choisis au hasard ou que l'absence de signal;
\item après $45$ minutes de repos, l'amélioration ne subsistait pas;
\item le blocage des deux types de récepteurs du glutamate, AMPA et NMDA, supprimait l'amélioration.
\end{itemize}

Le dernier point indique où réside ce qui est appris: dans les synapses \nt{GLUT}, et via le récepteur qui, à la section~\ref{sec:weightstore}, fonctionne comme détecteur de coïncidences entre l'entrée et la sortie. Le deuxième point montre ce qui manque: le changement rapide existe, la consolidation non, comme chez l'élève rapide du chapitre~\ref{sec:motorlearning} privé de l'élève lent. Quant au rôle du professeur, il a changé: l'ingénieur qui choisit le motif de courant a été remplacé par un programme, mais le professeur reste à l'extérieur de la boîte.

Parmi les expériences antérieures d'apprentissage de cultures sur matrices d'électrodes, nous n'en avons trouvé aucune qui utilise la dopamine: le signal d'apprentissage était partout une stimulation électrique.

\section{Qui enseigne à la culture}

\begin{table}[t]
\centering
\footnotesize
\setlength{\tabcolsep}{4pt}
\renewcommand{\arraystretch}{1.15}
\begin{tabular}{>{\raggedright}p{2.7cm}>{\raggedright}p{2.5cm}>{\raggedright}p{3.2cm}>{\raggedright\arraybackslash}p{4.4cm}}
\toprule
Expérience & Qui enseigne & Signal d'apprentissage & Ce que l'on sait \\
\midrule
arrêt de la stimulation à la bonne réponse & l'ingénieur & fin de la stimulation & la réponse se fixe (mesuré) \\
DishBrain (chap.~\ref{sec:dishbrain}) & l'ingénieur & courant prévisible ou aléatoire & hausse significative des échanges en une session seulement avec rétroaction complète, effet moindre avec le silence (mesuré); instable d'une session à l'autre; cause débattue \\
pendule sur chariot & programme d'apprentissage par renforcement & stimuli à haute fréquence choisis par le programme & mieux que le hasard, mais ne subsiste pas après repos (mesuré) \\
dopamine par éclair & règle \og plus d'impulsions = récompense\fg{} & dopamine libérée par la lumière & hausse des impulsions observée; rôle de la dopamine non démontré \\
synapse biohybride & cellules vivantes & dopamine des cellules & changement durable et retour du poids du composant (mesuré) \\
organoïdes à neurones \nt{DOPA} & --- & --- & source de dopamine présente; pas utilisée comme professeur \\
\bottomrule
\end{tabular}
\caption{Qui enseigne aux neurones vivants sur puce, et avec quoi.}
\label{tab:dishteachers}
\end{table}

Dans toutes les expériences où c'est la culture elle-même qui apprend, le professeur reste pour l'instant à l'extérieur (tableau~\ref{tab:dishteachers}): un ingénieur, un programme ou une règle fixée par l'ingénieur. La dopamine n'est entrée dans la boîte qu'une seule fois, et son rôle reste à démontrer. Monter dans la boîte un véritable canal de diffusion, avec un critique, une erreur et une trace comme au chapitre~\ref{sec:dofa}, est la prochaine étape, que personne n'a encore franchie.


\section{Bilan}

\begin{table}[t]
\centering
\footnotesize
\setlength{\tabcolsep}{4pt}
\renewcommand{\arraystretch}{1.15}
\begin{tabular}{>{\raggedright}p{2.8cm}>{\raggedright}p{4.2cm}>{\raggedright\arraybackslash}p{5.8cm}}
\toprule
Expérience & Ce que fait le réseau vivant & Ce que l'on sait \\
\midrule
réseau sans entrée & salves espacées d'une seconde à plusieurs minutes~\eqref{eq:burstinterval} & mesuré; le mécanisme par épuisement synaptique est le modèle admis \\
professeur qui se tait & la réponse voulue se fixe quand on coupe la stimulation & mesuré \\
partie de tennis & les séries de renvois s'allongent; hausse significative en une session seulement avec rétroaction complète & changement de comportement mesuré; apprentissage d'une session à l'autre instable; taille de l'effet et explication par la prédiction de l'entrée contestées \\
corps virtuel & déplacement vers la cible avec un motif de stimuli bien choisi & mesuré \\
réservoir & non-linéarité et mémoire pour la lecture~\eqref{eq:reservoir} & mesuré; apprentissage par l'erreur à l'extérieur, dans le silicium; les connexions de l'organoïde changent aussi \\
énergie & environ $10^{-4}$~W par million de neurones~\eqref{eq:dishpower} & estimation d'après~\eqref{eq:energy} \\
\bottomrule
\end{tabular}
\caption{Ce qu'ont montré les machines faites de neurones vivants.}
\label{tab:livingmachine}
\end{table}

La machine faite de neurones vivants (tableau~\ref{tab:livingmachine}) est un banc d'essai qui montre ce que savent faire un bus de données et un contrôle de flux sans registres de contrôle. Ils savent faire beaucoup, mais pas décider eux-mêmes de ce qui est bien. Sur le banc, ce rôle revient à l'ingénieur et à son motif de courant; dans le cerveau, aux quelques fils de diffusion auxquels est consacré le milieu du livre.

\section{Exercices}

\begin{exercise}[\lvlA]
\label{ex:21:1}
\emph{Intervalle entre salves.} Une salve vide la réserve jusqu'à zéro, puis $\tau_{\text{rec}}\,\dd x/\dd t=1-x$, et une nouvelle salve part à $x=x^*$; $\tau_{\text{rec}}=0.8$~s. (a)~Trouvez l'intervalle pour $x^*=0.9$ et $0.99$. (b)~Le seuil $x^*$ fluctue aléatoirement de $\pm0.005$ autour de $0.99$. Dans quelle plage tombent les intervalles?
\end{exercise}

\begin{exercise}[\lvlA]
\label{ex:21:2}
\emph{L'énergie du banc.} (a)~Quelle puissance l'estimation~\eqref{eq:dishpower} donne-t-elle pour le cerveau entier ($8.6\cdot10^{10}$ neurones)? (b)~Le banc enregistre $1024$ électrodes à $20$~kHz, sur $10$~bits par échantillon. À $1$~nJ par bit, combien de fois la transmission de ce flux coûte-t-elle plus cher que la puissance de la culture elle-même?
\end{exercise}

\begin{exercise}[\lvlB]
\label{ex:21:3}
\emph{Conception: supprimer les salves.} Une stimulation clairsemée par de nombreuses électrodes fait libérer à chaque synapse du neurotransmetteur à la fréquence $r_b$; la réserve suit $\tau_{\text{rec}}\,\dd x/\dd t=1-x-Ur_b\tau_{\text{rec}}x$, $U=0.5$, $\tau_{\text{rec}}=0.8$~s. Quelle est la plus petite $r_b$ pour laquelle la réserve n'atteint pas $x^*=0.9$; $0.99$, et de combien la synapse s'affaiblit-elle alors?
\end{exercise}

\begin{exercise}[\lvlB]
\label{ex:21:4}
\emph{La mémoire d'un réservoir ne dépasse pas son nombre de neurones.} Un réservoir à $N$ sorties $\mathbf r(t)$ reçoit une entrée blanche $u(t)$ de moyenne nulle et de variance unité; $m_k$ est le plus grand carré de la corrélation d'une lecture $\mathbf w_k\cdot\mathbf r(t)$ avec $u(t-k)$. Démontrez que la capacité mémoire $\mathrm{MC}=\sum_{k\ge0}m_k\le N$.
\end{exercise}

\begin{exercise}[\lvlB]
\label{ex:21:5}
\emph{Simulation: un réservoir en silicium.} Réservoir $\mathbf r(t+1)=\tanh(W\mathbf r(t)+\mathbf v\,u(t)+\mathbf c)$, $N=30$, $W$ aléatoire de rayon spectral $0.9$, $v_i\sim U(-0.5,0.5)$, $c_i\sim U(-0.2,0.2)$, $u\sim U(-1,1)$, lecture par régression ridge. Trouvez la capacité mémoire de l'exercice~\ref{ex:21:4} avec et sans $\tanh$. Quel réservoir retient le plus?
\end{exercise}

\begin{exercise}[\lvlB]
\label{ex:21:6}
\emph{Une lecture pour un réservoir vivant.} Un organoïde fournit $1024$ canaux et $P=240$ enregistrements d'entraînement de deux classes. (a)~Combien de caractéristiques $n$ garder pour qu'un étiquetage aléatoire soit séparable avec une probabilité d'au plus $1\%$ (exercice~\ref{ex:19:3})? (b)~De combien de sigmas $78\%$ sur $100$ enregistrements de test dépassent-ils le hasard?
\end{exercise}

\begin{exercise}[\lvlC]
\label{ex:21:7}
\emph{Recherche: un professeur sans canal de diffusion.} Proposez un protocole de stimulation qui réalise sur le banc la règle à trois facteurs du chapitre~\ref{sec:dofa} par $8$ à $1024$ électrodes, ainsi qu'un contrôle à lecture gelée qui distingue l'apprentissage des poids d'un simple décalage de l'état du réseau. Quel résultat réfuterait l'hypothèse d'apprentissage?
\end{exercise}

\chapter{DishBrain}
\chaptermark{DishBrain}
\label{sec:dishbrain}
\begin{chaptergoals}
\item comment DishBrain code la balle par lieu et fréquence, et lit la raquette en différence de comptes;
\item pourquoi un comptage sur $10\ms$ est bruité, si bien que le réseau ne peut déplacer que la moyenne;
\item comment secouer le réseau après les seuls ratés le fait dériver vers les réglages qui ratent moins;
\item ce qui a été mesuré, ce qui est débattu, et comment reconstruire le banc pour tester le mécanisme.
\end{chaptergoals}

\section{Un réseau en boîte de culture joue au tennis}

DishBrain, c'est le nom que les auteurs ont donné au banc sur lequel une culture de neurones vivants, cultivée sur une puce à électrodes, joue à un tennis simplifié à une seule raquette~\cite{Kagan2022}. C'est l'expérience la plus débattue sur une machine faite de neurones vivants (ces machines sont passées en revue au chapitre~\ref{sec:livingmachine}), et elle est particulièrement commode pour notre livre. Ici, le petit réseau est accessible en entier: chaque entrée est fixée par l'expérimentateur, chaque sortie est enregistrée. Toutes les questions des chapitres précédents (comment le signal est codé, qui enseigne, ce que voit la sonde) peuvent être posées à un réseau qui n'a ni yeux, ni muscles, ni neurones \nt{DOPA}. Analysons l'expérience comme un ingénieur analyse le banc d'un autre: schéma, entrée, sortie, professeur, résultat, controverses.

\section{Le banc}

Les neurones proviennent du cortex d'embryons de souris, environ $800\,000$ cellules par puce, ou sont obtenus à partir de cellules souches humaines, environ un million par puce. Pendant quelques semaines, ils poussent sur la puce, s'interconnectent et commencent à émettre spontanément des impulsions et des bouffées. Sous eux, sur une surface de $8$~mm$^2$, se trouvent $26\,000$ électrodes de platine; on peut en enregistrer jusqu'à $1024$ simultanément, et en pratique on stimule par huit d'entre elles.

Une boucle est refermée autour de la culture (fig.~\ref{fig:dishloop}). Un ordinateur simule le jeu: la balle vole, rebondit sur les murs, la raquette monte et descend. La position de la balle est convertie en courant sur huit électrodes \og sensorielles\fg{}. Les impulsions de deux zones \og motrices\fg{} sont comptées par fenêtres de $10\ms$ et déplacent la raquette. Après chaque renvoi ou raté, la culture reçoit un stimulus supplémentaire: la rétroaction. Les fenêtres de $10\ms$ sont le cycle de l'ordinateur du banc, et non celui de la culture: le réseau lui-même fonctionne toujours de façon asynchrone, comme tous les réseaux de ce livre.

\begin{figure}[t]
\centering
\begin{tikzpicture}[>=Stealth,
  blk/.style={draw,very thick,rounded corners=3pt,align=center,font=\small,minimum height=1.2cm},
  lab/.style={font=\scriptsize,align=center}]
\node[blk,fill=blue!6,text width=3.4cm] (C) at (0,0) {culture sur puce\\\scriptsize $\sim10^6$ neurones, $26\,000$ électrodes};
\node[blk,fill=orange!10,text width=3.0cm] (G) at (7.2,0) {ordinateur:\\partie de tennis};
\draw[->,very thick,blue!60!black] (G.north west) to[bend right=25] node[lab,above] {balle: \emph{lieu}, laquelle des 8 électrodes;\\\emph{fréquence}, 4--40 imp./s} (C.north east);
\draw[->,very thick,red!70!black] (C.south east) to[bend right=25] node[lab,below] {raquette: impulsions des zones\\\og haut\fg{} et \og bas\fg{} comptées sur $10\ms$} (G.south west);
\draw[->,thick,densely dashed,green!45!black] (G.east) -- ++(0.9,0) |- ++(0,-2.6) -| node[lab,pos=0.25,below] {renvoi: stimulus prévisible; raté: imprévisible} (C.south);
\end{tikzpicture}
\caption{La boucle DishBrain. Flèche bleue: l'entrée; la position de la balle est codée par un code de place et un code de fréquence (chapitre~\ref{sec:code}). Flèche rouge: la sortie; la différence des comptes des deux zones déplace la raquette. Flèche verte en tirets: la rétroaction après chaque échange, seul \og professeur\fg{} du banc.}
\label{fig:dishloop}
\end{figure}

\section{Entrée et sortie}

L'\emph{entrée} est codée à la fois par deux codes du chapitre~\ref{sec:code}. L'électrode stimulée, parmi les huit, indique à quelle hauteur se trouve la balle: c'est un code de place, comme celui des capteurs du chapitre~\ref{sec:sensors}. La cadence des impulsions de stimulation indique à quelle distance est la balle: $4$ par seconde quand elle est près du mur du fond, et davantage, linéairement, jusqu'à $40$ par seconde quand elle est près du mur de la raquette. C'est un code de fréquence; l'amplitude des impulsions de la balle est de $75$~mV. La culture ne \og connaît\fg{} d'avance aucun de ces codes: le code est fixé par l'expérimentateur, et le réseau doit apprendre à le lire.

La \emph{sortie} se lit comme dans les interfaces du chapitre~\ref{sec:debugging}. Les impulsions d'une zone motrice poussent la raquette vers le haut, celles de l'autre vers le bas; dans l'écriture la plus simple, pour chaque fenêtre,
\begin{empheq}[box=\keybox]{equation}
\Delta p \;=\; c\,\bigl(n_\uparrow - n_\downarrow\bigr),
\label{eq:dishreadout}
\end{empheq}
où $n_\uparrow$, $n_\downarrow$ sont les nombres d'impulsions des deux zones en $10\ms$, et $c$ un coefficient du banc. C'est le code à deux fils du chapitre~\ref{sec:glutcomputer}: le sens du mouvement n'est pas porté par le signe du signal, mais par celui des deux fils qui parle le plus fort.

Examinons le bruit d'une telle sortie. Si une zone émet au total $R$ impulsions par seconde, une fenêtre $T=10\ms$ en contient en moyenne $RT$, et d'après~\eqref{eq:rateerror} la dispersion relative vaut $1/\sqrt{RT}$. Pour $R=1000$ impulsions par seconde, cela fait environ $30$ pour cent dans chaque fenêtre; pour $R=100$, plus de cent. La raquette tremble inévitablement, et le réseau ne peut apprendre que d'une seule façon: en décalant la différence \emph{moyenne} des comptes dans le bon sens. Ce sont les mêmes lois qui limitent toute interface du chapitre~\ref{sec:debugging}, mais cette fois des deux côtés de la boucle.

\section{Un professeur sans dopamine}

Le plus intéressant, sur ce banc, c'est le professeur. Une culture de neurones du cortex ne contient pas de neurones \nt{DOPA}, et le banc n'envoie aucun signal \og mieux que prévu\fg{}. La rétroaction est d'une autre nature. Si le réseau renvoie la balle, un bref stimulus \emph{prévisible} est appliqué simultanément aux huit électrodes sensorielles: des impulsions de $75$~mV, $100$ par seconde pendant $0.1\,\text{s}$. S'il la rate, un stimulus que les auteurs qualifient d'\emph{imprévisible} dure quatre secondes: des impulsions de $150$~mV, $5$ par seconde, à des instants aléatoires et sur des électrodes tirées au hasard; suivent quatre secondes de pause sans stimulation, puis la balle repart~\cite{Kagan2022}.

Les auteurs partaient de l'hypothèse que le réseau agit de façon à rendre son entrée prévisible~\cite{Friston2010}. Pour un ingénieur, le sens est simple: la culture n'a qu'un moyen d'échapper à quatre secondes de bruit aléatoire, renvoyer la balle. Nous avons déjà rencontré la même idée dans l'économie d'impulsions du chapitre~\ref{sec:code} et dans la prédiction d'image du chapitre~\ref{sec:recognition}.

Mais la prévisibilité est-elle vraiment nécessaire ici? Un mécanisme plus grossier est possible. Supposons que le stimulus imprévisible après un raté se contente de \emph{brasser} les changements récents du réseau, tandis que le stimulus calme après un renvoi n'y touche pas. Décrivons le réglage du réseau par un vecteur $\boldsymbol\psi$, et soit $P_{\text{raté}}(\boldsymbol\psi)$ la probabilité de rater avec le réglage $\boldsymbol\psi$. Si les secousses sont symétriques (passer de $\boldsymbol\psi$ à $\boldsymbol\psi'$ est aussi probable que l'inverse), alors en régime établi le flux de départs de chaque réglage est équilibré par le flux d'arrivées, et la fraction du temps que le réseau passe dans le réglage $\boldsymbol\psi$ vaut
\begin{empheq}[box=\keybox]{equation}
\rho(\boldsymbol\psi) \;\propto\; \frac{1}{P_{\text{raté}}(\boldsymbol\psi)} .
\label{eq:dishdensity}
\end{empheq}
Un réglage qui rate dix fois moins souvent est conservé dix fois plus longtemps. Personne n'indique au réseau dans quel sens modifier ses poids: les bons réglages sont simplement détruits moins souvent.

Nous l'avons vérifié sur un modèle qui réduit la culture à deux nombres: la raquette arrive au point $a\,y+b$ quand la balle arrive à la hauteur $y$, et elle renvoie si l'écart est inférieur à $0.1$. Après un raté, les deux nombres reçoivent une secousse aléatoire; après un renvoi, ils restent inchangés. Sur $2000$ échanges, la proportion de balles renvoyées est passée de $0.21$ sur les deux cents premiers à $0.38$ sur les cinq cents derniers (quarante \og cultures\fg{}, fig.~\ref{fig:dishmodel}). Si l'on applique la secousse après \emph{chaque} échange, la rétroaction ne porte aucune information, et la proportion reste vers $0.12$; sans aucune secousse, elle reste à $0.19$. L'expérience sur le banc concorde: une hausse significative au cours d'une session n'apparaissait qu'avec la rétroaction complète. Avec du silence à la place des stimuli, la culture jouait tout de même mieux en fin de session que sans rétroaction, mais avec un effet moindre~\cite{Kagan2022}; notons que, dans cette condition aussi, un raté est suivi d'une longue pause et un renvoi d'une pause courte, de sorte que la différence entre les issues ne disparaît pas entièrement.

\begin{figure}[t]
\centering
\begin{tikzpicture}[x=1cm,y=5cm]
\draw[->,gray!70!black] (0,0) -- (10.6,0);
\draw[->,gray!70!black] (0,0) -- (0,0.5) node[above,font=\small,black] {proportion de renvois};
\foreach \v in {0.1,0.2,0.3,0.4}{\draw[gray!60!black] (-0.06,\v) -- (0.06,\v); \node[left,font=\scriptsize] at (0,\v) {\v};}
\foreach \x/\f/\l/\lab in {1.8/0.21/0.38/{bruit après\\un raté},5.2/0.13/0.11/{bruit après\\chaque échange},8.6/0.19/0.19/{sans bruit}}{
  \fill[blue!35] (\x-0.8,0) rectangle (\x-0.05,\f);
  \fill[red!60!black] (\x+0.05,0) rectangle (\x+0.8,\l);
  \node[below,font=\scriptsize,align=center] at (\x,-0.01) {\lab};
  \node[above,font=\scriptsize] at (\x-0.425,\f) {\f};
  \node[above,font=\scriptsize] at (\x+0.425,\l) {\l};}
\fill[blue!35] (7.4,0.47) rectangle (7.7,0.495); \node[right,font=\scriptsize] at (7.75,0.482) {200 premiers};
\fill[red!60!black] (7.4,0.41) rectangle (7.7,0.435); \node[right,font=\scriptsize] at (7.75,0.422) {500 derniers};
\end{tikzpicture}
\caption{Modèle du \og brassage après un raté\fg{} (\texttt{dishbrain\_model.py}): proportion de balles renvoyées au début et à la fin de $2000$ échanges, moyenne sur quarante cultures modèles. L'apprentissage n'a lieu que si la secousse suit un raté et non un renvoi, comme dans l'expérience, où la hausse significative en une session n'apparaissait qu'avec la rétroaction complète.}
\label{fig:dishmodel}
\end{figure}

\begin{keyblock}
\begin{proposition}[Le professeur brasseur]
\label{prop:dishscramble}
Si l'échec secoue le réseau et que le succès le laisse en paix, le réseau dérive de lui-même vers les réglages qui se trompent moins souvent: le temps passé dans un réglage est inversement proportionnel à la probabilité d'échec~\eqref{eq:dishdensity}. Un tel apprentissage n'exige ni signal de récompense, ni signe de ce signal, ni prévisibilité en tant que telle: il suffit que la rétroaction après un succès diffère de celle qui suit un échec. Le changement de comportement de la culture est mesuré; le mécanisme à l'œuvre dans la boîte (prédiction de l'entrée ou brassage) n'est pas établi.
\end{proposition}
\end{keyblock}

Comparez avec le chapitre~\ref{sec:dofa}. Là aussi, le bruit servait à la recherche, mais la dopamine avait un signe: l'erreur $\delta$ indiquait dans quel sens déplacer les poids, et le pas suivait le gradient~\eqref{eq:perturbation}. Ici, il n'y a pas de signe, seulement \og garder\fg{} ou \og secouer\fg{}. Une telle recherche est plus grossière et plus lente, mais elle exige moins du réseau: ni neurones \nt{DOPA}, ni critique, ni trace étalée sur des secondes.

\section{Ce qui a été mesuré et ce qui est débattu}

Chaque partie durait $20$ minutes; on comparait les $5$ premières minutes aux $15$ dernières. Chez les cultures à rétroaction complète, les séries de renvois se sont allongées et les ratés immédiats sont devenus plus rares; dans les contrôles sans cellules, sans jeu ou avec une \og raquette\fg{} pilotée par ordinateur, rien de tel. Les cultures de cellules humaines tenaient en moyenne plus longtemps que celles de souris. L'amélioration est statistiquement robuste (sur neuf cultures de souris et onze cultures humaines, dans une bonne centaine de parties par groupe), mais modeste: le réseau n'est pas devenu un bon joueur, il est devenu un peu meilleur que le hasard. Elle est robuste au sein d'une session; les auteurs n'ont pas obtenu d'apprentissage stable d'une session à l'autre sur plusieurs jours~\cite{Kagan2022}. Un travail ultérieur du même groupe a montré que, pendant le jeu, l'activité de la culture se rapproche du régime \emph{critique}, la frontière entre extinction et avalanche, cette même vie au bord du gouffre qu'au chapitre~\ref{sec:gaba}, et que plus elle s'en approche, meilleur est le jeu~\cite{Habibollahi2023}.

La principale controverse n'a pas porté sur les chiffres, mais sur les mots. Le titre de l'article affirmait que les neurones en boîte \og font preuve de sensibilité\fg{} (sentience), et un groupe de chercheurs a objecté qu'un changement de comportement en boucle fermée n'est encore ni intelligence ni sensation, et que les conclusions doivent être formulées plus modestement~\cite{Balci2023}. Du point de vue de l'ingénieur, deux questions restent ouvertes, et toutes deux portent sur le mécanisme. Première question: le réseau apprend-il, c'est-à-dire ses poids changent-ils durablement, ou la rétroaction ne fait-elle que décaler son état courant? Seconde question: la prévisibilité est-elle réellement nécessaire, ou toute différence entre la réponse au succès et la réponse à l'échec suffit-elle, comme dans la proposition~\ref{prop:dishscramble}? Un banc où chaque électrode est accessible permet de répondre aux deux, et c'est sans doute là sa principale valeur.

\section{Spécification du banc: comment le reproduire}
\label{sec:dishspec}

On trouvera ci-dessous tout ce qu'il faut pour reconstruire un tel banc: matériel, boucle, codage de l'entrée, lecture de la sortie, rétroaction et protocole expérimental. Chaque paramètre est marqué de sa source. \og Article\fg{}: la valeur est tirée du texte de l'étude~\cite{Kagan2022}. \og Choix\fg{}: l'article ne la donne pas, et il faudra la fixer soi-même pour reproduire l'expérience; nous proposons une valeur par défaut et indiquons comment la vérifier.

\subsection*{Matériel et culture}

\begin{center}
\footnotesize
\setlength{\tabcolsep}{4pt}
\renewcommand{\arraystretch}{1.15}
\begin{tabular}{>{\raggedright}p{4.0cm}>{\raggedright}p{6.9cm}>{\raggedright\arraybackslash}p{1.6cm}}
\toprule
Paramètre & Valeur & Source \\
\midrule
puce & matrice MaxOne: $26\,000$ électrodes de platine sur une surface de $8$~mm$^2$ & article \\
enregistrement & jusqu'à $1024$ canaux simultanés, $20\,000$ échantillons par seconde & article \\
stimulation & jusqu'à $32$ électrodes techniquement, $8$ indépendantes dans l'expérience & article \\
cellules & cortex d'embryon de souris; neurones humains issus de cellules souches (deux méthodes de culture); cellules HEK293T (non neuronales) comme contrôle & article \\
ensemencement & environ $10^6$ cellules par puce & article \\
âge de la culture lors de l'expérience & souris: à partir de $\sim10$ jours; humain: $14$--$24$ jours ou $\sim80$ jours selon la méthode de culture & article \\
\bottomrule
\end{tabular}
\end{center}

\subsection*{Boucle et temps}

La boucle fonctionne par fenêtres de $10\ms$ ($200$ échantillons): pendant chaque fenêtre, on compte les impulsions sur les électrodes des zones motrices, le jeu est mis à jour et la stimulation suivante est émise. Le délai entre une impulsion dans la culture et le stimulus de réponse est d'environ $5\ms$. Pendant le jeu, le signal de chaque électrode passe par un filtre passe-haut de Bessel d'ordre 2 de fréquence de coupure $100$~Hz; la valeur absolue du signal est lissée par un filtre passe-bas de Bessel d'ordre 1 à $1$~Hz, et le seuil de détection des impulsions est proportionnel à cette valeur absolue lissée. Le seuil de six écarts types du fond, l'article l'indique pour certaines mesures d'activité, et non pour le jeu. Pour que la stimulation ne soit pas comptée comme une réponse de la culture, tous les signaux sont masqués lorsque plus de $15$ grandes impulsions, au-dessus de $75$~mV, arrivent en même temps. Tout ceci: article.

\subsection*{Codage de la balle}

\begin{center}
\footnotesize
\setlength{\tabcolsep}{4pt}
\renewcommand{\arraystretch}{1.15}
\begin{tabular}{>{\raggedright}p{3.4cm}>{\raggedright}p{7.5cm}>{\raggedright\arraybackslash}p{1.6cm}}
\toprule
Paramètre & Valeur & Source \\
\midrule
zone sensorielle & $8$ électrodes disposées de sorte que des électrodes voisines correspondent à des positions voisines de la balle & article \\
code de place & le numéro d'électrode $k=1,\dots,8$ indique la position de la balle par rapport au centre de la raquette & article \\
quelle coordonnée & position verticale de la balle; pour la reproduction, par rapport à la raquette, $k=\lceil 8(y-p+\tfrac12)\rceil$ pour $y-p\in[-\tfrac12,\tfrac12]$ & article; formule: choix \\
code de fréquence & la fréquence de stimulation, de $4$ à $40$~imp./s, porte la distance à la raquette & article \\
sens de l'échelle & $4$~imp./s quand la balle est près du mur opposé, et linéairement jusqu'à $40$~imp./s près du mur de la raquette: $f=4+36\,(1-x/L)$, $x$ étant la distance au mur de la raquette et $L$ la largeur du court & article \\
forme d'impulsion & brève impulsion rectangulaire biphasique: d'abord la phase positive, puis la négative & article \\
amplitude & $75$~mV; choisie pour ne pas forcer la décharge des cellules inhibées & article \\
durée des phases & non indiquée; prendre le réglage standard du système de stimulation et ne pas le modifier pendant l'expérience & choix \\
\bottomrule
\end{tabular}
\end{center}

\subsection*{Lecture de la raquette}

L'article comporte deux zones motrices: les impulsions de la première déplacent la raquette vers le haut, celles de la seconde vers le bas; la contribution des zones est pondérée pour tenir compte de densités de cellules et d'activités différentes~\cite{Kagan2022}. La formule exacte n'est pas donnée. Pour la reproduction, nous prenons la règle~\eqref{eq:dishreadout}, $\Delta p=c\,(n_\uparrow-n_\downarrow)$ pour chaque fenêtre de $10\ms$, avec un poids qui égalise les zones selon leur activité moyenne au repos (choix): $n_\uparrow$ et $n_\downarrow$ sont divisés par le compte moyen de leur zone par fenêtre, mesuré avant le jeu. Le coefficient $c$ est choisi de sorte qu'une différence d'un compte moyen déplace la raquette de $1\%$ de la hauteur du court par fenêtre (choix); un avantage constant d'une zone fait alors traverser tout le court à la raquette en $1$~s.

L'emplacement des zones sur la puce n'est pas donné par des coordonnées dans le texte de l'article; il n'y a qu'un schéma. Pour la reproduction, trois groupes d'électrodes disjoints suffisent: les huit électrodes sensorielles au centre et, de part et d'autre, un groupe d'électrodes d'enregistrement, l'un pour \og haut\fg{}, l'autre pour \og bas\fg{} (choix).

\subsection*{Le jeu}

Les dimensions du court, la longueur de la raquette et la vitesse de la balle ne sont pas indiquées dans l'article; il est seulement dit que la balle se déplace à vitesse constante et rebondit sur les murs. Pour la reproduction (choix): court de $1\times1$, raquette de longueur $0.2$ sur le mur de droite (renvoi si $|y-p|<0.1$, comme dans le modèle~\texttt{models/dishbrain\_model.py}), la balle traverse le court en $1$~s. Un raté sans aucun renvoi dans l'échange compte comme un \og ace\fg{}; un échange de plus de trois renvois consécutifs est dit \og long\fg{} (article).

\subsection*{Rétroaction}

\begin{center}
\footnotesize
\setlength{\tabcolsep}{4pt}
\renewcommand{\arraystretch}{1.15}
\begin{tabular}{>{\raggedright}p{2.6cm}>{\raggedright}p{8.3cm}>{\raggedright\arraybackslash}p{1.6cm}}
\toprule
Événement & Ce qui est appliqué & Source \\
\midrule
renvoi & stimulus prévisible: les $8$ électrodes sensorielles simultanément, $75$~mV, $100$~imp./s, $100\ms$; la stimulation habituelle de la balle est interrompue pendant ce temps & article \\
raté & stimulus imprévisible: $150$~mV, $5$~imp./s, $4$~s, à des instants aléatoires sur des électrodes tirées au hasard parmi les $8$; puis $4$~s sans stimulation, et la balle repart dans une direction aléatoire & article \\
loi du hasard & non indiquée; pour la reproduction, instants des impulsions en flux de Poisson de fréquence moyenne $5$~imp./s & choix \\
\bottomrule
\end{tabular}
\end{center}

\subsection*{Protocole et contrôles}

Une session dure $20$~min; on compare les $5$ premières et les $15$ dernières minutes (article). Trois mesures du résultat: la longueur moyenne d'un échange, la proportion d'aces et la proportion d'échanges longs. Conditions expérimentales (article):
\begin{itemize}
\item \emph{stimulus}: boucle complète, telle que décrite ci-dessus;
\item \emph{silence}: au lieu du stimulus de renvoi ou de raté, la même durée sans aucune stimulation, puis la balle repart dans une direction aléatoire;
\item \emph{sans rétroaction}: après un raté, la balle rebondit simplement et continue sa course, et la stimulation de la position de la balle se poursuit;
\item \emph{repos}: le jeu tourne et la raquette suit l'activité, mais il n'y a aucune stimulation;
\item \emph{contrôles}: puce avec milieu seul; cellules HEK293T; \og culture in silico\fg{}, une simulation à la place des cellules.
\end{itemize}
Taille des échantillons dans la série principale: $9$ cultures de souris ($101$ sessions), $11$ cultures humaines ($138$), $6$ puces avec milieu seul ($80$), $20$ cultures au repos ($42$), $3$ simulations ($38$), soit $399$ sessions au total. Dans la série sur la rétroaction: $486$ sessions sur $3$ jours, à raison de $3$ sessions par jour, conditions \og stimulus\fg{}, \og silence\fg{} et \og sans rétroaction\fg{} ($n=27$, $17$ et $15$). La hausse de la longueur moyenne des échanges entre les $5$ premières minutes et les $15$ dernières n'a été obtenue qu'avec des cultures vivantes. Dans la série sur la rétroaction, une hausse significative au fil du temps n'apparaît que dans la condition \og stimulus\fg{}; en fin de session, la condition \og silence\fg{} dépassait tout de même \og sans rétroaction\fg{}, mais avec un effet moindre, et il n'y a pas eu d'apprentissage stable d'une session à l'autre sur les trois jours~\cite{Kagan2022}.

\subsection*{Le cycle du banc pas à pas}

Toutes les $10\ms$, l'ordinateur du banc effectue les opérations suivantes.
\begin{enumerate}
\item Il lit les nombres d'impulsions $n_\uparrow$, $n_\downarrow$ des deux zones motrices sur la fenêtre écoulée et les normalise par le compte moyen de la zone au repos.
\item Il déplace la raquette de $\Delta p=c\,(n_\uparrow-n_\downarrow)$ et la maintient dans les limites du court.
\item Il avance la balle d'un pas et la fait rebondir sur les murs du haut, du bas et de gauche.
\item Si la balle a atteint le mur de droite: pour $|y-p|<0.1$, c'est un renvoi, le compteur de l'échange augmente d'une unité et le stimulus de renvoi est appliqué ($100\ms$); sinon, c'est un raté, l'échange est enregistré, le stimulus de raté est appliqué ($4$~s), suivi de $4$~s de pause, et la balle repart.
\item Sinon, il choisit l'électrode $k$ d'après la position verticale de la balle et la fréquence $f$ d'après la distance au mur de la raquette, et applique à l'électrode $k$ des impulsions biphasiques de $75$~mV à la fréquence $f$ jusqu'à la fenêtre suivante.
\end{enumerate}
Cela suffit pour construire un banc compatible avec celui qui a été publié et pour tester la question centrale du chapitre: est-ce la prévisibilité qui enseigne à la culture, ou la simple différence entre la réponse au renvoi et celle au raté? Pour cela, il faut ajouter aux trois conditions de l'article une quatrième, qui n'y figure pas: après un raté, un stimulus \emph{prévisible} de même durée et de même énergie que le stimulus imprévisible. Si la culture apprend aussi dans ce cas, c'est la différence qui compte, et non la prévisibilité.


\section{Bilan}

\begin{table}[t]
\centering
\footnotesize
\setlength{\tabcolsep}{4pt}
\renewcommand{\arraystretch}{1.15}
\begin{tabular}{>{\raggedright}p{2.1cm}>{\raggedright}p{4.2cm}>{\raggedright}p{1.9cm}>{\raggedright\arraybackslash}p{4.6cm}}
\toprule
Bloc du banc & Fonctionnement & Chapitre du livre & Ce que l'on sait \\
\midrule
entrée & lieu (laquelle des 8 électrodes) et fréquence (4--40 imp./s) & \ref{sec:code}, \ref{sec:sensors} & fixée par l'expérimentateur \\
sortie & différence des comptes des deux zones en $10\ms$~\eqref{eq:dishreadout} & \ref{sec:glutcomputer}, \ref{sec:debugging} & fixée par l'expérimentateur; bruit selon~\eqref{eq:rateerror} \\
professeur & stimulus prévisible après un renvoi, imprévisible après un raté; pas de dopamine & \ref{sec:dofa} & hausse significative en une session seulement avec rétroaction complète; effet moindre avec le silence; instable d'une session à l'autre \\
mécanisme & prédiction de l'entrée ou brassage après un raté~\eqref{eq:dishdensity} & \ref{sec:dofa}, \ref{sec:recognition} & non établi; les deux expliquent le résultat \\
régime du réseau & rapprochement du régime critique pendant le jeu & \ref{sec:gaba} & lien avec la qualité du jeu mesuré \\
\og sensibilité\fg{} & affirmation des auteurs & --- & non démontrée; contestée \\
\bottomrule
\end{tabular}
\caption{DishBrain vu par un ingénieur.}
\label{tab:dishbrain}
\end{table}

DishBrain est une machine faite de neurones vivants sous sa forme la plus simple: l'entrée est fixée par un code, la sortie lue par un comptage, et le professeur est la différence entre la réponse au succès et la réponse à l'échec (tableau~\ref{tab:dishbrain}). Un réseau sans yeux, sans muscles et sans dopamine a appris à jouer un tout petit peu, et cela seul en fait un banc d'essai pour les idées du livre. Quel que soit le mécanisme (prédiction, brassage ou autre chose), la réponse sera trouvée comme le conseille le chapitre~\ref{sec:debugging}: avec une sonde, un signal de test et la mise hors circuit de certaines parties, sur une machine que l'on voit enfin en entier.

\section{Exercices}

\begin{exercise}[\lvlA]
\label{ex:22:1}
\emph{De combien décaler le compte.} Les deux zones émettent ensemble $R_\uparrow+R_\downarrow=2000$ impulsions par seconde, flux poissoniens. (a)~Dispersion relative du compte d'une zone en $10\ms$ pour $R=1000$ et $R=100$? (b)~En $1$~s de vol de la balle, les déplacements~\eqref{eq:dishreadout} s'additionnent. Plus petit $R_\uparrow-R_\downarrow$ donnant un déplacement moyen double de la dispersion?
\end{exercise}

\begin{exercise}[\lvlA]
\label{ex:22:2}
\emph{Combien d'échanges pour voir l'apprentissage.} Les échanges sont indépendants, la proportion de renvois est binomiale. Combien d'échanges faut-il dans chacune des deux périodes pour que la hausse de la proportion de renvois dépasse deux écarts types de la différence: (a)~de $0.20$ à $0.25$; (b)~de $0.21$ à $0.38$, comme dans le modèle?
\end{exercise}

\begin{exercise}[\lvlB]
\label{ex:22:3}
\emph{Démonstration de~\eqref{eq:dishdensity}.} Le réglage $\boldsymbol\psi$ parcourt un ensemble fini; en cas de raté (probabilité $P(\boldsymbol\psi)>0$), le réseau passe en $\boldsymbol\psi'$ avec une probabilité symétrique $K(\boldsymbol\psi,\boldsymbol\psi')$; en cas de renvoi, il reste. (a)~Démontrez que $\rho\propto1/P$ est la distribution stationnaire. (b)~Quelle est la proportion de ratés avec deux réglages où $P=0.9$ et $P=0.1$?
\end{exercise}

\begin{exercise}[\lvlB]
\label{ex:22:4}
\emph{La région absorbante du modèle.} La raquette arrive en $p=\min\bigl(1,\max(0,ay+b)\bigr)$, la balle à la hauteur $y\sim U(0,1)$, et il y a renvoi si $|p-y|<h=0.1$. (a)~Démontrez que pour $|b|<h$ et $|a-1+b|<h$ le réseau renvoie toutes les balles, et trouvez l'aire de cette région. (b)~Trouvez numériquement la proportion moyenne de renvois pour $a\sim U(-1,1)$, $b\sim U(0,1)$.
\end{exercise}

\begin{exercise}[\lvlB]
\label{ex:22:5}
\emph{Simulation: une clôture autour des réglages.} Dans une copie de \texttt{dishbrain\_model.py}, faites jouer $200$ cultures pendant $20\,000$ échanges chacune. Quelle est la proportion de renvois sur les $500$ derniers? Et si, après chaque secousse, on ramène le réglage dans le rectangle $a\in[-1,2]$, $b\in[-1,1]$? Expliquez la différence à l'aide de l'exercice~\ref{ex:22:3}.
\end{exercise}

\begin{exercise}[\lvlB]
\label{ex:22:6}
\emph{Conception du code d'entrée.} La balle est renvoyée si l'erreur de hauteur est inférieure à $h=0.1$; le réseau lit l'entrée en $T=0.25$~s et distingue environ $\sqrt{rT}$ niveaux jusqu'à la fréquence $r$; la stimulation ne dépasse pas $40$ impulsions par seconde. (a)~Un code de place sur huit électrodes suffit-il? (b)~Et la fréquence sur une seule électrode?
\end{exercise}

\begin{exercise}[\lvlC]
\label{ex:22:7}
\emph{Recherche: prédiction ou brassage?} Généralisez le modèle à $d$ nombres réglables ($p=\sum_{i<d}a_iy^i$). Comment le nombre d'échanges menant à $0.5$ de renvois croît-il avec $d$, par brassage et par la règle au signe de l'erreur~\eqref{eq:perturbation}? Quelle expérience sur le banc départagerait les deux hypothèses?
\end{exercise}

\appendix
\renewcommand{\chaptername}{\appendixname}
\chapter{Les substances qui ne définissent pas le type d'un neurone}
\label{sec:othermediators}

Au chapitre~\ref{sec:neurontypes}, nous avons classé les neurones selon leur neurotransmetteur principal et obtenu dix types (tableau~\ref{tab:types}). Il est temps d'ajouter les complications. Un neurone ne libère pas seulement son neurotransmetteur principal, et dans le cerveau, d'autres substances agissent en plus des neurotransmetteurs principaux. Elles modifient la façon dont le réseau transmet et stocke les signaux.

Outre le bus d'adresses et le bus de diffusion (section~\ref{sec:buses}), il en existe un troisième: le \emph{canal rétrograde}. Plusieurs substances ne sont pas libérées par l'émetteur mais par le \emph{destinataire}; elles traversent la fente en sens inverse, vers la sortie de l'émetteur, et modifient la quantité de neurotransmetteur que celui-ci libérera la fois suivante. Dans les réseaux de communication, cela ressemble à l'accusé de réception dont l'émetteur se sert pour ajuster sa transmission. C'est précisément ce canal qu'utilisent les neurones lorsqu'ils modifient les poids de leurs connexions: il informe la sortie de l'émetteur de l'activité du destinataire, c'est-à-dire du deuxième facteur des règles d'apprentissage du chapitre~\ref{sec:dofa}.

La liste complète des neurotransmetteurs connus est longue: plus d'une centaine de substances, dont la plupart sont de courtes chaînes protéiques, les neuropeptides. Mais toutes se rangent en quelques classes, et au sein de chaque classe les substances se comportent de façon semblable. Les substances qui ne sont jamais le neurotransmetteur principal sont réunies dans le tableau~\ref{tab:othermediators}, avec les trois questions que poserait un ingénieur: sur quel bus passe le signal, agit-il vite, et quel est d'ordinaire son signe. Elles ne définissent pas le type d'un neurone: soit elles sont libérées avec le neurotransmetteur principal, soit elles ne sont pas libérées par des neurones, soit elles circulent en sens inverse.

\begin{table}[t]
\centering
\footnotesize
\setlength{\tabcolsep}{4pt}
\renewcommand{\arraystretch}{1.12}
\begin{tabular}{>{\raggedright}p{2.25cm}>{\raggedright}p{2.8cm}>{\raggedright}p{1.8cm}>{\raggedright}p{1.5cm}>{\centering}p{0.8cm}>{\raggedright\arraybackslash}p{4.3cm}}
\toprule
Classe & Substance & Bus & Vitesse & Signe & Rôle dans le calcul \\
\midrule
Acides aminés & D-sérine & local & lente & ET & seconde clé du récepteur du glutamate: condition \og ET\fg{} \\
\midrule
Monoamines & amines traces & diffusion & lente & $\pm$ & réglage fin des canaux monoaminergiques \\
\midrule
\multirow{2}{2.25cm}{Purines}
 & ATP & adresse & rapide et lente & $+$ & neurotransmetteur associé; signal entre neurones et glie \\
 & adénosine & volumique & lente & $-$ & s'accumule avec l'activité: compteur de charge~\cite{Dunwiddie2001} \\
\midrule
\multirow{8}{2.25cm}{Neuropeptides (plus d'une centaine)}
 & enképhalines, endorphines, dynorphine & volumique & lente & $-$ & suppression de la transmission \\
 & substance P & volumique & lente & $+$ & renforcement lent \\
 & neuropeptide Y & volumique & lente & $-$ & inhibition lente \\
 & somatostatine & volumique & lente & $-$ & inhibition lente, accompagne le GABA \\
 & peptide vasoactif intestinal (VIP) & volumique & lente & $+$ & désinhibition: inhibition des neurones inhibiteurs \\
 & cholécystokinine & volumique & lente & $\pm$ & accompagne le GABA dans une partie des neurones inhibiteurs \\
 & orexine & diffusion & lente & $+$ & stabilisateur du régime de veille \\
 & ocytocine, vasopressine & diffusion & lente & $\pm$ & réglages globaux du comportement \\
\midrule
\multirow{2}{2.25cm}{Gaz}
 & monoxyde d'azote NO & rétrograde, volumique & lente & $\pm$ & traverse les membranes; signal rétrograde lors des changements de poids~\cite{Garthwaite2008} \\
 & CO, H$_2$S & volumique & lente & $\pm$ & rôle peu étudié \\
\midrule
Lipides & endocannabinoïdes (anandamide, 2-AG) & rétrograde & lente & $-$ & le destinataire réduit la libération de l'émetteur: rétroaction sur le poids~\cite{WilsonNicoll2001} \\
\bottomrule
\end{tabular}
\caption{Les substances qui ne définissent pas le type d'un neurone. \emph{Bus}, \emph{vitesse} et \emph{signe} comme dans le tableau~\ref{tab:types}; en outre, bus volumique: le neurotransmetteur diffuse dans l'espace extracellulaire autour du site de libération; rétrograde: il va du destinataire vers l'émetteur; local: il agit près du site de libération. \og ET\fg{} désigne un signal de validation: il ne produit pas de courant à lui seul, mais sans lui une autre entrée ne se déclenche pas, comme la seconde entrée d'une porte logique ET. Les neuropeptides sont presque toujours libérés avec le neurotransmetteur principal, et surtout à haute fréquence d'impulsions~\cite{vandenPol2012}.}
\label{tab:othermediators}
\end{table}

\chapter{Notations}
\label{sec:notation}

Les types de neurones sont désignés par les quatre premières lettres du nom anglais de leur neurotransmetteur principal: \nt{GLUT}, glutamate; \nt{GABA}, acide gamma-aminobutyrique; \nt{GLYC}, glycine; \nt{ACET}, acétylcholine; \nt{DOPA}, dopamine; \nt{SERO}, sérotonine; \nt{HIST}, histamine; \nt{NORA}, noradrénaline; \nt{ADRE}, adrénaline; \nt{ASPA}, aspartate (tableau~\ref{tab:types}).

Il y a moins de lettres que de grandeurs, et d'un chapitre à l'autre une même lettre désigne parfois des choses différentes. Dans le tableau ci-dessous, chaque signification a sa propre ligne; la colonne de droite indique la formule où elle est introduite.

\begin{center}
\footnotesize
\setlength{\tabcolsep}{4pt}
\renewcommand{\arraystretch}{1.12}
\begin{longtable}{>{\raggedright}p{2.0cm}>{\raggedright}p{9.6cm}>{\raggedright\arraybackslash}p{2.4cm}}
\toprule
Symbole & Signification & Introduit en \\
\midrule
\endfirsthead
\toprule
Symbole & Signification & Introduit en \\
\midrule
\endhead
\bottomrule
\endfoot
$a$ & pente d'une caractéristique de transfert linéaire, $f(u)=au$ & \eqref{eq:feedback} \\
$a$ & niveau d'\nt{ACET}: $0$ pour la lecture, $1$ pour l'écriture & \eqref{eq:acetnet} \\
$a_i$, $a$ & fatigue d'un groupe: dans le générateur de rythme; dans la bascule du sommeil lent & \eqref{eq:halfcentre}, \eqref{eq:updown} \\
$a_S$, $a_P$ & sensibilité du rythme cardiaque à l'accélérateur et au frein & \eqref{eq:gasbrake} \\
$A(r)$, $A_0$ & réponse d'une synapse à une impulsion à la fréquence $r$; pour une synapse reposée & \eqref{eq:depression} \\
$A_1$ & réponse du destinataire à une seule dose de neurotransmetteur & \eqref{eq:quantal} \\
$A$ & surface de la membrane du neurone, dendrites comprises & \eqref{eq:dendriteload} \\
$A_f$, $A_s$ & part de la correction conservée jusqu'au mouvement suivant, pour les processus rapide et lent & \eqref{eq:tworate} \\
$b_i$ & apport propre d'un stimulateur (pacemaker) & \eqref{eq:seroneuron} \\
$b$ & vitesse à laquelle le travail fatigue un groupe dans le générateur de rythme ou dans la bascule du sommeil & \eqref{eq:halfcentre}, \eqref{eq:updown} \\
$b$ & nombre de bits par échantillon numérisé & \eqref{eq:probedata} \\
$B_f$, $B_s$ & force avec laquelle la correction apprend de l'erreur, pour les processus rapide et lent & \eqref{eq:tworate} \\
$c$ & concentration du neurotransmetteur dans le volume & \eqref{eq:seroneuron} \\
$c$ & coefficient de corrélation des bruits de neurones voisins & \eqref{eq:correlated} \\
$c$ & coefficient avec lequel la différence des comptes déplace la raquette & \eqref{eq:dishreadout} \\
$c_f$ & réponse d'un neurone $C$: le plus grand des $s_{f,p}$ sur toutes les positions & \eqref{eq:scstage} \\
$c_m$ & capacité membranaire spécifique & \eqref{eq:dendriteload} \\
$C$ & coût de l'effort: une action de durée $\tau_a$ coûte $C/\tau_a$ & \eqref{eq:vigor} \\
$d'$ & discriminabilité de deux entrées: leur différence divisée par le bruit & \eqref{eq:gainsnr} \\
$d$, $d_j$ & retard sur le trajet du signal & \eqref{eq:glutneuron}, \eqref{eq:vstar} \\
$d$ & retard de la boucle de rétroaction avec le muscle & \eqref{eq:delaystable} \\
$D$ & coefficient de diffusion & \eqref{eq:diffusionlength} \\
$D$ & dénominateur des fréquences du réseau \nt{GLUT}--\nt{GABA} & \eqref{eq:isn} \\
$D$ & délai d'attente d'une récompense différée & \eqref{eq:patience} \\
$e$, $e_{ij}$ & trace d'éligibilité dans une synapse & \eqref{eq:threefactor} \\
$e$, $e_i$ & erreur de sortie ou de mouvement, arrivant par le fil d'erreur & \eqref{eq:lms}, \eqref{eq:adaptivefilter}, \eqref{eq:tworate} \\
$E_{\text{imp}}$ & énergie d'une impulsion, travail des synapses compris & \eqref{eq:energy} \\
$E$ & fréquence du groupe \nt{GLUT} & \eqref{eq:ei} \\
$E_E$ & potentiel d'équilibre d'une synapse excitatrice & \eqref{eq:shunt} \\
$f$, $F$ & caractéristique de transfert: fréquence en fonction de l'entrée & \eqref{eq:ratenet}, \eqref{eq:gain} \\
$f$, $f_0$ & fréquence cardiaque; rythme propre du stimulateur & \eqref{eq:gasbrake} \\
$f_s$ & fréquence d'échantillonnage d'une électrode & \eqref{eq:probedata} \\
$F(t)$, $F_1$ & force du muscle; force d'une secousse unique & \eqref{eq:forcesum} \\
$F_1$, $F_2$ & forces de deux muscles qui tirent une articulation en sens opposés & \eqref{eq:antagonists} \\
$g$ & sensibilité du récepteur à la concentration & \eqref{eq:seroneuron} \\
$g_L$, $g_E$, $g_I$ & conductances de fuite, d'excitation et d'inhibition & \eqref{eq:shunt} \\
$g$ & gain fixé par \nt{NORA} & \eqref{eq:gain}, \eqref{eq:machine} \\
$g$ & inhibition globale due au \nt{GABA} dans le réseau de mémoire & \eqref{eq:acetnet} \\
$g$ & gain du portillon de la douleur, fixé d'en haut & \eqref{eq:gate} \\
$g$ & gain d'un étage de la cascade hormonale & \eqref{eq:cascade} \\
$g_j$ & filtre $j$, avec sa propre constante de temps, à l'entrée du filtre adaptatif & \eqref{eq:adaptivefilter} \\
$G$ & gain de la boucle de rétroaction & \eqref{eq:feedback} \\
$G$ & vitesse de correction dans une boucle à retard & \eqref{eq:delaystable} \\
$h$, $h_I$ & entrée externe du groupe \nt{GLUT} et du groupe \nt{GABA} & \eqref{eq:ei}, \eqref{eq:isn} \\
$h(t)$ & secousse musculaire unique & \eqref{eq:twitch} \\
$H(t)$ & seuil haut de la pression de sommeil: en l'atteignant, la machine s'endort & \eqref{eq:sleeppressure} \\
$I$, $I_i$ & apport externe à un neurone ou à un groupe & \eqref{eq:ratenet}, \eqref{eq:wta}, \eqref{eq:updown} \\
$I$ & fréquence du groupe \nt{GABA} & \eqref{eq:ei} \\
$I$ & luminosité ou intensité du stimulus & \eqref{eq:photonnoise}, \eqref{eq:nakarushton} \\
$I$ & courant d'entrée qui charge la membrane & \eqref{eq:dendriteload} \\
$k$ & nombre d'impulsions dans une bouffée & \eqref{eq:burst} \\
$k$ & combien de fois l'accroissement doit dépasser le bruit & \eqref{eq:photonnoise} \\
$k$ & coefficient de rétroaction de la cascade hormonale & \eqref{eq:cascade} \\
$k_s$ & force avec laquelle la douleur alimente la sensibilisation & \eqref{eq:sensitization} \\
$K$ & nombre de caractéristiques du résultat & \eqref{eq:vectortd} \\
$K$ & raideur de l'articulation & \eqref{eq:antagonists} \\
$K$ & gain de boucle de la cascade hormonale, $K=g^3k$ & \eqref{eq:cascade}, \eqref{eq:cascadestable} \\
$K_\varphi$ & force de l'accrochage du générateur circadien au signal externe & \eqref{eq:pll} \\
$L$ & longueur de la ligne à retard & \eqref{eq:itd} \\
$L$ & longueur du fil venant du capteur de douleur & \eqref{eq:paintime} \\
$L(t)$ & seuil bas de la pression de sommeil: en l'atteignant, la machine se réveille & \eqref{eq:sleeppressure} \\
$M$ & mesure du lien causal entre un signe annonciateur et la récompense & \eqref{eq:retro} \\
$M_k$ & attente de la caractéristique $k$ & \eqref{eq:vectortd} \\
$M_{ik}$ & force de l'inhibition venant du neurone \nt{GABA} $k$ & \eqref{eq:machine} \\
$M$ & couple de l'articulation & \eqref{eq:antagonists} \\
$M$ & nombre de lignes d'entrée des mélangeurs & \eqref{eq:cover} \\
$n$, $\bar n$ & nombre d'impulsions dans une fenêtre; sa moyenne & \eqref{eq:rateerror} \\
$n$ & nombre d'entrées d'un élément à seuil & \eqref{eq:cover} \\
$n_\uparrow$, $n_\downarrow$ & compte des impulsions des zones \og haut\fg{} et \og bas\fg{} par fenêtre & \eqref{eq:dishreadout} \\
$N$ & nombre de neurones & \eqref{eq:correlated}, \eqref{eq:energy} \\
$N$ & nombre d'électrodes de la sonde & \eqref{eq:probedata} \\
$N_s$ & nombre de sites de libération entre deux neurones & \eqref{eq:quantal} \\
$p$ & probabilité de libération du neurotransmetteur par impulsion & \eqref{eq:burst} \\
$p$ & proportion de neurones actifs dans le réseau de mémoire & \eqref{eq:acetnet} \\
$p$ & perturbation que l'on apprend à compenser & \eqref{eq:tworate} \\
$P$ & puissance consacrée aux signaux & \eqref{eq:energy}, \eqref{eq:dishpower} \\
$P$ & incertitude de l'estimation stockée & \eqref{eq:kalman} \\
$P$ & activité du \og frein\fg{}: les fils \nt{ACET} vers le cœur & \eqref{eq:gasbrake} \\
$P_{\max}$ & nombre de motifs aléatoires qu'un élément à seuil peut séparer en deux classes & \eqref{eq:cover} \\
$P_{\text{raté}}$ & probabilité de raté & \eqref{eq:dishdensity} \\
$q$ & vitesse à laquelle le monde change & \eqref{eq:kalman} \\
$q_k$ & fréquence du neurone \nt{GABA} $k$ & \eqref{eq:machine} \\
$q_0$ & activité de fond de l'intermédiaire inhibiteur & \eqref{eq:gate} \\
$q$ & fréquence de l'intermédiaire inhibiteur dans le portillon de la douleur & \eqref{eq:gate} \\
$q$ & combien de fois l'unité motrice suivante est plus forte que la précédente & \eqref{eq:sizeprinciple} \\
$r$, $r_i$ & fréquence des impulsions d'un neurone & \eqref{eq:ratenet} \\
$r$, $r_t$ & récompense (flux de récompense) & \eqref{eq:rpe}, \eqref{eq:ctd} \\
$\mathbf r(t)$ & vecteur des réponses du réservoir & \eqref{eq:reservoir} \\
$R$, $\bar R$ & récompense reçue; son attente ou sa moyenne par unité de temps & \eqref{eq:perturbation}, \eqref{eq:vigor} \\
$R$ & variance du bruit à l'entrée du filtre & \eqref{eq:kalman} \\
$R$, $r_0$ & récompenses différée et immédiate & \eqref{eq:patience} \\
$R_{\max}$ & débit de transmission maximal, en bit/s & \eqref{eq:spikeentropy} \\
$r_{\max}$ & fréquence d'impulsions maximale & \eqref{eq:gain}, \eqref{eq:nakarushton} \\
$R_{\text{brut}}$, $R_{\text{imp}}$ & débit de données de la sonde: brut et impulsions seules, en bit/s & \eqref{eq:probedata} \\
$s_j$ & signe des sorties du neurone $j$ & \eqref{eq:dalesign} \\
$s_i$ & signe du récepteur du destinataire & \eqref{eq:seroneuron} \\
$s$, $\hat s$ & état du monde; son estimation & \eqref{eq:rpe}, \eqref{eq:kalman} \\
$s_{f,p}$ & réponse d'un neurone $S$ à la caractéristique $f$ en position $p$ & \eqref{eq:scstage} \\
$S$ & activité de l'\og accélérateur\fg{}: les fils \nt{NORA} vers le cœur & \eqref{eq:gasbrake} \\
$S$ & pression de sommeil & \eqref{eq:sleeppressure} \\
$t_1$ & instant de la première impulsion & \eqref{eq:latency} \\
$t_k$ & instant de la $k$-ième impulsion & \eqref{eq:forcesum} \\
$t_{f,i}$ & modèle de la caractéristique $f$ & \eqref{eq:scstage} \\
$T$ & fenêtre de comptage des impulsions; temps d'accumulation des photons & \eqref{eq:rateerror}, \eqref{eq:photonnoise} \\
$T_c$ & temps de montée d'une secousse musculaire & \eqref{eq:twitch} \\
$T_{\text{salve}}$ & intervalle entre les salves d'une culture & \eqref{eq:burstinterval} \\
$u$ & entrée totale d'un neurone & \eqref{eq:gain} \\
$u$, $u(t)$ & commande du cerveau: à l'entrée du filtre adaptatif; vers la cascade hormonale & \eqref{eq:adaptivefilter}, \eqref{eq:cascade} \\
$u(t)$ & entrée du réservoir & \eqref{eq:reservoir} \\
$U$ & niveau d'une entrée constante & \eqref{eq:latency} \\
$U$ & fraction de la réserve de neurotransmetteur libérée par impulsion & \eqref{eq:depression} \\
$v$ & vitesse du signal dans le fil; vitesse de déplacement de la tache & \eqref{eq:itd}, \eqref{eq:vstar} \\
$V$ & tension de membrane & \eqref{eq:glutneuron}, \eqref{eq:shunt} \\
$V$ & attente de la récompense future & \eqref{eq:rpe}, \eqref{eq:ctd} \\
$w$, $w_{ij}$, $W_{ij}$ & poids d'une connexion & \eqref{eq:glutneuron}, \eqref{eq:ratenet}, \eqref{eq:acetnet}, \eqref{eq:lms}, \eqref{eq:adaptivefilter} \\
$\mathbf w_k$ & poids des entrées du neurone $C$ numéro $k$ & \eqref{eq:traceinvariance} \\
$w_{q\mu}$, $w_{q\nu}$, $w_{z\mu}$ & poids des connexions dans le portillon de la douleur & \eqref{eq:gate} \\
$W_{\text{sort}}$ & poids de la lecture linéaire du réservoir & \eqref{eq:reservoir} \\
$x$, $y$ & entrées et sortie logiques & \eqref{eq:dualrail}, \eqref{eq:veto} \\
$x$, $y$ & activité de l'émetteur et du destinataire & \eqref{eq:hebb} \\
$x$ & position du détecteur sur la ligne à retard & \eqref{eq:itd} \\
$x_i$ & entrée venant des organes des sens & \eqref{eq:acetnet} \\
$x_p$ & entrée en position $p$ & \eqref{eq:scstage} \\
$\mathbf x$ & sorties des neurones $S$, entrée du neurone $C$ & \eqref{eq:traceinvariance} \\
$x_j$ & entrée $j$ du filtre adaptatif: la commande après le filtre $g_j$ & \eqref{eq:lms}, \eqref{eq:adaptivefilter} \\
$x_f$, $x_s$ & corrections des processus rapide et lent & \eqref{eq:tworate} \\
$x_1$, $x_2$, $x_3$ & niveaux aux trois étages de la cascade hormonale; $x_3$ est l'hormone finale & \eqref{eq:cascade} \\
$x^*$ & fraction de la réserve de neurotransmetteur reconstituée à laquelle démarre une nouvelle avalanche & \eqref{eq:burstinterval} \\
$y$ & entrée bruitée du filtre & \eqref{eq:kalman} \\
$y$, $\hat y$ & signal des capteurs; sa prédiction par le filtre adaptatif & \eqref{eq:adaptivefilter} \\
$y_k$, $\bar y_k$ & activité du neurone $C$ numéro $k$; sa trace & \eqref{eq:traceinvariance} \\
$y(t)$ & sortie de la lecture du réservoir & \eqref{eq:reservoir} \\
$z$ & fréquence de sortie du portillon de la douleur & \eqref{eq:gate} \\
\midrule
$\alpha_i^\pm$ & poids des erreurs positives et négatives & \eqref{eq:asymmetric} \\
$\beta$ & force de l'inhibition mutuelle & \eqref{eq:wta} \\
$\beta$ & force de l'inhibition de la sortie du portillon de la douleur & \eqref{eq:gate} \\
$\gamma$ & facteur d'actualisation & \eqref{eq:rpe} \\
$\delta(\cdot)$ & fonction delta: un saut instantané & \eqref{eq:glutneuron} \\
$\delta$, $\delta_t$ & erreur de prédiction de la récompense & \eqref{eq:rpe}, \eqref{eq:ctd} \\
$\delta t$ & différence des temps d'arrivée du son aux deux oreilles & \eqref{eq:itd} \\
$\Delta$, $\Delta_{\max}$ & intervalle entre impulsions; fenêtre de coïncidence & \eqref{eq:window} \\
$\Delta t$ & précision avec laquelle le destinataire distingue les temps & \eqref{eq:spikeentropy} \\
$\Delta F$ & accroissement de force apporté par l'unité motrice suivante & \eqref{eq:sizeprinciple} \\
$\Delta I$ & accroissement de luminosité perceptible & \eqref{eq:photonnoise} \\
$\Delta p$ & déplacement de la raquette par fenêtre & \eqref{eq:dishreadout} \\
$\Delta u$ & différence des entrées de deux groupes & \eqref{eq:gainsnr} \\
$\Delta V$ & élévation nécessaire de la tension de membrane & \eqref{eq:dendriteload} \\
$\Delta x$ & distance entre capteurs voisins & \eqref{eq:vstar} \\
$\varepsilon$ & différence relative des entrées & \eqref{eq:wtagain} \\
$\eta$ & vitesse d'apprentissage & \eqref{eq:plasticity}, \eqref{eq:lms}, \eqref{eq:traceinvariance} \\
$\eta$ & bruit dans le réseau après l'amplificateur & \eqref{eq:softmax} \\
$\mu$ & entrée tactile du portillon de la douleur & \eqref{eq:gate} \\
$\nu$ & entrée douloureuse & \eqref{eq:gate} \\
$\theta$ & seuil de déclenchement & \eqref{eq:window}, \eqref{eq:scstage} \\
$\kappa$ & quantité de neurotransmetteur par impulsion & \eqref{eq:seroneuron} \\
$\kappa_i$ & rapport des poids des erreurs positives et négatives & \eqref{eq:expectile} \\
$\kappa_m$ & lien entre la force d'un muscle et sa raideur & \eqref{eq:antagonists} \\
$\ell$ & distance sur laquelle le neurotransmetteur diffuse & \eqref{eq:diffusionlength} \\
$\ell$ & bras de levier avec lequel le muscle tire l'articulation & \eqref{eq:antagonists} \\
$\xi$ & fluctuation aléatoire de la sortie d'un neurone & \eqref{eq:perturbation} \\
$\rho(\boldsymbol\psi)$ & fraction du temps passé dans le réglage $\boldsymbol\psi$ & \eqref{eq:dishdensity} \\
$\sigma$ & dispersion, échelle du bruit & \eqref{eq:rateerror}, \eqref{eq:softmax} \\
$\sigma$ & luminosité à laquelle la réponse vaut la moitié du maximum & \eqref{eq:nakarushton} \\
$\sigma_{\text{avant}}$, $\sigma_{\text{après}}$ & bruit avant et après l'amplificateur & \eqref{eq:gainsnr} \\
$\tau$ & constante de temps de la membrane & \eqref{eq:glutneuron} \\
$\tau$ & constante de temps d'un étage de la cascade hormonale & \eqref{eq:cascade} \\
$\tau_{\text{eff}}$ & constante de temps d'un groupe qui s'auto-excite & \eqref{eq:perfectint} \\
$\tau_c$ & durée de vie du neurotransmetteur dans le volume & \eqref{eq:seroneuron} \\
$\tau_E$, $\tau_I$ & constantes de temps des groupes \nt{GLUT} et \nt{GABA} & \eqref{eq:ei} \\
$\tau_e$ & durée de vie de la trace d'éligibilité & \eqref{eq:threefactor} \\
$\tau_{\text{tr}}$ & durée de vie de la trace d'activité d'un neurone $C$ & \eqref{eq:traceinvariance} \\
$\tau_s$ & temps de retour de la sensibilité après une lésion & \eqref{eq:sensitization} \\
$\tau_r$, $\tau_d$ & temps d'accumulation de la pression de sommeil pendant la veille et de sa décharge pendant le sommeil & \eqref{eq:sleeppressure} \\
$\tau_{\text{rec}}$ & temps de reconstitution de la réserve de neurotransmetteur & \eqref{eq:depression} \\
$\tau_\gamma$ & horizon de planification & \eqref{eq:ctd} \\
$\tau_v$ & vitesse à laquelle l'attente se précise & \eqref{eq:ramp} \\
$\tau_a$ & temps de fatigue dans le générateur de rythme ou dans la bascule du sommeil & \eqref{eq:halfcentre}, \eqref{eq:updown} \\
$\tau_a^*$ & durée optimale d'exécution d'une action & \eqref{eq:vigor} \\
$\Phi$ & membre de droite de la règle d'apprentissage & \eqref{eq:plasticity} \\
$\Phi$ & flux de photons & \eqref{eq:photonnoise} \\
$\varphi_k$ & caractéristique $k$ du résultat obtenu & \eqref{eq:vectortd} \\
$\varphi$ & déphasage entre le générateur circadien et le signal externe & \eqref{eq:pll} \\
$\boldsymbol\psi$ & réglage du réseau dans le modèle DishBrain & \eqref{eq:dishdensity} \\
$\omega$, $\omega_0$ & fréquence du signal externe (la lumière); fréquence propre du générateur circadien & \eqref{eq:pll} \\
\end{longtable}
\end{center}

\chapter{Réponses}
\label{sec:answers}

Réponses brèves aux exercices de fin de chapitre. Les solutions complètes figurent dans un recueil séparé destiné aux enseignants.

\answer{ex:01:1}{(a)~$8\cdot10^{10}$ \nt{GLUT}: $10$ populations de la Terre; ceux de diffusion, $\approx1.9\cdot10^6$: une ville de la taille de Vienne. (b)~$\approx8\cdot10^{14}$ terminaisons, dont $\approx3\cdot10^{11}$ de diffusion, soit une part de $\approx4\cdot10^{-4}$, $\approx16$ fois plus grande que la part de ces neurones ($2.2\cdot10^{-5}$).}
\answer{ex:01:2}{(a)~$\tau_c\ln(c_0/c_*)\approx4.6\ms$. (b)~La ligne est libre si $T\ge\tau_c\ln101$: jusqu'à $\approx22$ libérations par seconde au lieu de $\approx220$.}
\answer{ex:01:3}{$V(s_k)=\gamma^{K-1-k}r$; $\delta_{\text{lampe}}=\gamma^Kr=re^{-T/\tau_\gamma}\approx0.60\,r$ quel que soit $\Delta$, $\tau_\gamma\approx1.95\,\text{s}$.}
\answer{ex:01:4}{$n^*=93$, $47$, $25$, $12$: $n^*\approx5/\alpha$ pour $\alpha$ petit.}
\answer{ex:01:5}{\nt{GLUT}: $1\to10\to10^4\to10^7$, soit $\approx10^7$ neurones, $4$ maillons, $4\ms$. \nt{DOPA}: $1\to10\to3.2\cdot10^5$, soit $\approx3\cdot10^5$ neurones, $3$ maillons, $30$ fois moins; c'est le même ordre de grandeur que le nombre de neurones \nt{DOPA}.}
\answer{ex:01:6}{$\lambda=f_EJ_E-f_IJ_I$, d'où $J_I=37.5$; rapport des courants $f_IJ_I/f_EJ_E=0.94$; avalanche ($\lambda\ge1$) dès que l'inhibition est affaiblie de seulement $6.7\,\%$.}
\answer{ex:01:7}{$n\ge57$ présentations pour chaque délai dans chaque condition. La même décroissance pourrait venir, par exemple, d'une imprécision de l'estimation interne du temps qui croît avec $T$.}

\answer{ex:02:1}{Rangée de $3.2$~mm de long au pas de $25\um$: $\approx129$ détecteurs. Se déclenche une bande de largeur $v\,\tau\ln2=1.7$~mm, soit $\approx69$ détecteurs, plus de la moitié de la rangée; la direction est le centre de la bande.}
\answer{ex:02:2}{Fenêtre $-\tau\ln\frac{w_2}{\theta-w_1}\le\delta\le\tau\ln\frac{w_1}{\theta-w_2}$, c.-à-d. $[-0.235,\,0.178]\ms$; centre $c=\frac\tau2\ln\frac{w_1(\theta-w_1)}{w_2(\theta-w_2)}=-28$~\textmu s. La direction se décale de $28$~\textmu s vers l'oreille la plus forte, près de trois pas de résolution.}
\answer{ex:02:3}{Il faut des $s_i$ tels que $s_iw_{ij}s_j\ge0$; ils existent si et seulement si les cycles sont pairs. L'inhibition mutuelle s'y ramène (pas d'oscillations durables), la boucle \nt{GLUT}--\nt{GABA} non (elle peut osciller).}
\answer{ex:02:4}{Six neurones $g_k=[x+y+z\ge k]$ et $\bar g_k=[\bar x+\bar y+\bar z\ge4-k]$, $k=1,2,3$; $C=g_2$, $\bar C=\bar g_2$, $S=[g_1+\bar g_2+g_3\ge2]$, $\bar S=[\bar g_1+g_2+\bar g_3\ge2]$: huit neurones. Avec des ET et des OU: $12$.}
\answer{ex:02:5}{$T_{\min}(0.6)=0.718\,\tau$; $I_c=0.225$; près du seuil, $T_{\min}\propto(I-I_c)^{-1/2}$.}
\answer{ex:02:6}{(a)~$500$ maillons, $5\cdot10^4$ neurones; dispersion $4.5\ms$ ($0.45\,\%$), erreur relative $\propto T^{-1/2}$. (b)~Un facteur aléatoire de retard commun à tous les maillons, avec une dispersion de $5\,\%$.}
\answer{ex:02:7}{Rafraîchir une fois toutes les $\tau_F\ln2=0.69\,\text{s}$: $4.3$ impulsions par seconde et par neurone contre $20$, soit $4.6$ fois plus économe ($4\cdot10^{-7}$ contre $2\cdot10^{-6}$~J à $10^{-10}$~J par impulsion). La bascule perd le bit; le condensateur le garde si le silence commence au plus tard $0.49\,\text{s}$ après le rafraîchissement (probabilité $0.71$).}

\answer{ex:03:1}{$35\to17.5\mV$ pour $g_E=g_L$; $56\to40\mV$ pour $g_E=4g_L$, une soustraction donnerait $38.5\mV$: le shunt divise $g_E$ par $3$. La fenêtre est réduite de moitié ($\tau_{\text{eff}}$ est divisée par $2$).}
\answer{ex:03:2}{$\operatorname{tr}=\frac{w_{EE}-1}{\tau_E}-\frac1{\tau_I}$, $\det=\frac{1-w_{EE}+w_{EI}w_{IE}}{\tau_E\tau_I}$; à la limite $\omega=\sqrt{\det}$, $\tau_I=20\ms$, période $73\ms$.}
\answer{ex:03:3}{$d_c=\tau_E\arccos(-1/G)/\sqrt{G^2-1}$, période $2\pi\tau_E/\sqrt{G^2-1}\in(2d_c,4d_c)$. $G=2$: $d_c=12.1\ms$, période $36\ms$; $G=5$: $d_c=3.6\ms$, période $12.8\ms$, en plein dans la bande des rythmes rapides.}
\answer{ex:03:4}{(a)~À $I_2=\beta I_1=2$ en montant et à $I_2=I_1/\beta=0.5$ en descendant: une boucle d'hystérésis de largeur $1.5$. (b)~De $\tau\ln10/(\beta-1)=2.3\tau$ par ordre de grandeur de $\varepsilon$.}
\answer{ex:03:5}{Trois intermédiaires, $d_k=5$, $10$, $15\ms$: les vetos doivent couvrir $15\ms$, à raison de $5\ms$ chacun.}
\answer{ex:03:6}{$n+M+1=3+4+1=8$ neurones. Avec deux: \nt{GABA} $=[x+y+z\ge2]$, sortie $=[x+y+z-2\,\nt{GABA}\ge1]$.}
\answer{ex:03:7}{$\binom84=70<100\le\binom94=126$: $9$ intermédiaires (théorème de Sperner). Sans connexions directes: $100$, car le rang de $\beta(J-I)$ vaut $n$.}
\answer{ex:03:8}{$h_c=\frac1{2w_{EE}}+\frac{w_{EI}w_{IE}-w_{EE}}{4w_{EE}^2}=\frac7{18}\approx0.39$; la simulation donne un changement de signe entre $h=0.38$ et $0.40$.}

\answer{ex:04:1}{(a)~$\approx1.1\cdot10^{11}$ bits/J. (b)~$\log_2\sqrt{10}\approx1.7$ bit contre $81$: $\approx50$ fois moins.}
\answer{ex:04:2}{$r^*=(P_0/E_{\text{imp}})\ln\bigl(1/(r^*\Delta t)\bigr)$, $P_0/E_{\text{imp}}=1.16\,\text{s}^{-1}$: $r^*\approx6$ impulsions par seconde, $7.4\cdot10^{10}$ bits/J.}
\answer{ex:04:3}{$\mathcal I=\frac{N}{1+(N-1)c}=9.9$ (limite $10$) contre $\mathcal I=\frac{N}{1-c}=1111$: un rapport de $\approx110$; la corrélation n'est nuisible que si elle ressemble au signal.}
\answer{ex:04:4}{$\theta\approx68.7$ et $19.9$ (en unités de $w$), $m=19$ et $m=10$: le destinataire rapide a besoin de presque deux fois moins d'entrées coïncidentes, bien que son entrée moyenne soit cinq fois plus petite.}
\answer{ex:04:5}{$U\tau_{\text{rec}}\ge0.725\,\text{s}$ et $\tau_{\text{rec}}(1-2U)\le0.05\,\text{s}$, ce qui exige $U\ge0.483$; le plus petit $\tau_{\text{rec}}=0.725\,\text{s}$ pour $U=1$ ($1.45\,\text{s}$ pour $U=0.5$).}
\answer{ex:04:6}{(a)~$\theta\,e/(e-1)\approx1.58\,\theta$ quels que soient $\tau$ et $\theta$. (b)~$14$ impulsions.}
\answer{ex:04:7}{(a)~$1.62$ bit par mot: $0.77$ pour le nombre d'impulsions, $0.84$ pour leur disposition. (b)~$1.13$ et $1.59$ bit: sur $30$ mots, l'estimation perd près d'un tiers.}

\answer{ex:05:1}{$\ell\approx17\um$; $\approx100$~jours: la diffusion générale dans le cerveau est assurée par la ramification du fil, la diffusion moléculaire ne fait qu'étaler chaque site de libération sur $\sim20\um$.}
\answer{ex:05:2}{$r=ab/(1+G)$, $S=1/(1+G)$ exactement; $+1.7\,\%$ au lieu de $+20\,\%$.}
\answer{ex:05:3}{$T^*=1.21\,\text{s}$ pour $G=2$ et $0.19\,\text{s}$ pour $G=9$: en pratique $G\sim2$--$4$, soit une réduction d'un facteur $3$ à $5$.}
\answer{ex:05:4}{$\mathrm{CV}=1/\sqrt{2\lambda\tau_c}\approx9\cdot10^{-4}$ pour une fréquence totale $\lambda$; un seul neurone demande $\tau_c=12.5\,\text{s}$.}
\answer{ex:05:5}{$c\approx0.40$ dans les deux cas; $\mathrm{CV}=0.050$ contre $0.158$, soit un facteur $\sqrt{10}$. Les fréquences des neurones restent proportionnelles aux $a_i$: seule la somme est régulée.}
\answer{ex:05:6}{$N_1=\overline{A\vee B}$, $N_2=\overline{A\vee N_1}$, $N_3=\overline{B\vee N_1}$, $N_4=\overline{N_2\vee N_3}$, XOR $=N_5=\overline{N_4}$; chaque commutation prend $\tau_c\ln2$, le chemin de $4$ portes $2.8\,\text{s}$ par XOR.}
\answer{ex:05:7}{(a)~$\tau_\gamma=6/\ln3\approx5.5\,\text{s}$. (b)~$\bar D<4\,\text{s}$ contre $D<2.2\,\text{s}$: l'imprévisibilité du délai rend plus patient.}
\answer{ex:05:8}{$k_2=1/\tau_d=10\,\text{s}^{-1}$, $k_1=1/\tau_d-1/\tau_\gamma=9.5\,\text{s}^{-1}$; $\tau_\gamma=1/(k_2-mk_1)$: doublement pour $+2.6\,\%$ (et pour $+5.3\,\%$ l'horizon devient infini).}

\answer{ex:06:1}{$e=(1-e^{-T_a/\tau_e})e^{-d/\tau_e}$. (a)~$0.110$ contre $4.5\cdot10^{-5}$, soit $\approx2.5\cdot10^3$ fois plus. (b)~$\tau_e^*=T_a/\ln(1+T_a/d)=0.59\,\text{s}$.}
\answer{ex:06:2}{Vitesse $\nu_x\nu_y(A_+\tau_+-A_-\tau_-)=0.8A_+$ par seconde, $\approx2900A_+$ en une heure; $\tau_-=\tau_+/0.6=33\ms$.}
\answer{ex:06:3}{$\mathbf e_1$ si $\mu^2+s_1^2>s_2^2$; ici $1.01>0.25$, et le neurone apprend $\mathbf e_1$, bien que la dispersion selon $\mathbf e_2$ soit $25$ fois plus grande: la règle trouve le vecteur principal de la matrice de corrélation, et non de covariance.}
\answer{ex:06:4}{$V=Re^{-(t_c+L-t)/\tau_\gamma}$ entre la lampe et le jus, donc $\dot V=V/\tau_\gamma$; pic d'aire $Re^{-L/\tau_\gamma}$: $0.61R$ et $0.78R$.}
\answer{ex:06:5}{Rapport signal sur bruit $1/(N+1)$, nombre d'essais $\propto N+1$: $5\cdot10^5$ essais $\approx1$~mois et $5\cdot10^{11}$ essais $\approx8\cdot10^4$~ans.}
\answer{ex:06:6}{(a)~$k_2=1/\tau_d$, $k_1=1/\tau_d-1/\tau_\gamma$. (b)~Fausse erreur $-(V/\tau_\gamma)\,\varepsilon/(1+\varepsilon)$, $\varepsilon=\tau_d/\tau_\gamma$, d'où $\tau_d\le0.105\,\text{s}$.}
\answer{ex:06:7}{$0.46$, $0.90$, $0.99$, $0.95$, et $0.70$ pour $d=3\,\text{s}$. Les points tombent sur une même courbe en fonction de la fraction de trace $F=(1-e^{-T_{\text{cue}}/\tau_e})e^{-d/\tau_e}$: apprentissage pour $F\gtrsim0.1$, choix aléatoire pour $F<10^{-2}$.}
\answer{ex:06:8}{$\eta^*=1/\bigl(2\sigma^2(N+2)\bigr)$, en $N+2$ essais; avec $m$ neurones fluctuants sur $N$, en $N(m+2)/m\ge N+2$: la parcimonie n'aide pas.}

\answer{ex:07:1}{$V=\kappa p/\bigl(\kappa p+(1-\kappa)(1-p)\bigr)$: $0.059$, $0.20$, $0.50$; la médiane vaut $0$, alors que le point~\eqref{eq:expectile} varie continûment avec $\kappa$.}
\answer{ex:07:2}{$p=1/3$, $x=0.75$, $\sigma_r=0.35$, $V(0.2)=0.083$.}
\answer{ex:07:3}{$V=\sqrt\kappa/(\sqrt\kappa+\sqrt{1-\kappa})$, de $0.25$ à $0.75$; dispersion $\propto\sqrt\eta$; pour les gouttes $V=0.25+0.5\kappa$, les distributions diffèrent de $0.05$ aux neurones extrêmes, il faut $\eta\lesssim0.01$.}
\answer{ex:07:4}{(a)~$J^*=2\sqrt{C\bar R}$. (b)~Surcoût $\bigl(k^{1/2}+k^{-1/2}\bigr)/2-1$: $6\,\%$ pour $k=2$ et $25\,\%$ pour $k=4$; une précision \og à un facteur deux près\fg{} suffit.}
\answer{ex:07:5}{Pic d'aire $R(e^{-1}-e^{-2})\approx0.23R$, soit $2.3\,\text{s}$ de rampe; si la dopamine signale $V$, il n'y a pas de pic, le niveau est simplement multiplié par $e$ d'un seul coup.}
\answer{ex:07:6}{Un événement déplace le poids de $\propto A\tau_f/(\tau_e+\tau_f)$, le niveau donne une erreur $s\tau_f$: $9\,\text{s}\le\tau_f\le17\,\text{s}$.}
\answer{ex:07:7}{Au départ, $\Delta P=0.80$, $M=0.90$. (a)~$\Delta P=0.74$ ($-8\,\%$), $M=0.43$ ($-52\,\%$); (b)~$\Delta P=0.40$ ($-50\,\%$), $M=0.80$ ($-11\,\%$): double dissociation.}
\answer{ex:07:8}{$N_{\text{eff}}\to2+\rho^2N\approx10^6$ essais: $10^4$ fois moins qu'avec un seul fil, mais une fuite de $1\,\%$ coûte encore un million d'essais.}

\answer{ex:08:1}{(a)~$\approx1.7\cdot10^6$. (b)~$g=\frac{0.9}{\sqrt{0.19}}\,\sigma_{\text{après}}/\sigma_{\text{avant}}\approx16.5$: la réponse ne dépend que du rapport des bruits.}
\answer{ex:08:2}{(a)~Valeurs propres $-(1\pm g\beta)/\tau$; la différence s'amortit en $\tau/(1-g\beta)$: $40\ms$ pour $g\beta=0.5$ (différence des entrées amplifiée trois fois), $200\ms$ pour $0.9$ ($19$ fois). (b)~$0.905$ et $1.105$; entre les deux frontières, seule l'entrée forte gagne.}
\answer{ex:08:3}{$P(k)=e^{gu_k/\sigma}/\sum_le^{gu_l/\sigma}$, c'est-à-dire~\eqref{eq:softmax}; $g=10\ln9\approx22$.}
\answer{ex:08:4}{$\bar\Delta=1/3$, $g^*\approx3\ln(1/h)$: $21$, $14$, $9$ contre $16$--$23$, $11$ et $8$ dans le modèle.}
\answer{ex:08:5}{$g^*$ va de $16$--$23$ pour $h=0.001$ à $6$--$8$ pour $h=0.2$; l'estimation $3\ln(1/h)$ est juste à un facteur $1.5$ près pour $h\le0.05$; pour $h\gtrsim\eta/10$, la règle $Q\leftarrow Q+\eta(R-Q)$ n'a pas le temps d'assimiler ce que l'exploration a trouvé, et $g^*$ reste bloqué vers $8$.}
\answer{ex:08:6}{$T_p=\tau\ln9=44\ms$ pour $g_{\text{lo}}=0$ (modèle: $44\ms$) et $70\ms$ pour $g_{\text{lo}}=0.25$: une bouffée de $2$--$3\tau$ avec un gain proche de zéro.}
\answer{ex:08:7}{(a)~Meilleur $g$ constant: $0.925$ ($g=8$), oracle: $0.929$ ($16/8$); il n'y a presque rien à gagner. (b)~Par exemple, des époques calmes à récompenses rares ($p\in[0,0.3]$) et des époques changeantes avec $h=0.1$: $0.850$ contre $0.872$; l'ajustement en récupère environ un quart pour $\eta=0.2$ et presque tout pour $\eta=0.05$.}

\answer{ex:09:1}{(a)~$P_\infty=0.2$, mémoire $20\,\text{s}$. (b)~$P_\infty=1$, mémoire $1\,\text{s}$. (c)~La mémoire est proportionnelle à l'amplitude du bruit: six fois.}
\answer{ex:09:2}{$P(t)=P_\infty\coth(t\sqrt{q/R})$. (a)~$\frac12\sqrt{R/q}=5\,\text{s}$, deux fois moins que la mémoire. (b)~$t=\frac12\ln21\,\sqrt{R/q}\approx15\,\text{s}$: c'est le temps pendant lequel \nt{ACET} doit rester élevé.}
\answer{ex:09:3}{Variance $qT/2+R/(2T)$; $T^*=\sqrt{R/q}$, variance $\sqrt{qR}=P_\infty$, comme pour le filtre~\eqref{eq:kalman}; croissance d'un facteur $(\kappa+1/\kappa)/2$: $1.25$ et $5.05$.}
\answer{ex:09:4}{$a>a_w=1-\theta/(mw)$: $0.15$ pour $w=0.041$, $0.22$ pour $w=0.045$.}
\answer{ex:09:5}{Les erreurs apparaissent vers $M\approx40$--$60$; pour $M=80$, $1.9$ neurone superflu; pour $M=100$, part du motif $0.86$; l'interférence des autres motifs a une variance $\approx(M-1)p(1-p)/N$, d'où $M_{\max}\approx80$, $M_{\max}\propto N/p$.}
\answer{ex:09:6}{Par exemple $\eta(a)=\eta\min\bigl(1,\max(0,(a-0.3)/0.6)\bigr)$. Avec $\eta a$, les trois neurones étrangers sont entrés dans le motif; avec $\eta(a)$, aucun, et l'écriture pour $a\ge0.9$ est la même (part $1.00$).}
\answer{ex:09:7}{(a)~$40$ répétitions; avec la fonction $\eta(a)$ de l'exercice~\ref{ex:09:6}, jamais. (b)~$200$ répétitions, $10$ nuits.}

\answer{ex:10:1}{Bus $\approx10^{12}$ bit/s ($11.4$ bits par impulsion), commande $\approx40$ bit/s; $2.5\cdot10^{10}\,\text{s}$, soit environ huit cents ans.}
\answer{ex:10:2}{$a>1/3$ (pour $a=0$, oscillations croissantes à $\approx67$~Hz). Non: la stabilité dépend du produit $g(1-a)$.}
\answer{ex:10:3}{Moyenne $\sigma^2R'x_j$, variance $(N+1)\,x_j^2\sigma^4R'^2$; rapport $1/(N+1)$, il faut $n\approx z^2(N+1)\propto N$ essais.}
\answer{ex:10:4}{$L_v\in[32.9,\,47.9]\ms$: délai de $35.4\ms$, fenêtre de $\pm7.5\ms$; fausses liaisons $2\Delta_{\max}r_ar_v\approx0.38$ par seconde.}
\answer{ex:10:5}{$0.098$ contre $0.180$ pour le meilleur $\alpha=0.2$: $45\,\%$ de moins.}
\answer{ex:10:6}{$n=\ln\bigl(\ln1.5/(0.6g)\bigr)/\ln(1-\alpha)$: $24$ et $4$; pour $g=2$, $11$ et $2$.}
\answer{ex:10:7}{$S'=S\bigl(1-4\eta\sigma^2+4\eta^2\sigma^4(G+2)\bigr)$, meilleur $\eta\sigma^2=1/(2(G+2))$; $\approx10^6$ essais, soit environ un mois.}

\answer{ex:11:1}{(a)~$T=k^2/\bigl(\Phi(\Delta I/I)^2\bigr)=9\,\text{s}$. (b)~$45$ capteurs, tache $\approx7\times7$, résolution $\approx7$ fois moins bonne.}
\answer{ex:11:2}{(a)~$2$--$3.3$ bits par impulsion, $22$--$34\,\%$ de la limite ($9.1$--$9.8$ bits). (b)~$\log_2\binom{400}{20}\approx111$ bits; en moyenne $1$ type commun.}
\answer{ex:11:3}{(a)~$\sigma/9\le I\le9\sigma$, $1.9$ ordre de grandeur. (b)~$6$ et $7$.}
\answer{ex:11:4}{(a)~$f<343/0.44\approx780$~Hz. (b)~$625$ impulsions, $125$ neurones.}
\answer{ex:11:5}{$21$ bits par événement, $4.8\cdot10^5$ événements/s, $7$ événements par pixel du bord: $136$ pixels/s. Caméra ordinaire: $1.25$ image/s.}
\answer{ex:11:6}{En logarithmes, $\approx1.0\,\text{s}$ dans les deux sens (théorie $0.96\tau$); linéairement, $0.2\,\text{s}$ vers le haut et $3.0\,\text{s}$ vers le bas (théorie $0.18\tau$ et $3.0\tau$).}
\answer{ex:11:7}{(a)~$\ln I$ suit une loi logistique de médiane $\ln\sigma$ et d'échelle $1$: écart $\pi/\sqrt3$ en logarithmes naturels, $0.79$ ordre de grandeur. (b)~$r=r_{\max}\Phi_I(I)^2$.}

\answer{ex:12:1}{$0.60$ (sans corrélation, ce serait $0.18$), limite $\sqrt{c/(rT)}\approx0.58$: cent neurones ne distinguent que \og oui/non\fg{}.}
\answer{ex:12:2}{(a)~$10^8$ impulsions, $\approx10\,\text{mJ}$: $0.3\,\%$ des $3$~J du cerveau entier. (b)~$3$~J, $\approx300$ fois plus.}
\answer{ex:12:3}{(a)~$\theta_A\in[2,3)$, $\theta_B\in[1,2)$; l'autre figure donne au plus $2$ et $1$. (b)~$10\cdot96^2\approx9.2\cdot10^4$.}
\answer{ex:12:4}{(a)~$T_d=0.85\,\tau_a=0.34\,\text{s}$. (b)~$T_d\propto\tau_a$ et ne dépend pas de $I$: $\tau_a\approx3.5\,\text{s}$ (simulation: $3.1\,\text{s}$).}
\answer{ex:12:5}{$p\approx0.17$. La proportion baisse lentement: $0.90$, $0.88$, $0.86$, $0.84$, $0.82$, $0.79$ pour $N=8$, $12$, $24$, $48$, $96$, $192$.}
\answer{ex:12:6}{Les dix essais sont sans erreur pour $\tau_{\text{tr}}=20$, $50$ et $200\ms$ (à $30\ms$, un essai sur dix échoue); à $5$--$10\ms$, environ $0.5$; à $0.5$--$2\,\text{s}$, la qualité tombe à $0.86$. Borne supérieure: la trace doit s'éteindre pendant la pause, $e^{-0.3\,\text{s}/\tau_{\text{tr}}}\ll1$.}
\answer{ex:12:7}{Un seul délai ne suffit pas: la fenêtre de $\pm5\ms$ ne couvre pas $\Delta x/v$ de $5$ à $20\ms$. Deux branches d'inhibition avec $5<d_1\le10\ms$ et $15\le d_2\le d_1+10\ms$.}
\answer{ex:12:8}{\texttt{two\_stage}: $\Delta=2$, égalité dès une inversion, changement de réponse à partir de deux. Compteur de uns: une seule inversion; $0.57$ pour $p=0.1$.}

\answer{ex:13:1}{(a)~$q=100^{1/99}\approx1.048$, pas $\approx4.8\%$; plage $\approx2.2\cdot10^3$ ($67$~dB). (b)~Pas de $22f_1$, soit $220\%$ à $10f_1$.}
\answer{ex:13:2}{$P_{\text{échec}}\approx1.8\cdot10^{-4}$ et $2.8\cdot10^{-16}$: environ $300$ échecs par jour pour $S=2$ et $5\cdot10^{-10}$ par jour pour $S=4$.}
\answer{ex:13:3}{(a)~$H(s)=eF_1T_c/(1+sT_c)^2$, filtre du second ordre de fréquence de coupure $1/(2\pi T_c)\approx3.2$~Hz. (b)~$r\approx22$~imp./s ($\approx20$ d'après le premier harmonique).}
\answer{ex:13:4}{(a)~Pour $Gd=1/e$, la racine $s=-1/d$ est double, et pour aucun autre $G$ la racine la plus à droite n'est plus à gauche. (b)~$d=30\ms$: $G\approx12$~s$^{-1}$, $\tau=30\ms$; $d=150\ms$: $G\approx2.5$~s$^{-1}$, $\tau=150\ms$; la constante de temps est égale au retard.}
\answer{ex:13:5}{(a)~$E\approx0.427$, $T\approx1.70\,\tau_a=0.68$~s (le modèle avec $\tau=20\ms$ donne $0.84$~s). (b)~L'équation ne contient pas $I$; le rythme existe pour $b>\beta-1$.}
\answer{ex:13:6}{Pour $b=0.8$, pas de rythme; pour $b=1.05$, $1.2$, $1.5$, $2$, $4$, $T\approx2.73$, $1.74$, $1.19$, $0.84$, $0.45$~s; pour tout $I$, $T=0.840$~s: l'entrée ne change que l'amplitude, pas le tempo.}
\answer{ex:13:7}{La substitution $s=u-a$ donne la vitesse d'amortissement maximale $a+1/d$ pour $G=e^{-1-ad}/d$; $a\ge33.3$~s$^{-1}$, $G=e^{-2}/d\approx4.5$~s$^{-1}$; plus la co-contraction est forte (plus $a$ est grand), plus il faut prendre $G$ petit.}
\answer{ex:13:8}{Problème ouvert. Par exemple, un seuil dans l'alimentation de la fatigue, $b[r_i-r_0]_+$, donne $I_{\min}=\beta b r_0/(b-\beta+1)\approx0.53$ pour $b=4$, $r_0=0.2$, et une période allant de $1.8$~s pour $I=0.6$ à $0.52$~s pour $I=3$.}

\answer{ex:14:1}{(a)~De $0.17$~s et de $1.5$~s. (b)~À une distance de moins de $11$~cm.}
\answer{ex:14:2}{Un toucher sans douleur ne passe pour aucun $\mu$ tant que $g\le\beta w_{q\mu}/w_{z\mu}=5$; un toucher $\mu=1$ passe pour $g>6$: une forte sensibilisation rend douloureux un contact ordinaire.}
\answer{ex:14:3}{(a)~$g^*=(1-k_sC)/(1-k_sA)$, stable pour $k_sA<1$. (b)~Sans toucher, pour $\nu\ge2$; avec toucher, dès $\nu\ge1.7$.}
\answer{ex:14:4}{$V(t)=-Pe^{-(T-t)/\tau_\gamma}$. (a)~$-Pe^{-T/\tau_\gamma}\approx-0.37P$. (b)~$+Pe^{-(T-t_e)/\tau_\gamma}\approx+0.61P$; elle croît avec $t_e$ jusqu'à $+P$ et enseigne à éviter la douleur.}
\answer{ex:14:5}{Pour $k_s=4$, pas de verrouillage (un second état stable existe pour $k_s\ge4.58$); $T_{\min}\approx4.35$, $1.40$, $0.65$, $0.32\,\tau_s$ pour $k_s=4.6$, $5$, $6$, $8$.}
\answer{ex:14:6}{$w_{q\nu}=4/9\approx0.444$, $q_0=2/9\approx0.222$, $w_{q\mu}=49/45\approx1.089$; un toucher sans douleur est sans danger jusqu'à $g=49/9\approx5.4$.}
\answer{ex:14:7}{(a)~Seuil $\nu_c(g)=\theta/g$ si $\nu_c\ge q_0/w_{q\nu}$, sinon $(\theta+\beta q_0)/(g+\beta w_{q\nu})$; $g^*\approx2.45$, $1.0$, $0.25$. (b)~Question ouverte.}

\answer{ex:15:1}{(a)~$2\cdot10^5$ par neurone, $6\cdot10^{12}$ pour le circuit. (b)~$\approx8$~bit/s par fil, au moins $2.5\cdot10^4$~s $\approx7$~h.}
\answer{ex:15:2}{Facteur $2$: corrélation entre voisins $0.943$, $\lambda$ de $4.18$ à $7.3\cdot10^{-4}$, rapport $\approx5.7\cdot10^3$; facteur $4$: corrélation $0.8$, rapport $\approx94$.}
\answer{ex:15:3}{(a)~$\lambda=0.988$ et $0.808$: $82$ et $4.7$ mouvements. (b)~$x_s^*=0.690$, $x_f^*=0.172$, somme $0.862$.}
\answer{ex:15:4}{(a)~$116$ mouvements. (b)~$19$. Un modèle à un seul processus, à somme nulle, n'apporte aucune économie.}
\answer{ex:15:5}{(a)~$G=50$~s$^{-1}$, contre une limite de $10.5$~s$^{-1}$ sans modèle. (b)~Non: pour $G=50$~s$^{-1}$, le retard critique est $d_c\approx103\ms$; pour $d=150\ms$, il faut $\hat k>0.445k$ (pour tous $G$ et $d$, $\hat k>k/2$).}
\answer{ex:15:6}{L'erreur passe sous $0.5\%$ en $6$--$30$~s, avec un plancher de $\approx(6$--$9)\cdot10^{-4}$ pour $7.5\cdot10^{-4}$ avec les meilleurs poids; pour $\eta=50$, les poids diffèrent encore des meilleurs de jusqu'à $0.2$ le long des directions mal conditionnées de $C$ (exercice~\ref{ex:15:2}).}
\answer{ex:15:7}{$\mathbf B=A\mathbf K$; $q_f/R\approx0.035$, $q_s/R\approx8.6\cdot10^{-4}$; avec $4q_s$, $B_f\approx0.088$, $B_s\approx0.047$: le processus lent apprend plus de deux fois plus vite, le rapide plus lentement.}

\answer{ex:16:1}{Le sang est un circuit RC avec $\tau=t_{1/2}/\ln2=14.4$~min; rapport du maximum au minimum $2^{T/t_{1/2}}=64$; harmonique fondamentale atténuée à $0.55$; rapport $2$ pour $t_{1/2}=T=60$~min.}
\answer{ex:16:2}{$\varphi^*=\arcsin\bigl((\omega-\omega_0)/K_\varphi\bigr)\approx41$~min ($\omega-\omega_0\approx0.047$~rad/j), temps de relaxation $1/(K_\varphi\cos\varphi^*)\approx3.9$~j; après un décalage de $+6$ et de $-6$~h, $7.7$ et $7.9$~j.}
\answer{ex:16:3}{$K_c(n)=\sec^n(\pi/n)$: $8$ et $2.37$, erreur de $11\%$ et $30\%$; quand $n\to\infty$, $K_c\to1$ et l'erreur $\to50\%$.}
\answer{ex:16:4}{Les commandes extrêmes (accélérateur $80$, frein $0$) donnent $f=120$ en $0.76$~s; \og accélérateur seul\fg{}: $2.33$~s, trois fois plus.}
\answer{ex:16:5}{$3\arctan(\omega\tau)+\omega d=\pi$, $K_c=(1+\omega^2\tau^2)^{3/2}$: $K_c=8$, $3.54$, $2.50$, $1.76$ et périodes $3.63$, $5.46$, $6.86$, $9.27\,\tau$; à $0.9K_c$ les oscillations s'amortissent, à $1.1K_c$ elles s'installent.}
\answer{ex:16:6}{$0.60\bar c$, $0.87\bar c$, $0.90\bar c$, limite $0.91\bar c$; à concentration constante, la réponse est nulle: un tel destinataire ne lit que les bouffées.}
\answer{ex:16:7}{Question ouverte. Avec la rétroaction, le rythme n'existe que pour $K>8$; sa période est proportionnelle au $\tau$ des étages et dépend à peine de $K$ ($3.64\,\tau$ pour $K=8.5$, $4.23\,\tau$ pour $K=40$): changer $K$ d'un facteur une fois et demie décale la période de $2$ à $4\%$, moins que l'erreur de mesure. Ce qui tranche, c'est l'ouverture de la boucle (une hormone finale constante à l'entrée du premier étage supprime un tel rythme) ou la réponse à un bref ajout d'hormone à différentes phases du cycle.}

\answer{ex:17:1}{$0.2$~W contre $81$~µW; le flux brut exigerait au plus $50$~pJ/bit.}
\answer{ex:17:2}{Pour $c=0.1$: $5.6$, $8.7$, $9.2$~bit/s, limite $9.3$~bit/s; pour $c=0$ et $N=1000$: $65$~bit/s.}
\answer{ex:17:3}{$B=\log_2N+P\log_2P+(1-P)\log_2\frac{1-P}{N-1}=4.87$, $4.37$, $3.29$~bits par caractère, soit $7.3$, $6.6$, $4.9$~bit/s; pour $10$~bit/s, il faut $2.3$~caractères par seconde.}
\answer{ex:17:4}{Variance de l'erreur $2b/(Nm^2T)+\sigma_\xi^2$ par composante; les contributions sont égales pour $N\approx90$; limite $\tfrac12\log_2(1+0.25/0.04)\approx1.4$~bit par composante et par fenêtre.}
\answer{ex:17:5}{La paire converge si $\eta_bd^2+\eta_db^2<2$, soit à peu près pour $\eta<1$: à $0.8$ elle converge, à $1.1$ oscillations entretenues de période $2$, à $1.5$ oscillations non périodiques, à $2$ divergence; il faut séparer les vitesses des élèves.}
\answer{ex:17:6}{Seuil $\approx4.4\sigma$; $102$~kbit/s, $0.10$~mW contre $205$~mW pour le flux brut, soit $2000$ fois moins; seuls $5.7\cdot10^{-5}$ des échantillons saturent (plus de $3$ impulsions).}
\answer{ex:17:7}{Question ouverte. Environ $46$~bit/s pour $\text{SNR}\approx0.9$ et environ $330$~bit/s avec un bruit indépendant. Si c'est un bruit semblable au signal qui gêne, l'information extraite sature quand le nombre de canaux augmente; si c'est la méconnaissance du code, elle augmente avec un meilleur décodeur (par exemple, qui tient compte des instants des impulsions).}

\answer{ex:18:1}{(a)~$0.9949\le w\le1.0049$, soit $|1-w|\lesssim0.005$. (b)~$K\ge400$ synapses.}
\answer{ex:18:2}{$(\omega_R\tau_w)^2=\sqrt{\alpha(\alpha+2+2\varrho)}/\varrho-1$, $\varrho=\tau_m/\tau_w$; $f_R\approx11.1$~Hz; $|Z(\omega_R)|/|Z(0)|=4.6$.}
\answer{ex:18:3}{(a)~$f=1/\bigl(\tau\ln\frac{\mu}{\mu-\theta}\bigr)$; pour $1$~Hz, il faut $\mu/\theta-1\approx3.7\cdot10^{-44}$. (b)~$T=\pi\tau/\sqrt\mu$, $f\approx3.2$~Hz: $f\propto\sqrt\mu$.}
\answer{ex:18:4}{L'intégrateur décharge pour $f\lesssim17.7$~Hz (filtre passe-bas), le résonateur seulement dans la bande $5.5$--$22$~Hz (filtre passe-bande).}
\answer{ex:18:5}{(a)~$T_{\min}=\frac{\tau}{w-1}\ln\frac{0.5}{0.3}\approx10.2\ms$. (b)~$B>w-1-0.2=0.3$.}
\answer{ex:18:6}{Constante de temps du réglage $\tau/(\eta\langle r^2\rangle)$, $\eta=\tau\ln10/60\,\text{s}\approx3.8\cdot10^{-3}$; pour $I\neq0$, le poids se décale: il faut retrancher de $e$ une copie de la commande $I/\tau$.}
\answer{ex:18:7}{Question ouverte. Le résonateur ne gagne qu'un facteur $1.35$ à $11$~Hz et perd un facteur $17$ à $1$~Hz; le principal avantage de la reconfiguration est le choix de la bande par le seuil (exercice~\ref{ex:18:4}).}

\answer{ex:19:1}{(a)~Une synapse donne $6.7\mV$ ($R=66.7$~M$\Omega$): pour atteindre le seuil tout court, il en faut au moins $4$ (trois ne suffisent qu'asymptotiquement); $8$ en $10\ms$; $14$ en $5\ms$. (b)~$0.1\ms\ll\tau_m$, correction due à la fuite $\sim0.25\%$.}
\answer{ex:19:2}{(a)~$1.75\cdot10^{10}$ (un quart des mélangeurs répond à tout), $4.4\cdot10^9$ et $0.06$ (pour $k=40$, presque jamais de réponse). (b)~D'un facteur $10^3$: $2\cdot10^5$ contre $200$.}
\answer{ex:19:3}{(a)~$C(2n,n)=2^{2n-1}$ par symétrie des coefficients binomiaux. (b)~$0.93$ et $0.09$; largeur de la transition $\sim\sqrt{2n}$ en $P$, relative $\sim1/\sqrt n$.}
\answer{ex:19:4}{Avec une seule dendrite, $V_{\text{corps}}<\kappa E_E<\theta$ quel que soit le nombre d'entrées; pour $g_0=0.5g_L$, il faut $n\ge3$ sur chaque dendrite.}
\answer{ex:19:5}{$I\ge C\sigma_V/\sigma_t=15$~nA, soit $150$ synapses synchrones, ce qui est irréaliste; au petit neurone, $1$~nA suffirait, et avec une synapse de $4$~nA la gigue vaut $5$~µs.}
\answer{ex:19:6}{Maximum à l'extrémité proche: $0.03\,\tau_m$ et $0.97$ ($L=0.2$); $0.32\,\tau_m$ et $0.66$ ($L=1$); $0.79\,\tau_m$ et $0.33$ ($L=2$). L'estimation~\eqref{eq:dendriteload} par la capacité totale n'est valable que pour $L\ll1$.}
\answer{ex:19:7}{Question ouverte. À $m$ fixé, le temps de réponse croît linéairement avec le nombre $P$ d'associations mémorisées; à fraction fixée $m=\varphi n$, il ne dépend plus de $n$, et la lenteur d'un grand neurone découle de la sommation de nombreuses entrées, et non de la taille en tant que telle.}

\answer{ex:20:1}{$t_{\text{v}}=\tau_r\ln\frac{1-L}{1-H}=16.8$~h, $t_{\text{s}}=\tau_d\ln\frac HL=5.8$~h, période $22.5$~h; c'est l'oscillation circadienne des seuils qui amène l'oscillateur à $24$~h: un verrouillage de phase~\eqref{eq:pll}.}
\answer{ex:20:2}{(a)~$L=0.111$; avec $L=0.092$, le sommeil durerait $9.6$~h. (b)~De $22\%$ ($1/0.82=1.22$), et non de $18\%$.}
\answer{ex:20:3}{$H=\frac{p-1}{p-1/q}=0.623$, $L=H/q=0.093$, où $p=e^{16/\tau_r}$, $q=e^{8/\tau_d}$. Après un réveil au bout de $4$~h, la phase est décalée pour toujours (endormissement $7.2$~h plus tôt); avec des seuils oscillants, elle revient en deux ou trois jours.}
\answer{ex:20:4}{(a)~$r_\pm=\frac12\bigl(1\pm\sqrt{1-4k/w}\bigr)=0.947,\ 0.053$; $a_{\text{h}}=3.51$, $a_{\text{b}}=0.49$. (b)~Activité $\approx0.33$~s, silence $\approx0.59$~s, période $\approx0.93$~s, fraction d'activité $0.36$.}
\answer{ex:20:5}{$a_{\text{h}}/r_+<b<a_{\text{b}}/r_-$: $3.71<b<9.20$. \nt{ACET} fait descendre $b$ sous $3.71$: l'intersection passe sur la branche haute, et c'est l'état actif ($r\approx1$) qui est stable.}
\answer{ex:20:6}{$4.9$, $6.8$, $8.8$, $5.5$, $8.9$~h pour $D=24$, $32$, $40$, $48$, $64$: la durée est fixée par la phase circadienne de l'endormissement, et non par la dette.}
\answer{ex:20:7}{Après B, l'erreur sur A vaut en moyenne $0.51$ (de $0.19$ à $1.08$ sur $20$ tirages; une réponse nulle donnerait $1$). Avec entrelacement, les deux erreurs sont inférieures à $10^{-4}$.}
\answer{ex:20:8}{Question ouverte. Une mise à l'échelle uniforme conserve les rapports des poids, mais ne conserve les réponses d'un neurone à seuil que si le seuil est mis à l'échelle du même facteur; les répétitions entrelacées modifient les rapports. La stabilité des gros contacts contredit une mise à l'échelle purement uniforme.}

\answer{ex:21:1}{(a)~$T=\tau_{\text{rec}}\ln\frac{1}{1-x^*}$: $1.84$~s et $3.68$~s. (b)~De $3.36$ à $4.24$~s: la sensibilité $\dd T/\dd x^*=\tau_{\text{rec}}/(1-x^*)=80$~s; quand $x^*\to1$, les salves deviennent irrégulières.}
\answer{ex:21:2}{(a)~$8.6$~W, comme dans~\eqref{eq:energy}. (b)~$205$~Mbit/s, soit $0.2$~W: $2\cdot10^3$ fois plus que la culture ($10^{-4}$~W).}
\answer{ex:21:3}{$x_{\text{st}}=1/(1+Ur_b\tau_{\text{rec}})<x^*$ pour $r_b>(1/x^*-1)/(U\tau_{\text{rec}})=0.28$~Hz si $x^*=0.9$ et $0.025$~Hz si $x^*=0.99$; la force de la synapse tombe à $x_{\text{st}}$, c'est-à-dire de $10\%$ et de $1\%$.}
\answer{ex:21:4}{Les $u(t-k)$ sont orthonormés, $m_k=\|P_Vu(t-k)\|^2$, où $P_V$ est le projecteur sur l'enveloppe linéaire des sorties, de dimension $N$; par l'inégalité de Bessel, $\sum_k m_k\le N$.}
\answer{ex:21:5}{$\mathrm{MC}\approx9$--$11$ avec $\tanh$ et $15$--$18$ pour le réservoir linéaire (trois graines), tous deux $\le30$: la non-linéarité se paie en mémoire.}
\answer{ex:21:6}{(a)~$n\le102$. (b)~De $5.6$.}
\answer{ex:21:7}{Problème ouvert. Contrôle clé: geler la lecture et vérifier si l'amélioration subsiste après une pause sans stimulation; sinon, c'est l'état qui a changé, et non les poids.}

\answer{ex:22:1}{(a)~$32\%$ et $100\%$. (b)~$R_\uparrow-R_\downarrow\ge2\sqrt{2000\,\text{Hz}/1\,\text{s}}\approx89$~Hz, soit seulement $4.5\%$ de la fréquence totale.}
\answer{ex:22:2}{(a)~$n\ge4(p_1q_1+p_2q_2)/\Delta^2\approx560$. (b)~$\approx56$.}
\answer{ex:22:3}{(a)~Bilan détaillé: $\rho PK$ est symétrique. (b)~La proportion de ratés est la moyenne harmonique $1/\langle1/P\rangle=0.18$, contre $0.5$ sans rétroaction.}
\answer{ex:22:4}{(a)~Un parallélogramme d'aire $4h^2=0.04$ (en tenant compte de la borne sur $p$, numériquement $0.045$). (b)~$\approx0.18$.}
\answer{ex:22:5}{$0.49$: une partie des cultures part dans une région où le raté est presque certain et y erre. Avec la borne: $1.00$; sur un ensemble borné, la densité $\propto1/P$ se concentre dans la région absorbante.}
\answer{ex:22:6}{(a)~Il faut $\ge5$ niveaux; $8$ électrodes suffisent. (b)~Non: il faudrait $r_{\max}\ge25/T=100$~Hz.}
\answer{ex:22:7}{Problème ouvert. Le temps de la recherche aléatoire croît avec $d$ à peu près comme le rapport des volumes $(L/h)^d$, celui de la recherche au signe de l'erreur, linéairement en $d$. Une expérience d'inversion départage les hypothèses: après un raté, un stimulus prévisible mais fort; après un renvoi, un stimulus imprévisible mais faible; il faut quelques dizaines de cultures par groupe.}


\chapter{Annexe expérimentale}
\label{sec:supplement}

Pour chaque chapitre, on trouvera ici ses principales affirmations et ce sur quoi elles reposent: l'expérience où elles ont été mesurées, ce qui a été mesuré exactement et où cela a été publié. Le statut, dans la colonne de droite, indique la solidité de l'affirmation: \emph{mesuré}, il existe une expérience directe; \emph{théorie}, conséquence mathématique établie dans le chapitre ou dans un travail classique; \emph{modèle}, résultat d'un calcul sur le modèle du livre (scripts dans le dossier \texttt{models/}); \emph{estimation}, ordre de grandeur sans mesure directe; \emph{hypothèse}, explication plausible que l'expérience n'a pas encore démontrée. Là où il n'existe pas d'expérience, nous le disons.

\section{Chapitre~\ref{sec:neurontypes}. Types de neurones}

Les nombres du chapitre reposent sur des comptages directs de cellules dans le cerveau humain, là où ils existent; la règle \og un neurone, un neurotransmetteur\fg, sur des atlas cellulaires et sur des expériences qui en ont aussi trouvé les exceptions; le rôle de \nt{DOPA}, sur des enregistrements d'impulsions; et les rôles des autres canaux de diffusion, sur des hypothèses.

{\footnotesize
\setlength{\tabcolsep}{3pt}\renewcommand{\arraystretch}{1.15}
\begin{longtable}{>{\raggedright}p{0.5cm}>{\raggedright}p{4.0cm}>{\raggedright}p{5.6cm}>{\raggedright}p{2.3cm}>{\raggedright\arraybackslash}p{1.7cm}}
\toprule
N\textsuperscript{o} & Affirmation & Expérience et ce qui est mesuré & Source & Statut \\
\midrule\endfirsthead
\toprule
N\textsuperscript{o} & Affirmation & Expérience et ce qui est mesuré & Source & Statut \\
\midrule\endhead
1 & Le cerveau humain compte environ $8.6\cdot10^{10}$ neurones, et presque tous les neurones \nt{GLUT} sont les petits neurones du cervelet (tabl.~\ref{tab:types}) & Fractionneur isotrope: le tissu est dissous en une suspension homogène de noyaux, et l'on compte les noyaux porteurs du marqueur neuronal NeuN. Chez l'homme adulte, $86.1\pm8.1$ milliards de neurones; le cortex cérébral n'en contient que $19\,\%$, la masse principale est dans le cervelet & \cite{Azevedo2009} & mesuré \\
2 & Presque chaque neurone a un seul neurotransmetteur principal (proposition~\ref{prop:dale}), mais il y a des exceptions & Atlas transcriptomique du cerveau de la souris: environ $4\cdot10^6$ cellules, $5322$ groupes; chaque groupe reçoit un neurotransmetteur d'après les gènes de synthèse et d'empaquetage. Les exceptions sont mesurées directement: en tranches, les neurones \nt{DOPA} libèrent aussi du GABA par le même transporteur vésiculaire et inhibent les neurones du striatum; dans une partie des axones \nt{DOPA}, le glutamate et la dopamine sortent de terminaisons différentes d'un même fil (microscopie électronique et optogénétique) & \cite{Strata1999,Yao2023,Tritsch2012,Zhang2015} & mesuré \\
3 & Toute la colonne de la matrice des poids est de même signe~\eqref{eq:dalesign} & Déduction du chapitre à partir de la proposition~\ref{prop:dale} et du fait que presque tous les récepteurs du glutamate sont excitateurs et ceux du GABA inhibiteurs. Il n'existe pas de vérification directe de \og tous les destinataires d'un neurone reçoivent un poids de même signe\fg{} comme règle générale; contre-exemples: la co-libération de GABA par les neurones \nt{DOPA} (ligne~2) et les destinataires dotés de récepteurs de signes différents (ligne~11) & déduction du chapitre & théorie \\
4 & De trois quarts à quatre cinquièmes des neurones du cortex sont \nt{GLUT}, d'un cinquième à un quart sont \nt{GABA} & Immunomarquage du GABA dans des colonnes de $50$\,\textmu m de large traversant toute l'épaisseur du cortex, dans $10$ aires corticales du singe: dans $8$ aires, les neurones GABA forment environ $25\,\%$ de tous les neurones, dans le cortex visuel primaire moins de $20\,\%$ & \cite{Hendry1987} & mesuré \\
5 & Un neurone \nt{GLUT} a environ $10^4$ sorties & Stéréologie post mortem: dans le néocortex de jeunes hommes, $1.64\cdot10^{14}$ synapses (microscopie électronique, dissecteur) et environ $2.3\cdot10^{10}$ neurones. Le rapport donne $\sim7\cdot10^3$ synapses par neurone; en moyenne, le nombre d'entrées égale le nombre de sorties & \cite{Tang2001synapses,Pakkenberg1997} & mesuré \\
6 & La sortie d'un neurone \nt{DOPA} se ramifie en $10^5$--$10^6$ terminaisons; c'est une estimation d'après la couverture de la cible, pas un comptage & Marquage viral de la membrane de neurones \nt{DOPA} isolés du rat et reconstruction complète de l'axone: un neurone couvre de $0.45$ à $5.7\,\%$ (en moyenne $2.7\,\%$) du volume du striatum. C'est la couverture qui est mesurée, pas le nombre de terminaisons: celui-ci en est déduit en ordre de grandeur, sans comptage direct. Pour \nt{SERO} et \nt{NORA}, les nombres de terminaisons sont des estimations grossières & \cite{Matsuda2009} & mesuré \\
7 & Les neurones de diffusion sont peu nombreux: \nt{DOPA} $\sim5\cdot10^5$ (le principal des deux groupes, borne inférieure), \nt{SERO} $\sim3\cdot10^5$ (somme de deux comptages partiels), \nt{HIST} $\sim6\cdot10^4$ (tabl.~\ref{tab:types}, fig.~\ref{fig:typepie}) & Comptages stéréologiques chez l'homme: $550\,000$ neurones pigmentés dans la substance noire (sans l'aire tegmentale ventrale); $235\,000$ neurones dans le noyau dorsal du raphé (tous les neurones du noyau) et encore $125\,000$ neurones sérotoninergiques dans le tegmentum du pont; environ $64\,000$ neurones histaminergiques de l'hypothalamus. Pour \nt{NORA}, \nt{GLYC}, \nt{ACET} et \nt{ADRE}, pas de comptage direct: le tableau donne des estimations grossières en ordre de grandeur & \cite{Pakkenberg1991,Baker1990,Baker1991,Airaksinen1991} & mesuré \\
8 & Le glutamate dans la fente monte jusqu'à environ une millimole par litre et retombe en $\sim1\ms$ & D'après le déplacement d'un bloqueur compétitif à dissociation rapide des récepteurs NMDA pendant la transmission synaptique (culture de neurones d'hippocampe): pic de $1.1$\,mM, constante de décroissance de $1.2\ms$ & \cite{Clements1992} & mesuré \\
9 & Dans le cortex en activité, les courants excitateur et inhibiteur sont presque équilibrés, et le neurone reste près du seuil & Théorie: un réseau à connexions fortes atteint de lui-même le régime équilibré. Expérience: dans le cortex préfrontal du furet in vivo, pendant les états \og up\fg, le potentiel d'inversion du courant synaptique reste vers $-37$\,mV pendant des centaines de millisecondes, alors que la conductance varie de $\sim21$\,nS: excitation et inhibition montent et descendent ensemble & \cite{vanVreeswijk1996,Haider2006} & mesuré \\
10 & Les neurones \nt{DOPA} annoncent l'erreur de prédiction de la récompense~\eqref{eq:rpe} (fig.~\ref{fig:rpe}) & Impulsions des neurones \nt{DOPA} du singe: bouffée au jus inattendu; après apprentissage, bouffée au signal et aucune réponse au jus prédit; creux de fréquence quand le jus promis n'est pas donné. Dans l'expérience quantitative, la fréquence est proportionnelle à la différence entre la récompense actuelle et la moyenne pondérée des précédentes, mais seulement pour les erreurs positives. L'équation~\eqref{eq:rpe} elle-même est l'algorithme des différences temporelles & \cite{Schultz1997,Bayer2005,SuttonBarto2018} & mesuré \\
11 & Le signe et la vitesse sont fixés par le récepteur: ionotropes, fractions de milliseconde; métabotropes, beaucoup plus lents (proposition~\ref{prop:receptor}) & Brèves impulsions de glutamate sur des fragments de membrane de neurones d'hippocampe: le courant des récepteurs ionotropes monte (de $20$ à $80\,\%$) en $0.2$--$0.6\ms$ et retombe en $2$--$3\ms$. Le potentiel inhibiteur lent des neurones pyramidaux est supprimé par un bloqueur sélectif du récepteur métabotrope du GABA. Deux familles de récepteurs de la dopamine à action opposée; dans le striatum, les neurones à récepteur D1 ont besoin de dopamine pour renforcer les synapses, ceux à D2 pour les affaiblir & \cite{Colquhoun1992,DutarNicoll1988,Beaulieu2011,Shen2008} & mesuré \\
12 & Un canal de diffusion n'est pas un fil unique, mais un faisceau de quelques lignes (section~\ref{sec:buses}) & Marquage viro-génétique des neurones sérotoninergiques du noyau dorsal du raphé de la souris: les neurones qui projettent vers l'amygdale et vers le cortex frontal occupent des places différentes du noyau, se ramifient vers des régions complémentaires et répondent de façon opposée à un stimulus désagréable. Revue: des sous-groupes de neurones du locus coeruleus (\nt{NORA}) codent des processus différents & \cite{Ren2018,Poe2020} & mesuré \\
13 & \nt{NORA} fixe le gain $g$ & Hypothèse du gain adaptatif; modèle de réseau dans lequel les catécholamines augmentent la pente de la caractéristique de transfert. Pas de mesure directe de \og $g$ croît avec la noradrénaline\fg{} à l'échelle du réseau entier; mesuré: le diamètre de la pupille suit les impulsions des neurones du locus coeruleus du singe & \cite{AstonJones2005,ServanSchreiber1990,Joshi2016} & hypothèse \\
14 & \nt{ACET} commute \og écriture/lecture\fg{} et fixe la vitesse d'apprentissage & Hypothèse. Appui: dans des tranches de cortex olfactif du rat, les agonistes de l'acétylcholine ($100$\,\textmu M) affaiblissent les synapses internes \og de rappel\fg{} de plus de $60\,\%$, et les synapses d'entrée de moins de $15\,\%$ & \cite{Hasselmo2006,HasselmoBower1992} & hypothèse \\
15 & \nt{SERO} fixe le facteur d'actualisation $\gamma$; les neurones sérotoninergiques travaillent régulièrement, à quelques impulsions par seconde, et se taisent presque pendant l'une des phases du sommeil & Hypothèse \og sérotonine $=\gamma$\fg. Argument le plus fort: l'activation optogénétique des neurones sérotoninergiques du noyau dorsal du raphé de la souris réduit le nombre d'abandons prématurés de l'attente d'une récompense différée. La fréquence des impulsions et le silence pendant le sommeil à mouvements oculaires rapides sont mesurés (revue d'enregistrements) & \cite{Doya2002,Miyazaki2014,JacobsAzmitia1992} & hypothèse \\
\bottomrule
\end{longtable}}

\section{Chapitre~\ref{sec:glutcomputer}. Ce que calculent les neurones \nt{GLUT}}

Les énoncés logiques du chapitre sont des théorèmes sur les circuits à poids positifs ou nuls, démontrés dans le chapitre lui-même; l'expérience confirme le modèle de neurone et le fait que les circuits construits d'après ces théorèmes (détecteur de coïncidences, ligne à retard, groupe auto-entretenu, chaîne) se rencontrent réellement dans le cerveau.

{\footnotesize
\setlength{\tabcolsep}{3pt}\renewcommand{\arraystretch}{1.15}
\begin{longtable}{>{\raggedright}p{0.5cm}>{\raggedright}p{4.0cm}>{\raggedright}p{5.6cm}>{\raggedright}p{2.3cm}>{\raggedright\arraybackslash}p{1.7cm}}
\toprule
N\textsuperscript{o} & Affirmation & Expérience et ce qui est mesuré & Source & Statut \\
\midrule\endfirsthead
\toprule
N\textsuperscript{o} & Affirmation & Expérience et ce qui est mesuré & Source & Statut \\
\midrule\endhead
1 & Le neurone est un circuit RC avec seuil et remise à zéro~\eqref{eq:glutneuron} & On a injecté dans le corps de neurones pyramidaux du cortex du rat (tranches) un courant bruité semblable au flux synaptique du cerveau vivant. La dépendance de la fréquence moyenne à la moyenne et à la dispersion du courant est décrite par un neurone intégrateur à fuite, si l'on ajoute une lente adaptation de la fréquence. Les plages de $\tau$ et des retards sont données dans le chapitre sans référence expérimentale & \cite{Rauch2003} & mesuré \\
2 & Le modèle~\eqref{eq:glutneuron} est une simplification: les branches d'entrée émettent elles-mêmes des impulsions locales & Enregistrement directement sur les dendrites de neurones pyramidaux du cortex visuel de la souris in vivo: un stimulus visuel déclenche des impulsions dendritiques locales dépendant des récepteurs NMDA; les supprimer élargit l'accord du neurone sur l'orientation & \cite{Smith2013} & mesuré \\
3 & Fenêtre de coïncidence $\Delta_{\max}=\tau\ln\frac{w}{\theta-w}$~\eqref{eq:window}; pour $\theta=1.5w$, elle vaut $0.7\tau$ & Conséquence du modèle~\eqref{eq:glutneuron}. Une mesure directe d'une fenêtre de sommation étroite existe pour les neurones à inhibition rapide (chapitre~\ref{sec:gaba}, ligne~3 de son tableau) & déduction du chapitre & théorie \\
4 & Un réseau fait seulement de neurones \nt{GLUT} ne calcule que des fonctions monotones des entrées (proposition~\ref{prop:monotone}) & Démonstration dans le chapitre: itération à partir du silence avec un membre de droite non décroissant. C'est un énoncé sur le modèle; aucune expérience ne l'a vérifié sur un réseau vivant privé d'inhibition & déduction du chapitre & théorie \\
5 & Les organes des sens fournissent une paire de fils: certaines cellules répondent au renforcement du stimulus, d'autres à son affaiblissement & Rétine: dès les cellules bipolaires, le signal des cônes se sépare en canaux ON et OFF. Blocage pharmacologique des cellules bipolaires ON chez le singe: les canaux restent séparés jusqu'au cortex; après le blocage, la détection d'une augmentation de lumière se dégrade, mais pas celle d'une diminution & \cite{Schiller1992} & mesuré \\
6 & Avec le code à deux fils, un réseau de neurones \nt{GLUT} calcule n'importe quelle fonction logique de façon asynchrone, sans cadence commune, au prix de deux fois plus de neurones (proposition~\ref{prop:dualrail}, \eqref{eq:dualrail}, fig.~\ref{fig:dualxor}) & Le code à deux fils et la détection de fin de calcul par le signal lui-même sont une technique standard des circuits asynchrones insensibles aux retards. Le circuit OU exclusif à six neurones est construit dans le chapitre. Aucune expérience ne montre que le cerveau calcule les fonctions logiques précisément en code à deux fils & \cite{Sparso2001} & théorie \\
7 & La plus grande différence de temps d'arrivée du son aux deux oreilles est chez l'homme de $\approx0.6\ms$, et le cerveau distingue environ dix microsecondes & Auditeurs dans une expérience à choix forcé entre deux réponses: seuil de discrimination de la différence de temps (75\,\% de réponses correctes) de $9$\,\textmu s pour un bruit dans la bande $150$--$1700$\,Hz, $11$\,\textmu s pour un son pur de $1$\,kHz, $28$\,\textmu s pour un clic. La limite de $0.6\ms$ est un calcul du chapitre, $0.22\,\text{m}/343\,\text{m/s}$ & \cite{KlumppEady1956} & mesuré \\
8 & Chez les oiseaux, la différence de temps se change en position grâce à des lignes à retard et des détecteurs de coïncidences~\eqref{eq:itd} (fig.~\ref{fig:itd}) & Effraie des clochers: les axones venus des deux oreilles entrent dans le noyau des détecteurs de coïncidences par des côtés opposés; le temps d'arrivée des impulsions le long de l'axone varie de façon ordonnée avec la profondeur, et l'étendue des retards à travers le noyau ($160$\,\textmu s sur $700$\,\textmu m) correspond à la plage des retards interauraux & \cite{Carr1990} & mesuré \\
9 & Chez les mammifères, on n'a pas trouvé de rangée de retards; la même tâche est résolue par des détecteurs de coïncidences avec une inhibition précisément calée venue de neurones \nt{GLYC} & Enregistrement in vivo dans l'olive supérieure médiane de la gerbille: les réponses ne forment pas de carte des directions; l'application de glycine et de son bloqueur par micropipette montre que l'inhibition \emph{glycinergique} précisément calée dans le temps fixe la plage utile des retards. L'inhibition vient ici de \nt{GLYC}, non de \nt{GABA} & \cite{Brand2002,Grothe2010} & mesuré \\
10 & Un groupe à forte rétroaction est une bascule; l'activité auto-entretenue persiste des secondes après le stimulus (fig.~\ref{fig:bistable}) & Cortex préfrontal du singe dans une tâche de saccade différée vers un point mémorisé (délai de $1$--$6$\,s): $87$ des $170$ neurones liés à la tâche changent de fréquence pendant le délai, et pour $79\,\%$ d'entre eux cela dépend de la direction du point mémorisé. Que cette activité soit entretenue par une rétroaction à deux états stables relève de la théorie & \cite{Funahashi1989,Wang2001} & mesuré \\
11 & Le bit s'écrit si l'apport dure plus de $\approx0.7\tau$ (fig.~\ref{fig:recurrentgroup}b) & Calcul de l'équation $\tau\,\dd r/\dd t=-r+f(wr+I)$ par la méthode d'Euler pour $w=1$, $I=0.6$; pas de script dans \texttt{models/}. Pas d'expérience & calcul du chapitre & modèle \\
12 & La mémoire peut être gardée dans la facilitation des synapses, presque sans impulsions & Théorie: le calcium résiduel dans les terminaisons garde l'information environ une seconde. Appui: dans le cortex préfrontal médian du furet (enregistrement de plusieurs neurones à la fois), il existe un sous-réseau de neurones pyramidaux reliés par des synapses facilitatrices à long effet rémanent. Revue des observations d'intervalles \og silencieux\fg{} pendant le maintien en mémoire. Quelle façon le cortex utilise et quand reste une question ouverte & \cite{Mongillo2008,WangMarkram2006,Stokes2015} & hypothèse \\
13 & La mémoire s'efface par fatigue: la réserve de neurotransmetteur s'épuise & Enregistrements appariés de neurones pyramidaux du cortex: la réponse de la synapse baisse lors d'impulsions fréquentes, la vitesse de la baisse est fixée par la probabilité de libération. Qu'elle éteigne précisément un groupe auto-entretenu n'est pas montré directement & \cite{Tsodyks1997} & mesuré \\
14 & Une chaîne de groupes est une minuterie déclenchée par un événement; chez les oiseaux chanteurs, les neurones se déclenchent strictement à tour de rôle & Enregistrement de neurones identifiés du centre vocal supérieur du diamant mandarin pendant le sommeil et le chant: chaque neurone qui projette vers le noyau moteur émet une seule brève bouffée à un instant précis du chant, et les neurones se déclenchent successivement & \cite{Hahnloser2002} & mesuré \\
\bottomrule
\end{longtable}}

\section{Chapitre~\ref{sec:gaba}. \nt{GLUT} et \nt{GABA} ensemble}

Les énoncés logiques et dynamiques du chapitre y sont déduits de l'analyse linéaire et de la loi d'Ohm; l'expérience montre que les circuits sur lesquels ils reposent (veto retardé, inhibition directe à la suite de l'excitation, stabilisation par l'inhibition, rythme de la boucle \nt{GLUT}--\nt{GABA}) fonctionnent réellement dans le cerveau.

{\footnotesize
\setlength{\tabcolsep}{3pt}\renewcommand{\arraystretch}{1.15}
\begin{longtable}{>{\raggedright}p{0.5cm}>{\raggedright}p{4.0cm}>{\raggedright}p{5.6cm}>{\raggedright}p{2.3cm}>{\raggedright\arraybackslash}p{1.7cm}}
\toprule
N\textsuperscript{o} & Affirmation & Expérience et ce qui est mesuré & Source & Statut \\
\midrule\endfirsthead
\toprule
N\textsuperscript{o} & Affirmation & Expérience et ce qui est mesuré & Source & Statut \\
\midrule\endhead
1 & Les neurones \nt{GABA} représentent d'un cinquième à un quart des neurones du cortex & Immunomarquage du GABA dans $10$ aires du cortex du singe: environ $25\,\%$ des neurones dans $8$ aires, moins de $20\,\%$ dans le cortex visuel primaire. Revue de la diversité des neurones inhibiteurs du cortex & \cite{Hendry1987,Markram2004} & mesuré \\
2 & Ce que fait l'inhibition dépend en grande partie du site d'arrivée: dendrite, corps, lieu de naissance de l'impulsion, terminaison du fil d'un autre neurone & Revue: les classes de neurones \nt{GABA} du cortex se distinguent par leur cible sur le destinataire. Enregistrement simultané du corps et d'une dendrite d'un neurone pyramidal: l'inhibition directe est bien plus forte sur le corps que sur les dendrites. L'inhibition sur la terminaison du fil d'un autre neurone, dans la moelle épinière, réduit la libération de neurotransmetteur d'une seule ligne d'entrée (revue de mesures) & \cite{Markram2004,Pouille2001,Rudomin1999} & mesuré \\
3 & Le changement de signe par un intermédiaire coûte un retard, environ une milliseconde; l'inhibition qui suit l'excitation rétrécit la fenêtre de coïncidence & Neurones pyramidaux de l'hippocampe du rat en tranches: l'inhibition par un intermédiaire suit l'excitation directe avec un court retard, et la fenêtre dans laquelle deux entrées excitatrices se somment jusqu'au seuil est inférieure à $2\ms$. Sur les dendrites, où cette inhibition est plus faible, la fenêtre est plus large & \cite{Pouille2001} & mesuré \\
4 & Le veto $a\wedge\bar b$~\eqref{eq:veto}, avec le OU et une constante égale à un, est fonctionnellement complet (proposition~\ref{prop:gabacomplete}) & Démonstration dans le chapitre. Résultat classique: un réseau d'éléments à seuil dotés d'entrées inhibitrices réalise n'importe quelle expression logique. Énoncé sur le modèle, pas d'expérience & \cite{McCullochPitts1943} & théorie \\
5 & Le veto retardé donne un détecteur de direction du mouvement de vitesse optimale $v^*=\Delta x/d$~\eqref{eq:vstar} & Rétine du lapin: la sélectivité à la direction est créée par une inhibition qui, pour un mouvement dans la direction \og nulle\fg, éteint l'excitation. L'enregistrement séparé des entrées excitatrice et inhibitrice des cellules sélectives a montré que les cellules amacrines en étoile (\nt{GABA}) les inhibent de façon asymétrique, seulement par les prolongements orientés du côté \og nul\fg. La formule de $v^*$ est une déduction du chapitre & \cite{Barlow1965,Fried2002} & mesuré \\
6 & L'inhibition shuntante divise la tension au lieu de soustraire~\eqref{eq:shunt} (fig.~\ref{fig:shunt}) & Loi d'Ohm pour un diviseur à trois conductances, avec un potentiel d'équilibre du chlorure proche du repos; pour un neurone sans impulsions & déduction du chapitre & théorie \\
7 & Sur la fréquence des impulsions, le shunt agit par soustraction, et la division vient de l'inhibition associée à un bruit de fond équilibré & Modèle: chez un neurone qui se déclenche, le courant à travers le shunt est presque constant, et l'effet sur la fréquence est soustractif. Expérience: dans des neurones pyramidaux du cortex somatosensoriel du rat (tranches), on a introduit par clamp dynamique des conductances de fond excitatrice et inhibitrice; leur augmentation simultanée et équilibrée divise la pente de la réponse sans décalage notable de la fréquence. Revue: la division par l'activité totale se rencontre dans beaucoup de régions du cerveau & \cite{HoltKoch1997,Chance2002,Carandini2012} & mesuré \\
8 & Pour $\beta>1$, l'inhibition mutuelle choisit un seul gagnant (proposition~\ref{prop:wta}, \eqref{eq:wta}) & Valeurs propres $-(1\pm\beta)/\tau$: déduction du chapitre. Les fréquences $0.73$ et $0.53$ pour $\beta=0.5$ sont le point fixe $r_{1,2}=(I_{1,2}-\beta I_{2,1})/(1-\beta^2)$; pas de script dans \texttt{models/}. Le chapitre ne cite pas d'expérience directe & déduction du chapitre & théorie \\
9 & Effacement de la mémoire par un signal via \nt{GABA}; deux groupes à inhibition mutuelle forment un verrou & Découle de la fig.~\ref{fig:bistable} et de la proposition~\ref{prop:wta}. Pas d'expérience & déduction du chapitre & théorie \\
10 & L'équilibre de la boucle \nt{GLUT}--\nt{GABA} est stable si l'inhibition est assez forte et assez rapide (proposition~\ref{prop:ei}) & Modèle à deux populations~\eqref{eq:ei}; les conditions (i) et (ii) portent sur le déterminant et la trace de sa matrice, déduits dans le chapitre & \cite{WilsonCowan1972} & théorie \\
11 & Pour $w_{EE}=1.5$ et $w_{EI}w_{IE}=2$, une inhibition rapide ($\tau_I=2\ms$) maintient $E^*=0.67h$, une inhibition lente ($30\ms$) déclenche des oscillations (fig.~\ref{fig:ei}) & Résolution numérique de~\eqref{eq:ei}, script \texttt{figures/make\_ei.py} & calcul & modèle \\
12 & Le cortex fonctionne en régime stabilisé par l'inhibition; signature: poussez les neurones \nt{GABA}, et leur fréquence baisse~\eqref{eq:isn} & La théorie a prédit la réponse paradoxale. Expérience: enregistrements intracellulaires dans le cortex visuel primaire du chat; quand la réponse est supprimée par un stimulus environnant, la conductance inhibitrice n'augmente pas, mais baisse avec la conductance excitatrice. L'excitation optogénétique des neurones PV de la souris dans les cortex visuel, somatosensoriel et moteur abaisse la fréquence des neurones inhibiteurs, y compris sans stimulus et sous anesthésie légère. Le coefficient $-1/3$ avec les valeurs de la fig.~\ref{fig:ei} est une déduction du chapitre & \cite{Tsodyks1997isn,Ozeki2009,Sanzeni2020} & mesuré \\
13 & La boucle \og \nt{GLUT} excite \nt{GABA}, \nt{GABA} inhibe \nt{GLUT}\fg{} engendre un rythme rapide de période $12$--$50\ms$ & La stimulation optogénétique des neurones \nt{GABA} rapides dans le cortex somatosensoriel de la souris, à des fréquences de $8$--$200$\,Hz, renforce sélectivement le rythme gamma ($20$--$80$\,Hz, période $12$--$50\ms$); la stimulation des neurones \nt{GLUT} ne renforce que des rythmes plus lents & \cite{Cardin2009,Buzsaki2012} & mesuré \\
\bottomrule
\end{longtable}}

\section{Chapitre~\ref{sec:code}. Le code Morse}

Les bornes du chapitre relèvent de la théorie de l'information et du calcul; sont mesurés l'endroit où naît le bruit dans le canal, la vitesse du cerveau et le coût d'une impulsion, tandis que la question du code que le cortex utilise réellement reste ouverte.

{\footnotesize
\setlength{\tabcolsep}{3pt}\renewcommand{\arraystretch}{1.15}
\begin{longtable}{>{\raggedright}p{0.5cm}>{\raggedright}p{4.0cm}>{\raggedright}p{5.6cm}>{\raggedright}p{2.3cm}>{\raggedright\arraybackslash}p{1.7cm}}
\toprule
N\textsuperscript{o} & Affirmation & Expérience et ce qui est mesuré & Source & Statut \\
\midrule\endfirsthead
\toprule
N\textsuperscript{o} & Affirmation & Expérience et ce qui est mesuré & Source & Statut \\
\midrule\endhead
1 & Les fréquences des neurones \nt{GLUT} du cortex vont d'une fraction à quelques dizaines d'impulsions par seconde, et la plupart du temps les neurones travaillent près de la borne inférieure & Enregistrement par une pipette accolée à la cellule (le choix du neurone ne dépend pas de son activité) dans le cortex auditif du rat éveillé: à tout instant, moins de $5\,\%$ des neurones dépassent $20$ impulsions par seconde; chez une partie des neurones, la fréquence de fond est inférieure à $0.01$ par seconde; la distribution des fréquences est log-normale & \cite{Hromadka2008} & mesuré \\
2 & Capacité d'un canal à impulsions $R_{\max}\approx r\log_2\frac{e}{r\Delta t}$~\eqref{eq:spikeentropy}, presque entièrement dans l'instant des impulsions (proposition~\ref{prop:capacity}) & Entropie d'une chaîne binaire de cases de longueur $\Delta t$. Valeurs du chapitre: $8$ bits par impulsion pour $r=10$ par seconde et $\Delta t=1\ms$, environ $5$ bits pour $\Delta t=10\ms$, moins de $2$ bits par seconde pour un compteur d'impulsions & \cite{MacKay1952} & théorie \\
3 & Les neurones sensoriels transmettent d'un à quelques bits par impulsion, dans les meilleurs cas jusqu'à la moitié de la limite & Reconstruction du stimulus à partir des impulsions de neurones sensoriels dont le stimulus est connu; synthèse des mesures dans le livre & \cite{Rieke1997} & mesuré \\
4 & Le générateur d'impulsions est précis; ce sont les variations rapides de l'entrée qui fixent l'instant précis & Neurones du cortex du rat en tranches: pour un courant répété à fluctuations rapides, les impulsions se répètent à mieux que $1\ms$ près; pour un courant constant, les instants se dispersent & \cite{Mainen1995} & mesuré \\
5 & La synapse n'est pas fiable: plus de la moitié des impulsions n'atteignent pas le destinataire, à cause du caractère aléatoire de la libération & Stimulation minimale des entrées de neurones pyramidaux de l'hippocampe (tranches): plus de la moitié des impulsions ne provoquent pas de réponse. Sont exclues les fluctuations du seuil de stimulation, les défauts de conduction aux ramifications de l'axone et la basse température de l'expérience; il reste la libération probabiliste & \cite{AllenStevens1994} & mesuré \\
6 & Le flux d'impulsions du cortex vivant paraît aléatoire: intervalles dispersés comme dans un flux de Poisson, variance du nombre d'impulsions proche de la moyenne; une partie du \og bruit\fg{} est un message sur les mouvements de l'animal & Enregistrements dans les aires visuelles V1 et MT du singe éveillé: coefficient de variation des intervalles proche de $1$, rapport de la variance du nombre d'impulsions à la moyenne comme pour un processus aléatoire; un modèle qui somme de petites entrées indépendantes donne une dispersion bien plus faible. Dans le cortex visuel de la souris, une part importante de la variabilité se prédit d'après l'enregistrement vidéo des mouvements de l'animal lui-même & \cite{Softky1993,Stringer2019} & mesuré \\
7 & Le code de fréquence est lent: erreur $1/\sqrt{rT}$~\eqref{eq:rateerror}, alors que le cerveau reconnaît une image en $\sim150\ms$ & La formule relève de la statistique de Poisson, déduction du chapitre. Expérience: des sujets décidaient si un animal figurait sur une photographie inconnue montrée pendant $20\ms$; les potentiels cérébraux pour \og oui\fg{} et \og non\fg{} divergent au bout d'environ $150\ms$. Le découpage de ce temps en une dizaine d'étapes de $10$--$20\ms$ est une estimation du chapitre, pas une mesure & \cite{Thorpe1996} & mesuré \\
8 & Le moyennage sur un groupe est limité par la corrélation: l'erreur ne baisse pas de plus d'un facteur $1/\sqrt c$~\eqref{eq:correlated} (proposition~\ref{prop:pooling}) & Paires de neurones de l'aire MT du singe enregistrés simultanément: coefficient de corrélation des nombres d'impulsions de $0.12$ en moyenne, et l'analyse théorique montre que même cette corrélation limite le gain du groupe. Théorie: seule la part de la corrélation qui ressemble au signal impose une limite. Modèle de la position opposée: la fréquence d'un groupe de $50$--$100$ neurones se lit en un seul intervalle entre impulsions, $10$--$50\ms$, et l'instant des impulsions isolées n'est pas utilisé dans le cortex & \cite{Zohary1994,MorenoBote2014,ShadlenNewsome1998} & mesuré \\
9 & Un nombre peut s'écrire dans l'instant de la première impulsion~\eqref{eq:latency}, dans l'ordre des impulsions d'un groupe ou dans la phase d'un rythme local & Rétine de la salamandre: la structure spatiale d'une image brièvement présentée est inscrite dans les délais relatifs des premières impulsions des neurones de sortie, presque indépendamment du contraste. Cellules de lieu de l'hippocampe du rat: en traversant \og leur\fg{} lieu, les impulsions se décalent de plus en plus tôt dans le cycle du rythme thêta ($7$--$12$\,Hz), avec un décalage de $100^\circ$ à $355^\circ$. Dans quelle mesure le cortex utilise l'instant des impulsions reste une question ouverte (ligne~8) & \cite{Gollisch2008,OKeefe1993} & mesuré \\
10 & Le code lu dépend de la constante de temps du destinataire~\eqref{eq:reader} (proposition~\ref{prop:reader}); dans le cortex actif, les neurones sont plus proches de détecteurs de coïncidences & La formule~\eqref{eq:reader} est une déduction du chapitre. Revue d'enregistrements in vivo: dans le cortex actif, la conductance de la membrane est plusieurs fois plus élevée qu'au repos, et la constante de temps d'autant plus courte. Que le cortex change de code précisément ainsi n'est pas montré directement & \cite{Destexhe2003} & hypothèse \\
11 & Une synapse à dépression transmet la variation de fréquence, au fond $\Delta r/r$~\eqref{eq:depression} (fig.~\ref{fig:depression}) & Enregistrements appariés de neurones pyramidaux du cortex: la vitesse de dépression est fixée par la probabilité de libération; avec une dépression lente, la réponse reflète la fréquence, avec une dépression rapide, les coïncidences. Modèle fondé sur la dépression mesurée dans le cortex du rat: des variations relatives égales de fréquence sur des entrées rapides et lentes donnent la même réponse, c'est-à-dire que la synapse transmet $\Delta r/r$. Les nombres $1.7\to2.2$, $6.7$ et $90\ms$ sont des calculs du chapitre & \cite{Tsodyks1997,Abbott1997depression} & mesuré \\
12 & La bouffée est un \og trait\fg: sa probabilité de se perdre, $(1-p)^k$~\eqref{eq:burst}, est faible, et la facilitation la réduit encore & Revue: aux synapses peu fiables, les impulsions isolées se perdent souvent, tandis que les bouffées sont transmises de façon fiable grâce à la facilitation; une bouffée provoque des modifications durables des synapses. Que le cerveau lise la bouffée comme un signe à part est une hypothèse & \cite{Lisman1997,AllenStevens1994} & hypothèse \\
13 & Les neurones \nt{GABA} rapides fixent le rythme et la \og ponctuation\fg; une autre classe encore inhibe les inhibiteurs et ouvre le portillon (proposition~\ref{prop:punctuation}) & Revue des propriétés des classes de neurones \nt{GABA} du cortex. Optogénétique: l'excitation rythmique des neurones \nt{GABA} rapides provoque un rythme gamma, et la réponse à un stimulus dépend de la phase du cycle. Dans le cortex visuel de la souris, les neurones PV s'inhibent mutuellement, les neurones SST inhibent toutes les autres classes, les neurones VIP surtout les SST. L'excitation des neurones VIP chez la souris éveillée lève l'inhibition des neurones \nt{GLUT}; ils sont activés par un signal de renforcement. Le partage des rôles comme règle générale est une hypothèse & \cite{Tremblay2016,Cardin2009,Pfeffer2013,Pi2013} & hypothèse \\
14 & L'énergie limite la fréquence moyenne à environ une impulsion par seconde~\eqref{eq:energy} & Budget énergétique de la substance grise du rongeur d'après des données anatomiques et physiologiques: impulsions et courants synaptiques représentent $47\,\%$ et $34\,\%$ de l'énergie de signalisation; au plus $\sim15\,\%$ des neurones peuvent être actifs simultanément. Pour le cortex humain: au plus $\sim1\,\%$ des neurones peuvent être fortement actifs simultanément. Le nombre de $\sim10^9$ ATP par impulsion et la puissance de $\sim10$\,W pour la signalisation sont des estimations du chapitre fondées sur ces calculs & \cite{Attwell2001,Lennie2003} & théorie \\
15 & Le code est parcimonieux et ajusté pour maximiser le nombre de bits par joule; le cortex échange de la précision contre de l'énergie (proposition~\ref{prop:economy}) & Hypothèse du code économe. Appui: chez des souris soumises à une restriction alimentaire, la conductance des récepteurs AMPA et la dépense synaptique d'ATP baissent dans le cortex visuel ($-29\,\%$), la fréquence des impulsions est conservée, et l'accord sur l'orientation s'élargit de $32\,\%$ & \cite{Barlow1961,Padamsey2022} & hypothèse \\
\bottomrule
\end{longtable}}

\section{Chapitre~\ref{sec:serocomputer}. Les neurones \nt{SERO}}

Les propriétés des neurones \nt{SERO} eux-mêmes (rythme, autoinhibition, signe fixé par le récepteur, sous-systèmes) sont mesurées directement; les calculs du chapitre (régulateur, filtre, logique) découlent de ses équations; $\tau_c$, le coefficient de diffusion de la sérotonine et le nombre de sites de libération sont des estimations; le rôle de la boucle comme régulateur de niveau dans le cerveau est une hypothèse (dans le noyau du raphé, l'autoinhibition est ponctuelle et ne produit que de courtes pauses), et le rôle de \nt{SERO} comme horizon est une hypothèse appuyée par des expériences comportementales.

{\footnotesize
\setlength{\tabcolsep}{3pt}\renewcommand{\arraystretch}{1.15}
\begin{longtable}{>{\raggedright}p{0.5cm}>{\raggedright}p{4.0cm}>{\raggedright}p{5.6cm}>{\raggedright}p{2.3cm}>{\raggedright\arraybackslash}p{1.7cm}}
\toprule
N\textsuperscript{o} & Affirmation & Expérience et ce qui est mesuré & Source & Statut \\
\midrule\endfirsthead
\toprule
N\textsuperscript{o} & Affirmation & Expérience et ce qui est mesuré & Source & Statut \\
\midrule\endhead
1 & Le signe de l'action de la sérotonine est choisi par le récepteur du destinataire (proposition~\ref{prop:receptor}) & Les récepteurs de la sérotonine se répartissent en familles selon leur pharmacologie et leur mécanisme: 5-HT$_{1A}$ ouvre des canaux potassiques par une protéine G et inhibe la cellule, 5-HT$_{2A}$ et le récepteur ionotrope 5-HT$_3$ l'excitent. Un même neurotransmetteur produit des réponses opposées dans des cellules différentes. & \cite{BarnesSharp1999} & mesuré \\
2 & À l'éveil, le neurone \nt{SERO} émet régulièrement quelques impulsions par seconde & Enregistrement de neurones du noyau du raphé dorsal chez le chat libre de ses mouvements: décharge rythmique lente de $2.8\pm0.2$ imp./s à l'éveil calme; en sommeil lent, la fréquence baisse de $34$--$68\,\%$, en sommeil paradoxal de $98\,\%$. Chez le rat anesthésié, $1.7\pm0.5$ imp./s, très régulièrement. & \cite{TrulsonJacobs1979, GagnonParent2014, JacobsAzmitia1992} & mesuré \\
3 & Le neurone \nt{SERO} est un pacemaker: pas besoin d'une poussée pour chaque impulsion, mais il lui faut un apport de fond régulier (en tranche, la noradrénaline par les récepteurs $\alpha_1$; dans l'organisme, \nt{NORA}) & En tranche de cerveau, les cellules déchargent lentement et régulièrement, chaque impulsion étant suivie d'une post-hyperpolarisation de $200$--$800\ms$. Les cellules spontanément actives sont moins nombreuses en tranche que dans l'organisme: le rythme régulier exige une entrée tonique de noradrénaline par les récepteurs $\alpha_1$, dont le blocage l'éteint. & \cite{VandermaelenAghajanian1983} & mesuré \\
4 & Presque tous les récepteurs sont métabotropes; la réponse est lente: elle monte et retombe en une seconde environ & Toutes les familles de récepteurs, sauf 5-HT$_3$, sont couplées à des protéines G. En tranche, une libération évoquée de sérotonine produit par 5-HT$_{1A}$ un courant inhibiteur qui monte et retombe en moins d'une seconde. & \cite{BarnesSharp1999, CourtneyFord2016} & mesuré \\
5 & Dans les régions cibles, la sérotonine est libérée dans le volume, et pas seulement dans la fente; dans le noyau du raphé lui-même, l'inhibition des voisins se fait point à point & Voltampérométrie cyclique rapide en tranche: la libération par impulsion est la même dans une région riche en synapses et dans le noyau du raphé, où il n'y a presque pas de synapses; elle ne dépend pas du blocage de la recapture ni des récepteurs; il y a bien moins de molécules par terminaison que de transporteurs et de récepteurs, et la concentration extracellulaire coïncide avec l'affinité du récepteur principal. Dans le noyau du raphé, l'inhibition par 5-HT$_{1A}$ se fait point à point: le transporteur empêche la sérotonine de s'accumuler hors de la fente. & \cite{Bunin1998, CourtneyFord2016} & mesuré \\
6 & La concentration selon~\eqref{eq:seroneuron} décroît par recapture; $\tau_c$ de l'ordre de la seconde est une estimation & Voltampérométrie dans le cerveau vivant: la sérotonine extracellulaire est limitée avant tout par la recapture et la dégradation, non par la synthèse. Les travaux cités ne mesurent pas $\tau_c$ directement; l'ordre de grandeur est une estimation du chapitre. & \cite{Hashemi2012serotonin} & mesuré; $\tau_c$: estimation \\
7 & NON et NON-OU sur un seul pacemaker; complétude de la logique \nt{SERO} (proposition~\ref{prop:serocomplete}, fig.~\ref{fig:serogates}) & Découle de~\eqref{eq:seroneuron}: un pacemaker inhibable est un NON-OU, et le NON-OU est fonctionnellement complet. Aucune expérience n'a construit de logique avec des neurones \nt{SERO}; la condition de volumes séparés n'est pas remplie dans le cerveau. & démonstration dans le chapitre & théorie \\
8 & La sérotonine inhibe le neurone \nt{SERO} lui-même par des récepteurs qu'il porte & Chez le rat anesthésié, des agonistes 5-HT$_{1A}$ administrés par voie intraveineuse et par iontophorèse inhibent la décharge des neurones du raphé dorsal avec la même relation dose-réponse que la sérotonine elle-même; les agonistes 5-HT$_{1B}$ n'agissent presque pas. En tranche, la sérotonine endogène des cellules voisines produit par 5-HT$_{1A}$ un courant et une courte pause dans la décharge, sans changer la fréquence moyenne. & \cite{Sprouse1987, CourtneyFord2016} & mesuré \\
9 & Dans le modèle, la boucle par un volume commun est un régulateur de niveau: la dispersion et la constante de temps sont divisées par $1+G$~\eqref{eq:feedback}; que la boucle fonctionne ainsi dans le cerveau est une hypothèse & Formule classique de l'amplificateur à rétroaction négative, redémontrée dans le chapitre. Le gain de boucle $G$ du système \nt{SERO} n'est pas mesuré. Ce qui est mesuré: en tranche, l'autoinhibition se fait point à point, produit une courte pause et ne change pas la fréquence moyenne. & \cite{Black1934feedback, CourtneyFord2016} & théorie; dans le cerveau: hypothèse \\
10 & La concentration est une moyenne glissante de la fréquence, un filtre passe-bas~\eqref{eq:lowpass}, fig.~\ref{fig:lowpass} & Solution de l'équation linéaire~\eqref{eq:seroneuron}; la figure est calculée par cette formule pour $\tau_c=1\,\text{s}$. Pas de modèle séparé dans \texttt{models/}. & démonstration dans le chapitre & théorie \\
11 & La tortuosité des espaces extracellulaires ralentit la diffusion d'une petite molécule d'un facteur $2.6$ environ & Méthode de la source ponctuelle: iontophorèse de tétraméthylammonium et enregistrement de sa concentration par une électrode sélective en tranche et dans le cerveau vivant. Tortuosité $\lambda\approx1.6$, c'est-à-dire un $D$ effectif $\lambda^2\approx2.6$ fois plus petit que le $D$ libre; fraction de volume extracellulaire voisine de $0.2$. Le coefficient de diffusion de la sérotonine elle-même dans le tissu n'est pas mesuré dans ces travaux; dans le chapitre, c'est une estimation. & \cite{Sykova2008} & mesuré \\
12 & La longueur d'un \og fil\fg{} indépendant est $\ell\sim\sqrt{D\tau}\approx20$~\textmu m~\eqref{eq:diffusionlength}; le nombre de fils est limité par le nombre de telles régions (proposition~\ref{prop:volumewires}) & Loi de la diffusion et substitution des estimations de $D$ et $\tau$ (lignes~6 et~11); la proposition~\ref{prop:volumewires} en est une conséquence. Aucune expérience directe ne compte les canaux indépendants de la transmission volumique. & démonstration dans le chapitre & théorie \\
13 & La sortie d'un seul neurone \nt{SERO} couvre d'immenses portions du cerveau; $\sim10^5$ sites de libération selon l'estimation & Reconstruction complète de neurones isolés du raphé dorsal chez le rat: tous les axones ascendants se ramifient abondamment, la longueur totale de l'axone d'une cellule atteint $18.7$~cm, les branches d'une même cellule atteignent le striatum, le cortex préfrontal et l'amygdale. Le nombre de sites de libération par cellule n'a pas été compté directement. & \cite{GagnonParent2014} & mesuré; $10^5$: estimation \\
14 & Les neurones \nt{SERO} se répartissent en quelques sous-systèmes aux sorties et aux réponses différentes & Marquage viro-génétique chez la souris: les neurones qui projettent vers l'amygdale et vers le cortex frontal ne se recouvrent presque pas par leurs ramifications, reçoivent des entrées différentes et répondent de façon opposée à un stimulus désagréable; activer ou inactiver chaque groupe modifie le comportement différemment. & \cite{Ren2018} & mesuré \\
15 & Condition de patience $D<\tau_\gamma\ln(R/r_0)$~\eqref{eq:patience} pour une dépréciation exponentielle; pour la dépréciation mesurée, hyperbolique, $D<\tau_\gamma(R/r_0-1)$ & Découle de la dépréciation $e^{-D/\tau_\gamma}$ ou $1/(1+D/\tau_\gamma)$; démonstration dans le chapitre. Singes, choix avec des délais de quelques secondes: la dépréciation est plus proche de l'hyperbole que de l'exponentielle; la réponse des neurones \nt{DOPA} au signal annonçant une récompense différée décroît elle aussi avec le délai selon une hyperbole. & démonstration dans le chapitre; \cite{KobayashiSchultz2008} & théorie \\
16 & \nt{SERO} fixe l'horizon $\tau_\gamma$ (section~\ref{sec:horizon}) & L'activation optogénétique des neurones \nt{SERO} du raphé dorsal chez la souris allonge l'attente d'une récompense différée et accroît la persévérance dans la recherche d'une récompense devenue rare. Chez l'homme, l'abaissement de la sérotonine (régime sans tryptophane) rend la dépréciation d'une récompense différée plus forte. Comment la concentration se transforme en $\tau_\gamma$, on ne le sait pas. & \cite{Doya2002, Miyazaki2014, Lottem2018, Schweighofer2008} & hypothèse \\
\bottomrule
\end{longtable}}

\section{Chapitre~\ref{sec:dofa}. Apprendre avec \nt{DOPA}}

Les mécanismes de la synapse (quanta, détecteur de coïncidences, calcium, fenêtres de plasticité) et le comportement de la dopamine comme erreur de prédiction sont mesurés directement; que les règles mènent au gradient et ce que coûte la diffusion sont des théorèmes; le résultat de l'apprentissage du choix est un calcul par le modèle du livre; le schéma qui calcule l'erreur est une hypothèse.

{\footnotesize
\setlength{\tabcolsep}{3pt}\renewcommand{\arraystretch}{1.15}
\begin{longtable}{>{\raggedright}p{0.5cm}>{\raggedright}p{4.0cm}>{\raggedright}p{5.6cm}>{\raggedright}p{2.3cm}>{\raggedright\arraybackslash}p{1.7cm}}
\toprule
N\textsuperscript{o} & Affirmation & Expérience et ce qui est mesuré & Source & Statut \\
\midrule\endfirsthead
\toprule
N\textsuperscript{o} & Affirmation & Expérience et ce qui est mesuré & Source & Statut \\
\midrule\endhead
1 & La rétropropagation de l'erreur, sous la forme qu'elle a dans les réseaux artificiels, n'a pas été trouvée dans le cerveau & Pas d'expérience directe: c'est une affirmation d'absence. Des revues analysent ce qu'exige l'algorithme (connexions de retour reproduisant les connexions aller, phase séparée) et les approximations biologiques proposées. & \cite{Lillicrap2020, WhittingtonBogacz2019} & hypothèse \\
2 & Le poids se décompose en nombre de sites de libération, probabilité de libération et réponse à un quantum: $w=N_s\,p\,A_1$~\eqref{eq:quantal} & Jonction neuromusculaire de grenouille à libération réduite: la réponse du destinataire varie par paliers multiples de la réponse miniature spontanée à un quantum; le nombre de quanta par impulsion suit une loi de Poisson. & \cite{delCastillo1954} & mesuré \\
3 & Le récepteur du destinataire est un détecteur de coïncidences; le calcium déclenche la modification du poids & Patch-clamp de neurones de souris: les ions Mg$^{2+}$ bouchent au potentiel de repos le canal ouvert par le glutamate et le libèrent lors de la dépolarisation. Tranches d'hippocampe: un chélateur du calcium dans le destinataire bloque la potentialisation à long terme, et la libération photo-induite de calcium dans la cellule renforce à elle seule la transmission. & \cite{Nowak1984, Malenka1988calcium} & mesuré \\
4 & N'importe quel côté exécute la décision: $A_1$ croît chez le destinataire, $p$ change chez l'émetteur & Le glutamate libéré par la lumière au-dessus d'une seule épine provoque une croissance rapide et durable de l'épine et du courant de ses récepteurs; la potentialisation s'accompagne de l'insertion de récepteurs AMPA. La dépolarisation du destinataire supprime temporairement la libération chez l'émetteur; l'effet est porté par des endocannabinoïdes qui retraversent la fente (montré sur des entrées \nt{GABA}). La dépression à court terme dépend de la probabilité de libération de l'émetteur. & \cite{Matsuzaki2004, Malinow2002, WilsonNicoll2001, Tsodyks1997} & mesuré \\
5 & La synapse se renforce quand l'émetteur et le destinataire sont actifs ensemble~\eqref{eq:hebb} & Lapin anesthésié: après de courtes séries d'impulsions sur la voie d'entrée de l'hippocampe, la réponse est renforcée pendant $30$~min à $10$~h. En tranche d'hippocampe, les impulsions du destinataire ne renforcent que les synapses actives au même moment. & \cite{Bliss1973LTP, Kelso1986hebb} & mesuré \\
6 & La règle~\eqref{eq:oja} trouve le vecteur propre principal des corrélations des entrées (proposition~\ref{prop:oja}) & Théorème; démonstration dans le chapitre. Aucune expérience ne montre un neurone ayant appris la composante principale de ses entrées; le mécanisme qui limite la croissance des poids dans la synapse n'est pas établi. & \cite{Oja1982} & théorie \\
7 & Hebb temporel: fenêtre d'environ $\pm20\ms$ (fig.~\ref{fig:stdp}); des séquences répétées font pousser des chaînes & Paires de neurones d'hippocampe en culture: une impulsion du destinataire dans les $20\ms$ qui suivent celle de l'émetteur donne un renforcement durable, dans les $20\ms$ qui précèdent, un affaiblissement; les deux dépendent des récepteurs NMDA. Le rapport $A_-=0.6A_+$ de la figure est une illustration. Que des chaînes poussent ainsi dans un réseau n'a pas été mesuré directement. & \cite{BiPoo1998} & mesuré \\
8 & Trace~\eqref{eq:threefactor}: la dopamine arrivée $1$--$2\,\text{s}$ après l'activité conjointe renforce la connexion, plus tard non (proposition~\ref{prop:threefactor}) & Tranches de striatum, stimulation optogénétique séparée des entrées glutamatergiques et dopaminergiques: l'épine ne grossit que si la dopamine arrive $0.3$--$2\,\text{s}$ après l'entrée glutamatergique; la fenêtre est fixée par la montée rapide et la dégradation de l'AMPc. Dans le cortex, l'appariement laisse des traces séparées pour le renforcement et l'affaiblissement, que les récepteurs des monoamines transforment en modifications durables dans une fenêtre de l'ordre de la seconde. & \cite{Yagishita2014, He2015eligibility} & mesuré \\
9 & Des fenêtres de plusieurs secondes existent aussi sans dopamine: le troisième signal est une forte excitation du destinataire & Souris vivante, champ CA1: un seul événement de plateau dans le destinataire renforce les entrées arrivées quelques secondes avant et après lui, et crée un champ de lieu en un seul passage. & \cite{Bittner2017} & mesuré \\
10 & Le pas moyen de la règle à trois facteurs suit le gradient de la récompense~\eqref{eq:perturbation} (proposition~\ref{prop:gradient}) & Théorème du gradient pour l'apprentissage par écarts aléatoires; dans le chapitre, établi pour un bruit faible. & \cite{Williams1992} & théorie \\
11 & Le bruit est une recherche: le réseau apprend de ses propres écarts aléatoires & Jeunes diamants mandarins: l'inactivation du noyau de sortie du circuit des ganglions de la base supprime brusquement et réversiblement la variabilité du chant. Diamants adultes: une perturbation est appliquée aux exécutions d'une syllabe dont la hauteur est d'un côté de la moyenne, et les oiseaux déplacent vite la hauteur de l'autre côté: ils apprennent d'après leur propre dispersion. & \cite{Olveczky2005, TumerBrainard2007} & mesuré \\
12 & Le choix entre deux options s'apprend avec $\tau_e=1\,\text{s}$ et $\delta=R-\bar R$; pas avec $\tau_e=50\ms$; sans soustraction, les poids croissent sans limite, et avec une grande vitesse d'apprentissage le réseau peut figer le plus mauvais choix & \texttt{models/dofa\_learning.py}, $\eta=0.05$, $500$ essais, $6$ graines. $\tau_e=1\,\text{s}$: part des choix de $A$ $0.65$--$0.79\to0.96$--$1.00$, poids $0.85$--$1.33$. $\tau_e=50\ms$: $0.38$--$0.55$, poids inchangés. $\delta=R$: poids du vainqueur $860$--$1190$. $\delta=R$, $\eta=0.5$, $3$ graines: l'une s'est figée sur $B$ ($0.04\to0$). Le script imprime les quatre cas; le recalcul concorde avec le chapitre. & \texttt{models/\allowbreak dofa\_\allowbreak learning.py} & modèle \\
13 & Le temps d'apprentissage par un seul signal commun croît proportionnellement au nombre de neurones bruités (proposition~\ref{prop:broadcastcost}) & Courbes d'apprentissage d'un réseau linéaire: la perturbation des neurones est plus lente que le gradient exact d'un facteur égal au nombre de neurones de sortie, la perturbation des poids encore d'un facteur égal au nombre d'entrées. & \cite{Werfel2005} & théorie \\
14 & Les sorties de \nt{DOPA} vont avant tout aux régions qui choisissent les actions & Reconstruction complète des axones de neurones \nt{DOPA} nigrostriés isolés chez le rat: les branches d'une cellule couvrent $0.45$--$5.7\,\%$ du volume du striatum (en moyenne $2.7\,\%$), et il n'y a presque pas de branches hors du striatum. & \cite{Matsuda2009} & mesuré \\
15 & Le cortex apprend à prédire son entrée, l'erreur est calculée sur place & Revue: le cortex sensoriel présente des réponses au désaccord entre l'entrée attendue et l'entrée reçue. Qu'il s'agisse d'une règle générale d'apprentissage du cortex n'est pas démontré. & \cite{Keller2018} & hypothèse \\
16 & La dopamine se comporte comme l'erreur~\eqref{eq:ctd}: pic, transfert sur le signal annonciateur, creux (fig.~\ref{fig:ctdcases}) & Neurones \nt{DOPA} du singe: pic pour un jus inattendu; après apprentissage, pic sur le signal annonciateur et pas de réponse au jus promis; creux quand le jus n'arrive pas. La forme continue~\eqref{eq:ctd} relève de la théorie. & \cite{Schultz1997, Doya2000} & mesuré \\
17 & Schéma qui calcule $\dd V/\dd t$ et soustrait l'attente & Souris, enregistrement et optogénétique dans l'aire tegmentale ventrale: les neurones \nt{DOPA} soustraient l'attente de la récompense; les neurones \nt{GABA} voisins les inhibent quand la récompense est attendue, et les exciter ou les inhiber modifie l'erreur. Le calcul de la dérivée n'a pas été montré. & \cite{Eshel2015arithmetic} & hypothèse \\
18 & Le neurone \nt{DOPA} est un pacemaker; les creux sont moins profonds que les pics & En tranche, les neurones \nt{DOPA} déchargent spontanément et très régulièrement, sur un potentiel pacemaker lent. Chez le singe, la fréquence suit quantitativement l'erreur positive, mais ne transmet que faiblement l'erreur négative. & \cite{GraceOnn1989, Bayer2005} & mesuré \\
19 & L'acteur et le critique (fig.~\ref{fig:actorcritic}) & L'algorithme vient de la théorie de l'apprentissage par renforcement. Chez l'homme (IRMf), le striatum ventral se comporte comme un critique avec ou sans choix, le striatum dorsal seulement quand il faut choisir des actions. & \cite{SuttonBarto2018, ODoherty2004actor} & hypothèse \\
20 & Le signe de $\delta$ dans la règle est inversé chez les neurones de la seconde famille de récepteurs & Tranches de striatum: chez les neurones à récepteurs D1, l'appariement produit un renforcement qui exige D1; chez les neurones à D2, un affaiblissement qui exige D2. & \cite{Shen2008} & mesuré \\
\bottomrule
\end{longtable}}

\section{Chapitre~\ref{sec:dofaalt}. Autres théories de \nt{DOPA}}

Chacune des images élargies s'appuie sur ses propres mesures directes de la dopamine (dispersion des points d'inversion, réponses au désagréable et au nouveau, montée lente, réponses au changement de goût), mais les formules qui les expliquent sont des théories, et pour deux d'entre elles les expériences se contredisent encore.

{\footnotesize
\setlength{\tabcolsep}{3pt}\renewcommand{\arraystretch}{1.15}
\begin{longtable}{>{\raggedright}p{0.5cm}>{\raggedright}p{4.0cm}>{\raggedright}p{5.6cm}>{\raggedright}p{2.3cm}>{\raggedright\arraybackslash}p{1.7cm}}
\toprule
N\textsuperscript{o} & Affirmation & Expérience et ce qui est mesuré & Source & Statut \\
\midrule\endfirsthead
\toprule
N\textsuperscript{o} & Affirmation & Expérience et ce qui est mesuré & Source & Statut \\
\midrule\endhead
1 & La règle asymétrique~\eqref{eq:asymmetric} converge vers le point~\eqref{eq:expectile}; pour $\kappa=1/2$ c'est la moyenne, pour une goutte de probabilité $1/2$, $V=\kappa$ (fig.~\ref{fig:expectiles}) & Le point~\eqref{eq:expectile} est un expectile, solution d'un problème de moindres carrés avec des poids différents pour les écarts positifs et négatifs; pour l'exemple du chapitre, démonstration dans le chapitre même. & \cite{NeweyPowell1987}; démonstration dans le chapitre & théorie \\
2 & Les différents neurones \nt{DOPA} ont des points d'inversion différents, ordonnés selon l'asymétrie (proposition~\ref{prop:expectile}) & Souris, neurones \nt{DOPA} de l'aire tegmentale ventrale marqués par optogénétique, récompenses de tailles différentes: le point où le pic fait place au creux diffère selon les neurones; il est plus haut chez les neurones aux réponses plus asymétriques; les réponses du groupe permettent de reconstituer la forme de la distribution des récompenses. Que ce soit la fonction principale de la dispersion est une hypothèse. & \cite{Dabney2020} & mesuré \\
3 & Une partie des neurones \nt{DOPA} répond au désagréable par un pic & Singes: certains neurones \nt{DOPA} sont excités par l'indice d'une récompense et inhibés par l'indice d'un jet d'air dans l'œil, d'autres sont excités par les deux; les groupes se trouvent dans des parties différentes du mésencéphale. & \cite{Matsumoto2009, BrombergMartin2010} & mesuré \\
4 & Le module de l'erreur ou la nouveauté fixe la vitesse d'apprentissage & Les travaux cités ne contiennent pas d'expérience directe où le signal d'importance des neurones \nt{DOPA} change la vitesse d'apprentissage; la revue propose un rôle de signal \og attention, c'est important\fg{}. & \cite{BrombergMartin2010} & hypothèse \\
5 & Un fil séparé pour l'axe du danger & Souris: les neurones \nt{DOPA} qui projettent vers l'extrémité postérieure du striatum répondent aux stimuli nouveaux et intenses, et non à la récompense; leur destruction affaiblit l'évitement d'un objet nouveau. & \cite{Menegas2018} & mesuré \\
6 & Temps d'action optimal $\tau_a^*=\sqrt{C/\bar R}$~\eqref{eq:vigor} & Minimum de la somme $C/\tau_a+\bar R\tau_a$; démonstration dans le chapitre. Dans le modèle d'origine, la vitesse des actions est fixée par le prix du temps, égal à la récompense moyenne. & \cite{Niv2007}; démonstration dans le chapitre & théorie \\
7 & Le niveau lent de dopamine porte $\bar R$ et fixe la vitesse des actions (proposition~\ref{prop:multiplex}) & Rats, microdialyse et voltampérométrie dans le noyau accumbens: le niveau de dopamine suit, de minute en minute, la fréquence des récompenses et la vitesse des actions. Chez l'homme, la L-DOPA renforce la dépendance du temps de réaction à la fréquence moyenne des récompenses. Que le niveau lent soit précisément $\bar R$ n'est pas démontré. & \cite{Hamid2016, Beierholm2013vigor} & hypothèse \\
8 & Le niveau lent a deux sources: la libération varie aux sorties sans que change la fréquence des impulsions, mais une montée lente existe aussi dans la fréquence elle-même & Rats, même tâche: la libération dans le noyau accumbens suit l'attente changeante de la récompense, mais pas la fréquence des impulsions des neurones \nt{DOPA} identifiés de l'aire tegmentale ventrale. Souris dans un couloir virtuel (ligne~10): la montée progressive à l'approche de la récompense existe aussi dans les impulsions de neurones \nt{DOPA} identifiés. & \cite{Mohebi2019, Kim2020} & mesuré \\
9 & La dopamine monte progressivement pendant que l'animal s'approche d'une récompense lointaine & Rats dans un labyrinthe, voltampérométrie dans le striatum: la concentration de dopamine augmente en quelques secondes à mesure que le but approche; la pente dépend de la distance et de la taille de la récompense. & \cite{Howe2013ramp, Hamid2016} & mesuré \\
10 & La montée est l'erreur $\delta$ quand l'estimation se précise~\eqref{eq:ramp}; un saut dans le couloir donne un pic (fig.~\ref{fig:ramp}) & La formule et la valeur $\delta=0.75\,V$ par seconde sont démontrées dans le chapitre. Expérience: souris dans un couloir virtuel, téléportation instantanée plus près du but; les impulsions des corps cellulaires, le calcium des corps et des terminaisons et la concentration dans le striatum répondent comme une erreur, et non comme l'attente $V$. Laquelle des deux explications vaut pour toutes les sorties est en discussion. & \cite{Kim2020} & mesuré \\
11 & Regard en arrière: la dopamine signale la mesure du lien causal~\eqref{eq:retro} & Souris, libération de dopamine dans le noyau accumbens dans des expériences où cette théorie et l'erreur~\eqref{eq:ctd} prédisent des choses différentes: dans plusieurs expériences, la libération a suivi la première. & \cite{Jeong2022} & hypothèse \\
12 & Au cours de l'apprentissage, la réponse dopaminergique se déplace progressivement de la récompense vers l'indice & Souris, signal des terminaisons des neurones \nt{DOPA} dans le striatum ventral au fil de l'apprentissage: le pic passe progressivement du moment de la récompense au moment de l'indice, comme l'exige l'apprentissage selon~\eqref{eq:ctd}. & \cite{Amo2022} & mesuré \\
13 & Les neurones \nt{DOPA} répondent au changement de goût de la récompense à valeur égale & Rats, enregistrement de neurones de l'aire tegmentale ventrale: le remplacement d'un jus par un autre de même valeur provoque une réponse, bien que l'erreur~\eqref{eq:ctd} soit nulle. & \cite{Takahashi2017} & mesuré \\
14 & Des pics artificiels enseignent le lien entre deux stimuli neutres & Rats, expérience de préconditionnement sensoriel sans récompense: l'activation optogénétique des neurones \nt{DOPA} lors du passage du premier stimulus au second crée le lien, leur inactivation empêche de l'apprendre. & \cite{Sharpe2017} & mesuré \\
15 & Vecteur d'erreurs par caractéristique~\eqref{eq:vectortd} & Théorie de l'erreur de prédiction des caractéristiques; la forme vectorielle n'a pas été testée directement. & \cite{Gardner2018} & hypothèse \\
16 & Dans une même tâche, différents neurones \nt{DOPA} portent différentes grandeurs & Souris dans un couloir virtuel, imagerie à deux photons des neurones \nt{DOPA}: certains sont liés à la récompense, d'autres à la vitesse et aux virages, d'autres encore aux indices et à la justesse du choix. & \cite{Engelhard2019} & mesuré \\
17 & Un faisceau au lieu d'un fil (proposition~\ref{prop:bundle}) & Synthèse des lignes~2--16: chaque décomposition est mesurée séparément; aucune expérience ne donne une image unifiée où elles seraient toutes les parties d'un même signal. & --- & hypothèse \\
\bottomrule
\end{longtable}}

\section{Chapitre~\ref{sec:nora}. \nt{NORA}}

L'organisation du système \nt{NORA}, ses deux modes, le lien avec la pupille (qui ne passe pas seulement par \nt{NORA}), l'augmentation du rapport du signal au fond et le U inversé du comportement sont mesurés directement; le gain multiplicatif $g$ est un modèle; que le gain change le régime du choix et agisse comme l'inverse d'une température sont des conclusions du chapitre; le bénéfice de l'ajustement du gain est un calcul par le modèle du livre; \og \nt{NORA}, bouton de l'exploration\fg{} et \og \nt{NORA}, interruption\fg{} sont des hypothèses.

{\footnotesize
\setlength{\tabcolsep}{3pt}\renewcommand{\arraystretch}{1.15}
\begin{longtable}{>{\raggedright}p{0.5cm}>{\raggedright}p{4.0cm}>{\raggedright}p{5.6cm}>{\raggedright}p{2.3cm}>{\raggedright\arraybackslash}p{1.7cm}}
\toprule
N\textsuperscript{o} & Affirmation & Expérience et ce qui est mesuré & Source & Statut \\
\midrule\endfirsthead
\toprule
N\textsuperscript{o} & Affirmation & Expérience et ce qui est mesuré & Source & Statut \\
\midrule\endhead
1 & L'homme possède quelques dizaines de milliers de neurones \nt{NORA}, presque tous en un seul groupe & Coupes sériées du tronc cérébral humain, marquage de la tyrosine hydroxylase: $53\,900$ cellules dans le locus coeruleus et encore $6\,260$ dans un sous-noyau voisin. & \cite{Baker1989LC} & mesuré \\
2 & Les sorties se répartissent dans presque tout le cerveau, mais les parties du groupe desservent en partie des régions différentes & Marquage viro-génétique chez la souris: les neurones \nt{NORA} du locus coeruleus reçoivent des entrées de vastes régions du cerveau et envoient des sorties vers de vastes régions du cerveau et de la moelle épinière. Chez le rat: le groupe qui projette vers l'amygdale renforce l'apprentissage de la peur, et un groupe distinct projetant vers le cortex préfrontal renforce son extinction. & \cite{Schwarz2015LC, Uematsu2017LC, Poe2020} & mesuré \\
3 & Mode tonique: impulsions régulières, d'une fraction d'impulsion à quelques impulsions par seconde, fréquence qui varie avec l'éveil & Rats, enregistrement de neurones du locus coeruleus pendant le cycle veille-sommeil: la fréquence est maximale à l'éveil, plus basse en sommeil lent, presque nulle en sommeil paradoxal; les variations de fréquence précèdent le changement d'état. & \cite{AstonJonesBloom1981} & mesuré \\
4 & Mode phasique: courte bouffée sur un événement important pour la tâche, à peu près au moment de la décision & Singes, tâche de détection d'une cible rare: tous les neurones du locus coeruleus répondent par une bouffée à la cible seulement, $\sim90\ms$ après elle et $\sim200\ms$ avant la réponse manuelle; quand la performance est mauvaise, les bouffées sont plus faibles. Dans une tâche de choix entre deux options, la bouffée est plus étroitement liée à la réponse qu'au stimulus, et elle est présente pour les réponses justes comme pour les erreurs. & \cite{AstonJones1994vigilance, Clayton2004LC} & mesuré \\
5 & Le diamètre de la pupille est un point de mesure imprécis de \nt{NORA}: il suit \nt{NORA}, mais aussi \nt{ACET} et d'autres régions & Singes: les impulsions du locus coeruleus, jusqu'aux impulsions isolées d'une seule cellule, prédisent la dilatation de la pupille; un lien semblable, mais plus tardif, existe dans les colliculus et le cortex cingulaire. Souris: les fluctuations de la pupille suivent l'activité des sorties noradrénergiques comme cholinergiques dans le cortex. & \cite{Joshi2016, Reimer2016} & mesuré \\
6 & La noradrénaline supprime l'activité de fond plus que l'activité évoquée; le gain multiplicatif~\eqref{eq:gain} est un modèle qui rend compte de cet effet & Singes éveillés, cortex auditif, iontophorèse de noradrénaline: l'activité de fond des neurones est plus supprimée que la réponse aux vocalisations des congénères; le rapport signal sur bruit augmente. La sigmoïde multiplicative~\eqref{eq:gain} est reprise d'un modèle de réseau; la multiplication littérale de la pente n'est pas mesurée. & \cite{Foote1975NE, ServanSchreiber1990} & mesuré; $g$: modèle \\
7 & Le gain n'augmente $d'$ que contre le bruit postérieur à l'amplificateur~\eqref{eq:gainsnr} (proposition~\ref{prop:gainnoise}) & Démonstration dans le chapitre; dans un modèle de réseau, le gain améliore la détection du signal. Aucune expérience ne sépare le bruit avant et après l'amplificateur pour tester cette distinction. & démonstration dans le chapitre; \cite{ServanSchreiber1990} & théorie \\
8 & La frontière $g\beta=1$ fait passer le réseau de choix de \og réfléchit\fg{} à \og a décidé\fg{}~\eqref{eq:wtagain} (proposition~\ref{prop:gainwta}, fig.~\ref{fig:pitchfork}) & Analyse linéaire de deux groupes qui s'inhibent mutuellement; démonstration dans le chapitre. Ce seuil n'a pas été mesuré directement dans le cerveau. & démonstration dans le chapitre & théorie \\
9 & Le gain est l'inverse de la température du choix~\eqref{eq:softmax} & Avec un bruit logistique après l'amplificateur, la probabilité de choix est un softmax avec $T=\sigma/g$; démonstration dans le chapitre. & démonstration dans le chapitre & théorie \\
10 & \nt{NORA} fixe combien chercher: forte activité tonique, petit $g$ effectif (exploration); bouffée phasique au moment de la décision, grand $g$ (décision) & Humains: avant un choix au hasard, la pupille de base est plus large qu'avant le choix de la meilleure option connue; les personnes à pupille plus large explorent davantage. Rats: le passage au mode de choix aléatoire est commandé par l'entrée \nt{NORA} dans le cortex cingulaire antérieur. Que la bouffée augmente le $g$ effectif n'est pas mesuré directement. & \cite{JepmaNieuwenhuis2011, Tervo2014stochastic, AstonJones2005} & hypothèse \\
11 & La récompense dépend de $g$ comme un U inversé; le meilleur $g$ est plus petit dans un monde qui change vite (proposition~\ref{prop:explore}, fig.~\ref{fig:invertedU}) & \texttt{models/nora\_explore.py}, $60$ simulations de $4000$ essais. Changement tous les $1000$ essais: meilleur $g=16$, part $0.976$; tous les $20$: $g=8$, part $0.886$; pour $g=0.5$, $0.77$. Le recalcul concorde avec le chapitre. & \texttt{models/\allowbreak nora\_\allowbreak explore.py} & modèle \\
12 & U inversé dans le comportement: meilleur résultat avec une activité tonique modérée et des bouffées nettes & Singes, même tâche qu'à la ligne~4: pendant les périodes de bonne performance, la fréquence tonique est modérée et les bouffées sur la cible sont nettes; avec une fréquence tonique élevée, il n'y a pas de bouffées et les fausses réponses sont plus nombreuses. Les variations de l'activité tonique et évoquée suivent les fluctuations de la qualité de la performance. & \cite{Usher1999LC, AstonJones2005} & mesuré \\
13 & \nt{NORA} signale l'incertitude inattendue, \nt{ACET} l'incertitude attendue & Théorie de l'inférence bayésienne à deux sortes d'incertitude; les travaux cités ne contiennent pas d'expérience directe séparant ces signaux selon le neurotransmetteur. & \cite{YuDayan2005} & hypothèse \\
14 & Après un changement brusque, les humains apprennent plus vite, et la pupille se dilate & Humains, tâche de prédiction avec changements soudains: les variations brèves de la pupille suivent l'estimation de la probabilité d'un changement et, avec la pupille de base, prédisent à quel point les nouvelles données modifient l'estimation (vitesse d'apprentissage); une variation de la pupille sans lien avec la lumière ni avec la tâche modifie cette vitesse. Que la pupille indique ici précisément \nt{NORA} est une conclusion tirée des données indirectes de la ligne~5. & \cite{Nassar2012} & mesuré \\
15 & La bouffée phasique fonctionne comme une interruption: le réseau remet son état à zéro & Aucune expérience directe ne montre une bouffée \nt{NORA} remettant à zéro l'état du réseau; seules les bouffées sur les événements importants sont mesurées (ligne~4). & \cite{BouretSara2005} & hypothèse \\
16 & L'ajustement de $g$ ou de la vitesse d'apprentissage par la surprise ne fait guère mieux que les meilleurs réglages constants & \texttt{models/nora\_explore.py}, monde à époques de $1000$ essais (changement tous les $1000$ et tous les $20$ essais). Meilleur $g$ constant, $g=8$: $0.925$; ajustement de $g$: jusqu'à $0.927$; ajustement de la vitesse d'apprentissage: jusqu'à $0.926$; vitesse constante $0.4$: $0.928$. Le recalcul concorde avec le chapitre. & \texttt{models/\allowbreak nora\_\allowbreak explore.py} & modèle \\
\bottomrule
\end{longtable}}

\section{Chapitre~\ref{sec:acet}. Ajoutons \nt{ACET}}

Les trois actions de l'acétylcholine sont mesurées sur tranches et dans le cerveau vivant, le conflit entre écriture et lecture est montré sur le modèle du livre, et que \nt{ACET} serve de ligne de validation d'écriture et signale l'incertitude attendue est une hypothèse en partie appuyée par l'expérience.

{\footnotesize
\setlength{\tabcolsep}{3pt}\renewcommand{\arraystretch}{1.15}
\begin{longtable}{>{\raggedright}p{0.5cm}>{\raggedright}p{4.0cm}>{\raggedright}p{5.6cm}>{\raggedright}p{2.3cm}>{\raggedright\arraybackslash}p{1.7cm}}
\toprule
N\textsuperscript{o} & Affirmation & Expérience et ce qui est mesuré & Source & Statut \\
\midrule\endfirsthead
\toprule
N\textsuperscript{o} & Affirmation & Expérience et ce qui est mesuré & Source & Statut \\
\midrule\endhead
1 & Les neurones \nt{ACET} du cerveau sont peu nombreux, réunis en quelques groupes, et leurs sorties se répartissent dans tout le cortex: c'est un canal de diffusion & Marquage par anticorps de l'enzyme de synthèse de l'acétylcholine et traçage rétrograde chez le rat: le cortex, l'hippocampe, l'amygdale, le bulbe olfactif et le thalamus reçoivent l'acétylcholine surtout de six groupes compacts de cellules du télencéphale basal et du tronc cérébral & \cite{Mesulam1983} & mesuré \\
2 & Le niveau d'acétylcholine dans le cortex est élevé à l'éveil et bas en sommeil lent & Microdialyse dans le cortex et l'hippocampe de chats libres de leurs mouvements, avec enregistrement EEG: la libération d'acétylcholine est supérieure de $\sim100\%$ à l'éveil calme et de $\sim175\%$ à l'éveil actif par rapport au sommeil lent; en sommeil paradoxal, dans le cortex, elle est la même qu'à l'éveil & \cite{Marrosu1995,Hasselmo1999} & mesuré \\
3 & Le sens du signal est fixé par le récepteur: l'acétylcholine a des récepteurs-canaux rapides et des récepteurs lents qui déclenchent des réactions dans la cellule (proposition~\ref{prop:receptor}) & Revue de mesures: l'action de l'acétylcholine sur l'excitabilité, la transmission et la plasticité dépend du lieu de libération, du sous-type de récepteur (nicotiniques: canaux ioniques; muscariniques: par des protéines G) et de la cellule destinataire & \cite{Picciotto2012} & mesuré \\
4 & Première action: affaiblir les connexions internes sans presque toucher aux entrées extérieures (facteur $1-a$ dans~\eqref{eq:acetnet}) & Tranches de cortex olfactif de rat: à $100$\,\textmu M d'agonistes de l'acétylcholine, les potentiels des connexions internes (couche Ib) baissent de plus de $60\%$, ceux de l'entrée externe (couche Ia) de moins de $15\%$, dans la même cellule; la suppression est présynaptique. Tranches d'hippocampe (CA1): $89\%$ contre $40\%$ pour la voie interne et la voie d'entrée & \cite{HasselmoBower1992,HasselmoSchnell1994} & mesuré \\
5 & Deuxième action: renforcer l'entrée (facteur $\frac{1+a}{2}$) & Tranches de néocortex: les récepteurs nicotiniques renforcent les synapses de l'entrée thalamique et n'agissent pas sur les synapses intracorticales; les muscariniques affaiblissent les deux types. L'acétylcholine augmente l'excitabilité des neurones (revue) & \cite{Gil1997,Hasselmo2006} & mesuré \\
6 & Troisième action: faciliter la modification des poids (vitesse $\eta a$) & Rat: la stimulation électrique du noyau basal (source de l'acétylcholine du cortex), appariée à un son, provoque chez l'animal adulte une forte réorganisation de la carte du cortex auditif en faveur du son présenté & \cite{KilgardMerzenich1998,Hasselmo2006} & mesuré \\
7 & La règle de Hebb pour motifs clairsemés, dans la seconde équation~\eqref{eq:acetnet}, donne une mémoire associative qui complète un motif à partir d'une partie & Calcul de la capacité d'un réseau à motifs clairsemés et règle de covariance: pour une faible fraction de neurones actifs $p$, le réseau stocke nettement plus de motifs que pour $p=1/2$ & \cite{Tsodyks1988} & théorie \\
8 & Écrire avec $a$ bas abîme la mémoire: le réseau écrit ce dont il se souvient, et non ce qu'il voit; pour $a=0.1$, les dégâts ne sont pas moindres que pour $a=0$ & \texttt{models/\allowbreak acet\_memory.py}: $200$ neurones, cinq motifs de $20$, un nouveau motif partage $10$ neurones avec l'un d'eux, $10$ simulations. Pour $a=0$, le nouveau motif n'est pas mémorisé, l'ancien entraîne $5.9$ neurones étrangers sur $10$; pour $a=0.1$, le nouveau motif est rappelé ($0.97$), mais les deux fusionnent: le nouveau entraîne $7.3$ neurones étrangers, l'ancien $7.9$; à partir de $a=0.2$, moins d'un & démonstration dans le chapitre & modèle \\
9 & Proposition~\ref{prop:writeread}: aucun niveau $a$ ne convient aussi bien à l'écriture qu'à la lecture (fig.~\ref{fig:acetsim}) & Même script, vitesse d'apprentissage $\eta a$: le nouveau motif est entièrement mémorisé en une présentation pour $a\ge0.9$, à $0.8$ pour $a=0.75$, pas du tout pour $a\le0.5$. Lecture d'après la moitié du motif: $1.0$ pour $a\le0.2$, $0.96$ pour $a=0.3$, $0.04$ pour $a=0.4$ & démonstration dans le chapitre & modèle \\
10 & Dans le cerveau, un seul fil de diffusion, \nt{ACET}, commute écriture et lecture & L'idée vient de modèles construits à partir de mesures sur tranches. L'expérience la plus directe est humaine: la scopolamine, qui bloque les récepteurs muscariniques, gêne l'apprentissage de nouvelles paires de mots, surtout celles qui recoupent d'anciennes, mais pas le rappel des paires déjà apprises. Les deux modes du réseau n'ont pas été mesurés directement & \cite{HasselmoAndersonBower1992,Atri2004} & hypothèse \\
11 & Le filtre~\eqref{eq:kalman} est optimal; le poids de l'entrée et la vitesse d'apprentissage sont une même grandeur $P/R$; en régime établi, $P=\sqrt{qR}$ et la mémoire vaut $\sqrt{R/q}$ (proposition~\ref{prop:kalman}) & Aucune expérience n'est requise: l'optimalité du filtre de Kalman--Bucy est un résultat classique de la théorie de l'estimation; le régime établi est démontré dans le chapitre & \cite{KalmanBucy1961} & théorie \\
12 & \nt{ACET} communique au réseau une grandeur semblable à $P$, l'incertitude attendue & Proposé dans un modèle probabiliste de l'attention. Il n'existe pas de mesure directe montrant que le niveau d'acétylcholine suit l'incertitude de l'estimation & \cite{YuDayan2005} & hypothèse \\
13 & Une partie des neurones \nt{ACET} répond à la récompense et à la punition en une vingtaine de millisecondes, d'autant plus fort que le renforcement est inattendu & Souris, neurones cholinergiques du télencéphale basal identifiés par optogénétique: réponse à la récompense et à la punition après $18\pm3$\,ms, amplitude croissant avec le caractère inattendu du renforcement, réponses semblables chez les neurones de deux noyaux & \cite{Hangya2015} & mesuré \\
14 & \nt{NORA} remarque un changement inattendu du monde, \nt{ACET} maintient le réseau en mode écriture & Division du travail proposée dans un modèle. Indirectement: chez l'homme, le diamètre de la pupille, lié à \nt{NORA}, suit les changements et la vitesse d'apprentissage. Le couple \nt{NORA}--\nt{ACET} n'a pas été testé directement & \cite{YuDayan2005,Nassar2012} & hypothèse \\
15 & En sommeil lent, le réseau en mode lecture rejoue ce qui a été écrit pendant la journée, et ces répétitions le recopient dans la mémoire à long terme & Enregistrements d'ensembles de neurones de l'hippocampe du rat: les cellules qui ont travaillé ensemble pendant la journée se déclenchent plus souvent ensemble pendant le sommeil lent qui suit, et dans le même ordre; c'est mesuré. Pour la recopie: supprimer les ondulations hippocampiques (ripples) après l'apprentissage dégrade la mémoire spatiale, et allonger les ondulations spontanées l'améliore; que la cause soit le bas niveau d'\nt{ACET} n'est pas démontré & \cite{WilsonMcNaughton1994,Skaggs1996,Girardeau2009,FernandezRuiz2019,Hasselmo1999} & hypothèse \\
\bottomrule
\end{longtable}}

\section{Chapitre~\ref{sec:synthesis}. Toute la machine}

Ce chapitre de synthèse n'apporte pas d'expériences nouvelles: chaque ligne s'appuie sur le chapitre où l'affirmation est examinée et sur les mêmes sources; le bus \nt{GLUT} et la commande du flux par \nt{GABA} sont solidement établis, tandis que les rôles des canaux de diffusion et leur fonctionnement conjoint sont pour l'essentiel une hypothèse.

{\footnotesize
\setlength{\tabcolsep}{3pt}\renewcommand{\arraystretch}{1.15}
\begin{longtable}{>{\raggedright}p{0.5cm}>{\raggedright}p{4.0cm}>{\raggedright}p{5.6cm}>{\raggedright}p{2.3cm}>{\raggedright\arraybackslash}p{1.7cm}}
\toprule
N\textsuperscript{o} & Affirmation & Expérience et ce qui est mesuré & Source & Statut \\
\midrule\endfirsthead
\toprule
N\textsuperscript{o} & Affirmation & Expérience et ce qui est mesuré & Source & Statut \\
\midrule\endhead
1 & Pour $8.6\cdot10^{10}$ neurones, il n'existe que quatre sortes de signaux de commande~\eqref{eq:machine} & Comptage des noyaux cellulaires dans un homogénat de cerveau (fractionneur isotrope): $86.1\pm8.1$ milliards de neurones chez l'homme adulte & \cite{Azevedo2009} & mesuré \\
2 & Proposition~\ref{prop:architecture}, deux premières couches: les données circulent sur le bus d'adresses \nt{GLUT}, le signe de la connexion est fixé par l'émetteur, les neurones \nt{GABA} orientent et normalisent le flux (chapitres~\ref{sec:neurontypes}, \ref{sec:gaba}) & Un seul neurotransmetteur principal par neurone (revue des mesures); l'inhibition directe rétrécit la fenêtre dans laquelle le neurone pyramidal additionne ses entrées; la division par l'activité totale est mesurée dans de nombreuses régions & \cite{Strata1999,Pouille2001,Carandini2012} & mesuré \\
3 & Les éléments ne sont pas fiables: une synapse perd la plupart des impulsions (tabl.~\ref{tab:computer}, chapitre~\ref{sec:code}) & Tranches d'hippocampe, stimulation minimale: plus de la moitié des impulsions ne provoquent aucune réponse dans CA1; les échecs du seuil et de la conduction axonale sont exclus, la cause est la libération probabiliste du neurotransmetteur & \cite{AllenStevens1994} & mesuré \\
4 & Le cerveau fonctionne avec $\sim20$ watts, en moyenne environ une impulsion par seconde et par neurone (chapitre~\ref{sec:code}); les ingénieurs n'ont pas encore construit une telle machine & Estimation~\eqref{eq:energy} à partir des coûts mesurés: environ $10^9$ molécules d'ATP par impulsion, travail des synapses compris (cortex de rongeur); une estimation détaillée pour le cortex donne une fréquence encore plus basse. État de la technique d'après une revue & \cite{Attwell2001,Lennie2003,MehonicKenyon2022} & théorie \\
5 & L'apprentissage de la plupart des synapses est le produit d'une trace locale et d'un seul nombre commun $\delta$ (deuxième équation de~\eqref{eq:machine}, chapitre~\ref{sec:dofa}) & Les neurones \nt{DOPA} du singe répondent à l'erreur de prédiction de la récompense; dans des tranches de striatum, la dopamine n'agrandit une épine que si elle arrive $0.3$--$2$\,s après l'entrée glutamatergique: la trace est mesurée & \cite{Schultz1997,Yagishita2014} & mesuré \\
6 & Un fil d'erreur propre à chaque sortie n'existe que dans le circuit d'apprentissage moteur (chapitre~\ref{sec:motorlearning}) & Cervelet: la coïncidence du signal de la fibre grimpante avec une entrée affaiblit cette entrée; chez le singe, les impulsions complexes des cellules de Purkinje portent la direction de l'erreur de mouvement et déclenchent une correction & \cite{Ito2001,Herzfeld2018} & mesuré \\
7 & Chaque canal de diffusion est un faisceau de quelques lignes (proposition~\ref{prop:bundle}) & Pour \nt{DOPA}: dans une même tâche, des neurones différents portent la récompense, le mouvement et les propriétés du stimulus; l'erreur rapide et le niveau lent de dopamine sont distincts. Pour \nt{SERO}, \nt{NORA}, \nt{ACET}, cette décomposition est moins étudiée & \cite{Engelhard2019,Hamid2016,Mohebi2019} & mesuré \\
8 & \nt{SERO}: régulateur de niveau et, par hypothèse, horizon $\tau_\gamma$ (chapitre~\ref{sec:serocomputer}) & Auto-inhibition: les agonistes de leurs propres récepteurs 5-HT$_{1A}$ inhibent les neurones sérotoninergiques du raphé (mesuré). L'horizon est proposé en théorie; l'expérience la plus forte: l'activation optogénétique de ces neurones chez la souris allonge l'attente d'une récompense différée & \cite{Sprouse1987,Doya2002,Miyazaki2014} & hypothèse \\
9 & \nt{NORA}: le gain $g$ choisit entre chercher et décider (chapitre~\ref{sec:nora}) & Hausse du rapport signal/bruit: chez le singe éveillé, la noradrénaline dans le cortex auditif supprime l'activité de fond plus fortement que la réponse aux vocalisations des congénères (mesuré); le gain multiplicatif est un modèle; le rôle dans le choix entre chercher et décider est une théorie & \cite{Foote1975NE,ServanSchreiber1990,AstonJones2005} & hypothèse \\
10 & \nt{ACET}: bascule entre écriture et lecture (chapitre~\ref{sec:acet}) & Les trois actions (affaiblissement des connexions internes, renforcement de l'entrée, facilitation de la plasticité) sont mesurées; la bascule des modes est soutenue indirectement par une expérience à la scopolamine chez l'homme & \cite{HasselmoBower1992,Gil1997,Atri2004} & hypothèse \\
11 & La formule~\eqref{eq:machine} décrit la machine entière & C'est un résumé des modèles des chapitres~\ref{sec:glutcomputer}--\ref{sec:acet}, chaque partie s'appuie sur son chapitre; l'ensemble n'a pas été vérifié par l'expérience & déduction du chapitre & hypothèse \\
12 & Les boutons fonctionnent de concert sans commande centrale: erreur, surprise, écriture, retour (fig.~\ref{fig:timeline}) & Pas d'expérience: le schéma est qualitatif. Maillons isolés: théorie de la division du travail entre \nt{NORA} et \nt{ACET}, et lien entre la pupille et la vitesse d'apprentissage chez l'homme & \cite{YuDayan2005,Nassar2012} & hypothèse \\
13 & L'apprentissage par un seul signal commun est d'autant plus lent que plus de neurones influent sur le résultat (proposition~\ref{prop:broadcastcost}) & Calcul des courbes d'apprentissage d'un réseau linéaire: l'apprentissage par perturbation des sorties est plus lent que la descente de gradient d'un facteur égal au nombre de sorties & \cite{Werfel2005} & théorie \\
14 & On ignore comment le cortex règle ses circuits multicouches; candidats: bouffées d'impulsions, calculs dans les branches du neurone, apprentissage auto-supervisé par prédiction & Bouffées comme porteuses de l'erreur: modèle; seul un réseau de $5$--$8$ couches a su reproduire la réponse d'un modèle détaillé de neurone cortical: modèle; potentiels calciques dans les branches des neurones corticaux humains, réalisant le OU exclusif: mesuré sur tranches; apprentissage auto-supervisé: revue des indices & \cite{Payeur2021,Beniaguev2021,Gidon2020,Keller2018} & hypothèse \\
\bottomrule
\end{longtable}}

\section{Chapitre~\ref{sec:sensors}. Les entrées de la machine}

La structure des capteurs est mesurée directement, par des enregistrements du courant de cellules isolées, des vibrations de la cochlée et par le comptage des cellules de la rétine; la limite de sensibilité et la formule du bruit de grenaille découlent de la physique; que les codes des capteurs soient ajustés pour porter le plus d'information est une hypothèse solidement appuyée.

{\footnotesize
\setlength{\tabcolsep}{3pt}\renewcommand{\arraystretch}{1.15}
\begin{longtable}{>{\raggedright}p{0.5cm}>{\raggedright}p{4.0cm}>{\raggedright}p{5.6cm}>{\raggedright}p{2.3cm}>{\raggedright\arraybackslash}p{1.7cm}}
\toprule
N\textsuperscript{o} & Affirmation & Expérience et ce qui est mesuré & Source & Statut \\
\midrule\endfirsthead
\toprule
N\textsuperscript{o} & Affirmation & Expérience et ce qui est mesuré & Source & Statut \\
\midrule\endhead
1 & Le capteur est un transducteur: le récepteur produit un courant, et la sortie des cellules sensorielles de l'œil et de l'oreille est du glutamate vers le bus \nt{GLUT}; les premières cellules ne produisent pas d'impulsions (fig.~\ref{fig:transducer}) & Les bâtonnets et les cônes de la tortue libèrent un acide aminé excitateur, détecté par des canaux glutamatergiques sur un fragment de membrane excisé. Cellules ciliées internes de la cochlée du rat: les courants dans les terminaisons des fibres nerveuses passent par des récepteurs AMPA du glutamate, et la fréquence des libérations croît avec la dépolarisation progressive de la cellule jusqu'à $150$\,Hz & \cite{CopenhagenJahr1989,GlowatzkiFuchs2002} & mesuré \\
2 & Dans l'obscurité, le capteur de lumière répond à un seul photon par un courant d'environ un picoampère & Électrode à succion sur le segment externe d'un bâtonnet de crapaud: les réponses aux éclairs faibles sont quantifiées, événement de $\approx1$\,pA ($5\%$ du courant d'obscurité) avec une dispersion de $0.2$\,pA, efficacité $0.5$. L'homme signale l'arrivée d'un seul photon plus souvent que le hasard & \cite{Baylor1979,Tinsley2016} & mesuré \\
3 & Le capteur de son détecte un déplacement inférieur au nanomètre & Méthode Mössbauer, membrane basilaire de la cochlée du cobaye au site $16$--$19$\,kHz: au seuil de réponse du nerf auditif, la vitesse de la membrane est d'environ $0.04$\,mm/s, soit un déplacement de $\approx0.35$\,nm (notre conversion). On a mesuré la membrane, et non le faisceau du capteur & \cite{Sellick1982,Hudspeth1989} & mesuré \\
4 & Dans l'obscurité, la précision est limitée par le bruit de grenaille des photons~\eqref{eq:photonnoise}; en lumière faible, l'accumulation est plus longue & La formule découle de la statistique de Poisson. Dans le bâtonnet, le bruit en lumière faible coïncide avec la somme d'événements quantiques indépendants; sur un fond lumineux, la réponse à un éclair est plus petite et plus brève. Le chiffre \og dixièmes de seconde\fg{} n'est pas vérifié séparément ici & \cite{Baylor1979} & théorie \\
5 & La plage d'entrée couvre dix ordres de grandeur pour la lumière et douze pour le son, la fréquence des impulsions moins de trois & Pour le son: la réponse de la membrane basilaire à la fréquence caractéristique croît avec une pente descendant jusqu'à $0.2$\,dB/dB dans la bande $40$--$80$\,dB; la compression reste visible au-delà de $100$\,dB. Les nombres d'ordres de grandeur eux-mêmes sont des estimations classiques, sans source dans le chapitre & \cite{Ruggero1997} & mesuré \\
6 & Réponse du capteur~\eqref{eq:nakarushton}; $\sigma$ suit la luminosité moyenne et fait glisser la courbe le long de l'axe logarithmique (fig.~\ref{fig:adaptation}) & La formule a d'abord été ajustée aux potentiels S de la rétine de poisson. Cônes de la tortue: les réponses obéissent à $V=I V_m/(I+\sigma)$ et couvrent $\sim3.5$ ordres de grandeur de luminosité; après $2$--$3$ minutes de fond, la courbe glisse horizontalement en gardant sa largeur. Lien avec la division par l'activité totale: revue & \cite{NakaRushton1966,NormannPerlman1979,Carandini2012} & mesuré \\
7 & Les capteurs transmettent les changements, et leurs codes sont ajustés pour porter le plus d'information par impulsion (proposition~\ref{prop:sensors}) & Le principe a été proposé en théorie. L'indice le plus fort: chez la mouche, la caractéristique \og contraste $\to$ réponse\fg{} des neurones de la première couche coïncide avec la distribution cumulée des contrastes des scènes naturelles, ce qui maximise l'information & \cite{Barlow1961,Laughlin1981} & hypothèse \\
8 & Les capteurs ON et OFF forment un code à deux fils; la caméra événementielle le reproduit dans le silicium & Revue d'expériences: les voies ON et OFF de la rétine se séparent dès la première synapse et peuvent être supprimées séparément. Caméra: $128\times128$ pixels sans images, plage de $120$\,dB, latence de $15$\,\textmu s & \cite{Schiller1992,Lichtsteiner2008,Gallego2022} & mesuré \\
9 & L'œil compte $\sim10^8$ capteurs de lumière et $\sim10^6$ neurones de sortie: compression d'un facteur $\sim100$ & Comptage sur des rétines humaines entières: en moyenne $92$ millions de bâtonnets et $4.6$ millions de cônes; de $0.7$ à $1.5$ million de cellules de sortie (ganglionnaires) selon les yeux & \cite{Curcio1990,CurcioAllen1990} & mesuré \\
10 & Chez la souris, l'œil a plus de trente canaux de sortie & Souris: imagerie calcique biphotonique de plus de $11\,000$ cellules de sortie et regroupement des réponses, nettement plus de $30$ types fonctionnels. Chez l'homme, le nombre n'est pas mesuré & \cite{Baden2016} & mesuré \\
11 & L'œil du cobaye transmet $\sim875\,000$ bit/s; rapporté à l'œil humain, cela donne $\sim10^7$ bit/s & Cobaye, matrice multiélectrode, mouvement dans des scènes naturelles: $6$--$13$ bit/s par cellule, $\sim875\,000$ bit/s pour $\sim10^5$ cellules de sortie. Pour l'homme ($\sim10^6$ cellules), $\sim10^7$ bit/s: extrapolation des auteurs & \cite{Koch2006} & mesuré \\
12 & L'oreille est un banc de filtres passe-bande avec une carte des fréquences presque logarithmique (fig.~\ref{fig:filterbank}) & Le lieu de vibration maximale de la membrane basilaire dépend de la fréquence; la fonction \og fréquence--lieu\fg{} est presque exponentielle et décrit sous une seule forme les cochlées de l'homme et d'animaux vivants & \cite{Greenwood1990,Robles2001} & mesuré \\
13 & Les résonateurs sont actifs: une partie des capteurs amplifie les vibrations faibles d'un facteur de plusieurs milliers en amplitude ($66$--$76$\,dB); le gain diminue quand le volume augmente & Vibrométrie laser à la base de la cochlée du chinchilla: pour des sons faibles à la fréquence caractéristique, gain de $66$--$76$\,dB par rapport à l'étrier; la mort réduit la sensibilité de $60$--$81$\,dB ($10^3$--$10^4$ fois en amplitude) et supprime la compression. La source de puissance est constituée par les cellules ciliées externes & \cite{Ruggero1997,Robles2001} & mesuré \\
14 & L'amplificateur de l'oreille est placé au bord de l'auto-oscillation & Calcul: près d'une bifurcation de Hopf, compression de la plage, accord fin et sons de combinaison sont inévitables, soit toutes les non-linéarités connues de l'audition. Que la cochlée soit réglée exactement ainsi n'est pas démontré directement & \cite{Eguiluz2000} & hypothèse \\
15 & Aux basses fréquences, les impulsions sont en phase avec le son (chez le cobaye, le verrouillage faiblit à partir de $\sim600$\,Hz et reste perceptible jusqu'à $\sim3.5$\,kHz), et la direction s'en déduit; aux hautes fréquences, il reste le code de position & Nerf auditif du cobaye: le verrouillage de phase commence à décliner vers $600$\,Hz et disparaît au-dessus de $3.5$\,kHz; chez le chat et le singe, la limite est environ une octave plus haut. Différence de temps d'arrivée aux deux oreilles: revue & \cite{PalmerRussell1986,Grothe2010} & mesuré \\
16 & Autres capteurs (tabl.~\ref{tab:sensors}): l'odorat a $\sim400$ types de récepteurs, un par neurone, et un code combinatoire; le goût passe en partie par l'ATP et la sérotonine; le toucher a des capteurs à adaptation rapide et lente; le chaud et le froid sont séparés & Souris: un récepteur reconnaît de nombreuses odeurs, une odeur active de nombreux récepteurs, des odeurs différentes activent des ensembles différents. Sans récepteurs de l'ATP, le nerf gustatif ne répond pas; les bourgeons du goût libèrent de la sérotonine. Enregistrement de fibres cutanées isolées de la main humaine: quatre types selon la vitesse d'adaptation. Le canal du froid est distinct de ceux du chaud & \cite{Mombaerts2004,Malnic1999,Finger2005,Huang2005,JohanssonVallbo1979,McKemy2002} & mesuré \\
\bottomrule
\end{longtable}}

\section{Chapitre~\ref{sec:recognition}. Reconnaissance}

La vitesse de reconnaissance, la structure des premiers étages, la lecture linéaire de la réponse et l'apprentissage sur la continuité du temps sont mesurés chez l'homme et le singe; la réduction de toute la chaîne à deux opérations et l'apprentissage sur la vidéo sont vérifiés sur les modèles du livre; les explications des illusions vont du bien fondé au discuté.

{\footnotesize
\setlength{\tabcolsep}{3pt}\renewcommand{\arraystretch}{1.15}
\begin{longtable}{>{\raggedright}p{0.5cm}>{\raggedright}p{4.0cm}>{\raggedright}p{5.6cm}>{\raggedright}p{2.3cm}>{\raggedright\arraybackslash}p{1.7cm}}
\toprule
N\textsuperscript{o} & Affirmation & Expérience et ce qui est mesuré & Source & Statut \\
\midrule\endfirsthead
\toprule
N\textsuperscript{o} & Affirmation & Expérience et ce qui est mesuré & Source & Statut \\
\midrule\endhead
1 & Sous translation, la comparaison pixel par pixel avec un modèle ne fait pas mieux que pile ou face; la paire d'étages~\eqref{eq:scstage} reconnaît la figure à n'importe quelle position & \texttt{models/\allowbreak recognition\_models.py}, \og rétine\fg{} de $24$ points, $4000$ essais: modèle, $0.49$, $0.51$, $0.52$ pour une proportion de points inversés de $0$, $0.1$, $0.2$; paire $S$ et $C$, $1.00$, $0.86$, $0.70$ & déduction du chapitre & modèle \\
2 & Un animal sur une image est reconnu en $\sim150$\,ms, en un seul passage à travers une dizaine d'étages de $10$--$20$\,ms & Homme, tâche \og animal ou pas\fg{} sur des photos présentées $20$\,ms: les potentiels évoqués des deux réponses divergent après $\sim150$\,ms. Singe: le temps de première réponse croît d'étage en étage (V1 d'abord, puis V2 et V4). Le nombre d'étages et \og zéro ou une impulsion par étage\fg{} sont déduits de ces chiffres & \cite{Thorpe1996,Schmolesky1998} & mesuré \\
3 & L'étage $S$ est un ET sur les parties, l'étage $C$ un maximum sur les positions (proposition~\ref{prop:conveyor}) & Chat, cortex visuel primaire: les cellules simples répondent à un bord d'orientation donnée à un endroit donné, les cellules complexes à la même orientation sur une zone plus grande. Enregistrement intracellulaire de cellules complexes: la réponse à deux barres est, pour beaucoup de cellules, proche de la plus grande des deux réponses, en moyenne très proche de l'opération max. Maximum par inhibition mutuelle: proposition~\ref{prop:wta}; division par l'activité totale: revue & \cite{HubelWiesel1962,Lampl2004,Carandini2012} & mesuré \\
4 & Plus haut dans la chaîne, les combinaisons sont plus grandes et plus complexes, et la réponse dépend de moins en moins de la position & Singe, simplification des stimuli jusqu'au \og trait critique\fg{}: les cellules de V4 et du cortex inférotemporal postérieur répondent à des combinaisons complexes de forme, de couleur et de texture avec un petit champ; dans le cortex inférotemporal antérieur, le champ est plus grand. L'alternance de $S$ et $C$ dans le modèle reproduit ce tableau & \cite{KobatakeTanaka1994,Riesenhuber1999} & mesuré \\
5 & Les réseaux convolutifs reproduisent cette alternance et prédisent le mieux les réponses des étages supérieurs; l'objet se lit linéairement à partir des fréquences d'un groupe de neurones & Réseau entraîné à reconnaître, sans ajustement au cerveau: la couche supérieure prédit les réponses des neurones du cortex inférotemporal du singe, les couches intermédiaires celles de V4. Un classifieur linéaire sur l'activité de $\sim100$ cellules prises au hasard, dans des fenêtres dès $12.5$\,ms, reconnaît l'objet et sa catégorie à différentes positions et tailles & \cite{Fukushima1980,Yamins2014,Hung2005} & mesuré \\
6 & La règle de Hebb avec trace~\eqref{eq:traceinvariance} apprend l'invariance sur une vidéo sans étiquettes si la trace dure au moins $\sim20$\,ms & L'idée des traits lents et de la règle avec trace relève de la théorie. \texttt{models/\allowbreak recognition\_models.py}, deux neurones $C$, $16$ entrées, $10$ essais, pause entre objets de $0.3$\,s. Pour $\tau_{\text{tr}}=5$\,ms, coïncidence avec la figure de $0.52$ en moyenne, pour $10$\,ms, $0.56$; pour $20$, $50$ et $200$\,ms, $1.00$ dans tous les essais (à $30$\,ms, un essai sur dix se trompe) & \cite{Foldiak1991,WiskottSejnowski2002} & modèle \\
7 & Dans le cerveau, l'invariance s'apprend sur la continuité du temps & Singe: on substituait discrètement l'objet pendant un saut de l'œil vers un endroit donné du champ; l'invariance de position des neurones du cortex inférotemporal se modifiait conformément à la substitution, de façon significative dès une heure. Qu'il s'agisse de la règle principale de l'apprentissage visuel est une hypothèse & \cite{LiDiCarlo2008} & mesuré \\
8 & Mouvement: détecteurs de direction, puis la même paire ET/OU sur de grandes zones & Détecteurs de direction dans la rétine du lapin; dans l'aire MT du singe, les $110$ neurones enregistrés sont tous sélectifs à la direction, et la réponse au mouvement y est plus forte que dans V1 & \cite{Barlow1965,Albright1984} & mesuré \\
9 & Les étages supérieurs envoient une prédiction vers le bas, l'erreur monte & Le modèle explique les effets extra-classiques des champs du cortex visuel primaire; certaines observations concordent (revue), mais ce n'est pas démontré comme principe général du pipeline & \cite{RaoBallard1999,Keller2018} & hypothèse \\
10 & Les images difficiles exigent des rétroactions et plusieurs tours & Singe: pour les images \og difficiles\fg{}, la décision dans le cortex inférotemporal se forme $\sim30$\,ms plus tard; ces réponses tardives sont mal prédites par un réseau direct, mieux par des réseaux récurrents ou plus profonds & \cite{Kar2019} & mesuré \\
11 & Une falsification invisible des pixels trompe le réseau; la vision humaine est bien plus robuste, mais pas invulnérable & Sur les réseaux: une petite perturbation choisie tout exprès change la réponse; l'exemple \og panda $\to$ gibbon\fg{} vient du deuxième article. Chez l'homme, en présentation brève, des falsifications qui se transfèrent bien d'un réseau à l'autre infléchissent le choix & \cite{Szegedy2014,Goodfellow2015,Elsayed2018} & mesuré \\
12 & Illusions d'attente: une entrée peu fiable cède à l'attente; ce qui est pâle bouge \og plus lentement\fg{}, la robe est vue de façons différentes (section~\ref{sec:illusions}) & Des réseaux de barres moins contrastés semblent bouger plus lentement. Enquête auprès de $1401$ personnes: la robe est vue selon trois variantes stables, un indice sur l'éclairage fait basculer la couleur; chez $\sim13\,000$ observateurs, c'est l'hypothèse sur l'éclairage qui décide. Un modèle bayésien avec attente de mouvements lents explique plusieurs illusions de mouvement & \cite{Thompson1982,LaferSousa2015,Wallisch2017,Weiss2002,Purves2011} & hypothèse \\
13 & Réglage du gain: cascade, inclinaison et contraste; la grille d'Hermann ne s'explique pas par l'inhibition des voisins & Les détecteurs de direction de la rétine du lapin, lors d'un mouvement prolongé, baissent leur fréquence et, après l'arrêt, descendent sous le niveau de fond. L'illusion d'inclinaison est reproduite par un modèle de division par l'activité des voisins. Des modifications de la grille qui ne changent pas la réponse des neurones de sortie de la rétine suppriment les taches: l'explication par inhibition des voisins est rejetée & \cite{BarlowHill1963,Schwartz2009,SchillerCarvey2005} & mesuré \\
14 & Compensation des délais: un objet en mouvement semble en avance sur l'éclair & L'illusion est mesurée. La psychophysique contredit à la fois le prolongement du mouvement et la différence de délais; une explication a posteriori est proposée, par les événements survenant $\sim80$\,ms après l'éclair. Le débat n'est pas tranché & \cite{Nijhawan1994,EaglemanSejnowski2000} & hypothèse \\
15 & Des \og serpents\fg{} immobiles tournent: la différence de délais~\eqref{eq:latency} est lue par les détecteurs de direction & L'illusion fonctionne aussi sans couleur; des paires d'éléments voisins la produisent isolément; des neurones corticaux du singe sélectifs à la direction donnent une réponse orientée à ces paires immobiles. Chez l'homme, la rotation est déclenchée par les microsaccades et les clignements & \cite{Conway2005,OteroMillan2012} & mesuré \\
16 & Images ambiguës: inhibition mutuelle, fatigue et bruit; les bascules sont surtout dues au bruit & Modèle de groupes de neurones en compétition: les bascules y sont causées par le \emph{bruit}, sans bruit il n'y en a pas; il explique la hausse de la fréquence des bascules avec l'intensité du stimulus. La fatigue y joue un rôle moindre & \cite{MorenoBote2007} & hypothèse \\
17 & Une partie des illusions découle de la tâche qu'apprend la machine & Un réseau entraîné à prédire l'image suivante d'une vidéo \og voit\fg{} la rotation d'anneaux immobiles et ne la voit pas sur les images témoins; des réseaux de débruitage et d'augmentation de la netteté cèdent aux illusions de couleur et de luminosité, comme l'homme & \cite{Watanabe2018,GomezVilla2020} & mesuré \\
\bottomrule
\end{longtable}}

\section{Chapitre~\ref{sec:muscles}. La sortie vers les muscles}

La structure de la sortie --- unités motrices, marge de sécurité de la synapse, activation par taille, paire de muscles, boucle d'étirement et générateurs de rythme --- est mesurée directement; la condition de stabilité de la boucle à retard est un théorème; le modèle interne du corps est une hypothèse bien fondée; les chiffres des figures sont calculés avec le modèle du livre.

{\footnotesize
\setlength{\tabcolsep}{3pt}\renewcommand{\arraystretch}{1.15}
\begin{longtable}{>{\raggedright}p{0.5cm}>{\raggedright}p{4.0cm}>{\raggedright}p{5.6cm}>{\raggedright}p{2.3cm}>{\raggedright\arraybackslash}p{1.7cm}}
\toprule
N\textsuperscript{o} & Affirmation & Expérience et ce qui est mesuré & Source & Statut \\
\midrule\endfirsthead
\toprule
N\textsuperscript{o} & Affirmation & Expérience et ce qui est mesuré & Source & Statut \\
\midrule\endhead
1 & Unité motrice: un neurone \nt{ACET} commande un groupe de fibres, de quelques centaines à deux mille environ; dans les muscles à réglage fin, les unités sont plus petites & Muscles humains, comptage des axones moteurs et des fibres musculaires: en moyenne environ $340$ fibres par neurone dans un muscle interosseux de la main, environ $410$ dans le muscle brachio-radial, environ $1930$ dans le gastrocnémien médial. Le chapitre ne donne pas de chiffre précis pour les plus petites unités. & \cite{Feinstein1955mu} & mesuré \\
2 & La synapse sur le muscle libère plusieurs fois plus de neurotransmetteur qu'il n'en faut pour déclencher la fibre & Revue des mesures sur les synapses neuromusculaires: la marge de sécurité (rapport du neurotransmetteur libéré au seuil) vaut en général $3$--$5$ chez les mammifères adultes; chez l'homme, elle repose davantage sur les replis de la membrane du destinataire que sur l'importance de la libération. & \cite{WoodSlater2001} & mesuré \\
3 & La secousse~\eqref{eq:twitch} monte en un temps $T_c$ de l'ordre de $50\ms$ & Unités motrices isolées de la main humaine lors d'un effort volontaire, force moyennée sur les impulsions de l'unité: temps de montée de $30$--$100\ms$, inférieur à $70\ms$ pour $80\,\%$ des unités. La fonction $(t/T_c)e^{1-t/T_c}$ est une approximation tirée d'un modèle de pool d'unités. & \cite{MilnerBrown1973rec, Fuglevand1993} & mesuré \\
4 & Somme des secousses~\eqref{eq:forcesum}: à $5$, $15$ et $40$ imp./s, force de $0.2$--$1.1F_1$, $1.8$--$2.2F_1$ et environ $5.4F_1$ (fig.~\ref{fig:tetanus}) & Calcul avec \texttt{models/\allowbreak muscle\_\allowbreak models.py}, $T_c=50\ms$: $0.21$--$1.10$, $1.76$--$2.19$ et $5.28$--$5.49\,F_1$ sur les dernières $0.3\,\text{s}$. L'addition linéaire est une simplification: le modèle ne contient pas la saturation d'un vrai muscle. & \texttt{models/\allowbreak muscle\_\allowbreak models.py} & modèle \\
5 & Les unités s'activent par ordre de taille; les plus faibles sont cent fois plus faibles que les plus fortes & Racines ventrales du chat: sous une entrée réflexe croissante, les premiers neurones à décharger sont ceux dont les impulsions sont petites dans la racine, c'est-à-dire les petits. Muscle interosseux de la main humaine: force de secousse des unités de $0.1$ à $10$~g, croissant presque linéairement avec la force à laquelle l'unité s'active. & \cite{Henneman1957, MilnerBrown1973rec} & mesuré \\
6 & Le pas de force est proportionnel à la force, $\Delta F/F\approx q-1$~\eqref{eq:sizeprinciple}: virgule flottante & Déduit dans le chapitre de la progression géométrique des forces. Concorde avec l'expérience: aux efforts élevés, bien moins d'unités nouvelles s'activent pour un même incrément de force. & déduction du chapitre; \cite{MilnerBrown1973rec} & théorie \\
7 & La dispersion de la force croît proportionnellement à la force & Effort isométrique volontaire d'un muscle du pouce: l'écart type de la force croît linéairement avec la force moyenne; pour le même effort provoqué par stimulation électrique, la dispersion est constante. Modèle de pool: la croissance linéaire vient de l'activation des unités par ordre de force. & \cite{Jones2002sdn} & mesuré \\
8 & La fluidité des mouvements découle d'un bruit dépendant du signal & Modèle: la trajectoire qui minimise la dispersion du point d'arrivée sous un tel bruit reproduit les trajectoires mesurées de l'œil et du bras. Aucune expérience directe ne distingue cette cause des autres. & \cite{HarrisWolpert1998} & hypothèse \\
9 & La différence des forces d'une paire de muscles fixe le couple, leur somme la raideur~\eqref{eq:antagonists} & Théorie du contrôle de l'impédance par co-contraction. Bras humain: la raideur de chaque articulation, mesurée par de petites poussées, est à peu près proportionnelle au couple qui s'y exerce; différents rapports de co-contraction changent la forme et l'orientation de l'ellipse de raideur de la main. & \cite{Hogan1984, GomiOsu1998stiff} & mesuré \\
10 & Le capteur d'étirement excite son propre neurone \nt{ACET} et, par un intermédiaire inhibiteur, coupe l'antagoniste (fig.~\ref{fig:antagonists}) & Moelle épinière du chat: un seul intermédiaire, excité par les fibres des capteurs d'étirement, produit des potentiels inhibiteurs monosynaptiques dans les neurones \nt{ACET} du muscle antagoniste, délai synaptique $0.28$--$0.42\ms$. Le neurotransmetteur de l'intermédiaire (la glycine) n'a pas été déterminé dans cette expérience. & \cite{Jankowska1972Ia} & mesuré \\
11 & Retard de la boucle: $25$--$30\ms$ pour la boucle spinale, une fois et demie à trois fois plus par le cortex, $100$--$200\ms$ par l'œil & Bras humain, poussée sur l'articulation: réponse musculaire courte après $20$--$45\ms$, longue après $45$--$105\ms$, volontaire après $120$--$180\ms$. Décalage de la position visible du doigt pendant le mouvement: correction après $160\ms$ en moyenne. & \cite{Pruszynski2008rapid, SaundersKnill2003vis} & mesuré \\
12 & $\dd x/\dd t=-Gx(t-d)$ est stable pour $Gd<\pi/2$~\eqref{eq:delaystable}; à la limite, la fréquence est $1/(4d)$ & Théorème sur les racines d'une équation à retard. Vérification avec \texttt{models/\allowbreak muscle\_\allowbreak models.py}, $d=30\ms$: à $3\,\text{s}$, amplitude de $3\cdot10^{-6}$ pour $G=40$, $0.13$ pour $G=50$, $38$ pour $G=55$ par seconde (limite $52$). & \cite{Hayes1950} & théorie \\
13 & Tremblement physiologique de $8$--$12$~Hz; la boucle d'étirement est l'une de ses sources & Revue des mesures: un petit tremblement vers $10$~Hz existe chez les personnes en bonne santé; son origine est mixte: oscillations centrales, propriétés des décharges des unités, résonance mécanique et résonance de la boucle réflexe, et la part dominante varie selon les conditions. & \cite{McAuleyMarsden2000} & hypothèse \\
14 & Le cerveau prédit le résultat d'une commande grâce à un modèle interne du corps & Un sujet tient une charge entre deux doigts et bouge le bras: sous charge inertielle, visqueuse et élastique, la force de préhension varie en même temps que la charge, et non avec le retard de la boucle; la charge est donc prédite. Une revue rassemble ces expériences; on n'observe pas directement le modèle lui-même dans les neurones. & \cite{FlanaganWing1997grip, WolpertGhahramani2000} & hypothèse \\
15 & Le rythme de la marche, de la respiration, de la mastication est produit par des réseaux locaux sous entrée constante & Revue d'expériences: un système nerveux isolé (moelle épinière, ganglions) produit un motif moteur rythmique sans entrée rythmique et sans rétroaction des muscles. & \cite{MarderBucher2001} & mesuré \\
16 & L'inhibition mutuelle avec fatigue~\eqref{eq:halfcentre} donne un rythme de période $0.84\,\text{s}\approx2\tau_a$ (fig.~\ref{fig:halfcentre}) & Théorème: un réseau à inhibition mutuelle et adaptation sans état stable oscille nécessairement sous entrée constante. Calcul avec \texttt{models/\allowbreak muscle\_\allowbreak models.py} ($I=1$, $\beta=2$, $b=2$, $\tau=20\ms$, $\tau_a=0.4\,\text{s}$): période de $0.84\,\text{s}$. Que les vrais générateurs soient construits ainsi n'est pas démontré. & \cite{Matsuoka1985osc} & modèle \\
\bottomrule
\end{longtable}}

\section{Chapitre~\ref{sec:pain}. La douleur}

Les capteurs de la douleur, les deux fils, le portillon de la moelle épinière, le bouton de volume descendant, la sensibilisation et la capture de l'attention sont mesurés directement; les formules du portillon et de la sensibilisation sont des modèles simplifiés du livre; la règle selon laquelle la machine choisit le volume de l'alarme est une hypothèse.

{\footnotesize
\setlength{\tabcolsep}{3pt}\renewcommand{\arraystretch}{1.15}
\begin{longtable}{>{\raggedright}p{0.5cm}>{\raggedright}p{4.0cm}>{\raggedright}p{5.6cm}>{\raggedright}p{2.3cm}>{\raggedright\arraybackslash}p{1.7cm}}
\toprule
N\textsuperscript{o} & Affirmation & Expérience et ce qui est mesuré & Source & Statut \\
\midrule\endfirsthead
\toprule
N\textsuperscript{o} & Affirmation & Expérience et ce qui est mesuré & Source & Statut \\
\midrule\endhead
1 & Seuil élevé: les seuils d'échauffement des différents capteurs de la douleur vont de $41^\circ$ à $46$--$48^\circ$ & Enregistrement de fibres isolées. Peau de singe: seuil médian de $41^\circ$ pour les capteurs lents (C), $46^\circ$ pour les capteurs rapides de type II, plus de $53^\circ$ pour le type I. Peau humaine: $40.7^\circ$ pour les capteurs C qui répondent aussi à la pression, $48^\circ$ pour ceux qui n'y répondent pas. & \cite{Treede1995heat, Weidner1999cfib} & mesuré \\
2 & Les capteurs de la douleur s'adaptent bien moins que les capteurs du toucher et deviennent plus sensibles après une lésion légère; un affaiblissement après un fort échauffement est mesuré pour un seul type & Brûlure légère de la peau ($50^\circ$, $100\,\text{s}$): après $5$--$10$~min, le seuil de douleur chez l'homme et le seuil des capteurs C chez le singe sont abaissés de $2$--$6^\circ$; les autres capteurs cutanés ne présentent pas ce décalage. Chez les capteurs rapides de type II, la réponse est supprimée après un fort échauffement. La comparaison de l'accoutumance avec les capteurs du toucher est une estimation générale, non étayée ici par une expérience propre. & \cite{LaMotte1982hyper, Treede1995heat} & mesuré \\
3 & Deux fils: rapide, $5$--$30$~m/s; lent, $0.5$--$2$~m/s; temps d'arrivée $t=L/v$~\eqref{eq:paintime} & Vitesse de conduction: capteurs rapides de la douleur du singe, en moyenne $15$ et $25$~m/s (deux types); capteurs C humains, $0.8$--$1.0$~m/s. Pour $L=1$~m, cela fait quelques dizaines de millisecondes et environ une seconde. & \cite{Treede1995heat, Weidner1999cfib} & mesuré \\
4 & La douleur arrive deux fois: d'abord rapide et précise, puis sourde et durable & Temps de réaction à l'échauffement de l'avant-bras humain: de $1100$ à $400\ms$ quand la température monte; un fort échauffement de la peau velue donne une double douleur, celui de la paume non. Seules des fibres plus rapides que $6$~m/s peuvent porter la première douleur; par la latence, elle correspond aux capteurs rapides de type II, absents de la paume. Le partage des rôles (\og où\fg{} et \og à quel point c'est grave\fg{}) est une interprétation. & \cite{CampbellLaMotte1983first, Treede1995heat} & mesuré \\
5 & Portillon: le neurone de sortie de la moelle reçoit la douleur et le toucher, l'intermédiaire inhibiteur est excité par le toucher (fig.~\ref{fig:gate}) & Le schéma du portillon a été proposé comme théorie. Souris, marquage génétique et suppression de groupes de neurones: les neurones excitateurs à somatostatine sont nécessaires à la douleur mécanique; les neurones inhibiteurs à dynorphine empêchent l'entrée des capteurs du toucher de les atteindre; sans ces neurones, un toucher léger provoque de la douleur. & \cite{MelzackWall1965, Duan2014} & mesuré \\
6 & Le toucher atténue la douleur & Homme, impulsions laser douloureuses sur la main: un toucher simultané réduit la douleur ressentie et la capacité à distinguer l'énergie des impulsions; l'effet faiblit avec la distance entre le point de toucher et le point de douleur au sein d'un même segment cutané. & \cite{Mancini2014touch} & mesuré \\
7 & Hypothèse de la théorie du portillon: une douleur forte éteint elle-même l'intermédiaire inhibiteur, et le portillon s'ouvre grand & Présenté dans le chapitre comme une hypothèse de la théorie originale. Le travail sur la souris décrit comment le portillon empêche le toucher d'entrer dans le canal de la douleur; l'extinction de l'intermédiaire par la douleur n'y est pas montrée. & \cite{MelzackWall1965} & hypothèse \\
8 & Portillon~\eqref{eq:gate}: seuil de $0.18$ sans toucher, $0.86$ avec toucher, $0.11$ pour $g=2$; pour $\nu=1$, le toucher réduit la sortie de $1$ à $0.25$, pour $\nu=2$ il est sans effet; le toucher seul ne passe pas (fig.~\ref{fig:gatecurves}) & Le calcul avec \texttt{models/\allowbreak pain\_\allowbreak gate.py} et les poids du texte reproduit tous ces nombres. & \texttt{models/\allowbreak pain\_\allowbreak gate.py} & modèle \\
9 & Les voies descendantes du tronc cérébral modifient le gain du portillon dans les deux sens & Rat, bulbe rachidien, enregistrement pendant l'échauffement de la queue: certains neurones déchargent avant le retrait (cellules \og ON\fg{}), d'autres se taisent (cellules \og OFF\fg{}); une faible stimulation de cette région supprime le réflexe. Revue: les mêmes circuits facilitent et inhibent la transmission de la douleur, et les opioïdes agissent par eux. & \cite{Fields1983rvm, Fields2004} & mesuré \\
10 & L'attente d'un soulagement atténue la douleur par le canal opioïde & Patients souffrant d'une douleur aiguë: chez ceux dont l'attente d'un soulagement avait réduit la douleur, le blocage des récepteurs opioïdes a ramené la douleur au niveau des autres; chez les autres, le blocage ne changeait pas la douleur. & \cite{Levine1978} & mesuré \\
11 & Le cerveau choisit le volume de l'alarme selon l'intérêt de l'interruption & Il n'existe pas d'expérience mesurant la règle du choix du volume. & --- & hypothèse \\
12 & Sensibilisation: après une lésion, la sortie du portillon augmente et le seuil baisse, en partie grâce à la moelle épinière; elle persiste après l'arrêt du stimulus nocif & Rat: après une lésion cutanée, le seuil du réflexe de flexion baisse et la réponse augmente; l'analyse électrophysiologique montre qu'une partie de cette hausse naît dans la moelle elle-même. Un stimulus nocif provoque des troubles durables de la sensibilité qui survivent au stimulus. Le chapitre ne donne pas le temps de décroissance exact. & \cite{Woolf1983} & mesuré \\
13 & La sensibilisation~\eqref{eq:sensitization} est une réaction positive; si la boucle est forte, l'alarme peut se verrouiller après la guérison & Découle du modèle~\eqref{eq:sensitization} (exercice en fin de chapitre). Aucune expérience ne montre qu'une douleur persistante après guérison soit une boucle verrouillée de ce type précis. & déduction du chapitre & hypothèse \\
14 & La douleur est une récompense négative, et l'apprentissage par la douleur suit l'erreur de différence temporelle & IRM, sujets humains, chaînes de signes annonciateurs de la douleur: l'activité du striatum ventral et de l'insula antérieure coïncide avec les signaux d'apprentissage par différence temporelle. Singe: une partie des neurones \nt{DOPA} est excitée par un souffle d'air désagréable et son signe annonciateur, une autre partie est inhibée. & \cite{Seymour2004, Matsumoto2009} & mesuré \\
15 & La douleur interrompt le travail en cours & Sujets humains, stimulus électrique douloureux et sondes sonores: la réponse à une sonde présentée $100\ms$ après le début de la douleur est nettement ralentie; à $1.5\,\text{s}$, il n'y a plus de différence avec le contrôle. Une revue rassemble ces expériences dans un modèle de capture de l'attention. & \cite{Crombez1996disrupt, EcclestonCrombez1999} & mesuré \\
16 & Le retrait se referme dans la moelle épinière & Rat: les neurones des couches profondes de la corne dorsale reproduisent le motif spatial du réflexe de retrait de muscles isolés; on ne parvient pas à les exciter de façon antidromique depuis la région cervicale: ce sont donc des intermédiaires spinaux. & \cite{Schouenborg1995wr} & mesuré \\
\bottomrule
\end{longtable}}

\section{Chapitre~\ref{sec:motorlearning}. L'apprentissage moteur}

La règle des moindres carrés et l'avantage d'un fil d'erreur séparé sont des théorèmes; la structure du circuit à fil d'erreur, l'affaiblissement par coïncidence, l'ajustement de la trace au retard, l'effet consécutif et le retour spontané de l'habileté sont mesurés; que le circuit entier fonctionne comme un apprentissage supervisé et construise un modèle interne est l'hypothèse dominante; les chiffres du modèle à deux processus sont calculés avec le modèle du livre.

{\footnotesize
\setlength{\tabcolsep}{3pt}\renewcommand{\arraystretch}{1.15}
\begin{longtable}{>{\raggedright}p{0.5cm}>{\raggedright}p{4.0cm}>{\raggedright}p{5.6cm}>{\raggedright}p{2.3cm}>{\raggedright\arraybackslash}p{1.7cm}}
\toprule
N\textsuperscript{o} & Affirmation & Expérience et ce qui est mesuré & Source & Statut \\
\midrule\endfirsthead
\toprule
N\textsuperscript{o} & Affirmation & Expérience et ce qui est mesuré & Source & Statut \\
\midrule\endhead
1 & La règle~\eqref{eq:lms} suit en moyenne le gradient du carré de l'erreur & Résultat de la théorie des filtres adaptatifs: une correction proportionnelle à l'entrée et à l'erreur de sortie est une descente stochastique du gradient de l'erreur quadratique moyenne. & \cite{WidrowHoff1960} & théorie \\
2 & Avec un fil d'erreur par sortie, la vitesse ne dépend pas du nombre de sorties; avec un seul nombre commun, elle baisse avec le nombre de neurones bruités (proposition~\ref{prop:teachers}) & Courbes d'apprentissage d'un réseau linéaire: l'apprentissage par perturbations aléatoires des neurones est plus lent que le gradient exact d'un facteur égal au nombre de neurones perturbés. & \cite{Werfel2005} & théorie \\
3 & Le circuit à fil d'erreur fonctionne comme un apprentissage supervisé (proposition~\ref{prop:errorwire}) & Théories déduites de la structure du circuit. Les expériences des lignes 7--10 concordent avec elles; que le circuit entier fonctionne exactement ainsi n'est pas démontré. & \cite{Marr1969, Albus1971} & hypothèse \\
4 & Couche d'expansion: environ $7\cdot10^{10}$ petits neurones \nt{GLUT}, plus que dans tout le reste du cerveau & Comptage des noyaux dans une suspension de tissu dissous: $69\cdot10^9$ neurones dans le cervelet humain sur $86\cdot10^9$ dans tout le cerveau. Stéréologie: $101\cdot10^9$ petits neurones dans le cervelet d'hommes âgés. & \cite{Azevedo2009, Andersen1992cereb} & mesuré \\
5 & Augmenter la dimension de l'entrée rend linéaire la dépendance voulue & Théorème sur le nombre de dichotomies: une somme pondérée de $n$ entrées avec seuil réalise un étiquetage arbitraire d'environ $2n$ vecteurs en position générale. Que le cervelet exploite précisément cela fait partie de l'hypothèse de la ligne~3. & \cite{Cover1965} & théorie \\
6 & Environ $3\cdot10^7$ neurones \nt{GABA} de sortie, chacun avec de l'ordre de $10^5$ entrées & Stéréologie du cervelet humain: $30.5\cdot10^6$ neurones de sortie. Microscopie électronique, rat: environ $1.75\cdot10^5$ entrées de la couche d'expansion par neurone de sortie. Le neurotransmetteur inhibiteur des neurones de sortie: dans la revue. & \cite{Andersen1992cereb, NapperHarvey1988, Ito2001} & mesuré \\
7 & Exactement un fil d'erreur par neurone de sortie, environ une fois par seconde, avec chaque fois une décharge puissante & Revue d'enregistrements: chaque neurone de sortie reçoit une seule entrée \nt{GLUT} grimpante, qui provoque une impulsion complexe à une fréquence de l'ordre de $1$~Hz. & \cite{Ito2001} & mesuré \\
8 & La probabilité de la décharge dépend de la direction de l'erreur de mouvement & Singe, région oculomotrice du cervelet, adaptation des saccades: la direction de l'erreur visuelle fixe la probabilité de l'impulsion complexe, sa grandeur fixe son instant; la décharge modifie les impulsions simples à l'essai suivant et déplace le mouvement le long de l'erreur préférée de la cellule. & \cite{Herzfeld2018} & mesuré \\
9 & La coïncidence de la décharge d'erreur avec une entrée active affaiblit cette entrée ($-\eta e_i x_j$) & Revue d'expériences: la stimulation conjointe d'une entrée de la couche d'expansion et du fil d'erreur provoque un affaiblissement durable de cette entrée; la stimulation de l'entrée seule, non. & \cite{Ito2001} & mesuré \\
10 & La durée de la trace est ajustée au retard de l'erreur: pour un retard de $100\ms$, c'est l'entrée active $100\ms$ avant qui s'affaiblit le plus & Tranches et souris vivantes: dans la région chargée de l'apprentissage oculomoteur, l'affaiblissement est étroitement accordé sur un intervalle d'environ $120\ms$ entre l'entrée et l'erreur, soit le temps que met le signal d'erreur dans cette tâche; dans une autre région, différentes cellules sont accordées sur différents intervalles. & \cite{Suvrathan2016} & mesuré \\
11 & Le circuit est un filtre adaptatif~\eqref{eq:adaptivefilter} qui apprend un modèle interne du corps & Modèle déduit de la structure du circuit. Aucune expérience directe n'a mesuré le filtre appris comme fonction de transfert du corps. & \cite{Fujita1982} & hypothèse \\
12 & La prédiction soustrait les conséquences de nos propres mouvements, c'est pourquoi on ne peut pas se chatouiller soi-même & Sujets humains: un toucher par soi-même semble moins chatouilleux que le même toucher par autrui; en IRM, le cortex somatosensoriel répond plus faiblement au toucher par soi-même, et le cervelet est moins actif quand le mouvement touche la peau. L'explication par soustraction de la prédiction est une hypothèse. & \cite{Blakemore1998} & hypothèse \\
13 & Sous une force extérieure, le raté diminue, et après sa suppression la main rate de l'autre côté & Des sujets tendent le bras en tenant la poignée d'un robot dans un champ de forces: les trajectoires redeviennent droites; après la suppression du champ, l'effet consécutif est presque l'image en miroir des premières déformations; il se transfère aussi à des régions non entraînées de l'espace. & \cite{ShadmehrMussaIvaldi1994} & mesuré \\
14 & Deux processus apprennent~\eqref{eq:tworate}; après une perturbation inverse, l'habileté revient d'elle-même & Sujets humains, champ de forces, puis bref champ inverse et essais à erreur bloquée: la correction revient d'elle-même vers l'ancienne; le modèle à deux processus le reproduit, le modèle à un processus non. & \cite{Smith2006} & mesuré \\
15 & Chiffres de la fig.~\ref{fig:tworate}: $0.84$ (lent $0.63$) en $200$ mouvements; $6$ inverses; rapide $-0.53$, lent $0.46$; retour jusqu'à $0.37$ en $40$ mouvements & Calcul avec \texttt{models/\allowbreak motor\_\allowbreak learning.py}, $A_f=0.92$, $B_f=0.1$, $A_s=0.996$, $B_s=0.02$: $0.840$ ($0.634$ et $0.205$), $6$ mouvements, $-0.531$ et $0.457$, pic de $0.371$ au $40$\textsuperscript{e} mouvement. & \texttt{models/\allowbreak motor\_\allowbreak learning.py} & modèle \\
16 & Deux processus sont nécessaires parce que le corps change à des vitesses différentes & Théorie: l'estimation optimale sous des perturbations à plusieurs durées de vie donne plusieurs filtres de constantes de temps différentes et explique le retour spontané et le réapprentissage accéléré. & \cite{Kording2007} & hypothèse \\
\bottomrule
\end{longtable}}

\section{Chapitre~\ref{sec:chemoutput}. La sortie chimique}

Les deux fils vers le cœur et leurs vitesses différentes, la contre-réaction dans les cascades hormonales, le code par fréquence des bouffées, le générateur circadien et son recalage par la lumière sont mesurés directement; la limite $K<8$ est un théorème de l'automatique démontré dans le chapitre; les pulsations engendrées par la rétroaction de la cascade elle-même sont une hypothèse soutenue par une expérience solide.

{\footnotesize
\setlength{\tabcolsep}{3pt}\renewcommand{\arraystretch}{1.15}
\begin{longtable}{>{\raggedright}p{0.5cm}>{\raggedright}p{4.0cm}>{\raggedright}p{5.6cm}>{\raggedright}p{2.3cm}>{\raggedright\arraybackslash}p{1.7cm}}
\toprule
N\textsuperscript{o} & Affirmation & Expérience et ce qui est mesuré & Source & Statut \\
\midrule\endfirsthead
\toprule
N\textsuperscript{o} & Affirmation & Expérience et ce qui est mesuré & Source & Statut \\
\midrule\endhead
1 & Le stimulateur du cœur bat de lui-même environ $100$ fois par minute; au repos, le frein est enfoncé, et la fréquence est en général de l'ordre de $60$--$70$ & Humains, blocage complet des deux fils vers le cœur: la fréquence propre dépasse la fréquence de repos et décroît avec l'âge; chez l'adulte, elle est de l'ordre de $100$ battements par minute. Au repos, c'est donc le frein qui domine. La fréquence de repos de $60$--$70$ est une valeur typique, non citée séparément. & \cite{JoseCollison1970ihr} & mesuré \\
2 & L'accélérateur \nt{NORA} et le frein \nt{ACET} sont des filtres passe-bas, le frein est plus rapide~\eqref{eq:gasbrake} & Chien, stimulation modulée en fréquence du nerf vague ou du nerf sympathique cardiaque, et fonction de transfert vers la fréquence auriculaire: les deux voies sont des filtres passe-bas; la voie sympathique a une fréquence de coupure bien plus basse et un retard pur d'environ $1.7\,\text{s}$ en plus. Les caractéristiques dépendent du tonus moyen, si bien que la formule linéaire~\eqref{eq:gasbrake} est une approximation. & \cite{Berger1989} & mesuré \\
3 & Le frein est rapide parce qu'\nt{ACET} ouvre le canal potassique sans second messager, alors que \nt{NORA} passe par une chaîne de réactions & Fragments de membrane de cellules auriculaires: le canal potassique ouvert par l'acétylcholine est ouvert par les sous-unités $\beta\gamma$ de la protéine G couplée au récepteur, sans second messager intracellulaire. La voie de \nt{NORA} par l'AMPc n'est pas citée ici. & \cite{Logothetis1987gbg} & mesuré \\
4 & Le sang fait le tour du corps en une minute environ; \nt{ADRE} circulant atteint tous les organes munis du récepteur & Estimation générale d'ordre de grandeur; formulée ainsi dans le chapitre, sans source citée. & --- & estimation \\
5 & La cascade hormonale est bouclée par une contre-réaction: l'hormone finale freine les premiers étages & Revue d'expériences: les hormones corticosurrénaliennes répriment la libération d'ACTH par l'hypophyse et de l'hormone de libération par l'hypothalamus dans trois gammes de temps: rapide, intermédiaire et lente. & \cite{KellerWood1984fb} & mesuré \\
6 & Trois étages identiques sont stables pour $K<8$~\eqref{eq:cascadestable}; à la limite, la période vaut $2\pi\tau/\sqrt3\approx3.6\tau$ & Démontré dans le chapitre: à la pulsation $\sqrt3/\tau$, trois étages donnent un déphasage de $180^\circ$ et une atténuation d'un facteur $8$. & démonstration dans le chapitre & théorie \\
7 & Erreur sur le niveau $1/(1+K)$: un tel régulateur ne tient donc pas mieux qu'environ $10\,\%$ (proposition~\ref{prop:cascade}) & Conséquence de la ligne~6 pour une rétroaction proportionnelle. L'axe réel a une rétroaction sur plusieurs étages et dans plusieurs gammes de temps (ligne~5): la limite se rapporte donc au modèle~\eqref{eq:cascade} et n'est pas mesurée. & démonstration dans le chapitre & théorie \\
8 & Pour $K=4$ et $7$, le niveau se stabilise; pour $K=12$, pulsations régulières de période $3.7$--$3.9\,\tau$ (fig.~\ref{fig:cascade}) & Calcul avec \texttt{models/\allowbreak chem\_\allowbreak models.py}: pour $K=12$, périodes successives $3.9$, $3.7$, $3.8$, $3.7\,\tau$; pour $K=4$, amplitude en fin de calcul $0.003$. & \texttt{models/\allowbreak chem\_\allowbreak models.py} & modèle \\
9 & L'hormone du stress sort par pulsations environ une fois par heure; les pulsations naissent de la rétroaction elle-même, sans générateur séparé & Rats, perfusion constante de l'hormone de libération: l'ACTH et la corticostérone oscillent à une fréquence physiologique, proche d'une heure; une entrée rythmique venant de l'hypothalamus n'est pas nécessaire aux pulsations. Un modèle hypophyse--surrénale à rétroaction retardée produit de telles oscillations. & \cite{Walker2012glc, Walker2010} & hypothèse \\
10 & Le destinataire lit la fréquence des bouffées, non le niveau & Singes privés de leur propre libération de l'hormone de libération des gonadotrophines: une perfusion constante ne rétablit pas la sécrétion des hormones hypophysaires, des bouffées toutes les heures la rétablissent; le retour à la perfusion constante éteint de nouveau la réponse. La fatigue du récepteur vue comme filtre passe-haut est une interprétation. & \cite{Belchetz1978} & mesuré \\
11 & Des fréquences de bouffées différentes agissent différemment sur les diverses cellules d'une même glande & Mêmes singes: porter les bouffées à $2$--$5$ par heure fait baisser les deux hormones hypophysaires; les espacer à une toutes les $3$~heures fait baisser l'une (LH) et monter l'autre (FSH), si bien que leur rapport est fixé par la fréquence. & \cite{Wildt1981gnrh} & mesuré \\
12 & Le rythme circadien naît d'une contre-réaction intracellulaire retardée de plusieurs heures & Revue: les protéines des gènes d'horloge répriment avec retard la transcription de leurs propres gènes; la boucle transcription--traduction donne une période d'environ un jour. & \cite{TakahashiJS2017} & mesuré \\
13 & Le générateur principal est un groupe de quelques dizaines de milliers de neurones qui synchronise le reste du corps & Revue: le noyau suprachiasmatique est le générateur circadien principal; ses neurones isolés en culture sont des générateurs autonomes, et le réseau les synchronise; chez la souris, environ $10^4$ neurones de chaque côté. & \cite{Welsh2010scn} & mesuré \\
14 & La période propre chez l'humain est de $24.18$~heures & Humains en lumière faible, avec un horaire découplé du jour: période des rythmes de la mélatonine, de la température et du cortisol en moyenne de $24.18$~heures chez les jeunes comme chez les âgés, avec une faible dispersion. & \cite{Czeisler1999} & mesuré \\
15 & La lumière recale le générateur comme une boucle à verrouillage de phase~\eqref{eq:pll}; il faut un décalage d'environ $11$~min par jour & Humains, une seule exposition à une lumière vive de $6.7$~heures à différentes heures: courbe de réponse de phase de type~1, d'amplitude $5$~heures; retard avant le minimum de la température corporelle, avance après. L'équation~\eqref{eq:pll} est le modèle standard; $11$~min $=0.18$~heure est de l'arithmétique. & \cite{Khalsa2003prc} & mesuré \\
16 & Sans lumière, pas de verrouillage: le rythme dérive vers sa période propre & Personnes totalement aveugles, sans perception de la lumière, vivant en société normale: chez environ la moitié d'entre elles, les rythmes de la mélatonine et du cortisol suivent leur propre période, différente de celle du jour. & \cite{Sack1992blind} & mesuré \\
\bottomrule
\end{longtable}}

\section{Chapitre~\ref{sec:debugging}. Déboguer la machine}

Ce qu'entend une électrode, la faible dimension de l'activité, le fonctionnement des interfaces à grilles rigides, l'apprentissage conjoint du cerveau et du décodeur, et les points visuels produits par stimulation sont mesurés dans des expériences publiées avec comité de lecture; les estimations des flux de données sont l'arithmétique du chapitre; sur l'appareil de Neuralink implanté chez l'humain, pour autant que nous ayons pu le vérifier, les données ont surtout été publiées par la société elle-même.

{\footnotesize
\setlength{\tabcolsep}{3pt}\renewcommand{\arraystretch}{1.15}
\begin{longtable}{>{\raggedright}p{0.5cm}>{\raggedright}p{4.0cm}>{\raggedright}p{5.6cm}>{\raggedright}p{2.3cm}>{\raggedright\arraybackslash}p{1.7cm}}
\toprule
N\textsuperscript{o} & Affirmation & Expérience et ce qui est mesuré & Source & Statut \\
\midrule\endfirsthead
\toprule
N\textsuperscript{o} & Affirmation & Expérience et ce qui est mesuré & Source & Statut \\
\midrule\endhead
1 & L'électrode entend les impulsions des neurones situés dans un rayon d'une centaine de micromètres, en général d'un seul à quelques-uns; on les distingue par leur forme & Rat, région CA1, enregistrement intracellulaire et extracellulaire simultané: la forme de l'impulsion extracellulaire reproduit la largeur et l'amplitude de l'impulsion intracellulaire. Estimation: une tétrode pourrait distinguer $60$--$100$ neurones, alors qu'en pratique on en isole environ six. & \cite{Henze2000extra} & mesuré \\
2 & Le cerveau compte $8.6\cdot10^{10}$ neurones; la sonde en entend environ un cent-millionième & Comptage des noyaux dans une suspension de tissu dissous: $86\cdot10^9$ neurones dans le cerveau humain. La fraction est l'arithmétique du chapitre. & \cite{Azevedo2009} & mesuré \\
3 & L'activité de centaines de neurones du cortex moteur tient dans une dizaine ou deux de variables communes & Revue d'enregistrements dans le cortex moteur du singe: l'activité de la population est décrite par quelques variables latentes communes à de nombreux neurones. & \cite{Gallego2017} & mesuré \\
4 & Le bruit des voisins est corrélé, et une centaine de neurones porte presque autant qu'un millier (proposition~\ref{prop:pooling}) & Cortex visuel du singe: coefficient de corrélation du bruit entre neurones voisins d'environ $0.1$, et le gain de la moyenne sature vers une centaine de neurones. Théorie des corrélations qui limitent l'information. & \cite{Zohary1994, MorenoBote2014} & mesuré \\
5 & Flux brut de $2\cdot10^8$~bit/s contre $8\cdot10^4$~bit/s en impulsions~\eqref{eq:probedata} & Arithmétique du chapitre à partir de la capacité du flux d'impulsions~\eqref{eq:spikeentropy}. Les paramètres de numérisation sont réels: la puce de Neuralink échantillonne à $19.3$~kHz sur $10$~bits par canal. & démonstration dans le chapitre; \cite{Rieke1997, Musk2019} & théorie \\
6 & Premier système Neuralink: les puces amplifient et numérisent, le flux brut part par câble, un programme externe détecte les impulsions; détection sur puce, boîtier fermé, radio et alimentation sans fil: pour l'appareil destiné à l'humain, selon la société & Article de la société: tout le flux brut passe par un câble USB-C, et un programme détecte les impulsions en temps réel hors de la puce; compression embarquée, alimentation et transmission sans fil sont présentées comme un projet pour les appareils cliniques. Pour l'appareil implanté chez l'humain, seulement des communications de la société. Que la détection des impulsions sur puce soit nécessaire est une conclusion du chapitre tirée de~\eqref{eq:probedata}. & \cite{Musk2019}; rapports de Neuralink, pas d'article à comité de lecture & mesuré (article de la société); chez l'humain: rapport de la société \\
7 & Jusqu'à $3072$ électrodes sur $96$ fils de $32$; un robot insère les fils en évitant les vaisseaux & Article de la société: $3072$ électrodes sur $96$ fils de $32$; le robot insère $6$ fils ($192$ électrodes) par minute; chez le rat porteur d'une grille implantée à long terme, jusqu'à $70\,\%$ des électrodes entendent des impulsions. & \cite{Musk2019} & mesuré (article de la société) \\
8 & Chez l'humain, environ mille canaux sur $64$ fils; une partie des fils s'est déplacée; curseur de l'ordre de $10$~bit/s & Communications de la société seulement. Une recherche dans PubMed n'a trouvé aucun article à comité de lecture présentant des résultats chez l'humain, ce qui concorde avec la réserve faite dans le texte du chapitre. & rapports de Neuralink; pas d'article à comité de lecture & mesuré (rapport de la société) \\
9 & Une interface sur une grille d'une centaine d'électrodes pilote un curseur & Personne paralysée, $96$ électrodes dans le cortex moteur primaire: l'intention de bouger le bras modifie les impulsions trois ans après la lésion; le décodeur fournit un \og curseur neuronal\fg{} (courrier, télévision), la commande d'une prothèse de main et d'un bras robotisé. & \cite{Hochberg2006} & mesuré \\
10 & Écriture \og par la main mentale\fg{}: environ $90$ caractères par minute, moins de $8$~bit/s & Personne à la main paralysée, tentatives d'écriture manuscrite, décodeur récurrent: $90$ caractères par minute avec une précision de $94.1\,\%$ en temps réel, plus de $99\,\%$ après correction automatique. L'estimation en bits est l'arithmétique du chapitre. & \cite{Willett2021} & mesuré \\
11 & Les tentatives de parole sont converties en texte à environ $60$ mots par minute & Personne ayant perdu une parole intelligible, grilles dans le cortex: $62$ mots par minute; $9.1\,\%$ d'erreurs sur les mots avec un vocabulaire de $50$ mots, $23.8\,\%$ avec un vocabulaire de $125\,000$. & \cite{Willett2023} & mesuré \\
12 & L'interface n'extrait qu'une petite part de la capacité: quelques dizaines de bits par seconde pour un neurone, quelques milliers pour une centaine (proposition~\ref{prop:probe}) & La borne supérieure~\eqref{eq:spikeentropy} est théorique; la comparaison avec les lignes~8 et~10 est une conclusion du chapitre. Que le goulot d'étranglement soit la faible dimension et la méconnaissance du code, et non le fil, n'a pas été vérifié par une expérience directe. & \cite{Rieke1997} & théorie \\
13 & Le cerveau et le décodeur apprennent l'un de l'autre & Singes pilotant un curseur; le décodeur s'adapte en boucle fermée pendant que les neurones réapprennent: la précision augmente, l'habileté se conserve et souffre moins des mouvements du bras propre. Le conseil de séparer les vitesses d'apprentissage est une règle d'ingénieur, sans expérience propre dans le chapitre. & \cite{Orsborn2014coad} & mesuré \\
14 & Le courant envoyé par une électrode excite les neurones et les fils environnants sans distinction de type & Cortex de souris et de chat, imagerie calcique à deux photons: même un faible courant excite des neurones épars, dispersés jusqu'à des millimètres, par excitation directe des axones dans un volume de quelques dizaines de micromètres de diamètre. L'excitation est donc non sélective, mais pas continue. & \cite{Histed2009stim} & mesuré \\
15 & La stimulation du cortex visuel allume des points; jusqu'à $16$ points à la fois permettent de reconnaître certaines lettres et les contours des objets & Personne aveugle, $96$ électrodes dans le cortex visuel pendant $6$ mois: seuil d'une électrode $66.8\pm36.5$~µA; la stimulation simultanée de groupes (jusqu'à $16$ électrodes, limite du stimulateur) abaisse le seuil et produit des motifs distincts; certaines lettres et les contours des objets sont reconnus. & \cite{Fernandez2021} & mesuré \\
16 & Le second axe de Neuralink est un appareil qui envoie des signaux dans le cortex visuel & Seulement des communications de la société sur son développement; aucune donnée à comité de lecture sur son fonctionnement n'a été trouvée. & rapports de Neuralink & hypothèse \\
17 & On peut mesurer les concentrations des neurotransmetteurs de diffusion par un capteur chimique dans le tissu & Tranches de cerveau de rat, voltampérométrie cyclique rapide sur fibre de carbone: libération et recapture de la sérotonine mesurées; la substance sort de la synapse. & \cite{Bunin1998} & mesuré \\
\bottomrule
\end{longtable}}

\section{Chapitre~\ref{sec:eetypes}. Les neurones comme composants}

Les modes de fonctionnement des cellules individuelles sont mesurés directement, par un courant injecté par électrode et un enregistrement de la tension; les mathématiques de l'intégrateur et de la division en classes relèvent de la théorie classique, tandis que l'usage par le cerveau de la reconfiguration des modes pour calculer reste une hypothèse.

{\footnotesize
\setlength{\tabcolsep}{3pt}\renewcommand{\arraystretch}{1.15}
\begin{longtable}{>{\raggedright}p{0.5cm}>{\raggedright}p{4.0cm}>{\raggedright}p{5.6cm}>{\raggedright}p{2.3cm}>{\raggedright\arraybackslash}p{1.7cm}}
\toprule
N\textsuperscript{o} & Affirmation & Expérience et ce qui est mesuré & Source & Statut \\
\midrule\endfirsthead
\toprule
N\textsuperscript{o} & Affirmation & Expérience et ce qui est mesuré & Source & Statut \\
\midrule\endhead
1 & Résonateur: un courant lent qui s'oppose aux variations de tension fait du neurone un filtre passe-bande dont le maximum se situe entre quelques hertz et quelques dizaines de hertz (section~\ref{sec:resonator}) & En tranches, on injectait dans des neurones \nt{GLUT} un courant sinusoïdal de fréquence croissante et on mesurait l'impédance $|Z(f)|$: maximum à $2$--$5$~Hz à $33^\circ$C (environ $7$~Hz à $38^\circ$C); la résonance est créée par des courants lents potassique et cationique, et amplifiée par le courant sodique persistant. Une revue décrit la même méthode pour de nombreuses classes de cellules & \cite{HuVervaekeStorm2002,HutcheonYarom2000} & mesuré \\
2 & Le courant lent se comporte comme une inductance; avec la capacité de la membrane, il forme un circuit oscillant & La réponse sous le seuil d'un axone à un courant a été mesurée et décrite par des équations de conductance linéarisées; le schéma équivalent contient une inductance \og phénoménologique\fg{} dont la valeur dépend de la tension & \cite{Mauro1970} & théorie \\
3 & Constante de temps d'un groupe à rétroaction $\tau_{\text{eff}}=\tau/(1-w)$~\eqref{eq:perfectint}; pour $\tau=0.1$~s et $w=0.995$, $20$~s & Déduite dans le chapitre de l'équation linéaire du groupe; proposée comme mécanisme de maintien de la position de l'œil dans un travail théorique & démonstration dans le chapitre; \cite{Seung1996} & théorie \\
4 & L'intégrateur de position de l'œil maintient son entrée grâce au réseau, et non grâce à une propriété d'une cellule isolée & Enregistrement intracellulaire chez le poisson, dans le noyau qui maintient la position de l'œil: après une rotation rapide, la fréquence du neurone change par palier et se maintient; un bref courant injecté dans une seule cellule ne provoque qu'un décalage passager, tandis que les paliers de tension persistent même sous le seuil: ce sont les entrées synaptiques qui les entretiennent & \cite{Aksay2001} & mesuré \\
5 & L'intégrateur sans fuite a besoin d'un réglage permanent par apprentissage; déréglé, il oublie ou part en saturation & On entraînait des poissons avec une rétroaction visuelle proportionnelle à $\pm$ la position de l'œil: le maintien devenait instable ou fuyant, et la constante de temps des neurones chutait de plus d'un ordre de grandeur, jusqu'à $<1$~s; la rétroaction visuelle normale rétablissait peu à peu la stabilité & \cite{Major2004} & mesuré \\
6 & Neurone bistable: une brève entrée enclenche un plateau durable, l'inhibition l'arrête; la propriété est activée par la sérotonine (section~\ref{sec:bistable}) & Enregistrement intracellulaire chez le chat de neurones \nt{ACET} qui commandent les muscles: une brève excitation synaptique fait passer le neurone en activité durable, une brève inhibition le remet à zéro; chez le chat spinal, l'état apparaît après administration d'un précurseur de la sérotonine & \cite{Hounsgaard1988} & mesuré \\
7 & Deux classes: les intégrateurs (fréquence croissant à partir de zéro) et les résonateurs (fréquence finie d'emblée au seuil) & Classes déduites de la théorie des bifurcations. Expérience: en tranches de cortex, les neurones \nt{GLUT} se mettent à décharger à une fréquence aussi basse qu'on veut (type~1), les neurones \nt{GABA} rapides démarrent d'un bond à fréquence finie (type~2) & \cite{Izhikevich2007,Tateno2004} & mesuré \\
8 & Un même neurone est intégrateur ou détecteur de coïncidences: l'inhibition raccourcit la fenêtre de sommation & En tranches d'hippocampe, la fenêtre dans laquelle les entrées excitatrices s'additionnent jusqu'au seuil est inférieure à $2$~ms; elle est raccourcie par l'inhibition directe qui suit immédiatement l'excitation; dans les dendrites, où l'inhibition est plus faible, la fenêtre est plus large & \cite{Pouille2001} & mesuré \\
9 & Ligne à retard faite d'axones (tableau~\ref{tab:eetypes}) & Chez la chouette effraie, on a mesuré les retards dans les axones qui vont aux détecteurs de coïncidences: des longueurs d'axone différentes donnent des retards différents, et les détecteurs sont accordés sur la différence des temps d'arrivée du son & \cite{Carr1990} & mesuré \\
10 & Stimulateur pendant la veille, cellule silencieuse pendant le sommeil & Les neurones \nt{SERO} déchargent de façon presque régulière le jour, ralentissent en sommeil lent et se taisent presque en sommeil paradoxal (détails au chapitre~\ref{sec:sleep}) & \cite{JacobsAzmitia1992,McGintyHarper1976} & mesuré \\
11 & Le composant est fixé par les canaux ioniques et par les canaux de diffusion qui les ajustent (proposition~\ref{prop:eetypes}) & Montré sur de petits réseaux: les substances régulatrices modifient les propriétés propres des neurones et la force des synapses, si bien qu'un même circuit produit des rythmes différents & \cite{Marder2012} & mesuré \\
12 & Le cerveau exploite la reconfiguration des modes pour calculer (proposition~\ref{prop:eetypes}) & Il n'existe aucune expérience directe où la reconfiguration du mode d'une cellule serait nécessaire pour résoudre une tâche; seulement des expériences où les régulateurs modifient la sortie du circuit & \cite{Marder2012} & hypothèse \\
\bottomrule
\end{longtable}}

\section{Chapitre~\ref{sec:dendrites}. Le nombre de dendrites}

La structure des classes et le nombre d'entrées sont mesurés par l'anatomie et par des enregistrements électriques; l'estimation~\eqref{eq:dendriteload} relève de la simple physique du condensateur, tandis que l'explication computationnelle du nombre d'entrées s'appuie sur la théorie et sur des modèles.

{\footnotesize
\setlength{\tabcolsep}{3pt}\renewcommand{\arraystretch}{1.15}
\begin{longtable}{>{\raggedright}p{0.5cm}>{\raggedright}p{4.0cm}>{\raggedright}p{5.6cm}>{\raggedright}p{2.3cm}>{\raggedright\arraybackslash}p{1.7cm}}
\toprule
N\textsuperscript{o} & Affirmation & Expérience et ce qui est mesuré & Source & Statut \\
\midrule\endfirsthead
\toprule
N\textsuperscript{o} & Affirmation & Expérience et ce qui est mesuré & Source & Statut \\
\midrule\endhead
1 & Capacité spécifique de la membrane $c_m\approx1$~µF/cm$^2$ & On prélevait la membrane du corps d'un neurone dans une pipette, on appliquait un échelon de tension et, d'après la décroissance du courant capacitif, on trouvait la capacité totale, puis on la divisait par la surface: $0.9$~µF/cm$^2$ pour trois classes de neurones & \cite{Gentet2000} & mesuré \\
2 & Temps de charge $t\approx c_mA\Delta V/I$~\eqref{eq:dendriteload}: un grand neurone a besoin de dizaines d'entrées simultanées, une seule grosse synapse suffit à un petit & Estimation par la loi de charge d'un condensateur; les nombres ($20$ et $300$~pF, $0.1$ et quelques nA) sont des ordres de grandeur & démonstration dans le chapitre & théorie \\
3 & Une seule synapse énorme délivre un courant de quelques nanoampères et amène aussitôt un petit neurone au seuil & Enregistrement simultané de la terminaison et du destinataire en tranche: courant d'une synapse de $2$--$13$~nA, chaque impulsion d'entrée provoque une réponse suprasliminaire, avec un délai d'environ $0.4$~ms à $36^\circ$C; le destinataire émet \nt{GLYC} & \cite{Borst1995,BorstSoria2012} & mesuré \\
4 & Zéro dendrite: dans les neurones sensoriels du toucher et de la douleur, l'impulsion va du capteur à la moelle épinière sans passer par le corps cellulaire & La structure (un axone qui se divise en deux près du corps) est établie par l'anatomie. Que la conduction n'exige pas l'excitabilité du corps a été montré sur un modèle validé d'un tel neurone: sans corps excitable, l'impulsion franchit la bifurcation avec la même fiabilité & \cite{AmirDevor2003} & théorie \\
5 & Une dendrite: le neurone olfactif porte un seul type de récepteur et transmet une seule ligne d'entrée & Revue d'expériences sur les gènes des récepteurs: chaque neurone olfactif exprime un seul gène de récepteur parmi une grande famille & \cite{Mombaerts2004} & mesuré \\
6 & Deux dendrites: dans les détecteurs de direction du son, les entrées des oreilles gauche et droite se posent sur des dendrites différentes & Anatomie et enregistrements chez les mammifères, réunis dans une revue: les entrées des deux oreilles sont réparties sur deux dendrites opposées & \cite{Grothe2010} & mesuré \\
7 & Répartir les entrées sur deux dendrites rend le \og ET\fg{} plus strict & Montré sur un modèle: la saturation~\eqref{eq:shunt} sur une seule dendrite empêche les entrées d'un même côté d'atteindre le seuil, et la détection des coïncidences s'améliore. Il n'existe pas d'expérience directe où l'on aurait modifié la disposition des entrées & \cite{AgmonSnir1998} & théorie \\
8 & Environ $7\cdot10^{10}$ petits neurones \nt{GLUT} dans le circuit d'apprentissage moteur, avec $\sim4$ entrées chacun & Comptage des cellules par fractionnement isotrope: environ $69$~milliards de neurones dans cette partie du cerveau humain. Enregistrement d'une telle cellule chez le rat vivant: un contact déclenche des bouffées de courants discrets venant de quelques entrées, et la cellule décharge quand ils s'additionnent & \cite{Azevedo2009,Chadderton2004} & mesuré \\
9 & Un élément à seuil à $n$ entrées sépare au plus $P_{\max}\approx2n$ motifs aléatoires~\eqref{eq:cover}; pour le neurone de sortie du circuit d'apprentissage moteur, la borne est $\sim2\cdot10^5$, et l'estimation réaliste $\sim4\cdot10^4$ associations & La borne est un résultat classique de géométrie combinatoire. L'estimation réaliste vient d'une théorie de la capacité avec des poids uniquement positifs et une marge contre le bruit, appliquée au nombre mesuré d'entrées de ce neurone & \cite{Cover1965,Brunel2004} & théorie \\
10 & Les mélangeurs à peu d'entrées servent à déployer l'entrée; \og quatre\fg{} est proche de l'optimum & Théorie: la dimension de la représentation que fournit une couche de mélangeurs est maximale à peu près pour le nombre d'entrées observé en anatomie; l'hypothèse initiale de l'expansion figure dans des travaux classiques & \cite{Marr1969,Albus1971,LitwinKumar2017} & hypothèse \\
11 & Grands arbres: $10^4$--$10^5$ entrées & Comptage des synapses sur des images de microscopie électronique: $\approx1.7\cdot10^5$ synapses sur un seul neurone \nt{GABA} de sortie du circuit d'apprentissage moteur chez le rat; environ $30\,000$ entrées excitatrices et $1\,700$ inhibitrices sur un neurone \nt{GLUT} de l'hippocampe & \cite{NapperHarvey1988,Megias2001} & mesuré \\
12 & Les branches d'un grand arbre sont des sous-circuits distincts dotés de leurs propres seuils; le neurone est un petit réseau multicouche & Stimulation ponctuelle de deux sites sur de fines dendrites: les entrées d'une même branche s'additionnent selon une sigmoïde, celles de branches différentes linéairement. Dans les dendrites de neurones du cortex humain, on a trouvé des impulsions calciques graduées qui permettent à une seule cellule de résoudre un problème non linéairement séparable. Modèle: pour reproduire un seul neurone, il faut un réseau profond & \cite{Polsky2004,Gidon2020,Beniaguev2021} & mesuré \\
13 & Dendrites-sorties: chaque branche d'un neurone \nt{GABA} de la rétine est un détecteur de direction distinct & Imagerie calcique à deux photons dans des branches isolées: le signal d'une branche est sélectif au mouvement allant du corps cellulaire vers la pointe, alors que la tension du corps ne l'est pas & \cite{Euler2002} & mesuré \\
14 & Le nombre d'entrées suit la tâche (proposition~\ref{prop:dendrites}) & La structure des classes est mesurée (lignes 3--13); que le nombre d'entrées soit choisi pour la vitesse ou pour la classification n'a été montré que par la théorie et des modèles, pour certaines classes & \cite{LitwinKumar2017,AgmonSnir1998} & hypothèse \\
\bottomrule
\end{longtable}}

\section{Chapitre~\ref{sec:sleep}. Le sommeil}

Les modes de sommeil, les répétitions et la diminution des synapses sont mesurés par des enregistrements de neurones, par microdialyse et par microscopie électronique; l'horaire du sommeil et la bascule lente sont décrits par les modèles du livre, tandis que la fonction du sommeil relève surtout d'hypothèses.

{\footnotesize
\setlength{\tabcolsep}{3pt}\renewcommand{\arraystretch}{1.15}
\begin{longtable}{>{\raggedright}p{0.5cm}>{\raggedright}p{4.0cm}>{\raggedright}p{5.6cm}>{\raggedright}p{2.3cm}>{\raggedright\arraybackslash}p{1.7cm}}
\toprule
N\textsuperscript{o} & Affirmation & Expérience et ce qui est mesuré & Source & Statut \\
\midrule\endfirsthead
\toprule
N\textsuperscript{o} & Affirmation & Expérience et ce qui est mesuré & Source & Statut \\
\midrule\endhead
1 & Pendant le sommeil, les neurones du cortex ne se taisent pas; c'est le mode de fonctionnement qui change & Longs enregistrements de neurones du cortex chez le rat: en sommeil lent, la fréquence reste non nulle, et des périodes d'activité alternent avec des périodes de silence; après une longue veille, la fréquence est plus élevée dans tous les états, après le sommeil elle est plus basse & \cite{Vyazovskiy2009} & mesuré \\
2 & \nt{ACET} est élevé le jour et en sommeil paradoxal, bas en sommeil lent (tableau~\ref{tab:sleepmodes}) & Microdialyse dans le cortex chez des chats se déplaçant librement: pendant la veille, la libération d'acétylcholine est supérieure de $100$--$175\%$ à celle du sommeil lent, et en sommeil paradoxal elle est à peu près au niveau de la veille calme & \cite{Marrosu1995} & mesuré \\
3 & \nt{NORA} et \nt{SERO} s'apaisent en sommeil lent et se taisent presque en sommeil paradoxal & Enregistrement de neurones \nt{NORA} isolés chez le rat et de neurones \nt{SERO} chez le chat selon les stades du sommeil: la fréquence est maximale pendant la veille, plus basse en sommeil lent et presque nulle en sommeil paradoxal & \cite{AstonJonesBloom1981,McGintyHarper1976} & mesuré \\
4 & En sommeil paradoxal, les muscles sont coupés par une inhibition via \nt{GABA} et \nt{GLYC} & Chez le rat, on mesurait l'activité des neurones \nt{ACET} du muscle masticateur et on leur appliquait directement des bloqueurs: la paralysie n'est levée que si l'on bloque à la fois les deux types de récepteurs \nt{GABA} et les récepteurs \nt{GLYC} & \cite{BrooksPeever2012} & mesuré \\
5 & La pression de sommeil $S$ est un condensateur à deux seuils~\eqref{eq:sleeppressure}, $\tau_r=18.2$~h, $\tau_d=4.2$~h & Modèle à deux processus; les constantes de temps sont tirées de la façon dont la puissance des ondes lentes de l'EEG croît avec la veille et décroît pendant le sommeil, et la forme des seuils de l'heure du réveil spontané & \cite{Borbely1982,Daan1984} & théorie \\
6 & La machine trouve d'elle-même $16.1$~h de veille et $8.0$~h de sommeil (fig.~\ref{fig:twoprocess}) & Calcul selon~\eqref{eq:sleeppressure} avec des seuils oscillants, script \texttt{models/sleep\_models.py}: $16.1$~h de veille et $7.9$--$8.0$~h de sommeil & --- & modèle \\
7 & Après $40$~h sans sommeil, $S=0.90$, mais le sommeil ne dure que $8.8$~h: la dette se rembourse par la profondeur, et non par la durée & Modèle: \texttt{models/sleep\_models.py} donne $S=0.90$ et $8.8$~h. Expérience: après $40.5$~h sans sommeil, chez huit personnes, la puissance des ondes lentes de l'EEG est significativement plus élevée que d'habitude lors de la première nuit de récupération & \cite{Borbely1981} & modèle \\
8 & Physiquement, $S$ est l'adénosine & Microdialyse chez le chat: l'adénosine extracellulaire, dans l'une des régions qui commandent la veille, augmente pendant la veille, augmente encore en cas de privation de sommeil et baisse pendant le sommeil. Que $S$ soit uniquement l'adénosine n'est pas démontré & \cite{PorkkaHeiskanen1997,Dunwiddie2001} & hypothèse \\
9 & En sommeil lent, le cortex bascule: quelques centaines de ms d'activité, puis le silence, moins d'une fois par seconde ($<1$~Hz, le plus souvent $0.3$--$0.4$~Hz sous anesthésie) & Enregistrement intracellulaire chez le chat anesthésié: dépolarisation de $0.8$--$1.5$~s et longue hyperpolarisation, rythme $<1$~Hz, le plus souvent $0.3$--$0.4$~Hz; la même alternance d'activité et de silence s'observe dans le sommeil naturel (ligne~1) & \cite{Steriade1993,Vyazovskiy2009} & mesuré \\
10 & La bascule est un groupe à rétroaction et fatigue~\eqref{eq:updown}; pour $b=5$, période de $1.1$~s (valeur du modèle) et activité environ $35\%$ du temps; pour $b=1$, fonctionnement régulier (fig.~\ref{fig:updown}) & Calcul, script \texttt{models/sleep\_models.py}: période $1.11$~s, fraction du temps d'activité $0.35$ (moyenne sur un nombre entier de périodes); pour $b=1$, fraction d'activité $1.0$. La période du modèle est plus courte que la période mesurée (ligne~9) & --- & modèle \\
11 & \nt{ACET} arrête la bascule et fait passer le cortex à un fonctionnement régulier & La stimulation d'un noyau de neurones \nt{ACET} chez le chat bloquait le rythme lent dans $79\%$ des neurones du cortex et les faisait passer à un fonctionnement régulier, surtout par disparition des longues hyperpolarisations. L'acétylcholine supprime les courants potassiques lents qui créent la fatigue & \cite{SteriadeAmzica1993,McCormick1992} & mesuré \\
12 & Les bouffées de rythme à $7$--$15$~Hz naissent d'une boucle \nt{GLUT}--\nt{GABA} entre le cortex et un nœud relais & Dans des tranches du nœud relais visuel du furet, les bouffées de rythme sont créées par les bouffées de réponse des cellules relais à l'inhibition venant d'un noyau \nt{GABA} voisin, qu'elles excitent elles-mêmes & \cite{vonKrosigk1993,SteriadeMcCormickSejnowski1993} & mesuré \\
13 & Les répétitions de la journée se font en accéléré: environ $20$ fois plus vite dans l'hippocampe, $6$--$7$ fois dans le cortex & En sommeil lent, les séquences répétées de neurones de l'hippocampe sont comprimées dans le temps. Après une course sur une piste, les séquences de cellules de lieu se répétaient dans le même ordre par brèves bouffées de $\sim100$~ms, comprimées environ $20$ fois; dans le cortex, compression de $6$--$7$ fois & \cite{Nadasdy1999,LeeWilson2002,Euston2007} & mesuré \\
14 & Renforcer la coordination des répétitions avec le cortex améliore la mémoire (lien causal) & Chez le rat, après l'apprentissage, une stimulation calée sur les répétitions renforçait leur couplage avec les ondes lentes et les bouffées de rythme du cortex: le lendemain, la mémoire était élevée, alors que chez les témoins elle restait au niveau du hasard & \cite{Maingret2016} & mesuré \\
15 & La fonction des répétitions est l'apprentissage entrelacé de la mémoire lente & Théorie des deux systèmes d'apprentissage et calculs sur des réseaux artificiels: l'apprentissage séquentiel abîme l'ancien, l'apprentissage entrelacé non. Il n'existe pas d'expérience directe sur le cerveau montrant que les répétitions servent précisément à cela & \cite{McClelland1995} & hypothèse \\
16 & Après le sommeil, les synapses diminuent proportionnellement à leur taille; les grosses changent à peine & Microscopie électronique tridimensionnelle de $6920$ synapses du cortex de souris: surface de contact inférieure de $\sim18\%$ après le sommeil par rapport à la veille, diminution proportionnelle à la taille, les grosses synapses stables ne changent pas & \cite{deVivo2017} & mesuré \\
17 & La tâche principale du sommeil est la renormalisation des poids & Hypothèse de l'homéostasie synaptique; les lignes 1 et 16 s'accordent avec elle, mais ne prouvent pas qu'il s'agit de la fonction principale & \cite{TononiCirelli2014} & hypothèse \\
18 & Le sommeil paradoxal est un essai sur banc du modèle du monde & Le mode (tableau~\ref{tab:sleepmodes}) est mesuré; que le réseau fasse tourner un modèle du monde est une supposition théorique, sans expérience & \cite{Hobson2009} & hypothèse \\
19 & Lavage du cerveau pendant le sommeil: question ouverte & Une expérience: pendant le sommeil, l'espace intercellulaire est $60\%$ plus large et le nettoyage plus rapide. Une autre, avec une autre méthode de mesure: pendant le sommeil et sous anesthésie, le nettoyage est nettement plus lent. Les résultats se contredisent & \cite{Xie2013,Miao2024} & hypothèse \\
20 & Sommeil local: des zones du cortex d'un animal éveillé sombrent brièvement dans le silence & Chez le rat, après une longue veille, les neurones d'une zone du cortex se taisent brièvement, comme en sommeil lent, avec une onde lente dans l'EEG local; à ces moments-là, le rat se trompe plus souvent dans la tâche & \cite{Vyazovskiy2011} & mesuré \\
\bottomrule
\end{longtable}}

\section{Chapitre~\ref{sec:livingmachine}. Une machine faite de neurones vivants}

Le comportement des cultures sur matrices d'électrodes est mesuré dans des expériences directes d'enregistrement et de stimulation; le mécanisme des salves s'appuie sur le modèle d'épuisement synaptique et sur des expériences en microcultures, et les affirmations sur la dopamine en boîte de culture sur des expériences isolées, dont certaines ne sont, de l'aveu même des auteurs, qu'une démonstration.

{\footnotesize
\setlength{\tabcolsep}{3pt}\renewcommand{\arraystretch}{1.15}
\begin{longtable}{>{\raggedright}p{0.5cm}>{\raggedright}p{4.0cm}>{\raggedright}p{5.6cm}>{\raggedright}p{2.3cm}>{\raggedright\arraybackslash}p{1.7cm}}
\toprule
N\textsuperscript{o} & Affirmation & Expérience et ce qui est mesuré & Source & Statut \\
\midrule\endfirsthead
\toprule
N\textsuperscript{o} & Affirmation & Expérience et ce qui est mesuré & Source & Statut \\
\midrule\endhead
1 & Culture sur puce: surtout des \nt{GLUT}, une part plus faible de \nt{GABA} qui dépend de la préparation; environ $6\%$ dans les cultures d'hippocampe & Un réseau de milliers de neurones sur matrice d'électrodes est une expérience standard (ligne~3). La part des neurones \nt{GABA} dans les cultures d'hippocampe a été comptée par marquage du GABA: environ $6\%$ des neurones. Pour les cultures de cortex, nous ne donnons pas de chiffre vérifié & \cite{Benson1994} & mesuré \\
2 & Organoïdes: tissu nerveux tridimensionnel d'environ un millimètre, issu de cellules souches humaines & À partir de cellules souches, on a cultivé des organoïdes tridimensionnels comportant plusieurs régions cérébrales, dont un cortex avec des couches de progéniteurs et des neurones matures & \cite{Lancaster2013} & mesuré \\
3 & Sans entrée, la culture s'embrase en salves espacées d'une seconde à plusieurs minutes, selon la culture & $58$ cultures de cortex, d'une densité de $3\,000$ à $50\,000$ neurones, ont été enregistrées pendant cinq semaines ($963$ enregistrements d'une demi-heure): après deux semaines, les salves de tout le réseau dominent dans presque toutes les cultures; les intervalles entre salves vont de $1$ à $300$~s et dépendent fortement de la préparation & \cite{Wagenaar2006} & mesuré \\
4 & Une salve s'achève par épuisement du neurotransmetteur, avec un intervalle $T_{\text{salve}}\approx\tau_{\text{rec}}\ln\frac{1}{1-x^*}$~\eqref{eq:burstinterval}; $2$--$4$~s est l'estimation du modèle avec $\tau_{\text{rec}}=0.8$~s, la récupération mesurée est plus lente & La formule est établie dans le chapitre. Modèle: un réseau à synapses dépressives produit spontanément des salves de tout le réseau. Expérience sur des microcultures de $4$--$30$ neurones: la durée d'une salve dépend du temps écoulé depuis la précédente, et l'épuisement des vésicules de neurotransmetteur supprime les salves; conclusion des auteurs: la salve est limitée par la réserve de neurotransmetteur. La réserve s'y reconstitue en quelques dizaines de secondes, ce qui, avec la même formule, donne des intervalles de dizaines de secondes à plusieurs minutes & démonstration dans le chapitre; \cite{Tsodyks2000,CohenSegal2011,Tsodyks1997} & théorie \\
5 & \og Le professeur qui se tait\fg{}: la stimulation est coupée quand la réponse voulue apparaît, et la réponse se fixe & La culture était stimulée à $0.3$--$1$~Hz jusqu'à ce que la réponse choisie apparaisse $50\pm10$~ms après le stimulus, puis la stimulation était coupée pendant $5$~min; après quelques cycles, la réponse voulue était évoquée d'emblée; des courbes d'apprentissage ont été tracées & \cite{Shahaf2001} & mesuré \\
6 & La culture joue au tennis; hausse significative des séries en une session seulement avec rétroaction complète & Détails au chapitre~\ref{sec:dishbrain} et dans son annexe: avec du silence à la place du stimulus, la culture joue aussi mieux en fin de session que sans rétroaction, mais avec un effet moindre; l'apprentissage d'une session à l'autre est instable & \cite{Kagan2022} & mesuré \\
7 & Que la culture apprenne à prédire son entrée reste une hypothèse; les grands mots sur sa \og sensibilité\fg{} sont contestés & Le principe de l'énergie libre est une proposition théorique; aucune expérience directe ne le distingue d'explications plus simples. Dans un article de réponse, trente chercheurs ont souligné qu'un changement de comportement dans la boucle ne démontre ni intelligence ni sensations & \cite{Friston2010,Balci2023} & hypothèse \\
8 & La culture dirige un corps virtuel vers une cible avec un motif de stimuli bien choisi & Un programme choisissait lui-même les motifs de stimulation d'entraînement selon le résultat courant; en quelques dizaines de minutes, le réseau atteignait les états visés, et la plasticité durait $80$~min après $2$~h d'entraînement. Les mêmes stimuli, appliqués sans lien avec le comportement, restaient sans effet & \cite{Bakkum2008} & mesuré \\
9 & Réservoir: seule la lecture linéaire $y=W_{\text{sort}}\mathbf r$~\eqref{eq:reservoir} est entraînée & Théorie du calcul par réservoir: tout système non linéaire à mémoire évanescente, plus une sortie linéaire entraînée & \cite{Maass2002,JaegerHaas2004} & théorie \\
10 & Un organoïde utilisé comme réservoir distingue les locuteurs avec une précision d'environ $78\%$ (selon les auteurs); la lecture apprend par l'erreur, et les connexions de l'organoïde changent sans professeur & Des sons de parole de plusieurs locuteurs étaient injectés à l'organoïde sous forme de motifs de courant sur une matrice d'électrodes; la lecture entraînée distinguait le locuteur avec une précision d'environ $78\%$, selon les auteurs. Ceux-ci rapportent aussi que la connectivité fonctionnelle de l'organoïde changeait pendant l'entraînement, et parlent d'apprentissage non supervisé dans l'organoïde lui-même & \cite{Cai2023} & mesuré \\
11 & Un million de neurones dépense environ $10^{-4}$~W en impulsions~\eqref{eq:dishpower} & Estimation: le coût d'une impulsion, $10^{-10}$~J, tiré du budget énergétique du cortex~\eqref{eq:energy}, multiplié par le nombre d'impulsions; il n'existe pas de mesure directe de la puissance d'une culture & démonstration dans le chapitre; \cite{Attwell2001} & théorie \\
12 & Dopamine par éclair de lumière: $1$~mM de dopamine en cage, $240$ répétitions, le nombre d'impulsions évoquées augmentait & Sur une plateforme d'organoïdes à distance, la dopamine enfermée dans une cage photosensible ($1$~mM) était libérée par ultraviolet ($365$~nm, $0.8$~s) quand le stimulus provoquait plus d'impulsions; en $240$ étapes (environ $5$~h), le nombre d'impulsions a globalement augmenté. Les auteurs écrivent qu'une seule expérience ne permet pas d'établir le lien avec la dopamine; il n'y a pas de contrôles & \cite{Jordan2024} & hypothèse \\
13 & La dopamine de cellules vivantes modifie le poids d'un transistor organique de façon durable et réversible & Des cellules sécrétant de la dopamine ont été cultivées au-dessus d'un composant neuromorphique organique; l'oxydation de la dopamine sur l'électrode modifie la conductance du canal; un changement durable du poids et son retour ont été démontrés & \cite{Keene2020} & mesuré \\
14 & Des organoïdes contenant des neurones \nt{DOPA} fonctionnels existent; leur dopamine n'a pas été utilisée comme professeur & Dans des organoïdes issus de cellules souches humaines, on a trouvé des neurones \nt{DOPA} matures électriquement actifs et une production de dopamine. Nous n'avons trouvé dans la littérature aucune expérience où cette dopamine enseigne à un autre réseau & \cite{Jo2016} & mesuré \\
15 & Un organoïde tient un pendule en équilibre: les stimuli d'apprentissage choisis par le programme font mieux que le hasard, sans consolidation, via les synapses \nt{GLUT} & Organoïdes de cortex de souris en boucle fermée sur la tâche du \og pendule inversé sur chariot\fg{}: pour la plupart des organoïdes, les signaux choisis par apprentissage par renforcement ont donné un meilleur résultat que des signaux aléatoires ou l'absence de signal; après $45$~min de repos, l'amélioration ne subsistait pas; le blocage des récepteurs AMPA et NMDA la supprimait & \cite{Robbins2026} & mesuré \\
16 & Sans registres de contrôle, le réseau du banc ne peut pas décider lui-même de ce qui compte comme un succès (proposition~\ref{prop:livingmachine}) & Dans toutes les expériences des lignes 5--15, le signal d'apprentissage est fixé par l'expérimentateur ou par un programme; aucune expérience ne montre une culture qui élaborerait elle-même un critère de succès. C'est une conclusion tirée d'une absence, et non une expérience & --- & hypothèse \\
\bottomrule
\end{longtable}}

\section{Chapitre~\ref{sec:dishbrain}. DishBrain}

La conception du banc et les résultats du jeu proviennent d'une seule étude expérimentale et de sa suite; les formules du bruit et de la dérive vers les bons réglages sont établies dans le chapitre et vérifiées sur le modèle du livre, tandis que le mécanisme d'apprentissage dans la boîte n'est pas établi.

{\footnotesize
\setlength{\tabcolsep}{3pt}\renewcommand{\arraystretch}{1.15}
\begin{longtable}{>{\raggedright}p{0.5cm}>{\raggedright}p{4.0cm}>{\raggedright}p{5.6cm}>{\raggedright}p{2.3cm}>{\raggedright\arraybackslash}p{1.7cm}}
\toprule
N\textsuperscript{o} & Affirmation & Expérience et ce qui est mesuré & Source & Statut \\
\midrule\endfirsthead
\toprule
N\textsuperscript{o} & Affirmation & Expérience et ce qui est mesuré & Source & Statut \\
\midrule\endhead
1 & Banc: environ $800\,000$ cellules de cortex de souris ou environ $10^6$ cellules humaines par puce; $26\,000$ électrodes sur $8$~mm$^2$, enregistrement jusqu'à $1024$, stimulation par $8$ & Méthodes de l'étude: puce MaxOne, $26\,000$ électrodes de platine sur $8$~mm$^2$, enregistrement de $1024$ canaux à $20\,000$ échantillons par seconde; techniquement, on peut stimuler jusqu'à $32$ électrodes, $8$ ont pu l'être indépendamment; les cultures de souris ont été utilisées à partir d'environ $10$~jours & \cite{Kagan2022} & mesuré \\
2 & Boucle cadencée à $10$~ms: les impulsions des zones motrices déplacent la raquette (fig.~\ref{fig:dishloop}) & Les impulsions sont comptées par fenêtres de $10$~ms; le délai entre une impulsion et le stimulus de réponse, d'environ $5$~ms, est fixé par le tampon du matériel & \cite{Kagan2022} & mesuré \\
3 & Entrée: code de place (laquelle des $8$ électrodes) et fréquence de $4$ à $40$~imp./s & Dans l'expérience principale, la fréquence de stimulation croît linéairement de $4$~Hz quand la balle est près du mur du fond à $40$~Hz près de la raquette; l'amplitude des impulsions de la balle est de $75$~mV; les impulsions sont rectangulaires biphasiques & \cite{Kagan2022} & mesuré \\
4 & Sortie $\Delta p=c\,(n_\uparrow-n_\downarrow)$~\eqref{eq:dishreadout}, bruit du compte par fenêtre $1/\sqrt{RT}$ & La règle de différence entre deux zones est décrite dans l'étude (avec des poids des zones selon leur activité), la formule exacte n'est pas donnée. L'estimation du bruit découle du comptage poissonien et est établie dans le chapitre & \cite{Kagan2022}; démonstration dans le chapitre & théorie \\
5 & Professeur: après un renvoi, stimulus prévisible de $75$~mV, $100$~imp./s, $0.1$~s; après un raté, stimulus imprévisible de $150$~mV, $5$~imp./s, $4$~s en des lieux et instants aléatoires, puis $4$~s de pause & Méthodes: après un renvoi, $75$~mV, $100$~Hz, $100$~ms sur les $8$ électrodes à la fois; après un raté, $150$~mV, $5$~Hz sur des électrodes aléatoires à des instants aléatoires pendant $4$~s, puis $4$~s sans stimulation. La culture ne contient pas de dopamine & \cite{Kagan2022} & mesuré \\
6 & Le réseau agit de façon à rendre son entrée prévisible & Le principe de l'énergie libre est une proposition théorique dont les auteurs ont tiré le protocole; l'expérience est compatible avec lui, mais ne le distingue pas de mécanismes plus simples (lignes~7--8) & \cite{Friston2010,Kagan2022} & hypothèse \\
7 & Si l'échec secoue le réseau et non le succès, le temps passé dans un réglage vaut $\rho\propto1/P_{\text{raté}}$~\eqref{eq:dishdensity} (proposition~\ref{prop:dishscramble}) & Découle de l'équilibre des flux dans une chaîne de Markov à secousses symétriques, établi dans le chapitre. Aucun test direct sur une culture & démonstration dans le chapitre & théorie \\
8 & Le modèle du \og brassage après un raté\fg{} apprend: proportion de renvois $0.21\to0.38$; avec une secousse après chaque échange $\approx0.12$, sans secousse $0.19$ (fig.~\ref{fig:dishmodel}) & Calcul, script \texttt{models/dishbrain\_model.py}, $40$ cultures modèles de $2000$ échanges chacune: $0.21\to0.38$; $0.13\to0.11$; $0.19\to0.19$ & --- & modèle \\
9 & Les séries de renvois ne s'allongent que chez les cultures vivantes en boucle complète; les contrôles n'apprennent pas & $399$ sessions de $20$~min: $9$ cultures de souris ($101$ sessions), $11$ humaines ($138$), $6$ puces avec milieu seul ($80$), $20$ cultures au repos ($42$), $3$ simulations ($38$). La longueur moyenne d'un échange sur les $15$ dernières minutes n'est supérieure à celle des $5$ premières que chez les cultures vivantes; les cellules non neuronales (HEK293T) et les puces avec milieu seul ne montrent aucun effet & \cite{Kagan2022} & mesuré \\
10 & Hausse significative en une session seulement avec rétroaction complète; avec du silence à la place du stimulus, le jeu est meilleur en fin de session que sans rétroaction, mais l'effet est moindre & $486$ sessions sur $3$ jours: \og stimulus\fg{} ($n=27$), \og silence\fg{} ($17$), \og sans rétroaction\fg{} ($15$). La hausse au fil du temps n'est significative que dans la condition \og stimulus\fg{}; en fin de session, la condition \og silence\fg{} dépassait \og sans rétroaction\fg{} avec un effet moindre & \cite{Kagan2022} & mesuré \\
11 & L'effet est faible; savoir si le réseau apprend durablement reste une question ouverte & Au sein d'une session, l'amélioration est robuste, mais, de l'aveu même des auteurs, aucun apprentissage stable d'une session à l'autre n'a été observé sur plusieurs jours & \cite{Kagan2022} & mesuré \\
12 & Pendant le jeu, le réseau se rapproche du régime critique, et plus il s'en approche, meilleur est le jeu & Analyse des avalanches d'activité des mêmes cultures: une entrée structurée rapproche le réseau du régime critique, et la qualité du jeu est corrélée à cette proximité; mais la criticité seule, sans information sur les conséquences des actions, ne suffit pas à l'apprentissage & \cite{Habibollahi2023} & mesuré \\
13 & La \og sensibilité\fg{} de la culture n'est pas démontrée & Article de réponse de trente chercheurs: un changement de comportement en boucle fermée ne prouve ni intelligence ni sensations; aucune expérience ne le démontre & \cite{Balci2023} & hypothèse \\
14 & Quatrième contrôle proposé: un stimulus prévisible après un raté, de même durée et de même énergie (section~\ref{sec:dishspec}) & Cette expérience n'a pas été réalisée; c'est une proposition du livre & --- & hypothèse \\
\bottomrule
\end{longtable}}


\begin{thebibliography}{99}
\bibitem{Azevedo2009} F. A. C. Azevedo et al., \emph{Equal numbers of neuronal and nonneuronal cells make the human brain an isometrically scaled-up primate brain}, J. Comp. Neurol. \textbf{513}, 532 (2009).
\bibitem{Schultz1997} W. Schultz, P. Dayan, and P. R. Montague, \emph{A neural substrate of prediction and reward}, Science \textbf{275}, 1593 (1997).
\bibitem{SuttonBarto2018} R. S. Sutton and A. G. Barto, \emph{Reinforcement Learning: An Introduction}, 2nd ed. (MIT Press, Cambridge, MA, 2018).
\bibitem{Doya2002} K. Doya, \emph{Metalearning and neuromodulation}, Neural Networks \textbf{15}, 495 (2002).
\bibitem{Strata1999} P. Strata and R. Harvey, \emph{Dale's principle}, Brain Res. Bull. \textbf{50}, 349 (1999).
\bibitem{vanVreeswijk1996} C. van Vreeswijk and H. Sompolinsky, \emph{Chaos in neuronal networks with balanced excitatory and inhibitory activity}, Science \textbf{274}, 1724 (1996).
\bibitem{AstonJones2005} G. Aston-Jones and J. D. Cohen, \emph{An integrative theory of locus coeruleus-norepinephrine function: adaptive gain and optimal performance}, Annu. Rev. Neurosci. \textbf{28}, 403 (2005).
\bibitem{Beaulieu2011} J.-M. Beaulieu and R. R. Gainetdinov, \emph{The physiology, signaling, and pharmacology of dopamine receptors}, Pharmacol. Rev. \textbf{63}, 182 (2011).
\bibitem{Sparso2001} J. Spars{\o} and S. Furber (eds.), \emph{Principles of Asynchronous Circuit Design: A Systems Perspective} (Kluwer, Boston, 2001).
\bibitem{Carr1990} C. E. Carr and M. Konishi, \emph{A circuit for detection of interaural time differences in the brain stem of the barn owl}, J. Neurosci. \textbf{10}, 3227 (1990).
\bibitem{Wang2001} X.-J. Wang, \emph{Synaptic reverberation underlying mnemonic persistent activity}, Trends Neurosci. \textbf{24}, 455 (2001).
\bibitem{Hahnloser2002} R. H. R. Hahnloser, A. A. Kozhevnikov, and M. S. Fee, \emph{An ultra-sparse code underlies the generation of neural sequences in a songbird}, Nature \textbf{419}, 65 (2002).
\bibitem{Barlow1965} H. B. Barlow and W. R. Levick, \emph{The mechanism of directionally selective units in rabbit's retina}, J. Physiol. \textbf{178}, 477 (1965).
\bibitem{Carandini2012} M. Carandini and D. J. Heeger, \emph{Normalization as a canonical neural computation}, Nat. Rev. Neurosci. \textbf{13}, 51 (2012).
\bibitem{Pouille2001} F. Pouille and M. Scanziani, \emph{Enforcement of temporal fidelity in pyramidal cells by somatic feed-forward inhibition}, Science \textbf{293}, 1159 (2001).
\bibitem{Tsodyks1997isn} M. V. Tsodyks, W. E. Skaggs, T. J. Sejnowski, and B. L. McNaughton, \emph{Paradoxical effects of external modulation of inhibitory interneurons}, J. Neurosci. \textbf{17}, 4382 (1997).
\bibitem{Buzsaki2012} G. Buzs\'aki and X.-J. Wang, \emph{Mechanisms of gamma oscillations}, Annu. Rev. Neurosci. \textbf{35}, 203 (2012).
\bibitem{Rieke1997} F. Rieke, D. Warland, R. de Ruyter van Steveninck, and W. Bialek, \emph{Spikes: Exploring the Neural Code} (MIT Press, Cambridge, MA, 1997).
\bibitem{Mainen1995} Z. F. Mainen and T. J. Sejnowski, \emph{Reliability of spike timing in neocortical neurons}, Science \textbf{268}, 1503 (1995).
\bibitem{Softky1993} W. R. Softky and C. Koch, \emph{The highly irregular firing of cortical cells is inconsistent with temporal integration of random EPSPs}, J. Neurosci. \textbf{13}, 334 (1993).
\bibitem{Thorpe1996} S. Thorpe, D. Fize, and C. Marlot, \emph{Speed of processing in the human visual system}, Nature \textbf{381}, 520 (1996).
\bibitem{Zohary1994} E. Zohary, M. N. Shadlen, and W. T. Newsome, \emph{Correlated neuronal discharge rate and its implications for psychophysical performance}, Nature \textbf{370}, 140 (1994).
\bibitem{MorenoBote2014} R. Moreno-Bote et al., \emph{Information-limiting correlations}, Nat. Neurosci. \textbf{17}, 1410 (2014).
\bibitem{Gollisch2008} T. Gollisch and M. Meister, \emph{Rapid neural coding in the retina with relative spike latencies}, Science \textbf{319}, 1108 (2008).
\bibitem{OKeefe1993} J. O'Keefe and M. L. Recce, \emph{Phase relationship between hippocampal place units and the EEG theta rhythm}, Hippocampus \textbf{3}, 317 (1993).
\bibitem{Destexhe2003} A. Destexhe, M. Rudolph, and D. Par\'e, \emph{The high-conductance state of neocortical neurons in vivo}, Nat. Rev. Neurosci. \textbf{4}, 739 (2003).
\bibitem{Tsodyks1997} M. V. Tsodyks and H. Markram, \emph{The neural code between neocortical pyramidal neurons depends on neurotransmitter release probability}, Proc. Natl. Acad. Sci. USA \textbf{94}, 719 (1997).
\bibitem{Lisman1997} J. E. Lisman, \emph{Bursts as a unit of neural information: making unreliable synapses reliable}, Trends Neurosci. \textbf{20}, 38 (1997).
\bibitem{Attwell2001} D. Attwell and S. B. Laughlin, \emph{An energy budget for signaling in the grey matter of the brain}, J. Cereb. Blood Flow Metab. \textbf{21}, 1133 (2001).
\bibitem{Lennie2003} P. Lennie, \emph{The cost of cortical computation}, Curr. Biol. \textbf{13}, 493 (2003).
\bibitem{Yagishita2014} S. Yagishita et al., \emph{A critical time window for dopamine actions on the structural plasticity of dendritic spines}, Science \textbf{345}, 1616 (2014).
\bibitem{Werfel2005} J. Werfel, X. Xie, and H. S. Seung, \emph{Learning curves for stochastic gradient descent in linear feedforward networks}, Neural Comput. \textbf{17}, 2699 (2005).
\bibitem{Doya2000} K. Doya, \emph{Reinforcement learning in continuous time and space}, Neural Comput. \textbf{12}, 219 (2000).
\bibitem{Oja1982} E. Oja, \emph{Simplified neuron model as a principal component analyzer}, J. Math. Biol. \textbf{15}, 267 (1982).
\bibitem{BiPoo1998} G.-Q. Bi and M.-M. Poo, \emph{Synaptic modifications in cultured hippocampal neurons: dependence on spike timing, synaptic strength, and postsynaptic cell type}, J. Neurosci. \textbf{18}, 10464 (1998).
\bibitem{Williams1992} R. J. Williams, \emph{Simple statistical gradient-following algorithms for connectionist reinforcement learning}, Mach. Learn. \textbf{8}, 229 (1992).
\bibitem{Bayer2005} H. M. Bayer and P. W. Glimcher, \emph{Midbrain dopamine neurons encode a quantitative reward prediction error signal}, Neuron \textbf{47}, 129 (2005).
\bibitem{Shen2008} W. Shen, M. Flajolet, P. Greengard, and D. J. Surmeier, \emph{Dichotomous dopaminergic control of striatal synaptic plasticity}, Science \textbf{321}, 848 (2008).
\bibitem{Dabney2020} W. Dabney, Z. Kurth-Nelson, N. Uchida, C. K. Starkweather, D. Hassabis, R. Munos, and M. Botvinick, \emph{A distributional code for value in dopamine-based reinforcement learning}, Nature \textbf{577}, 671 (2020).
\bibitem{Matsumoto2009} M. Matsumoto and O. Hikosaka, \emph{Two types of dopamine neuron distinctly convey positive and negative motivational signals}, Nature \textbf{459}, 837 (2009).
\bibitem{BrombergMartin2010} E. S. Bromberg-Martin, M. Matsumoto, and O. Hikosaka, \emph{Dopamine in motivational control: rewarding, aversive, and alerting}, Neuron \textbf{68}, 815 (2010).
\bibitem{Menegas2018} W. Menegas, K. Akiti, R. Amo, N. Uchida, and M. Watabe-Uchida, \emph{Dopamine neurons projecting to the posterior striatum reinforce avoidance of threatening stimuli}, Nat. Neurosci. \textbf{21}, 1421 (2018).
\bibitem{Niv2007} Y. Niv, N. D. Daw, D. Joel, and P. Dayan, \emph{Tonic dopamine: opportunity costs and the control of response vigor}, Psychopharmacology \textbf{191}, 507 (2007).
\bibitem{Mohebi2019} A. Mohebi et al., \emph{Dissociable dopamine dynamics for learning and motivation}, Nature \textbf{570}, 65 (2019).
\bibitem{Hamid2016} A. A. Hamid et al., \emph{Mesolimbic dopamine signals the value of work}, Nat. Neurosci. \textbf{19}, 117 (2016).
\bibitem{Kim2020} H. R. Kim, A. N. Malik et al., \emph{A unified framework for dopamine signals across timescales}, Cell \textbf{183}, 1600 (2020).
\bibitem{Jeong2022} H. Jeong et al., \emph{Mesolimbic dopamine release conveys causal associations}, Science \textbf{378}, eabq6740 (2022).
\bibitem{Amo2022} R. Amo et al., \emph{A gradual temporal shift of dopamine responses mirrors the progression of temporal difference error in machine learning}, Nat. Neurosci. \textbf{25}, 1082 (2022).
\bibitem{Takahashi2017} Y. K. Takahashi et al., \emph{Dopamine neurons respond to errors in the prediction of sensory features of expected rewards}, Neuron \textbf{95}, 1395 (2017).
\bibitem{Sharpe2017} M. J. Sharpe et al., \emph{Dopamine transients are sufficient and necessary for acquisition of model-based associations}, Nat. Neurosci. \textbf{20}, 735 (2017).
\bibitem{Gardner2018} M. P. H. Gardner, G. Schoenbaum, and S. J. Gershman, \emph{Rethinking dopamine as generalized prediction error}, Proc. R. Soc. B \textbf{285}, 20181645 (2018).
\bibitem{Engelhard2019} B. Engelhard et al., \emph{Specialized coding of sensory, motor and cognitive variables in VTA dopamine neurons}, Nature \textbf{570}, 509 (2019).
\bibitem{Clements1992} J. D. Clements, R. A. J. Lester, G. Tong, C. E. Jahr, and G. L. Westbrook, \emph{The time course of glutamate in the synaptic cleft}, Science \textbf{258}, 1498 (1992).
\bibitem{Hasselmo2006} M. E. Hasselmo, \emph{The role of acetylcholine in learning and memory}, Curr. Opin. Neurobiol. \textbf{16}, 710 (2006).
\bibitem{JacobsAzmitia1992} B. L. Jacobs and E. C. Azmitia, \emph{Structure and function of the brain serotonin system}, Physiol. Rev. \textbf{72}, 165 (1992).
\bibitem{Schiller1992} P. H. Schiller, \emph{The ON and OFF channels of the visual system}, Trends Neurosci. \textbf{15}, 86 (1992).
\bibitem{Grothe2010} B. Grothe, M. Pecka, and D. McAlpine, \emph{Mechanisms of sound localization in mammals}, Physiol. Rev. \textbf{90}, 983 (2010).
\bibitem{Markram2004} H. Markram et al., \emph{Interneurons of the neocortical inhibitory system}, Nat. Rev. Neurosci. \textbf{5}, 793 (2004).
\bibitem{Joshi2016} S. Joshi, Y. Li, R. M. Kalwani, and J. I. Gold, \emph{Relationships between pupil diameter and neuronal activity in the locus coeruleus, colliculi, and cingulate cortex}, Neuron \textbf{89}, 221 (2016).
\bibitem{ServanSchreiber1990} D. Servan-Schreiber, H. Printz, and J. D. Cohen, \emph{A network model of catecholamine effects: gain, signal-to-noise ratio, and behavior}, Science \textbf{249}, 892 (1990).
\bibitem{YuDayan2005} A. J. Yu and P. Dayan, \emph{Uncertainty, neuromodulation, and attention}, Neuron \textbf{46}, 681 (2005).
\bibitem{Nassar2012} M. R. Nassar et al., \emph{Rational regulation of learning dynamics by pupil-linked arousal systems}, Nat. Neurosci. \textbf{15}, 1040 (2012).
\bibitem{BouretSara2005} S. Bouret and S. J. Sara, \emph{Network reset: a simplified overarching theory of locus coeruleus noradrenaline function}, Trends Neurosci. \textbf{28}, 574 (2005).
\bibitem{Hebb1949} D. O. Hebb, \emph{The Organization of Behavior: A Neuropsychological Theory} (Wiley, New York, 1949).
\bibitem{MacKay1952} D. M. MacKay and W. S. McCulloch, \emph{The limiting information capacity of a neuronal link}, Bull. Math. Biophys. \textbf{14}, 127 (1952).
\bibitem{WilsonCowan1972} H. R. Wilson and J. D. Cowan, \emph{Excitatory and inhibitory interactions in localized populations of model neurons}, Biophys. J. \textbf{12}, 1 (1972).
\bibitem{AllenStevens1994} C. Allen and C. F. Stevens, \emph{An evaluation of causes for unreliability of synaptic transmission}, Proc. Natl. Acad. Sci. USA \textbf{91}, 10380 (1994).
\bibitem{Barlow1961} H. B. Barlow, \emph{Possible principles underlying the transformations of sensory messages}, in \emph{Sensory Communication}, ed. W. A. Rosenblith (MIT Press, Cambridge, MA, 1961), pp.~217--234.
\bibitem{Cardin2009} J. A. Cardin et al., \emph{Driving fast-spiking cells induces gamma rhythm and controls sensory responses}, Nature \textbf{459}, 663 (2009).
\bibitem{Lillicrap2020} T. P. Lillicrap, A. Santoro, L. Marris, C. J. Akerman, and G. Hinton, \emph{Backpropagation and the brain}, Nat. Rev. Neurosci. \textbf{21}, 335 (2020).
\bibitem{Fremaux2016} N. Fr\'emaux and W. Gerstner, \emph{Neuromodulated spike-timing-dependent plasticity, and theory of three-factor learning rules}, Front. Neural Circuits \textbf{9}, 85 (2016).
\bibitem{Haider2006} B. Haider, A. Duque, A. R. Hasenstaub, and D. A. McCormick, \emph{Neocortical network activity in vivo is generated through a dynamic balance of excitation and inhibition}, J. Neurosci. \textbf{26}, 4535 (2006).
\bibitem{HoltKoch1997} G. R. Holt and C. Koch, \emph{Shunting inhibition does not have a divisive effect on firing rates}, Neural Comput. \textbf{9}, 1001 (1997).
\bibitem{Chance2002} F. S. Chance, L. F. Abbott, and A. D. Reyes, \emph{Gain modulation from background synaptic input}, Neuron \textbf{35}, 773 (2002).
\bibitem{Ozeki2009} H. Ozeki, I. M. Finn, E. S. Schaffer, K. D. Miller, and D. Ferster, \emph{Inhibitory stabilization of the cortical network underlies visual surround suppression}, Neuron \textbf{62}, 578 (2009).
\bibitem{HasselmoBower1992} M. E. Hasselmo and J. M. Bower, \emph{Cholinergic suppression specific to intrinsic not afferent fiber synapses in rat piriform (olfactory) cortex}, J. Neurophysiol. \textbf{67}, 1222 (1992).
\bibitem{HasselmoAndersonBower1992} M. E. Hasselmo, B. P. Anderson, and J. M. Bower, \emph{Cholinergic modulation of cortical associative memory function}, J. Neurophysiol. \textbf{67}, 1230 (1992).
\bibitem{Hasselmo1999} M. E. Hasselmo, \emph{Neuromodulation: acetylcholine and memory consolidation}, Trends Cogn. Sci. \textbf{3}, 351 (1999).
\bibitem{Tsodyks1988} M. V. Tsodyks and M. V. Feigel'man, \emph{The enhanced storage capacity in neural networks with low activity level}, Europhys. Lett. \textbf{6}, 101 (1988).
\bibitem{KalmanBucy1961} R. E. Kalman and R. S. Bucy, \emph{New results in linear filtering and prediction theory}, J. Basic Eng. \textbf{83}, 95 (1961).
\bibitem{WilsonMcNaughton1994} M. A. Wilson and B. L. McNaughton, \emph{Reactivation of hippocampal ensemble memories during sleep}, Science \textbf{265}, 676 (1994).
\bibitem{BarnesSharp1999} N. M. Barnes and T. Sharp, \emph{A review of central 5-HT receptors and their function}, Neuropharmacology \textbf{38}, 1083 (1999).
\bibitem{Bunin1998} M. A. Bunin and R. M. Wightman, \emph{Quantitative evaluation of 5-hydroxytryptamine (serotonin) neuronal release and uptake: an investigation of extrasynaptic transmission}, J. Neurosci. \textbf{18}, 4854 (1998).
\bibitem{Sprouse1987} J. S. Sprouse and G. K. Aghajanian, \emph{Electrophysiological responses of serotoninergic dorsal raphe neurons to 5-HT$_{1A}$ and 5-HT$_{1B}$ agonists}, Synapse \textbf{1}, 3 (1987).
\bibitem{Sykova2008} E. Syková and C. Nicholson, \emph{Diffusion in brain extracellular space}, Physiol. Rev. \textbf{88}, 1277 (2008).
\bibitem{WilsonNicoll2001} R. I. Wilson and R. A. Nicoll, \emph{Endogenous cannabinoids mediate retrograde signalling at hippocampal synapses}, Nature \textbf{410}, 588 (2001).
\bibitem{Garthwaite2008} J. Garthwaite, \emph{Concepts of neural nitric oxide-mediated transmission}, Eur. J. Neurosci. \textbf{27}, 2783 (2008).
\bibitem{vandenPol2012} A. N. van den Pol, \emph{Neuropeptide transmission in brain circuits}, Neuron \textbf{76}, 98 (2012).
\bibitem{Dunwiddie2001} T. V. Dunwiddie and S. A. Masino, \emph{The role and regulation of adenosine in the central nervous system}, Annu. Rev. Neurosci. \textbf{24}, 31 (2001).
\bibitem{Skaggs1996} W. E. Skaggs and B. L. McNaughton, \emph{Replay of neuronal firing sequences in rat hippocampus during sleep following spatial experience}, Science \textbf{271}, 1870 (1996).
\bibitem{Baylor1979} D. A. Baylor, T. D. Lamb, and K.-W. Yau, \emph{Responses of retinal rods to single photons}, J. Physiol. \textbf{288}, 613 (1979).
\bibitem{Hudspeth1989} A. J. Hudspeth, \emph{How the ear's works work}, Nature \textbf{341}, 397 (1989).
\bibitem{NakaRushton1966} K. I. Naka and W. A. H. Rushton, \emph{S-potentials from colour units in the retina of fish (Cyprinidae)}, J. Physiol. \textbf{185}, 536 (1966).
\bibitem{Lichtsteiner2008} P. Lichtsteiner, C. Posch, and T. Delbruck, \emph{A $128\times128$ 120 dB 15 $\mu$s latency asynchronous temporal contrast vision sensor}, IEEE J. Solid-State Circuits \textbf{43}, 566 (2008).
\bibitem{Koch2006} K. Koch et al., \emph{How much the eye tells the brain}, Curr. Biol. \textbf{16}, 1428 (2006).
\bibitem{Robles2001} L. Robles and M. A. Ruggero, \emph{Mechanics of the mammalian cochlea}, Physiol. Rev. \textbf{81}, 1305 (2001).
\bibitem{Mombaerts2004} P. Mombaerts, \emph{Genes and ligands for odorant, vomeronasal and taste receptors}, Nat. Rev. Neurosci. \textbf{5}, 263 (2004).
\bibitem{WoodSlater2001} S. J. Wood and C. R. Slater, \emph{Safety factor at the neuromuscular junction}, Prog. Neurobiol. \textbf{64}, 393 (2001).
\bibitem{Henneman1957} E. Henneman, \emph{Relation between size of neurons and their susceptibility to discharge}, Science \textbf{126}, 1345 (1957).
\bibitem{HarrisWolpert1998} C. M. Harris and D. M. Wolpert, \emph{Signal-dependent noise determines motor planning}, Nature \textbf{394}, 780 (1998).
\bibitem{Hogan1984} N. Hogan, \emph{Adaptive control of mechanical impedance by coactivation of antagonist muscles}, IEEE Trans. Autom. Control \textbf{29}, 681 (1984).
\bibitem{McAuleyMarsden2000} J. H. McAuley and C. D. Marsden, \emph{Physiological and pathological tremors and rhythmic central motor control}, Brain \textbf{123}, 1545 (2000).
\bibitem{WolpertGhahramani2000} D. M. Wolpert and Z. Ghahramani, \emph{Computational principles of movement neuroscience}, Nat. Neurosci. \textbf{3}, 1212 (2000).
\bibitem{MarderBucher2001} E. Marder and D. Bucher, \emph{Central pattern generators and the control of rhythmic movements}, Curr. Biol. \textbf{11}, R986 (2001).
\bibitem{Miyazaki2014} K. W. Miyazaki et al., \emph{Optogenetic activation of dorsal raphe serotonin neurons enhances patience for future rewards}, Curr. Biol. \textbf{24}, 2033 (2014).
\bibitem{Fuglevand1993} A. J. Fuglevand, D. A. Winter, and A. E. Patla, \emph{Models of recruitment and rate coding organization in motor-unit pools}, J. Neurophysiol. \textbf{70}, 2470 (1993).
\bibitem{Curcio1990} C. A. Curcio, K. R. Sloan, R. E. Kalina, and A. E. Hendrickson, \emph{Human photoreceptor topography}, J. Comp. Neurol. \textbf{292}, 497 (1990).
\bibitem{Hayes1950} N. D. Hayes, \emph{Roots of the transcendental equation associated with a certain difference-differential equation}, J. London Math. Soc. \textbf{25}, 226 (1950).
\bibitem{MelzackWall1965} R. Melzack and P. D. Wall, \emph{Pain mechanisms: a new theory}, Science \textbf{150}, 971 (1965).
\bibitem{Fields2004} H. Fields, \emph{State-dependent opioid control of pain}, Nat. Rev. Neurosci. \textbf{5}, 565 (2004).
\bibitem{Levine1978} J. D. Levine, N. C. Gordon, and H. L. Fields, \emph{The mechanism of placebo analgesia}, Lancet \textbf{312}, 654 (1978).
\bibitem{Woolf1983} C. J. Woolf, \emph{Evidence for a central component of post-injury pain hypersensitivity}, Nature \textbf{306}, 686 (1983).
\bibitem{Seymour2004} B. Seymour et al., \emph{Temporal difference models describe higher-order learning in humans}, Nature \textbf{429}, 664 (2004).
\bibitem{EcclestonCrombez1999} C. Eccleston and G. Crombez, \emph{Pain demands attention: a cognitive-affective model of the interruptive function of pain}, Psychol. Bull. \textbf{125}, 356 (1999).
\bibitem{WidrowHoff1960} B. Widrow and M. E. Hoff, \emph{Adaptive switching circuits}, IRE WESCON Convention Record, part 4, 96 (1960).
\bibitem{Marr1969} D. Marr, \emph{A theory of cerebellar cortex}, J. Physiol. \textbf{202}, 437 (1969).
\bibitem{Albus1971} J. S. Albus, \emph{A theory of cerebellar function}, Math. Biosci. \textbf{10}, 25 (1971).
\bibitem{Cover1965} T. M. Cover, \emph{Geometrical and statistical properties of systems of linear inequalities with applications in pattern recognition}, IEEE Trans. Electron. Comput. \textbf{EC-14}, 326 (1965).
\bibitem{Ito2001} M. Ito, \emph{Cerebellar long-term depression: characterization, signal transduction, and functional roles}, Physiol. Rev. \textbf{81}, 1143 (2001).
\bibitem{Suvrathan2016} A. Suvrathan, H. L. Payne, and J. L. Raymond, \emph{Timing rules for synaptic plasticity matched to behavioral function}, Neuron \textbf{92}, 959 (2016).
\bibitem{Fujita1982} M. Fujita, \emph{Adaptive filter model of the cerebellum}, Biol. Cybern. \textbf{45}, 195 (1982).
\bibitem{Blakemore1998} S.-J. Blakemore, D. M. Wolpert, and C. D. Frith, \emph{Central cancellation of self-produced tickle sensation}, Nat. Neurosci. \textbf{1}, 635 (1998).
\bibitem{Smith2006} M. A. Smith, A. Ghazizadeh, and R. Shadmehr, \emph{Interacting adaptive processes with different timescales underlie short-term motor learning}, PLoS Biol. \textbf{4}, e179 (2006).
\bibitem{Kording2007} K. P. K\"ording, J. B. Tenenbaum, and R. Shadmehr, \emph{The dynamics of memory as a consequence of optimal adaptation to a changing body}, Nat. Neurosci. \textbf{10}, 779 (2007).
\bibitem{Yao2023} Z. Yao et al., \emph{A high-resolution transcriptomic and spatial atlas of cell types in the whole mouse brain}, Nature \textbf{624}, 317 (2023).
\bibitem{Sanzeni2020} A. Sanzeni, B. Akitake, H. C. Goldbach, C. E. Leedy, N. Brunel, and M. H. Histed, \emph{Inhibition stabilization is a widespread property of cortical networks}, eLife \textbf{9}, e54875 (2020).
\bibitem{Padamsey2022} Z. Padamsey, D. Katsanevaki, N. Dupuy, and N. L. Rochefort, \emph{Neocortex saves energy by reducing coding precision during food scarcity}, Neuron \textbf{110}, 280 (2022).
\bibitem{Stringer2019} C. Stringer, M. Pachitariu, N. Steinmetz, C. B. Reddy, M. Carandini, and K. D. Harris, \emph{Spontaneous behaviors drive multidimensional, brainwide activity}, Science \textbf{364}, eaav7893 (2019).
\bibitem{Bittner2017} K. C. Bittner, A. D. Milstein, C. Grienberger, S. Romani, and J. C. Magee, \emph{Behavioral time scale synaptic plasticity underlies CA1 place fields}, Science \textbf{357}, 1033 (2017).
\bibitem{WhittingtonBogacz2019} J. C. R. Whittington and R. Bogacz, \emph{Theories of error back-propagation in the brain}, Trends Cogn. Sci. \textbf{23}, 235 (2019).
\bibitem{Reimer2016} J. Reimer et al., \emph{Pupil fluctuations track rapid changes in adrenergic and cholinergic activity in cortex}, Nat. Commun. \textbf{7}, 13289 (2016).
\bibitem{Beniaguev2021} D. Beniaguev, I. Segev, and M. London, \emph{Single cortical neurons as deep artificial neural networks}, Neuron \textbf{109}, 2727 (2021).
\bibitem{Gidon2020} A. Gidon et al., \emph{Dendritic action potentials and computation in human layer 2/3 cortical neurons}, Science \textbf{367}, 83 (2020).
\bibitem{Gallego2022} G. Gallego et al., \emph{Event-based vision: a survey}, IEEE Trans. Pattern Anal. Mach. Intell. \textbf{44}, 154 (2022).
\bibitem{Berger1989} R. D. Berger, J. P. Saul, and R. J. Cohen, \emph{Transfer function analysis of autonomic regulation. I. Canine atrial rate response}, Am. J. Physiol. \textbf{256}, H142 (1989).
\bibitem{Walker2010} J. J. Walker, J. R. Terry, and S. L. Lightman, \emph{Origin of ultradian pulsatility in the hypothalamic--pituitary--adrenal axis}, Proc. R. Soc. B \textbf{277}, 1627 (2010).
\bibitem{Belchetz1978} P. E. Belchetz, T. M. Plant, Y. Nakai, E. J. Keogh, and E. Knobil, \emph{Hypophysial responses to continuous and intermittent delivery of hypothalamic gonadotropin-releasing hormone}, Science \textbf{202}, 631 (1978).
\bibitem{Czeisler1999} C. A. Czeisler et al., \emph{Stability, precision, and near-24-hour period of the human circadian pacemaker}, Science \textbf{284}, 2177 (1999).
\bibitem{TakahashiJS2017} J. S. Takahashi, \emph{Transcriptional architecture of the mammalian circadian clock}, Nat. Rev. Genet. \textbf{18}, 164 (2017).
\bibitem{Musk2019} E. Musk and Neuralink, \emph{An integrated brain-machine interface platform with thousands of channels}, J. Med. Internet Res. \textbf{21}, e16194 (2019).
\bibitem{Hochberg2006} L. R. Hochberg et al., \emph{Neuronal ensemble control of prosthetic devices by a human with tetraplegia}, Nature \textbf{442}, 164 (2006).
\bibitem{Willett2021} F. R. Willett et al., \emph{High-performance brain-to-text communication via handwriting}, Nature \textbf{593}, 249 (2021).
\bibitem{Willett2023} F. R. Willett et al., \emph{A high-performance speech neuroprosthesis}, Nature \textbf{620}, 1031 (2023).
\bibitem{Gallego2017} J. A. Gallego, M. G. Perich, L. E. Miller, and S. A. Solla, \emph{Neural manifolds for the control of movement}, Neuron \textbf{94}, 978 (2017).
\bibitem{Fernandez2021} E. Fern\'andez et al., \emph{Visual percepts evoked with an intracortical 96-channel microelectrode array inserted in human occipital cortex}, J. Clin. Invest. \textbf{131}, e151331 (2021).
\bibitem{Tritsch2012} N. X. Tritsch, J. B. Ding, and B. L. Sabatini, \emph{Dopaminergic neurons inhibit striatal output through non-canonical release of GABA}, Nature \textbf{490}, 262 (2012).
\bibitem{Zhang2015} S. Zhang et al., \emph{Dopaminergic and glutamatergic microdomains in a subset of rodent mesoaccumbens axons}, Nat. Neurosci. \textbf{18}, 386 (2015).
\bibitem{Mongillo2008} G. Mongillo, O. Barak, and M. Tsodyks, \emph{Synaptic theory of working memory}, Science \textbf{319}, 1543 (2008).
\bibitem{Stokes2015} M. G. Stokes, \emph{`Activity-silent' working memory in prefrontal cortex: a dynamic coding framework}, Trends Cogn. Sci. \textbf{19}, 394 (2015).
\bibitem{Smith2013} S. L. Smith, I. T. Smith, T. Branco, and M. H\"ausser, \emph{Dendritic spikes enhance stimulus selectivity in cortical neurons in vivo}, Nature \textbf{503}, 115 (2013).
\bibitem{Baden2016} T. Baden et al., \emph{The functional diversity of retinal ganglion cells in the mouse}, Nature \textbf{529}, 345 (2016).
\bibitem{Lottem2018} E. Lottem et al., \emph{Activation of serotonin neurons promotes active persistence in a probabilistic foraging task}, Nat. Commun. \textbf{9}, 1000 (2018).
\bibitem{Payeur2021} A. Payeur, J. Guerguiev, F. Zenke, B. A. Richards, and R. Naud, \emph{Burst-dependent synaptic plasticity can coordinate learning in hierarchical circuits}, Nat. Neurosci. \textbf{24}, 1010 (2021).
\bibitem{MehonicKenyon2022} A. Mehonic and A. J. Kenyon, \emph{Brain-inspired computing needs a master plan}, Nature \textbf{604}, 255 (2022).
\bibitem{Poe2020} G. R. Poe et al., \emph{Locus coeruleus: a new look at the blue spot}, Nat. Rev. Neurosci. \textbf{21}, 644 (2020).
\bibitem{Hangya2015} B. Hangya, S. P. Ranade, M. Lorenc, and A. Kepecs, \emph{Central cholinergic neurons are rapidly recruited by reinforcement feedback}, Cell \textbf{162}, 1155 (2015).
\bibitem{FernandezRuiz2019} A. Fern\'andez-Ruiz et al., \emph{Long-duration hippocampal sharp wave ripples improve memory}, Science \textbf{364}, 1082 (2019).
\bibitem{Herzfeld2018} D. J. Herzfeld, Y. Kojima, R. Soetedjo, and R. Shadmehr, \emph{Encoding of error and learning to correct that error by the Purkinje cells of the cerebellum}, Nat. Neurosci. \textbf{21}, 736 (2018).
\bibitem{Duan2014} B. Duan et al., \emph{Identification of spinal circuits transmitting and gating mechanical pain}, Cell \textbf{159}, 1417 (2014).
\bibitem{Tremblay2016} R. Tremblay, S. Lee, and B. Rudy, \emph{GABAergic interneurons in the neocortex: from cellular properties to circuits}, Neuron \textbf{91}, 260 (2016).
\bibitem{Ren2018} J. Ren et al., \emph{Anatomically defined and functionally distinct dorsal raphe serotonin sub-systems}, Cell \textbf{175}, 472 (2018).
\bibitem{Pfeffer2013} C. K. Pfeffer, M. Xue, M. He, Z. J. Huang, and M. Scanziani, \emph{Inhibition of inhibition in visual cortex: the logic of connections between molecularly distinct interneurons}, Nat. Neurosci. \textbf{16}, 1068 (2013).
\bibitem{Pi2013} H.-J. Pi et al., \emph{Cortical interneurons that specialize in disinhibitory control}, Nature \textbf{503}, 521 (2013).
\bibitem{Keller2018} G. B. Keller and T. D. Mrsic-Flogel, \emph{Predictive processing: a canonical cortical computation}, Neuron \textbf{100}, 424 (2018).
\bibitem{HubelWiesel1962} D. H. Hubel and T. N. Wiesel, \emph{Receptive fields, binocular interaction and functional architecture in the cat's visual cortex}, J. Physiol. \textbf{160}, 106 (1962).
\bibitem{Riesenhuber1999} M. Riesenhuber and T. Poggio, \emph{Hierarchical models of object recognition in cortex}, Nat. Neurosci. \textbf{2}, 1019 (1999).
\bibitem{Fukushima1980} K. Fukushima, \emph{Neocognitron: a self-organizing neural network model for a mechanism of pattern recognition unaffected by shift in position}, Biol. Cybern. \textbf{36}, 193 (1980).
\bibitem{Yamins2014} D. L. K. Yamins et al., \emph{Performance-optimized hierarchical models predict neural responses in higher visual cortex}, Proc. Natl. Acad. Sci. USA \textbf{111}, 8619 (2014).
\bibitem{Hung2005} C. P. Hung, G. Kreiman, T. Poggio, and J. J. DiCarlo, \emph{Fast readout of object identity from macaque inferior temporal cortex}, Science \textbf{310}, 863 (2005).
\bibitem{WiskottSejnowski2002} L. Wiskott and T. J. Sejnowski, \emph{Slow feature analysis: unsupervised learning of invariances}, Neural Comput. \textbf{14}, 715 (2002).
\bibitem{Foldiak1991} P. F\"oldi\'ak, \emph{Learning invariance from transformation sequences}, Neural Comput. \textbf{3}, 194 (1991).
\bibitem{LiDiCarlo2008} N. Li and J. J. DiCarlo, \emph{Unsupervised natural experience rapidly alters invariant object representation in visual cortex}, Science \textbf{321}, 1502 (2008).
\bibitem{RaoBallard1999} R. P. N. Rao and D. H. Ballard, \emph{Predictive coding in the visual cortex: a functional interpretation of some extra-classical receptive-field effects}, Nat. Neurosci. \textbf{2}, 79 (1999).
\bibitem{Kar2019} K. Kar, J. Kubilius, K. Schmidt, E. B. Issa, and J. J. DiCarlo, \emph{Evidence that recurrent circuits are critical to the ventral stream's execution of core object recognition behavior}, Nat. Neurosci. \textbf{22}, 974 (2019).
\bibitem{Szegedy2014} C. Szegedy et al., \emph{Intriguing properties of neural networks}, in \emph{Proc. ICLR} (2014), arXiv:1312.6199.
\bibitem{Weiss2002} Y. Weiss, E. P. Simoncelli, and E. H. Adelson, \emph{Motion illusions as optimal percepts}, Nat. Neurosci. \textbf{5}, 598 (2002).
\bibitem{Purves2011} D. Purves, W. T. Wojtach, and R. B. Lotto, \emph{Understanding vision in wholly empirical terms}, Proc. Natl. Acad. Sci. USA \textbf{108}, Suppl.~3, 15588 (2011).
\bibitem{LaferSousa2015} R. Lafer-Sousa, K. L. Hermann, and B. R. Conway, \emph{Striking individual differences in color perception uncovered by `the dress' photograph}, Curr. Biol. \textbf{25}, R545 (2015).
\bibitem{Wallisch2017} P. Wallisch, \emph{Illumination assumptions account for individual differences in the perceptual interpretation of a profoundly ambiguous stimulus in the color domain: ``The dress''}, J. Vis. \textbf{17}(4), 5 (2017).
\bibitem{Schwartz2009} O. Schwartz, T. J. Sejnowski, and P. Dayan, \emph{Perceptual organization in the tilt illusion}, J. Vis. \textbf{9}(4), 19 (2009).
\bibitem{SchillerCarvey2005} P. H. Schiller and C. E. Carvey, \emph{The Hermann grid illusion revisited}, Perception \textbf{34}, 1375 (2005).
\bibitem{Nijhawan1994} R. Nijhawan, \emph{Motion extrapolation in catching}, Nature \textbf{370}, 256 (1994).
\bibitem{Conway2005} B. R. Conway, A. Kitaoka, A. Yazdanbakhsh, C. C. Pack, and M. S. Livingstone, \emph{Neural basis for a powerful static motion illusion}, J. Neurosci. \textbf{25}, 5651 (2005).
\bibitem{OteroMillan2012} J. Otero-Millan, S. L. Macknik, and S. Martinez-Conde, \emph{Microsaccades and blinks trigger illusory rotation in the ``rotating snakes'' illusion}, J. Neurosci. \textbf{32}, 6043 (2012).
\bibitem{MorenoBote2007} R. Moreno-Bote, J. Rinzel, and N. Rubin, \emph{Noise-induced alternations in an attractor network model of perceptual bistability}, J. Neurophysiol. \textbf{98}, 1125 (2007).
\bibitem{Watanabe2018} E. Watanabe, A. Kitaoka, K. Sakamoto, M. Yasugi, and K. Tanaka, \emph{Illusory motion reproduced by deep neural networks trained for prediction}, Front. Psychol. \textbf{9}, 345 (2018).
\bibitem{GomezVilla2020} A. G\'omez-Villa, A. Mart\'in, J. Vazquez-Corral, M. Bertalm\'io, and J. Malo, \emph{Color illusions also deceive CNNs for low-level vision tasks: analysis and implications}, Vision Res. \textbf{176}, 156 (2020).
\bibitem{delCastillo1954} J. del Castillo and B. Katz, \emph{Quantal components of the end-plate potential}, J. Physiol. \textbf{124}, 560 (1954).
\bibitem{Nowak1984} L. Nowak, P. Bregestovski, P. Ascher, A. Herbet, and A. Prochiantz, \emph{Magnesium gates glutamate-activated channels in mouse central neurones}, Nature \textbf{307}, 462 (1984).
\bibitem{Malinow2002} R. Malinow and R. C. Malenka, \emph{AMPA receptor trafficking and synaptic plasticity}, Annu. Rev. Neurosci. \textbf{25}, 103 (2002).
\bibitem{Matsuzaki2004} M. Matsuzaki, N. Honkura, G. C. R. Ellis-Davies, and H. Kasai, \emph{Structural basis of long-term potentiation in single dendritic spines}, Nature \textbf{429}, 761 (2004).
\bibitem{Rudomin1999} P. Rudomin and R. F. Schmidt, \emph{Presynaptic inhibition in the vertebrate spinal cord revisited}, Exp. Brain Res. \textbf{129}, 1 (1999).
\bibitem{HutcheonYarom2000} B. Hutcheon and Y. Yarom, \emph{Resonance, oscillation and the intrinsic frequency preferences of neurons}, Trends Neurosci. \textbf{23}, 216 (2000).
\bibitem{Seung1996} H. S. Seung, \emph{How the brain keeps the eyes still}, Proc. Natl. Acad. Sci. USA \textbf{93}, 13339 (1996).
\bibitem{Hounsgaard1988} J. Hounsgaard, H. Hultborn, B. Jespersen, and O. Kiehn, \emph{Bistability of alpha-motoneurones in the decerebrate cat and in the acute spinal cat after intravenous 5-hydroxytryptophan}, J. Physiol. \textbf{405}, 345 (1988).
\bibitem{Izhikevich2007} E. M. Izhikevich, \emph{Dynamical Systems in Neuroscience: The Geometry of Excitability and Bursting} (MIT Press, Cambridge, MA, 2007).
\bibitem{AgmonSnir1998} H. Agmon-Snir, C. E. Carr, and J. Rinzel, \emph{The role of dendrites in auditory coincidence detection}, Nature \textbf{393}, 268 (1998).
\bibitem{BorstSoria2012} J. G. G. Borst and J. Soria van Hoeve, \emph{The calyx of Held synapse: from model synapse to auditory relay}, Annu. Rev. Physiol. \textbf{74}, 199 (2012).
\bibitem{Euler2002} T. Euler, P. B. Detwiler, and W. Denk, \emph{Directionally selective calcium signals in dendrites of starburst amacrine cells}, Nature \textbf{418}, 845 (2002).
\bibitem{AstonJonesBloom1981} G. Aston-Jones and F. E. Bloom, \emph{Activity of norepinephrine-containing locus coeruleus neurons in behaving rats anticipates fluctuations in the sleep-waking cycle}, J. Neurosci. \textbf{1}, 876 (1981).
\bibitem{BrooksPeever2012} P. L. Brooks and J. H. Peever, \emph{Identification of the transmitter and receptor mechanisms responsible for REM sleep paralysis}, J. Neurosci. \textbf{32}, 9785 (2012).
\bibitem{Borbely1982} A. A. Borb\'ely, \emph{A two process model of sleep regulation}, Hum. Neurobiol. \textbf{1}, 195 (1982).
\bibitem{PorkkaHeiskanen1997} T. Porkka-Heiskanen et al., \emph{Adenosine: a mediator of the sleep-inducing effects of prolonged wakefulness}, Science \textbf{276}, 1265 (1997).
\bibitem{Steriade1993} M. Steriade, A. N\'u\~nez, and F. Amzica, \emph{A novel slow ($<1$ Hz) oscillation of neocortical neurons in vivo: depolarizing and hyperpolarizing components}, J. Neurosci. \textbf{13}, 3252 (1993).
\bibitem{McCormick1992} D. A. McCormick, \emph{Neurotransmitter actions in the thalamus and cerebral cortex and their role in neuromodulation of thalamocortical activity}, Prog. Neurobiol. \textbf{39}, 337 (1992).
\bibitem{SteriadeMcCormickSejnowski1993} M. Steriade, D. A. McCormick, and T. J. Sejnowski, \emph{Thalamocortical oscillations in the sleeping and aroused brain}, Science \textbf{262}, 679 (1993).
\bibitem{Nadasdy1999} Z. N\'adasdy et al., \emph{Replay and time compression of recurring spike sequences in the hippocampus}, J. Neurosci. \textbf{19}, 9497 (1999).
\bibitem{Maingret2016} N. Maingret et al., \emph{Hippocampo-cortical coupling mediates memory consolidation during sleep}, Nat. Neurosci. \textbf{19}, 959 (2016).
\bibitem{McClelland1995} J. L. McClelland, B. L. McNaughton, and R. C. O'Reilly, \emph{Why there are complementary learning systems in the hippocampus and neocortex}, Psychol. Rev. \textbf{102}, 419 (1995).
\bibitem{TononiCirelli2014} G. Tononi and C. Cirelli, \emph{Sleep and the price of plasticity: from synaptic and cellular homeostasis to memory consolidation and integration}, Neuron \textbf{81}, 12 (2014).
\bibitem{deVivo2017} L. de Vivo et al., \emph{Ultrastructural evidence for synaptic scaling across the wake/sleep cycle}, Science \textbf{355}, 507 (2017).
\bibitem{Xie2013} L. Xie et al., \emph{Sleep drives metabolite clearance from the adult brain}, Science \textbf{342}, 373 (2013).
\bibitem{Miao2024} A. Miao et al., \emph{Brain clearance is reduced during sleep and anesthesia}, Nat. Neurosci. \textbf{27}, 1046 (2024).
\bibitem{Vyazovskiy2011} V. V. Vyazovskiy et al., \emph{Local sleep in awake rats}, Nature \textbf{472}, 443 (2011).
\bibitem{Lancaster2013} M. A. Lancaster et al., \emph{Cerebral organoids model human brain development and microcephaly}, Nature \textbf{501}, 373 (2013).
\bibitem{Jordan2024} F. D. Jordan et al., \emph{Open and remotely accessible Neuroplatform for research in wetware computing}, Front. Artif. Intell. \textbf{7}, 1376042 (2024).
\bibitem{Smirnova2023} L. Smirnova et al., \emph{Organoid intelligence (OI): the new frontier in biocomputing and intelligence-in-a-dish}, Front. Sci. \textbf{1}, 1017235 (2023).
\bibitem{Wagenaar2006} D. A. Wagenaar, J. Pine, and S. M. Potter, \emph{An extremely rich repertoire of bursting patterns during the development of cortical cultures}, BMC Neurosci. \textbf{7}, 11 (2006).
\bibitem{Shahaf2001} G. Shahaf and S. Marom, \emph{Learning in networks of cortical neurons}, J. Neurosci. \textbf{21}, 8782 (2001).
\bibitem{Kagan2022} B. J. Kagan et al., \emph{In vitro neurons learn and exhibit sentience when embodied in a simulated game-world}, Neuron \textbf{110}, 3952 (2022).
\bibitem{Friston2010} K. Friston, \emph{The free-energy principle: a unified brain theory?}, Nat. Rev. Neurosci. \textbf{11}, 127 (2010).
\bibitem{Balci2023} F. Balci et al., \emph{A response to claims of emergent intelligence and sentience in a dish}, Neuron \textbf{111}, 604 (2023).
\bibitem{Bakkum2008} D. J. Bakkum, Z. C. Chao, and S. M. Potter, \emph{Spatio-temporal electrical stimuli shape behavior of an embodied cortical network in a goal-directed learning task}, J. Neural Eng. \textbf{5}, 310 (2008).
\bibitem{Maass2002} W. Maass, T. Natschl\"ager, and H. Markram, \emph{Real-time computing without stable states: a new framework for neural computation based on perturbations}, Neural Comput. \textbf{14}, 2531 (2002).
\bibitem{JaegerHaas2004} H. Jaeger and H. Haas, \emph{Harnessing nonlinearity: predicting chaotic systems and saving energy in wireless communication}, Science \textbf{304}, 78 (2004).
\bibitem{Cai2023} H. Cai et al., \emph{Brain organoid reservoir computing for artificial intelligence}, Nat. Electron. \textbf{6}, 1032 (2023).
\bibitem{Habibollahi2023} F. Habibollahi, B. J. Kagan, A. N. Burkitt, and C. French, \emph{Critical dynamics arise during structured information presentation within embodied in vitro neuronal networks}, Nat. Commun. \textbf{14}, 5287 (2023).
\bibitem{Robbins2026} A. Robbins et al., \emph{Goal-directed learning in cortical organoids}, Cell Rep. \textbf{45}, 116984 (2026).
\bibitem{Keene2020} S. T. Keene et al., \emph{A biohybrid synapse with neurotransmitter-mediated plasticity}, Nat. Mater. \textbf{19}, 969 (2020).
\bibitem{Jo2016} J. Jo et al., \emph{Midbrain-like organoids from human pluripotent stem cells contain functional dopaminergic and neuromelanin-producing neurons}, Cell Stem Cell \textbf{19}, 248 (2016).
\bibitem{Hendry1987} S. H. Hendry, H. D. Schwark, E. G. Jones, and J. Yan, \emph{Numbers and proportions of GABA-immunoreactive neurons in different areas of monkey cerebral cortex}, J. Neurosci. \textbf{7}, 1503 (1987).
\bibitem{Tang2001synapses} Y. Tang, J. R. Nyengaard, D. M. G. De Groot, and H. J. G. Gundersen, \emph{Total regional and global number of synapses in the human brain neocortex}, Synapse \textbf{41}, 258 (2001).
\bibitem{Pakkenberg1997} B. Pakkenberg and H. J. G. Gundersen, \emph{Neocortical neuron number in humans: effect of sex and age}, J. Comp. Neurol. \textbf{384}, 312 (1997).
\bibitem{Matsuda2009} W. Matsuda, T. Furuta, K. C. Nakamura, H. Hioki, F. Fujiyama, R. Arai, and T. Kaneko, \emph{Single nigrostriatal dopaminergic neurons form widely spread and highly dense axonal arborizations in the neostriatum}, J. Neurosci. \textbf{29}, 444 (2009).
\bibitem{Pakkenberg1991} B. Pakkenberg, A. M{\o}ller, H. J. G. Gundersen, A. Mouritzen Dam, and H. Pakkenberg, \emph{The absolute number of nerve cells in substantia nigra in normal subjects and in patients with Parkinson's disease estimated with an unbiased stereological method}, J. Neurol. Neurosurg. Psychiatry \textbf{54}, 30 (1991).
\bibitem{Baker1990} K. G. Baker, G. M. Halliday, and I. T\"ork, \emph{Cytoarchitecture of the human dorsal raphe nucleus}, J. Comp. Neurol. \textbf{301}, 147 (1990).
\bibitem{Baker1991} K. G. Baker, G. M. Halliday, P. Halasz, J.-P. Hornung, L. B. Geffen, R. G. H. Cotton, and I. T\"ork, \emph{Cytoarchitecture of serotonin-synthesizing neurons in the pontine tegmentum of the human brain}, Synapse \textbf{7}, 301 (1991).
\bibitem{Airaksinen1991} M. S. Airaksinen, A. Paetau, L. Palj\"arvi, K. Reinikainen, P. Riekkinen, R. Suomalainen, and P. Panula, \emph{Histamine neurons in human hypothalamus: anatomy in normal and Alzheimer diseased brains}, Neuroscience \textbf{44}, 465 (1991).
\bibitem{Colquhoun1992} D. Colquhoun, P. Jonas, and B. Sakmann, \emph{Action of brief pulses of glutamate on AMPA/kainate receptors in patches from different neurones of rat hippocampal slices}, J. Physiol. \textbf{458}, 261 (1992).
\bibitem{DutarNicoll1988} P. Dutar and R. A. Nicoll, \emph{A physiological role for GABA$_B$ receptors in the central nervous system}, Nature \textbf{332}, 156 (1988).

\bibitem{Rauch2003} A. Rauch, G. La Camera, H.-R. L\"uscher, W. Senn, and S. Fusi, \emph{Neocortical pyramidal cells respond as integrate-and-fire neurons to in vivo-like input currents}, J. Neurophysiol. \textbf{90}, 1598 (2003).
\bibitem{KlumppEady1956} R. G. Klumpp and H. R. Eady, \emph{Some measurements of interaural time difference thresholds}, J. Acoust. Soc. Am. \textbf{28}, 859 (1956).
\bibitem{Brand2002} A. Brand, O. Behrend, T. Marquardt, D. McAlpine, and B. Grothe, \emph{Precise inhibition is essential for microsecond interaural time difference coding}, Nature \textbf{417}, 543 (2002).
\bibitem{Funahashi1989} S. Funahashi, C. J. Bruce, and P. S. Goldman-Rakic, \emph{Mnemonic coding of visual space in the monkey's dorsolateral prefrontal cortex}, J. Neurophysiol. \textbf{61}, 331 (1989).
\bibitem{WangMarkram2006} Y. Wang, H. Markram, P. H. Goodman, T. K. Berger, J. Ma, and P. S. Goldman-Rakic, \emph{Heterogeneity in the pyramidal network of the medial prefrontal cortex}, Nat. Neurosci. \textbf{9}, 534 (2006).

\bibitem{McCullochPitts1943} W. S. McCulloch and W. Pitts, \emph{A logical calculus of the ideas immanent in nervous activity}, Bull. Math. Biophys. \textbf{5}, 115 (1943).
\bibitem{Fried2002} S. I. Fried, T. A. M\"unch, and F. S. Werblin, \emph{Mechanisms and circuitry underlying directional selectivity in the retina}, Nature \textbf{420}, 411 (2002).

\bibitem{Hromadka2008} T. Hrom\'adka, M. R. DeWeese, and A. M. Zador, \emph{Sparse representation of sounds in the unanesthetized auditory cortex}, PLoS Biol. \textbf{6}, e16 (2008).
\bibitem{Abbott1997depression} L. F. Abbott, J. A. Varela, K. Sen, and S. B. Nelson, \emph{Synaptic depression and cortical gain control}, Science \textbf{275}, 221 (1997).
\bibitem{ShadlenNewsome1998} M. N. Shadlen and W. T. Newsome, \emph{The variable discharge of cortical neurons: implications for connectivity, computation, and information coding}, J. Neurosci. \textbf{18}, 3870 (1998).

\bibitem{TrulsonJacobs1979} M. E. Trulson and B. L. Jacobs, \emph{Raphe unit activity in freely moving cats: correlation with level of behavioral arousal}, Brain Res. \textbf{163}, 135 (1979).
\bibitem{GagnonParent2014} D. Gagnon and M. Parent, \emph{Distribution of VGLUT3 in highly collateralized axons from the rat dorsal raphe nucleus as revealed by single-neuron reconstructions}, PLoS One \textbf{9}, e87709 (2014).
\bibitem{VandermaelenAghajanian1983} C. P. Vandermaelen and G. K. Aghajanian, \emph{Electrophysiological and pharmacological characterization of serotonergic dorsal raphe neurons recorded extracellularly and intracellularly in rat brain slices}, Brain Res. \textbf{289}, 109 (1983).
\bibitem{CourtneyFord2016} N. A. Courtney and C. P. Ford, \emph{Mechanisms of 5-HT$_{1A}$ receptor-mediated transmission in dorsal raphe serotonin neurons}, J. Physiol. \textbf{594}, 953 (2016).
\bibitem{Hashemi2012serotonin} P. Hashemi et al., \emph{Brain dopamine and serotonin differ in regulation and its consequences}, Proc. Natl. Acad. Sci. USA \textbf{109}, 11510 (2012).
\bibitem{Black1934feedback} H. S. Black, \emph{Stabilized feedback amplifiers}, Bell Syst. Tech. J. \textbf{13}, 1 (1934).
\bibitem{KobayashiSchultz2008} S. Kobayashi and W. Schultz, \emph{Influence of reward delays on responses of dopamine neurons}, J. Neurosci. \textbf{28}, 7837 (2008).
\bibitem{Schweighofer2008} N. Schweighofer et al., \emph{Low-serotonin levels increase delayed reward discounting in humans}, J. Neurosci. \textbf{28}, 4528 (2008).

\bibitem{Malenka1988calcium} R. C. Malenka, J. A. Kauer, R. S. Zucker, and R. A. Nicoll, \emph{Postsynaptic calcium is sufficient for potentiation of hippocampal synaptic transmission}, Science \textbf{242}, 81 (1988).
\bibitem{Bliss1973LTP} T. V. P. Bliss and T. L{\o}mo, \emph{Long-lasting potentiation of synaptic transmission in the dentate area of the anaesthetized rabbit following stimulation of the perforant path}, J. Physiol. \textbf{232}, 331 (1973).
\bibitem{Kelso1986hebb} S. R. Kelso, A. H. Ganong, and T. H. Brown, \emph{Hebbian synapses in hippocampus}, Proc. Natl. Acad. Sci. USA \textbf{83}, 5326 (1986).
\bibitem{He2015eligibility} K. He et al., \emph{Distinct eligibility traces for LTP and LTD in cortical synapses}, Neuron \textbf{88}, 528 (2015).
\bibitem{Olveczky2005} B. P. \"Olveczky, A. S. Andalman, and M. S. Fee, \emph{Vocal experimentation in the juvenile songbird requires a basal ganglia circuit}, PLoS Biol. \textbf{3}, e153 (2005).
\bibitem{TumerBrainard2007} E. C. Tumer and M. S. Brainard, \emph{Performance variability enables adaptive plasticity of `crystallized' adult birdsong}, Nature \textbf{450}, 1240 (2007).
\bibitem{Eshel2015arithmetic} N. Eshel et al., \emph{Arithmetic and local circuitry underlying dopamine prediction errors}, Nature \textbf{525}, 243 (2015).
\bibitem{GraceOnn1989} A. A. Grace and S.-P. Onn, \emph{Morphology and electrophysiological properties of immunocytochemically identified rat dopamine neurons recorded in vitro}, J. Neurosci. \textbf{9}, 3463 (1989).
\bibitem{ODoherty2004actor} J. O'Doherty et al., \emph{Dissociable roles of ventral and dorsal striatum in instrumental conditioning}, Science \textbf{304}, 452 (2004).

\bibitem{NeweyPowell1987} W. K. Newey and J. L. Powell, \emph{Asymmetric least squares estimation and testing}, Econometrica \textbf{55}, 819 (1987).
\bibitem{Beierholm2013vigor} U. Beierholm et al., \emph{Dopamine modulates reward-related vigor}, Neuropsychopharmacology \textbf{38}, 1495 (2013).
\bibitem{Howe2013ramp} M. W. Howe, P. L. Tierney, S. G. Sandberg, P. E. M. Phillips, and A. M. Graybiel, \emph{Prolonged dopamine signalling in striatum signals proximity and value of distant rewards}, Nature \textbf{500}, 575 (2013).

\bibitem{Baker1989LC} K. G. Baker, I. T\"ork, J.-P. Hornung, and P. Halasz, \emph{The human locus coeruleus complex: an immunohistochemical and three dimensional reconstruction study}, Exp. Brain Res. \textbf{77}, 257 (1989).
\bibitem{Schwarz2015LC} L. A. Schwarz et al., \emph{Viral-genetic tracing of the input--output organization of a central noradrenaline circuit}, Nature \textbf{524}, 88 (2015).
\bibitem{Uematsu2017LC} A. Uematsu et al., \emph{Modular organization of the brainstem noradrenaline system coordinates opposing learning states}, Nat. Neurosci. \textbf{20}, 1602 (2017).
\bibitem{AstonJones1994vigilance} G. Aston-Jones, J. Rajkowski, P. Kubiak, and T. Alexinsky, \emph{Locus coeruleus neurons in monkey are selectively activated by attended cues in a vigilance task}, J. Neurosci. \textbf{14}, 4467 (1994).
\bibitem{Clayton2004LC} E. C. Clayton, J. Rajkowski, J. D. Cohen, and G. Aston-Jones, \emph{Phasic activation of monkey locus ceruleus neurons by simple decisions in a forced-choice task}, J. Neurosci. \textbf{24}, 9914 (2004).
\bibitem{Foote1975NE} S. L. Foote, R. Freedman, and A. P. Oliver, \emph{Effects of putative neurotransmitters on neuronal activity in monkey auditory cortex}, Brain Res. \textbf{86}, 229 (1975).
\bibitem{JepmaNieuwenhuis2011} M. Jepma and S. Nieuwenhuis, \emph{Pupil diameter predicts changes in the exploration--exploitation trade-off: evidence for the adaptive gain theory}, J. Cogn. Neurosci. \textbf{23}, 1587 (2011).
\bibitem{Tervo2014stochastic} D. G. R. Tervo et al., \emph{Behavioral variability through stochastic choice and its gating by anterior cingulate cortex}, Cell \textbf{159}, 21 (2014).
\bibitem{Usher1999LC} M. Usher, J. D. Cohen, D. Servan-Schreiber, J. Rajkowski, and G. Aston-Jones, \emph{The role of locus coeruleus in the regulation of cognitive performance}, Science \textbf{283}, 549 (1999).

\bibitem{Mesulam1983} M.-M. Mesulam, E. J. Mufson, B. H. Wainer, and A. I. Levey, \emph{Central cholinergic pathways in the rat: an overview based on an alternative nomenclature (Ch1--Ch6)}, Neuroscience \textbf{10}, 1185 (1983).
\bibitem{Marrosu1995} F. Marrosu et al., \emph{Microdialysis measurement of cortical and hippocampal acetylcholine release during sleep-wake cycle in freely moving cats}, Brain Res. \textbf{671}, 329 (1995).
\bibitem{Picciotto2012} M. R. Picciotto, M. J. Higley, and Y. S. Mineur, \emph{Acetylcholine as a neuromodulator: cholinergic signaling shapes nervous system function and behavior}, Neuron \textbf{76}, 116 (2012).
\bibitem{HasselmoSchnell1994} M. E. Hasselmo and E. Schnell, \emph{Laminar selectivity of the cholinergic suppression of synaptic transmission in rat hippocampal region CA1: computational modeling and brain slice physiology}, J. Neurosci. \textbf{14}, 3898 (1994).
\bibitem{Gil1997} Z. Gil, B. W. Connors, and Y. Amitai, \emph{Differential regulation of neocortical synapses by neuromodulators and activity}, Neuron \textbf{19}, 679 (1997).
\bibitem{KilgardMerzenich1998} M. P. Kilgard and M. M. Merzenich, \emph{Cortical map reorganization enabled by nucleus basalis activity}, Science \textbf{279}, 1714 (1998).
\bibitem{Atri2004} A. Atri et al., \emph{Blockade of central cholinergic receptors impairs new learning and increases proactive interference in a word paired-associate memory task}, Behav. Neurosci. \textbf{118}, 223 (2004).
\bibitem{Girardeau2009} G. Girardeau, K. Benchenane, S. I. Wiener, G. Buzs\'aki, and M. B. Zugaro, \emph{Selective suppression of hippocampal ripples impairs spatial memory}, Nat. Neurosci. \textbf{12}, 1222 (2009).

\input{supplement/bib_10}
\bibitem{CopenhagenJahr1989} D. R. Copenhagen and C. E. Jahr, \emph{Release of endogenous excitatory amino acids from turtle photoreceptors}, Nature \textbf{341}, 536 (1989).
\bibitem{GlowatzkiFuchs2002} E. Glowatzki and P. A. Fuchs, \emph{Transmitter release at the hair cell ribbon synapse}, Nat. Neurosci. \textbf{5}, 147 (2002).
\bibitem{Tinsley2016} J. N. Tinsley et al., \emph{Direct detection of a single photon by humans}, Nat. Commun. \textbf{7}, 12172 (2016).
\bibitem{Sellick1982} P. M. Sellick, R. Patuzzi, and B. M. Johnstone, \emph{Measurement of basilar membrane motion in the guinea pig using the M\"ossbauer technique}, J. Acoust. Soc. Am. \textbf{72}, 131 (1982).
\bibitem{Ruggero1997} M. A. Ruggero, N. C. Rich, A. Recio, S. S. Narayan, and L. Robles, \emph{Basilar-membrane responses to tones at the base of the chinchilla cochlea}, J. Acoust. Soc. Am. \textbf{101}, 2151 (1997).
\bibitem{NormannPerlman1979} R. A. Normann and I. Perlman, \emph{The effects of background illumination on the photoresponses of red and green cones}, J. Physiol. \textbf{286}, 491 (1979).
\bibitem{Laughlin1981} S. Laughlin, \emph{A simple coding procedure enhances a neuron's information capacity}, Z. Naturforsch. C \textbf{36}, 910 (1981).
\bibitem{CurcioAllen1990} C. A. Curcio and K. A. Allen, \emph{Topography of ganglion cells in human retina}, J. Comp. Neurol. \textbf{300}, 5 (1990).
\bibitem{Greenwood1990} D. D. Greenwood, \emph{A cochlear frequency-position function for several species --- 29 years later}, J. Acoust. Soc. Am. \textbf{87}, 2592 (1990).
\bibitem{Eguiluz2000} V. M. Egu\'iluz, M. Ospeck, Y. Choe, A. J. Hudspeth, and M. O. Magnasco, \emph{Essential nonlinearities in hearing}, Phys. Rev. Lett. \textbf{84}, 5232 (2000).
\bibitem{PalmerRussell1986} A. R. Palmer and I. J. Russell, \emph{Phase-locking in the cochlear nerve of the guinea-pig and its relation to the receptor potential of inner hair-cells}, Hear. Res. \textbf{24}, 1 (1986).
\bibitem{Malnic1999} B. Malnic, J. Hirono, T. Sato, and L. B. Buck, \emph{Combinatorial receptor codes for odors}, Cell \textbf{96}, 713 (1999).
\bibitem{Finger2005} T. E. Finger et al., \emph{ATP signaling is crucial for communication from taste buds to gustatory nerves}, Science \textbf{310}, 1495 (2005).
\bibitem{Huang2005} Y.-J. Huang et al., \emph{Mouse taste buds use serotonin as a neurotransmitter}, J. Neurosci. \textbf{25}, 843 (2005).
\bibitem{JohanssonVallbo1979} R. S. Johansson and A. B. Vallbo, \emph{Tactile sensibility in the human hand: relative and absolute densities of four types of mechanoreceptive units in glabrous skin}, J. Physiol. \textbf{286}, 283 (1979).
\bibitem{McKemy2002} D. D. McKemy, W. M. Neuhausser, and D. Julius, \emph{Identification of a cold receptor reveals a general role for TRP channels in thermosensation}, Nature \textbf{416}, 52 (2002).

\bibitem{Schmolesky1998} M. T. Schmolesky et al., \emph{Signal timing across the macaque visual system}, J. Neurophysiol. \textbf{79}, 3272 (1998).
\bibitem{Lampl2004} I. Lampl, D. Ferster, T. Poggio, and M. Riesenhuber, \emph{Intracellular measurements of spatial integration and the MAX operation in complex cells of the cat primary visual cortex}, J. Neurophysiol. \textbf{92}, 2704 (2004).
\bibitem{KobatakeTanaka1994} E. Kobatake and K. Tanaka, \emph{Neuronal selectivities to complex object features in the ventral visual pathway of the macaque cerebral cortex}, J. Neurophysiol. \textbf{71}, 856 (1994).
\bibitem{Albright1984} T. D. Albright, \emph{Direction and orientation selectivity of neurons in visual area MT of the macaque}, J. Neurophysiol. \textbf{52}, 1106 (1984).
\bibitem{Goodfellow2015} I. J. Goodfellow, J. Shlens, and C. Szegedy, \emph{Explaining and harnessing adversarial examples}, in \emph{Proc. ICLR} (2015), arXiv:1412.6572.
\bibitem{Elsayed2018} G. F. Elsayed et al., \emph{Adversarial examples that fool both computer vision and time-limited humans}, in \emph{Advances in Neural Information Processing Systems 31} (2018), arXiv:1802.08195.
\bibitem{Thompson1982} P. Thompson, \emph{Perceived rate of movement depends on contrast}, Vision Res. \textbf{22}, 377 (1982).
\bibitem{BarlowHill1963} H. B. Barlow and R. M. Hill, \emph{Evidence for a physiological explanation of the waterfall phenomenon and figural after-effects}, Nature \textbf{200}, 1345 (1963).
\bibitem{EaglemanSejnowski2000} D. M. Eagleman and T. J. Sejnowski, \emph{Motion integration and postdiction in visual awareness}, Science \textbf{287}, 2036 (2000).

\bibitem{Feinstein1955mu} B. Feinstein, B. Lindeg{\aa}rd, E. Nyman, and G. Wohlfart, \emph{Morphologic studies of motor units in normal human muscles}, Acta Anat. \textbf{23}, 127 (1955).
\bibitem{MilnerBrown1973rec} H. S. Milner-Brown, R. B. Stein, and R. Yemm, \emph{The orderly recruitment of human motor units during voluntary isometric contractions}, J. Physiol. \textbf{230}, 359 (1973).
\bibitem{Jones2002sdn} K. E. Jones, A. F. de C. Hamilton, and D. M. Wolpert, \emph{Sources of signal-dependent noise during isometric force production}, J. Neurophysiol. \textbf{88}, 1533 (2002).
\bibitem{GomiOsu1998stiff} H. Gomi and R. Osu, \emph{Task-dependent viscoelasticity of human multijoint arm and its spatial characteristics for interaction with environments}, J. Neurosci. \textbf{18}, 8965 (1998).
\bibitem{Jankowska1972Ia} E. Jankowska and W. J. Roberts, \emph{Synaptic actions of single interneurones mediating reciprocal Ia inhibition of motoneurones}, J. Physiol. \textbf{222}, 623 (1972).
\bibitem{Pruszynski2008rapid} J. A. Pruszynski, I. Kurtzer, and S. H. Scott, \emph{Rapid motor responses are appropriately tuned to the metrics of a visuospatial task}, J. Neurophysiol. \textbf{100}, 224 (2008).
\bibitem{SaundersKnill2003vis} J. A. Saunders and D. C. Knill, \emph{Humans use continuous visual feedback from the hand to control fast reaching movements}, Exp. Brain Res. \textbf{152}, 341 (2003).
\bibitem{FlanaganWing1997grip} J. R. Flanagan and A. M. Wing, \emph{The role of internal models in motion planning and control: evidence from grip force adjustments during movements of hand-held loads}, J. Neurosci. \textbf{17}, 1519 (1997).
\bibitem{Matsuoka1985osc} K. Matsuoka, \emph{Sustained oscillations generated by mutually inhibiting neurons with adaptation}, Biol. Cybern. \textbf{52}, 367 (1985).

\bibitem{Treede1995heat} R.-D. Treede, R. A. Meyer, S. N. Raja, and J. N. Campbell, \emph{Evidence for two different heat transduction mechanisms in nociceptive primary afferents innervating monkey skin}, J. Physiol. \textbf{483}, 747 (1995).
\bibitem{Weidner1999cfib} C. Weidner et al., \emph{Functional attributes discriminating mechano-insensitive and mechano-responsive C nociceptors in human skin}, J. Neurosci. \textbf{19}, 10184 (1999).
\bibitem{CampbellLaMotte1983first} J. N. Campbell and R. H. LaMotte, \emph{Latency to detection of first pain}, Brain Res. \textbf{266}, 203 (1983).
\bibitem{LaMotte1982hyper} R. H. LaMotte, J. G. Thalhammer, H. E. Torebj{\"o}rk, and C. J. Robinson, \emph{Peripheral neural mechanisms of cutaneous hyperalgesia following mild injury by heat}, J. Neurosci. \textbf{2}, 765 (1982).
\bibitem{Mancini2014touch} F. Mancini, T. Nash, G. D. Iannetti, and P. Haggard, \emph{Pain relief by touch: a quantitative approach}, Pain \textbf{155}, 635 (2014).
\bibitem{Fields1983rvm} H. L. Fields, J. Bry, I. Hentall, and G. Zorman, \emph{The activity of neurons in the rostral medulla of the rat during withdrawal from noxious heat}, J. Neurosci. \textbf{3}, 2545 (1983).
\bibitem{Crombez1996disrupt} G. Crombez, C. Eccleston, F. Baeyens, and P. Eelen, \emph{The disruptive nature of pain: an experimental investigation}, Behav. Res. Ther. \textbf{34}, 911 (1996).
\bibitem{Schouenborg1995wr} J. Schouenborg, H.-R. Weng, J. Kalliom{\"a}ki, and H. Holmberg, \emph{A survey of spinal dorsal horn neurones encoding the spatial organization of withdrawal reflexes in the rat}, Exp. Brain Res. \textbf{106}, 19 (1995).

\bibitem{Andersen1992cereb} B. B. Andersen, L. Korbo, and B. Pakkenberg, \emph{A quantitative study of the human cerebellum with unbiased stereological techniques}, J. Comp. Neurol. \textbf{326}, 549 (1992).
\bibitem{ShadmehrMussaIvaldi1994} R. Shadmehr and F. A. Mussa-Ivaldi, \emph{Adaptive representation of dynamics during learning of a motor task}, J. Neurosci. \textbf{14}, 3208 (1994).

\bibitem{JoseCollison1970ihr} A. D. Jose and D. Collison, \emph{The normal range and determinants of the intrinsic heart rate in man}, Cardiovasc. Res. \textbf{4}, 160 (1970).
\bibitem{Logothetis1987gbg} D. E. Logothetis, Y. Kurachi, J. Galper, E. J. Neer, and D. E. Clapham, \emph{The $\beta\gamma$ subunits of GTP-binding proteins activate the muscarinic K$^+$ channel in heart}, Nature \textbf{325}, 321 (1987).
\bibitem{KellerWood1984fb} M. E. Keller-Wood and M. F. Dallman, \emph{Corticosteroid inhibition of ACTH secretion}, Endocr. Rev. \textbf{5}, 1 (1984).
\bibitem{Walker2012glc} J. J. Walker et al., \emph{The origin of glucocorticoid hormone oscillations}, PLoS Biol. \textbf{10}, e1001341 (2012).
\bibitem{Wildt1981gnrh} L. Wildt et al., \emph{Frequency and amplitude of gonadotropin-releasing hormone stimulation and gonadotropin secretion in the rhesus monkey}, Endocrinology \textbf{109}, 376 (1981).
\bibitem{Welsh2010scn} D. K. Welsh, J. S. Takahashi, and S. A. Kay, \emph{Suprachiasmatic nucleus: cell autonomy and network properties}, Annu. Rev. Physiol. \textbf{72}, 551 (2010).
\bibitem{Khalsa2003prc} S. B. S. Khalsa, M. E. Jewett, C. Cajochen, and C. A. Czeisler, \emph{A phase response curve to single bright light pulses in human subjects}, J. Physiol. \textbf{549}, 945 (2003).
\bibitem{Sack1992blind} R. L. Sack, A. J. Lewy, M. L. Blood, L. D. Keith, and H. Nakagawa, \emph{Circadian rhythm abnormalities in totally blind people: incidence and clinical significance}, J. Clin. Endocrinol. Metab. \textbf{75}, 127 (1992).

\bibitem{Henze2000extra} D. A. Henze et al., \emph{Intracellular features predicted by extracellular recordings in the hippocampus in vivo}, J. Neurophysiol. \textbf{84}, 390 (2000).
\bibitem{Histed2009stim} M. H. Histed, V. Bonin, and R. C. Reid, \emph{Direct activation of sparse, distributed populations of cortical neurons by electrical microstimulation}, Neuron \textbf{63}, 508 (2009).
\bibitem{Orsborn2014coad} A. L. Orsborn et al., \emph{Closed-loop decoder adaptation shapes neural plasticity for skillful neuroprosthetic control}, Neuron \textbf{82}, 1380 (2014).

\bibitem{HuVervaekeStorm2002} H. Hu, K. Vervaeke, and J. F. Storm, \emph{Two forms of electrical resonance at theta frequencies, generated by M-current, h-current and persistent Na$^+$ current in rat hippocampal pyramidal cells}, J. Physiol. \textbf{545}, 783 (2002).
\bibitem{Mauro1970} A. Mauro, F. Conti, F. Dodge, and R. Schor, \emph{Subthreshold behavior and phenomenological impedance of the squid giant axon}, J. Gen. Physiol. \textbf{55}, 497 (1970).
\bibitem{Aksay2001} E. Aksay et al., \emph{In vivo intracellular recording and perturbation of persistent activity in a neural integrator}, Nat. Neurosci. \textbf{4}, 184 (2001).
\bibitem{Major2004} G. Major et al., \emph{Plasticity and tuning of the time course of analog persistent firing in a neural integrator}, Proc. Natl. Acad. Sci. USA \textbf{101}, 7745 (2004).
\bibitem{Tateno2004} T. Tateno, A. Harsch, and H. P. C. Robinson, \emph{Threshold firing frequency-current relationships of neurons in rat somatosensory cortex: type 1 and type 2 dynamics}, J. Neurophysiol. \textbf{92}, 2283 (2004).
\bibitem{Marder2012} E. Marder, \emph{Neuromodulation of neuronal circuits: back to the future}, Neuron \textbf{76}, 1 (2012).
\bibitem{McGintyHarper1976} D. J. McGinty and R. M. Harper, \emph{Dorsal raphe neurons: depression of firing during sleep in cats}, Brain Res. \textbf{101}, 569 (1976).

\bibitem{Gentet2000} L. J. Gentet, G. J. Stuart, and J. D. Clements, \emph{Direct measurement of specific membrane capacitance in neurons}, Biophys. J. \textbf{79}, 314 (2000).
\bibitem{Borst1995} J. G. G. Borst, F. Helmchen, and B. Sakmann, \emph{Pre- and postsynaptic whole-cell recordings in the medial nucleus of the trapezoid body of the rat}, J. Physiol. \textbf{489}, 825 (1995).
\bibitem{AmirDevor2003} R. Amir and M. Devor, \emph{Electrical excitability of the soma of sensory neurons is required for spike invasion of the soma, but not for through-conduction}, Biophys. J. \textbf{84}, 2181 (2003).
\bibitem{Chadderton2004} P. Chadderton, T. W. Margrie, and M. H\"ausser, \emph{Integration of quanta in cerebellar granule cells during sensory processing}, Nature \textbf{428}, 856 (2004).
\bibitem{LitwinKumar2017} A. Litwin-Kumar et al., \emph{Optimal degrees of synaptic connectivity}, Neuron \textbf{93}, 1153 (2017).
\bibitem{NapperHarvey1988} R. M. A. Napper and R. J. Harvey, \emph{Number of parallel fiber synapses on an individual Purkinje cell in the cerebellum of the rat}, J. Comp. Neurol. \textbf{274}, 168 (1988).
\bibitem{Megias2001} M. Meg\'{\i}as, Z. Emri, T. F. Freund, and A. I. Guly\'as, \emph{Total number and distribution of inhibitory and excitatory synapses on hippocampal CA1 pyramidal cells}, Neuroscience \textbf{102}, 527 (2001).
\bibitem{Polsky2004} A. Polsky, B. W. Mel, and J. Schiller, \emph{Computational subunits in thin dendrites of pyramidal cells}, Nat. Neurosci. \textbf{7}, 621 (2004).
\bibitem{Brunel2004} N. Brunel et al., \emph{Optimal information storage and the distribution of synaptic weights: perceptron versus Purkinje cell}, Neuron \textbf{43}, 745 (2004).

\bibitem{Vyazovskiy2009} V. V. Vyazovskiy et al., \emph{Cortical firing and sleep homeostasis}, Neuron \textbf{63}, 865 (2009).
\bibitem{Daan1984} S. Daan, D. G. Beersma, and A. A. Borb\'ely, \emph{Timing of human sleep: recovery process gated by a circadian pacemaker}, Am. J. Physiol. \textbf{246}, R161 (1984).
\bibitem{Borbely1981} A. A. Borb\'ely et al., \emph{Sleep deprivation: effect on sleep stages and EEG power density in man}, Electroencephalogr. Clin. Neurophysiol. \textbf{51}, 483 (1981).
\bibitem{SteriadeAmzica1993} M. Steriade, F. Amzica, and A. Nu\~nez, \emph{Cholinergic and noradrenergic modulation of the slow ($\approx0.3$ Hz) oscillation in neocortical cells}, J. Neurophysiol. \textbf{70}, 1385 (1993).
\bibitem{vonKrosigk1993} M. von Krosigk, T. Bal, and D. A. McCormick, \emph{Cellular mechanisms of a synchronized oscillation in the thalamus}, Science \textbf{261}, 361 (1993).
\bibitem{LeeWilson2002} A. K. Lee and M. A. Wilson, \emph{Memory of sequential experience in the hippocampus during slow wave sleep}, Neuron \textbf{36}, 1183 (2002).
\bibitem{Euston2007} D. R. Euston, M. Tatsuno, and B. L. McNaughton, \emph{Fast-forward playback of recent memory sequences in prefrontal cortex during sleep}, Science \textbf{318}, 1147 (2007).
\bibitem{Hobson2009} J. A. Hobson, \emph{REM sleep and dreaming: towards a theory of protoconsciousness}, Nat. Rev. Neurosci. \textbf{10}, 803 (2009).

\bibitem{CohenSegal2011} D. Cohen and M. Segal, \emph{Network bursts in hippocampal microcultures are terminated by exhaustion of vesicle pools}, J. Neurophysiol. \textbf{106}, 2314 (2011).
\bibitem{Tsodyks2000} M. Tsodyks, A. Uziel, and H. Markram, \emph{Synchrony generation in recurrent networks with frequency-dependent synapses}, J. Neurosci. \textbf{20}, RC50 (2000).
\bibitem{Benson1994} D. L. Benson, F. H. Watkins, O. Steward, and G. Banker, \emph{Characterization of GABAergic neurons in hippocampal cell cultures}, J. Neurocytol. \textbf{23}, 279 (1994).

\input{supplement/bib_22}
\end{thebibliography}


\end{document}
